\documentclass[11pt,a4paper]{report}
\usepackage{fontspec,amsmath,amssymb,amsthm,microtype,xcolor}
\setmainfont{lmroman10-regular.otf}[BoldFont=lmroman10-bold.otf,ItalicFont=lmroman10-italic.otf,BoldItalicFont=lmroman10-bolditalic.otf]
\usepackage[most]{tcolorbox}
\usepackage{mathrsfs}
\providecommand{\lt}{<}
\providecommand{\gt}{>}
\setmonofont{DejaVuSansMono.ttf}[BoldFont=DejaVuSansMono-Bold.ttf,ItalicFont=DejaVuSansMono-Oblique.ttf,BoldItalicFont=DejaVuSansMono-BoldOblique.ttf,Scale=0.86]
\newfontfamily\leanfallback{STIXTwoMath-Regular.otf}[Scale=0.86]
\definecolor{leanrule}{gray}{0.62}
\definecolor{leanback}{gray}{0.962}
\definecolor{leantitle}{gray}{0.22}
\definecolor{leanmuted}{gray}{0.38}
\newtcolorbox{leanbox}[1]{enhanced,breakable,colback=leanback,colframe=leanrule,
  boxrule=0.4pt,arc=0.6mm,left=2.4mm,right=2.4mm,top=1.8mm,bottom=1.8mm,
  before skip=0.8em,after skip=1em,fonttitle=\small\bfseries,coltitle=leantitle,
  colbacktitle=leanback,titlerule=0.3pt,title={#1}}
\newlength{\leancharwidth}
\newenvironment{leancode}{\par\smallskip\begingroup\ttfamily\footnotesize\raggedright
  \settowidth{\leancharwidth}{M}\setlength{\parskip}{0pt}\setlength{\parindent}{0pt}
  \frenchspacing}{\par\endgroup\smallskip}
\newcommand{\leanline}[2]{\hangindent=\dimexpr#1\leancharwidth+4\leancharwidth\relax
  \hangafter=1\noindent\hspace*{#1\leancharwidth}#2\par}
\newcommand{\leansource}[1]{\par{\footnotesize\color{leanmuted}#1}\par}
\newcommand{\leanlabel}[1]{\par\smallskip\noindent{\small\bfseries #1}\par}
\usepackage[margin=25mm]{geometry}
\usepackage[unicode,hidelinks]{hyperref}
\hypersetup{pdfauthor={LangevinLean Project},bookmarksdepth=2}
\setlength{\emergencystretch}{3em}
\setlength{\parskip}{0.4em}
\setlength{\parindent}{0pt}
\setcounter{tocdepth}{1}
\makeatletter
\renewcommand*\l@section{\@dottedtocline{1}{1.5em}{3.8em}}
\renewcommand*\@pnumwidth{2.5em}
\renewcommand*\@tocrmarg{3.5em}
\makeatother
\setcounter{secnumdepth}{2}
\allowdisplaybreaks[2]
\raggedbottom
\title{\texorpdfstring{Training Latent Diffusion Models with Interacting Particle Algorithms}{Training Latent Diffusion Models with Interacting Particle Algorithms}\\[0.6em]\large Mathematical proofs}
\author{LangevinLean Project}
\date{Draft report LL-08\\[0.4em]\small Content revision e7fa56ba1732}
\begin{document}
\pagenumbering{roman}
\maketitle
\chapter*{About this document}
\addcontentsline{toc}{chapter}{About this document}
Chapter 1 states the assumptions and the main results. Chapter 2 proves them, following the argument from the hypotheses to the conclusion in order. Chapter 3 proves the supporting lemmas, each before its first use. The proofs are written for readers with an undergraduate background in analysis, linear algebra and probability; definitions are recalled where they are first used, and each inequality is justified.

\section*{Mathematics and its Lean formalization}
Every result is followed by a grey box headed \emph{Lean}. The box reproduces the exact statement of the corresponding Lean declaration, read from the source file and line range named in the box, and then explains how the Lean proof follows the mathematical argument: which hypothesis encodes which assumption, and which formal step carries out which step of the proof. The mathematics can be read without the boxes; the boxes connect each claim to the statement that Lean has checked.

\section*{Checking the proofs}
The results are verified by Lean 4 with the Mathlib library, at the versions pinned in the repository \url{https://github.com/odakyildiz/LangevinLean}. After installing the Lean toolchain manager elan and cloning the repository, run from its root:
\begin{leancode}
\leanline{0}{lake\ exe\ cache\ get}
\leanline{0}{lake\ build\ LatentDiffusion}
\leanline{0}{lake\ env\ lean\ LatentDiffusion/Check.lean}
\end{leancode}
The second command checks every proof in the library. The third prints the axioms on which the main theorems depend; they are only the standard axioms of propositional extensionality, choice and quotient soundness.

Prepared by the LangevinLean Project with AI assistance, these notes form part of an ongoing effort to formalize Langevin algorithms comprehensively. This is draft work and may contain mistakes.

\section*{Literature and scope}
\begin{enumerate}
\item Tim Y. J. Wang, Juan Kuntz, and O. Deniz Akyildiz (2026). \emph{Training Latent Diffusion Models with Interacting Particle Algorithms}. Proceedings of The 29th International Conference on Artificial Intelligence and Statistics, Proceedings of Machine Learning Research 300, 2791–2799. \url{https://proceedings.mlr.press/v300/wang26f.html}.
\end{enumerate}
\tableofcontents
\clearpage
\pagenumbering{arabic}
\chapter{Assumptions and main results}
\section{Assumptions of the paper}
\subsection{\texorpdfstring{A1 — Model differentiability}{A1 — Model differentiability}}



Put \(a=(\theta,\phi)\), \(\tilde\rho_a(z)=\tilde p_a(x,z)\), and \(\tilde\pi_a(z)=\tilde\rho_a(z)/A_a\). For every \(z\), the maps \(a\mapsto\tilde\pi_a(z)\) and \(a\mapsto A_a\) are differentiable. For every \(a\), the spatial map \(z\mapsto\tilde\pi_a(z)\) is \(C^2\), and \(\tilde\rho_a(z)>0\) everywhere.




Here \(A_a\) is the model's normalizing constant. The source contract records these conditions; it does not establish differentiation under integrals or boundary integration by parts.

\begin{leanbox}{Lean: \texttt{normalizedPosterior}, \texttt{ModelRegularity}}
\leansource{\texttt{normalizedPosterior} \quad LatentDiffusion/\allowbreak{}SourceAssumptions.lean, lines 31–32}
\begin{leancode}
\leanline{0}{def\ normalizedPosterior\ (ρ\ :\ P\ →\ Z\ →\ ℝ)\ (A\ :\ P\ →\ ℝ)\ (p\ :\ P)\ (z\ :\ Z)\ :\ ℝ\ :=}
\leanline{2}{ρ\ p\ z\ /\ A\ p}
\end{leancode}
\leansource{\texttt{ModelRegularity} \quad LatentDiffusion/\allowbreak{}SourceAssumptions.lean, lines 36–41}
\begin{leancode}
\leanline{0}{structure\ ModelRegularity\ (ρ\ :\ P\ →\ Z\ →\ ℝ)\ (A\ :\ P\ →\ ℝ)\ :\ Prop\ where}
\leanline{2}{posterior\_parameter\_differentiable\ :}
\leanline{4}{∀\ z,\ Differentiable\ ℝ\ (fun\ p\ ↦\ normalizedPosterior\ ρ\ A\ p\ z)}
\leanline{2}{evidence\_differentiable\ :\ Differentiable\ ℝ\ A}
\leanline{2}{posterior\_spatial\_C2\ :\ ∀\ p,\ ContDiff\ ℝ\ 2\ (normalizedPosterior\ ρ\ A\ p)}
\leanline{2}{tilted\_positive\ :\ ∀\ p\ z,\ 0\ <\ ρ\ p\ z}
\end{leancode}
\end{leanbox}
\subsection{\texorpdfstring{A2 — Classical solutions for every initial law}{A2 — Classical solutions for every initial law}}



For every parameter \(a_0\) and every collection \(\mu_0^{1:M}\in\mathcal P_2(Z)^M\), the source assumes a classical solution of its gradient flow with those initial values. The parameter path is \(C^1\) on \([0,\infty)\). Each law remains in \(\mathcal P_2(Z)\) and has a positive Lebesgue density \(q_t^m(z)\), jointly \(C^{1,2}\) in time and space on \([0,\infty)\times Z\).




\textbf{Source issue.} The universal quantifier includes singular initial laws, while the density and positivity conditions include \(t=0\). This premise is displayed literally and is not claimed to follow from A1, A3 or A4. The source audit records the resulting limitation.

\begin{leanbox}{Lean: \texttt{ClassicalDensityRegularity}, \texttt{ClassicalSolution}, \texttt{SolutionRegularity}}
\leansource{\texttt{ClassicalDensityRegularity} \quad LatentDiffusion/\allowbreak{}SourceAssumptions.lean, lines 49–60}
\begin{leancode}
\leanline{0}{def\ ClassicalDensityRegularity\ \{M\ :\ ℕ\}\ (q\ :\ Fin\ M\ →\ ℝ\ →\ Z\ →\ ℝ)\ :\ Prop\ :=}
\leanline{2}{∃\ (dt\ :\ Fin\ M\ →\ ℝ\ →\ Z\ →\ ℝ)}
\leanline{4}{(dz\ :\ Fin\ M\ →\ ℝ\ →\ Z\ →\ (Z\ →L[ℝ]\ ℝ))}
\leanline{4}{(dzz\ :\ Fin\ M\ →\ ℝ\ →\ Z\ →\ (Z\ →L[ℝ]\ (Z\ →L[ℝ]\ ℝ))),}
\leanline{4}{∀\ m,}
\leanline{6}{ContinuousOn\ (fun\ tz\ :\ ℝ\ ×\ Z\ ↦\ q\ m\ tz.1\ tz.2)\ timeSpaceDomain\ ∧}
\leanline{6}{ContinuousOn\ (fun\ tz\ :\ ℝ\ ×\ Z\ ↦\ dt\ m\ tz.1\ tz.2)\ timeSpaceDomain\ ∧}
\leanline{6}{ContinuousOn\ (fun\ tz\ :\ ℝ\ ×\ Z\ ↦\ dz\ m\ tz.1\ tz.2)\ timeSpaceDomain\ ∧}
\leanline{6}{ContinuousOn\ (fun\ tz\ :\ ℝ\ ×\ Z\ ↦\ dzz\ m\ tz.1\ tz.2)\ timeSpaceDomain\ ∧}
\leanline{6}{(∀\ t,\ 0\ ≤\ t\ →\ ∀\ z,\ HasDerivWithinAt\ (fun\ s\ ↦\ q\ m\ s\ z)\ (dt\ m\ t\ z)\ (Ici\ 0)\ t)\ ∧}
\leanline{6}{(∀\ t,\ 0\ ≤\ t\ →\ ∀\ z,\ HasFDerivAt\ (q\ m\ t)\ (dz\ m\ t\ z)\ z)\ ∧}
\leanline{6}{(∀\ t,\ 0\ ≤\ t\ →\ ∀\ z,\ HasFDerivAt\ (dz\ m\ t)\ (dzz\ m\ t\ z)\ z)}
\end{leancode}
\leansource{\texttt{ClassicalSolution} \quad LatentDiffusion/\allowbreak{}SourceAssumptions.lean, lines 67–81}
\begin{leancode}
\leanline{0}{structure\ ClassicalSolution\ \{M\ :\ ℕ\}\ (ν\ :\ Measure\ Z)}
\leanline{4}{(flowEquation\ :\ (ℝ\ →\ P)\ →\ (Fin\ M\ →\ ℝ\ →\ ProbabilityMeasure\ Z)\ →\ Prop)}
\leanline{4}{(p₀\ :\ P)\ (μ₀\ :\ Fin\ M\ →\ ProbabilityMeasure\ Z)\ where}
\leanline{2}{parameter\ :\ ℝ\ →\ P}
\leanline{2}{law\ :\ Fin\ M\ →\ ℝ\ →\ ProbabilityMeasure\ Z}
\leanline{2}{density\ :\ Fin\ M\ →\ ℝ\ →\ Z\ →\ ℝ}
\leanline{2}{initial\_parameter\ :\ parameter\ 0\ =\ p₀}
\leanline{2}{initial\_law\ :\ ∀\ m,\ law\ m\ 0\ =\ μ₀\ m}
\leanline{2}{parameter\_C1\ :\ ContDiffOn\ ℝ\ 1\ parameter\ (Ici\ 0)}
\leanline{2}{density\_C12\ :\ ClassicalDensityRegularity\ density}
\leanline{2}{density\_positive\ :\ ∀\ m\ t,\ 0\ ≤\ t\ →\ ∀\ z,\ 0\ <\ density\ m\ t\ z}
\leanline{2}{density\_realizes\ :\ ∀\ m\ t,\ 0\ ≤\ t\ →}
\leanline{4}{(law\ m\ t\ :\ Measure\ Z)\ =\ ν.withDensity\ (fun\ z\ ↦\ ENNReal.ofReal\ (density\ m\ t\ z))}
\leanline{2}{finite\_second\_moment\ :\ ∀\ m\ t,\ 0\ ≤\ t\ →\ MemLp\ (fun\ z\ :\ Z\ ↦\ z)\ 2\ (law\ m\ t\ :\ Measure\ Z)}
\leanline{2}{solves\ :\ flowEquation\ parameter\ law}
\end{leancode}
\leansource{\texttt{SolutionRegularity} \quad LatentDiffusion/\allowbreak{}SourceAssumptions.lean, lines 86–90}
\begin{leancode}
\leanline{0}{def\ SolutionRegularity\ \{M\ :\ ℕ\}\ (ν\ :\ Measure\ Z)}
\leanline{4}{(flowEquation\ :\ (ℝ\ →\ P)\ →\ (Fin\ M\ →\ ℝ\ →\ ProbabilityMeasure\ Z)\ →\ Prop)\ :\ Prop\ :=}
\leanline{2}{∀\ (p₀\ :\ P)\ (μ₀\ :\ Fin\ M\ →\ ProbabilityMeasure\ Z),}
\leanline{4}{(∀\ m,\ MemLp\ (fun\ z\ :\ Z\ ↦\ z)\ 2\ (μ₀\ m\ :\ Measure\ Z))\ →}
\leanline{6}{Nonempty\ (ClassicalSolution\ ν\ flowEquation\ p₀\ μ₀)}
\end{leancode}
\leanlabel{How the Lean proof follows the argument}
The Lean contract uses a supplied reference measure \(\nu\), instantiated by Lebesgue measure, and the actual flow equation as an explicit predicate. Its \(C^{1,2}\) contract requires continuous first time derivatives and first and second spatial derivatives, including right time derivatives at zero.
\end{leanbox}
\subsection{\texorpdfstring{A3 — Joint strong log-concavity}{A3 — Joint strong log-concavity}}



For every datum \(m\), the tilted joint density is \(\lambda\)-strongly log-concave in the full pair \((a,z)\), with one \(\lambda>0\). Equivalently, for \(U_m(a,z)=-\log\tilde p_a(x^m,z)\),



\[U_m(sx+(1-s)y)\le sU_m(x)+(1-s)U_m(y)-\tfrac\lambda2s(1-s)\|x-y\|^2,\quad 0\le s\le1.\]



The particle-average geometry used below is normalized in the latent coordinates. The same constant does not hold in an unweighted replicated Euclidean norm without accounting for that normalization.

\begin{leanbox}{Lean: \texttt{JointStrongLogConcavity}}
\leansource{\texttt{JointStrongLogConcavity} \quad LatentDiffusion/\allowbreak{}SourceAssumptions.lean, lines 98–100}
\begin{leancode}
\leanline{0}{structure\ JointStrongLogConcavity\ (U\ :\ I\ →\ H\ →\ ℝ)\ (m\ :\ ℝ)\ :\ Prop\ where}
\leanline{2}{positive\ :\ 0\ <\ m}
\leanline{2}{strong\ :\ ∀\ i,\ StrongConvexOn\ univ\ m\ (U\ i)}
\end{leancode}
\end{leanbox}
\subsection{\texorpdfstring{A4 — Joint Lipschitz gradient}{A4 — Joint Lipschitz gradient}}



Every \(U_m\) is differentiable and, for one \(L>0\),



\[\|\nabla U_m(a,z)-\nabla U_m(a',z')\|\le L\|(a,z)-(a',z')\|.\]



Negating the log tilted density leaves the Lipschitz constant unchanged. The source's discretization theorem additionally assumes Lipschitz gradients of the tilting penalty \(\mathcal R\); the numerical results below instead expose the complete force and its approximation error directly.

\begin{leanbox}{Lean: \texttt{LipschitzGradient}}
\leansource{\texttt{LipschitzGradient} \quad LatentDiffusion/\allowbreak{}SourceAssumptions.lean, lines 104–107}
\begin{leancode}
\leanline{0}{structure\ LipschitzGradient\ (U\ :\ I\ →\ H\ →\ ℝ)\ (L\ :\ ℝ≥0)\ :\ Prop\ where}
\leanline{2}{positive\ :\ 0\ <\ L}
\leanline{2}{differentiable\ :\ ∀\ i,\ Differentiable\ ℝ\ (U\ i)}
\leanline{2}{lipschitz\ :\ ∀\ i,\ LipschitzWith\ L\ (gradient\ (U\ i))}
\end{leancode}
\end{leanbox}
\section{Extra assumptions}
\subsection{\texorpdfstring{R1 — Positive normalized densities and the finite-KL domain}{R1 — Positive normalized densities and the finite-KL domain}}


Let \(M\in\mathbb N\), \(M>0\), let \(\nu\) be the reference measure on the latent space, and let \(a\) range over the parameter domain. For every datum \(m\), assume that \(\rho_a^m:Z\to(0,\infty)\) is measurable, \(\rho_a^m\in L^1(\nu)\), and
        \[0<Z_m(a):=\int\rho_a^m\,d\nu<\infty,\qquad
          \pi_a^m(dz):=Z_m(a)^{-1}\rho_a^m(z)\,\nu(dz).\]
        For a trial family \(q=(q^m)_{m=1}^M\), the domain condition is
        \(\mathrm{KL}(q^m\|\pi_a^m)<\infty\) for every \(m\). If minimization is restricted to a law domain \(\mathcal D\), assume \((\pi_a^m)_m\in\mathcal D\) for every \(a\). For the paper's domain take \(\mathcal D=\mathcal P_2(Z)^M\), so \(\int\|z\|^2\,d\pi_a^m<\infty\) is required. The real Radon–Nikodym density \(f_m=dq^m/d\pi_a^m\), its normalization and \(f_m\log f_m\in L^1(\pi_a^m)\) follow from finite KL. The literal density \(q_m=f_m\rho_a^m/Z_m(a)\) and the energy identity are then derived. This variational domain does not impose smoothness on arbitrary trial laws. The classical-flow statements additionally require the positive differentiable initial density specified in R2–R3; a Dirac initial law does not satisfy it.

\begin{leanbox}{Lean: \texttt{DensityModel}, \texttt{DensityModel.\allowbreak{}literal\_\allowbreak{}integral}, \texttt{DensityModel.\allowbreak{}literal\_\allowbreak{}minimizers\_\allowbreak{}on\_\allowbreak{}iff}}
\leansource{\texttt{DensityModel} \quad LatentDiffusion/\allowbreak{}RegularVariational.lean, lines 24–29}
\begin{leancode}
\leanline{0}{structure\ DensityModel\ (A\ Z\ :\ Type*)\ [MeasurableSpace\ Z]\ (M\ :\ ℕ)\ (ν\ :\ Measure\ Z)\ where}
\leanline{2}{density\ :\ A\ →\ Fin\ M\ →\ Z\ →\ ℝ}
\leanline{2}{measurable\ :\ ∀\ a\ i,\ Measurable\ (density\ a\ i)}
\leanline{2}{positive\ :\ ∀\ a\ i\ z,\ 0\ <\ density\ a\ i\ z}
\leanline{2}{integrable\ :\ ∀\ a\ i,\ Integrable\ (density\ a\ i)\ ν}
\leanline{2}{mass\_positive\ :\ ∀\ a\ i,\ 0\ <\ ∫\ z,\ density\ a\ i\ z\ ∂ν}
\end{leancode}
\leansource{\texttt{DensityModel.\allowbreak{}literal\_\allowbreak{}integral} \quad LatentDiffusion/\allowbreak{}RegularVariational.lean, lines 87–117}
\begin{leancode}
\leanline{0}{theorem\ DensityModel.literal\_integral\ (D\ :\ DensityModel\ A\ Z\ M\ ν)\ (a\ :\ A)}
\leanline{4}{(q\ :\ Fin\ M\ →\ ProbabilityMeasure\ Z)\ (hadm\ :\ D.family.Admissible\ a\ q)\ (i\ :\ Fin\ M)\ :}
\leanline{4}{Density.densityMeasure\ ν\ (D.trialDensity\ a\ q\ i)\ =\ (q\ i\ :\ Measure\ Z)\ ∧}
\leanline{4}{Integrable\ (fun\ z\ ↦\ D.trialDensity\ a\ q\ i\ z\ *}
\leanline{6}{Real.log\ (D.trialDensity\ a\ q\ i\ z\ /\ D.density\ a\ i\ z))\ ν\ ∧}
\leanline{4}{(∫\ z,\ D.trialDensity\ a\ q\ i\ z\ *\ Real.log\ (D.trialDensity\ a\ q\ i\ z\ /\ D.density\ a\ i\ z)\ ∂ν)\ =}
\leanline{6}{(InformationTheory.klDiv\ (q\ i\ :\ Measure\ Z)\ (D.posterior\ a\ i\ :\ Measure\ Z)).toReal\ -}
\leanline{8}{Real.log\ (D.mass\ a\ i)}
\end{leancode}
\leansource{\texttt{DensityModel.\allowbreak{}literal\_\allowbreak{}minimizers\_\allowbreak{}on\_\allowbreak{}iff} \quad LatentDiffusion/\allowbreak{}RegularVariational.lean, lines 148–176}
\begin{leancode}
\leanline{0}{theorem\ DensityModel.literal\_minimizers\_on\_iff\ (D\ :\ DensityModel\ A\ Z\ M\ ν)\ (hM\ :\ 0\ <\ M)}
\leanline{4}{(domain\ :\ (Fin\ M\ →\ ProbabilityMeasure\ Z)\ →\ Prop)}
\leanline{4}{(hposterior\ :\ ∀\ a,\ domain\ (D.posterior\ a))\ (a\ :\ A)\ (q\ :\ Fin\ M\ →\ ProbabilityMeasure\ Z)\ :}
\leanline{4}{(D.family.Admissible\ a\ q\ ∧\ domain\ q\ ∧\ ∀\ b\ r,\ D.family.Admissible\ b\ r\ →\ domain\ r\ →}
\leanline{6}{D.literalEnergy\ a\ q\ ≤\ D.literalEnergy\ b\ r)\ ↔}
\leanline{4}{(∀\ b,\ D.family.logEvidence\ b\ ≤\ D.family.logEvidence\ a)\ ∧\ q\ =\ D.posterior\ a}
\end{leancode}
\end{leanbox}
\subsection{\texorpdfstring{R2 — Quantified domination for differentiation under the integral}{R2 — Quantified domination for differentiation under the integral}}


For each differentiated normalizer at a parameter \(a_0\), assume a single neighbourhood \(V\ni a_0\) and \(g\in L^1(\nu)\) such that, outside one \(\nu\)-null set, for every \(b\in V\), \(U_m(\cdot,z)\) is differentiable at \(b\) and
        \[e^{-U_m(b,z)}\|\nabla_aU_m(b,z)\|\le g(z).\]
        Require \(e^{-U_m(a_0,\cdot)}\in L^1(\nu)\), almost-everywhere strong measurability of \(z\mapsto e^{-U_m(b,z)}\) for all \(b\) in a neighbourhood of \(a_0\), and of \(z\mapsto e^{-U_m(a_0,z)}\nabla_aU_m(a_0,z)\). These are the exact domination inputs for the normalizing integral.




For an energy derivative at \(t_0\), write \(U_t(z)=U_m(a(t),z)\), \(\dot U_t(z)=D_aU_m(a(t),z)[\dot a(t)]\), and
        \[H_t(z)=q_t(z)(\log q_t(z)+U_t(z)),\qquad
          D_t(z)=\dot q_t(z)(\log q_t(z)+U_t(z)+1)+q_t(z)\dot U_t(z).\]
        Assume one neighbourhood \(J\ni t_0\) and \(g\in L^1(\nu)\) such that, outside one null set and for every \(t\in J\), \(q_t>0\), \(\partial_tq_t=\dot q_t\), and \(|D_t|\le g\). Require \(H_{t_0}\in L^1(\nu)\), almost-everywhere strong measurability of \(D_{t_0}\), and of \(H_t\) for all \(t\) in a neighbourhood of \(t_0\). For a direct potential derivative require, outside one null set, \(\partial_tU_t(z)=\dot U_t(z)\) for every \(t\in J\); the chain-rule constructor instead uses differentiability of \(a\) at every such time and of \(U_m(\cdot,z)\) at the corresponding parameter. In the flow and directional statements also require \(q_t\nabla_aU_m(a(t),\cdot)\in L^1(\nu)\). These quantified conditions, rather than an assumed derivative of an integral, define R2.

\begin{leanbox}{Lean: \texttt{LocalEnergyDomination}, \texttt{LocalEnergyDomination.\allowbreak{}along\_\allowbreak{}parameter\_\allowbreak{}path}, \texttt{LocalNormalizerRegularity} and 1 more}
\leansource{\texttt{LocalEnergyDomination} \quad LatentDiffusion/\allowbreak{}RegularFlowCalculus.lean, lines 45–60}
\begin{leancode}
\leanline{0}{structure\ LocalEnergyDomination\ \{X\ :\ Type*\}\ [MeasurableSpace\ X]}
\leanline{4}{(ν\ :\ Measure\ X)\ (q\ qdot\ U\ Udot\ :\ ℝ\ →\ X\ →\ ℝ)\ (t\ :\ ℝ)\ where}
\leanline{2}{neighborhood\ :\ Set\ ℝ}
\leanline{2}{neighborhood\_mem\ :\ neighborhood\ ∈\ {\leanfallback 𝓝}\ t}
\leanline{2}{bound\ :\ X\ →\ ℝ}
\leanline{2}{measurable\ :\ ∀ᶠ\ u\ in\ {\leanfallback 𝓝}\ t,}
\leanline{4}{AEStronglyMeasurable\ (fun\ x\ ↦\ q\ u\ x\ *\ (Real.log\ (q\ u\ x)\ +\ U\ u\ x))\ ν}
\leanline{2}{integrable\ :\ Integrable\ (fun\ x\ ↦\ q\ t\ x\ *\ (Real.log\ (q\ t\ x)\ +\ U\ t\ x))\ ν}
\leanline{2}{derivative\_measurable\ :\ AEStronglyMeasurable}
\leanline{4}{(fun\ x\ ↦\ qdot\ t\ x\ *\ (Real.log\ (q\ t\ x)\ +\ U\ t\ x\ +\ 1)\ +\ q\ t\ x\ *\ Udot\ t\ x)\ ν}
\leanline{2}{derivative\_bound\ :\ ∀ᵐ\ x\ ∂ν,\ ∀\ u\ ∈\ neighborhood,}
\leanline{4}{‖qdot\ u\ x\ *\ (Real.log\ (q\ u\ x)\ +\ U\ u\ x\ +\ 1)\ +\ q\ u\ x\ *\ Udot\ u\ x‖\ ≤\ bound\ x}
\leanline{2}{bound\_integrable\ :\ Integrable\ bound\ ν}
\leanline{2}{positive\ :\ ∀ᵐ\ x\ ∂ν,\ ∀\ u\ ∈\ neighborhood,\ 0\ <\ q\ u\ x}
\leanline{2}{density\_derivative\ :\ ∀ᵐ\ x\ ∂ν,\ ∀\ u\ ∈\ neighborhood,}
\leanline{4}{HasDerivAt\ (fun\ v\ ↦\ q\ v\ x)\ (qdot\ u\ x)\ u}
\end{leancode}
\leansource{\texttt{LocalEnergyDomination.\allowbreak{}along\_\allowbreak{}parameter\_\allowbreak{}path} \quad LatentDiffusion/\allowbreak{}RegularFlowCalculus.lean, lines 480–491}
\begin{leancode}
\leanline{0}{def\ LocalEnergyDomination.along\_parameter\_path}
\leanline{4}{(ν\ :\ Measure\ X)\ (q\ qdot\ :\ ℝ\ →\ X\ →\ ℝ)\ (U\ :\ A\ →\ X\ →\ ℝ)}
\leanline{4}{(p\ pdot\ :\ ℝ\ →\ A)\ (t\ :\ ℝ)}
\leanline{4}{(data\ :\ LocalEnergyDomination\ ν\ q\ qdot\ (fun\ u\ ↦\ U\ (p\ u))}
\leanline{6}{(fun\ u\ x\ ↦\ fderiv\ ℝ\ (fun\ a\ ↦\ U\ a\ x)\ (p\ u)\ (pdot\ u))\ t)}
\leanline{4}{(hp\ :\ ∀\ u\ ∈\ data.neighborhood,\ HasDerivAt\ p\ (pdot\ u)\ u)}
\leanline{4}{(hU\ :\ ∀\ x\ u,\ u\ ∈\ data.neighborhood\ →\ DifferentiableAt\ ℝ\ (fun\ a\ ↦\ U\ a\ x)\ (p\ u))\ :}
\leanline{4}{LocalEnergyRegularity\ ν\ q\ qdot\ (fun\ u\ ↦\ U\ (p\ u))}
\leanline{6}{(fun\ u\ x\ ↦\ fderiv\ ℝ\ (fun\ a\ ↦\ U\ a\ x)\ (p\ u)\ (pdot\ u))\ t\ where}
\leanline{2}{toLocalEnergyDomination\ :=\ data}
\leanline{2}{potential\_derivative\ :=\ ae\_of\_all\ \_\ fun\ x\ u\ hu\ ↦}
\leanline{4}{(hU\ x\ u\ hu).hasFDerivAt.comp\_hasDerivAt\ u\ (hp\ u\ hu)}
\end{leancode}
\leansource{\texttt{LocalNormalizerRegularity} \quad LatentDiffusion/\allowbreak{}RegularEvidence.lean, lines 43–55}
\begin{leancode}
\leanline{0}{structure\ LocalNormalizerRegularity\ (ν\ :\ Measure\ Z)\ (U\ :\ A\ →\ Z\ →\ ℝ)\ (a\ :\ A)\ where}
\leanline{2}{neighborhood\ :\ Set\ A}
\leanline{2}{neighborhood\_mem\ :\ neighborhood\ ∈\ {\leanfallback 𝓝}\ a}
\leanline{2}{bound\ :\ Z\ →\ ℝ}
\leanline{2}{measurable\ :\ ∀ᶠ\ b\ in\ {\leanfallback 𝓝}\ a,\ AEStronglyMeasurable\ (fun\ z\ ↦\ Real.exp\ (-U\ b\ z))\ ν}
\leanline{2}{density\_integrable\ :\ Integrable\ (fun\ z\ ↦\ Real.exp\ (-U\ a\ z))\ ν}
\leanline{2}{gradient\_measurable\ :\ AEStronglyMeasurable}
\leanline{4}{(fun\ z\ ↦\ Real.exp\ (-U\ a\ z)\ •\ gradient\ (fun\ b\ ↦\ U\ b\ z)\ a)\ ν}
\leanline{2}{derivative\_bound\ :\ ∀ᵐ\ z\ ∂ν,\ ∀\ b\ ∈\ neighborhood,}
\leanline{4}{Real.exp\ (-U\ b\ z)\ *\ ‖gradient\ (fun\ c\ ↦\ U\ c\ z)\ b‖\ ≤\ bound\ z}
\leanline{2}{bound\_integrable\ :\ Integrable\ bound\ ν}
\leanline{2}{differentiable\ :\ ∀ᵐ\ z\ ∂ν,\ ∀\ b\ ∈\ neighborhood,}
\leanline{4}{DifferentiableAt\ ℝ\ (fun\ c\ ↦\ U\ c\ z)\ b}
\end{leancode}
\leansource{\texttt{sum\_\allowbreak{}posterior\_\allowbreak{}parameter\_\allowbreak{}gradient\_\allowbreak{}eq\_\allowbreak{}zero} \quad LatentDiffusion/\allowbreak{}RegularEvidence.lean, lines 185–193}
\begin{leancode}
\leanline{0}{theorem\ sum\_posterior\_parameter\_gradient\_eq\_zero\ \{M\ :\ ℕ\}\ (hM\ :\ 0\ <\ M)}
\leanline{4}{(ν\ :\ Measure\ Z)\ (U\ :\ Fin\ M\ →\ A\ →\ Z\ →\ ℝ)\ (a\ :\ A)}
\leanline{4}{(data\ :\ ∀\ i,\ LocalNormalizerRegularity\ ν\ (U\ i)\ a)}
\leanline{4}{(hmeas\ :\ ∀\ i,\ Measurable\ (fun\ z\ ↦\ Real.exp\ (-U\ i\ a\ z)))}
\leanline{4}{(hZ\ :\ ∀\ i,\ 0\ <\ normalizer\ ν\ (U\ i)\ a)}
\leanline{4}{(hmax\ :\ IsLocalMax\ (logEvidence\ ν\ U)\ a)\ :}
\leanline{4}{∑\ i,\ (∫\ z,\ gradient\ (fun\ b\ ↦\ U\ i\ b\ z)\ a\ ∂posterior\ ν\ (U\ i)\ a)\ =\ 0}
\end{leancode}
\end{leanbox}
\subsection{\texorpdfstring{R3 — Classical equations, differentiability and spatial integrability}{R3 — Classical equations, differentiability and spatial integrability}}


On the Euclidean latent space \(Z=\mathbb R^{d_z}\), put \(\psi_t^m=\log q_t^m+U_m(a(t),\cdot)\) and \(J_t^m=q_t^m\nabla_z\psi_t^m\). The density-dynamics contract requires, for every \(t\in\mathbb R\) and \(m\), \(q_t^m(z)>0\) for every \(z\), \(q_t^m,U_m(a(t),\cdot)\in C^1(Z)\), differentiability of \(a\), differentiability of \(U_m(\cdot,z)\) at every parameter, the R2 energy domination, and \(q_t^m\nabla_aU_m(a(t),\cdot)\in L^1(dz)\). For \(t\ge0\), impose
        \[\dot a(t)=-\frac1M\sum_m\int q_t^m(z)\nabla_aU_m(a(t),z)\,dz,
          \qquad \partial_tq_t^m=\nabla_z\!\cdot J_t^m\quad\text{a.e.}\]
        For every spatial coordinate \(j\), require \(J_{t,j}^m\) differentiable and all three functions
        \[J_{t,j}^m(\psi_t^m+1),\qquad
          (\partial_jJ_{t,j}^m)(\psi_t^m+1),\qquad
          J_{t,j}^m\partial_j\psi_t^m\]
        in \(L^1(dz)\). This is the precise no-boundary-term hypothesis. For a first variation along \(a_s=(1-s)a+sb\), \(T_s=(1-s)\mathrm{id}+sT\), impose the same R2 domination at \(s=0\), the equation \(\dot q_0=-\nabla\cdot(q_0(T-\mathrm{id}))\), differentiability of every component of \(q_0(T-\mathrm{id})\), and these three integrability conditions with flux \(q_0(T-\mathrm{id})\) and test \(\log q_0+U_m(a,\cdot)+1\). Require also \(U_m(a_s,\cdot)\in L^1(q_s^m)\) for every \(s\in[0,1]\). The gradient-pairing statement imposes the corresponding continuity equation and three integrable products for each velocity under consideration. Fisher integrability and de Bruijn are conclusions of these equations and hypotheses.

\begin{leanbox}{Lean: \texttt{NoBoundaryFlux}, \texttt{integral\_\allowbreak{}divergence\_\allowbreak{}mul}, \texttt{normalized\_\allowbreak{}product\_\allowbreak{}gradient\_\allowbreak{}representation} and 3 more}
\leansource{\texttt{NoBoundaryFlux} \quad LatentDiffusion/\allowbreak{}RegularFlowCalculus.lean, lines 199–207}
\begin{leancode}
\leanline{0}{structure\ NoBoundaryFlux\ (V\ :\ E\ →\ E)\ (ψ\ :\ E\ →\ ℝ)\ :\ Prop\ where}
\leanline{2}{flux\_differentiable\ :\ ∀\ i,\ Differentiable\ ℝ\ (fluxComponent\ V\ i)}
\leanline{2}{derivative\_flux\_integrable\ :\ ∀\ i,\ Integrable}
\leanline{4}{(fun\ x\ ↦\ fderiv\ ℝ\ (fluxComponent\ V\ i)\ x\ (stdOrthonormalBasis\ ℝ\ E\ i)\ *\ ψ\ x)}
\leanline{4}{volume}
\leanline{2}{derivative\_test\_integrable\ :\ ∀\ i,\ Integrable}
\leanline{4}{(fun\ x\ ↦\ fluxComponent\ V\ i\ x\ *\ fderiv\ ℝ\ ψ\ x\ (stdOrthonormalBasis\ ℝ\ E\ i))}
\leanline{4}{volume}
\leanline{2}{product\_integrable\ :\ ∀\ i,\ Integrable\ (fun\ x\ ↦\ fluxComponent\ V\ i\ x\ *\ ψ\ x)\ volume}
\end{leancode}
\leansource{\texttt{integral\_\allowbreak{}divergence\_\allowbreak{}mul} \quad LatentDiffusion/\allowbreak{}RegularFlowCalculus.lean, lines 251–273}
\begin{leancode}
\leanline{0}{theorem\ integral\_divergence\_mul}
\leanline{4}{(V\ :\ E\ →\ E)\ (ψ\ :\ E\ →\ ℝ)\ (hψ\ :\ Differentiable\ ℝ\ ψ)}
\leanline{4}{(htail\ :\ NoBoundaryFlux\ V\ ψ)\ :}
\leanline{4}{(∫\ x,\ fluxDivergence\ V\ x\ *\ ψ\ x\ ∂volume)\ =}
\leanline{6}{-(∫\ x,\ inner\ ℝ\ (V\ x)\ (gradient\ ψ\ x)\ ∂volume)}
\end{leancode}
\leansource{\texttt{normalized\_\allowbreak{}product\_\allowbreak{}gradient\_\allowbreak{}representation} \quad LatentDiffusion/\allowbreak{}RegularFlowCalculus.lean, lines 548–590}
\begin{leancode}
\leanline{0}{theorem\ normalized\_product\_gradient\_representation}
\leanline{4}{(q\ qdot\ :\ Fin\ M\ →\ ℝ\ →\ E\ →\ ℝ)\ (U\ :\ Fin\ M\ →\ A\ →\ E\ →\ ℝ)}
\leanline{4}{(p\ pdot\ :\ ℝ\ →\ A)\ (t\ :\ ℝ)\ (v\ :\ Fin\ M\ →\ E\ →\ E)}
\leanline{4}{(data\ :\ ∀\ i,\ LocalEnergyDomination\ volume\ (q\ i)\ (qdot\ i)\ (fun\ u\ ↦\ U\ i\ (p\ u))}
\leanline{6}{(fun\ u\ x\ ↦\ fderiv\ ℝ\ (fun\ a\ ↦\ U\ i\ a\ x)\ (p\ u)\ (pdot\ u))\ t)}
\leanline{4}{(hp\ :\ ∀\ u,\ HasDerivAt\ p\ (pdot\ u)\ u)}
\leanline{4}{(hparam\ :\ ∀\ i\ x\ a,\ DifferentiableAt\ ℝ\ (fun\ b\ ↦\ U\ i\ b\ x)\ a)}
\leanline{4}{(hpos\ :\ ∀\ i\ x,\ 0\ <\ q\ i\ t\ x)\ (hq\ :\ ∀\ i,\ ContDiff\ ℝ\ 1\ (q\ i\ t))}
\leanline{4}{(hU\ :\ ∀\ i,\ ContDiff\ ℝ\ 1\ (U\ i\ (p\ t)))}
\leanline{4}{(hgrad\ :\ ∀\ i,\ Integrable}
\leanline{6}{(fun\ x\ ↦\ q\ i\ t\ x\ •\ gradient\ (fun\ b\ ↦\ U\ i\ b\ x)\ (p\ t))\ volume)}
\leanline{4}{(hcontinuity\ :\ ∀\ i,\ ∀ᵐ\ x\ ∂volume,}
\leanline{6}{qdot\ i\ t\ x\ =\ -fluxDivergence\ (fun\ y\ ↦\ q\ i\ t\ y\ •\ v\ i\ y)\ x)}
\leanline{4}{(htail\ :\ ∀\ i,\ NoBoundaryFlux\ (fun\ y\ ↦\ q\ i\ t\ y\ •\ v\ i\ y)}
\leanline{6}{(fun\ x\ ↦\ chemicalPotential\ (q\ i\ t)\ (U\ i\ (p\ t))\ x\ +\ 1))\ :}
\leanline{4}{HasDerivAt\ (fun\ u\ ↦\ normalizedDensityEnergy\ volume\ (fun\ i\ ↦\ q\ i\ u)\ U\ (p\ u))}
\leanline{6}{(normalizedProductGradientPairing\ volume\ (fun\ i\ ↦\ q\ i\ t)\ U\ (p\ t)\ (pdot\ t)\ v)\ t}
\end{leancode}
\leansource{\texttt{normalized\_\allowbreak{}product\_\allowbreak{}deBruijn} \quad LatentDiffusion/\allowbreak{}RegularFlowCalculus.lean, lines 613–663}
\begin{leancode}
\leanline{0}{theorem\ normalized\_product\_deBruijn}
\leanline{4}{(q\ qdot\ :\ Fin\ M\ →\ ℝ\ →\ E\ →\ ℝ)\ (U\ :\ Fin\ M\ →\ A\ →\ E\ →\ ℝ)}
\leanline{4}{(p\ pdot\ :\ ℝ\ →\ A)\ (t\ :\ ℝ)}
\leanline{4}{(data\ :\ ∀\ i,\ LocalEnergyDomination\ volume\ (q\ i)\ (qdot\ i)\ (fun\ u\ ↦\ U\ i\ (p\ u))}
\leanline{6}{(fun\ u\ x\ ↦\ fderiv\ ℝ\ (fun\ a\ ↦\ U\ i\ a\ x)\ (p\ u)\ (pdot\ u))\ t)}
\leanline{4}{(hp\ :\ ∀\ u,\ HasDerivAt\ p\ (pdot\ u)\ u)}
\leanline{4}{(hparam\ :\ ∀\ i\ x\ a,\ DifferentiableAt\ ℝ\ (fun\ b\ ↦\ U\ i\ b\ x)\ a)}
\leanline{4}{(hpos\ :\ ∀\ i\ x,\ 0\ <\ q\ i\ t\ x)\ (hq\ :\ ∀\ i,\ ContDiff\ ℝ\ 1\ (q\ i\ t))}
\leanline{4}{(hU\ :\ ∀\ i,\ ContDiff\ ℝ\ 1\ (U\ i\ (p\ t)))}
\leanline{4}{(hgrad\ :\ ∀\ i,\ Integrable}
\leanline{6}{(fun\ x\ ↦\ q\ i\ t\ x\ •\ gradient\ (fun\ b\ ↦\ U\ i\ b\ x)\ (p\ t))\ volume)}
\leanline{4}{(hode\ :\ pdot\ t\ =\ -averagedParameterForce\ volume\ (fun\ i\ ↦\ q\ i\ t)\ U\ (p\ t))}
\leanline{4}{(hpde\ :\ ∀\ i,\ ∀ᵐ\ x\ ∂volume,\ qdot\ i\ t\ x\ =}
\leanline{6}{fluxDivergence\ (fokkerPlanckFlux\ (q\ i\ t)\ (U\ i\ (p\ t)))\ x)}
\leanline{4}{(htail\ :\ ∀\ i,\ NoBoundaryFlux\ (fokkerPlanckFlux\ (q\ i\ t)\ (U\ i\ (p\ t)))}
\leanline{6}{(fun\ x\ ↦\ chemicalPotential\ (q\ i\ t)\ (U\ i\ (p\ t))\ x\ +\ 1))\ :}
\leanline{4}{(∀\ i,\ Integrable}
\leanline{6}{(fun\ x\ ↦\ q\ i\ t\ x\ *\ ‖gradient\ (chemicalPotential\ (q\ i\ t)\ (U\ i\ (p\ t)))\ x‖\ \textasciicircum{}\ 2)}
\leanline{6}{volume)\ ∧}
\leanline{4}{HasDerivAt\ (fun\ u\ ↦\ normalizedDensityEnergy\ volume\ (fun\ i\ ↦\ q\ i\ u)\ U\ (p\ u))}
\leanline{6}{(-normalizedSlopeSquared\ volume\ (fun\ i\ ↦\ q\ i\ t)\ U\ (p\ t))\ t}
\end{leancode}
\leansource{\texttt{RegularDensityDynamics} \quad LatentDiffusion/\allowbreak{}RegularFlowTalagrand.lean, lines 36–57}
\begin{leancode}
\leanline{0}{structure\ RegularDensityDynamics\ (M\ :\ ℕ)\ (U\ :\ Fin\ M\ →\ A\ →\ Z\ →\ ℝ)\ where}
\leanline{2}{parameter\ :\ ℝ\ →\ A}
\leanline{2}{parameterDerivative\ :\ ℝ\ →\ A}
\leanline{2}{density\ :\ Fin\ M\ →\ ℝ\ →\ Z\ →\ ℝ}
\leanline{2}{densityDerivative\ :\ Fin\ M\ →\ ℝ\ →\ Z\ →\ ℝ}
\leanline{2}{parameter\_differentiable\ :\ ∀\ t,\ HasDerivAt\ parameter\ (parameterDerivative\ t)\ t}
\leanline{2}{potential\_parameter\_differentiable\ :\ ∀\ i\ z\ a,\ DifferentiableAt\ ℝ\ (fun\ b\ ↦\ U\ i\ b\ z)\ a}
\leanline{2}{positive\ :\ ∀\ i\ t\ z,\ 0\ <\ density\ i\ t\ z}
\leanline{2}{density\_C1\ :\ ∀\ i\ t,\ ContDiff\ ℝ\ 1\ (density\ i\ t)}
\leanline{2}{potential\_C1\ :\ ∀\ i\ t,\ ContDiff\ ℝ\ 1\ (U\ i\ (parameter\ t))}
\leanline{2}{parameter\_gradient\_integrable\ :\ ∀\ t\ i,}
\leanline{4}{Integrable\ (fun\ z\ ↦\ density\ i\ t\ z\ •\ gradient\ (fun\ a\ ↦\ U\ i\ a\ z)\ (parameter\ t))\ volume}
\leanline{2}{domination\ :\ ∀\ t\ i,\ LocalEnergyDomination\ volume\ (density\ i)\ (densityDerivative\ i)}
\leanline{4}{(fun\ u\ ↦\ U\ i\ (parameter\ u))}
\leanline{4}{(fun\ u\ z\ ↦\ fderiv\ ℝ\ (fun\ a\ ↦\ U\ i\ a\ z)\ (parameter\ u)\ (parameterDerivative\ u))\ t}
\leanline{2}{parameter\_equation\ :\ ∀\ t,\ 0\ ≤\ t\ →\ parameterDerivative\ t\ =}
\leanline{4}{-averagedParameterForce\ volume\ (fun\ i\ ↦\ density\ i\ t)\ U\ (parameter\ t)}
\leanline{2}{density\_equation\ :\ ∀\ t,\ 0\ ≤\ t\ →\ ∀\ i,\ ∀ᵐ\ z\ ∂volume,\ densityDerivative\ i\ t\ z\ =}
\leanline{4}{fluxDivergence\ (fokkerPlanckFlux\ (density\ i\ t)\ (U\ i\ (parameter\ t)))\ z}
\leanline{2}{integrable\_tails\ :\ ∀\ t\ i,\ NoBoundaryFlux}
\leanline{4}{(fokkerPlanckFlux\ (density\ i\ t)\ (U\ i\ (parameter\ t)))}
\leanline{4}{(fun\ z\ ↦\ chemicalPotential\ (density\ i\ t)\ (U\ i\ (parameter\ t))\ z\ +\ 1)}
\end{leancode}
\leansource{\texttt{DirectionalRegularity} \quad LatentDiffusion/\allowbreak{}RegularDirectional.lean, lines 43–61}
\begin{leancode}
\leanline{0}{structure\ DirectionalRegularity}
\leanline{4}{(D\ :\ (i\ :\ Fin\ M)\ →\ RegularDisplacement\ volume\ (μ\ i)\ d)}
\leanline{4}{(U\ :\ Fin\ M\ →\ A\ →\ Z\ →\ ℝ)\ (a\ b\ :\ A)\ where}
\leanline{2}{densityDerivative\ :\ Fin\ M\ →\ ℝ\ →\ Z\ →\ ℝ}
\leanline{2}{domination\ :\ ∀\ i,\ LocalEnergyDomination\ volume\ (D\ i).density\ (densityDerivative\ i)}
\leanline{4}{(fun\ s\ ↦\ U\ i\ (parameterInterpolate\ a\ b\ s))}
\leanline{4}{(fun\ s\ z\ ↦\ fderiv\ ℝ\ (fun\ p\ ↦\ U\ i\ p\ z)\ (parameterInterpolate\ a\ b\ s)\ (b-a))\ 0}
\leanline{2}{parameter\_differentiable\ :\ ∀\ i\ z\ p,\ DifferentiableAt\ ℝ\ (fun\ v\ ↦\ U\ i\ v\ z)\ p}
\leanline{2}{density\_C1\ :\ ∀\ i,\ ContDiff\ ℝ\ 1\ ((D\ i).density\ 0)}
\leanline{2}{potential\_C1\ :\ ∀\ i,\ ContDiff\ ℝ\ 1\ (U\ i\ a)}
\leanline{2}{parameter\_gradient\_integrable\ :\ ∀\ i,\ Integrable}
\leanline{4}{(fun\ z\ ↦\ (D\ i).density\ 0\ z\ •\ gradient\ (fun\ p\ ↦\ U\ i\ p\ z)\ a)\ volume}
\leanline{2}{continuity\ :\ ∀\ i,\ ∀ᵐ\ z\ ∂volume,\ densityDerivative\ i\ 0\ z\ =}
\leanline{4}{-fluxDivergence\ (fun\ w\ ↦\ (D\ i).density\ 0\ w\ •\ ((D\ i).transport\ w-w))\ z}
\leanline{2}{integrable\_tails\ :\ ∀\ i,\ NoBoundaryFlux}
\leanline{4}{(fun\ z\ ↦\ (D\ i).density\ 0\ z\ •\ ((D\ i).transport\ z-z))}
\leanline{4}{(fun\ z\ ↦\ chemicalPotential\ ((D\ i).density\ 0)\ (U\ i\ a)\ z\ +\ 1)}
\leanline{2}{potential\_integrable\ :\ ∀\ i\ s,\ s\ ∈\ Icc\ (0\ :\ ℝ)\ 1\ →}
\leanline{4}{Integrable\ (U\ i\ (parameterInterpolate\ a\ b\ s))\ ((D\ i).law\ s)}
\end{leancode}
\end{leanbox}
\subsection{\texorpdfstring{R4 — Density change of variables and certified optimal maps}{R4 — Density change of variables and certified optimal maps}}


For each displacement argument supply a measurable map \(T:Z\to Z\), differentiable at every point, and \(T_s=(1-s)\mathrm{id}+sT\). For every \(z\), require an orthonormal eigenbasis of \(DT(z)\) with eigenvalues \(\lambda_j(z)>0\). For every \(s\in[0,1]\), supply a measurable everywhere positive density \(q_s\) satisfying
        \[(T_s)_\#\mu=q_s\,dz,\qquad
          q_s(T_s z)=\frac{q_0(z)}{\prod_{j=1}^{d_z}(1-s+s\lambda_j(z))},\qquad
          \log(q_s\circ T_s)\in L^1(\mu).\]
        For each datum the tangent dataset contract further requires \(q_0\in C^1(Z)\), \(\nabla\log q_0\in L^2(\mu)\), \(\mathrm{id}-T\in L^2(\mu)\), and \(z\mapsto\nabla_{(a,z)}U_m(a,z)\in L^2(\mu)\). With \(a_s=(1-s)a+sb\), require \(U_m(a_s,T_s\cdot)\in L^1(\mu)\) for every \(s\in[0,1]\). For \(w=T-\mathrm{id}\), every \(w_j\) is differentiable and \(w_jq_0\), \((\partial_jw_j)q_0\), \(w_j\partial_jq_0\) belong to \(L^1(dz)\). These are the transport-score integration-by-parts inputs, distinct from the flux/\allowbreak{}test pair in R3.




Whenever a map from \(\mu\) to \(\eta\) is used to identify an exact Wasserstein distance, require \(T_\#\mu=\eta\), \(z\mapsto z-Tz\in L^2(\mu)\), and, for every probability coupling \(\gamma\in\Gamma(\mu,\eta)\),
        \[\left(\int\|z-Tz\|^2\,d\mu\right)^{1/2}
          \le\left(\int\|x-y\|^2\,d\gamma(x,y)\right)^{1/2}.\]
        Here coupling costs use the extended nonnegative integral, allowing \(+\infty\). The constant-speed geodesic statements additionally require such an optimal map between \(q_s\) and \(q_t\) for every \(s,t\in[0,1]\). Their converse assumes endpoint uniqueness: every optimal \(S:q_0\to q_1\) satisfies \(S=T\), \(q_0\)-a.e. The gradient-map version also supplies a convex function \(\Phi:Z\to\mathbb R\) with \(\nabla\Phi(z)=T(z)\) at every \(z\). These are assumptions on the chosen transport data. Entropy displacement convexity and the score tangent inequality are derived from them and R3.

\begin{leanbox}{Lean: \texttt{RegularDisplacement}, \texttt{DatasetRegularity}, \texttt{OptimalMap} and 3 more}
\leansource{\texttt{RegularDisplacement} \quad LatentDiffusion/\allowbreak{}RegularTransport.lean, lines 100–119}
\begin{leancode}
\leanline{0}{structure\ RegularDisplacement\ (ν\ μ\ :\ Measure\ Z)\ (d\ :\ ℕ)\ where}
\leanline{2}{transport\ :\ Z\ →\ Z}
\leanline{2}{measurable\_transport\ :\ Measurable\ transport}
\leanline{2}{density\ :\ ℝ\ →\ Z\ →\ ℝ}
\leanline{2}{measurable\_density\ :\ ∀\ s\ ∈\ Icc\ (0\ :\ ℝ)\ 1,\ Measurable\ (density\ s)}
\leanline{2}{positive\_density\ :\ ∀\ s\ ∈\ Icc\ (0\ :\ ℝ)\ 1,\ ∀\ z,\ 0\ <\ density\ s\ z}
\leanline{2}{density\_law\ :\ ∀\ s\ ∈\ Icc\ (0\ :\ ℝ)\ 1,}
\leanline{4}{μ.map\ (interpolate\ transport\ s)\ =\ ν.withDensity\ (fun\ z\ ↦\ ENNReal.ofReal\ (density\ s\ z))}
\leanline{2}{eigenvalue\ :\ Z\ →\ Fin\ d\ →\ ℝ}
\leanline{2}{positive\_eigenvalue\ :\ ∀\ z\ i,\ 0\ <\ eigenvalue\ z\ i}
\leanline{2}{jacobian\ :\ Z\ →\ Z\ →L[ℝ]\ Z}
\leanline{2}{derivative\_transport\ :\ ∀\ z,\ HasFDerivAt\ transport\ (jacobian\ z)\ z}
\leanline{2}{spectral\_basis\ :\ Z\ →\ OrthonormalBasis\ (Fin\ d)\ ℝ\ Z}
\leanline{2}{jacobian\_eigenvector\ :\ ∀\ z\ i,}
\leanline{4}{jacobian\ z\ (spectral\_basis\ z\ i)\ =\ eigenvalue\ z\ i\ •\ spectral\_basis\ z\ i}
\leanline{2}{density\_jacobian\ :\ ∀\ s\ ∈\ Icc\ (0\ :\ ℝ)\ 1,\ ∀\ z,}
\leanline{4}{density\ s\ (interpolate\ transport\ s\ z)\ =}
\leanline{6}{density\ 0\ z\ /\ spectralJacobian\ (eigenvalue\ z)\ s}
\leanline{2}{integrable\_log\ :\ ∀\ s\ ∈\ Icc\ (0\ :\ ℝ)\ 1,}
\leanline{4}{Integrable\ (fun\ z\ ↦\ Real.log\ (density\ s\ (interpolate\ transport\ s\ z)))\ μ}
\end{leancode}
\leansource{\texttt{DatasetRegularity} \quad LatentDiffusion/\allowbreak{}RegularTransportFirstOrder.lean, lines 528–553}
\begin{leancode}
\leanline{0}{structure\ DatasetRegularity\ (M\ d\ :\ ℕ)\ (p₀\ p₁\ :\ P)\ (rho\ :\ ℝ)\ where}
\leanline{2}{law\ :\ Fin\ M\ →\ ProbabilityMeasure\ Z}
\leanline{2}{displacement\ :\ (m\ :\ Fin\ M)\ →\ RegularDisplacement\ volume\ (law\ m\ :\ Measure\ Z)\ d}
\leanline{2}{potential\ :\ Fin\ M\ →\ P\ →\ Z\ →\ ℝ}
\leanline{2}{potential\_strong\ :\ ∀\ m,\ StrongConvexOn\ univ\ rho}
\leanline{4}{(fun\ w\ :\ PBGSO.ProductTransport.Joint\ P\ Z\ ↦\ potential\ m\ w.fst\ w.snd)}
\leanline{2}{potential\_measurable\ :\ ∀\ m\ p,\ Measurable\ (potential\ m\ p)}
\leanline{2}{potential\_integrable\_initial\ :\ ∀\ m,\ Integrable\ (potential\ m\ p₀)\ (law\ m\ :\ Measure\ Z)}
\leanline{2}{potential\_integrable\_final\ :\ ∀\ m,}
\leanline{4}{Integrable\ (fun\ z\ ↦\ potential\ m\ p₁\ ((displacement\ m).transport\ z))\ (law\ m\ :\ Measure\ Z)}
\leanline{2}{potential\_integrable\_path\ :\ ∀\ m\ s,\ s\ ∈\ Icc\ (0\ :\ ℝ)\ 1\ →}
\leanline{4}{Integrable\ (fun\ z\ ↦\ potential\ m\ (parameterInterpolate\ p₀\ p₁\ s)}
\leanline{6}{(interpolate\ (displacement\ m).transport\ s\ z))\ (law\ m\ :\ Measure\ Z)}
\leanline{2}{potential\_gradient\ :\ Fin\ M\ →\ Z\ →\ PBGSO.ProductTransport.Joint\ P\ Z}
\leanline{2}{potential\_hasGradientAt\ :\ ∀\ m\ z,\ HasGradientAt}
\leanline{4}{(fun\ w\ :\ PBGSO.ProductTransport.Joint\ P\ Z\ ↦\ potential\ m\ w.fst\ w.snd)}
\leanline{4}{(potential\_gradient\ m\ z)\ (PBGSO.ProductTransport.embed\ p₀\ z)}
\leanline{2}{potential\_gradient\_memLp\ :\ ∀\ m,\ MemLp\ (potential\_gradient\ m)\ 2\ (law\ m\ :\ Measure\ Z)}
\leanline{2}{displacement\_memLp\ :\ ∀\ m,}
\leanline{4}{MemLp\ (fun\ z\ ↦\ z\ -\ (displacement\ m).transport\ z)\ 2\ (law\ m\ :\ Measure\ Z)}
\leanline{2}{score\_memLp\ :\ ∀\ m,}
\leanline{4}{MemLp\ (fun\ z\ ↦\ gradient\ (fun\ z\ ↦\ Real.log\ ((displacement\ m).density\ 0\ z))\ z)}
\leanline{6}{2\ (law\ m\ :\ Measure\ Z)}
\leanline{2}{density\_C1\ :\ ∀\ m,\ ContDiff\ ℝ\ 1\ ((displacement\ m).density\ 0)}
\leanline{2}{integrable\_tails\ :\ ∀\ m,\ RegularFlowCalculus.NoBoundaryFlux}
\leanline{4}{(fun\ z\ ↦\ (displacement\ m).transport\ z\ -\ z)\ ((displacement\ m).density\ 0)}
\end{leancode}
\leansource{\texttt{OptimalMap} \quad LatentDiffusion/\allowbreak{}RegularTransport.lean, lines 427–433}
\begin{leancode}
\leanline{0}{structure\ OptimalMap\ (μ\ ν\ :\ Measure\ Z)\ where}
\leanline{2}{transport\ :\ Z\ →\ Z}
\leanline{2}{measurable\_transport\ :\ Measurable\ transport}
\leanline{2}{pushes\ :\ μ.map\ transport\ =\ ν}
\leanline{2}{finite\_cost\ :\ MemLp\ (fun\ z\ ↦\ z\ -\ transport\ z)\ 2\ μ}
\leanline{2}{minimal\ :\ ∀\ c\ :\ ULA.Coupling\ μ\ ν,}
\leanline{4}{eLpNorm\ (fun\ z\ ↦\ z\ -\ transport\ z)\ 2\ μ\ ≤\ ULA.couplingCost\ c.joint}
\end{leancode}
\leansource{\texttt{RegularBrenierMap} \quad LatentDiffusion/\allowbreak{}RegularTransportGeodesic.lean, lines 74–77}
\begin{leancode}
\leanline{0}{structure\ RegularBrenierMap\ (μ\ ν\ :\ Measure\ Z)\ extends\ OptimalMap\ μ\ ν\ where}
\leanline{2}{potential\ :\ Z\ →\ ℝ}
\leanline{2}{potential\_convex\ :\ ConvexOn\ ℝ\ univ\ potential}
\leanline{2}{hasGradientAt\ :\ ∀\ z,\ HasGradientAt\ potential\ (transport\ z)\ z}
\end{leancode}
\leansource{\texttt{entropy\_\allowbreak{}displacement\_\allowbreak{}convex} \quad LatentDiffusion/\allowbreak{}RegularTransport.lean, lines 193–204}
\begin{leancode}
\leanline{0}{theorem\ entropy\_displacement\_convex\ (s\ :\ ℝ)\ (hs\ :\ s\ ∈\ Icc\ (0\ :\ ℝ)\ 1)\ :}
\leanline{4}{D.entropy\ s\ ≤\ (1\ -\ s)\ *\ D.entropy\ 0\ +\ s\ *\ D.entropy\ 1}
\end{leancode}
\leansource{\texttt{entropy\_\allowbreak{}tangent} \quad LatentDiffusion/\allowbreak{}RegularTransportFirstOrder.lean, lines 70–92}
\begin{leancode}
\leanline{0}{theorem\ entropy\_tangent}
\leanline{4}{(hq\ :\ ContDiff\ ℝ\ 1\ (D.density\ 0))}
\leanline{4}{(htail\ :\ RegularFlowCalculus.NoBoundaryFlux\ (fun\ z\ ↦\ D.transport\ z\ -\ z)\ (D.density\ 0))\ :}
\leanline{4}{(∫\ z,\ inner\ ℝ\ (D.transport\ z\ -\ z)\ (gradient\ (fun\ z\ ↦\ Real.log\ (D.density\ 0\ z))\ z)\ ∂μ)\ ≤}
\leanline{6}{D.entropy\ 1\ -\ D.entropy\ 0}
\end{leancode}
\end{leanbox}
\subsection{\texorpdfstring{R5 — Differentiable characteristics and optimal coupling anchors}{R5 — Differentiable characteristics and optimal coupling anchors}}


For each datum choose a probability space \((\Omega_m,P_m)\), maps \(Y_m,v_m:\mathbb R\to L^2(P_m;Z)\), and a curve \(x:\mathbb R\to X\), continuous on \([0,\infty)\), in the R6 chart. For every \(t\in\mathbb R\), require
        \[\operatorname{Law}(Y_m(t))=q_{x(t)}^m\,dz,\qquad
          \frac{d}{dt}Y_m(t)=v_m(t)\text{ in }L^2(P_m;Z),\qquad
          v_m(t)=-\nabla_z(\log q_t^m+U_m(a(t),\cdot))(Y_m(t))\quad P_m\text{-a.s.}\]
        Match the chart parameter and density exactly to those of the R3 dynamics: \(a(x(t))=a(t)\), \(q_{x(t)}^m=q_t^m\). For every \(s,t\in\mathbb R\) and \(m\), supply \(A_{s,t}^m\in L^2(P_m;Z)\) with law \(q_{x(s)}^m\,dz\), such that for every \(\gamma\in\Gamma(\operatorname{Law}(Y_m(t)),q_{x(s)}^m dz)\),
        \[\|Y_m(t)-A_{s,t}^m\|_{L^2(P_m)}
          \le\left(\int\|z-w\|^2\,d\gamma(z,w)\right)^{1/2}.\]
        The coupling cost on the right is an extended nonnegative integral. The scalar integrated-speed lemma additionally assumes Bochner interval integrability \(v_m\in L^1([s,t];L^2(P_m;Z))\). The product statement requires the entire lifted velocity \(u\mapsto(\dot a(u),(v_m(u)/\sqrt M)_m)\) Bochner interval integrable on \([s,t]\). The product-space realization uses \(Y_m/\sqrt M\). Distance increments, upper-right slopes and integrated speed bounds follow from these hypotheses; an ordinary distance derivative at a target crossing is not required.

\begin{leanbox}{Lean: \texttt{OptimalAnchor}, \texttt{distance\_\allowbreak{}integral\_\allowbreak{}bound}, \texttt{distance\_\allowbreak{}right\_\allowbreak{}slope} and 2 more}
\leansource{\texttt{OptimalAnchor} \quad LatentDiffusion/\allowbreak{}RegularMetricFlow.lean, lines 21–25}
\begin{leancode}
\leanline{0}{structure\ OptimalAnchor\ (Y\ :\ Lp\ E\ 2\ P)\ (target\ :\ Measure\ E)\ where}
\leanline{2}{value\ :\ Lp\ E\ 2\ P}
\leanline{2}{law\ :\ P.map\ value\ =\ target}
\leanline{2}{minimal\ :\ ∀\ c\ :\ ULA.Coupling\ (P.map\ Y)\ target,}
\leanline{4}{ENNReal.ofReal\ ‖Y\ -\ value‖\ ≤\ ULA.couplingCost\ c.joint}
\end{leancode}
\leansource{\texttt{distance\_\allowbreak{}integral\_\allowbreak{}bound} \quad LatentDiffusion/\allowbreak{}RegularProductMetricFlow.lean, lines 176–199}
\begin{leancode}
\leanline{0}{theorem\ distance\_integral\_bound}
\leanline{4}{(a\ adot\ :\ ℝ\ →\ P)\ (Y\ v\ :\ ℝ\ →\ ∀\ m,\ Lp\ Z\ 2\ (Q\ m))\ (astar\ :\ P)}
\leanline{4}{(target\ :\ Fin\ M\ →\ Measure\ Z)}
\leanline{4}{(anchor\ :\ ∀\ t\ m,\ RegularMetricFlow.OptimalAnchor\ (Q\ m)\ (Y\ t\ m)\ (target\ m))}
\leanline{4}{(s\ t\ :\ ℝ)\ (hst\ :\ s\ ≤\ t)}
\leanline{4}{(ha\ :\ ∀\ u\ ∈\ uIcc\ s\ t,\ HasDerivAt\ a\ (adot\ u)\ u)}
\leanline{4}{(hY\ :\ ∀\ u\ ∈\ uIcc\ s\ t,\ ∀\ m,\ HasDerivAt\ (fun\ r\ ↦\ Y\ r\ m)\ (v\ u\ m)\ u)}
\leanline{4}{(hv\ :\ IntervalIntegrable\ (fun\ u\ ↦\ lift\ Q\ (adot\ u)\ (v\ u))\ volume\ s\ t)\ :}
\leanline{4}{|distance\ Q\ a\ Y\ astar\ target\ t\ -\ distance\ Q\ a\ Y\ astar\ target\ s|\ ≤}
\leanline{6}{∫\ u\ in\ s..t,\ Real.sqrt\ (‖adot\ u‖\ \textasciicircum{}\ 2\ +\ (M\ :\ ℝ)⁻¹\ *\ ∑\ m,\ ‖v\ u\ m‖\ \textasciicircum{}\ 2)}
\end{leancode}
\leansource{\texttt{distance\_\allowbreak{}right\_\allowbreak{}slope} \quad LatentDiffusion/\allowbreak{}RegularProductMetricFlow.lean, lines 149–170}
\begin{leancode}
\leanline{0}{theorem\ distance\_right\_slope}
\leanline{4}{(a\ :\ ℝ\ →\ P)\ (Y\ :\ ℝ\ →\ ∀\ m,\ Lp\ Z\ 2\ (Q\ m))\ (astar\ :\ P)}
\leanline{4}{(target\ :\ Fin\ M\ →\ Measure\ Z)}
\leanline{4}{(anchor\ :\ ∀\ t\ m,\ RegularMetricFlow.OptimalAnchor\ (Q\ m)\ (Y\ t\ m)\ (target\ m))}
\leanline{4}{(t\ :\ ℝ)\ (adot\ :\ P)\ (v\ :\ ∀\ m,\ Lp\ Z\ 2\ (Q\ m))}
\leanline{4}{(ha\ :\ HasDerivAt\ a\ adot\ t)\ (hY\ :\ ∀\ m,\ HasDerivAt\ (fun\ s\ ↦\ Y\ s\ m)\ (v\ m)\ t)}
\leanline{4}{(r\ :\ ℝ)\ (hr\ :\ Real.sqrt\ (‖adot‖\ \textasciicircum{}\ 2\ +\ (M\ :\ ℝ)⁻¹\ *\ ∑\ m,\ ‖v\ m‖\ \textasciicircum{}\ 2)\ <\ r)\ :}
\leanline{4}{∀ᶠ\ s\ in\ {\leanfallback 𝓝}[>]\ t,\ slope\ (distance\ Q\ a\ Y\ astar\ target)\ t\ s\ <\ r}
\end{leancode}
\leansource{\texttt{kinetic\_\allowbreak{}energy\_\allowbreak{}eq\_\allowbreak{}fisher} \quad LatentDiffusion/\allowbreak{}RegularProductMetricFlow.lean, lines 239–251}
\begin{leancode}
\leanline{0}{theorem\ kinetic\_energy\_eq\_fisher}
\leanline{4}{(Y\ v\ :\ ∀\ m,\ Lp\ Z\ 2\ (Q\ m))\ (field\ :\ Fin\ M\ →\ Z\ →\ Z)\ (force\ adot\ :\ P)}
\leanline{4}{(hadot\ :\ adot\ =\ -force)\ (hmeas\ :\ ∀\ m,\ Measurable\ (field\ m))}
\leanline{4}{(hvelocity\ :\ ∀\ m,\ (v\ m\ :\ Ω\ m\ →\ Z)\ =ᵐ[Q\ m]}
\leanline{6}{(fun\ ω\ ↦\ -(field\ m\ ((Y\ m)\ ω))))\ :}
\leanline{4}{‖adot‖\ \textasciicircum{}\ 2\ +\ (M\ :\ ℝ)⁻¹\ *\ ∑\ m,\ ‖v\ m‖\ \textasciicircum{}\ 2\ =}
\leanline{6}{‖force‖\ \textasciicircum{}\ 2\ +\ (M\ :\ ℝ)⁻¹\ *\ ∑\ m,}
\leanline{8}{∫\ z,\ ‖field\ m\ z‖\ \textasciicircum{}\ 2\ ∂(Q\ m).map\ (Y\ m)}
\end{leancode}
\leansource{\texttt{CharacteristicRealization} \quad LatentDiffusion/\allowbreak{}RegularFlowTalagrand.lean, lines 114–127}
\begin{leancode}
\leanline{0}{structure\ CharacteristicRealization}
\leanline{4}{(D\ :\ ProductDensityChart\ X\ A\ Z\ M\ U)\ (S\ :\ RegularDensityDynamics\ M\ U)\ where}
\leanline{2}{state\ :\ ℝ\ →\ X}
\leanline{2}{state\_continuous\ :\ ContinuousOn\ state\ (Ici\ 0)}
\leanline{2}{parameter\_eq\ :\ ∀\ t,\ D.parameter\ (state\ t)\ =\ S.parameter\ t}
\leanline{2}{density\_eq\ :\ ∀\ t\ i,\ (D.chart\ i).density\ (state\ t)\ =\ S.density\ i\ t}
\leanline{2}{characteristic\ :\ ℝ\ →\ ∀\ i,\ Lp\ Z\ 2\ (Q\ i)}
\leanline{2}{velocity\ :\ ℝ\ →\ ∀\ i,\ Lp\ Z\ 2\ (Q\ i)}
\leanline{2}{law\_eq\ :\ ∀\ t\ i,\ (Q\ i).map\ (characteristic\ t\ i)\ =\ (D.law\ (state\ t)\ i\ :\ Measure\ Z)}
\leanline{2}{derivative\ :\ ∀\ t\ i,\ HasDerivAt\ (fun\ u\ ↦\ characteristic\ u\ i)\ (velocity\ t\ i)\ t}
\leanline{2}{velocity\_eq\ :\ ∀\ t\ i,\ (velocity\ t\ i\ :\ Ω\ i\ →\ Z)\ =ᵐ[Q\ i]}
\leanline{4}{(fun\ ω\ ↦\ -(S.score\ t\ i\ (characteristic\ t\ i\ ω)))}
\leanline{2}{anchor\ :\ ∀\ s\ t\ i,\ RegularMetricFlow.OptimalAnchor\ (Q\ i)\ (characteristic\ t\ i)}
\leanline{4}{(D.law\ (state\ s)\ i\ :\ Measure\ Z)}
\end{leancode}
\end{leanbox}
\subsection{\texorpdfstring{R6 — A complete density chart with the exact transport metric}{R6 — A complete density chart with the exact transport metric}}


Let \((X,d_X)\) be a complete metric space with a parameter map \(a:X\to A\), positive measurable densities \(q_x^m\), and laws \(\mu_x^m=q_x^m dz\in\mathcal P_2(Z)\). For every \(x,y\in X\), require the exact identity
        \[d_X(x,y)^2=\|a(x)-a(y)\|^2+
          \frac1M\sum_m W_2(\mu_x^m,\mu_y^m)^2.\]
        For every \(m,z\), assume both \(x\mapsto q_x^m(z)\) and \(x\mapsto U_m(a(x),z)\) continuous. Require spatial measurability of each density and potential. For every \(x\in X\) and \(m\), assume a neighbourhood \(V_{x,m}\) and \(g_{x,m}\in L^1(dz)\) such that for every \(y\in V_{x,m}\), almost every \(z\),
        \[\left|q_y^m(z)\log q_y^m(z)+q_y^m(z)U_m(a(y),z)\right|\le g_{x,m}(z).\]
        These assumptions imply continuity of \(F(x)=M^{-1}\sum_m\int(q_x^m\log q_x^m+q_x^mU_m(a(x),\cdot))\,dz\). For the minimizing-limit statement, assume a finite lower bound \(F_\star\le F(x)\) for every \(x\), a Cauchy sequence \(x_n\) in this same chart, and \(F(x_n)\to F_\star\). A minimizing limit in \(X\) is then proved. The flow application constructs the Cauchy sequence from \(x(t_n)\), \(t_n\to\infty\). Pointwise density continuity and completeness of this chosen chart are extra assumptions; neither is inferred from membership in \(\mathcal P_2\).

\begin{leanbox}{Lean: \texttt{ProductDensityChart}, \texttt{ProductDensityChart.\allowbreak{}continuous\_\allowbreak{}energy}, \texttt{ProductDensityChart.\allowbreak{}lowerSemicontinuous\_\allowbreak{}energy} and 1 more}
\leansource{\texttt{ProductDensityChart} \quad LatentDiffusion/\allowbreak{}RegularDensityDomain.lean, lines 32–45}
\begin{leancode}
\leanline{0}{structure\ ProductDensityChart\ (X\ A\ Z\ :\ Type*)\ [MetricSpace\ X]}
\leanline{4}{[NormedAddCommGroup\ A]\ [NormedSpace\ ℝ\ A]}
\leanline{4}{[NormedAddCommGroup\ Z]\ [InnerProductSpace\ ℝ\ Z]\ [FiniteDimensional\ ℝ\ Z]}
\leanline{4}{[MeasurableSpace\ Z]\ [BorelSpace\ Z]\ (M\ :\ ℕ)\ (U\ :\ Fin\ M\ →\ A\ →\ Z\ →\ ℝ)\ where}
\leanline{2}{parameter\ :\ X\ →\ A}
\leanline{2}{chart\ :\ Fin\ M\ →\ RegularVariational.DominatedDensityChart\ X\ Z\ volume}
\leanline{2}{law\ :\ X\ →\ Fin\ M\ →\ ProbabilityMeasure\ Z}
\leanline{2}{finite\_second\_moment\ :\ ∀\ x\ i,\ MemLp\ (fun\ z\ :\ Z\ ↦\ z)\ 2\ (law\ x\ i\ :\ Measure\ Z)}
\leanline{2}{density\_law\ :\ ∀\ x\ i,\ (law\ x\ i\ :\ Measure\ Z)\ =}
\leanline{4}{volume.withDensity\ (fun\ z\ ↦\ ENNReal.ofReal\ ((chart\ i).density\ x\ z))}
\leanline{2}{positive\ :\ ∀\ x\ i\ z,\ 0\ <\ (chart\ i).density\ x\ z}
\leanline{2}{potential\_eq\ :\ ∀\ x\ i\ z,\ (chart\ i).potential\ x\ z\ =\ U\ i\ (parameter\ x)\ z}
\leanline{2}{metric\_eq\ :\ ∀\ x\ y,\ dist\ x\ y\ \textasciicircum{}\ 2\ =\ ‖parameter\ x-parameter\ y‖\ \textasciicircum{}\ 2\ +}
\leanline{4}{(M\ :\ ℝ)⁻¹\ *\ ∑\ i,\ ULA.wasserstein2Real\ (law\ x\ i\ :\ Measure\ Z)\ (law\ y\ i\ :\ Measure\ Z)\ \textasciicircum{}\ 2}
\end{leancode}
\leansource{\texttt{ProductDensityChart.\allowbreak{}continuous\_\allowbreak{}energy} \quad LatentDiffusion/\allowbreak{}RegularDensityDomain.lean, lines 66–72}
\begin{leancode}
\leanline{0}{theorem\ ProductDensityChart.continuous\_energy\ :\ Continuous\ D.energy}
\end{leancode}
\leansource{\texttt{ProductDensityChart.\allowbreak{}lowerSemicontinuous\_\allowbreak{}energy} \quad LatentDiffusion/\allowbreak{}RegularDensityDomain.lean, lines 74–75}
\begin{leancode}
\leanline{0}{theorem\ ProductDensityChart.lowerSemicontinuous\_energy\ :\ LowerSemicontinuous\ D.energy}
\end{leancode}
\leansource{\texttt{ProductDensityChart.\allowbreak{}cauchy\_\allowbreak{}minimizer} \quad LatentDiffusion/\allowbreak{}RegularDensityDomain.lean, lines 79–85}
\begin{leancode}
\leanline{0}{theorem\ ProductDensityChart.cauchy\_minimizer\ [CompleteSpace\ X]}
\leanline{4}{(Fstar\ :\ ℝ)\ (hstar\ :\ ∀\ x,\ Fstar\ ≤\ D.energy\ x)\ (x\ :\ ℕ\ →\ X)}
\leanline{4}{(hCauchy\ :\ CauchySeq\ x)\ (hvalue\ :\ Tendsto\ (fun\ n\ ↦\ D.energy\ (x\ n))\ atTop\ ({\leanfallback 𝓝}\ Fstar))\ :}
\leanline{4}{∃\ xstar,\ Tendsto\ x\ atTop\ ({\leanfallback 𝓝}\ xstar)\ ∧\ D.energy\ xstar\ =\ Fstar\ ∧}
\leanline{6}{∀\ y,\ D.energy\ xstar\ ≤\ D.energy\ y}
\end{leancode}
\end{leanbox}
\subsection{\texorpdfstring{R7 — Brownian inputs, posterior initialization and score moments}{R7 — Brownian inputs, posterior initialization and score moments}}


Write \(K=MN>0\) and use the normalized cloud norm \(\|(a,z)\|_K^2=\|a\|^2+K^{-1}\sum_{i=1}^K\|z_i\|^2\). On a probability space \((\Omega,P)\), supply \(Z_0\in L^2(P;Z^K)\) whose law is the actual Gibbs law
        \[\operatorname{Law}(Z_0)(dz)=
          \frac{\exp(-\sum_{i=1}^K U_i(a_\star,z_i))}{\int_{Z^K}\exp(-\sum_{i=1}^K U_i(a_\star,w_i))\,dw}\,dz,\qquad
          \operatorname{Law}(Z_{0,i})=\pi_{a_\star}^i.\]
        Here \(a_\star\) is the proved evidence optimizer, the datum potentials are repeated across their \(N\) copies, and R2 normalizer domination holds at \(a_\star\). Require the parameter scores \(S_i=\nabla_aU_i(a_\star,Z_{0,i})\) to be independent across all \(i\). For the explicit particle-count rate assume
        \[\mathbb E\|S_i-\mathbb ES_i\|^2\le V_0\quad(1\le i\le K).\]
        The hypotheses of the rate theorem supply one \(V_0\); a claim uniform in \(N\) requires that same bound for every \(N\).




Supply real coordinate processes \(B_{i,j}(t)\), \(t\ge0\), whose every finite-dimensional law is centered Gaussian with covariance \(\mathbb E[B_{i,j}(s)B_{i,j}(t)]=\min(s,t)\). In particular, \(B_{i,j}(0)=0\) almost surely and \(B_{i,j}(t)-B_{i,j}(s)\sim N(0,t-s)\) for \(0\le s\le t\). For every \(s>0\), the entire family \(\Delta B_{n,i,j}^{(s)}=B_{i,j}((n+1)s)-B_{i,j}(ns)\), indexed by \((n,i,j)\), must be independent, and \(Z_0\) must be independent of the complete sequence \((\Delta B_{n,i,j}^{(s)})_{n,i,j}\). The continuous comparisons use \(L^2(P)\)-continuous curves solving the actual Bochner integral equations for their displayed forces and these same Brownian inputs. For example, the full particle curve satisfies
        \[X_t-X_0=-\int_0^t F(X_r)\,dr+
          (0,\sqrt2(B_t-B_0)),\qquad t\ge0,\]
        The force components are \[F_a(a,z)=K^{-1}\sum_i\nabla_aU_i(a,z_i),\qquad F_{z_i}(a,z)=\nabla_zU_i(a,z_i),\] and norms are measured by \(\|\cdot\|_K\). The frozen curve starts at \((a_\star,Z_0)\); its existence and stationarity are derived from the Gibbs initialization and independent increments. A spatial comparison that omits an initial-error term assumes identical coupled initial states; otherwise it retains that term. The nonstationary mean-field comparison separately requires independent parameter-score copies at every \(t\ge0\) and \(\mathbb E\|S_i(t)-\mathbb ES_i(t)\|^2\le V_0\) for every such \(t\) and \(i\). The direct temporal comparison assumes \(\|X_t\|_{L^2(P;\|\cdot\|_K)}\le R_K\) for every \(t\ge0\). Constants uniform in the particle count require \(\sup_KR_K<\infty\) and \(\sup_K\|F_K(0)\|_K<\infty\) in that branch.




For the full-batch Euler bound require an arbitrary \(X_0\in L^2(P)\) and the exact stable-step restrictions
        \[h>0,\qquad \lambda>0,\qquad L\ge0,\qquad \lambda h\le1,\qquad hL^2\le\lambda.\]
        A count-independent transient requires \(\sup_N\|X_0-(a_\star,Z_0)\|_{L^2(P;\|\cdot\|_K)}<\infty\). The endpoint derives its frozen-force moment bound from A3–A4 and the Gibbs law; it does not assume a temporal error recurrence. The asserted continuous solution class is continuity in \(L^2(P)\) with the displayed integral identity.

\begin{leanbox}{Lean: \texttt{sourceReference}, \texttt{exists\_\allowbreak{}spatial\_\allowbreak{}comparison}, \texttt{temporal\_\allowbreak{}comparison} and 5 more}
\leansource{\texttt{sourceReference} \quad LatentDiffusion/\allowbreak{}RegularStochastic.lean, lines 861–884}
\begin{leancode}
\leanline{0}{def\ sourceReference}
\leanline{4}{(hK\ :\ 0\ <\ K)\ (U\ :\ Fin\ K\ →\ Cloud.Joint\ dθ\ dz\ →\ ℝ)}
\leanline{4}{(m\ :\ ℝ)\ (L\ :\ ℝ≥0)\ (hm\ :\ 0\ <\ m)\ (hdiff\ :\ ∀\ i,\ Differentiable\ ℝ\ (U\ i))}
\leanline{4}{(hsc\ :\ ∀\ i,\ StrongConvexOn\ Set.univ\ m\ (U\ i))}
\leanline{4}{(hLip\ :\ ∀\ i,\ LipschitzWith\ L\ (gradient\ (U\ i)))\ (astar\ :\ Parameter\ dθ)}
\leanline{4}{(initial\ :\ Lp\ (Latent\ K\ dz)\ 2\ P)}
\leanline{4}{(hinitial\ :\ P.map\ initial\ =\ ULA.gibbsMeasure\ (frozenPotential\ U\ astar))}
\leanline{4}{(B\ :\ (Fin\ K\ ×\ Fin\ dz)\ →\ ℝ≥0\ →\ Ω\ →\ ℝ)}
\leanline{4}{(hB\ :\ ∀\ i,\ IsPreBrownianReal\ (B\ i)\ P)}
\leanline{4}{(hscalar\ :\ ∀\ h\ :\ ℝ,\ 0\ <\ h\ →\ iIndepFun\ (fun\ p\ :\ ℕ\ ×\ (Fin\ K\ ×\ Fin\ dz)\ =>\ fun\ ω\ =>}
\leanline{6}{ULA.Hilbert.brownianEulerNoise\ P\ B\ hB\ h\ p.1\ ω\ p.2)\ P)}
\leanline{4}{(hind\ :\ ∀\ h\ :\ ℝ,\ 0\ <\ h\ →\ IndepFun\ initial}
\leanline{6}{(fun\ ω\ n\ =>\ ULA.Hilbert.brownianEulerNoise\ P\ B\ hB\ h\ n\ ω)\ P)\ :}
\leanline{4}{FrozenPosteriorSolution\ P\ (Cloud.meanForce\_regular\ hK\ U\ m\ L\ hdiff\ hsc\ hLip)\ B\ hB\ :=\ by}
\leanline{2}{let\ V\ :=\ frozenPotential\ U\ astar}
\leanline{2}{let\ hv\ :=\ frozenPotential\_gradient\_lipschitz\ hK\ U\ L\ hdiff\ hLip\ astar}
\leanline{2}{let\ Z\ :=\ latentFlow\ P\ V\ L\ hv\ initial\ B\ hB}
\leanline{2}{have\ hs\ :=\ latentFlow\_spec\ P\ V\ L\ hv\ initial\ B\ hB}
\leanline{2}{have\ hzero\ :\ Z\ 0\ =\ initial\ :=\ hs.2.1}
\leanline{2}{exact\ FrozenPosteriorSolution.of\_latent\_gibbs\ P}
\leanline{4}{(Cloud.meanForce\_regular\ hK\ U\ m\ L\ hdiff\ hsc\ hLip)\ V\ m\ L\ hm}
\leanline{4}{(frozenPotential\_differentiable\ U\ hdiff\ astar)\ (frozenPotential\_strongConvex\ hK\ U\ m\ hsc\ astar)}
\leanline{4}{hv\ astar\ (frozenPotential\_force\_identity\ hK\ U\ hdiff\ astar)\ B\ hB\ Z\ hs.1\ hs.2.2}
\leanline{4}{(by\ rw\ [hzero];\ exact\ hinitial)\ hscalar\ (by\ rw\ [hzero];\ exact\ hind)}
\end{leancode}
\leansource{\texttt{exists\_\allowbreak{}spatial\_\allowbreak{}comparison} \quad LatentDiffusion/\allowbreak{}RegularMeanField.lean, lines 189–234}
\begin{leancode}
\leanline{0}{theorem\ exists\_spatial\_comparison}
\leanline{4}{(hK\ :\ 0\ <\ K)\ (U\ :\ Fin\ K\ →\ Cloud.Joint\ dθ\ dz\ →\ ℝ)}
\leanline{4}{(m\ :\ ℝ)\ (L\ :\ ℝ≥0)\ (hm\ :\ 0\ <\ m)\ (hdiff\ :\ ∀\ i,\ Differentiable\ ℝ\ (U\ i))}
\leanline{4}{(hsc\ :\ ∀\ i,\ StrongConvexOn\ Set.univ\ m\ (U\ i))}
\leanline{4}{(hLip\ :\ ∀\ i,\ LipschitzWith\ L\ (gradient\ (U\ i)))}
\leanline{4}{(Y\ :\ ℝ\ →\ Lp\ (Cloud.State\ K\ dθ\ dz)\ 2\ P)\ (hY\ :\ Continuous\ Y)}
\leanline{4}{(B\ :\ ℝ\ →\ Lp\ (Cloud.State\ K\ dθ\ dz)\ 2\ P)\ (hB\ :\ Continuous\ B)\ (hB0\ :\ B\ 0=0)}
\leanline{4}{(hYe\ :\ ∀\ t,\ 0\ ≤\ t\ →\ Y\ t-Y\ 0\ =}
\leanline{6}{-(∫\ r\ in\ 0..t,\ force\ P\ (Cloud.meanForce\_regular\ hK\ U\ m\ L\ hdiff\ hsc\ hLip)\ (Y\ r))+B\ t)}
\leanline{4}{(V\ :\ ℝ)\ (\_hV\ :\ 0\ ≤\ V)}
\leanline{4}{(hind\ :\ ∀\ t,\ 0\ ≤\ t\ →\ iIndepFun\ (fun\ i\ ω\ =>\ scoreLp\ P\ U\ L\ hLip\ (Y\ t)\ i\ ω)\ P)}
\leanline{4}{(hvar\ :\ ∀\ t,\ 0\ ≤\ t\ →\ ∀\ i,\ ‖scoreLp\ P\ U\ L\ hLip\ (Y\ t)\ i\ -}
\leanline{6}{constantL2\ P\ (∫\ ω,\ scoreLp\ P\ U\ L\ hLip\ (Y\ t)\ i\ ω\ ∂P)‖\textasciicircum{}2\ ≤\ V)\ :}
\leanline{4}{∃\ X\ :\ ℝ\ →\ Lp\ (Cloud.State\ K\ dθ\ dz)\ 2\ P,\ Continuous\ X\ ∧\ X\ 0=Y\ 0\ ∧}
\leanline{6}{(∀\ t,\ 0\ ≤\ t\ →\ X\ t-X\ 0\ =}
\leanline{8}{-(∫\ r\ in\ 0..t,\ particleForce\ P\ hK\ U\ m\ L\ hdiff\ hsc\ hLip\ (X\ r))+B\ t)\ ∧}
\leanline{6}{∀\ t,\ 0\ ≤\ t\ →\ ‖X\ t-Y\ t‖\ ≤\ Real.sqrt\ (V/K)/m*(1-Real.exp\ (-m*t))}
\end{leancode}
\leansource{\texttt{temporal\_\allowbreak{}comparison} \quad LatentDiffusion/\allowbreak{}RegularMeanField.lean, lines 253–301}
\begin{leancode}
\leanline{0}{theorem\ temporal\_comparison}
\leanline{4}{(hK\ :\ 0\ <\ K)\ (U\ :\ Fin\ K\ →\ Cloud.Joint\ dθ\ dz\ →\ ℝ)}
\leanline{4}{(m\ h\ R\ :\ ℝ)\ (L\ :\ ℝ≥0)\ (hR\ :\ 0\ ≤\ R)}
\leanline{4}{(hdiff\ :\ ∀\ i,\ Differentiable\ ℝ\ (U\ i))}
\leanline{4}{(hsc\ :\ ∀\ i,\ StrongConvexOn\ Set.univ\ m\ (U\ i))}
\leanline{4}{(hLip\ :\ ∀\ i,\ LipschitzWith\ L\ (gradient\ (U\ i)))\ (hs\ :\ StableStep\ m\ L\ h)}
\leanline{4}{(B\ :\ (Fin\ K\ ×\ Fin\ dz)\ →\ ℝ≥0\ →\ Ω\ →\ ℝ)}
\leanline{4}{(hB\ :\ ∀\ i,\ IsPreBrownianReal\ (B\ i)\ P)}
\leanline{4}{(X\ :\ ℝ\ →\ Lp\ (Cloud.State\ K\ dθ\ dz)\ 2\ P)\ (hX\ :\ Continuous\ X)}
\leanline{4}{(hXe\ :\ ∀\ t,\ 0\ ≤\ t\ →\ X\ t-X\ 0\ =}
\leanline{6}{-(∫\ r\ in\ 0..t,\ particleForce\ P\ hK\ U\ m\ L\ hdiff\ hsc\ hLip\ (X\ r))+physicalNoise\ P\ B\ hB\ t)}
\leanline{4}{(hMoment\ :\ ∀\ t,\ 0\ ≤\ t\ →\ ‖X\ t‖\ ≤\ R)\ (n\ :\ ℕ)\ :}
\leanline{4}{‖eulerTrajectory\ (particleForce\ P\ hK\ U\ m\ L\ hdiff\ hsc\ hLip)\ h}
\leanline{6}{(fun\ k\ =>\ physicalNoise\ P\ B\ hB\ ((k\ :\ ℝ)*h+h)-physicalNoise\ P\ B\ hB\ ((k\ :\ ℝ)*h))\ (X\ 0)\ n-}
\leanline{6}{X\ ((n\ :\ ℝ)*h)‖\ ≤}
\leanline{6}{(L\ :\ ℝ)/m\ *\ (h*(‖particleForce\ P\ hK\ U\ m\ L\ hdiff\ hsc\ hLip\ 0‖+(L\ :\ ℝ)*R)+}
\leanline{8}{(4/3\ :\ ℝ)*Real.sqrt\ (2*(dz\ :\ ℝ))*Real.sqrt\ h)}
\end{leancode}
\leansource{\texttt{independent\_\allowbreak{}average\_\allowbreak{}norm\_\allowbreak{}sq} \quad LatentDiffusion/\allowbreak{}RegularStochastic.lean, lines 214–256}
\begin{leancode}
\leanline{0}{theorem\ independent\_average\_norm\_sq\ \{E\ :\ Type*\}}
\leanline{4}{[NormedAddCommGroup\ E]\ [InnerProductSpace\ ℝ\ E]\ [CompleteSpace\ E]}
\leanline{4}{[MeasurableSpace\ E]\ [BorelSpace\ E]}
\leanline{4}{(g\ :\ Fin\ K\ →\ Lp\ E\ 2\ P)}
\leanline{4}{(hind\ :\ iIndepFun\ (fun\ i\ ω\ =>\ g\ i\ ω)\ P)}
\leanline{4}{(hopt\ :\ ∑\ i,\ (∫\ ω,\ g\ i\ ω\ ∂P)\ =\ 0)\ :}
\leanline{4}{‖MiniBatchVariance.average\ g‖\ \textasciicircum{}\ 2\ =\ (1\ /\ (K\ :\ ℝ))\ \textasciicircum{}\ 2\ *}
\leanline{6}{∑\ i,\ ‖g\ i\ -\ constantL2\ P\ (∫\ ω,\ g\ i\ ω\ ∂P)‖\ \textasciicircum{}\ 2}
\end{leancode}
\leansource{\texttt{full\_\allowbreak{}batch\_\allowbreak{}parameter\_\allowbreak{}bound} \quad LatentDiffusion/\allowbreak{}RegularStochastic.lean, lines 482–528}
\begin{leancode}
\leanline{0}{theorem\ full\_batch\_parameter\_bound}
\leanline{4}{(hK\ :\ 0\ <\ K)\ (U\ :\ Fin\ K\ →\ Cloud.Joint\ dθ\ dz\ →\ ℝ)}
\leanline{4}{(m\ h\ :\ ℝ)\ (L\ :\ ℝ≥0)\ (hdiff\ :\ ∀\ i,\ Differentiable\ ℝ\ (U\ i))}
\leanline{4}{(hsc\ :\ ∀\ i,\ StrongConvexOn\ Set.univ\ m\ (U\ i))}
\leanline{4}{(hLip\ :\ ∀\ i,\ LipschitzWith\ L\ (gradient\ (U\ i)))\ (hs\ :\ StableStep\ m\ L\ h)}
\leanline{4}{(B\ :\ (Fin\ K\ ×\ Fin\ dz)\ →\ ℝ≥0\ →\ Ω\ →\ ℝ)}
\leanline{4}{(hB\ :\ ∀\ i,\ IsPreBrownianReal\ (B\ i)\ P)}
\leanline{4}{(S\ :\ FrozenPosteriorSolution\ P\ (Cloud.meanForce\_regular\ hK\ U\ m\ L\ hdiff\ hsc\ hLip)\ B\ hB)}
\leanline{4}{(hind\ :\ iIndepFun\ (fun\ i\ ω\ =>\ scoreLp\ P\ U\ L\ hLip\ (S.path\ 0)\ i\ ω)\ P)}
\leanline{4}{(hopt\ :\ ∑\ i,\ (∫\ ω,\ scoreLp\ P\ U\ L\ hLip\ (S.path\ 0)\ i\ ω\ ∂P)\ =\ 0)}
\leanline{4}{(x₀\ :\ Lp\ (Cloud.State\ K\ dθ\ dz)\ 2\ P)\ (astar\ :\ Parameter\ dθ)}
\leanline{4}{(hstar\ :\ parameterOutput\ P\ (S.path\ 0)\ =\ constantL2\ P\ astar)\ (n\ :\ ℕ)\ :}
\leanline{4}{‖parameterOutput\ P}
\leanline{6}{(eulerTrajectory\ (particleForce\ P\ hK\ U\ m\ L\ hdiff\ hsc\ hLip)\ h}
\leanline{8}{(fun\ k\ =>\ physicalNoise\ P\ B\ hB\ ((k\ :\ ℝ)*h+h)\ -\ physicalNoise\ P\ B\ hB\ ((k\ :\ ℝ)*h))\ x₀\ n)\ -}
\leanline{6}{constantL2\ P\ astar‖\ ≤}
\leanline{6}{(1\ -\ m*h/2)\textasciicircum{}n\ *\ ‖x₀\ -\ S.path\ 0‖\ +}
\leanline{6}{(2/m)\ *\ Real.sqrt\ ((1/(K\ :\ ℝ))\textasciicircum{}2\ *\ scoreVarianceBudget\ P\ U\ L\ hLip\ (S.path\ 0))\ +}
\leanline{6}{(L\ :\ ℝ)/m\ *\ (h\ *\ ‖frozenForce\ P\ (Cloud.meanForce\_regular\ hK\ U\ m\ L\ hdiff\ hsc\ hLip)}
\leanline{8}{(S.path\ 0)‖\ +\ (4/3\ :\ ℝ)*physicalNoiseScale\ K\ dz*Real.sqrt\ h)}
\end{leancode}
\leansource{\texttt{full\_\allowbreak{}batch\_\allowbreak{}optimizer\_\allowbreak{}rms\_\allowbreak{}rate} \quad LatentDiffusion/\allowbreak{}RegularSourceTheorem.lean, lines 167–209}
\begin{leancode}
\leanline{0}{theorem\ full\_batch\_optimizer\_rms\_rate}
\leanline{4}{\{M\ N\ :\ ℕ\}\ (hMN\ :\ 0\ <\ M*N)}
\leanline{4}{(U\ :\ Fin\ (M*N)\ →\ Cloud.Joint\ dθ\ dz\ →\ ℝ)}
\leanline{4}{(m\ h\ :\ ℝ)\ (L\ :\ ℝ≥0)\ (hm\ :\ 0\ <\ m)}
\leanline{4}{(hdiff\ :\ ∀\ i,\ Differentiable\ ℝ\ (U\ i))}
\leanline{4}{(hsc\ :\ ∀\ i,\ StrongConvexOn\ Set.univ\ m\ (U\ i))}
\leanline{4}{(hLip\ :\ ∀\ i,\ LipschitzWith\ L\ (gradient\ (U\ i)))\ (hs\ :\ StableStep\ m\ L\ h)}
\leanline{4}{(initial\ :\ Lp\ (Latent\ (M*N)\ dz)\ 2\ P)}
\leanline{4}{(hinitial\ :\ P.map\ initial\ =\ ULA.gibbsMeasure}
\leanline{6}{(frozenPotential\ U\ (optimizerParameter\ hMN\ U\ m\ hm\ hdiff\ hsc)))}
\leanline{4}{(hposterior\ :\ ∀\ i,\ P.map\ (fun\ ω\ =>\ physicalParticleCLM\ i}
\leanline{8}{((packLp\ P\ (optimizerParameter\ hMN\ U\ m\ hm\ hdiff\ hsc)\ initial)\ ω))\ =}
\leanline{6}{RegularEvidence.posterior\ volume\ (latentPotential\ U\ i)\ (optimizerParameter\ hMN\ U\ m\ hm\ hdiff\ hsc))}
\leanline{4}{(data\ :\ ∀\ i,\ RegularEvidence.LocalNormalizerRegularity\ volume\ (latentPotential\ U\ i)}
\leanline{6}{(optimizerParameter\ hMN\ U\ m\ hm\ hdiff\ hsc))}
\leanline{4}{(hindScores\ :\ iIndepFun\ (fun\ i\ ω\ =>\ scoreLp\ P\ U\ L\ hLip}
\leanline{6}{(packLp\ P\ (optimizerParameter\ hMN\ U\ m\ hm\ hdiff\ hsc)\ initial)\ i\ ω)\ P)}
\leanline{4}{(B\ :\ (Fin\ (M*N)\ ×\ Fin\ dz)\ →\ ℝ≥0\ →\ Ω\ →\ ℝ)}
\leanline{4}{(hB\ :\ ∀\ i,\ IsPreBrownianReal\ (B\ i)\ P)}
\leanline{4}{(hscalar\ :\ ∀\ s\ :\ ℝ,\ 0\ <\ s\ →\ iIndepFun\ (fun\ p\ :\ ℕ\ ×\ (Fin\ (M*N)\ ×\ Fin\ dz)\ =>\ fun\ ω\ =>}
\leanline{6}{ULA.Hilbert.brownianEulerNoise\ P\ B\ hB\ s\ p.1\ ω\ p.2)\ P)}
\leanline{4}{(hind\ :\ ∀\ s\ :\ ℝ,\ 0\ <\ s\ →\ IndepFun\ initial}
\leanline{6}{(fun\ ω\ n\ =>\ ULA.Hilbert.brownianEulerNoise\ P\ B\ hB\ s\ n\ ω)\ P)}
\leanline{4}{(V₀\ :\ ℝ)}
\leanline{4}{(hvar\ :\ ∀\ i,\ ‖scoreLp\ P\ U\ L\ hLip}
\leanline{8}{(packLp\ P\ (optimizerParameter\ hMN\ U\ m\ hm\ hdiff\ hsc)\ initial)\ i-}
\leanline{6}{constantL2\ P\ (∫\ ω,\ scoreLp\ P\ U\ L\ hLip}
\leanline{8}{(packLp\ P\ (optimizerParameter\ hMN\ U\ m\ hm\ hdiff\ hsc)\ initial)\ i\ ω\ ∂P)‖\textasciicircum{}2\ ≤\ V₀)}
\leanline{4}{(x₀\ :\ Lp\ (Cloud.State\ (M*N)\ dθ\ dz)\ 2\ P)\ (n\ :\ ℕ)\ :}
\leanline{4}{Real.sqrt\ (∫\ ω,\ ‖IPLA.parameterCLM}
\leanline{6}{(eulerTrajectory\ (Cloud.meanForce\ U)\ h}
\leanline{8}{(fun\ k\ =>\ noiseIncrement\ P\ B\ hB\ h\ k\ ω)\ (x₀\ ω)\ n)-}
\leanline{6}{optimizerParameter\ hMN\ U\ m\ hm\ hdiff\ hsc‖\textasciicircum{}2\ ∂P)\ ≤}
\leanline{6}{(1-m*h/2)\textasciicircum{}n\ *\ ‖x₀-packLp\ P\ (optimizerParameter\ hMN\ U\ m\ hm\ hdiff\ hsc)\ initial‖\ +}
\leanline{6}{(2/m)\ *\ Real.sqrt\ (V₀/((M*N\ :\ ℕ)\ :\ ℝ))\ +}
\leanline{6}{(L\ :\ ℝ)/m\ *\ (h*Real.sqrt\ ((L\ :\ ℝ)*(dz\ :\ ℝ))+}
\leanline{8}{(4/3\ :\ ℝ)*Real.sqrt\ (2*(dz\ :\ ℝ))*Real.sqrt\ h)}
\end{leancode}
\leansource{\texttt{optimizer\_\allowbreak{}replication} \quad LatentDiffusion/\allowbreak{}RegularSourceTheorem.lean, lines 250–262}
\begin{leancode}
\leanline{0}{theorem\ optimizer\_replication\ \{M\ N\ :\ ℕ\}\ (hM\ :\ 0\ <\ M)\ (hN\ :\ 0\ <\ N)}
\leanline{4}{(U\ :\ Fin\ M\ →\ Cloud.Joint\ dθ\ dz\ →\ ℝ)\ (m\ :\ ℝ)\ (hm\ :\ 0\ <\ m)}
\leanline{4}{(hdiff\ :\ ∀\ i,\ Differentiable\ ℝ\ (U\ i))}
\leanline{4}{(hsc\ :\ ∀\ i,\ StrongConvexOn\ Set.univ\ m\ (U\ i))\ :}
\leanline{4}{optimizerParameter\ (Nat.mul\_pos\ hM\ hN)\ (replicatePotentials\ (N\ :=\ N)\ U)\ m\ hm}
\leanline{6}{(fun\ i\ =>\ hdiff\ i.divNat)\ (fun\ i\ =>\ hsc\ i.divNat)\ =}
\leanline{8}{optimizerParameter\ hM\ U\ m\ hm\ hdiff\ hsc}
\end{leancode}
\leansource{\texttt{sum\_\allowbreak{}posterior\_\allowbreak{}parameter\_\allowbreak{}gradient\_\allowbreak{}eq\_\allowbreak{}zero} \quad LatentDiffusion/\allowbreak{}RegularEvidence.lean, lines 185–193}
\begin{leancode}
\leanline{0}{theorem\ sum\_posterior\_parameter\_gradient\_eq\_zero\ \{M\ :\ ℕ\}\ (hM\ :\ 0\ <\ M)}
\leanline{4}{(ν\ :\ Measure\ Z)\ (U\ :\ Fin\ M\ →\ A\ →\ Z\ →\ ℝ)\ (a\ :\ A)}
\leanline{4}{(data\ :\ ∀\ i,\ LocalNormalizerRegularity\ ν\ (U\ i)\ a)}
\leanline{4}{(hmeas\ :\ ∀\ i,\ Measurable\ (fun\ z\ ↦\ Real.exp\ (-U\ i\ a\ z)))}
\leanline{4}{(hZ\ :\ ∀\ i,\ 0\ <\ normalizer\ ν\ (U\ i)\ a)}
\leanline{4}{(hmax\ :\ IsLocalMax\ (logEvidence\ ν\ U)\ a)\ :}
\leanline{4}{∑\ i,\ (∫\ z,\ gradient\ (fun\ b\ ↦\ U\ i\ b\ z)\ a\ ∂posterior\ ν\ (U\ i)\ a)\ =\ 0}
\end{leancode}
\end{leanbox}
\section{Density model and analytic chain}
\subsection{\texorpdfstring{V1 — Normalized free energy and finite KL domain}{V1 — Normalized free energy and finite KL domain}}



Let \(M>0\), let \(\ell(a)\) be the average tilted log evidence, and let \(\pi_a^m\) be the actual normalized posterior probability laws. Set



\[\mathcal F(a,q)=-\ell(a)+\frac1M\sum_{m=1}^M\mathrm{KL}(q^m\|\pi_a^m),.\]



For every datum require \(\mathrm{KL}(q^m\|\pi_a^m)<\infty\) before conversion to real numbers.



KL here is the actual measure-theoretic divergence. The density-normalization result below proves the log-density integral identity under explicit normalization and entropy-integrability conditions. Those conditions are not inferred from pointwise smoothness.

\begin{leanbox}{Lean: \texttt{PosteriorFamily}, \texttt{Admissible}, \texttt{freeEnergy} and 1 more}
\leansource{\texttt{PosteriorFamily} \quad LatentDiffusion/\allowbreak{}Variational.lean, lines 22–24}
\begin{leancode}
\leanline{0}{structure\ PosteriorFamily\ (A\ Z\ :\ Type*)\ [MeasurableSpace\ Z]\ (M\ :\ ℕ)\ where}
\leanline{2}{logEvidence\ :\ A\ →\ ℝ}
\leanline{2}{posterior\ :\ A\ →\ Fin\ M\ →\ ProbabilityMeasure\ Z}
\end{leancode}
\leansource{\texttt{Admissible} \quad LatentDiffusion/\allowbreak{}Variational.lean, lines 30–31}
\begin{leancode}
\leanline{0}{def\ Admissible\ (a\ :\ A)\ (q\ :\ Fin\ M\ →\ ProbabilityMeasure\ Z)\ :\ Prop\ :=}
\leanline{2}{∀\ m,\ InformationTheory.klDiv\ (q\ m\ :\ Measure\ Z)\ (P.posterior\ a\ m\ :\ Measure\ Z)\ ≠\ ⊤}
\end{leancode}
\leansource{\texttt{freeEnergy} \quad LatentDiffusion/\allowbreak{}Variational.lean, lines 35–37}
\begin{leancode}
\leanline{0}{def\ freeEnergy\ (a\ :\ A)\ (q\ :\ Fin\ M\ →\ ProbabilityMeasure\ Z)\ :\ ℝ\ :=}
\leanline{2}{-P.logEvidence\ a\ +\ (M\ :\ ℝ)⁻¹\ *}
\leanline{4}{∑\ m,\ (InformationTheory.klDiv\ (q\ m\ :\ Measure\ Z)\ (P.posterior\ a\ m\ :\ Measure\ Z)).toReal}
\end{leancode}
\leansource{\texttt{IsGlobalMin} \quad LatentDiffusion/\allowbreak{}Variational.lean, lines 39–40}
\begin{leancode}
\leanline{0}{def\ IsGlobalMin\ (a\ :\ A)\ (q\ :\ Fin\ M\ →\ ProbabilityMeasure\ Z)\ :\ Prop\ :=}
\leanline{2}{P.Admissible\ a\ q\ ∧\ ∀\ b\ r,\ P.Admissible\ b\ r\ →\ P.freeEnergy\ a\ q\ ≤\ P.freeEnergy\ b\ r}
\end{leancode}
\end{leanbox}
\subsection{\texorpdfstring{Classical flow — Extra assumptions and equations with proved entropy calculus}{Classical flow — Extra assumptions and equations with proved entropy calculus}}



The orange \(R\) extra assumptions specify positive density paths, their actual time derivatives, common-neighbourhood integral domination, regular transport charts and integrable spatial tails. The parameter and density equations are



\[\dot a_t=-\frac1M\sum_m\int\nabla_a U_m(a_t,z)q_t^m(z)\,dz,\qquad\partial_tq_t^m=\nabla\!\cdot\!\bigl(q_t^m\nabla_z(\log q_t^m+U_m(a_t,\cdot))\bigr).\]

\begin{leanbox}{Lean: \texttt{ClassicalFlow}, \texttt{deBruijn}, \texttt{extendedLSI} and 4 more}
\leansource{\texttt{ClassicalFlow} \quad LatentDiffusion/\allowbreak{}RegularConvergence.lean, lines 170–201}
\begin{leancode}
\leanline{0}{structure\ ClassicalFlow\ (M\ d\ :\ ℕ)\ (rho\ :\ ℝ)\ (astar\ :\ A)}
\leanline{4}{(U\ :\ Fin\ M\ →\ A\ →\ Z\ →\ ℝ)\ where}
\leanline{2}{parameter\ :\ ℝ\ →\ A}
\leanline{2}{parameterDerivative\ :\ ℝ\ →\ A}
\leanline{2}{density\ :\ Fin\ M\ →\ ℝ\ →\ Z\ →\ ℝ}
\leanline{2}{densityDerivative\ :\ Fin\ M\ →\ ℝ\ →\ Z\ →\ ℝ}
\leanline{2}{targetDensity\ :\ Fin\ M\ →\ Z\ →\ ℝ}
\leanline{2}{transport\ :\ (t\ :\ ℝ)\ →\ DatasetRegularity\ (Z\ :=\ Z)\ M\ d\ (parameter\ t)\ astar\ rho}
\leanline{2}{potential\_eq\ :\ ∀\ t,\ (transport\ t).potential\ =\ U}
\leanline{2}{density\_eq\ :\ ∀\ t\ i,\ ((transport\ t).displacement\ i).density\ 0\ =\ density\ i\ t}
\leanline{2}{target\_density\_eq\ :\ ∀\ t\ i,\ ((transport\ t).displacement\ i).density\ 1\ =\ targetDensity\ i}
\leanline{2}{parameter\_differentiable\ :\ ∀\ t,\ HasDerivAt\ parameter\ (parameterDerivative\ t)\ t}
\leanline{2}{potential\_parameter\_differentiable\ :\ ∀\ i\ z\ a,\ DifferentiableAt\ ℝ\ (fun\ b\ ↦\ U\ i\ b\ z)\ a}
\leanline{2}{positive\ :\ ∀\ i\ t\ z,\ 0\ <\ density\ i\ t\ z}
\leanline{2}{density\_C1\ :\ ∀\ i\ t,\ ContDiff\ ℝ\ 1\ (density\ i\ t)}
\leanline{2}{potential\_C1\ :\ ∀\ i\ t,\ ContDiff\ ℝ\ 1\ (U\ i\ (parameter\ t))}
\leanline{2}{parameter\_gradient\ :\ ∀\ t\ i\ z,\ HasGradientAt\ (fun\ a\ ↦\ U\ i\ a\ z)}
\leanline{4}{(((transport\ t).potential\_gradient\ i\ z).fst)\ (parameter\ t)}
\leanline{2}{latent\_gradient\ :\ ∀\ t\ i\ z,\ HasGradientAt\ (U\ i\ (parameter\ t))}
\leanline{4}{(((transport\ t).potential\_gradient\ i\ z).snd)\ z}
\leanline{2}{parameter\_gradient\_integrable\ :\ ∀\ t\ i,}
\leanline{4}{Integrable\ (fun\ z\ ↦\ density\ i\ t\ z\ •\ gradient\ (fun\ a\ ↦\ U\ i\ a\ z)\ (parameter\ t))\ volume}
\leanline{2}{domination\ :\ ∀\ t\ i,\ LocalEnergyDomination\ volume\ (density\ i)\ (densityDerivative\ i)}
\leanline{4}{(fun\ u\ ↦\ U\ i\ (parameter\ u))}
\leanline{4}{(fun\ u\ z\ ↦\ fderiv\ ℝ\ (fun\ a\ ↦\ U\ i\ a\ z)\ (parameter\ u)\ (parameterDerivative\ u))\ t}
\leanline{2}{parameter\_equation\ :\ ∀\ t,\ 0\ ≤\ t\ →\ parameterDerivative\ t\ =}
\leanline{4}{-averagedParameterForce\ volume\ (fun\ i\ ↦\ density\ i\ t)\ U\ (parameter\ t)}
\leanline{2}{density\_equation\ :\ ∀\ t,\ 0\ ≤\ t\ →\ ∀\ i,\ ∀ᵐ\ z\ ∂volume,\ densityDerivative\ i\ t\ z\ =}
\leanline{4}{fluxDivergence\ (fokkerPlanckFlux\ (density\ i\ t)\ (U\ i\ (parameter\ t)))\ z}
\leanline{2}{integrable\_tails\ :\ ∀\ t\ i,\ NoBoundaryFlux}
\leanline{4}{(fokkerPlanckFlux\ (density\ i\ t)\ (U\ i\ (parameter\ t)))}
\leanline{4}{(fun\ z\ ↦\ chemicalPotential\ (density\ i\ t)\ (U\ i\ (parameter\ t))\ z\ +\ 1)}
\end{leancode}
\leansource{\texttt{deBruijn} \quad LatentDiffusion/\allowbreak{}RegularConvergence.lean, lines 349–355}
\begin{leancode}
\leanline{0}{theorem\ deBruijn\ (t\ :\ ℝ)\ (ht\ :\ 0\ ≤\ t)\ :\ HasDerivAt\ S.energy\ (-S.fisher\ t)\ t}
\end{leancode}
\leansource{\texttt{extendedLSI} \quad LatentDiffusion/\allowbreak{}RegularConvergence.lean, lines 359–362}
\begin{leancode}
\leanline{0}{theorem\ extendedLSI\ (hM\ :\ 0\ <\ M)\ (hrho\ :\ 0\ <\ rho)\ (t\ :\ ℝ)\ :}
\leanline{4}{2*rho*(S.energy\ t-(S.transport\ t).freeEnergy\ 1)\ ≤\ S.fisher\ t}
\end{leancode}
\leansource{\texttt{energy\_\allowbreak{}decay} \quad LatentDiffusion/\allowbreak{}RegularConvergence.lean, lines 368–381}
\begin{leancode}
\leanline{0}{theorem\ energy\_decay\ (hM\ :\ 0\ <\ M)\ (hrho\ :\ 0\ <\ rho)\ (T\ :\ ℝ)}
\leanline{4}{(t\ :\ ℝ)\ (ht\ :\ t\ ∈\ Icc\ 0\ T)\ :}
\leanline{4}{S.energy\ t-S.targetEnergy\ ≤\ (S.energy\ 0-S.targetEnergy)*Real.exp\ (-2*rho*t)}
\end{leancode}
\leansource{\texttt{RegularDensityDynamics} \quad LatentDiffusion/\allowbreak{}RegularFlowTalagrand.lean, lines 36–57}
\begin{leancode}
\leanline{0}{structure\ RegularDensityDynamics\ (M\ :\ ℕ)\ (U\ :\ Fin\ M\ →\ A\ →\ Z\ →\ ℝ)\ where}
\leanline{2}{parameter\ :\ ℝ\ →\ A}
\leanline{2}{parameterDerivative\ :\ ℝ\ →\ A}
\leanline{2}{density\ :\ Fin\ M\ →\ ℝ\ →\ Z\ →\ ℝ}
\leanline{2}{densityDerivative\ :\ Fin\ M\ →\ ℝ\ →\ Z\ →\ ℝ}
\leanline{2}{parameter\_differentiable\ :\ ∀\ t,\ HasDerivAt\ parameter\ (parameterDerivative\ t)\ t}
\leanline{2}{potential\_parameter\_differentiable\ :\ ∀\ i\ z\ a,\ DifferentiableAt\ ℝ\ (fun\ b\ ↦\ U\ i\ b\ z)\ a}
\leanline{2}{positive\ :\ ∀\ i\ t\ z,\ 0\ <\ density\ i\ t\ z}
\leanline{2}{density\_C1\ :\ ∀\ i\ t,\ ContDiff\ ℝ\ 1\ (density\ i\ t)}
\leanline{2}{potential\_C1\ :\ ∀\ i\ t,\ ContDiff\ ℝ\ 1\ (U\ i\ (parameter\ t))}
\leanline{2}{parameter\_gradient\_integrable\ :\ ∀\ t\ i,}
\leanline{4}{Integrable\ (fun\ z\ ↦\ density\ i\ t\ z\ •\ gradient\ (fun\ a\ ↦\ U\ i\ a\ z)\ (parameter\ t))\ volume}
\leanline{2}{domination\ :\ ∀\ t\ i,\ LocalEnergyDomination\ volume\ (density\ i)\ (densityDerivative\ i)}
\leanline{4}{(fun\ u\ ↦\ U\ i\ (parameter\ u))}
\leanline{4}{(fun\ u\ z\ ↦\ fderiv\ ℝ\ (fun\ a\ ↦\ U\ i\ a\ z)\ (parameter\ u)\ (parameterDerivative\ u))\ t}
\leanline{2}{parameter\_equation\ :\ ∀\ t,\ 0\ ≤\ t\ →\ parameterDerivative\ t\ =}
\leanline{4}{-averagedParameterForce\ volume\ (fun\ i\ ↦\ density\ i\ t)\ U\ (parameter\ t)}
\leanline{2}{density\_equation\ :\ ∀\ t,\ 0\ ≤\ t\ →\ ∀\ i,\ ∀ᵐ\ z\ ∂volume,\ densityDerivative\ i\ t\ z\ =}
\leanline{4}{fluxDivergence\ (fokkerPlanckFlux\ (density\ i\ t)\ (U\ i\ (parameter\ t)))\ z}
\leanline{2}{integrable\_tails\ :\ ∀\ t\ i,\ NoBoundaryFlux}
\leanline{4}{(fokkerPlanckFlux\ (density\ i\ t)\ (U\ i\ (parameter\ t)))}
\leanline{4}{(fun\ z\ ↦\ chemicalPotential\ (density\ i\ t)\ (U\ i\ (parameter\ t))\ z\ +\ 1)}
\end{leancode}
\leansource{\texttt{CharacteristicRealization} \quad LatentDiffusion/\allowbreak{}RegularFlowTalagrand.lean, lines 114–127}
\begin{leancode}
\leanline{0}{structure\ CharacteristicRealization}
\leanline{4}{(D\ :\ ProductDensityChart\ X\ A\ Z\ M\ U)\ (S\ :\ RegularDensityDynamics\ M\ U)\ where}
\leanline{2}{state\ :\ ℝ\ →\ X}
\leanline{2}{state\_continuous\ :\ ContinuousOn\ state\ (Ici\ 0)}
\leanline{2}{parameter\_eq\ :\ ∀\ t,\ D.parameter\ (state\ t)\ =\ S.parameter\ t}
\leanline{2}{density\_eq\ :\ ∀\ t\ i,\ (D.chart\ i).density\ (state\ t)\ =\ S.density\ i\ t}
\leanline{2}{characteristic\ :\ ℝ\ →\ ∀\ i,\ Lp\ Z\ 2\ (Q\ i)}
\leanline{2}{velocity\ :\ ℝ\ →\ ∀\ i,\ Lp\ Z\ 2\ (Q\ i)}
\leanline{2}{law\_eq\ :\ ∀\ t\ i,\ (Q\ i).map\ (characteristic\ t\ i)\ =\ (D.law\ (state\ t)\ i\ :\ Measure\ Z)}
\leanline{2}{derivative\ :\ ∀\ t\ i,\ HasDerivAt\ (fun\ u\ ↦\ characteristic\ u\ i)\ (velocity\ t\ i)\ t}
\leanline{2}{velocity\_eq\ :\ ∀\ t\ i,\ (velocity\ t\ i\ :\ Ω\ i\ →\ Z)\ =ᵐ[Q\ i]}
\leanline{4}{(fun\ ω\ ↦\ -(S.score\ t\ i\ (characteristic\ t\ i\ ω)))}
\leanline{2}{anchor\ :\ ∀\ s\ t\ i,\ RegularMetricFlow.OptimalAnchor\ (Q\ i)\ (characteristic\ t\ i)}
\leanline{4}{(D.law\ (state\ s)\ i\ :\ Measure\ Z)}
\end{leancode}
\leansource{\texttt{exists\_\allowbreak{}minimizer\_\allowbreak{}talagrand} \quad LatentDiffusion/\allowbreak{}RegularFlowTalagrand.lean, lines 196–217}
\begin{leancode}
\leanline{0}{theorem\ exists\_minimizer\_talagrand\ [CompleteSpace\ X]}
\leanline{4}{(Fstar\ rho\ :\ ℝ)\ (hrho\ :\ 0\ <\ rho)\ (hstar\ :\ ∀\ x,\ Fstar\ ≤\ D.energy\ x)}
\leanline{4}{(hxLSI\ :\ ∀\ x,\ 2\ *\ rho\ *\ (D.energy\ x-Fstar)\ ≤}
\leanline{6}{normalizedSlopeSquared\ volume\ (fun\ i\ ↦\ (D.chart\ i).density\ x)\ U\ (D.parameter\ x))\ :}
\leanline{4}{∃\ xstar,\ Tendsto\ (fun\ n\ :\ ℕ\ ↦\ R.state\ n)\ atTop\ ({\leanfallback 𝓝}\ xstar)\ ∧}
\leanline{6}{D.energy\ xstar\ =\ Fstar\ ∧\ (∀\ x,\ D.energy\ xstar\ ≤\ D.energy\ x)\ ∧}
\leanline{6}{∀\ t,\ 0\ ≤\ t\ →\ rho\ *\ dist\ (R.state\ t)\ xstar\ \textasciicircum{}\ 2\ ≤}
\leanline{8}{2\ *\ (D.energy\ (R.state\ t)-Fstar)}
\end{leancode}
\leanlabel{How the Lean proof follows the argument}
\texttt{Classical\allowbreak{}Flow} records these primitive equations and extra-assumption data. The displayed theorems derive the actual energy derivative, normalized extended LSI and exponential energy decay. For the source xLSI-to-Talagrand implication, the separate bare \texttt{Regular\allowbreak{}Density\allowbreak{}Dynamics} omits A3 entirely: actual characteristics and coupling anchors derive its metric speed, and xLSI is the sole functional-inequality premise. Dissipation and speed are conclusions of the connected theorem.
\end{leanbox}
\subsection{\texorpdfstring{Density bridge — The actual log-density integral equals KL minus log evidence}{Density bridge — The actual log-density integral equals KL minus log evidence}}



Let \(\rho>0\) be measurable and integrable relative to \(\nu\), with \(Z=\int\rho\,\mathrm d\nu>0\). Put \(\pi=(\rho/Z)\nu\). Let \(r\ge0\) be measurable and integrable under \(\pi\), with \(\int r\,\mathrm d\pi=1\) and \(r\log r\in L^1(\pi)\). Set \(q_\nu=r\rho/Z\) and \(q=q_\nu\nu\).



\[\int q_\nu\log\!\left(\frac{q_\nu}{\rho}\right)\,\mathrm d\nu=\int\log\!\left(\frac{q_\nu}{\rho}\right)\,\mathrm dq=\mathrm{KL}(q\|\pi)-\log Z.\]



The theorem constructs both actual probability measures, proves that KL is finite, and proves integrability of the log ratio and this identity. It permits the relative density to vanish. The regular variational module additionally realizes every finite-KL trial law by its actual Radon–Nikodym density.

\begin{leanbox}{Lean: \texttt{density\_\allowbreak{}freeEnergy\_\allowbreak{}eq\_\allowbreak{}kl}, \texttt{density\_\allowbreak{}weighted\_\allowbreak{}freeEnergy\_\allowbreak{}eq\_\allowbreak{}kl}}
\leansource{\texttt{density\_\allowbreak{}freeEnergy\_\allowbreak{}eq\_\allowbreak{}kl} \quad LatentDiffusion/\allowbreak{}Density.lean, lines 117–141}
\begin{leancode}
\leanline{0}{theorem\ density\_freeEnergy\_eq\_kl\ (ν\ :\ Measure\ E)\ (ρ\ r\ :\ E\ →\ ℝ)\ (Z\ :\ ℝ)}
\leanline{4}{(hρmeas\ :\ Measurable\ ρ)\ (hrmeas\ :\ Measurable\ r)}
\leanline{4}{(hρ\ :\ ∀\ x,\ 0\ <\ ρ\ x)\ (hr\ :\ ∀\ x,\ 0\ ≤\ r\ x)\ (hZ\ :\ 0\ <\ Z)}
\leanline{4}{(hρint\ :\ Integrable\ ρ\ ν)\ (hZeq\ :\ (∫\ x,\ ρ\ x\ ∂ν)\ =\ Z)}
\leanline{4}{(hrint\ :\ Integrable\ r\ (posteriorMeasure\ ν\ ρ\ Z))}
\leanline{4}{(hrmass\ :\ (∫\ x,\ r\ x\ ∂posteriorMeasure\ ν\ ρ\ Z)\ =\ 1)}
\leanline{4}{(hentropy\ :\ Integrable\ (fun\ x\ ↦\ r\ x\ *\ Real.log\ (r\ x))\ (posteriorMeasure\ ν\ ρ\ Z))\ :}
\leanline{4}{let\ π\ :=\ posteriorMeasure\ ν\ ρ\ Z}
\leanline{4}{let\ q\ :=\ densityMeasure\ ν\ (trialDensity\ ρ\ r\ Z)}
\leanline{4}{IsProbabilityMeasure\ π\ ∧\ IsProbabilityMeasure\ q\ ∧}
\leanline{6}{InformationTheory.klDiv\ q\ π\ ≠\ ⊤\ ∧}
\leanline{6}{Integrable\ (fun\ x\ ↦\ Real.log\ (trialDensity\ ρ\ r\ Z\ x\ /\ ρ\ x))\ q\ ∧}
\leanline{6}{(∫\ x,\ Real.log\ (trialDensity\ ρ\ r\ Z\ x\ /\ ρ\ x)\ ∂q)\ =}
\leanline{8}{(InformationTheory.klDiv\ q\ π).toReal\ -\ Real.log\ Z}
\end{leancode}
\leansource{\texttt{density\_\allowbreak{}weighted\_\allowbreak{}freeEnergy\_\allowbreak{}eq\_\allowbreak{}kl} \quad LatentDiffusion/\allowbreak{}Density.lean, lines 146–176}
\begin{leancode}
\leanline{0}{theorem\ density\_weighted\_freeEnergy\_eq\_kl\ (ν\ :\ Measure\ E)\ (ρ\ r\ :\ E\ →\ ℝ)\ (Z\ :\ ℝ)}
\leanline{4}{(hρmeas\ :\ Measurable\ ρ)\ (hrmeas\ :\ Measurable\ r)}
\leanline{4}{(hρ\ :\ ∀\ x,\ 0\ <\ ρ\ x)\ (hr\ :\ ∀\ x,\ 0\ ≤\ r\ x)\ (hZ\ :\ 0\ <\ Z)}
\leanline{4}{(hρint\ :\ Integrable\ ρ\ ν)\ (hZeq\ :\ (∫\ x,\ ρ\ x\ ∂ν)\ =\ Z)}
\leanline{4}{(hrint\ :\ Integrable\ r\ (posteriorMeasure\ ν\ ρ\ Z))}
\leanline{4}{(hrmass\ :\ (∫\ x,\ r\ x\ ∂posteriorMeasure\ ν\ ρ\ Z)\ =\ 1)}
\leanline{4}{(hentropy\ :\ Integrable\ (fun\ x\ ↦\ r\ x\ *\ Real.log\ (r\ x))\ (posteriorMeasure\ ν\ ρ\ Z))\ :}
\leanline{4}{let\ π\ :=\ posteriorMeasure\ ν\ ρ\ Z}
\leanline{4}{let\ f\ :=\ trialDensity\ ρ\ r\ Z}
\leanline{4}{let\ q\ :=\ densityMeasure\ ν\ f}
\leanline{4}{InformationTheory.klDiv\ q\ π\ ≠\ ⊤\ ∧}
\leanline{6}{Integrable\ (fun\ x\ ↦\ f\ x\ *\ Real.log\ (f\ x\ /\ ρ\ x))\ ν\ ∧}
\leanline{6}{(∫\ x,\ f\ x\ *\ Real.log\ (f\ x\ /\ ρ\ x)\ ∂ν)\ =}
\leanline{8}{(InformationTheory.klDiv\ q\ π).toReal\ -\ Real.log\ Z}
\end{leancode}
\end{leanbox}
\subsection{\texorpdfstring{Source objective — Tilting adds the expected penalty}{Source objective — Tilting adds the expected penalty}}



For measurable \(\rho\ge0\) and \(R\ge0\), define \(\widetilde\rho=\rho e^{-R}\). If \(\rho\) is integrable, then \(\widetilde\rho\) is integrable and \(\int\widetilde\rho\le\int\rho\). When the tilted normalizer is positive, \(\log\int\widetilde\rho\le\log\int\rho\).




For \(q\ge0\), \(\rho>0\), with both terms on the right integrable, the literal density free energy satisfies



\[\int q\log\frac{q}{\rho e^{-R}}\,d\nu=\int q\log\frac q\rho\,d\nu+\int qR\,d\nu\ \ge\ \int q\log\frac q\rho\,d\nu.\]



This proves the source's density-level penalty decomposition and evidence order. The identification of \(R\) with a conditional path-space KL and the path-law chain rule remain separate model obligations. No differentiation of integrals is used in this identity.

\begin{leanbox}{Lean: \texttt{tilted\_\allowbreak{}freeEnergy\_\allowbreak{}decomposition}, \texttt{tilted\_\allowbreak{}freeEnergy\_\allowbreak{}le}, \texttt{tilted\_\allowbreak{}evidence\_\allowbreak{}le} and 1 more}
\leansource{\texttt{tilted\_\allowbreak{}freeEnergy\_\allowbreak{}decomposition} \quad LatentDiffusion/\allowbreak{}SourceResults.lean, lines 158–168}
\begin{leancode}
\leanline{0}{theorem\ tilted\_freeEnergy\_decomposition}
\leanline{4}{\{Z\ :\ Type*\}\ [MeasurableSpace\ Z]\ (ν\ :\ Measure\ Z)\ (q\ ρ\ R\ :\ Z\ →\ ℝ)}
\leanline{4}{(hq\ :\ ∀\ z,\ 0\ ≤\ q\ z)\ (hρ\ :\ ∀\ z,\ 0\ <\ ρ\ z)}
\leanline{4}{(hfree\ :\ Integrable\ (fun\ z\ ↦\ q\ z\ *\ Real.log\ (q\ z\ /\ ρ\ z))\ ν)}
\leanline{4}{(hpenalty\ :\ Integrable\ (fun\ z\ ↦\ q\ z\ *\ R\ z)\ ν)\ :}
\leanline{4}{Integrable\ (fun\ z\ ↦\ q\ z\ *\ Real.log\ (q\ z\ /\ tiltedDensity\ ρ\ R\ z))\ ν\ ∧}
\leanline{6}{densityFreeEnergy\ ν\ q\ (tiltedDensity\ ρ\ R)\ =}
\leanline{8}{densityFreeEnergy\ ν\ q\ ρ\ +\ ∫\ z,\ q\ z\ *\ R\ z\ ∂ν}
\end{leancode}
\leansource{\texttt{tilted\_\allowbreak{}freeEnergy\_\allowbreak{}le} \quad LatentDiffusion/\allowbreak{}SourceResults.lean, lines 172–179}
\begin{leancode}
\leanline{0}{theorem\ tilted\_freeEnergy\_le}
\leanline{4}{\{Z\ :\ Type*\}\ [MeasurableSpace\ Z]\ (ν\ :\ Measure\ Z)\ (q\ ρ\ R\ :\ Z\ →\ ℝ)}
\leanline{4}{(hq\ :\ ∀\ z,\ 0\ ≤\ q\ z)\ (hρ\ :\ ∀\ z,\ 0\ <\ ρ\ z)\ (hR\ :\ ∀\ z,\ 0\ ≤\ R\ z)}
\leanline{4}{(hfree\ :\ Integrable\ (fun\ z\ ↦\ q\ z\ *\ Real.log\ (q\ z\ /\ ρ\ z))\ ν)}
\leanline{4}{(hpenalty\ :\ Integrable\ (fun\ z\ ↦\ q\ z\ *\ R\ z)\ ν)\ :}
\leanline{4}{densityFreeEnergy\ ν\ q\ ρ\ ≤\ densityFreeEnergy\ ν\ q\ (tiltedDensity\ ρ\ R)}
\end{leancode}
\leansource{\texttt{tilted\_\allowbreak{}evidence\_\allowbreak{}le} \quad LatentDiffusion/\allowbreak{}SourceResults.lean, lines 184–200}
\begin{leancode}
\leanline{0}{theorem\ tilted\_evidence\_le}
\leanline{4}{\{Z\ :\ Type*\}\ [MeasurableSpace\ Z]\ (ν\ :\ Measure\ Z)\ (ρ\ R\ :\ Z\ →\ ℝ)}
\leanline{4}{(hρmeas\ :\ Measurable\ ρ)\ (hRmeas\ :\ Measurable\ R)}
\leanline{4}{(hρ\ :\ ∀\ z,\ 0\ ≤\ ρ\ z)\ (hR\ :\ ∀\ z,\ 0\ ≤\ R\ z)\ (hint\ :\ Integrable\ ρ\ ν)\ :}
\leanline{4}{Integrable\ (tiltedDensity\ ρ\ R)\ ν\ ∧}
\leanline{6}{(∫\ z,\ tiltedDensity\ ρ\ R\ z\ ∂ν)\ ≤\ ∫\ z,\ ρ\ z\ ∂ν}
\end{leancode}
\leansource{\texttt{tilted\_\allowbreak{}logEvidence\_\allowbreak{}le} \quad LatentDiffusion/\allowbreak{}SourceResults.lean, lines 204–210}
\begin{leancode}
\leanline{0}{theorem\ tilted\_logEvidence\_le}
\leanline{4}{\{Z\ :\ Type*\}\ [MeasurableSpace\ Z]\ (ν\ :\ Measure\ Z)\ (ρ\ R\ :\ Z\ →\ ℝ)}
\leanline{4}{(hρmeas\ :\ Measurable\ ρ)\ (hRmeas\ :\ Measurable\ R)}
\leanline{4}{(hρ\ :\ ∀\ z,\ 0\ ≤\ ρ\ z)\ (hR\ :\ ∀\ z,\ 0\ ≤\ R\ z)\ (hint\ :\ Integrable\ ρ\ ν)}
\leanline{4}{(hpositive\ :\ 0\ <\ ∫\ z,\ tiltedDensity\ ρ\ R\ z\ ∂ν)\ :}
\leanline{4}{Real.log\ (∫\ z,\ tiltedDensity\ ρ\ R\ z\ ∂ν)\ ≤\ Real.log\ (∫\ z,\ ρ\ z\ ∂ν)}
\end{leancode}
\end{leanbox}
\section{Results of the paper in corrected form}
\subsection{\texorpdfstring{Variational characterization of the tilted likelihood}{Variational characterization of the tilted likelihood}}


Let \(M\ge1\), let \(\rho_u^m(z)=\widetilde p_u(x^m,z)\),
        and suppose \(0\lt A_u^m=\int\rho_u^m\,dz\lt\infty\). Set
        \(\pi_u^m=\rho_u^m/A_u^m\) and
        \(\widetilde\ell(u)=M^{-1}\sum_m\log A_u^m\). On the domain where the
        density integrals and relative entropies are finite, and the normalized
        posteriors are admissible comparison laws,


        \[\widetilde F(u,q)=-\widetilde\ell(u)
          +\frac1M\sum_m\mathrm{KL}(q^m\Vert\pi_u^m).\]


Consequently \((u_\star,q_\star)\) is a global minimizer exactly when
        \(u_\star\) globally maximizes \(\widetilde\ell\) and
        \(q_\star^m=\pi_{u_\star}^m\) for every datum. In particular \(u\)
        maximizes \(\widetilde\ell\) if and only if it is the parameter
        component of some global variational minimizer.

\begin{leanbox}{Lean: \texttt{DensityModel.\allowbreak{}literal\_\allowbreak{}minimizers\_\allowbreak{}on\_\allowbreak{}iff}, \texttt{existsUnique\_\allowbreak{}p2\_\allowbreak{}literalOptimizer}}
\leansource{\texttt{DensityModel.\allowbreak{}literal\_\allowbreak{}minimizers\_\allowbreak{}on\_\allowbreak{}iff} \quad LatentDiffusion/\allowbreak{}RegularVariational.lean, lines 148–176}
\begin{leancode}
\leanline{0}{theorem\ DensityModel.literal\_minimizers\_on\_iff\ (D\ :\ DensityModel\ A\ Z\ M\ ν)\ (hM\ :\ 0\ <\ M)}
\leanline{4}{(domain\ :\ (Fin\ M\ →\ ProbabilityMeasure\ Z)\ →\ Prop)}
\leanline{4}{(hposterior\ :\ ∀\ a,\ domain\ (D.posterior\ a))\ (a\ :\ A)\ (q\ :\ Fin\ M\ →\ ProbabilityMeasure\ Z)\ :}
\leanline{4}{(D.family.Admissible\ a\ q\ ∧\ domain\ q\ ∧\ ∀\ b\ r,\ D.family.Admissible\ b\ r\ →\ domain\ r\ →}
\leanline{6}{D.literalEnergy\ a\ q\ ≤\ D.literalEnergy\ b\ r)\ ↔}
\leanline{4}{(∀\ b,\ D.family.logEvidence\ b\ ≤\ D.family.logEvidence\ a)\ ∧\ q\ =\ D.posterior\ a}
\end{leancode}
\leansource{\texttt{existsUnique\_\allowbreak{}p2\_\allowbreak{}literalOptimizer} \quad LatentDiffusion/\allowbreak{}RegularModel.lean, lines 108–133}
\begin{leancode}
\leanline{0}{theorem\ existsUnique\_p2\_literalOptimizer\ (U\ :\ Fin\ M\ →\ IPLA.Joint\ dθ\ dz\ →\ ℝ)}
\leanline{4}{(hM\ :\ 0\ <\ M)\ (m\ :\ ℝ)\ (hm\ :\ 0\ <\ m)}
\leanline{4}{(hdiff\ :\ ∀\ i,\ Differentiable\ ℝ\ (U\ i))}
\leanline{4}{(hsc\ :\ ∀\ i,\ StrongConvexOn\ univ\ m\ (U\ i))\ :}
\leanline{4}{let\ D\ :=\ densityModel\ U\ m\ hm\ hdiff\ hsc}
\leanline{4}{∃!\ p\ :\ EuclideanSpace\ ℝ\ (Fin\ dθ)\ ×}
\leanline{8}{(Fin\ M\ →\ ProbabilityMeasure\ (EuclideanSpace\ ℝ\ (Fin\ dz))),}
\leanline{6}{D.family.Admissible\ p.1\ p.2\ ∧}
\leanline{6}{(∀\ i,\ MemLp\ (fun\ z\ :\ EuclideanSpace\ ℝ\ (Fin\ dz)\ ↦\ z)\ 2\ (p.2\ i\ :\ Measure\ \_))\ ∧}
\leanline{6}{∀\ b\ q,\ D.family.Admissible\ b\ q\ →}
\leanline{8}{(∀\ i,\ MemLp\ (fun\ z\ :\ EuclideanSpace\ ℝ\ (Fin\ dz)\ ↦\ z)\ 2\ (q\ i\ :\ Measure\ \_))\ →}
\leanline{8}{D.literalEnergy\ p.1\ p.2\ ≤\ D.literalEnergy\ b\ q}
\end{leancode}
\leanlabel{How the Lean proof follows the argument}
\textbf{Correction.} The normalizer integrates the tilted density over
        \(z=z_0\), the latent variable on which it was defined in
        \texttt{eq:ptildedef} (main, lines 100–102). The occurrence of
        \(z_{0:K}\) in the proposition's normalizer is inconsistent with that
        definition. Positivity, finiteness, admissibility and integrability are
        explicit here. On the paper's \(\mathcal P_2\) domain the normalized
        posteriors must in particular have finite second moments; positive
        finite normalizers alone do not guarantee this. Arbitrary extended-integral conventions are not silently
        converted to real arithmetic. This proposition concerns the tilted
        likelihood. The subsequent inference about the original likelihood
        requires the separate zero-tilt correction below.
\end{leanbox}
\subsection{\texorpdfstring{Unique tilted optimizer and exponential parameter convergence}{Unique tilted optimizer and exponential parameter convergence}}


Let \(M\ge1\), and let the actual tilted potentials \(U_m=-\log\rho^m\) be differentiable and satisfy A3. Under the explicitly stated extra assumptions R1–R4, the average tilted evidence has a unique maximizer \(u_\star=(\theta_\star,\phi_\star)\). The classical parameter/\allowbreak{}density equations satisfy


 \[\|u_t-u_\star\|\le\sqrt{\frac{2[\widetilde F(u_0,q_0)-\widetilde F_\star]}{\lambda}}e^{-\lambda t},\qquad t\ge0.\]


The potential normalizers and finite-second-moment posteriors are constructed from their actual integrals. Energy dissipation, the entropy transport tangent, extended LSI and the decay estimate are proved from the equations and the orange extra assumptions.



\textbf{Correction.} A2 must be interpreted on an admissible initial class,
        or require smooth densities only for positive time. Its literal
        quantification over all \(\mathcal P_2\) initial laws is contradicted by
        Dirac laws. Entropy differentiation and transport control need the
        explicit analytic hypotheses above. The averaged joint potential is
        strongly convex in the normalized product norm
        \(\|u\|^2+M^{-1}\sum_m\|z^m\|^2\); the sketch at main lines 168–174
        must use that norm when retaining the constant \(\lambda\).

\begin{leanbox}{Lean: \texttt{existsUnique\_\allowbreak{}variationalOptimizer}, \texttt{existsUnique\_\allowbreak{}optimizer\_\allowbreak{}and\_\allowbreak{}flow\_\allowbreak{}convergence}}
\leansource{\texttt{existsUnique\_\allowbreak{}variationalOptimizer} \quad LatentDiffusion/\allowbreak{}RegularModel.lean, lines 66–74}
\begin{leancode}
\leanline{0}{theorem\ existsUnique\_variationalOptimizer\ (U\ :\ Fin\ M\ →\ IPLA.Joint\ dθ\ dz\ →\ ℝ)}
\leanline{4}{(hM\ :\ 0\ <\ M)\ (m\ :\ ℝ)\ (hm\ :\ 0\ <\ m)}
\leanline{4}{(hdiff\ :\ ∀\ i,\ Differentiable\ ℝ\ (U\ i))}
\leanline{4}{(hsc\ :\ ∀\ i,\ StrongConvexOn\ univ\ m\ (U\ i))\ :}
\leanline{4}{∃!\ p\ :\ EuclideanSpace\ ℝ\ (Fin\ dθ)\ ×}
\leanline{8}{(Fin\ M\ →\ ProbabilityMeasure\ (EuclideanSpace\ ℝ\ (Fin\ dz))),}
\leanline{6}{(densityModel\ U\ m\ hm\ hdiff\ hsc).family.IsGlobalMin\ p.1\ p.2}
\end{leancode}
\leansource{\texttt{existsUnique\_\allowbreak{}optimizer\_\allowbreak{}and\_\allowbreak{}flow\_\allowbreak{}convergence} \quad LatentDiffusion/\allowbreak{}RegularConvergence.lean, lines 478–517}
\begin{leancode}
\leanline{0}{theorem\ existsUnique\_optimizer\_and\_flow\_convergence\ \{dθ\ dz\ :\ ℕ\}}
\leanline{4}{(U\ :\ Fin\ M\ →\ IPLA.Joint\ dθ\ dz\ →\ ℝ)\ (hM\ :\ 0\ <\ M)}
\leanline{4}{(rho\ :\ ℝ)\ (hrho\ :\ 0\ <\ rho)}
\leanline{4}{(hdiff\ :\ ∀\ i,\ Differentiable\ ℝ\ (U\ i))}
\leanline{4}{(hsc\ :\ ∀\ i,\ StrongConvexOn\ univ\ rho\ (U\ i))\ :}
\leanline{4}{let\ V\ :=\ fun\ i\ a\ z\ ↦\ U\ i\ (IPLA.pair\ a\ z)}
\leanline{4}{let\ ell\ :=\ fun\ a\ ↦\ (M\ :\ ℝ)⁻¹\ *\ ∑\ i,\ Real.log\ (∫\ z,\ Real.exp\ (-V\ i\ a\ z))}
\leanline{4}{∃!\ astar\ :\ EuclideanSpace\ ℝ\ (Fin\ dθ),}
\leanline{6}{(∀\ a,\ ell\ a\ ≤\ ell\ astar)\ ∧}
\leanline{6}{∀\ S\ :\ ClassicalFlow\ M\ dz\ rho\ astar\ V,}
\leanline{8}{(∀\ i\ z,\ S.targetDensity\ i\ z\ =\ Real.exp\ (-V\ i\ astar\ z)/}
\leanline{10}{(∫\ w,\ Real.exp\ (-V\ i\ astar\ w)))\ →}
\leanline{8}{∀\ t,\ 0\ ≤\ t\ →\ ‖S.parameter\ t-astar‖\ ≤}
\leanline{10}{Real.sqrt\ (2*(S.energy\ 0-S.targetEnergy)/rho)*Real.exp\ (-rho*t)}
\end{leancode}
\end{leanbox}
\subsection{\texorpdfstring{Nonasymptotic error of the full-batch Euler scheme}{Nonasymptotic error of the full-batch Euler scheme}}


Here \(C_{h,N}\) is independent of \(T\), \(C\) is independent of
        \((h,N,T)\), and \(0\lt\rho\lt1\) may depend on \(h\). The checked
        contraction factor is \(\rho=1-\lambda h/2\), with
        \(\rho^T\le e^{-\lambda hT/2}\). Uniform constants require the stated initialization and
        moment controls. With \(K=MN\), latent dimension \(d_z\) and the R7
        per-copy score variance bound \(V_0\), the checked bound has
        \(C=\|X_0-S_0\|_{L^2}\), the normalized initial distance from the
        Euler initialization \(X_0\) to the posterior-initialized reference
        \(S_0\), and
        \(C_{h,N}=\frac2\lambda\sqrt{V_0/K}+\frac L\lambda\bigl(h\sqrt{Ld_z}+\frac43\sqrt{2d_z}\sqrt h\bigr)\).



\textbf{Correction.} The temporal lemma contains an extra square-root
        factor; its corrected general estimate is \(C\sqrt h\), not the
        \(O(h)\) accidentally written there. The empirical comparison needs a
        specified coupling, and the spatial comparison needs compatible
        initialization and uniform moments. This theorem concerns the full
        batch method, as the paper itself states at main line 314. It gives no
        guarantee for Algorithm 1's changed minibatch temperature or the adaptive
        algorithm.

\begin{leanbox}{Lean: \texttt{full\_\allowbreak{}batch\_\allowbreak{}optimizer\_\allowbreak{}rms\_\allowbreak{}rate}, \texttt{optimizer\_\allowbreak{}replication}}
\leansource{\texttt{full\_\allowbreak{}batch\_\allowbreak{}optimizer\_\allowbreak{}rms\_\allowbreak{}rate} \quad LatentDiffusion/\allowbreak{}RegularSourceTheorem.lean, lines 167–209}
\begin{leancode}
\leanline{0}{theorem\ full\_batch\_optimizer\_rms\_rate}
\leanline{4}{\{M\ N\ :\ ℕ\}\ (hMN\ :\ 0\ <\ M*N)}
\leanline{4}{(U\ :\ Fin\ (M*N)\ →\ Cloud.Joint\ dθ\ dz\ →\ ℝ)}
\leanline{4}{(m\ h\ :\ ℝ)\ (L\ :\ ℝ≥0)\ (hm\ :\ 0\ <\ m)}
\leanline{4}{(hdiff\ :\ ∀\ i,\ Differentiable\ ℝ\ (U\ i))}
\leanline{4}{(hsc\ :\ ∀\ i,\ StrongConvexOn\ Set.univ\ m\ (U\ i))}
\leanline{4}{(hLip\ :\ ∀\ i,\ LipschitzWith\ L\ (gradient\ (U\ i)))\ (hs\ :\ StableStep\ m\ L\ h)}
\leanline{4}{(initial\ :\ Lp\ (Latent\ (M*N)\ dz)\ 2\ P)}
\leanline{4}{(hinitial\ :\ P.map\ initial\ =\ ULA.gibbsMeasure}
\leanline{6}{(frozenPotential\ U\ (optimizerParameter\ hMN\ U\ m\ hm\ hdiff\ hsc)))}
\leanline{4}{(hposterior\ :\ ∀\ i,\ P.map\ (fun\ ω\ =>\ physicalParticleCLM\ i}
\leanline{8}{((packLp\ P\ (optimizerParameter\ hMN\ U\ m\ hm\ hdiff\ hsc)\ initial)\ ω))\ =}
\leanline{6}{RegularEvidence.posterior\ volume\ (latentPotential\ U\ i)\ (optimizerParameter\ hMN\ U\ m\ hm\ hdiff\ hsc))}
\leanline{4}{(data\ :\ ∀\ i,\ RegularEvidence.LocalNormalizerRegularity\ volume\ (latentPotential\ U\ i)}
\leanline{6}{(optimizerParameter\ hMN\ U\ m\ hm\ hdiff\ hsc))}
\leanline{4}{(hindScores\ :\ iIndepFun\ (fun\ i\ ω\ =>\ scoreLp\ P\ U\ L\ hLip}
\leanline{6}{(packLp\ P\ (optimizerParameter\ hMN\ U\ m\ hm\ hdiff\ hsc)\ initial)\ i\ ω)\ P)}
\leanline{4}{(B\ :\ (Fin\ (M*N)\ ×\ Fin\ dz)\ →\ ℝ≥0\ →\ Ω\ →\ ℝ)}
\leanline{4}{(hB\ :\ ∀\ i,\ IsPreBrownianReal\ (B\ i)\ P)}
\leanline{4}{(hscalar\ :\ ∀\ s\ :\ ℝ,\ 0\ <\ s\ →\ iIndepFun\ (fun\ p\ :\ ℕ\ ×\ (Fin\ (M*N)\ ×\ Fin\ dz)\ =>\ fun\ ω\ =>}
\leanline{6}{ULA.Hilbert.brownianEulerNoise\ P\ B\ hB\ s\ p.1\ ω\ p.2)\ P)}
\leanline{4}{(hind\ :\ ∀\ s\ :\ ℝ,\ 0\ <\ s\ →\ IndepFun\ initial}
\leanline{6}{(fun\ ω\ n\ =>\ ULA.Hilbert.brownianEulerNoise\ P\ B\ hB\ s\ n\ ω)\ P)}
\leanline{4}{(V₀\ :\ ℝ)}
\leanline{4}{(hvar\ :\ ∀\ i,\ ‖scoreLp\ P\ U\ L\ hLip}
\leanline{8}{(packLp\ P\ (optimizerParameter\ hMN\ U\ m\ hm\ hdiff\ hsc)\ initial)\ i-}
\leanline{6}{constantL2\ P\ (∫\ ω,\ scoreLp\ P\ U\ L\ hLip}
\leanline{8}{(packLp\ P\ (optimizerParameter\ hMN\ U\ m\ hm\ hdiff\ hsc)\ initial)\ i\ ω\ ∂P)‖\textasciicircum{}2\ ≤\ V₀)}
\leanline{4}{(x₀\ :\ Lp\ (Cloud.State\ (M*N)\ dθ\ dz)\ 2\ P)\ (n\ :\ ℕ)\ :}
\leanline{4}{Real.sqrt\ (∫\ ω,\ ‖IPLA.parameterCLM}
\leanline{6}{(eulerTrajectory\ (Cloud.meanForce\ U)\ h}
\leanline{8}{(fun\ k\ =>\ noiseIncrement\ P\ B\ hB\ h\ k\ ω)\ (x₀\ ω)\ n)-}
\leanline{6}{optimizerParameter\ hMN\ U\ m\ hm\ hdiff\ hsc‖\textasciicircum{}2\ ∂P)\ ≤}
\leanline{6}{(1-m*h/2)\textasciicircum{}n\ *\ ‖x₀-packLp\ P\ (optimizerParameter\ hMN\ U\ m\ hm\ hdiff\ hsc)\ initial‖\ +}
\leanline{6}{(2/m)\ *\ Real.sqrt\ (V₀/((M*N\ :\ ℕ)\ :\ ℝ))\ +}
\leanline{6}{(L\ :\ ℝ)/m\ *\ (h*Real.sqrt\ ((L\ :\ ℝ)*(dz\ :\ ℝ))+}
\leanline{8}{(4/3\ :\ ℝ)*Real.sqrt\ (2*(dz\ :\ ℝ))*Real.sqrt\ h)}
\end{leancode}
\leansource{\texttt{optimizer\_\allowbreak{}replication} \quad LatentDiffusion/\allowbreak{}RegularSourceTheorem.lean, lines 250–262}
\begin{leancode}
\leanline{0}{theorem\ optimizer\_replication\ \{M\ N\ :\ ℕ\}\ (hM\ :\ 0\ <\ M)\ (hN\ :\ 0\ <\ N)}
\leanline{4}{(U\ :\ Fin\ M\ →\ Cloud.Joint\ dθ\ dz\ →\ ℝ)\ (m\ :\ ℝ)\ (hm\ :\ 0\ <\ m)}
\leanline{4}{(hdiff\ :\ ∀\ i,\ Differentiable\ ℝ\ (U\ i))}
\leanline{4}{(hsc\ :\ ∀\ i,\ StrongConvexOn\ Set.univ\ m\ (U\ i))\ :}
\leanline{4}{optimizerParameter\ (Nat.mul\_pos\ hM\ hN)\ (replicatePotentials\ (N\ :=\ N)\ U)\ m\ hm}
\leanline{6}{(fun\ i\ =>\ hdiff\ i.divNat)\ (fun\ i\ =>\ hsc\ i.divNat)\ =}
\leanline{8}{optimizerParameter\ hM\ U\ m\ hm\ hdiff\ hsc}
\end{leancode}
\leanlabel{How the Lean proof follows the argument}
For the full-batch exact-gradient Euler–Maruyama
        recursion in \texttt{eq:disc\_\allowbreak{}theta}–\texttt{eq:disc\_\allowbreak{}par}, the source
        assumes A1–A4 and Lipschitz gradients of \(R\). For sufficiently small
        \(h>0\) it states


        \[\bigl(\mathbb E\|u_T^{N,h}-u_\star\|^2\bigr)^{1/2}
          \le C_{h,N}+C\rho^T,\qquad
          C_{h,N}=O(\sqrt h+N^{-1/2}),\quad T\in\mathbb N.\]
\end{leanbox}
\subsection{\texorpdfstring{First variations of the free energy}{First variations of the free energy}}


For one datum, write
        \(V_u(z)=\mathcal L_D(\theta,z)-\log p_\phi(x\mid z)\) and
        \(\widetilde F(u,q)=\int[V_u+\log q]\,q\,dz\). For admissible
        variations \(r\) with \(\int r\,dz=0\), and parameter directions for
        which differentiation may pass through the integrals,


        \[\delta_q\widetilde F=\log q+V_u\pmod{\mathbb R},\qquad
        \nabla_\theta\widetilde F=\int\nabla_\theta\mathcal L_D\,dq,\qquad
        \nabla_\phi\widetilde F=-\int\nabla_\phi\log p_\phi(x\mid z)\,dq.\]


Equivalently \(D_q\widetilde F[r]=\int(\log q+V_u)r\,dz\).
        The derivative of \(q\log q\) contains \(+1\); its integral against a
        probability tangent direction vanishes.



\textbf{Correction.} The density first variation is a cotangent equivalence class modulo constants, as the source's own tangent construction specifies. An integrable dominating derivative or another valid interchange argument is needed for both the parameter Taylor remainder and entropy differentiation.

\begin{leanbox}{Lean: \texttt{parameter\_\allowbreak{}freeEnergy\_\allowbreak{}hasFDerivAt}, \texttt{density\_\allowbreak{}firstVariation\_\allowbreak{}of\_\allowbreak{}zero\_\allowbreak{}mass}}
\leansource{\texttt{parameter\_\allowbreak{}freeEnergy\_\allowbreak{}hasFDerivAt} \quad LatentDiffusion/\allowbreak{}SourceCalculus.lean, lines 44–67}
\begin{leancode}
\leanline{0}{theorem\ parameter\_freeEnergy\_hasFDerivAt}
\leanline{4}{(ν\ :\ Measure\ Z)\ (q\ :\ Z\ →\ ℝ)\ (U\ :\ A\ →\ Z\ →\ ℝ)}
\leanline{4}{(DU\ :\ A\ →\ Z\ →\ A\ →L[ℝ]\ ℝ)\ (a₀\ :\ A)\ (s\ :\ Set\ A)\ (bound\ :\ Z\ →\ ℝ)}
\leanline{4}{(hs\ :\ s\ ∈\ {\leanfallback 𝓝}\ a₀)\ (hq\ :\ Measurable\ q)}
\leanline{4}{(hU\ :\ ∀ᶠ\ a\ in\ {\leanfallback 𝓝}\ a₀,\ Measurable\ (U\ a))}
\leanline{4}{(hbase\ :\ Integrable\ (fun\ z\ ↦\ q\ z\ *\ (Real.log\ (q\ z)\ +\ U\ a₀\ z))\ ν)}
\leanline{4}{(hDU\ :\ AEStronglyMeasurable\ (fun\ z\ ↦\ q\ z\ •\ DU\ a₀\ z)\ ν)}
\leanline{4}{(hbound\ :\ ∀ᵐ\ z\ ∂ν,\ ∀\ a\ ∈\ s,\ |q\ z|\ *\ ‖DU\ a\ z‖\ ≤\ bound\ z)}
\leanline{4}{(hbound\_int\ :\ Integrable\ bound\ ν)}
\leanline{4}{(hdiff\ :\ ∀ᵐ\ z\ ∂ν,\ ∀\ a\ ∈\ s,\ HasFDerivAt\ (fun\ b\ ↦\ U\ b\ z)\ (DU\ a\ z)\ a)\ :}
\leanline{4}{HasFDerivAt\ (fun\ a\ ↦\ densityFreeEnergy\ ν\ q\ (U\ a))}
\leanline{6}{(∫\ z,\ q\ z\ •\ DU\ a₀\ z\ ∂ν)\ a₀}
\end{leancode}
\leansource{\texttt{density\_\allowbreak{}firstVariation\_\allowbreak{}of\_\allowbreak{}zero\_\allowbreak{}mass} \quad LatentDiffusion/\allowbreak{}SourceCalculus.lean, lines 118–148}
\begin{leancode}
\leanline{0}{theorem\ density\_firstVariation\_of\_zero\_mass}
\leanline{4}{(ν\ :\ Measure\ Z)\ (q\ r\ U\ :\ Z\ →\ ℝ)\ (s\ :\ Set\ ℝ)\ (bound\ :\ Z\ →\ ℝ)}
\leanline{4}{(hs\ :\ s\ ∈\ {\leanfallback 𝓝}\ (0\ :\ ℝ))\ (hq\ :\ Measurable\ q)\ (hr\ :\ Measurable\ r)}
\leanline{4}{(hU\ :\ Measurable\ U)}
\leanline{4}{(hbase\ :\ Integrable\ (fun\ z\ ↦\ q\ z\ *\ (Real.log\ (q\ z)\ +\ U\ z))\ ν)}
\leanline{4}{(hpositive\ :\ ∀ᵐ\ z\ ∂ν,\ ∀\ t\ ∈\ s,\ 0\ <\ q\ z\ +\ t\ *\ r\ z)}
\leanline{4}{(hbound\ :\ ∀ᵐ\ z\ ∂ν,\ ∀\ t\ ∈\ s,}
\leanline{6}{|r\ z\ *\ (Real.log\ (q\ z\ +\ t\ *\ r\ z)\ +\ U\ z\ +\ 1)|\ ≤\ bound\ z)}
\leanline{4}{(hbound\_int\ :\ Integrable\ bound\ ν)}
\leanline{4}{(hr\_int\ :\ Integrable\ r\ ν)\ (hmass\ :\ (∫\ z,\ r\ z\ ∂ν)\ =\ 0)\ :}
\leanline{4}{Integrable\ (fun\ z\ ↦\ r\ z\ *\ (Real.log\ (q\ z)\ +\ U\ z))\ ν\ ∧}
\leanline{6}{HasDerivAt\ (fun\ t\ ↦\ densityFreeEnergy\ ν\ (fun\ z\ ↦\ q\ z\ +\ t\ *\ r\ z)\ U)}
\leanline{8}{(∫\ z,\ r\ z\ *\ (Real.log\ (q\ z)\ +\ U\ z)\ ∂ν)\ 0}
\end{leancode}
\end{leanbox}
\subsection{\texorpdfstring{Gradient under the normalized Euclidean–Wasserstein metric}{Gradient under the normalized Euclidean–Wasserstein metric}}


Equip the product state with
        \(d^2((u,q),(v,r))=\|u-v\|^2+M^{-1}\sum_mW_2(q^m,r^m)^2\), and let
        \(V_u^m=\mathcal L_D(\theta,\cdot)-\log p_\phi(x^m\mid\cdot)\).
        Under the explicit first-variation and continuity-equation
        extra assumptions R2–R3, the parameter and density gradients are


        \[\nabla_u\widetilde F=\frac1M\sum_m\int\nabla_u V_u^m\,dq^m,
        \qquad
        \nabla_{q^m}\widetilde F=-\nabla_z\!\cdot
           \bigl(q^m\nabla_z(\log q^m+V_u^m)\bigr).\]


Thus the descent flow has
        \(\dot u=-M^{-1}\sum_m\int\nabla_uV_u^m\,dq^m\) and
        \(\partial_tq^m=\nabla_z\cdot(q^m\nabla_zV_u^m)+\Delta q^m\).



\textbf{Correction.} The inverse metric in each law block contributes
        \(M\), cancelling the objective's \(1/M\). There is consequently no
        additional \(1/M\) in the density drift or diffusion. These formulas
        use the descent loss \(V\), avoiding the inconsistent concise-loss sign
        in main lines 222–226.

\begin{leanbox}{Lean: \texttt{normalized\_\allowbreak{}product\_\allowbreak{}gradient\_\allowbreak{}representation}}
\leansource{\texttt{normalized\_\allowbreak{}product\_\allowbreak{}gradient\_\allowbreak{}representation} \quad LatentDiffusion/\allowbreak{}RegularFlowCalculus.lean, lines 548–590}
\begin{leancode}
\leanline{0}{theorem\ normalized\_product\_gradient\_representation}
\leanline{4}{(q\ qdot\ :\ Fin\ M\ →\ ℝ\ →\ E\ →\ ℝ)\ (U\ :\ Fin\ M\ →\ A\ →\ E\ →\ ℝ)}
\leanline{4}{(p\ pdot\ :\ ℝ\ →\ A)\ (t\ :\ ℝ)\ (v\ :\ Fin\ M\ →\ E\ →\ E)}
\leanline{4}{(data\ :\ ∀\ i,\ LocalEnergyDomination\ volume\ (q\ i)\ (qdot\ i)\ (fun\ u\ ↦\ U\ i\ (p\ u))}
\leanline{6}{(fun\ u\ x\ ↦\ fderiv\ ℝ\ (fun\ a\ ↦\ U\ i\ a\ x)\ (p\ u)\ (pdot\ u))\ t)}
\leanline{4}{(hp\ :\ ∀\ u,\ HasDerivAt\ p\ (pdot\ u)\ u)}
\leanline{4}{(hparam\ :\ ∀\ i\ x\ a,\ DifferentiableAt\ ℝ\ (fun\ b\ ↦\ U\ i\ b\ x)\ a)}
\leanline{4}{(hpos\ :\ ∀\ i\ x,\ 0\ <\ q\ i\ t\ x)\ (hq\ :\ ∀\ i,\ ContDiff\ ℝ\ 1\ (q\ i\ t))}
\leanline{4}{(hU\ :\ ∀\ i,\ ContDiff\ ℝ\ 1\ (U\ i\ (p\ t)))}
\leanline{4}{(hgrad\ :\ ∀\ i,\ Integrable}
\leanline{6}{(fun\ x\ ↦\ q\ i\ t\ x\ •\ gradient\ (fun\ b\ ↦\ U\ i\ b\ x)\ (p\ t))\ volume)}
\leanline{4}{(hcontinuity\ :\ ∀\ i,\ ∀ᵐ\ x\ ∂volume,}
\leanline{6}{qdot\ i\ t\ x\ =\ -fluxDivergence\ (fun\ y\ ↦\ q\ i\ t\ y\ •\ v\ i\ y)\ x)}
\leanline{4}{(htail\ :\ ∀\ i,\ NoBoundaryFlux\ (fun\ y\ ↦\ q\ i\ t\ y\ •\ v\ i\ y)}
\leanline{6}{(fun\ x\ ↦\ chemicalPotential\ (q\ i\ t)\ (U\ i\ (p\ t))\ x\ +\ 1))\ :}
\leanline{4}{HasDerivAt\ (fun\ u\ ↦\ normalizedDensityEnergy\ volume\ (fun\ i\ ↦\ q\ i\ u)\ U\ (p\ u))}
\leanline{6}{(normalizedProductGradientPairing\ volume\ (fun\ i\ ↦\ q\ i\ t)\ U\ (p\ t)\ (pdot\ t)\ v)\ t}
\end{leancode}
\end{leanbox}
\subsection{\texorpdfstring{de Bruijn identity for the joint flow}{de Bruijn identity for the joint flow}}


For a solution with finite energy, finite score work,
        justified differentiation under the integrals, and vanishing boundary
        terms or a valid cutoff limit, set


        \[I(u,q)=\|\nabla_u\widetilde F(u,q)\|^2+
          \frac1M\sum_m\int
          \left\|\nabla_z\log\frac{q^m}{\rho_u^m}\right\|^2dq^m.\]


The source identity, on that analytic domain, is


        \[\frac d{dt}\widetilde F(u_t,q_t)=-I(u_t,q_t),\qquad t>0.\]

\textbf{Correction.} Classical local smoothness in A2 does not by itself supply entropy integrability, a dominating time derivative, finite Fisher information or integration by parts at infinity. The argument at supplement lines 503–524 needs those analytic justifications. A consistent initial-law version of A2 is also required.

\begin{leanbox}{Lean: \texttt{normalized\_\allowbreak{}product\_\allowbreak{}deBruijn}}
\leansource{\texttt{normalized\_\allowbreak{}product\_\allowbreak{}deBruijn} \quad LatentDiffusion/\allowbreak{}RegularFlowCalculus.lean, lines 613–663}
\begin{leancode}
\leanline{0}{theorem\ normalized\_product\_deBruijn}
\leanline{4}{(q\ qdot\ :\ Fin\ M\ →\ ℝ\ →\ E\ →\ ℝ)\ (U\ :\ Fin\ M\ →\ A\ →\ E\ →\ ℝ)}
\leanline{4}{(p\ pdot\ :\ ℝ\ →\ A)\ (t\ :\ ℝ)}
\leanline{4}{(data\ :\ ∀\ i,\ LocalEnergyDomination\ volume\ (q\ i)\ (qdot\ i)\ (fun\ u\ ↦\ U\ i\ (p\ u))}
\leanline{6}{(fun\ u\ x\ ↦\ fderiv\ ℝ\ (fun\ a\ ↦\ U\ i\ a\ x)\ (p\ u)\ (pdot\ u))\ t)}
\leanline{4}{(hp\ :\ ∀\ u,\ HasDerivAt\ p\ (pdot\ u)\ u)}
\leanline{4}{(hparam\ :\ ∀\ i\ x\ a,\ DifferentiableAt\ ℝ\ (fun\ b\ ↦\ U\ i\ b\ x)\ a)}
\leanline{4}{(hpos\ :\ ∀\ i\ x,\ 0\ <\ q\ i\ t\ x)\ (hq\ :\ ∀\ i,\ ContDiff\ ℝ\ 1\ (q\ i\ t))}
\leanline{4}{(hU\ :\ ∀\ i,\ ContDiff\ ℝ\ 1\ (U\ i\ (p\ t)))}
\leanline{4}{(hgrad\ :\ ∀\ i,\ Integrable}
\leanline{6}{(fun\ x\ ↦\ q\ i\ t\ x\ •\ gradient\ (fun\ b\ ↦\ U\ i\ b\ x)\ (p\ t))\ volume)}
\leanline{4}{(hode\ :\ pdot\ t\ =\ -averagedParameterForce\ volume\ (fun\ i\ ↦\ q\ i\ t)\ U\ (p\ t))}
\leanline{4}{(hpde\ :\ ∀\ i,\ ∀ᵐ\ x\ ∂volume,\ qdot\ i\ t\ x\ =}
\leanline{6}{fluxDivergence\ (fokkerPlanckFlux\ (q\ i\ t)\ (U\ i\ (p\ t)))\ x)}
\leanline{4}{(htail\ :\ ∀\ i,\ NoBoundaryFlux\ (fokkerPlanckFlux\ (q\ i\ t)\ (U\ i\ (p\ t)))}
\leanline{6}{(fun\ x\ ↦\ chemicalPotential\ (q\ i\ t)\ (U\ i\ (p\ t))\ x\ +\ 1))\ :}
\leanline{4}{(∀\ i,\ Integrable}
\leanline{6}{(fun\ x\ ↦\ q\ i\ t\ x\ *\ ‖gradient\ (chemicalPotential\ (q\ i\ t)\ (U\ i\ (p\ t)))\ x‖\ \textasciicircum{}\ 2)}
\leanline{6}{volume)\ ∧}
\leanline{4}{HasDerivAt\ (fun\ u\ ↦\ normalizedDensityEnergy\ volume\ (fun\ i\ ↦\ q\ i\ u)\ U\ (p\ u))}
\leanline{6}{(-normalizedSlopeSquared\ volume\ (fun\ i\ ↦\ q\ i\ t)\ U\ (p\ t))\ t}
\end{leancode}
\end{leanbox}
\subsection{\texorpdfstring{Geodesics in the normalized product space}{Geodesics in the normalized product space}}


On R4, assume finite-cost optimal maps exist
        between every pair of intermediate laws and the endpoint optimal maps
        \(T_m:q^m\to r^m\) are unique up to \(q^m\)-null sets.
        A path is a constant-speed geodesic for the normalized product distance
        if and only if, for the same \(s\in[0,1]\),


        \[u_s=(1-s)u+sv,\qquad q_s^m=((1-s)\mathrm{id}+sT_m)_\#q^m.\]


For the R4 regular Brenier-map data, \(T_m=\nabla f^m\) with
        a convex potential and its actual gradient certificate, giving the
        source's gradient-map formula. Both directions of the product
        characterization are proved.



\textbf{Correction.} Use the constant-speed endpoint parameterization in all factors. The corrected formal domain supplies finite-cost regular optimal maps between intermediate laws and uniqueness of endpoint maps; the gradient formulation additionally supplies the convex-potential gradient certificate. These are explicit R4 transport assumptions.

\begin{leanbox}{Lean: \texttt{product\_\allowbreak{}geodesic\_\allowbreak{}iff\_\allowbreak{}optimalInterpolation}, \texttt{product\_\allowbreak{}geodesic\_\allowbreak{}iff\_\allowbreak{}gradientInterpolation}}
\leansource{\texttt{product\_\allowbreak{}geodesic\_\allowbreak{}iff\_\allowbreak{}optimalInterpolation} \quad LatentDiffusion/\allowbreak{}RegularTransportGeodesic.lean, lines 518–533}
\begin{leancode}
\leanline{0}{theorem\ product\_geodesic\_iff\_optimalInterpolation\ (hM\ :\ 0\ <\ M)}
\leanline{4}{(a\ :\ ℝ\ →\ P)\ (q\ :\ ℝ\ →\ Fin\ M\ →\ Measure\ Z)}
\leanline{4}{(T\ :\ ∀\ m,\ OptimalMap\ (q\ 0\ m)\ (q\ 1\ m))}
\leanline{4}{(hunique\ :\ ∀\ m,\ ∀\ S\ :\ OptimalMap\ (q\ 0\ m)\ (q\ 1\ m),\ S.transport\ =ᵐ[q\ 0\ m]\ (T\ m).transport)}
\leanline{4}{(hregular\ :\ ∀\ m,\ OptimalMapRegularity\ (fun\ t\ ↦\ q\ t\ m))\ :}
\leanline{4}{IsProductGeodesic\ a\ q\ ↔}
\leanline{6}{(∀\ s\ ∈\ Icc\ (0\ :\ ℝ)\ 1,\ a\ s\ =\ parameterInterpolate\ (a\ 0)\ (a\ 1)\ s)\ ∧}
\leanline{8}{∀\ m\ s,\ s\ ∈\ Icc\ (0\ :\ ℝ)\ 1\ →\ q\ s\ m\ =\ optimalInterpolation\ (T\ m)\ s}
\end{leancode}
\leansource{\texttt{product\_\allowbreak{}geodesic\_\allowbreak{}iff\_\allowbreak{}gradientInterpolation} \quad LatentDiffusion/\allowbreak{}RegularTransportGeodesic.lean, lines 537–551}
\begin{leancode}
\leanline{0}{theorem\ product\_geodesic\_iff\_gradientInterpolation\ (hM\ :\ 0\ <\ M)}
\leanline{4}{(a\ :\ ℝ\ →\ P)\ (q\ :\ ℝ\ →\ Fin\ M\ →\ Measure\ Z)}
\leanline{4}{(T\ :\ ∀\ m,\ RegularBrenierMap\ (q\ 0\ m)\ (q\ 1\ m))}
\leanline{4}{(hunique\ :\ ∀\ m,\ ∀\ S\ :\ OptimalMap\ (q\ 0\ m)\ (q\ 1\ m),}
\leanline{6}{S.transport\ =ᵐ[q\ 0\ m]\ (T\ m).transport)}
\leanline{4}{(hregular\ :\ ∀\ m,\ OptimalMapRegularity\ (fun\ t\ ↦\ q\ t\ m))\ :}
\leanline{4}{IsProductGeodesic\ a\ q\ ↔}
\leanline{6}{(∀\ s\ ∈\ Icc\ (0\ :\ ℝ)\ 1,\ a\ s\ =\ parameterInterpolate\ (a\ 0)\ (a\ 1)\ s)\ ∧}
\leanline{8}{∀\ m\ s,\ s\ ∈\ Icc\ (0\ :\ ℝ)\ 1\ →\ q\ s\ m\ =\ (q\ 0\ m).map}
\leanline{10}{(fun\ z\ ↦\ (1\ -\ s)\ •\ z\ +\ s\ •\ gradient\ (T\ m).potential\ z)}
\end{leancode}
\end{leanbox}
\subsection{\texorpdfstring{Strong geodesic convexity of the tilted free energy}{Strong geodesic convexity of the tilted free energy}}


Suppose each
        \(U_m(u,z)=-\log\rho_u^m(z)\) is jointly
        \(\lambda\)-strongly convex. For a constant-speed product
        Wasserstein geodesic \(\gamma_s\) between states \(a,b\), on a
        finite, integrable R1/\allowbreak{}R4 density-energy domain, the corrected result is


        \[\widetilde F(\gamma_s)\le(1-s)\widetilde F(a)
          +s\widetilde F(b)-\frac\lambda2s(1-s)d(a,b)^2,
          \qquad 0\le s\le1.\]


The potential term uses joint strong convexity along optimal
        displacement interpolations; the negative entropy term uses
        displacement convexity.



\textbf{Correction.} The normalized product distance is essential for the constant to remain independent of \(M\). The corrected theorem uses finite entropy and integrable potentials along the full displacement interpolation; finite second moments alone do not imply those conditions.

\begin{leanbox}{Lean: \texttt{entropy\_\allowbreak{}displacement\_\allowbreak{}convex}, \texttt{displacement\_\allowbreak{}strongConvex}}
\leansource{\texttt{entropy\_\allowbreak{}displacement\_\allowbreak{}convex} \quad LatentDiffusion/\allowbreak{}RegularTransport.lean, lines 193–204}
\begin{leancode}
\leanline{0}{theorem\ entropy\_displacement\_convex\ (s\ :\ ℝ)\ (hs\ :\ s\ ∈\ Icc\ (0\ :\ ℝ)\ 1)\ :}
\leanline{4}{D.entropy\ s\ ≤\ (1\ -\ s)\ *\ D.entropy\ 0\ +\ s\ *\ D.entropy\ 1}
\end{leancode}
\leansource{\texttt{displacement\_\allowbreak{}strongConvex} \quad LatentDiffusion/\allowbreak{}RegularTransportFirstOrder.lean, lines 731–747}
\begin{leancode}
\leanline{0}{theorem\ displacement\_strongConvex\ (hM\ :\ 0\ <\ M)\ (hrho\ :\ 0\ ≤\ rho)}
\leanline{4}{(s\ :\ ℝ)\ (hs\ :\ s\ ∈\ Icc\ (0\ :\ ℝ)\ 1)\ :}
\leanline{4}{R.freeEnergy\ s\ ≤\ (1\ -\ s)\ *\ R.freeEnergy\ 0\ +\ s\ *\ R.freeEnergy\ 1-}
\leanline{6}{rho\ /\ 2\ *\ s\ *\ (1\ -\ s)\ *\ R.distanceSq}
\end{leancode}
\end{leanbox}
\subsection{\texorpdfstring{First-order lower bound along a transport geodesic}{First-order lower bound along a transport geodesic}}


Let \(\gamma_s\) run from \((u,q)\) to
        \((v,r)\), with initial optimal transport velocities
        \(w_m=T_m-\mathrm{id}\). Write
        \(s_m=\nabla_z\log(q^m/\rho_u^m)\). On R2–R4, use locally dominated
        time derivatives, the initial continuity equation for these actual
        transport velocities, and integrable spatial tails. Then


        \[\liminf_{s\downarrow0}
          \frac{\widetilde F(\gamma_s)-\widetilde F(\gamma_0)}s
          \ge\langle v-u,\nabla_u\widetilde F(u,q)\rangle
          +\frac1M\sum_m\int\langle w_m,s_m\rangle\,dq^m.\]

\textbf{Correction.} Almost-everywhere differentiability of a density alone is insufficient. A finite-energy entropy-subgradient domain, finite score norm or another condition giving integrable transport pairings is required. The velocity formulation explicitly identifies the Wasserstein tangent pairing used in the source.

\begin{leanbox}{Lean: \texttt{lower\_\allowbreak{}directional\_\allowbreak{}liminf\_\allowbreak{}law\_\allowbreak{}pairing}}
\leansource{\texttt{lower\_\allowbreak{}directional\_\allowbreak{}liminf\_\allowbreak{}law\_\allowbreak{}pairing} \quad LatentDiffusion/\allowbreak{}RegularDirectional.lean, lines 178–183}
\begin{leancode}
\leanline{0}{theorem\ lower\_directional\_liminf\_law\_pairing\ :}
\leanline{4}{(lawPairing\ (D\ :=\ D)\ (U\ :=\ U)\ (a\ :=\ a)\ (b\ :=\ b)\ :\ EReal)\ ≤}
\leanline{6}{Filter.liminf\ (fun\ s\ :\ ℝ\ ↦}
\leanline{8}{(slope\ (lawEnergy\ (D\ :=\ D)\ (U\ :=\ U)\ (a\ :=\ a)\ (b\ :=\ b))\ 0\ s\ :\ EReal))\ ({\leanfallback 𝓝}[>]\ 0)}
\end{leancode}
\end{leanbox}
\subsection{\texorpdfstring{Strong log-concavity implies extended log-Sobolev control}{Strong log-concavity implies extended log-Sobolev control}}


For \(M\ge1\), jointly \(\lambda\)-strongly convex tilted potentials with \(\lambda>0\), and the finite-energy density/\allowbreak{}transport domain R1, R3 and R4,


 \[2\lambda[\widetilde F(u,q)-\widetilde F_\star]\le I(u,q).\]


Sum the actual entropy and potential tangent pairings first. Completing one square for the averaged parameter force and the normalized latent squares proves the inequality with exactly the source Fisher information. The posterior optimizer is independently constructed before it is used.



\textbf{Correction.} The proof chooses an admissible minimizer at supplement line 587 before its existence has been established in that chain. The checked marginal/\allowbreak{}posterior construction supplies it first. The entropy and potential tangent inequalities are then proved from the explicit extra assumptions.

\begin{leanbox}{Lean: \texttt{existsUnique\_\allowbreak{}variationalOptimizer}, \texttt{extendedLSI}}
\leansource{\texttt{existsUnique\_\allowbreak{}variationalOptimizer} \quad LatentDiffusion/\allowbreak{}RegularModel.lean, lines 66–74}
\begin{leancode}
\leanline{0}{theorem\ existsUnique\_variationalOptimizer\ (U\ :\ Fin\ M\ →\ IPLA.Joint\ dθ\ dz\ →\ ℝ)}
\leanline{4}{(hM\ :\ 0\ <\ M)\ (m\ :\ ℝ)\ (hm\ :\ 0\ <\ m)}
\leanline{4}{(hdiff\ :\ ∀\ i,\ Differentiable\ ℝ\ (U\ i))}
\leanline{4}{(hsc\ :\ ∀\ i,\ StrongConvexOn\ univ\ m\ (U\ i))\ :}
\leanline{4}{∃!\ p\ :\ EuclideanSpace\ ℝ\ (Fin\ dθ)\ ×}
\leanline{8}{(Fin\ M\ →\ ProbabilityMeasure\ (EuclideanSpace\ ℝ\ (Fin\ dz))),}
\leanline{6}{(densityModel\ U\ m\ hm\ hdiff\ hsc).family.IsGlobalMin\ p.1\ p.2}
\end{leancode}
\leansource{\texttt{extendedLSI} \quad LatentDiffusion/\allowbreak{}RegularTransportFirstOrder.lean, lines 661–718}
\begin{leancode}
\leanline{0}{theorem\ extendedLSI\ (hM\ :\ 0\ <\ M)\ (hrho\ :\ 0\ <\ rho)\ :}
\leanline{4}{2\ *\ rho\ *\ (R.freeEnergy\ 0\ -\ R.freeEnergy\ 1)\ ≤\ R.fisherInformation}
\end{leancode}
\end{leanbox}
\subsection{\texorpdfstring{Metric-speed control of distance along the flow}{Metric-speed control of distance along the flow}}


Under the R3 equations and R5 characteristic
        realization, the actual normalized product-distance upper-right slope
        to every fixed state \(b\) is at most \(\sqrt{I(a_t)}\).
        With integrable characteristic velocities, the proved integrated
        estimate is, for \(s\le t\),


        \[|d(a_t,b)-d(a_s,b)|\le\int_s^t\sqrt{I(a_r)}\,dr.\]


Where the ordinary distance derivative exists it is also at most
        \(\sqrt{I(a_t)}\). Both conclusions follow from the actual coupling
        costs and characteristic derivatives; a speed inequality is not an
        assumption.



\textbf{Correction.} The source says the ordinary derivative exists and satisfies the bound at every positive time. Even a smooth Euclidean path can have a corner in its distance to a fixed point when it crosses that point. An almost-everywhere metric-derivative statement, or an upper-Dini-derivative formulation, supplies the integrated estimate actually needed by the proof.

\begin{leanbox}{Lean: \texttt{distance\_\allowbreak{}right\_\allowbreak{}slope}, \texttt{kinetic\_\allowbreak{}energy\_\allowbreak{}eq\_\allowbreak{}fisher}, \texttt{distance\_\allowbreak{}integral\_\allowbreak{}bound}}
\leansource{\texttt{distance\_\allowbreak{}right\_\allowbreak{}slope} \quad LatentDiffusion/\allowbreak{}RegularProductMetricFlow.lean, lines 149–170}
\begin{leancode}
\leanline{0}{theorem\ distance\_right\_slope}
\leanline{4}{(a\ :\ ℝ\ →\ P)\ (Y\ :\ ℝ\ →\ ∀\ m,\ Lp\ Z\ 2\ (Q\ m))\ (astar\ :\ P)}
\leanline{4}{(target\ :\ Fin\ M\ →\ Measure\ Z)}
\leanline{4}{(anchor\ :\ ∀\ t\ m,\ RegularMetricFlow.OptimalAnchor\ (Q\ m)\ (Y\ t\ m)\ (target\ m))}
\leanline{4}{(t\ :\ ℝ)\ (adot\ :\ P)\ (v\ :\ ∀\ m,\ Lp\ Z\ 2\ (Q\ m))}
\leanline{4}{(ha\ :\ HasDerivAt\ a\ adot\ t)\ (hY\ :\ ∀\ m,\ HasDerivAt\ (fun\ s\ ↦\ Y\ s\ m)\ (v\ m)\ t)}
\leanline{4}{(r\ :\ ℝ)\ (hr\ :\ Real.sqrt\ (‖adot‖\ \textasciicircum{}\ 2\ +\ (M\ :\ ℝ)⁻¹\ *\ ∑\ m,\ ‖v\ m‖\ \textasciicircum{}\ 2)\ <\ r)\ :}
\leanline{4}{∀ᶠ\ s\ in\ {\leanfallback 𝓝}[>]\ t,\ slope\ (distance\ Q\ a\ Y\ astar\ target)\ t\ s\ <\ r}
\end{leancode}
\leansource{\texttt{kinetic\_\allowbreak{}energy\_\allowbreak{}eq\_\allowbreak{}fisher} \quad LatentDiffusion/\allowbreak{}RegularProductMetricFlow.lean, lines 239–251}
\begin{leancode}
\leanline{0}{theorem\ kinetic\_energy\_eq\_fisher}
\leanline{4}{(Y\ v\ :\ ∀\ m,\ Lp\ Z\ 2\ (Q\ m))\ (field\ :\ Fin\ M\ →\ Z\ →\ Z)\ (force\ adot\ :\ P)}
\leanline{4}{(hadot\ :\ adot\ =\ -force)\ (hmeas\ :\ ∀\ m,\ Measurable\ (field\ m))}
\leanline{4}{(hvelocity\ :\ ∀\ m,\ (v\ m\ :\ Ω\ m\ →\ Z)\ =ᵐ[Q\ m]}
\leanline{6}{(fun\ ω\ ↦\ -(field\ m\ ((Y\ m)\ ω))))\ :}
\leanline{4}{‖adot‖\ \textasciicircum{}\ 2\ +\ (M\ :\ ℝ)⁻¹\ *\ ∑\ m,\ ‖v\ m‖\ \textasciicircum{}\ 2\ =}
\leanline{6}{‖force‖\ \textasciicircum{}\ 2\ +\ (M\ :\ ℝ)⁻¹\ *\ ∑\ m,}
\leanline{8}{∫\ z,\ ‖field\ m\ z‖\ \textasciicircum{}\ 2\ ∂(Q\ m).map\ (Y\ m)}
\end{leancode}
\leansource{\texttt{distance\_\allowbreak{}integral\_\allowbreak{}bound} \quad LatentDiffusion/\allowbreak{}RegularProductMetricFlow.lean, lines 176–199}
\begin{leancode}
\leanline{0}{theorem\ distance\_integral\_bound}
\leanline{4}{(a\ adot\ :\ ℝ\ →\ P)\ (Y\ v\ :\ ℝ\ →\ ∀\ m,\ Lp\ Z\ 2\ (Q\ m))\ (astar\ :\ P)}
\leanline{4}{(target\ :\ Fin\ M\ →\ Measure\ Z)}
\leanline{4}{(anchor\ :\ ∀\ t\ m,\ RegularMetricFlow.OptimalAnchor\ (Q\ m)\ (Y\ t\ m)\ (target\ m))}
\leanline{4}{(s\ t\ :\ ℝ)\ (hst\ :\ s\ ≤\ t)}
\leanline{4}{(ha\ :\ ∀\ u\ ∈\ uIcc\ s\ t,\ HasDerivAt\ a\ (adot\ u)\ u)}
\leanline{4}{(hY\ :\ ∀\ u\ ∈\ uIcc\ s\ t,\ ∀\ m,\ HasDerivAt\ (fun\ r\ ↦\ Y\ r\ m)\ (v\ u\ m)\ u)}
\leanline{4}{(hv\ :\ IntervalIntegrable\ (fun\ u\ ↦\ lift\ Q\ (adot\ u)\ (v\ u))\ volume\ s\ t)\ :}
\leanline{4}{|distance\ Q\ a\ Y\ astar\ target\ t\ -\ distance\ Q\ a\ Y\ astar\ target\ s|\ ≤}
\leanline{6}{∫\ u\ in\ s..t,\ Real.sqrt\ (‖adot\ u‖\ \textasciicircum{}\ 2\ +\ (M\ :\ ℝ)⁻¹\ *\ ∑\ m,\ ‖v\ u\ m‖\ \textasciicircum{}\ 2)}
\end{leancode}
\end{leanbox}
\subsection{\texorpdfstring{Lower semicontinuity of the tilted energy}{Lower semicontinuity of the tilted energy}}


On the R6 density chart equipped with the
        actual normalized product transport metric, assume pointwise
        continuous densities and potentials and a common local integrable
        envelope for \(q\log q+qU\). The actual normalized energy


        \[\widetilde F(x)=\frac1M\sum_m\int q_x^m(z)
          [\log q_x^m(z)+U_m(u_x,z)]\,dz\]


is continuous, hence lower semicontinuous. In particular, for
        \(x_n\to x\) in this chart,


        \[\widetilde F(x_n)\longrightarrow\widetilde F(x),\qquad
          \widetilde F(x)\le\liminf_n\widetilde F(x_n).\]


Dominated convergence proves this statement on the explicitly
        identified probability-law domain.



\textbf{Correction.} Joint continuity of \(R\) is used at supplement line 665 but is not an explicit displayed assumption. The corrected proof instead states the sufficient R6 pointwise density/\allowbreak{}potential continuity and local integrable domination for the total energy. The limiting argument uses a complete admissible R6 domain; differentiable densities alone do not form a complete Wasserstein space.

\begin{leanbox}{Lean: \texttt{ProductDensityChart.\allowbreak{}lowerSemicontinuous\_\allowbreak{}energy}}
\leansource{\texttt{ProductDensityChart.\allowbreak{}lowerSemicontinuous\_\allowbreak{}energy} \quad LatentDiffusion/\allowbreak{}RegularDensityDomain.lean, lines 74–75}
\begin{leancode}
\leanline{0}{theorem\ ProductDensityChart.lowerSemicontinuous\_energy\ :\ LowerSemicontinuous\ D.energy}
\end{leancode}
\end{leanbox}
\subsection{\texorpdfstring{A Cauchy flow sequence with optimal energies reaches a minimizer}{A Cauchy flow sequence with optimal energies reaches a minimizer}}


The corrected compactness step is: let the
        ambient metric space be complete and the energy \(E\) be lower
        semicontinuous with finite infimum \(E_\star\). If
        \(t_n\uparrow\infty\), \((a_{t_n})\) is Cauchy, and
        \(E(a_{t_n})\to E_\star\), then there is \(a_\star\) such that


        \[a_{t_n}\longrightarrow a_\star,
          \qquad E(a_\star)=E_\star.\]


The Cauchy property yields a limit in the complete ambient space;
        lower semicontinuity and optimal-energy convergence identify it as a
        minimizer.



\textbf{Correction.} The printed hypothesis gives an arbitrary Cauchy sequence unrelated to the flow, then concludes convergence of a flow subsequence. The proof needs the Cauchy sequence to be the stated flow sequence and its energies to approach the infimum. Completeness of all \(\mathcal P_2\) does not ensure that its limit retains a differentiable density; the energy domain and admissibility of the limit need separate treatment.

\begin{leanbox}{Lean: \texttt{ProductDensityChart.\allowbreak{}cauchy\_\allowbreak{}minimizer}, \texttt{metric\_\allowbreak{}flow\_\allowbreak{}talagrand}}
\leansource{\texttt{ProductDensityChart.\allowbreak{}cauchy\_\allowbreak{}minimizer} \quad LatentDiffusion/\allowbreak{}RegularDensityDomain.lean, lines 79–85}
\begin{leancode}
\leanline{0}{theorem\ ProductDensityChart.cauchy\_minimizer\ [CompleteSpace\ X]}
\leanline{4}{(Fstar\ :\ ℝ)\ (hstar\ :\ ∀\ x,\ Fstar\ ≤\ D.energy\ x)\ (x\ :\ ℕ\ →\ X)}
\leanline{4}{(hCauchy\ :\ CauchySeq\ x)\ (hvalue\ :\ Tendsto\ (fun\ n\ ↦\ D.energy\ (x\ n))\ atTop\ ({\leanfallback 𝓝}\ Fstar))\ :}
\leanline{4}{∃\ xstar,\ Tendsto\ x\ atTop\ ({\leanfallback 𝓝}\ xstar)\ ∧\ D.energy\ xstar\ =\ Fstar\ ∧}
\leanline{6}{∀\ y,\ D.energy\ xstar\ ≤\ D.energy\ y}
\end{leancode}
\leansource{\texttt{metric\_\allowbreak{}flow\_\allowbreak{}talagrand} \quad LatentDiffusion/\allowbreak{}RegularTalagrand.lean, lines 203–279}
\begin{leancode}
\leanline{0}{theorem\ metric\_flow\_talagrand}
\leanline{4}{\{X\ :\ Type*\}\ [MetricSpace\ X]\ [CompleteSpace\ X]}
\leanline{4}{(F\ :\ X\ →\ ℝ)\ (Fstar\ :\ ℝ)\ (x\ :\ ℝ\ →\ X)\ (I\ :\ ℝ\ →\ ℝ)\ (rho\ :\ ℝ)}
\leanline{4}{(hrho\ :\ 0\ <\ rho)\ (hF\ :\ LowerSemicontinuous\ F)\ (hstar\ :\ ∀\ y,\ Fstar\ ≤\ F\ y)}
\leanline{4}{(hx\ :\ ContinuousOn\ x\ (Ici\ 0))}
\leanline{4}{(henergy\ :\ ContinuousOn\ (fun\ t\ ↦\ F\ (x\ t)\ -\ Fstar)\ (Ici\ 0))}
\leanline{4}{(hderiv\ :\ ∀\ t,\ 0\ ≤\ t\ →\ HasDerivWithinAt\ (fun\ u\ ↦\ F\ (x\ u)\ -\ Fstar)}
\leanline{6}{(-I\ t)\ (Ici\ t)\ t)}
\leanline{4}{(hI\ :\ ∀\ t,\ 0\ ≤\ t\ →\ 0\ ≤\ I\ t)}
\leanline{4}{(hlsi\ :\ ∀\ t,\ 0\ ≤\ t\ →\ 2\ *\ rho\ *\ (F\ (x\ t)\ -\ Fstar)\ ≤\ I\ t)}
\leanline{4}{(hspeed\ :\ ∀\ s\ t,\ 0\ ≤\ s\ →\ 0\ ≤\ t\ →\ ∀\ r,\ Real.sqrt\ (I\ t)\ <\ r\ →}
\leanline{6}{∀ᶠ\ u\ in\ {\leanfallback 𝓝}[>]\ t,\ slope\ (fun\ v\ ↦\ dist\ (x\ s)\ (x\ v))\ t\ u\ <\ r)\ :}
\leanline{4}{∃\ xstar,\ Tendsto\ (fun\ n\ :\ ℕ\ ↦\ x\ n)\ atTop\ ({\leanfallback 𝓝}\ xstar)\ ∧}
\leanline{6}{F\ xstar\ =\ Fstar\ ∧\ (∀\ y,\ F\ xstar\ ≤\ F\ y)\ ∧}
\leanline{6}{∀\ t,\ 0\ ≤\ t\ →\ rho\ *\ dist\ (x\ t)\ xstar\ \textasciicircum{}\ 2\ ≤\ 2\ *\ (F\ (x\ t)\ -\ Fstar)}
\end{leancode}
\end{leanbox}
\subsection{\texorpdfstring{Extended log-Sobolev control implies distance to the minimizing set}{Extended log-Sobolev control implies distance to the minimizing set}}


For the given state \(a=(u,q)\), the intended
        extended Talagrand inequality is


        \[\lambda\,d(a,\mathcal M_\star)^2
          \le2[\widetilde F(a)-\widetilde F_\star],\qquad
          d(a,\mathcal M_\star)=\inf_{b\in\mathcal M_\star}d(a,b).\]


A corrected flow proof requires a consistent global flow,
        dissipation, the metric-speed bound, extended LSI, lower
        semicontinuity and the preceding complete-space minimizer argument.
        Integrating the speed gives a total remaining path length at most
        \(\sqrt{2[\widetilde F(a)-\widetilde F_\star]/\lambda}\).

\begin{leanbox}{Lean: \texttt{exists\_\allowbreak{}minimizer\_\allowbreak{}talagrand}}
\leansource{\texttt{exists\_\allowbreak{}minimizer\_\allowbreak{}talagrand} \quad LatentDiffusion/\allowbreak{}RegularFlowTalagrand.lean, lines 196–217}
\begin{leancode}
\leanline{0}{theorem\ exists\_minimizer\_talagrand\ [CompleteSpace\ X]}
\leanline{4}{(Fstar\ rho\ :\ ℝ)\ (hrho\ :\ 0\ <\ rho)\ (hstar\ :\ ∀\ x,\ Fstar\ ≤\ D.energy\ x)}
\leanline{4}{(hxLSI\ :\ ∀\ x,\ 2\ *\ rho\ *\ (D.energy\ x-Fstar)\ ≤}
\leanline{6}{normalizedSlopeSquared\ volume\ (fun\ i\ ↦\ (D.chart\ i).density\ x)\ U\ (D.parameter\ x))\ :}
\leanline{4}{∃\ xstar,\ Tendsto\ (fun\ n\ :\ ℕ\ ↦\ R.state\ n)\ atTop\ ({\leanfallback 𝓝}\ xstar)\ ∧}
\leanline{6}{D.energy\ xstar\ =\ Fstar\ ∧\ (∀\ x,\ D.energy\ xstar\ ≤\ D.energy\ x)\ ∧}
\leanline{6}{∀\ t,\ 0\ ≤\ t\ →\ rho\ *\ dist\ (R.state\ t)\ xstar\ \textasciicircum{}\ 2\ ≤}
\leanline{8}{2\ *\ (D.energy\ (R.state\ t)-Fstar)}
\end{leancode}
\leanlabel{How the Lean proof follows the argument}
\textbf{Correction.} The definition \texttt{defn:x\allowbreak{}Talagrand} at
        supplement lines 614–620 takes a redundant infimum over the starting
        state itself. If a minimizer exists, that literal right-hand side is
        zero. Remove this infimum and keep the given starting state fixed; the
        only infimum is over the minimizing set. Section-wide A1/\allowbreak{}A2 and the
        analytic lemmas are also used, even though the theorem's local sentence
        mentions only xLSI.
\end{leanbox}
\subsection{\texorpdfstring{Full product-state exponential convergence}{Full product-state exponential convergence}}


The appendix's stronger form of the main
        convergence theorem asserts a unique minimizer
        \(a_\star=(u_\star,\pi_{u_\star}^{1:M})\) and


        \[d(a_t,a_\star)\le
          \sqrt{\frac{2[\widetilde F(a_t)-\widetilde F_\star]}\lambda}
          \le\sqrt{\frac{2[\widetilde F(a_0)-\widetilde F_\star]}\lambda}
             e^{-\lambda t},\qquad t>0.\]


Here \(d^2=\|u-v\|^2+M^{-1}\sum_mW_2(q^m,r^m)^2\). The first
        inequality uses the corrected Talagrand statement; the second uses
        dissipation, xLSI and Grönwall.



\textbf{Correction.} Read A2 on a consistent initial class and supply the analytic entropy/\allowbreak{}transport hypotheses, the corrected Talagrand definition, and the corrected Cauchy-flow argument. An independent optimizer construction avoids the circular use of a minimizer in the source xLSI proof.

\begin{leanbox}{Lean: \texttt{existsUnique\_\allowbreak{}optimizer\_\allowbreak{}and\_\allowbreak{}flow\_\allowbreak{}convergence}, \texttt{full\_\allowbreak{}product\_\allowbreak{}exponential\_\allowbreak{}convergence}, \texttt{distanceSq\_\allowbreak{}eq\_\allowbreak{}posterior}}
\leansource{\texttt{existsUnique\_\allowbreak{}optimizer\_\allowbreak{}and\_\allowbreak{}flow\_\allowbreak{}convergence} \quad LatentDiffusion/\allowbreak{}RegularConvergence.lean, lines 478–517}
\begin{leancode}
\leanline{0}{theorem\ existsUnique\_optimizer\_and\_flow\_convergence\ \{dθ\ dz\ :\ ℕ\}}
\leanline{4}{(U\ :\ Fin\ M\ →\ IPLA.Joint\ dθ\ dz\ →\ ℝ)\ (hM\ :\ 0\ <\ M)}
\leanline{4}{(rho\ :\ ℝ)\ (hrho\ :\ 0\ <\ rho)}
\leanline{4}{(hdiff\ :\ ∀\ i,\ Differentiable\ ℝ\ (U\ i))}
\leanline{4}{(hsc\ :\ ∀\ i,\ StrongConvexOn\ univ\ rho\ (U\ i))\ :}
\leanline{4}{let\ V\ :=\ fun\ i\ a\ z\ ↦\ U\ i\ (IPLA.pair\ a\ z)}
\leanline{4}{let\ ell\ :=\ fun\ a\ ↦\ (M\ :\ ℝ)⁻¹\ *\ ∑\ i,\ Real.log\ (∫\ z,\ Real.exp\ (-V\ i\ a\ z))}
\leanline{4}{∃!\ astar\ :\ EuclideanSpace\ ℝ\ (Fin\ dθ),}
\leanline{6}{(∀\ a,\ ell\ a\ ≤\ ell\ astar)\ ∧}
\leanline{6}{∀\ S\ :\ ClassicalFlow\ M\ dz\ rho\ astar\ V,}
\leanline{8}{(∀\ i\ z,\ S.targetDensity\ i\ z\ =\ Real.exp\ (-V\ i\ astar\ z)/}
\leanline{10}{(∫\ w,\ Real.exp\ (-V\ i\ astar\ w)))\ →}
\leanline{8}{∀\ t,\ 0\ ≤\ t\ →\ ‖S.parameter\ t-astar‖\ ≤}
\leanline{10}{Real.sqrt\ (2*(S.energy\ 0-S.targetEnergy)/rho)*Real.exp\ (-rho*t)}
\end{leancode}
\leansource{\texttt{full\_\allowbreak{}product\_\allowbreak{}exponential\_\allowbreak{}convergence} \quad LatentDiffusion/\allowbreak{}RegularConvergence.lean, lines 416–450}
\begin{leancode}
\leanline{0}{theorem\ full\_product\_exponential\_convergence\ (hM\ :\ 0\ <\ M)\ (hrho\ :\ 0\ <\ rho)}
\leanline{4}{(hweight\ :\ ∀\ i\ a,\ Integrable\ (fun\ z\ ↦\ Real.exp\ (-U\ i\ a\ z))\ volume)}
\leanline{4}{(hZ\ :\ ∀\ i\ a,\ 0\ <\ ∫\ z,\ Real.exp\ (-U\ i\ a\ z))}
\leanline{4}{(hmax\ :\ ∀\ a,\ S.evidence\ a\ ≤\ S.evidence\ astar)}
\leanline{4}{(hpost\ :\ ∀\ i\ z,\ S.targetDensity\ i\ z\ =}
\leanline{6}{Real.exp\ (-U\ i\ astar\ z)/(∫\ w,\ Real.exp\ (-U\ i\ astar\ w)))}
\leanline{4}{(t\ :\ ℝ)\ (ht\ :\ 0\ ≤\ t)\ :}
\leanline{4}{Real.sqrt\ (S.transport\ t).distanceSq\ ≤\ Real.sqrt\ (2*(S.energy\ t-S.targetEnergy)/rho)\ ∧}
\leanline{4}{Real.sqrt\ (2*(S.energy\ t-S.targetEnergy)/rho)\ ≤}
\leanline{6}{Real.sqrt\ (2*(S.energy\ 0-S.targetEnergy)/rho)*Real.exp\ (-rho*t)}
\end{leancode}
\leansource{\texttt{distanceSq\_\allowbreak{}eq\_\allowbreak{}posterior} \quad LatentDiffusion/\allowbreak{}RegularConvergence.lean, lines 292–299}
\begin{leancode}
\leanline{0}{theorem\ distanceSq\_eq\_posterior}
\leanline{4}{(hpost\ :\ ∀\ i\ z,\ S.targetDensity\ i\ z\ =}
\leanline{6}{Real.exp\ (-U\ i\ astar\ z)\ /\ (∫\ w,\ Real.exp\ (-U\ i\ astar\ w)))\ (t\ :\ ℝ)\ :}
\leanline{4}{(S.transport\ t).distanceSq\ =\ ‖S.parameter\ t-astar‖\ \textasciicircum{}\ 2\ +\ (M\ :\ ℝ)⁻¹\ *\ ∑\ i,}
\leanline{6}{ULA.wasserstein2Real\ ((S.transport\ t).law\ i\ :\ Measure\ Z)}
\leanline{8}{(Density.posteriorMeasure\ volume\ (fun\ z\ ↦\ Real.exp\ (-U\ i\ astar\ z))}
\leanline{10}{(∫\ w,\ Real.exp\ (-U\ i\ astar\ w)))\ \textasciicircum{}\ 2}
\end{leancode}
\end{leanbox}
\subsection{\texorpdfstring{Continuous mean-field solution and empirical-cloud comparison}{Continuous mean-field solution and empirical-cloud comparison}}


Under R7, construct a continuous
        \(L^2\) Brownian mean-field integral solution \(Y_t\), with
        deterministic initial parameter, and a synchronous
        stationary posterior reference \(S_t\). The reference parameter is the
        actual optimizer \(u_\star\), and every physical reference particle for
        datum \(m\) has law \(\pi_{u_\star}^m\). For their per-datum empirical
        laws \(Q_t^{m,N},Q_{\star,t}^{m,N}\), the corrected comparison is


        \[\left(\mathbb E\left[\|u_t-u_\star\|^2+
          \frac1M\sum_m W_2(Q_t^{m,N},Q_{\star,t}^{m,N})^2\right]
          \right)^{1/2}
          \le\|Y_t-S_t\|_{L^2}\le e^{-\lambda t}\|Y_0-S_0\|_{L^2}.\]


The cloud norm averages physical squared particle distances by
        \(MN\). Both the empirical statistic's measurability/\allowbreak{}integrability and
        its bound by this actual coupling norm are proved.



\textbf{Correction.} Arbitrary independent target samples do not satisfy the identity equating their pairwise \(L^2\) cost to \(W_2\) at supplement line 781. The corrected theorem uses a constructed synchronous stationary posterior reference and an explicit initial coupling norm as its prefactor. The solution class is continuous \(L^2\) Brownian integral solutions. Replicating each datum potential \(N\) times identifies the target with the original \(M\)-datum optimizer; different data may have different posterior laws.

\begin{leanbox}{Lean: \texttt{exists\_\allowbreak{}meanField}, \texttt{meanField\_\allowbreak{}actual\_\allowbreak{}empirical\_\allowbreak{}convergence}, \texttt{empiricalProduct\_\allowbreak{}rms\_\allowbreak{}le\_\allowbreak{}l2} and 1 more}
\leansource{\texttt{exists\_\allowbreak{}meanField} \quad LatentDiffusion/\allowbreak{}RegularMeanField.lean, lines 86–102}
\begin{leancode}
\leanline{0}{theorem\ exists\_meanField}
\leanline{4}{(hK\ :\ 0\ <\ K)\ (U\ :\ Fin\ K\ →\ Cloud.Joint\ dθ\ dz\ →\ ℝ)}
\leanline{4}{(m\ :\ ℝ)\ (L\ :\ ℝ≥0)\ (hdiff\ :\ ∀\ i,\ Differentiable\ ℝ\ (U\ i))}
\leanline{4}{(hsc\ :\ ∀\ i,\ StrongConvexOn\ Set.univ\ m\ (U\ i))}
\leanline{4}{(hLip\ :\ ∀\ i,\ LipschitzWith\ L\ (gradient\ (U\ i)))}
\leanline{4}{(initial\ :\ Lp\ (Cloud.State\ K\ dθ\ dz)\ 2\ P)}
\leanline{4}{(B\ :\ (Fin\ K\ ×\ Fin\ dz)\ →\ ℝ≥0\ →\ Ω\ →\ ℝ)}
\leanline{4}{(hB\ :\ ∀\ i,\ IsPreBrownianReal\ (B\ i)\ P)\ :}
\leanline{4}{∃\ Y\ :\ ℝ\ →\ Lp\ (Cloud.State\ K\ dθ\ dz)\ 2\ P,\ Continuous\ Y\ ∧\ Y\ 0\ =\ initial\ ∧}
\leanline{6}{∀\ t,\ 0\ ≤\ t\ →\ Y\ t\ -\ Y\ 0\ =}
\leanline{8}{-(∫\ r\ in\ 0..t,\ force\ P\ (Cloud.meanForce\_regular\ hK\ U\ m\ L\ hdiff\ hsc\ hLip)\ (Y\ r))\ +}
\leanline{10}{physicalNoise\ P\ B\ hB\ t}
\end{leancode}
\leansource{\texttt{meanField\_\allowbreak{}actual\_\allowbreak{}empirical\_\allowbreak{}convergence} \quad LatentDiffusion/\allowbreak{}RegularSourceTheorem.lean, lines 512–564}
\begin{leancode}
\leanline{0}{theorem\ meanField\_actual\_empirical\_convergence}
\leanline{4}{\{M\ N\ :\ ℕ\}\ (hK\ :\ 0\ <\ M\ *\ N)\ (U\ :\ Fin\ (M\ *\ N)\ →\ Cloud.Joint\ dθ\ dz\ →\ ℝ)}
\leanline{4}{(m\ :\ ℝ)\ (L\ :\ ℝ≥0)\ (hm\ :\ 0\ <\ m)}
\leanline{4}{(hdiff\ :\ ∀\ i,\ Differentiable\ ℝ\ (U\ i))}
\leanline{4}{(hsc\ :\ ∀\ i,\ StrongConvexOn\ Set.univ\ m\ (U\ i))}
\leanline{4}{(hLip\ :\ ∀\ i,\ LipschitzWith\ L\ (gradient\ (U\ i)))}
\leanline{4}{(initial\ :\ Lp\ (Latent\ (M\ *\ N)\ dz)\ 2\ P)}
\leanline{4}{(hinitial\ :\ P.map\ initial\ =\ ULA.gibbsMeasure}
\leanline{6}{(frozenPotential\ U\ (optimizerParameter\ hK\ U\ m\ hm\ hdiff\ hsc)))}
\leanline{4}{(hposterior\ :\ ∀\ i,\ P.map\ (fun\ ω\ =>\ physicalParticleCLM\ i}
\leanline{8}{((packLp\ P\ (optimizerParameter\ hK\ U\ m\ hm\ hdiff\ hsc)\ initial)\ ω))\ =}
\leanline{6}{RegularEvidence.posterior\ volume\ (latentPotential\ U\ i)}
\leanline{8}{(optimizerParameter\ hK\ U\ m\ hm\ hdiff\ hsc))}
\leanline{4}{(data\ :\ ∀\ i,\ RegularEvidence.LocalNormalizerRegularity\ volume\ (latentPotential\ U\ i)}
\leanline{6}{(optimizerParameter\ hK\ U\ m\ hm\ hdiff\ hsc))}
\leanline{4}{(B\ :\ (Fin\ (M\ *\ N)\ ×\ Fin\ dz)\ →\ ℝ≥0\ →\ Ω\ →\ ℝ)}
\leanline{4}{(hB\ :\ ∀\ i,\ IsPreBrownianReal\ (B\ i)\ P)}
\leanline{4}{(hscalar\ :\ ∀\ s\ :\ ℝ,\ 0\ <\ s\ →\ iIndepFun\ (fun\ p\ :\ ℕ\ ×\ (Fin\ (M\ *\ N)\ ×\ Fin\ dz)\ =>\ fun\ ω\ =>}
\leanline{6}{ULA.Hilbert.brownianEulerNoise\ P\ B\ hB\ s\ p.1\ ω\ p.2)\ P)}
\leanline{4}{(hind\ :\ ∀\ s\ :\ ℝ,\ 0\ <\ s\ →\ IndepFun\ initial}
\leanline{6}{(fun\ ω\ n\ =>\ ULA.Hilbert.brownianEulerNoise\ P\ B\ hB\ s\ n\ ω)\ P)}
\leanline{4}{(Y\ :\ ℝ\ →\ Lp\ (Cloud.State\ (M\ *\ N)\ dθ\ dz)\ 2\ P)\ (hY\ :\ Continuous\ Y)}
\leanline{4}{(hYe\ :\ ∀\ t,\ 0\ ≤\ t\ →\ Y\ t\ -\ Y\ 0\ =}
\leanline{6}{-(∫\ r\ in\ 0..t,\ RegularMeanField.force\ P}
\leanline{8}{(Cloud.meanForce\_regular\ hK\ U\ m\ L\ hdiff\ hsc\ hLip)\ (Y\ r))\ +\ physicalNoise\ P\ B\ hB\ t)}
\leanline{4}{(a₀\ :\ Parameter\ dθ)\ (ha₀\ :\ parameterOutput\ P\ (Y\ 0)\ =\ constantL2\ P\ a₀)}
\leanline{4}{:\ ∃\ S\ :\ FrozenPosteriorSolution\ P}
\leanline{8}{(Cloud.meanForce\_regular\ hK\ U\ m\ L\ hdiff\ hsc\ hLip)\ B\ hB,}
\leanline{6}{S.path\ 0\ =\ packLp\ P\ (optimizerParameter\ hK\ U\ m\ hm\ hdiff\ hsc)\ initial\ ∧}
\leanline{6}{(∀\ t,\ 0\ ≤\ t\ →\ parameterOutput\ P\ (S.path\ t)\ =}
\leanline{8}{constantL2\ P\ (optimizerParameter\ hK\ U\ m\ hm\ hdiff\ hsc))\ ∧}
\leanline{6}{(∀\ t,\ 0\ ≤\ t\ →\ ∀\ i,\ P.map\ (fun\ ω\ =>\ physicalParticleCLM\ i\ (S.path\ t\ ω))\ =}
\leanline{8}{RegularEvidence.posterior\ volume\ (latentPotential\ U\ i)}
\leanline{10}{(optimizerParameter\ hK\ U\ m\ hm\ hdiff\ hsc))\ ∧}
\leanline{6}{(∀\ t,\ 0\ ≤\ t\ →\ ‖Y\ t\ -\ S.path\ t‖\ ≤}
\leanline{8}{Real.exp\ (-m\ *\ t)\ *\ ‖Y\ 0\ -\ packLp\ P\ (optimizerParameter\ hK\ U\ m\ hm\ hdiff\ hsc)\ initial‖)\ ∧}
\leanline{6}{∀\ t,\ 0\ ≤\ t\ →\ Real.sqrt}
\leanline{8}{(∫\ ω,\ empiricalProductDistanceSq\ (Y\ t\ ω)\ (S.path\ t\ ω)\ ∂P)\ ≤}
\leanline{10}{Real.exp\ (-m\ *\ t)\ *\ ‖Y\ 0\ -\ packLp\ P\ (optimizerParameter\ hK\ U\ m\ hm\ hdiff\ hsc)\ initial‖}
\end{leancode}
\leansource{\texttt{empiricalProduct\_\allowbreak{}rms\_\allowbreak{}le\_\allowbreak{}l2} \quad LatentDiffusion/\allowbreak{}RegularSourceTheorem.lean, lines 404–425}
\begin{leancode}
\leanline{0}{theorem\ empiricalProduct\_rms\_le\_l2\ \{M\ N\ :\ ℕ\}\ (hM\ :\ 0\ <\ M)\ (hN\ :\ 0\ <\ N)}
\leanline{4}{(X\ Y\ :\ Lp\ (Cloud.State\ (M\ *\ N)\ dθ\ dz)\ 2\ P)\ :}
\leanline{4}{Real.sqrt\ (∫\ ω,\ empiricalProductDistanceSq\ (X\ ω)\ (Y\ ω)\ ∂P)\ ≤\ ‖X\ -\ Y‖}
\end{leancode}
\leansource{\texttt{optimizer\_\allowbreak{}replication} \quad LatentDiffusion/\allowbreak{}RegularSourceTheorem.lean, lines 250–262}
\begin{leancode}
\leanline{0}{theorem\ optimizer\_replication\ \{M\ N\ :\ ℕ\}\ (hM\ :\ 0\ <\ M)\ (hN\ :\ 0\ <\ N)}
\leanline{4}{(U\ :\ Fin\ M\ →\ Cloud.Joint\ dθ\ dz\ →\ ℝ)\ (m\ :\ ℝ)\ (hm\ :\ 0\ <\ m)}
\leanline{4}{(hdiff\ :\ ∀\ i,\ Differentiable\ ℝ\ (U\ i))}
\leanline{4}{(hsc\ :\ ∀\ i,\ StrongConvexOn\ Set.univ\ m\ (U\ i))\ :}
\leanline{4}{optimizerParameter\ (Nat.mul\_pos\ hM\ hN)\ (replicatePotentials\ (N\ :=\ N)\ U)\ m\ hm}
\leanline{6}{(fun\ i\ =>\ hdiff\ i.divNat)\ (fun\ i\ =>\ hsc\ i.divNat)\ =}
\leanline{8}{optimizerParameter\ hM\ U\ m\ hm\ hdiff\ hsc}
\end{leancode}
\end{leanbox}
\subsection{\texorpdfstring{Uniform finite-particle approximation error}{Uniform finite-particle approximation error}}


Couple the interacting particle system and the
        mean-field copies with the same Brownian motions and matching initial
        states, so \(\xi_0^N=0\). Let


        \[\xi_t^N=\|u_t^N-u_t\|^2+
          \frac1{MN}\sum_{m,n}\|\bar Z_t^{m,n}-Z_t^{m,n}\|^2.\]


Under A3/\allowbreak{}A4 and a uniform centered-force moment bound supplying a
        finite constant \(C_{\mathrm{sp}}\), the corrected source estimate is


        \[(\mathbb E\xi_t^N)^{1/2}\le
          (1-e^{-\lambda t})\frac{C_{\mathrm{sp}}}{\lambda\sqrt N}.\]


With a per-copy centered score variance bound \(V_0\) over the
        \(K=MN\) copies, the checked constant is
        \(C_{\mathrm{sp}}=\sqrt{V_0/M}\), so the bound equals
        \((1-e^{-\lambda t})\sqrt{V_0/K}/\lambda\).
        It is \(O(N^{-1/2})\) uniformly in time when the initial coupling
        error is zero. Independence from \(M\) requires a moment bound uniform
        in the data, or a bounded averaged replacement.



\textbf{Correction.} The statement omits the matched initialization that the proof uses at supplement line 837. The heterogeneous initial laws and force moments must be specified; a single undeclared \(Z_0\) does not give a uniform data-independent constant. Any origin-based moment expression must also control the drift's offset or be centered about an equilibrium.

\begin{leanbox}{Lean: \texttt{exists\_\allowbreak{}spatial\_\allowbreak{}comparison}}
\leansource{\texttt{exists\_\allowbreak{}spatial\_\allowbreak{}comparison} \quad LatentDiffusion/\allowbreak{}RegularMeanField.lean, lines 189–234}
\begin{leancode}
\leanline{0}{theorem\ exists\_spatial\_comparison}
\leanline{4}{(hK\ :\ 0\ <\ K)\ (U\ :\ Fin\ K\ →\ Cloud.Joint\ dθ\ dz\ →\ ℝ)}
\leanline{4}{(m\ :\ ℝ)\ (L\ :\ ℝ≥0)\ (hm\ :\ 0\ <\ m)\ (hdiff\ :\ ∀\ i,\ Differentiable\ ℝ\ (U\ i))}
\leanline{4}{(hsc\ :\ ∀\ i,\ StrongConvexOn\ Set.univ\ m\ (U\ i))}
\leanline{4}{(hLip\ :\ ∀\ i,\ LipschitzWith\ L\ (gradient\ (U\ i)))}
\leanline{4}{(Y\ :\ ℝ\ →\ Lp\ (Cloud.State\ K\ dθ\ dz)\ 2\ P)\ (hY\ :\ Continuous\ Y)}
\leanline{4}{(B\ :\ ℝ\ →\ Lp\ (Cloud.State\ K\ dθ\ dz)\ 2\ P)\ (hB\ :\ Continuous\ B)\ (hB0\ :\ B\ 0=0)}
\leanline{4}{(hYe\ :\ ∀\ t,\ 0\ ≤\ t\ →\ Y\ t-Y\ 0\ =}
\leanline{6}{-(∫\ r\ in\ 0..t,\ force\ P\ (Cloud.meanForce\_regular\ hK\ U\ m\ L\ hdiff\ hsc\ hLip)\ (Y\ r))+B\ t)}
\leanline{4}{(V\ :\ ℝ)\ (\_hV\ :\ 0\ ≤\ V)}
\leanline{4}{(hind\ :\ ∀\ t,\ 0\ ≤\ t\ →\ iIndepFun\ (fun\ i\ ω\ =>\ scoreLp\ P\ U\ L\ hLip\ (Y\ t)\ i\ ω)\ P)}
\leanline{4}{(hvar\ :\ ∀\ t,\ 0\ ≤\ t\ →\ ∀\ i,\ ‖scoreLp\ P\ U\ L\ hLip\ (Y\ t)\ i\ -}
\leanline{6}{constantL2\ P\ (∫\ ω,\ scoreLp\ P\ U\ L\ hLip\ (Y\ t)\ i\ ω\ ∂P)‖\textasciicircum{}2\ ≤\ V)\ :}
\leanline{4}{∃\ X\ :\ ℝ\ →\ Lp\ (Cloud.State\ K\ dθ\ dz)\ 2\ P,\ Continuous\ X\ ∧\ X\ 0=Y\ 0\ ∧}
\leanline{6}{(∀\ t,\ 0\ ≤\ t\ →\ X\ t-X\ 0\ =}
\leanline{8}{-(∫\ r\ in\ 0..t,\ particleForce\ P\ hK\ U\ m\ L\ hdiff\ hsc\ hLip\ (X\ r))+B\ t)\ ∧}
\leanline{6}{∀\ t,\ 0\ ≤\ t\ →\ ‖X\ t-Y\ t‖\ ≤\ Real.sqrt\ (V/K)/m*(1-Real.exp\ (-m*t))}
\end{leancode}
\end{leanbox}
\subsection{\texorpdfstring{Uniform Euler–Maruyama time-discretization error}{Uniform Euler–Maruyama time-discretization error}}


Use matching initial states, Brownian increments
        indexed by both \((m,n)\), and uniform normalized moment controls.
        Under A3/\allowbreak{}A4 and the explicit stable-step restriction
        \(0\lt h,\quad\lambda h\le1,\quad hL^2\le\lambda\), the corrected order asserted by the source's
        main theorem is


        \[d_{\mathrm{RMS}}((u_k^{N,h},\bar Q_k^{N,h}),
                    (u_{kh}^{N},\bar Q_{kh}^{N}))\le C\sqrt h,
          \qquad k\in\mathbb N,\]


where \(C\) is independent of \(k,N\) provided the moment controls
        and the normalized drift-at-origin bound are uniform in \(N\).
        The explicit drift radius is \(G=\|F(0)\|+LR\), where
        \(\sup_t\|X_t\|_{L^2}\le R\).



\textbf{Correction.} The printed lemma simultaneously writes
        \(C(h)=O(\sqrt h)\) and error \(\le\sqrt h\,C(h)\)
        (supplement, lines 847–850), which would assert \(O(h)\). Retain one
        square root: either \(C(h)=O(1)\) in that display or error
        \(\le C(h)=O(\sqrt h)\). The noise \(W_k^n\) at line 845 must retain
        the datum index \(m\), and the heterogeneous drift keeps its superscript
        \(\widetilde\ell^m\). The checked force-level proof uses the stated
        stronger stable-step restriction; it does not claim the full printed
        step-size interval without the corresponding sharper gradient estimate.

\begin{leanbox}{Lean: \texttt{temporal\_\allowbreak{}comparison}}
\leansource{\texttt{temporal\_\allowbreak{}comparison} \quad LatentDiffusion/\allowbreak{}RegularMeanField.lean, lines 253–301}
\begin{leancode}
\leanline{0}{theorem\ temporal\_comparison}
\leanline{4}{(hK\ :\ 0\ <\ K)\ (U\ :\ Fin\ K\ →\ Cloud.Joint\ dθ\ dz\ →\ ℝ)}
\leanline{4}{(m\ h\ R\ :\ ℝ)\ (L\ :\ ℝ≥0)\ (hR\ :\ 0\ ≤\ R)}
\leanline{4}{(hdiff\ :\ ∀\ i,\ Differentiable\ ℝ\ (U\ i))}
\leanline{4}{(hsc\ :\ ∀\ i,\ StrongConvexOn\ Set.univ\ m\ (U\ i))}
\leanline{4}{(hLip\ :\ ∀\ i,\ LipschitzWith\ L\ (gradient\ (U\ i)))\ (hs\ :\ StableStep\ m\ L\ h)}
\leanline{4}{(B\ :\ (Fin\ K\ ×\ Fin\ dz)\ →\ ℝ≥0\ →\ Ω\ →\ ℝ)}
\leanline{4}{(hB\ :\ ∀\ i,\ IsPreBrownianReal\ (B\ i)\ P)}
\leanline{4}{(X\ :\ ℝ\ →\ Lp\ (Cloud.State\ K\ dθ\ dz)\ 2\ P)\ (hX\ :\ Continuous\ X)}
\leanline{4}{(hXe\ :\ ∀\ t,\ 0\ ≤\ t\ →\ X\ t-X\ 0\ =}
\leanline{6}{-(∫\ r\ in\ 0..t,\ particleForce\ P\ hK\ U\ m\ L\ hdiff\ hsc\ hLip\ (X\ r))+physicalNoise\ P\ B\ hB\ t)}
\leanline{4}{(hMoment\ :\ ∀\ t,\ 0\ ≤\ t\ →\ ‖X\ t‖\ ≤\ R)\ (n\ :\ ℕ)\ :}
\leanline{4}{‖eulerTrajectory\ (particleForce\ P\ hK\ U\ m\ L\ hdiff\ hsc\ hLip)\ h}
\leanline{6}{(fun\ k\ =>\ physicalNoise\ P\ B\ hB\ ((k\ :\ ℝ)*h+h)-physicalNoise\ P\ B\ hB\ ((k\ :\ ℝ)*h))\ (X\ 0)\ n-}
\leanline{6}{X\ ((n\ :\ ℝ)*h)‖\ ≤}
\leanline{6}{(L\ :\ ℝ)/m\ *\ (h*(‖particleForce\ P\ hK\ U\ m\ L\ hdiff\ hsc\ hLip\ 0‖+(L\ :\ ℝ)*R)+}
\leanline{8}{(4/3\ :\ ℝ)*Real.sqrt\ (2*(dz\ :\ ℝ))*Real.sqrt\ h)}
\end{leancode}
\end{leanbox}
\subsection{\texorpdfstring{Appendix full-batch parameter error theorem}{Appendix full-batch parameter error theorem}}


This is the appendix repetition of the main
        full-batch error theorem, with \(\widetilde\theta=(\theta,\phi)\):


        \[\bigl(\mathbb E\|\widetilde\theta_T^{N,h}
                    -\widetilde\theta_\star\|^2\bigr)^{1/2}
          \le C_{h,N}+C\rho^T,\qquad
          C_{h,N}=O(\sqrt h+N^{-1/2}).\]


Its proof combines continuous convergence, spatial approximation and
        temporal approximation by the triangle inequality. It retains the
        initialization, coupling, integrability and moment requirements of those
        component estimates.



\textbf{Correction.} Use the corrected temporal order, empirical coupling and initialization hypotheses above. The source theorem concerns exact full-batch gradients; subsampled, reweighted, annealed or adaptive updates require a separate theorem for their actual recursions.

\begin{leanbox}{Lean: \texttt{full\_\allowbreak{}batch\_\allowbreak{}optimizer\_\allowbreak{}rms\_\allowbreak{}rate}, \texttt{optimizer\_\allowbreak{}replication}}
\leansource{\texttt{full\_\allowbreak{}batch\_\allowbreak{}optimizer\_\allowbreak{}rms\_\allowbreak{}rate} \quad LatentDiffusion/\allowbreak{}RegularSourceTheorem.lean, lines 167–209}
\begin{leancode}
\leanline{0}{theorem\ full\_batch\_optimizer\_rms\_rate}
\leanline{4}{\{M\ N\ :\ ℕ\}\ (hMN\ :\ 0\ <\ M*N)}
\leanline{4}{(U\ :\ Fin\ (M*N)\ →\ Cloud.Joint\ dθ\ dz\ →\ ℝ)}
\leanline{4}{(m\ h\ :\ ℝ)\ (L\ :\ ℝ≥0)\ (hm\ :\ 0\ <\ m)}
\leanline{4}{(hdiff\ :\ ∀\ i,\ Differentiable\ ℝ\ (U\ i))}
\leanline{4}{(hsc\ :\ ∀\ i,\ StrongConvexOn\ Set.univ\ m\ (U\ i))}
\leanline{4}{(hLip\ :\ ∀\ i,\ LipschitzWith\ L\ (gradient\ (U\ i)))\ (hs\ :\ StableStep\ m\ L\ h)}
\leanline{4}{(initial\ :\ Lp\ (Latent\ (M*N)\ dz)\ 2\ P)}
\leanline{4}{(hinitial\ :\ P.map\ initial\ =\ ULA.gibbsMeasure}
\leanline{6}{(frozenPotential\ U\ (optimizerParameter\ hMN\ U\ m\ hm\ hdiff\ hsc)))}
\leanline{4}{(hposterior\ :\ ∀\ i,\ P.map\ (fun\ ω\ =>\ physicalParticleCLM\ i}
\leanline{8}{((packLp\ P\ (optimizerParameter\ hMN\ U\ m\ hm\ hdiff\ hsc)\ initial)\ ω))\ =}
\leanline{6}{RegularEvidence.posterior\ volume\ (latentPotential\ U\ i)\ (optimizerParameter\ hMN\ U\ m\ hm\ hdiff\ hsc))}
\leanline{4}{(data\ :\ ∀\ i,\ RegularEvidence.LocalNormalizerRegularity\ volume\ (latentPotential\ U\ i)}
\leanline{6}{(optimizerParameter\ hMN\ U\ m\ hm\ hdiff\ hsc))}
\leanline{4}{(hindScores\ :\ iIndepFun\ (fun\ i\ ω\ =>\ scoreLp\ P\ U\ L\ hLip}
\leanline{6}{(packLp\ P\ (optimizerParameter\ hMN\ U\ m\ hm\ hdiff\ hsc)\ initial)\ i\ ω)\ P)}
\leanline{4}{(B\ :\ (Fin\ (M*N)\ ×\ Fin\ dz)\ →\ ℝ≥0\ →\ Ω\ →\ ℝ)}
\leanline{4}{(hB\ :\ ∀\ i,\ IsPreBrownianReal\ (B\ i)\ P)}
\leanline{4}{(hscalar\ :\ ∀\ s\ :\ ℝ,\ 0\ <\ s\ →\ iIndepFun\ (fun\ p\ :\ ℕ\ ×\ (Fin\ (M*N)\ ×\ Fin\ dz)\ =>\ fun\ ω\ =>}
\leanline{6}{ULA.Hilbert.brownianEulerNoise\ P\ B\ hB\ s\ p.1\ ω\ p.2)\ P)}
\leanline{4}{(hind\ :\ ∀\ s\ :\ ℝ,\ 0\ <\ s\ →\ IndepFun\ initial}
\leanline{6}{(fun\ ω\ n\ =>\ ULA.Hilbert.brownianEulerNoise\ P\ B\ hB\ s\ n\ ω)\ P)}
\leanline{4}{(V₀\ :\ ℝ)}
\leanline{4}{(hvar\ :\ ∀\ i,\ ‖scoreLp\ P\ U\ L\ hLip}
\leanline{8}{(packLp\ P\ (optimizerParameter\ hMN\ U\ m\ hm\ hdiff\ hsc)\ initial)\ i-}
\leanline{6}{constantL2\ P\ (∫\ ω,\ scoreLp\ P\ U\ L\ hLip}
\leanline{8}{(packLp\ P\ (optimizerParameter\ hMN\ U\ m\ hm\ hdiff\ hsc)\ initial)\ i\ ω\ ∂P)‖\textasciicircum{}2\ ≤\ V₀)}
\leanline{4}{(x₀\ :\ Lp\ (Cloud.State\ (M*N)\ dθ\ dz)\ 2\ P)\ (n\ :\ ℕ)\ :}
\leanline{4}{Real.sqrt\ (∫\ ω,\ ‖IPLA.parameterCLM}
\leanline{6}{(eulerTrajectory\ (Cloud.meanForce\ U)\ h}
\leanline{8}{(fun\ k\ =>\ noiseIncrement\ P\ B\ hB\ h\ k\ ω)\ (x₀\ ω)\ n)-}
\leanline{6}{optimizerParameter\ hMN\ U\ m\ hm\ hdiff\ hsc‖\textasciicircum{}2\ ∂P)\ ≤}
\leanline{6}{(1-m*h/2)\textasciicircum{}n\ *\ ‖x₀-packLp\ P\ (optimizerParameter\ hMN\ U\ m\ hm\ hdiff\ hsc)\ initial‖\ +}
\leanline{6}{(2/m)\ *\ Real.sqrt\ (V₀/((M*N\ :\ ℕ)\ :\ ℝ))\ +}
\leanline{6}{(L\ :\ ℝ)/m\ *\ (h*Real.sqrt\ ((L\ :\ ℝ)*(dz\ :\ ℝ))+}
\leanline{8}{(4/3\ :\ ℝ)*Real.sqrt\ (2*(dz\ :\ ℝ))*Real.sqrt\ h)}
\end{leancode}
\leansource{\texttt{optimizer\_\allowbreak{}replication} \quad LatentDiffusion/\allowbreak{}RegularSourceTheorem.lean, lines 250–262}
\begin{leancode}
\leanline{0}{theorem\ optimizer\_replication\ \{M\ N\ :\ ℕ\}\ (hM\ :\ 0\ <\ M)\ (hN\ :\ 0\ <\ N)}
\leanline{4}{(U\ :\ Fin\ M\ →\ Cloud.Joint\ dθ\ dz\ →\ ℝ)\ (m\ :\ ℝ)\ (hm\ :\ 0\ <\ m)}
\leanline{4}{(hdiff\ :\ ∀\ i,\ Differentiable\ ℝ\ (U\ i))}
\leanline{4}{(hsc\ :\ ∀\ i,\ StrongConvexOn\ Set.univ\ m\ (U\ i))\ :}
\leanline{4}{optimizerParameter\ (Nat.mul\_pos\ hM\ hN)\ (replicatePotentials\ (N\ :=\ N)\ U)\ m\ hm}
\leanline{6}{(fun\ i\ =>\ hdiff\ i.divNat)\ (fun\ i\ =>\ hsc\ i.divNat)\ =}
\leanline{8}{optimizerParameter\ hM\ U\ m\ hm\ hdiff\ hsc}
\end{leancode}
\end{leanbox}
\section{New assumptions}
\subsection{\texorpdfstring{N1 — Monotone, Lipschitz force}{N1 — Monotone, Lipschitz force}}



Let \(H\) be a real inner-product space. The positive force \(F\), equal to minus the drift, satisfies for every \(x,y\in H\)



\[\lambda\|x-y\|^2\le\langle F(x)-F(y),x-y\rangle,\qquad \|F(x)-F(y)\|\le L\|x-y\|,\quad L\ge0.\]



This is an explicit assumption on the complete force of the parameter and particle system in the chosen norm. Applying a per-datapoint hypothesis to that complete system requires its own lifting argument.

\begin{leanbox}{Lean: \texttt{ForceRegularity}}
\leansource{\texttt{ForceRegularity} \quad LatentDiffusion/\allowbreak{}Numerics.lean, lines 28–30}
\begin{leancode}
\leanline{0}{structure\ ForceRegularity\ (F\ :\ H\ →\ H)\ (m\ :\ ℝ)\ (L\ :\ ℝ≥0)\ :\ Prop\ where}
\leanline{2}{monotone\ :\ ∀\ x\ y,\ m\ *\ ‖x\ -\ y‖\ \textasciicircum{}\ 2\ ≤\ inner\ ℝ\ (F\ x\ -\ F\ y)\ (x\ -\ y)}
\leanline{2}{lipschitz\ :\ ∀\ x\ y,\ ‖F\ x\ -\ F\ y‖\ ≤\ (L\ :\ ℝ)\ *\ ‖x\ -\ y‖}
\end{leancode}
\end{leanbox}
\subsection{\texorpdfstring{N2 — Stable step size}{N2 — Stable step size}}



Require



\[\lambda>0,\qquad h>0,\qquad\lambda h\le1,\qquad hL^2\le\lambda.\]



For \(L>0\), the familiar sufficient condition is \(0<h\le\min\{1/\lambda,\lambda/L^2\}\). These restrictions depend on the complete force constants.

\begin{leanbox}{Lean: \texttt{StableStep}, \texttt{stableStep\_\allowbreak{}of\_\allowbreak{}min}}
\leansource{\texttt{StableStep} \quad LatentDiffusion/\allowbreak{}Numerics.lean, lines 33–37}
\begin{leancode}
\leanline{0}{structure\ StableStep\ (m\ :\ ℝ)\ (L\ :\ ℝ≥0)\ (h\ :\ ℝ)\ :\ Prop\ where}
\leanline{2}{positive\ :\ 0\ <\ h}
\leanline{2}{monotonicity\_positive\ :\ 0\ <\ m}
\leanline{2}{scaled\_monotonicity\ :\ m\ *\ h\ ≤\ 1}
\leanline{2}{scaled\_lipschitz\ :\ h\ *\ (L\ :\ ℝ)\ \textasciicircum{}\ 2\ ≤\ m}
\end{leancode}
\leansource{\texttt{stableStep\_\allowbreak{}of\_\allowbreak{}min} \quad LatentDiffusion/\allowbreak{}Numerics.lean, lines 82–91}
\begin{leancode}
\leanline{0}{theorem\ stableStep\_of\_min\ (m\ h\ :\ ℝ)\ (L\ :\ ℝ≥0)\ (hm\ :\ 0\ <\ m)}
\leanline{4}{(hL\ :\ 0\ <\ (L\ :\ ℝ))\ (hh\ :\ 0\ <\ h)}
\leanline{4}{(hsmall\ :\ h\ ≤\ min\ (1\ /\ m)\ (m\ /\ (L\ :\ ℝ)\ \textasciicircum{}\ 2))\ :\ StableStep\ m\ L\ h}
\end{leancode}
\end{leanbox}
\subsection{\texorpdfstring{N3 — Force error and reference consistency}{N3 — Force error and reference consistency}}



Write \(E_F(x,\xi)=x-hF(x)+\xi\), and let \(X_{n+1}=E_G(X_n,\xi_n)\) be the implemented scheme. A reference path \(y_n\) uses the same additive increments. For every step require



\[\|G(X_n)-F(X_n)\|\le B,\qquad \|y_{n+1}-E_F(y_n,\xi_n)\|\le D h\sqrt h,\qquad B,D\ge0.\]



The numerical theorem is pathwise. An application to an exact SDE skeleton must establish the displayed local defect bound. Brownian increments, stochastic independence, moment estimates and propagation of chaos are not inferred from this premise.

\begin{leanbox}{Lean: \texttt{trajectory\_\allowbreak{}error\_\allowbreak{}bound}, \texttt{trajectory\_\allowbreak{}error\_\allowbreak{}bound}}
\leansource{\texttt{trajectory\_\allowbreak{}error\_\allowbreak{}bound} \quad LatentDiffusion/\allowbreak{}Numerics.lean, lines 143–163}
\begin{leancode}
\leanline{0}{theorem\ trajectory\_error\_bound\ (F\ G\ :\ H\ →\ H)\ (m\ h\ B\ D\ :\ ℝ)\ (L\ :\ ℝ≥0)}
\leanline{4}{(hF\ :\ ForceRegularity\ F\ m\ L)\ (hs\ :\ StableStep\ m\ L\ h)\ (hB\ :\ 0\ ≤\ B)\ (hD\ :\ 0\ ≤\ D)}
\leanline{4}{(ξ\ :\ ℕ\ →\ H)\ (x₀\ :\ H)\ (y\ :\ ℕ\ →\ H)}
\leanline{4}{(hforce\ :\ ∀\ n,\ ‖G\ (eulerTrajectory\ G\ h\ ξ\ x₀\ n)\ -\ F\ (eulerTrajectory\ G\ h\ ξ\ x₀\ n)‖\ ≤\ B)}
\leanline{4}{(hdefect\ :\ ∀\ n,\ ‖y\ (n\ +\ 1)\ -\ eulerStep\ F\ h\ (y\ n)\ (ξ\ n)‖\ ≤\ D\ *\ h\ *\ Real.sqrt\ h)}
\leanline{4}{(n\ :\ ℕ)\ :}
\leanline{4}{‖eulerTrajectory\ G\ h\ ξ\ x₀\ n\ -\ y\ n‖\ ≤}
\leanline{6}{(1\ -\ m\ *\ h\ /\ 2)\ \textasciicircum{}\ n\ *\ ‖x₀\ -\ y\ 0‖\ +}
\leanline{8}{(2\ /\ m)\ *\ (B\ +\ D\ *\ Real.sqrt\ h)}
\end{leancode}
\leansource{\texttt{trajectory\_\allowbreak{}error\_\allowbreak{}bound} \quad LatentDiffusion/\allowbreak{}Numerics.lean, lines 143–163}
\begin{leancode}
\leanline{0}{theorem\ trajectory\_error\_bound\ (F\ G\ :\ H\ →\ H)\ (m\ h\ B\ D\ :\ ℝ)\ (L\ :\ ℝ≥0)}
\leanline{4}{(hF\ :\ ForceRegularity\ F\ m\ L)\ (hs\ :\ StableStep\ m\ L\ h)\ (hB\ :\ 0\ ≤\ B)\ (hD\ :\ 0\ ≤\ D)}
\leanline{4}{(ξ\ :\ ℕ\ →\ H)\ (x₀\ :\ H)\ (y\ :\ ℕ\ →\ H)}
\leanline{4}{(hforce\ :\ ∀\ n,\ ‖G\ (eulerTrajectory\ G\ h\ ξ\ x₀\ n)\ -\ F\ (eulerTrajectory\ G\ h\ ξ\ x₀\ n)‖\ ≤\ B)}
\leanline{4}{(hdefect\ :\ ∀\ n,\ ‖y\ (n\ +\ 1)\ -\ eulerStep\ F\ h\ (y\ n)\ (ξ\ n)‖\ ≤\ D\ *\ h\ *\ Real.sqrt\ h)}
\leanline{4}{(n\ :\ ℕ)\ :}
\leanline{4}{‖eulerTrajectory\ G\ h\ ξ\ x₀\ n\ -\ y\ n‖\ ≤}
\leanline{6}{(1\ -\ m\ *\ h\ /\ 2)\ \textasciicircum{}\ n\ *\ ‖x₀\ -\ y\ 0‖\ +}
\leanline{8}{(2\ /\ m)\ *\ (B\ +\ D\ *\ Real.sqrt\ h)}
\end{leancode}
\end{leanbox}
\subsection{\texorpdfstring{N4 — Primitive integral-path bounds}{N4 — Primitive integral-path bounds}}



Let \(H\) be complete. On each step, an exact segment \(Y_n(t)\), \(0\le t\le h\), starts at \(y_n\), ends at \(y_{n+1}\), and satisfies



\[Y_n(t)-Y_n(0)=-\int_0^t F(Y_n(s))\,\mathrm ds+W_n(t),\qquad W_n(h)=\xi_n.\]



Assume the drift is integrable on each segment and, uniformly in \(n\),



\[\|F(Y_n(t))\|\le G,\qquad\|W_n(t)\|\le N\sqrt t,\qquad G,N\ge0.\]



These are primitive integral and moment bounds. In an \(L^2\) realization, the norms are root mean squares. An application must supply the SDE reference and verify these uniform estimates. The pointwise-force lift below establishes the required \(L^2\) force monotonicity and Lipschitz bounds from N1.

\begin{leanbox}{Lean: \texttt{IntegralReference}}
\leansource{\texttt{IntegralReference} \quad LatentDiffusion/\allowbreak{}Numerics.lean, lines 273–282}
\begin{leancode}
\leanline{0}{structure\ IntegralReference\ (F\ :\ H\ →\ H)\ (h\ G\ N\ :\ ℝ)\ (y\ :\ ℕ\ →\ H)}
\leanline{4}{(ξ\ :\ ℕ\ →\ H)\ (Y\ W\ :\ ℕ\ →\ ℝ\ →\ H)\ :\ Prop\ where}
\leanline{2}{starts\ :\ ∀\ n,\ Y\ n\ 0\ =\ y\ n}
\leanline{2}{ends\ :\ ∀\ n,\ Y\ n\ h\ =\ y\ (n\ +\ 1)}
\leanline{2}{noise\_end\ :\ ∀\ n,\ W\ n\ h\ =\ ξ\ n}
\leanline{2}{integrable\ :\ ∀\ n,\ IntervalIntegrable\ (fun\ t\ ↦\ F\ (Y\ n\ t))\ volume\ 0\ h}
\leanline{2}{drift\_radius\ :\ ∀\ n\ t,\ t\ ∈\ Set.Icc\ 0\ h\ →\ ‖F\ (Y\ n\ t)‖\ ≤\ G}
\leanline{2}{noise\_growth\ :\ ∀\ n\ t,\ t\ ∈\ Set.Icc\ 0\ h\ →\ ‖W\ n\ t‖\ ≤\ N\ *\ Real.sqrt\ t}
\leanline{2}{equation\ :\ ∀\ n\ t,\ t\ ∈\ Set.Icc\ 0\ h\ →}
\leanline{4}{Y\ n\ t\ -\ Y\ n\ 0\ =\ -(∫\ s\ in\ 0..t,\ F\ (Y\ n\ s))\ +\ W\ n\ t}
\end{leancode}
\end{leanbox}
\subsection{\texorpdfstring{N5 — Centered stochastic force innovations}{N5 — Centered stochastic force innovations}}



Compare \(X_{n+1}=E_F(X_n,\xi_n)+h\zeta_n\) with \(Y_{n+1}=E_F(Y_n,\xi_n)\), and put \(A_n=E_F(X_n,\xi_n)-E_F(Y_n,\xi_n)\). In the chosen Hilbert space require



\[\langle A_n,\zeta_n\rangle=0,\qquad\|\zeta_n\|^2\le\sigma^2\quad(n\ge0).\]



In \(L^2(\Omega,P;E)\), these are precisely a zero expected cross term and a second-moment bound. Conditional centering given the states and shared increment suffices when justified. The premise states the needed orthogonality explicitly.

\begin{leanbox}{Lean: \texttt{HilbertInnovationControl}}
\leansource{\texttt{HilbertInnovationControl} \quad LatentDiffusion/\allowbreak{}Numerics.lean, lines 348–353}
\begin{leancode}
\leanline{0}{structure\ HilbertInnovationControl\ (F\ :\ H\ →\ H)\ (h\ σ\ :\ ℝ)}
\leanline{4}{(ξ\ ζ\ :\ ℕ\ →\ H)\ (x₀\ y₀\ :\ H)\ :\ Prop\ where}
\leanline{2}{orthogonal\ :\ ∀\ n,\ inner\ ℝ}
\leanline{4}{(eulerStep\ F\ h\ (eulerTrajectory\ F\ h\ (fun\ k\ ↦\ ξ\ k\ +\ h\ •\ ζ\ k)\ x₀\ n)\ (ξ\ n)\ -}
\leanline{6}{eulerStep\ F\ h\ (eulerTrajectory\ F\ h\ ξ\ y₀\ n)\ (ξ\ n))\ (ζ\ n)\ =\ 0}
\leanline{2}{variance\ :\ ∀\ n,\ ‖ζ\ n‖\ \textasciicircum{}\ 2\ ≤\ σ\ \textasciicircum{}\ 2}
\end{leancode}
\end{leanbox}
\subsection{\texorpdfstring{N6 — Frozen reference and empirical-force residual}{N6 — Frozen reference and empirical-force residual}}



In the genuine space \(L^2(\Omega;\text{scaled cloud})\), a reference force \(R\) is \(L\)-Lipschitz and admits actual segments satisfying N4, with nonnegative drift radius \(G\) and noise bound \(W\sqrt t\). It need not contract in the fixed parameter coordinate.




The actual physical parameter projection \(C\) has norm at most one and satisfies \(Cy_n=a_\star\) almost surely for every \(n\). The norms in N4 are root mean squares. The construction and stationarity of a frozen posterior process, and its uniform drift/\allowbreak{}noise moment bounds, must be supplied by the application. The theorem accepts any such fixed reference parameter. An inference guarantee to the paper's tilted maximizer additionally identifies this \(a_\star\) with the optimizer proved in Part I.

\begin{leanbox}{Lean: \texttt{particle\_\allowbreak{}oracle\_\allowbreak{}parameter\_\allowbreak{}bound}, \texttt{particle\_\allowbreak{}oracle\_\allowbreak{}parameter\_\allowbreak{}bound}, \texttt{particle\_\allowbreak{}oracle\_\allowbreak{}parameter\_\allowbreak{}bound}}
\leansource{\texttt{particle\_\allowbreak{}oracle\_\allowbreak{}parameter\_\allowbreak{}bound} \quad LatentDiffusion/\allowbreak{}MainResults.lean, lines 47–85}
\begin{leancode}
\leanline{0}{theorem\ particle\_oracle\_parameter\_bound}
\leanline{4}{(hK\ :\ 0\ <\ K)\ (U\ :\ Fin\ K\ →\ Cloud.Joint\ dθ\ dz\ →\ ℝ)}
\leanline{4}{(m\ h\ G\ Wbound\ V\ σ\ :\ ℝ)\ (L\ :\ ℝ≥0)}
\leanline{4}{(hdiff\ :\ ∀\ i,\ Differentiable\ ℝ\ (U\ i))}
\leanline{4}{(hsc\ :\ ∀\ i,\ StrongConvexOn\ Set.univ\ m\ (U\ i))}
\leanline{4}{(hLip\ :\ ∀\ i,\ LipschitzWith\ L\ (gradient\ (U\ i)))}
\leanline{4}{(hs\ :\ StableStep\ m\ L\ h)\ (hG\ :\ 0\ ≤\ G)\ (hW\ :\ 0\ ≤\ Wbound)\ (\_hV\ :\ 0\ ≤\ V)}
\leanline{4}{(R\ :\ Lp\ (Cloud.State\ K\ dθ\ dz)\ 2\ P\ →\ Lp\ (Cloud.State\ K\ dθ\ dz)\ 2\ P)}
\leanline{4}{(hR\ :\ LipschitzWith\ L\ R)}
\leanline{4}{(y\ ξ\ ζ\ :\ ℕ\ →\ Lp\ (Cloud.State\ K\ dθ\ dz)\ 2\ P)}
\leanline{4}{(Y\ W\ :\ ℕ\ →\ ℝ\ →\ Lp\ (Cloud.State\ K\ dθ\ dz)\ 2\ P)}
\leanline{4}{(x₀\ :\ Lp\ (Cloud.State\ K\ dθ\ dz)\ 2\ P)\ (astar\ :\ EuclideanSpace\ ℝ\ (Fin\ dθ))}
\leanline{4}{(hr\ :\ IntegralReference\ R\ h\ G\ Wbound\ y\ ξ\ Y\ W)}
\leanline{4}{(fluctuation\ :\ ℕ\ →\ Fin\ K\ →\ Lp\ (Cloud.State\ K\ dθ\ dz)\ 2\ P)}
\leanline{4}{(hres\ :\ ∀\ n,\ particleForce\ P\ hK\ U\ m\ L\ hdiff\ hsc\ hLip\ (y\ n)\ -\ R\ (y\ n)\ =}
\leanline{6}{MiniBatchVariance.average\ (fluctuation\ n))}
\leanline{4}{(horth\ :\ ∀\ n\ i\ j,\ i\ ≠\ j\ →\ inner\ ℝ\ (fluctuation\ n\ i)\ (fluctuation\ n\ j)\ =\ 0)}
\leanline{4}{(hvariance\ :\ ∀\ n\ i,\ ‖fluctuation\ n\ i‖\ \textasciicircum{}\ 2\ ≤\ V)}
\leanline{4}{(horacle\ :\ HilbertInnovationControl\ (particleForce\ P\ hK\ U\ m\ L\ hdiff\ hsc\ hLip)}
\leanline{6}{h\ σ\ ξ\ ζ\ x₀\ x₀)}
\leanline{4}{(htarget\ :\ ∀\ n,\ parameterOutput\ P\ (y\ n)\ =\ constantL2\ P\ astar)\ (n\ :\ ℕ)\ :}
\leanline{4}{‖parameterOutput\ P}
\leanline{8}{(eulerTrajectory\ (particleForce\ P\ hK\ U\ m\ L\ hdiff\ hsc\ hLip)\ h}
\leanline{10}{(fun\ k\ =>\ ξ\ k\ +\ h\ •\ ζ\ k)\ x₀\ n)\ -\ constantL2\ P\ astar‖\ ≤}
\leanline{6}{(1\ -\ m\ *\ h\ /\ 2)\ \textasciicircum{}\ n\ *\ ‖x₀\ -\ y\ 0‖\ +}
\leanline{8}{(2\ /\ m)\ *\ Real.sqrt\ (V\ /\ K)\ +}
\leanline{8}{(L\ :\ ℝ)\ /\ m\ *\ (h\ *\ G\ +\ (4\ /\ 3\ :\ ℝ)\ *\ Wbound\ *\ Real.sqrt\ h)\ +}
\leanline{8}{Real.sqrt\ (h\ *\ σ\ \textasciicircum{}\ 2\ /\ m)}
\end{leancode}
\leansource{\texttt{particle\_\allowbreak{}oracle\_\allowbreak{}parameter\_\allowbreak{}bound} \quad LatentDiffusion/\allowbreak{}MainResults.lean, lines 47–85}
\begin{leancode}
\leanline{0}{theorem\ particle\_oracle\_parameter\_bound}
\leanline{4}{(hK\ :\ 0\ <\ K)\ (U\ :\ Fin\ K\ →\ Cloud.Joint\ dθ\ dz\ →\ ℝ)}
\leanline{4}{(m\ h\ G\ Wbound\ V\ σ\ :\ ℝ)\ (L\ :\ ℝ≥0)}
\leanline{4}{(hdiff\ :\ ∀\ i,\ Differentiable\ ℝ\ (U\ i))}
\leanline{4}{(hsc\ :\ ∀\ i,\ StrongConvexOn\ Set.univ\ m\ (U\ i))}
\leanline{4}{(hLip\ :\ ∀\ i,\ LipschitzWith\ L\ (gradient\ (U\ i)))}
\leanline{4}{(hs\ :\ StableStep\ m\ L\ h)\ (hG\ :\ 0\ ≤\ G)\ (hW\ :\ 0\ ≤\ Wbound)\ (\_hV\ :\ 0\ ≤\ V)}
\leanline{4}{(R\ :\ Lp\ (Cloud.State\ K\ dθ\ dz)\ 2\ P\ →\ Lp\ (Cloud.State\ K\ dθ\ dz)\ 2\ P)}
\leanline{4}{(hR\ :\ LipschitzWith\ L\ R)}
\leanline{4}{(y\ ξ\ ζ\ :\ ℕ\ →\ Lp\ (Cloud.State\ K\ dθ\ dz)\ 2\ P)}
\leanline{4}{(Y\ W\ :\ ℕ\ →\ ℝ\ →\ Lp\ (Cloud.State\ K\ dθ\ dz)\ 2\ P)}
\leanline{4}{(x₀\ :\ Lp\ (Cloud.State\ K\ dθ\ dz)\ 2\ P)\ (astar\ :\ EuclideanSpace\ ℝ\ (Fin\ dθ))}
\leanline{4}{(hr\ :\ IntegralReference\ R\ h\ G\ Wbound\ y\ ξ\ Y\ W)}
\leanline{4}{(fluctuation\ :\ ℕ\ →\ Fin\ K\ →\ Lp\ (Cloud.State\ K\ dθ\ dz)\ 2\ P)}
\leanline{4}{(hres\ :\ ∀\ n,\ particleForce\ P\ hK\ U\ m\ L\ hdiff\ hsc\ hLip\ (y\ n)\ -\ R\ (y\ n)\ =}
\leanline{6}{MiniBatchVariance.average\ (fluctuation\ n))}
\leanline{4}{(horth\ :\ ∀\ n\ i\ j,\ i\ ≠\ j\ →\ inner\ ℝ\ (fluctuation\ n\ i)\ (fluctuation\ n\ j)\ =\ 0)}
\leanline{4}{(hvariance\ :\ ∀\ n\ i,\ ‖fluctuation\ n\ i‖\ \textasciicircum{}\ 2\ ≤\ V)}
\leanline{4}{(horacle\ :\ HilbertInnovationControl\ (particleForce\ P\ hK\ U\ m\ L\ hdiff\ hsc\ hLip)}
\leanline{6}{h\ σ\ ξ\ ζ\ x₀\ x₀)}
\leanline{4}{(htarget\ :\ ∀\ n,\ parameterOutput\ P\ (y\ n)\ =\ constantL2\ P\ astar)\ (n\ :\ ℕ)\ :}
\leanline{4}{‖parameterOutput\ P}
\leanline{8}{(eulerTrajectory\ (particleForce\ P\ hK\ U\ m\ L\ hdiff\ hsc\ hLip)\ h}
\leanline{10}{(fun\ k\ =>\ ξ\ k\ +\ h\ •\ ζ\ k)\ x₀\ n)\ -\ constantL2\ P\ astar‖\ ≤}
\leanline{6}{(1\ -\ m\ *\ h\ /\ 2)\ \textasciicircum{}\ n\ *\ ‖x₀\ -\ y\ 0‖\ +}
\leanline{8}{(2\ /\ m)\ *\ Real.sqrt\ (V\ /\ K)\ +}
\leanline{8}{(L\ :\ ℝ)\ /\ m\ *\ (h\ *\ G\ +\ (4\ /\ 3\ :\ ℝ)\ *\ Wbound\ *\ Real.sqrt\ h)\ +}
\leanline{8}{Real.sqrt\ (h\ *\ σ\ \textasciicircum{}\ 2\ /\ m)}
\end{leancode}
\leansource{\texttt{particle\_\allowbreak{}oracle\_\allowbreak{}parameter\_\allowbreak{}bound} \quad LatentDiffusion/\allowbreak{}MainResults.lean, lines 47–85}
\begin{leancode}
\leanline{0}{theorem\ particle\_oracle\_parameter\_bound}
\leanline{4}{(hK\ :\ 0\ <\ K)\ (U\ :\ Fin\ K\ →\ Cloud.Joint\ dθ\ dz\ →\ ℝ)}
\leanline{4}{(m\ h\ G\ Wbound\ V\ σ\ :\ ℝ)\ (L\ :\ ℝ≥0)}
\leanline{4}{(hdiff\ :\ ∀\ i,\ Differentiable\ ℝ\ (U\ i))}
\leanline{4}{(hsc\ :\ ∀\ i,\ StrongConvexOn\ Set.univ\ m\ (U\ i))}
\leanline{4}{(hLip\ :\ ∀\ i,\ LipschitzWith\ L\ (gradient\ (U\ i)))}
\leanline{4}{(hs\ :\ StableStep\ m\ L\ h)\ (hG\ :\ 0\ ≤\ G)\ (hW\ :\ 0\ ≤\ Wbound)\ (\_hV\ :\ 0\ ≤\ V)}
\leanline{4}{(R\ :\ Lp\ (Cloud.State\ K\ dθ\ dz)\ 2\ P\ →\ Lp\ (Cloud.State\ K\ dθ\ dz)\ 2\ P)}
\leanline{4}{(hR\ :\ LipschitzWith\ L\ R)}
\leanline{4}{(y\ ξ\ ζ\ :\ ℕ\ →\ Lp\ (Cloud.State\ K\ dθ\ dz)\ 2\ P)}
\leanline{4}{(Y\ W\ :\ ℕ\ →\ ℝ\ →\ Lp\ (Cloud.State\ K\ dθ\ dz)\ 2\ P)}
\leanline{4}{(x₀\ :\ Lp\ (Cloud.State\ K\ dθ\ dz)\ 2\ P)\ (astar\ :\ EuclideanSpace\ ℝ\ (Fin\ dθ))}
\leanline{4}{(hr\ :\ IntegralReference\ R\ h\ G\ Wbound\ y\ ξ\ Y\ W)}
\leanline{4}{(fluctuation\ :\ ℕ\ →\ Fin\ K\ →\ Lp\ (Cloud.State\ K\ dθ\ dz)\ 2\ P)}
\leanline{4}{(hres\ :\ ∀\ n,\ particleForce\ P\ hK\ U\ m\ L\ hdiff\ hsc\ hLip\ (y\ n)\ -\ R\ (y\ n)\ =}
\leanline{6}{MiniBatchVariance.average\ (fluctuation\ n))}
\leanline{4}{(horth\ :\ ∀\ n\ i\ j,\ i\ ≠\ j\ →\ inner\ ℝ\ (fluctuation\ n\ i)\ (fluctuation\ n\ j)\ =\ 0)}
\leanline{4}{(hvariance\ :\ ∀\ n\ i,\ ‖fluctuation\ n\ i‖\ \textasciicircum{}\ 2\ ≤\ V)}
\leanline{4}{(horacle\ :\ HilbertInnovationControl\ (particleForce\ P\ hK\ U\ m\ L\ hdiff\ hsc\ hLip)}
\leanline{6}{h\ σ\ ξ\ ζ\ x₀\ x₀)}
\leanline{4}{(htarget\ :\ ∀\ n,\ parameterOutput\ P\ (y\ n)\ =\ constantL2\ P\ astar)\ (n\ :\ ℕ)\ :}
\leanline{4}{‖parameterOutput\ P}
\leanline{8}{(eulerTrajectory\ (particleForce\ P\ hK\ U\ m\ L\ hdiff\ hsc\ hLip)\ h}
\leanline{10}{(fun\ k\ =>\ ξ\ k\ +\ h\ •\ ζ\ k)\ x₀\ n)\ -\ constantL2\ P\ astar‖\ ≤}
\leanline{6}{(1\ -\ m\ *\ h\ /\ 2)\ \textasciicircum{}\ n\ *\ ‖x₀\ -\ y\ 0‖\ +}
\leanline{8}{(2\ /\ m)\ *\ Real.sqrt\ (V\ /\ K)\ +}
\leanline{8}{(L\ :\ ℝ)\ /\ m\ *\ (h\ *\ G\ +\ (4\ /\ 3\ :\ ℝ)\ *\ Wbound\ *\ Real.sqrt\ h)\ +}
\leanline{8}{Real.sqrt\ (h\ *\ σ\ \textasciicircum{}\ 2\ /\ m)}
\end{leancode}
\end{leanbox}
\subsection{\texorpdfstring{N7 — Primitive centered-copy residual and variance}{N7 — Primitive centered-copy residual and variance}}



Let \(K=MN>0\). The difference between the actual interacting force \(F_K\) and the reference force along \(y_n\) is represented by actual square-integrable fluctuations:



\[F_K(y_n)-R(y_n)=\frac1K\sum_{i=1}^K g_{n,i},\quad\langle g_{n,i},g_{n,j}\rangle_{L^2}=0\ (i\ne j),\quad\|g_{n,i}\|_{L^2}^2\le V,\quad V\ge0.\]



These are a residual identity and primitive cross-moment and variance conditions. They permit heterogeneous data. A posterior-cloud application must prove that its chosen independent-copy reference has this representation and these moments. The \(K^{-1/2}\) force-residual bound is derived below.

\begin{leanbox}{Lean: \texttt{particle\_\allowbreak{}oracle\_\allowbreak{}parameter\_\allowbreak{}bound}, \texttt{particle\_\allowbreak{}oracle\_\allowbreak{}parameter\_\allowbreak{}bound}, \texttt{particle\_\allowbreak{}oracle\_\allowbreak{}parameter\_\allowbreak{}bound}}
\leansource{\texttt{particle\_\allowbreak{}oracle\_\allowbreak{}parameter\_\allowbreak{}bound} \quad LatentDiffusion/\allowbreak{}MainResults.lean, lines 47–85}
\begin{leancode}
\leanline{0}{theorem\ particle\_oracle\_parameter\_bound}
\leanline{4}{(hK\ :\ 0\ <\ K)\ (U\ :\ Fin\ K\ →\ Cloud.Joint\ dθ\ dz\ →\ ℝ)}
\leanline{4}{(m\ h\ G\ Wbound\ V\ σ\ :\ ℝ)\ (L\ :\ ℝ≥0)}
\leanline{4}{(hdiff\ :\ ∀\ i,\ Differentiable\ ℝ\ (U\ i))}
\leanline{4}{(hsc\ :\ ∀\ i,\ StrongConvexOn\ Set.univ\ m\ (U\ i))}
\leanline{4}{(hLip\ :\ ∀\ i,\ LipschitzWith\ L\ (gradient\ (U\ i)))}
\leanline{4}{(hs\ :\ StableStep\ m\ L\ h)\ (hG\ :\ 0\ ≤\ G)\ (hW\ :\ 0\ ≤\ Wbound)\ (\_hV\ :\ 0\ ≤\ V)}
\leanline{4}{(R\ :\ Lp\ (Cloud.State\ K\ dθ\ dz)\ 2\ P\ →\ Lp\ (Cloud.State\ K\ dθ\ dz)\ 2\ P)}
\leanline{4}{(hR\ :\ LipschitzWith\ L\ R)}
\leanline{4}{(y\ ξ\ ζ\ :\ ℕ\ →\ Lp\ (Cloud.State\ K\ dθ\ dz)\ 2\ P)}
\leanline{4}{(Y\ W\ :\ ℕ\ →\ ℝ\ →\ Lp\ (Cloud.State\ K\ dθ\ dz)\ 2\ P)}
\leanline{4}{(x₀\ :\ Lp\ (Cloud.State\ K\ dθ\ dz)\ 2\ P)\ (astar\ :\ EuclideanSpace\ ℝ\ (Fin\ dθ))}
\leanline{4}{(hr\ :\ IntegralReference\ R\ h\ G\ Wbound\ y\ ξ\ Y\ W)}
\leanline{4}{(fluctuation\ :\ ℕ\ →\ Fin\ K\ →\ Lp\ (Cloud.State\ K\ dθ\ dz)\ 2\ P)}
\leanline{4}{(hres\ :\ ∀\ n,\ particleForce\ P\ hK\ U\ m\ L\ hdiff\ hsc\ hLip\ (y\ n)\ -\ R\ (y\ n)\ =}
\leanline{6}{MiniBatchVariance.average\ (fluctuation\ n))}
\leanline{4}{(horth\ :\ ∀\ n\ i\ j,\ i\ ≠\ j\ →\ inner\ ℝ\ (fluctuation\ n\ i)\ (fluctuation\ n\ j)\ =\ 0)}
\leanline{4}{(hvariance\ :\ ∀\ n\ i,\ ‖fluctuation\ n\ i‖\ \textasciicircum{}\ 2\ ≤\ V)}
\leanline{4}{(horacle\ :\ HilbertInnovationControl\ (particleForce\ P\ hK\ U\ m\ L\ hdiff\ hsc\ hLip)}
\leanline{6}{h\ σ\ ξ\ ζ\ x₀\ x₀)}
\leanline{4}{(htarget\ :\ ∀\ n,\ parameterOutput\ P\ (y\ n)\ =\ constantL2\ P\ astar)\ (n\ :\ ℕ)\ :}
\leanline{4}{‖parameterOutput\ P}
\leanline{8}{(eulerTrajectory\ (particleForce\ P\ hK\ U\ m\ L\ hdiff\ hsc\ hLip)\ h}
\leanline{10}{(fun\ k\ =>\ ξ\ k\ +\ h\ •\ ζ\ k)\ x₀\ n)\ -\ constantL2\ P\ astar‖\ ≤}
\leanline{6}{(1\ -\ m\ *\ h\ /\ 2)\ \textasciicircum{}\ n\ *\ ‖x₀\ -\ y\ 0‖\ +}
\leanline{8}{(2\ /\ m)\ *\ Real.sqrt\ (V\ /\ K)\ +}
\leanline{8}{(L\ :\ ℝ)\ /\ m\ *\ (h\ *\ G\ +\ (4\ /\ 3\ :\ ℝ)\ *\ Wbound\ *\ Real.sqrt\ h)\ +}
\leanline{8}{Real.sqrt\ (h\ *\ σ\ \textasciicircum{}\ 2\ /\ m)}
\end{leancode}
\leansource{\texttt{particle\_\allowbreak{}oracle\_\allowbreak{}parameter\_\allowbreak{}bound} \quad LatentDiffusion/\allowbreak{}MainResults.lean, lines 47–85}
\begin{leancode}
\leanline{0}{theorem\ particle\_oracle\_parameter\_bound}
\leanline{4}{(hK\ :\ 0\ <\ K)\ (U\ :\ Fin\ K\ →\ Cloud.Joint\ dθ\ dz\ →\ ℝ)}
\leanline{4}{(m\ h\ G\ Wbound\ V\ σ\ :\ ℝ)\ (L\ :\ ℝ≥0)}
\leanline{4}{(hdiff\ :\ ∀\ i,\ Differentiable\ ℝ\ (U\ i))}
\leanline{4}{(hsc\ :\ ∀\ i,\ StrongConvexOn\ Set.univ\ m\ (U\ i))}
\leanline{4}{(hLip\ :\ ∀\ i,\ LipschitzWith\ L\ (gradient\ (U\ i)))}
\leanline{4}{(hs\ :\ StableStep\ m\ L\ h)\ (hG\ :\ 0\ ≤\ G)\ (hW\ :\ 0\ ≤\ Wbound)\ (\_hV\ :\ 0\ ≤\ V)}
\leanline{4}{(R\ :\ Lp\ (Cloud.State\ K\ dθ\ dz)\ 2\ P\ →\ Lp\ (Cloud.State\ K\ dθ\ dz)\ 2\ P)}
\leanline{4}{(hR\ :\ LipschitzWith\ L\ R)}
\leanline{4}{(y\ ξ\ ζ\ :\ ℕ\ →\ Lp\ (Cloud.State\ K\ dθ\ dz)\ 2\ P)}
\leanline{4}{(Y\ W\ :\ ℕ\ →\ ℝ\ →\ Lp\ (Cloud.State\ K\ dθ\ dz)\ 2\ P)}
\leanline{4}{(x₀\ :\ Lp\ (Cloud.State\ K\ dθ\ dz)\ 2\ P)\ (astar\ :\ EuclideanSpace\ ℝ\ (Fin\ dθ))}
\leanline{4}{(hr\ :\ IntegralReference\ R\ h\ G\ Wbound\ y\ ξ\ Y\ W)}
\leanline{4}{(fluctuation\ :\ ℕ\ →\ Fin\ K\ →\ Lp\ (Cloud.State\ K\ dθ\ dz)\ 2\ P)}
\leanline{4}{(hres\ :\ ∀\ n,\ particleForce\ P\ hK\ U\ m\ L\ hdiff\ hsc\ hLip\ (y\ n)\ -\ R\ (y\ n)\ =}
\leanline{6}{MiniBatchVariance.average\ (fluctuation\ n))}
\leanline{4}{(horth\ :\ ∀\ n\ i\ j,\ i\ ≠\ j\ →\ inner\ ℝ\ (fluctuation\ n\ i)\ (fluctuation\ n\ j)\ =\ 0)}
\leanline{4}{(hvariance\ :\ ∀\ n\ i,\ ‖fluctuation\ n\ i‖\ \textasciicircum{}\ 2\ ≤\ V)}
\leanline{4}{(horacle\ :\ HilbertInnovationControl\ (particleForce\ P\ hK\ U\ m\ L\ hdiff\ hsc\ hLip)}
\leanline{6}{h\ σ\ ξ\ ζ\ x₀\ x₀)}
\leanline{4}{(htarget\ :\ ∀\ n,\ parameterOutput\ P\ (y\ n)\ =\ constantL2\ P\ astar)\ (n\ :\ ℕ)\ :}
\leanline{4}{‖parameterOutput\ P}
\leanline{8}{(eulerTrajectory\ (particleForce\ P\ hK\ U\ m\ L\ hdiff\ hsc\ hLip)\ h}
\leanline{10}{(fun\ k\ =>\ ξ\ k\ +\ h\ •\ ζ\ k)\ x₀\ n)\ -\ constantL2\ P\ astar‖\ ≤}
\leanline{6}{(1\ -\ m\ *\ h\ /\ 2)\ \textasciicircum{}\ n\ *\ ‖x₀\ -\ y\ 0‖\ +}
\leanline{8}{(2\ /\ m)\ *\ Real.sqrt\ (V\ /\ K)\ +}
\leanline{8}{(L\ :\ ℝ)\ /\ m\ *\ (h\ *\ G\ +\ (4\ /\ 3\ :\ ℝ)\ *\ Wbound\ *\ Real.sqrt\ h)\ +}
\leanline{8}{Real.sqrt\ (h\ *\ σ\ \textasciicircum{}\ 2\ /\ m)}
\end{leancode}
\leansource{\texttt{particle\_\allowbreak{}oracle\_\allowbreak{}parameter\_\allowbreak{}bound} \quad LatentDiffusion/\allowbreak{}MainResults.lean, lines 47–85}
\begin{leancode}
\leanline{0}{theorem\ particle\_oracle\_parameter\_bound}
\leanline{4}{(hK\ :\ 0\ <\ K)\ (U\ :\ Fin\ K\ →\ Cloud.Joint\ dθ\ dz\ →\ ℝ)}
\leanline{4}{(m\ h\ G\ Wbound\ V\ σ\ :\ ℝ)\ (L\ :\ ℝ≥0)}
\leanline{4}{(hdiff\ :\ ∀\ i,\ Differentiable\ ℝ\ (U\ i))}
\leanline{4}{(hsc\ :\ ∀\ i,\ StrongConvexOn\ Set.univ\ m\ (U\ i))}
\leanline{4}{(hLip\ :\ ∀\ i,\ LipschitzWith\ L\ (gradient\ (U\ i)))}
\leanline{4}{(hs\ :\ StableStep\ m\ L\ h)\ (hG\ :\ 0\ ≤\ G)\ (hW\ :\ 0\ ≤\ Wbound)\ (\_hV\ :\ 0\ ≤\ V)}
\leanline{4}{(R\ :\ Lp\ (Cloud.State\ K\ dθ\ dz)\ 2\ P\ →\ Lp\ (Cloud.State\ K\ dθ\ dz)\ 2\ P)}
\leanline{4}{(hR\ :\ LipschitzWith\ L\ R)}
\leanline{4}{(y\ ξ\ ζ\ :\ ℕ\ →\ Lp\ (Cloud.State\ K\ dθ\ dz)\ 2\ P)}
\leanline{4}{(Y\ W\ :\ ℕ\ →\ ℝ\ →\ Lp\ (Cloud.State\ K\ dθ\ dz)\ 2\ P)}
\leanline{4}{(x₀\ :\ Lp\ (Cloud.State\ K\ dθ\ dz)\ 2\ P)\ (astar\ :\ EuclideanSpace\ ℝ\ (Fin\ dθ))}
\leanline{4}{(hr\ :\ IntegralReference\ R\ h\ G\ Wbound\ y\ ξ\ Y\ W)}
\leanline{4}{(fluctuation\ :\ ℕ\ →\ Fin\ K\ →\ Lp\ (Cloud.State\ K\ dθ\ dz)\ 2\ P)}
\leanline{4}{(hres\ :\ ∀\ n,\ particleForce\ P\ hK\ U\ m\ L\ hdiff\ hsc\ hLip\ (y\ n)\ -\ R\ (y\ n)\ =}
\leanline{6}{MiniBatchVariance.average\ (fluctuation\ n))}
\leanline{4}{(horth\ :\ ∀\ n\ i\ j,\ i\ ≠\ j\ →\ inner\ ℝ\ (fluctuation\ n\ i)\ (fluctuation\ n\ j)\ =\ 0)}
\leanline{4}{(hvariance\ :\ ∀\ n\ i,\ ‖fluctuation\ n\ i‖\ \textasciicircum{}\ 2\ ≤\ V)}
\leanline{4}{(horacle\ :\ HilbertInnovationControl\ (particleForce\ P\ hK\ U\ m\ L\ hdiff\ hsc\ hLip)}
\leanline{6}{h\ σ\ ξ\ ζ\ x₀\ x₀)}
\leanline{4}{(htarget\ :\ ∀\ n,\ parameterOutput\ P\ (y\ n)\ =\ constantL2\ P\ astar)\ (n\ :\ ℕ)\ :}
\leanline{4}{‖parameterOutput\ P}
\leanline{8}{(eulerTrajectory\ (particleForce\ P\ hK\ U\ m\ L\ hdiff\ hsc\ hLip)\ h}
\leanline{10}{(fun\ k\ =>\ ξ\ k\ +\ h\ •\ ζ\ k)\ x₀\ n)\ -\ constantL2\ P\ astar‖\ ≤}
\leanline{6}{(1\ -\ m\ *\ h\ /\ 2)\ \textasciicircum{}\ n\ *\ ‖x₀\ -\ y\ 0‖\ +}
\leanline{8}{(2\ /\ m)\ *\ Real.sqrt\ (V\ /\ K)\ +}
\leanline{8}{(L\ :\ ℝ)\ /\ m\ *\ (h\ *\ G\ +\ (4\ /\ 3\ :\ ℝ)\ *\ Wbound\ *\ Real.sqrt\ h)\ +}
\leanline{8}{Real.sqrt\ (h\ *\ σ\ \textasciicircum{}\ 2\ /\ m)}
\end{leancode}
\end{leanbox}
\section{New theorems}
\subsection{\texorpdfstring{New result — Two explicit minibatch temperature corrections}{New result — Two explicit minibatch temperature corrections}}



\textbf{Random scan:} select both an uninflated drift and its \(\sqrt{2h}\) noise with probability \(p\), leaving the particle fixed otherwise. Its scalar second-moment fixed point is



\[v_{\mathrm{scan}}=\frac1{a(1-ah/2)}.\]



\textbf{Full noise:} retain the inverse-inclusion drift and replace the printed noise by \(\sqrt{2h}\,W\) on every particle. Its second-moment fixed point is



\[v_{\mathrm{full}}=\frac1{a(1-ah/(2p))}.\]



The denominators must be nonzero in the algebraic identities. For meaningful positive stable scalar moments take \(a>0\) and respectively \(0<ah<2\) or \(0<ah<2p\). These exact quadratic calculations expose the remaining Euler bias and the \(1/p\) cost of random drift. They do not assert invariant-law existence for the adaptive nonlinear algorithm.

\begin{leanbox}{Lean: \texttt{randomScan\_\allowbreak{}stationary\_\allowbreak{}variance}, \texttt{fullNoise\_\allowbreak{}stationary\_\allowbreak{}variance}}
\leansource{\texttt{randomScan\_\allowbreak{}stationary\_\allowbreak{}variance} \quad LatentDiffusion/\allowbreak{}BatchAudit.lean, lines 234–245}
\begin{leancode}
\leanline{0}{theorem\ randomScan\_stationary\_variance\ (p\ h\ a\ :\ ℝ)\ (ha\ :\ a\ ≠\ 0)}
\leanline{4}{(hs\ :\ 1\ -\ a\ *\ h\ /\ 2\ ≠\ 0)\ :}
\leanline{4}{(1\ -\ 2\ *\ p\ *\ a\ *\ h\ +\ p\ *\ a\ \textasciicircum{}\ 2\ *\ h\ \textasciicircum{}\ 2)\ *}
\leanline{8}{(1\ /\ (a\ *\ (1\ -\ a\ *\ h\ /\ 2)))\ +\ 2\ *\ p\ *\ h\ =}
\leanline{6}{1\ /\ (a\ *\ (1\ -\ a\ *\ h\ /\ 2))}
\end{leancode}
\leansource{\texttt{fullNoise\_\allowbreak{}stationary\_\allowbreak{}variance} \quad LatentDiffusion/\allowbreak{}BatchAudit.lean, lines 265–281}
\begin{leancode}
\leanline{0}{theorem\ fullNoise\_stationary\_variance\ (p\ h\ a\ :\ ℝ)}
\leanline{4}{(hp\ :\ p\ ≠\ 0)\ (ha\ :\ a\ ≠\ 0)\ (hs\ :\ 1\ -\ a\ *\ h\ /\ (2\ *\ p)\ ≠\ 0)\ :}
\leanline{4}{paperCoefficient\ p\ h\ a\ *\ (1\ /\ (a\ *\ (1\ -\ a\ *\ h\ /\ (2\ *\ p))))\ +\ 2\ *\ h\ =}
\leanline{6}{1\ /\ (a\ *\ (1\ -\ a\ *\ h\ /\ (2\ *\ p)))}
\end{leancode}
\end{leanbox}
\subsection{\texorpdfstring{New result — Quantitative robustness to diffusion tilting}{New result — Quantitative robustness to diffusion tilting}}



Let the true evidence \(\ell\) be \(\lambda\)-strongly concave, \(\lambda>0\), with global maximizer \(a_\star\). Assume a uniform one-sided tilting gap and an approximate tilted optimizer:



\[0\le\ell(a)-\tilde\ell(a)\le\delta,\qquad\tilde\ell(a)\le\tilde\ell(\hat a)+\varepsilon\quad\text{for every }a,\qquad\delta,\varepsilon\ge0.\]



Then both the true-objective and true-parameter errors are controlled:



\[\ell(a_\star)-\ell(\hat a)\le\delta+\varepsilon,\qquad\|\hat a-a_\star\|\le\sqrt{\frac{2(\delta+\varepsilon)}\lambda}.\]



This quantifies the source's expressiveness discussion with an explicit model-bias assumption; it does not infer a uniform evidence gap from a pointwise training penalty.

\begin{leanbox}{Lean: \texttt{tilted\_\allowbreak{}objective\_\allowbreak{}bias}, \texttt{tilted\_\allowbreak{}parameter\_\allowbreak{}bias}}
\leansource{\texttt{tilted\_\allowbreak{}objective\_\allowbreak{}bias} \quad LatentDiffusion/\allowbreak{}Marginal.lean, lines 213–221}
\begin{leancode}
\leanline{0}{theorem\ tilted\_objective\_bias}
\leanline{4}{\{A\ :\ Type*\}\ (ell\ tilted\ :\ A\ →\ ℝ)\ (delta\ epsilon\ :\ ℝ)}
\leanline{4}{(hgap\ :\ ∀\ x,\ 0\ ≤\ ell\ x\ -\ tilted\ x\ ∧\ ell\ x\ -\ tilted\ x\ ≤\ delta)}
\leanline{4}{(star\ approx\ :\ A)\ (happrox\ :\ ∀\ x,\ tilted\ x\ ≤\ tilted\ approx\ +\ epsilon)\ :}
\leanline{4}{ell\ star\ -\ ell\ approx\ ≤\ delta\ +\ epsilon}
\end{leancode}
\leansource{\texttt{tilted\_\allowbreak{}parameter\_\allowbreak{}bias} \quad LatentDiffusion/\allowbreak{}Marginal.lean, lines 226–244}
\begin{leancode}
\leanline{0}{theorem\ tilted\_parameter\_bias}
\leanline{4}{\{H\ :\ Type*\}\ [NormedAddCommGroup\ H]\ [NormedSpace\ ℝ\ H]}
\leanline{4}{(ell\ tilted\ :\ H\ →\ ℝ)\ (rho\ delta\ epsilon\ :\ ℝ)\ (hrho\ :\ 0\ <\ rho)}
\leanline{4}{(hdelta\ :\ 0\ ≤\ delta)\ (hepsilon\ :\ 0\ ≤\ epsilon)}
\leanline{4}{(hsc\ :\ StrongConvexOn\ univ\ rho\ (fun\ x\ =>\ -ell\ x))}
\leanline{4}{(hgap\ :\ ∀\ x,\ 0\ ≤\ ell\ x\ -\ tilted\ x\ ∧\ ell\ x\ -\ tilted\ x\ ≤\ delta)}
\leanline{4}{(star\ approx\ :\ H)\ (hstar\ :\ ∀\ x,\ ell\ x\ ≤\ ell\ star)}
\leanline{4}{(happrox\ :\ ∀\ x,\ tilted\ x\ ≤\ tilted\ approx\ +\ epsilon)\ :}
\leanline{4}{‖approx\ -\ star‖\ ≤\ Real.sqrt\ (2\ *\ (delta\ +\ epsilon)\ /\ rho)}
\end{leancode}
\end{leanbox}
\subsection{\texorpdfstring{New result — Sharp Euler contraction}{New result — Sharp Euler contraction}}



Under N1 and \(h\ge0\), the same increment cancels in two Euler updates:



\[\|E_F(x,\xi)-E_F(y,\xi)\|^2\le(1-2\lambda h+L^2h^2)\|x-y\|^2.\]



Under N2 this gives, with \(q=1-\lambda h/2\in[0,1)\),



\[\|E_F(x,\xi)-E_F(y,\xi)\|\le q\|x-y\|.\]



The result holds for arbitrary common increments and the actual update map.

\begin{leanbox}{Lean: \texttt{euler\_\allowbreak{}sq\_\allowbreak{}contraction}, \texttt{euler\_\allowbreak{}contraction}}
\leansource{\texttt{euler\_\allowbreak{}sq\_\allowbreak{}contraction} \quad LatentDiffusion/\allowbreak{}Numerics.lean, lines 50–56}
\begin{leancode}
\leanline{0}{theorem\ euler\_sq\_contraction\ (F\ :\ H\ →\ H)\ (m\ h\ :\ ℝ)\ (L\ :\ ℝ≥0)}
\leanline{4}{(hF\ :\ ForceRegularity\ F\ m\ L)\ (hh\ :\ 0\ ≤\ h)\ (x\ y\ ξ\ :\ H)\ :}
\leanline{4}{‖eulerStep\ F\ h\ x\ ξ\ -\ eulerStep\ F\ h\ y\ ξ‖\ \textasciicircum{}\ 2\ ≤}
\leanline{6}{(1\ -\ 2\ *\ m\ *\ h\ +\ (L\ :\ ℝ)\ \textasciicircum{}\ 2\ *\ h\ \textasciicircum{}\ 2)\ *\ ‖x\ -\ y‖\ \textasciicircum{}\ 2}
\end{leancode}
\leansource{\texttt{euler\_\allowbreak{}contraction} \quad LatentDiffusion/\allowbreak{}Numerics.lean, lines 69–79}
\begin{leancode}
\leanline{0}{theorem\ euler\_contraction\ (F\ :\ H\ →\ H)\ (m\ h\ :\ ℝ)\ (L\ :\ ℝ≥0)}
\leanline{4}{(hF\ :\ ForceRegularity\ F\ m\ L)\ (hs\ :\ StableStep\ m\ L\ h)\ (x\ y\ ξ\ :\ H)\ :}
\leanline{4}{‖eulerStep\ F\ h\ x\ ξ\ -\ eulerStep\ F\ h\ y\ ξ‖\ ≤}
\leanline{6}{(1\ -\ m\ *\ h\ /\ 2)\ *\ ‖x\ -\ y‖}
\end{leancode}
\end{leanbox}
\subsection{\texorpdfstring{New result — Count-independent full-batch contraction}{New result — Count-independent full-batch contraction}}



Let \(K=MN>0\) count all datum–particle pairs. Give every coordinate its own datum potential \(U_i\), each satisfying A3–A4 with the same \(\lambda,L\). Use scaled state \(s=(a,z_1/\sqrt K,\ldots,z_K/\sqrt K)\), and put



\[F_K(s)=\frac1K\nabla_s\sum_{i=1}^K U_i(a,\sqrt K\,s_i),\qquad\|s-\bar s\|^2=\|a-\bar a\|^2+\frac1K\sum_i\|z_i-\bar z_i\|^2.\]



The full force \(F_K\) is \(\lambda\)-strongly monotone and \(L\)-Lipschitz with exactly the original pair constants. Thus, under N2, the actual Euler map contracts synchronously:



\[\|E_{F_K}(s,\xi)-E_{F_K}(\bar s,\xi)\|\le(1-\lambda h/2)\|s-\bar s\|.\]



The step restriction and contraction constant are independent of both data and particle counts. The coordinate-gradient identities are checked from the full lifted potential.

\begin{leanbox}{Lean: \texttt{meanForce\_\allowbreak{}regular}, \texttt{full\_\allowbreak{}batch\_\allowbreak{}contraction}}
\leansource{\texttt{meanForce\_\allowbreak{}regular} \quad LatentDiffusion/\allowbreak{}Cloud.lean, lines 233–253}
\begin{leancode}
\leanline{0}{theorem\ meanForce\_regular\ \{N\ dθ\ dx\ :\ ℕ\}\ (hN\ :\ 0\ <\ N)}
\leanline{4}{(U\ :\ Fin\ N\ →\ Joint\ dθ\ dx\ →\ ℝ)\ (m\ :\ ℝ)\ (L\ :\ ℝ≥0)}
\leanline{4}{(hdiff\ :\ ∀\ i,\ Differentiable\ ℝ\ (U\ i))}
\leanline{4}{(hsc\ :\ ∀\ i,\ StrongConvexOn\ Set.univ\ m\ (U\ i))}
\leanline{4}{(hLip\ :\ ∀\ i,\ LipschitzWith\ L\ (gradient\ (U\ i)))\ :}
\leanline{4}{LatentDiffusion.ForceRegularity\ (meanForce\ U)\ m\ L}
\end{leancode}
\leansource{\texttt{full\_\allowbreak{}batch\_\allowbreak{}contraction} \quad LatentDiffusion/\allowbreak{}Cloud.lean, lines 258–268}
\begin{leancode}
\leanline{0}{theorem\ full\_batch\_contraction\ \{N\ dθ\ dx\ :\ ℕ\}\ (hN\ :\ 0\ <\ N)}
\leanline{4}{(U\ :\ Fin\ N\ →\ Joint\ dθ\ dx\ →\ ℝ)\ (m\ h\ :\ ℝ)\ (L\ :\ ℝ≥0)}
\leanline{4}{(hdiff\ :\ ∀\ i,\ Differentiable\ ℝ\ (U\ i))}
\leanline{4}{(hsc\ :\ ∀\ i,\ StrongConvexOn\ Set.univ\ m\ (U\ i))}
\leanline{4}{(hLip\ :\ ∀\ i,\ LipschitzWith\ L\ (gradient\ (U\ i)))}
\leanline{4}{(hs\ :\ LatentDiffusion.StableStep\ m\ L\ h)\ (x\ y\ ξ\ :\ State\ N\ dθ\ dx)\ :}
\leanline{4}{‖LatentDiffusion.eulerStep\ (meanForce\ U)\ h\ x\ ξ\ -}
\leanline{6}{LatentDiffusion.eulerStep\ (meanForce\ U)\ h\ y\ ξ‖\ ≤}
\leanline{8}{(1\ -\ m\ *\ h\ /\ 2)\ *\ ‖x\ -\ y‖}
\end{leancode}
\end{leanbox}
\subsection{\texorpdfstring{New result — Uniform discretization and model-bias budget}{New result — Uniform discretization and model-bias budget}}



Under N1–N3, for every integer \(n\ge0\),



\[\begin{aligned}\|X_n-y_n\|&\le q^n\|X_0-y_0\|+(hB+Dh\sqrt h)\frac{1-q^n}{1-q}\\&\le q^n\|X_0-y_0\|+\frac2\lambda(B+D\sqrt h).\end{aligned}\]



The bound separates the initial error, implemented-force bias and temporal consistency. Its constants are explicit and the error floor is uniform in the number of steps. Particle-count or dimension dependence can enter through \(\lambda,L,B,D\).

\begin{leanbox}{Lean: \texttt{trajectory\_\allowbreak{}error\_\allowbreak{}exact}, \texttt{trajectory\_\allowbreak{}error\_\allowbreak{}bound}}
\leansource{\texttt{trajectory\_\allowbreak{}error\_\allowbreak{}exact} \quad LatentDiffusion/\allowbreak{}Numerics.lean, lines 122–139}
\begin{leancode}
\leanline{0}{theorem\ trajectory\_error\_exact\ (F\ G\ :\ H\ →\ H)\ (m\ h\ B\ D\ :\ ℝ)\ (L\ :\ ℝ≥0)}
\leanline{4}{(hF\ :\ ForceRegularity\ F\ m\ L)\ (hs\ :\ StableStep\ m\ L\ h)}
\leanline{4}{(ξ\ :\ ℕ\ →\ H)\ (x₀\ :\ H)\ (y\ :\ ℕ\ →\ H)}
\leanline{4}{(hforce\ :\ ∀\ n,\ ‖G\ (eulerTrajectory\ G\ h\ ξ\ x₀\ n)\ -\ F\ (eulerTrajectory\ G\ h\ ξ\ x₀\ n)‖\ ≤\ B)}
\leanline{4}{(hdefect\ :\ ∀\ n,\ ‖y\ (n\ +\ 1)\ -\ eulerStep\ F\ h\ (y\ n)\ (ξ\ n)‖\ ≤\ D\ *\ h\ *\ Real.sqrt\ h)}
\leanline{4}{(n\ :\ ℕ)\ :}
\leanline{4}{‖eulerTrajectory\ G\ h\ ξ\ x₀\ n\ -\ y\ n‖\ ≤}
\leanline{6}{(1\ -\ m\ *\ h\ /\ 2)\ \textasciicircum{}\ n\ *\ ‖x₀\ -\ y\ 0‖\ +}
\leanline{8}{(h\ *\ B\ +\ D\ *\ h\ *\ Real.sqrt\ h)\ *}
\leanline{10}{(1\ -\ (1\ -\ m\ *\ h\ /\ 2)\ \textasciicircum{}\ n)\ /\ (1\ -\ (1\ -\ m\ *\ h\ /\ 2))}
\end{leancode}
\leansource{\texttt{trajectory\_\allowbreak{}error\_\allowbreak{}bound} \quad LatentDiffusion/\allowbreak{}Numerics.lean, lines 143–163}
\begin{leancode}
\leanline{0}{theorem\ trajectory\_error\_bound\ (F\ G\ :\ H\ →\ H)\ (m\ h\ B\ D\ :\ ℝ)\ (L\ :\ ℝ≥0)}
\leanline{4}{(hF\ :\ ForceRegularity\ F\ m\ L)\ (hs\ :\ StableStep\ m\ L\ h)\ (hB\ :\ 0\ ≤\ B)\ (hD\ :\ 0\ ≤\ D)}
\leanline{4}{(ξ\ :\ ℕ\ →\ H)\ (x₀\ :\ H)\ (y\ :\ ℕ\ →\ H)}
\leanline{4}{(hforce\ :\ ∀\ n,\ ‖G\ (eulerTrajectory\ G\ h\ ξ\ x₀\ n)\ -\ F\ (eulerTrajectory\ G\ h\ ξ\ x₀\ n)‖\ ≤\ B)}
\leanline{4}{(hdefect\ :\ ∀\ n,\ ‖y\ (n\ +\ 1)\ -\ eulerStep\ F\ h\ (y\ n)\ (ξ\ n)‖\ ≤\ D\ *\ h\ *\ Real.sqrt\ h)}
\leanline{4}{(n\ :\ ℕ)\ :}
\leanline{4}{‖eulerTrajectory\ G\ h\ ξ\ x₀\ n\ -\ y\ n‖\ ≤}
\leanline{6}{(1\ -\ m\ *\ h\ /\ 2)\ \textasciicircum{}\ n\ *\ ‖x₀\ -\ y\ 0‖\ +}
\leanline{8}{(2\ /\ m)\ *\ (B\ +\ D\ *\ Real.sqrt\ h)}
\end{leancode}
\end{leanbox}
\subsection{\texorpdfstring{New result — Local consistency derived from integral dynamics}{New result — Local consistency derived from integral dynamics}}



Under N1 and N4 with \(h\ge0\), the local Euler defect is derived, rather than assumed:



\[\|y_{n+1}-E_F(y_n,\xi_n)\|\le L\left(\frac{h^2G}{2}+\frac23h\sqrt h\,N\right).\]



Under N2 and an implemented-force error \(\|A(X_n)-F(X_n)\|\le B\), \(B\ge0\), this yields the uniform bound



\[\|X_n-y_n\|\le q^n\|X_0-y_0\|+\frac{2B}{\lambda}+\frac L\lambda\left(hG+\frac43N\sqrt h\right).\]



The same theorem applies in any complete real Hilbert space. Its \(O(h+\sqrt h)\) term follows from the displayed integral equation and uniform drift and noise bounds.

\begin{leanbox}{Lean: \texttt{integral\_\allowbreak{}reference\_\allowbreak{}local\_\allowbreak{}defect}, \texttt{integral\_\allowbreak{}reference\_\allowbreak{}error\_\allowbreak{}bound}}
\leansource{\texttt{integral\_\allowbreak{}reference\_\allowbreak{}local\_\allowbreak{}defect} \quad LatentDiffusion/\allowbreak{}Numerics.lean, lines 286–303}
\begin{leancode}
\leanline{0}{theorem\ integral\_reference\_local\_defect\ (F\ :\ H\ →\ H)\ (m\ h\ G\ N\ :\ ℝ)\ (L\ :\ ℝ≥0)}
\leanline{4}{(hF\ :\ ForceRegularity\ F\ m\ L)\ (hh\ :\ 0\ ≤\ h)}
\leanline{4}{(y\ :\ ℕ\ →\ H)\ (ξ\ :\ ℕ\ →\ H)\ (Y\ W\ :\ ℕ\ →\ ℝ\ →\ H)}
\leanline{4}{(hr\ :\ IntegralReference\ F\ h\ G\ N\ y\ ξ\ Y\ W)\ (n\ :\ ℕ)\ :}
\leanline{4}{‖y\ (n\ +\ 1)\ -\ eulerStep\ F\ h\ (y\ n)\ (ξ\ n)‖\ ≤}
\leanline{6}{(L\ :\ ℝ)\ *\ (h\ \textasciicircum{}\ 2\ *\ G\ /\ 2\ +\ (2\ /\ 3\ :\ ℝ)\ *\ h\ *\ Real.sqrt\ h\ *\ N)}
\end{leancode}
\leansource{\texttt{integral\_\allowbreak{}reference\_\allowbreak{}error\_\allowbreak{}bound} \quad LatentDiffusion/\allowbreak{}Numerics.lean, lines 310–337}
\begin{leancode}
\leanline{0}{theorem\ integral\_reference\_error\_bound\ (F\ A\ :\ H\ →\ H)\ (m\ h\ B\ G\ N\ :\ ℝ)\ (L\ :\ ℝ≥0)}
\leanline{4}{(hF\ :\ ForceRegularity\ F\ m\ L)\ (hs\ :\ StableStep\ m\ L\ h)}
\leanline{4}{(hB\ :\ 0\ ≤\ B)\ (hG\ :\ 0\ ≤\ G)\ (hN\ :\ 0\ ≤\ N)}
\leanline{4}{(y\ :\ ℕ\ →\ H)\ (ξ\ :\ ℕ\ →\ H)\ (Y\ W\ :\ ℕ\ →\ ℝ\ →\ H)\ (x₀\ :\ H)}
\leanline{4}{(hr\ :\ IntegralReference\ F\ h\ G\ N\ y\ ξ\ Y\ W)}
\leanline{4}{(hforce\ :\ ∀\ n,\ ‖A\ (eulerTrajectory\ A\ h\ ξ\ x₀\ n)\ -\ F\ (eulerTrajectory\ A\ h\ ξ\ x₀\ n)‖\ ≤\ B)}
\leanline{4}{(n\ :\ ℕ)\ :}
\leanline{4}{‖eulerTrajectory\ A\ h\ ξ\ x₀\ n\ -\ y\ n‖\ ≤}
\leanline{6}{(1\ -\ m\ *\ h\ /\ 2)\ \textasciicircum{}\ n\ *\ ‖x₀\ -\ y\ 0‖\ +}
\leanline{8}{2\ *\ B\ /\ m\ +\ (L\ :\ ℝ)\ /\ m\ *\ (h\ *\ G\ +\ (4\ /\ 3\ :\ ℝ)\ *\ N\ *\ Real.sqrt\ h)}
\end{leancode}
\end{leanbox}
\subsection{\texorpdfstring{New result — Stochastic-oracle mean-square error floor}{New result — Stochastic-oracle mean-square error floor}}



Under N1, N2 and N5,



\[\|X_n-Y_n\|^2\le(1-\lambda h)^n\|X_0-Y_0\|^2+\frac{h\sigma^2}{\lambda}.\]



For the actual Hilbert space \(L^2(\Omega,P;E)\) of square-integrable random variables on an arbitrary probability space, this is the genuine expectation bound



\[\mathbb E\|X_n-Y_n\|^2\le(1-\lambda h)^n\mathbb E\|X_0-Y_0\|^2+\frac{h\sigma^2}{\lambda}.\]



The force monotonicity and Lipschitz bounds in this version are explicit \(L^2\) premise. The theorem accepts Gaussian innovations and does not replace nonintegrable expectations by totalized real values.

\begin{leanbox}{Lean: \texttt{hilbert\_\allowbreak{}innovation\_\allowbreak{}sq\_\allowbreak{}bound}, \texttt{l2\_\allowbreak{}innovation\_\allowbreak{}mean\_\allowbreak{}square\_\allowbreak{}bound}}
\leansource{\texttt{hilbert\_\allowbreak{}innovation\_\allowbreak{}sq\_\allowbreak{}bound} \quad LatentDiffusion/\allowbreak{}Numerics.lean, lines 357–383}
\begin{leancode}
\leanline{0}{theorem\ hilbert\_innovation\_sq\_bound\ (F\ :\ H\ →\ H)\ (m\ h\ σ\ :\ ℝ)\ (L\ :\ ℝ≥0)}
\leanline{4}{(hF\ :\ ForceRegularity\ F\ m\ L)\ (hs\ :\ StableStep\ m\ L\ h)}
\leanline{4}{(ξ\ ζ\ :\ ℕ\ →\ H)\ (x₀\ y₀\ :\ H)\ (hc\ :\ HilbertInnovationControl\ F\ h\ σ\ ξ\ ζ\ x₀\ y₀)}
\leanline{4}{(n\ :\ ℕ)\ :}
\leanline{4}{‖eulerTrajectory\ F\ h\ (fun\ k\ ↦\ ξ\ k\ +\ h\ •\ ζ\ k)\ x₀\ n\ -\ eulerTrajectory\ F\ h\ ξ\ y₀\ n‖\ \textasciicircum{}\ 2\ ≤}
\leanline{6}{(1\ -\ m\ *\ h)\ \textasciicircum{}\ n\ *\ ‖x₀\ -\ y₀‖\ \textasciicircum{}\ 2\ +\ h\ *\ σ\ \textasciicircum{}\ 2\ /\ m}
\end{leancode}
\leansource{\texttt{l2\_\allowbreak{}innovation\_\allowbreak{}mean\_\allowbreak{}square\_\allowbreak{}bound} \quad LatentDiffusion/\allowbreak{}Numerics.lean, lines 405–414}
\begin{leancode}
\leanline{0}{theorem\ l2\_innovation\_mean\_square\_bound}
\leanline{4}{(F\ :\ Lp\ E\ 2\ P\ →\ Lp\ E\ 2\ P)\ (m\ h\ σ\ :\ ℝ)\ (L\ :\ ℝ≥0)}
\leanline{4}{(hF\ :\ ForceRegularity\ F\ m\ L)\ (hs\ :\ StableStep\ m\ L\ h)}
\leanline{4}{(ξ\ ζ\ :\ ℕ\ →\ Lp\ E\ 2\ P)\ (x₀\ y₀\ :\ Lp\ E\ 2\ P)}
\leanline{4}{(hc\ :\ HilbertInnovationControl\ F\ h\ σ\ ξ\ ζ\ x₀\ y₀)\ (n\ :\ ℕ)\ :}
\leanline{4}{(∫\ ω,\ ‖(eulerTrajectory\ F\ h\ (fun\ k\ ↦\ ξ\ k\ +\ h\ •\ ζ\ k)\ x₀\ n\ -}
\leanline{6}{eulerTrajectory\ F\ h\ ξ\ y₀\ n)\ ω‖\ \textasciicircum{}\ 2\ ∂P)\ ≤}
\leanline{8}{(1\ -\ m\ *\ h)\ \textasciicircum{}\ n\ *\ (∫\ ω,\ ‖(x₀\ -\ y₀)\ ω‖\ \textasciicircum{}\ 2\ ∂P)\ +\ h\ *\ σ\ \textasciicircum{}\ 2\ /\ m}
\end{leancode}
\end{leanbox}
\subsection{\texorpdfstring{New result — Genuine probability-space lift of the pointwise force}{New result — Genuine probability-space lift of the pointwise force}}



On any probability space, a pointwise force \(F:E\to E\) satisfying N1 induces a genuine operator \(\mathscr F:L^2(\Omega;E)\to L^2(\Omega;E)\) with



\[(\mathscr F X)(\omega)=F(X(\omega))\quad\text{almost surely}.\]



The lifted force inherits the same monotonicity and Lipschitz constants. Its Euler iterates agree almost surely with the pointwise recursion. Therefore the stochastic-oracle expectation bound holds from the monotonicity and Lipschitz bounds of the original pointwise force, together with N2 and the explicit innovation condition N5:



\[\mathbb E\|X_n-Y_n\|^2\le(1-\lambda h)^n\mathbb E\|X_0-Y_0\|^2+\frac{h\sigma^2}{\lambda}.\]



No separate lifted-force monotonicity or Lipschitz assumption or condition \(F(0)=0\) is needed.

\begin{leanbox}{Lean: \texttt{liftedForce\_\allowbreak{}coe}, \texttt{liftedForce\_\allowbreak{}regular}, \texttt{lifted\_\allowbreak{}eulerTrajectory\_\allowbreak{}coe} and 1 more}
\leansource{\texttt{liftedForce\_\allowbreak{}coe} \quad LatentDiffusion/\allowbreak{}Lift.lean, lines 48–58}
\begin{leancode}
\leanline{0}{theorem\ liftedForce\_coe\ \{F\ :\ E\ →\ E\}\ \{m\ :\ ℝ\}\ \{L\ :\ ℝ≥0\}}
\leanline{4}{(hF\ :\ ForceRegularity\ F\ m\ L)\ (X\ :\ Lp\ E\ 2\ P)\ :}
\leanline{4}{liftedForce\ P\ hF\ X\ =ᵐ[P]\ fun\ ω\ ↦\ F\ (X\ ω)}
\end{leancode}
\leansource{\texttt{liftedForce\_\allowbreak{}regular} \quad LatentDiffusion/\allowbreak{}Lift.lean, lines 62–79}
\begin{leancode}
\leanline{0}{theorem\ liftedForce\_regular\ \{F\ :\ E\ →\ E\}\ \{m\ :\ ℝ\}\ \{L\ :\ ℝ≥0\}}
\leanline{4}{(hF\ :\ ForceRegularity\ F\ m\ L)\ :\ ForceRegularity\ (liftedForce\ P\ hF)\ m\ L}
\end{leancode}
\leansource{\texttt{lifted\_\allowbreak{}eulerTrajectory\_\allowbreak{}coe} \quad LatentDiffusion/\allowbreak{}Lift.lean, lines 95–107}
\begin{leancode}
\leanline{0}{theorem\ lifted\_eulerTrajectory\_coe\ \{F\ :\ E\ →\ E\}\ \{m\ :\ ℝ\}\ \{L\ :\ ℝ≥0\}}
\leanline{4}{(hF\ :\ ForceRegularity\ F\ m\ L)\ (h\ :\ ℝ)\ (ξ\ :\ ℕ\ →\ Lp\ E\ 2\ P)}
\leanline{4}{(x₀\ :\ Lp\ E\ 2\ P)\ (n\ :\ ℕ)\ :}
\leanline{4}{eulerTrajectory\ (liftedForce\ P\ hF)\ h\ ξ\ x₀\ n\ =ᵐ[P]}
\leanline{6}{fun\ ω\ ↦\ eulerTrajectory\ F\ h\ (fun\ k\ ↦\ ξ\ k\ ω)\ (x₀\ ω)\ n}
\end{leancode}
\leansource{\texttt{pointwise\_\allowbreak{}force\_\allowbreak{}mean\_\allowbreak{}square\_\allowbreak{}bound} \quad LatentDiffusion/\allowbreak{}Lift.lean, lines 111–120}
\begin{leancode}
\leanline{0}{theorem\ pointwise\_force\_mean\_square\_bound\ (F\ :\ E\ →\ E)\ (m\ h\ σ\ :\ ℝ)\ (L\ :\ ℝ≥0)}
\leanline{4}{(hF\ :\ ForceRegularity\ F\ m\ L)\ (hs\ :\ StableStep\ m\ L\ h)}
\leanline{4}{(ξ\ ζ\ :\ ℕ\ →\ Lp\ E\ 2\ P)\ (x₀\ y₀\ :\ Lp\ E\ 2\ P)}
\leanline{4}{(hc\ :\ HilbertInnovationControl\ (liftedForce\ P\ hF)\ h\ σ\ ξ\ ζ\ x₀\ y₀)\ (n\ :\ ℕ)\ :}
\leanline{4}{(∫\ ω,\ ‖(eulerTrajectory\ (liftedForce\ P\ hF)\ h}
\leanline{6}{(fun\ k\ ↦\ ξ\ k\ +\ h\ •\ ζ\ k)\ x₀\ n\ -}
\leanline{6}{eulerTrajectory\ (liftedForce\ P\ hF)\ h\ ξ\ y₀\ n)\ ω‖\ \textasciicircum{}\ 2\ ∂P)\ ≤}
\leanline{8}{(1\ -\ m\ *\ h)\ \textasciicircum{}\ n\ *\ (∫\ ω,\ ‖(x₀\ -\ y₀)\ ω‖\ \textasciicircum{}\ 2\ ∂P)\ +\ h\ *\ σ\ \textasciicircum{}\ 2\ /\ m}
\end{leancode}
\end{leanbox}
\subsection{\texorpdfstring{New result — Heterogeneous particle-average variance}{New result — Heterogeneous particle-average variance}}



Let \(K>0\), and let \(g_1,\ldots,g_K\) be vectors in a real Hilbert space, such as centered random force errors in \(L^2\). If \(\langle g_i,g_j\rangle=0\) for \(i\ne j\), then



\[\left\|\frac1K\sum_i g_i\right\|^2=\frac1{K^2}\sum_i\|g_i\|^2\le\frac VK\quad\text{when }\|g_i\|^2\le V.\]



Hence a residual represented by this average has \(B\le\sqrt{V/K}\). The copies may have different laws, as occurs for different data. Orthogonality is the explicit cross-moment premise; independence and centering must be verified when used to establish it.

\begin{leanbox}{Lean: \texttt{mean\_\allowbreak{}sq\_\allowbreak{}sum\_\allowbreak{}of\_\allowbreak{}orthogonal}, \texttt{mean\_\allowbreak{}sum\_\allowbreak{}of\_\allowbreak{}orthogonal\_\allowbreak{}le}}
\leansource{\texttt{mean\_\allowbreak{}sq\_\allowbreak{}sum\_\allowbreak{}of\_\allowbreak{}orthogonal} \quad LatentDiffusion/\allowbreak{}MiniBatchVariance.lean, lines 41–45}
\begin{leancode}
\leanline{0}{theorem\ mean\_sq\_sum\_of\_orthogonal\ \{K\ :\ ℕ\}\ (g\ :\ Fin\ K\ →\ H)}
\leanline{4}{(hc\ :\ ∀\ i\ j,\ i\ ≠\ j\ →\ inner\ ℝ\ (g\ i)\ (g\ j)\ =\ 0)\ :}
\leanline{4}{‖average\ g‖\ \textasciicircum{}\ 2\ =\ (1\ /\ (K\ :\ ℝ))\ \textasciicircum{}\ 2\ *\ ∑\ i,\ ‖g\ i‖\ \textasciicircum{}\ 2}
\end{leancode}
\leansource{\texttt{mean\_\allowbreak{}sum\_\allowbreak{}of\_\allowbreak{}orthogonal\_\allowbreak{}le} \quad LatentDiffusion/\allowbreak{}MiniBatchVariance.lean, lines 64–69}
\begin{leancode}
\leanline{0}{theorem\ mean\_sum\_of\_orthogonal\_le\ \{K\ :\ ℕ\}\ (hK\ :\ 0\ <\ K)}
\leanline{4}{(g\ :\ Fin\ K\ →\ H)\ (V\ :\ ℝ)}
\leanline{4}{(hc\ :\ ∀\ i\ j,\ i\ ≠\ j\ →\ inner\ ℝ\ (g\ i)\ (g\ j)\ =\ 0)}
\leanline{4}{(hV\ :\ ∀\ i,\ ‖g\ i‖\ \textasciicircum{}\ 2\ ≤\ V)\ :\ ‖average\ g‖\ ≤\ Real.sqrt\ (V\ /\ K)}
\end{leancode}
\end{leanbox}
\subsection{\texorpdfstring{New result — Exact minibatch finite-population correction}{New result — Exact minibatch finite-population correction}}



Let \(M>1\), \(b>0\). An actual finite probability space carries batches of size \(b\), with inclusion probabilities



\[\mathbb P(i\in\mathcal B)=\frac bM,\qquad\mathbb P(i,j\in\mathcal B)=\frac{b(b-1)}{M(M-1)}\quad(i\ne j).\]



For arbitrary heterogeneous per-datum gradients \(g_i\), put \(\bar g=M^{-1}\sum_i g_i\). The actual raw-gradient oracle obeys



\[\mathbb E\left\|\frac1b\sum_{i\in\mathcal B}g_i-\bar g\right\|^2=\frac{M-b}{b(M-1)}\,\frac1M\sum_i\|g_i-\bar g\|^2.\]



For \(b\le M\), a population dispersion bound \(M^{-1}\sum_i\|g_i-\bar g\|^2\le V\) therefore gives the derived oracle variance \(\sigma^2\le (M-b)V/[b(M-1)]\). The theorem derives this from the primitive first and pair inclusion laws, including zero sampling variance at full batch. Applying this identity to the whole-cloud stochastic force requires its actual scaled blocks and inverse-inclusion factors: the dispersion \(V\) may depend on both \(M\) and \(N\), so a count-independent \(\sigma^2\) bound does not follow automatically.

\begin{leanbox}{Lean: \texttt{BatchInclusionLaw}, \texttt{minibatch\_\allowbreak{}gradient\_\allowbreak{}variance}, \texttt{minibatch\_\allowbreak{}oracle\_\allowbreak{}variance\_\allowbreak{}bound}}
\leansource{\texttt{BatchInclusionLaw} \quad LatentDiffusion/\allowbreak{}MiniBatchVariance.lean, lines 124–129}
\begin{leancode}
\leanline{0}{structure\ BatchInclusionLaw\ \{M\ :\ ℕ\}\ (B\ :\ ℕ)\ (batch\ :\ Ω\ →\ Finset\ (Fin\ M))\ :\ Prop\ where}
\leanline{2}{size\ :\ ∀\ ω,\ (batch\ ω).card\ =\ B}
\leanline{2}{first\ :\ ∀\ i,\ (∫\ ω,\ inclusion\ batch\ ω\ i\ ∂P)\ =\ (B\ :\ ℝ)\ /\ M}
\leanline{2}{pair\ :\ ∀\ i\ j,\ i\ ≠\ j\ →}
\leanline{4}{(∫\ ω,\ inclusion\ batch\ ω\ i\ *\ inclusion\ batch\ ω\ j\ ∂P)\ =}
\leanline{6}{(B\ :\ ℝ)\ *\ ((B\ :\ ℝ)\ -\ 1)\ /\ ((M\ :\ ℝ)\ *\ ((M\ :\ ℝ)\ -\ 1))}
\end{leancode}
\leansource{\texttt{minibatch\_\allowbreak{}gradient\_\allowbreak{}variance} \quad LatentDiffusion/\allowbreak{}MiniBatchVariance.lean, lines 209–217}
\begin{leancode}
\leanline{0}{theorem\ minibatch\_gradient\_variance\ \{M\ B\ :\ ℕ\}\ (hM\ :\ 1\ <\ M)\ (hB\ :\ 0\ <\ B)}
\leanline{4}{(batch\ :\ Ω\ →\ Finset\ (Fin\ M))\ (hbatch\ :\ BatchInclusionLaw\ P\ B\ batch)}
\leanline{4}{(g\ :\ Fin\ M\ →\ H)\ :}
\leanline{4}{(∫\ ω,\ ‖batchAverage\ B\ batch\ g\ ω\ -\ average\ g‖\ \textasciicircum{}\ 2\ ∂P)\ =}
\leanline{6}{((M\ :\ ℝ)\ -\ B)\ /\ ((B\ :\ ℝ)\ *\ ((M\ :\ ℝ)\ -\ 1))\ *}
\leanline{8}{((1\ /\ (M\ :\ ℝ))\ *\ ∑\ i,\ ‖g\ i\ -\ average\ g‖\ \textasciicircum{}\ 2)}
\end{leancode}
\leansource{\texttt{minibatch\_\allowbreak{}oracle\_\allowbreak{}variance\_\allowbreak{}bound} \quad LatentDiffusion/\allowbreak{}MiniBatchVariance.lean, lines 225–239}
\begin{leancode}
\leanline{0}{theorem\ minibatch\_oracle\_variance\_bound\ \{M\ B\ :\ ℕ\}}
\leanline{4}{(hM\ :\ 1\ <\ M)\ (hB\ :\ 0\ <\ B)\ (hBM\ :\ B\ ≤\ M)}
\leanline{4}{(batch\ :\ Ω\ →\ Finset\ (Fin\ M))\ (hbatch\ :\ BatchInclusionLaw\ P\ B\ batch)}
\leanline{4}{(g\ :\ Fin\ M\ →\ H)\ (V\ :\ ℝ)}
\leanline{4}{(hV\ :\ ((1\ /\ (M\ :\ ℝ))\ *\ ∑\ i,\ ‖g\ i\ -\ average\ g‖\ \textasciicircum{}\ 2)\ ≤\ V)\ :}
\leanline{4}{(∫\ ω,\ ‖batchAverage\ B\ batch\ g\ ω\ -\ average\ g‖\ \textasciicircum{}\ 2\ ∂P)\ ≤}
\leanline{6}{oracleVariance\ M\ B\ V}
\end{leancode}
\end{leanbox}
\subsection{\texorpdfstring{New main bound — Uniform inference error with a stochastic oracle}{New main bound — Uniform inference error with a stochastic oracle}}



Let \(K=MN>0\), and impose A3–A4 on each actual datum potential and N2, N5–N7 on the reference and innovations. Construct the complete interacting force \(F_K\) from those potentials, and use its actual noisy Euler update \(X_{n+1}=X_n-hF_K(X_n)+\xi_n+h\zeta_n\). Write \(a_n=CX_n\) for the physical parameter and \(q=1-\lambda h/2\). For every \(n\ge0\),



\[\begin{aligned}\left(\mathbb E\|a_n-a_\star\|^2\right)^{1/2}\le{}&q^n\|X_0-y_0\|_{L^2}+\frac2\lambda\sqrt{\frac VK}\\&+\frac L\lambda\left(hG+\frac43W\sqrt h\right)+\sqrt{\frac{h\sigma^2}{\lambda}}.\end{aligned}\]



The force monotonicity and Lipschitz bounds are derived from the original per-datum potentials, including the exact pointwise-to-\(L^2\) lift. The output norm is proved equal to the physical parameter root mean-square error. The four explicit terms quantify initialization, finite-particle averaging, temporal discretization and the centered stochastic oracle. The finite-population identity below calibrates oracle variance from batch dispersion.




The bound is uniform in the number of steps. Its target \(a_\star\) is the fixed reference parameter specified in N6; interpreting it as convergence to the paper's optimizer requires that identification and a corresponding frozen posterior reference. It retains explicit reference existence, moment and innovation premises; it does not claim convergence of the unrestricted Adam/\allowbreak{}preconditioned implementation.

\begin{leanbox}{Lean: \texttt{particle\_\allowbreak{}oracle\_\allowbreak{}parameter\_\allowbreak{}bound}, \texttt{parameterOutput\_\allowbreak{}rms}}
\leansource{\texttt{particle\_\allowbreak{}oracle\_\allowbreak{}parameter\_\allowbreak{}bound} \quad LatentDiffusion/\allowbreak{}MainResults.lean, lines 47–85}
\begin{leancode}
\leanline{0}{theorem\ particle\_oracle\_parameter\_bound}
\leanline{4}{(hK\ :\ 0\ <\ K)\ (U\ :\ Fin\ K\ →\ Cloud.Joint\ dθ\ dz\ →\ ℝ)}
\leanline{4}{(m\ h\ G\ Wbound\ V\ σ\ :\ ℝ)\ (L\ :\ ℝ≥0)}
\leanline{4}{(hdiff\ :\ ∀\ i,\ Differentiable\ ℝ\ (U\ i))}
\leanline{4}{(hsc\ :\ ∀\ i,\ StrongConvexOn\ Set.univ\ m\ (U\ i))}
\leanline{4}{(hLip\ :\ ∀\ i,\ LipschitzWith\ L\ (gradient\ (U\ i)))}
\leanline{4}{(hs\ :\ StableStep\ m\ L\ h)\ (hG\ :\ 0\ ≤\ G)\ (hW\ :\ 0\ ≤\ Wbound)\ (\_hV\ :\ 0\ ≤\ V)}
\leanline{4}{(R\ :\ Lp\ (Cloud.State\ K\ dθ\ dz)\ 2\ P\ →\ Lp\ (Cloud.State\ K\ dθ\ dz)\ 2\ P)}
\leanline{4}{(hR\ :\ LipschitzWith\ L\ R)}
\leanline{4}{(y\ ξ\ ζ\ :\ ℕ\ →\ Lp\ (Cloud.State\ K\ dθ\ dz)\ 2\ P)}
\leanline{4}{(Y\ W\ :\ ℕ\ →\ ℝ\ →\ Lp\ (Cloud.State\ K\ dθ\ dz)\ 2\ P)}
\leanline{4}{(x₀\ :\ Lp\ (Cloud.State\ K\ dθ\ dz)\ 2\ P)\ (astar\ :\ EuclideanSpace\ ℝ\ (Fin\ dθ))}
\leanline{4}{(hr\ :\ IntegralReference\ R\ h\ G\ Wbound\ y\ ξ\ Y\ W)}
\leanline{4}{(fluctuation\ :\ ℕ\ →\ Fin\ K\ →\ Lp\ (Cloud.State\ K\ dθ\ dz)\ 2\ P)}
\leanline{4}{(hres\ :\ ∀\ n,\ particleForce\ P\ hK\ U\ m\ L\ hdiff\ hsc\ hLip\ (y\ n)\ -\ R\ (y\ n)\ =}
\leanline{6}{MiniBatchVariance.average\ (fluctuation\ n))}
\leanline{4}{(horth\ :\ ∀\ n\ i\ j,\ i\ ≠\ j\ →\ inner\ ℝ\ (fluctuation\ n\ i)\ (fluctuation\ n\ j)\ =\ 0)}
\leanline{4}{(hvariance\ :\ ∀\ n\ i,\ ‖fluctuation\ n\ i‖\ \textasciicircum{}\ 2\ ≤\ V)}
\leanline{4}{(horacle\ :\ HilbertInnovationControl\ (particleForce\ P\ hK\ U\ m\ L\ hdiff\ hsc\ hLip)}
\leanline{6}{h\ σ\ ξ\ ζ\ x₀\ x₀)}
\leanline{4}{(htarget\ :\ ∀\ n,\ parameterOutput\ P\ (y\ n)\ =\ constantL2\ P\ astar)\ (n\ :\ ℕ)\ :}
\leanline{4}{‖parameterOutput\ P}
\leanline{8}{(eulerTrajectory\ (particleForce\ P\ hK\ U\ m\ L\ hdiff\ hsc\ hLip)\ h}
\leanline{10}{(fun\ k\ =>\ ξ\ k\ +\ h\ •\ ζ\ k)\ x₀\ n)\ -\ constantL2\ P\ astar‖\ ≤}
\leanline{6}{(1\ -\ m\ *\ h\ /\ 2)\ \textasciicircum{}\ n\ *\ ‖x₀\ -\ y\ 0‖\ +}
\leanline{8}{(2\ /\ m)\ *\ Real.sqrt\ (V\ /\ K)\ +}
\leanline{8}{(L\ :\ ℝ)\ /\ m\ *\ (h\ *\ G\ +\ (4\ /\ 3\ :\ ℝ)\ *\ Wbound\ *\ Real.sqrt\ h)\ +}
\leanline{8}{Real.sqrt\ (h\ *\ σ\ \textasciicircum{}\ 2\ /\ m)}
\end{leancode}
\leansource{\texttt{parameterOutput\_\allowbreak{}rms} \quad LatentDiffusion/\allowbreak{}MainResults.lean, lines 89–104}
\begin{leancode}
\leanline{0}{theorem\ parameterOutput\_rms\ (X\ :\ Lp\ (Cloud.State\ K\ dθ\ dz)\ 2\ P)}
\leanline{4}{(astar\ :\ EuclideanSpace\ ℝ\ (Fin\ dθ))\ :}
\leanline{4}{‖parameterOutput\ P\ X\ -\ constantL2\ P\ astar‖\ =}
\leanline{6}{Real.sqrt\ (∫\ ω,\ ‖IPLA.parameterCLM\ (X\ ω)\ -\ astar‖\ \textasciicircum{}\ 2\ ∂P)}
\end{leancode}
\end{leanbox}
\subsection{\texorpdfstring{New result — Explicit choices for any target accuracy}{New result — Explicit choices for any target accuracy}}



Let \(\varepsilon>0\) and \(R_0=\|X_0-y_0\|_{L^2}\). Under the preceding main bound, the following explicit conditions suffice for parameter RMS error at most \(\varepsilon\):



\[\begin{gathered}\frac{\lambda h n}{2}\ge\log\max\!\left\{1,\frac{5R_0}{\varepsilon}\right\},\qquad100V\le\lambda^2\varepsilon^2K,\\5LhG\le\lambda\varepsilon,\qquad400L^2W^2h\le9\lambda^2\varepsilon^2,\qquad25h\sigma^2\le\lambda\varepsilon^2.\end{gathered}\]



Each of the five initialization, particle, drift, noise-modulus and oracle terms receives \(\varepsilon/5\). The conditions are division-free, including zero drift, noise or variance constants. The burn-in formula also covers zero initial error. The probabilistic/\allowbreak{}reference premises remain those of the main theorem.

\begin{leanbox}{Lean: \texttt{burnin}, \texttt{budget\_\allowbreak{}le\_\allowbreak{}epsilon}}
\leansource{\texttt{burnin} \quad LatentDiffusion/\allowbreak{}Accuracy.lean, lines 25–43}
\begin{leancode}
\leanline{0}{theorem\ burnin\ (m\ h\ R0\ epsilon\ :\ ℝ)\ (hR\ :\ 0\ ≤\ R0)\ (he\ :\ 0\ <\ epsilon)}
\leanline{4}{(hq\ :\ 0\ ≤\ 1\ -\ m\ *\ h\ /\ 2)\ (n\ :\ ℕ)}
\leanline{4}{(hn\ :\ Real.log\ (max\ 1\ (5\ *\ R0\ /\ epsilon))\ ≤\ m\ *\ h\ *\ n\ /\ 2)\ :}
\leanline{4}{(1\ -\ m\ *\ h\ /\ 2)\ \textasciicircum{}\ n\ *\ R0\ ≤\ epsilon\ /\ 5}
\end{leancode}
\leansource{\texttt{budget\_\allowbreak{}le\_\allowbreak{}epsilon} \quad LatentDiffusion/\allowbreak{}Accuracy.lean, lines 53–82}
\begin{leancode}
\leanline{0}{theorem\ budget\_le\_epsilon}
\leanline{4}{(m\ h\ R0\ V\ L\ G\ W\ sigma\ epsilon\ K\ E\ :\ ℝ)\ (n\ :\ ℕ)}
\leanline{4}{(hm\ :\ 0\ <\ m)\ (hh\ :\ 0\ ≤\ h)\ (hR\ :\ 0\ ≤\ R0)\ (hV\ :\ 0\ ≤\ V)}
\leanline{4}{(hL\ :\ 0\ ≤\ L)\ (\_hG\ :\ 0\ ≤\ G)\ (hW\ :\ 0\ ≤\ W)\ (hK\ :\ 0\ <\ K)\ (he\ :\ 0\ <\ epsilon)}
\leanline{4}{(hq\ :\ 0\ ≤\ 1\ -\ m\ *\ h\ /\ 2)}
\leanline{4}{(hn\ :\ Real.log\ (max\ 1\ (5\ *\ R0\ /\ epsilon))\ ≤\ m\ *\ h\ *\ n\ /\ 2)}
\leanline{4}{(hparticles\ :\ 100\ *\ V\ ≤\ m\ \textasciicircum{}\ 2\ *\ epsilon\ \textasciicircum{}\ 2\ *\ K)}
\leanline{4}{(hdrift\ :\ 5\ *\ L\ *\ h\ *\ G\ ≤\ m\ *\ epsilon)}
\leanline{4}{(hnoise\ :\ 400\ *\ L\ \textasciicircum{}\ 2\ *\ W\ \textasciicircum{}\ 2\ *\ h\ ≤\ 9\ *\ m\ \textasciicircum{}\ 2\ *\ epsilon\ \textasciicircum{}\ 2)}
\leanline{4}{(horacle\ :\ 25\ *\ h\ *\ sigma\ \textasciicircum{}\ 2\ ≤\ m\ *\ epsilon\ \textasciicircum{}\ 2)}
\leanline{4}{(hbudget\ :\ E\ ≤\ (1\ -\ m\ *\ h\ /\ 2)\ \textasciicircum{}\ n\ *\ R0\ +\ (2\ /\ m)\ *\ Real.sqrt\ (V\ /\ K)\ +}
\leanline{6}{L\ /\ m\ *\ (h\ *\ G\ +\ (4\ /\ 3\ :\ ℝ)\ *\ W\ *\ Real.sqrt\ h)\ +}
\leanline{6}{Real.sqrt\ (h\ *\ sigma\ \textasciicircum{}\ 2\ /\ m))\ :\ E\ ≤\ epsilon}
\end{leancode}
\end{leanbox}
\chapter{Mathematical argument}
\section{\texorpdfstring{Variational characterization of the evidence}{Variational characterization of the evidence}}


This section proves the variational characterization of the evidence. Let \(M\ge1\) be the number of data, let \(\mathcal A\) be an arbitrary set of parameters and let \(Z\) be a measurable space of latent variables. A \emph{posterior family} consists of a real function \(\ell:\mathcal A\to\mathbb R\), the log evidence, and, for each parameter \(a\in\mathcal A\) and each datum \(m\in\{1,\dots,M\}\), a probability measure \(\pi_a^m\) on \(Z\), the posterior of datum \(m\). For a collection \(q=(q^1,\dots,q^M)\) of probability measures on \(Z\), the free energy is


\[F(a,q)=-\ell(a)+\frac1M\sum_{m=1}^M\mathrm{KL}(q^m\|\pi_a^m),\]


where \(\mathrm{KL}\) denotes relative entropy. We prove that \(F(a,q)\ge-\ell(a)\), with equality exactly when \(q^m=\pi_a^m\) for every \(m\), and deduce the corrected form of the paper's Proposition 1: a pair \((a_\star,q_\star)\) minimizes \(F\) among all pairs with finite relative entropies if and only if \(a_\star\) maximizes \(\ell\) and \(q_\star^m=\pi_{a_\star}^m\) for every \(m\). No existence of a maximizer is asserted; existence and uniqueness are proved later from strong convexity of the marginal potential.




In the application \(\mathcal A\) is the space of parameters \(a=(\theta,\phi)\), \(\pi_a^m\) is the normalized tilted posterior \(Z_m(a)^{-1}e^{-U_m(a,z)}\,dz\) of datum \(m\), and \(\ell(a)=M^{-1}\sum_m\log Z_m(a)\) is the average tilted log evidence. The results hold for every parameter set and every measurable latent space, in particular for tilted models. The last step shows that a zero tilting penalty at a maximizer of the tilted evidence does not identify a maximizer of the original evidence, and gives a sufficient condition that does.



\textbf{Claim.} Let \(M\ge1\), \(a\in\mathcal A\), and let \(q\) be a collection of \(M\) probability measures on \(Z\). Then \((a,q)\) is a global minimizer of \(F\) if and only if


\[\ell(b)\le\ell(a)\ \text{ for every }b\in\mathcal A\qquad\text{and}\qquad q^m=\pi_a^m\ \text{ for every }m.\]


\textbf{Proof.} (\(\Rightarrow\)) Let \((a,q)\) be a global minimizer; in particular \(q\) is admissible at \(a\). The posterior collection \(\pi_a\) is admissible at \(a\), so minimality and the posterior identity of the second step give


\[F(a,q)\le F(a,\pi_a)=-\ell(a).\]


The lower bound of the second step gives \(F(a,q)\ge-\ell(a)\); hence \(F(a,q)=-\ell(a)\), and the previous step yields \(q=\pi_a\). Now let \(b\in\mathcal A\). Since \(\pi_b\) is admissible at \(b\), minimality gives


\[-\ell(a)=F(a,q)\le F(b,\pi_b)=-\ell(b),\]


that is, \(\ell(b)\le\ell(a)\).




(\(\Leftarrow\)) Suppose \(a\) maximizes \(\ell\) and \(q=\pi_a\). Then \(q\) is admissible at \(a\). For every \(b\) and every \(r\) admissible at \(b\), the posterior identity, the maximality of \(a\) and the lower bound at \(b\) give


\[F(a,\pi_a)=-\ell(a)\le-\ell(b)\le F(b,r).\]


Thus \((a,\pi_a)\) is a global minimizer. \(\square\)




The statement asserts neither that \(\ell\) has a maximizer nor that a maximizer is unique; both questions are settled later for strongly convex potentials. The normalized posterior collection is the unique optimal variational collection at each maximizer.

\begin{leanbox}{Lean: \texttt{global\_\allowbreak{}minimizers\_\allowbreak{}iff}}
\leansource{\texttt{global\_\allowbreak{}minimizers\_\allowbreak{}iff} \quad LatentDiffusion/\allowbreak{}Variational.lean, lines 90–106}
\begin{leancode}
\leanline{0}{theorem\ global\_minimizers\_iff\ (hM\ :\ 0\ <\ M)\ (a\ :\ A)}
\leanline{4}{(q\ :\ Fin\ M\ →\ ProbabilityMeasure\ Z)\ :}
\leanline{4}{P.IsGlobalMin\ a\ q\ ↔}
\leanline{6}{(∀\ b,\ P.logEvidence\ b\ ≤\ P.logEvidence\ a)\ ∧\ q\ =\ P.posterior\ a}
\end{leancode}
\leanlabel{How the Lean proof follows the argument}
\texttt{global\_\allowbreak{}minimizers\_\allowbreak{}iff} states the equivalence for \texttt{P.\allowbreak{}Is\allowbreak{}Global\allowbreak{}Min a q}; the right-hand side is the conjunction of \texttt{\ensuremath{\forall} b, P.\allowbreak{}log\allowbreak{}Evidence b \ensuremath{\le} P.\allowbreak{}log\allowbreak{}Evidence a} and \texttt{q = P.\allowbreak{}posterior a}.



\begin{enumerate}


\item Forward direction: \texttt{rintro \ensuremath{\langle}hq, hmin\ensuremath{\rangle}} splits the hypothesis into admissibility and minimality. \texttt{he} is the identity \(F(a,q)=-\ell(a)\), obtained by \texttt{le\_\allowbreak{}antisymm} from \texttt{hmin a (P.\allowbreak{}posterior a) (P.\allowbreak{}posterior\_\allowbreak{}admissible a)}, rewritten with \texttt{P.\allowbreak{}posterior\_\allowbreak{}free\allowbreak{}Energy}, and from \texttt{P.\allowbreak{}lower\_\allowbreak{}bound a q}. The second component is \texttt{(P.\allowbreak{}equality\_\allowbreak{}iff h\allowbreak{}M a q hq).\allowbreak{}mp he}. For the first component, \texttt{hb} applies minimality against \((b,\pi_b)\), and after rewriting with \texttt{he} and \texttt{P.\allowbreak{}posterior\_\allowbreak{}free\allowbreak{}Energy}, \texttt{linarith} turns \(-\ell(a)\le-\ell(b)\) into \(\ell(b)\le\ell(a)\).

\item Reverse direction: \texttt{rintro \ensuremath{\langle}hmax, rfl\ensuremath{\rangle}} substitutes \(q=\pi_a\). Admissibility is \texttt{P.\allowbreak{}posterior\_\allowbreak{}admissible a}. For an admissible pair \((b,r)\) the chain of the proof is \texttt{(neg\_\allowbreak{}le\_\allowbreak{}neg (hmax b)).\allowbreak{}trans (P.\allowbreak{}lower\_\allowbreak{}bound b r)}, after rewriting \(F(a,\pi_a)\) by \texttt{P.\allowbreak{}posterior\_\allowbreak{}free\allowbreak{}Energy}.

\end{enumerate}
\end{leanbox}
\section{\texorpdfstring{Unique maximizer of the marginal evidence}{Unique maximizer of the marginal evidence}}


For each datum \(m\in\{1,\dots,M\}\) let \(U_m:\mathbb R^{d_\theta}\times\mathbb R^{d_z}\to\mathbb R\) be the joint potential, the negative logarithm of the tilted joint density of the parameter \(a\) and the latent variable \(z\). The pair space carries the Euclidean norm \(\|(a,z)\|^2=\|a\|^2+\|z\|^2\). Assumption A3 states that each \(U_m\) is differentiable and \(\lambda\)-strongly convex with one constant \(\lambda>0\):


\[U_m(sx+(1-s)y)\le sU_m(x)+(1-s)U_m(y)-\frac\lambda2s(1-s)\|x-y\|^2\qquad(x,y\in\mathbb R^{d_\theta+d_z},\ s\in[0,1]).\]


The normalizer and the marginal potential of datum \(m\) are


\[Z_m(a)=\int_{\mathbb R^{d_z}}e^{-U_m(a,z)}\,dz,\qquad V_m(a)=-\log Z_m(a),\qquad V(a)=\frac1M\sum_{m=1}^MV_m(a)=-\ell(a),\]


so \(V\) is minus the average log evidence \(\ell\). This section proves that \(V\) is \(\lambda\)-strongly convex, tends to \(+\infty\) at infinity and has exactly one minimizer \(a_\star\); that for every posterior family with log evidence \(-V\) the free energy has exactly one global minimizer \((a_\star,\pi_{a_\star})\); and that \(V\) grows quadratically away from \(a_\star\). The last step quantifies the effect of the tilting on the optimizer. The unique minimizer is the target of all convergence results of this report.



\textbf{Claim.} Let \(M\ge1\), \(\lambda>0\), and let every \(U_m\) be differentiable and \(\lambda\)-strongly convex. Then there is exactly one \(a_\star\in\mathbb R^{d_\theta}\) with \(V(a_\star)\le V(a)\) for all \(a\). Equivalently, the average log evidence \(\ell=-V\) has exactly one maximizer.




\textbf{Proof.} \emph{Continuity.} Each \(V_m\) is continuous (companion report: the normalizer is continuous and positive, and \(\log\) is continuous on \((0,\infty)\)); hence \(V\), a constant multiple of a finite sum, is continuous.




\emph{Existence.} Fix \(a_0\). By the previous step there is \(R\) with \(V(a)>V(a_0)\) for \(\|a\|\ge R\). The closed ball \(\bar B_R=\{\|a\|\le R\}\) is compact (Heine–Borel) and nonempty, since \(a_0\in\bar B_R\). A continuous function on a nonempty compact set attains its minimum, say at \(a_\star\in\bar B_R\). Then \(V(a_\star)\le V(a)\) for \(a\in\bar B_R\), and for \(\|a\|\ge R\) we have \(V(a)>V(a_0)\ge V(a_\star)\). Thus \(a_\star\) is a global minimizer.




\emph{Uniqueness.} Let \(a\) and \(b\) both be global minimizers, so \(V(a)=V(b)=\min V\). Strong convexity with \(\alpha=\beta=\frac12\) gives


\[V\Bigl(\frac{a+b}2\Bigr)\le\frac{V(a)+V(b)}2-\frac\lambda8\|a-b\|^2=\min V-\frac\lambda8\|a-b\|^2.\]


If \(a\ne b\), the right side is strictly smaller than \(\min V\) because \(\lambda>0\), which is impossible. Hence \(a=b\). \(\square\)

\begin{leanbox}{Lean: \texttt{existsUnique\_\allowbreak{}averageMarginalMinimizer}}
\leansource{\texttt{existsUnique\_\allowbreak{}averageMarginalMinimizer} \quad LatentDiffusion/\allowbreak{}Marginal.lean, lines 62–75}
\begin{leancode}
\leanline{0}{theorem\ existsUnique\_averageMarginalMinimizer}
\leanline{4}{(U\ :\ Fin\ M\ →\ IPLA.Joint\ dθ\ dz\ →\ ℝ)\ (hM\ :\ 0\ <\ M)}
\leanline{4}{(rho\ :\ ℝ)\ (hrho\ :\ 0\ <\ rho)\ (hdiff\ :\ ∀\ i,\ Differentiable\ ℝ\ (U\ i))}
\leanline{4}{(hsc\ :\ ∀\ i,\ StrongConvexOn\ univ\ rho\ (U\ i))\ :}
\leanline{4}{∃!\ θ,\ ∀\ θ',\ averageMarginalPotential\ U\ θ\ ≤\ averageMarginalPotential\ U\ θ'}
\end{leancode}
\leanlabel{How the Lean proof follows the argument}
The conclusion \texttt{\ensuremath{\exists}! \ensuremath{\theta}, \ensuremath{\forall} \ensuremath{\theta}', average\allowbreak{}Marginal\allowbreak{}Potential U \ensuremath{\theta} \ensuremath{\le} average\allowbreak{}Marginal\allowbreak{}Potential U \ensuremath{\theta}'} is “there exists a unique global minimizer”.



\begin{enumerate}


\item \texttt{hc} proves continuity from \texttt{IPLA.\allowbreak{}marginal\allowbreak{}Potential\_\allowbreak{}continuous}, \texttt{continuous\_\allowbreak{}finset\allowbreak{}Sum} and multiplication by a constant.

\item \texttt{hc.\allowbreak{}exists\_\allowbreak{}forall\_\allowbreak{}le} is Mathlib's \texttt{Continuous.\allowbreak{}exists\_\allowbreak{}forall\_\allowbreak{}le}: a continuous function that tends to \(+\infty\) along \texttt{cocompact} has a global minimum. Its coercivity input is the previous step. This packages the compactness argument of the proof.

\item Uniqueness: \texttt{Strong\allowbreak{}Convex\allowbreak{}On.\allowbreak{}strict\allowbreak{}Convex\allowbreak{}On} turns strong convexity with \texttt{hrho : 0 < rho} into strict convexity, and \texttt{Strict\allowbreak{}Convex\allowbreak{}On.\allowbreak{}eq\_\allowbreak{}of\_\allowbreak{}is\allowbreak{}Min\allowbreak{}On} states that a strictly convex function has at most one minimizer on a convex set; its proof is the midpoint argument above.

\end{enumerate}
\end{leanbox}
\section{\texorpdfstring{The density model and its unique optimizer}{The density model and its unique optimizer}}


This section constructs the tilted latent model from the joint potentials \(U_1,\dots,U_M\) of assumption A3 and identifies its optimizer. For a parameter \(a\in\mathbb R^{d_\theta}\) and a datum \(m\), the unnormalized posterior density on \(\mathbb R^{d_z}\) is \(\rho_a^m(z)=e^{-U_m(a,z)}\), its normalizer is \(Z_m(a)=\int\rho_a^m\,dz\), and the posterior is the probability measure


\[\pi_a^m(dz)=\frac{\rho_a^m(z)}{Z_m(a)}\,dz.\]


The log evidence is \(\ell(a)=M^{-1}\sum_m\log Z_m(a)=-V(a)\), with \(V\) the average marginal potential of the previous section. We show that measurability, positivity, integrability and positivity of the normalizers follow from A3, so the construction requires no further hypothesis; that the free energy \(F(a,q)=-\ell(a)+M^{-1}\sum_m\mathrm{KL}(q^m\|\pi_a^m)\) of this model has exactly one global minimizer \((a_\star,\pi_{a_\star})\); and that every posterior has a finite second moment. The last step transfers uniqueness to the paper's literal density integral


\[\widetilde F(a,q)=\frac1M\sum_{m=1}^M\int_{\mathbb R^{d_z}}\tilde q^m(z)\log\frac{\tilde q^m(z)}{\rho_a^m(z)}\,dz\]


on the domain of collections \(q\) with finite relative entropies and finite second moments; here \(\tilde q^m\) is the Lebesgue density of \(q^m\).




This domain is the corrected form of the paper's \(\mathcal P_2\) domain (extra assumption R1).




The optimizer \(a_\star\) is the target of both main convergence theorems.



\textbf{Claim.} Let \(M\ge1\), \(\lambda>0\), and let every \(U_m\) be differentiable and \(\lambda\)-strongly convex. For the density model of the previous step there is exactly one pair \((a,q)\), with \(a\in\mathbb R^{d_\theta}\) and \(q\) a collection of \(M\) probability measures on \(\mathbb R^{d_z}\), that is a global minimizer of


\[F(a,q)=-\ell(a)+\frac1M\sum_{m=1}^M\mathrm{KL}(q^m\|\pi_a^m)\]


over admissible pairs (pairs with finite relative entropies). This pair is \((a_\star,\pi_{a_\star})\), where \(a_\star\) is the unique minimizer of \(V\).




\textbf{Proof.} The marginal section proves: if \(M\ge1\), \(\lambda>0\), every \(U_m\) is differentiable and \(\lambda\)-strongly convex, and \((\ell,\pi)\) is a posterior family with \(\ell=-V\), then \(F\) has exactly one global minimizer, \((a_\star,\pi_{a_\star})\). The posterior family of the density model satisfies \(\ell=-V\) by the previous step. \(\square\)

\begin{leanbox}{Lean: \texttt{existsUnique\_\allowbreak{}variationalOptimizer}}
\leansource{\texttt{existsUnique\_\allowbreak{}variationalOptimizer} \quad LatentDiffusion/\allowbreak{}RegularModel.lean, lines 66–74}
\begin{leancode}
\leanline{0}{theorem\ existsUnique\_variationalOptimizer\ (U\ :\ Fin\ M\ →\ IPLA.Joint\ dθ\ dz\ →\ ℝ)}
\leanline{4}{(hM\ :\ 0\ <\ M)\ (m\ :\ ℝ)\ (hm\ :\ 0\ <\ m)}
\leanline{4}{(hdiff\ :\ ∀\ i,\ Differentiable\ ℝ\ (U\ i))}
\leanline{4}{(hsc\ :\ ∀\ i,\ StrongConvexOn\ univ\ m\ (U\ i))\ :}
\leanline{4}{∃!\ p\ :\ EuclideanSpace\ ℝ\ (Fin\ dθ)\ ×}
\leanline{8}{(Fin\ M\ →\ ProbabilityMeasure\ (EuclideanSpace\ ℝ\ (Fin\ dz))),}
\leanline{6}{(densityModel\ U\ m\ hm\ hdiff\ hsc).family.IsGlobalMin\ p.1\ p.2}
\end{leancode}
\leanlabel{How the Lean proof follows the argument}
\texttt{exists\allowbreak{}Unique\_\allowbreak{}variational\allowbreak{}Optimizer} is a one-line application of \texttt{Latent\allowbreak{}Diffusion.\allowbreak{}exists\allowbreak{}Unique\_\allowbreak{}variational\allowbreak{}Optimizer} from the marginal section to the family \texttt{(density\allowbreak{}Model U m hm hdiff hsc).\allowbreak{}family}, with the identification \texttt{density\allowbreak{}Model\_\allowbreak{}log\allowbreak{}Evidence U m hm hdiff hsc} as the hypothesis \(\ell=-V\). The strong-convexity constant is named \texttt{m} in this file.
\end{leanbox}
\section{\texorpdfstring{De Bruijn identity for the joint flow}{De Bruijn identity for the joint flow}}


This section proves the de Bruijn identity for the joint parameter–density flow of the paper,.




The extra assumptions are R2 and R3.




The latent space \(E\) is a finite-dimensional real inner product space with Lebesgue measure, and the parameter space \(A\) is a complete real inner product space. For probability densities \(q^1,\dots,q^M\) on \(E\), potentials \(U_m:A\times E\to\mathbb R\) and a parameter \(a\), define the normalized density energy, the averaged parameter force and the Fisher information


\[\mathcal E(a,q)=\frac1M\sum_{m=1}^M\int_Eq^m(\log q^m+U_m(a,\cdot))\,dz,\qquad G(a,q)=\frac1M\sum_{m=1}^M\int_Eq^m\,\nabla_aU_m(a,\cdot)\,dz,\]
\[I(a,q)=\|G(a,q)\|^2+\frac1M\sum_{m=1}^M\int_Eq^m\,\bigl\|\nabla_z\bigl(\log q^m+U_m(a,\cdot)\bigr)\bigr\|^2\,dz.\]


When \(q^m\) is the density of a probability measure \(Q^m\) and the integrals are finite, \(\int q^m(\log q^m+U_m(a,\cdot))\,dz=\mathrm{KL}(Q^m\|\pi_a^m)-\log Z_m(a)\), so \(\mathcal E(a,q)\) is the free energy \(F(a,Q)\) of the variational section, written as the paper's literal integral. The flow is the parameter equation \(\dot a_t=-G(a_t,q_t)\) together with the Fokker–Planck equations \(\partial_tq_t^m=\nabla\cdot\bigl(q_t^m\nabla(\log q_t^m+U_m(a_t,\cdot))\bigr)\). The main result is


\[\frac{d}{dt}\,\mathcal E(a_t,q_t)=-I(a_t,q_t).\]


Its proof has three ingredients: differentiation under the integral sign, justified by local domination; Lebesgue integration by parts on the unbounded space \(E\), justified by integrability of three explicit products; and the chain rule in the parameter. Along the way we obtain the gradient of \(\mathcal E\) for the normalized product metric \(d^2=\|a-b\|^2+M^{-1}\sum_mW_2(q^m,r^m)^2\): its parameter component is \(G\) and its \(m\)-th density component is \(\nabla(\log q^m+U_m)\), with no factor \(1/M\). The identity is used in the convergence theorem for the classical flow; the integration-by-parts identities are also used for the entropy tangent inequality of the transport section.



The Fokker–Planck flux of a density \(q\) and a potential \(U\) is \(J_{q,U}=q\,\nabla(\log q+U)\). The Fisher information \(I(a,q)\) is defined in the summary; it is nonnegative whenever every \(q^m\ge0\), since it is a squared norm plus an average of integrals of nonnegative functions.




\textbf{Claim.} Assume hypotheses (1)–(4) of the previous step. Assume further the parameter equation and the Fokker–Planck equations at time \(t\),


\[\dot p(t)=-G(p(t),q_t),\qquad\dot q^m_t=\nabla\cdot J_{q^m_t,U_m(p(t),\cdot)}\ \text{ almost everywhere }(m=1,\dots,M),\]


and the no-boundary-flux condition for \(\bigl(J_{q^m_t,U_m(p(t),\cdot)},\mu_m+1\bigr)\) for every \(m\). Then



\begin{enumerate}


\item every function \(q^m_t\,\|\nabla\mu_m\|^2\) is integrable, and

\item \(u\mapsto\mathcal E(p(u),q_u)\) is differentiable at \(t\) with derivative \(-I(p(t),q_t)\).

\end{enumerate}



\textbf{Proof.} Put \(v_m=-\nabla\mu_m\). Then \(q^m_tv_m=-J^m\), where \(J^m\) is the flux of datum \(m\). By linearity of the divergence, \(-\nabla\cdot(q^m_tv_m)=\nabla\cdot J^m=\dot q^m_t\) almost everywhere, so the continuity equation of the previous step holds. The no-boundary-flux condition is invariant under \(V\mapsto-V\), so it holds for \((q^m_tv_m,\mu_m+1)\). The previous step therefore gives the derivative


\[\Pi=\langle\dot p(t),G\rangle+\frac1M\sum_m\int q^m_t\langle-\nabla\mu_m,\nabla\mu_m\rangle\,dz=-\|G\|^2-\frac1M\sum_m\int q^m_t\|\nabla\mu_m\|^2\,dz=-I,\]


using the parameter equation in the first term. For (1), the integration-by-parts step shows that \(\langle J^m,\nabla(\mu_m+1)\rangle\) is integrable; since \(\nabla(\mu_m+1)=\nabla\mu_m\), this function is \(q^m_t\langle\nabla\mu_m,\nabla\mu_m\rangle=q^m_t\|\nabla\mu_m\|^2\). \(\square\)




The identity is proved from the equations; no entropy-dissipation inequality is assumed. Statement (1) shows that the Fisher information is a finite integral at every time where the hypotheses hold.

\begin{leanbox}{Lean: \texttt{normalized\_\allowbreak{}product\_\allowbreak{}deBruijn}}
\leansource{\texttt{normalized\_\allowbreak{}product\_\allowbreak{}deBruijn} \quad LatentDiffusion/\allowbreak{}RegularFlowCalculus.lean, lines 613–663}
\begin{leancode}
\leanline{0}{theorem\ normalized\_product\_deBruijn}
\leanline{4}{(q\ qdot\ :\ Fin\ M\ →\ ℝ\ →\ E\ →\ ℝ)\ (U\ :\ Fin\ M\ →\ A\ →\ E\ →\ ℝ)}
\leanline{4}{(p\ pdot\ :\ ℝ\ →\ A)\ (t\ :\ ℝ)}
\leanline{4}{(data\ :\ ∀\ i,\ LocalEnergyDomination\ volume\ (q\ i)\ (qdot\ i)\ (fun\ u\ ↦\ U\ i\ (p\ u))}
\leanline{6}{(fun\ u\ x\ ↦\ fderiv\ ℝ\ (fun\ a\ ↦\ U\ i\ a\ x)\ (p\ u)\ (pdot\ u))\ t)}
\leanline{4}{(hp\ :\ ∀\ u,\ HasDerivAt\ p\ (pdot\ u)\ u)}
\leanline{4}{(hparam\ :\ ∀\ i\ x\ a,\ DifferentiableAt\ ℝ\ (fun\ b\ ↦\ U\ i\ b\ x)\ a)}
\leanline{4}{(hpos\ :\ ∀\ i\ x,\ 0\ <\ q\ i\ t\ x)\ (hq\ :\ ∀\ i,\ ContDiff\ ℝ\ 1\ (q\ i\ t))}
\leanline{4}{(hU\ :\ ∀\ i,\ ContDiff\ ℝ\ 1\ (U\ i\ (p\ t)))}
\leanline{4}{(hgrad\ :\ ∀\ i,\ Integrable}
\leanline{6}{(fun\ x\ ↦\ q\ i\ t\ x\ •\ gradient\ (fun\ b\ ↦\ U\ i\ b\ x)\ (p\ t))\ volume)}
\leanline{4}{(hode\ :\ pdot\ t\ =\ -averagedParameterForce\ volume\ (fun\ i\ ↦\ q\ i\ t)\ U\ (p\ t))}
\leanline{4}{(hpde\ :\ ∀\ i,\ ∀ᵐ\ x\ ∂volume,\ qdot\ i\ t\ x\ =}
\leanline{6}{fluxDivergence\ (fokkerPlanckFlux\ (q\ i\ t)\ (U\ i\ (p\ t)))\ x)}
\leanline{4}{(htail\ :\ ∀\ i,\ NoBoundaryFlux\ (fokkerPlanckFlux\ (q\ i\ t)\ (U\ i\ (p\ t)))}
\leanline{6}{(fun\ x\ ↦\ chemicalPotential\ (q\ i\ t)\ (U\ i\ (p\ t))\ x\ +\ 1))\ :}
\leanline{4}{(∀\ i,\ Integrable}
\leanline{6}{(fun\ x\ ↦\ q\ i\ t\ x\ *\ ‖gradient\ (chemicalPotential\ (q\ i\ t)\ (U\ i\ (p\ t)))\ x‖\ \textasciicircum{}\ 2)}
\leanline{6}{volume)\ ∧}
\leanline{4}{HasDerivAt\ (fun\ u\ ↦\ normalizedDensityEnergy\ volume\ (fun\ i\ ↦\ q\ i\ u)\ U\ (p\ u))}
\leanline{6}{(-normalizedSlopeSquared\ volume\ (fun\ i\ ↦\ q\ i\ t)\ U\ (p\ t))\ t}
\end{leancode}
\leanlabel{How the Lean proof follows the argument}
\begin{enumerate}


\item \texttt{fokker\allowbreak{}Planck\allowbreak{}Flux q U} is \texttt{fun x \ensuremath{\mapsto} q x \ensuremath{\bullet} gradient (chemical\allowbreak{}Potential q U) x}. \texttt{normalized\allowbreak{}Slope\allowbreak{}Squared} is \(I\), and \texttt{normalized\allowbreak{}Slope\allowbreak{}Squared\_\allowbreak{}nonneg} proves \(I\ge0\) from \texttt{sq\_\allowbreak{}nonneg} and \texttt{integral\_\allowbreak{}nonneg}.

\item \texttt{normalized\_\allowbreak{}product\_\allowbreak{}de\allowbreak{}Bruijn}: \texttt{hode} is the parameter equation, \texttt{hpde} the Fokker–Planck equations and \texttt{htail} the no-boundary-flux conditions. The velocity \texttt{v} is \(-\nabla\mu_m\); \texttt{hvel} is \(q v_m=-J^m\), \texttt{hcont} derives the continuity equation with \texttt{flux\allowbreak{}Divergence\_\allowbreak{}neg}, and \texttt{htailv} is \texttt{(htail i).\allowbreak{}neg}.

\item \texttt{hd} applies \texttt{normalized\_\allowbreak{}product\_\allowbreak{}gradient\_\allowbreak{}representation}, and \texttt{hrate} evaluates the pairing as \(-I\) by substituting \texttt{hode} and simplifying \(\langle-x,x\rangle=-\|x\|^2\) (\texttt{inner\_\allowbreak{}neg\_\allowbreak{}left}, \texttt{real\_\allowbreak{}inner\_\allowbreak{}self\_\allowbreak{}eq\_\allowbreak{}norm\_\allowbreak{}sq}).

\item Integrability (1) is \texttt{(htail i).\allowbreak{}integrable\_\allowbreak{}flux\_\allowbreak{}gradient}, rewritten with \texttt{gradient\_\allowbreak{}add\_\allowbreak{}one}, \texttt{inner\_\allowbreak{}smul\_\allowbreak{}left} and \texttt{real\_\allowbreak{}inner\_\allowbreak{}self\_\allowbreak{}eq\_\allowbreak{}norm\_\allowbreak{}sq}.

\end{enumerate}
\end{leanbox}
\section{\texorpdfstring{Strong displacement convexity of the free energy}{Strong displacement convexity of the free energy}}


This section proves the strong geodesic convexity of the tilted free energy along displacement interpolations, the paper's Lemma on geodesic convexity,.




The extra transport assumptions are R4.




Let \(Z\) be a finite-dimensional real inner product space of dimension \(d\), let \(\mu\) be a probability measure on \(Z\) (the initial law) and let \(T:Z\to Z\) be a transport map. The displacement interpolation is \(T_s(z)=(1-s)z+sT(z)\), \(s\in[0,1]\), and its law is the pushforward \(\mu_s=\mu\circ T_s^{-1}\). For one datum, a parameter interpolation \(p_s=(1-s)p_0+sp_1\) and a joint potential \(U(p,z)\), the free energy along the interpolation is


\[\mathcal F(s)=H(s)+\int_ZU(p_s,z)\,\mu_s(dz),\qquad H(s)=\int_Z\log\rho_s\,d\mu_s,\]


where \(\rho_s\) is the Lebesgue density of \(\mu_s\) and \(H\) is the negative Boltzmann entropy.




The extra transport assumptions are positive densities, positive eigenvalues of the Jacobian \(DT\) in an orthonormal eigenbasis, the change-of-variables formula \(\rho_s(T_s(z))=\rho_0(z)/\det\bigl((1-s)I+sDT(z)\bigr)\), and integrable logarithms; no convexity or transport inequality is among them.




If \(U\) is jointly \(\lambda\)-strongly convex, we prove


\[\mathcal F(s)\le(1-s)\mathcal F(0)+s\mathcal F(1)-\frac\lambda2s(1-s)\bigl(\|p_0-p_1\|^2+W_2(\mu_0,\mu_1)^2\bigr),\]


where \(W_2\) is the Wasserstein distance. The entropy part follows from concavity of \(\log\det\); the potential part from joint strong convexity. Consequences are a Talagrand-type quadratic growth at a minimizer and the tangent inequality used for the extended log-Sobolev inequality in the next section.



\textbf{Wasserstein distance.} For probability measures \(\mu,\eta\) on \(Z\), a coupling is a probability measure \(\gamma\) on \(Z\times Z\) with marginals \(\mu\) and \(\eta\), and


\[W_2(\mu,\eta)=\inf_\gamma\Bigl(\int\|x-y\|^2\,\gamma(dx,dy)\Bigr)^{1/2}\in[0,\infty].\]


\textbf{Claim.} Under the hypotheses of the previous step, for \(s\in[0,1]\):



\begin{enumerate}


\item \(\mathcal F(s)\le(1-s)\mathcal F(0)+s\mathcal F(1)-\frac\lambda2s(1-s)\bigl(\|p_0-p_1\|^2+\|\mathrm{id}-T\|_{L^2(\mu)}^2\bigr)\);

\item \(W_2(\mu,\mu_1)^2\le\|\mathrm{id}-T\|_{L^2(\mu)}^2\), and in particular \(W_2(\mu,\mu_1)<\infty\);

\item if moreover \(\lambda\ge0\), then \(\mathcal F(s)\le(1-s)\mathcal F(0)+s\mathcal F(1)-\frac\lambda2s(1-s)\bigl(\|p_0-p_1\|^2+W_2(\mu,\mu_1)^2\bigr)\).

\end{enumerate}



\textbf{Proof.} (1) Add the entropy inequality of the third step to the potential inequality of the previous step. At the endpoints, \(\mathcal F(0)=H(0)+\int U(p_0,\cdot)\,d\mu\) because \(\mu_0=\mu\), and \(\mathcal F(1)=H(1)+\int U(p_1,T(z))\,\mu(dz)\) by the pushforward formula for \(\mu_1=\mu\circ T^{-1}\).




(2) The law of the pair \((z,T(z))\) under \(\mu\) is a coupling of \(\mu\) and \(\mu_1\), and its cost is \(\int\|z-T(z)\|^2\,\mu(dz)\). The infimum over couplings is at most this value, which is finite by square integrability.




(3) When \(\lambda\ge0\) and \(s\in[0,1]\), the coefficient \(c=\frac\lambda2s(1-s)\) is nonnegative. By (2), \(c\bigl(\|p_0-p_1\|^2+W_2(\mu,\mu_1)^2\bigr)\le c\bigl(\|p_0-p_1\|^2+\|\mathrm{id}-T\|_{L^2(\mu)}^2\bigr)\), so the right side of (1) is at most the right side of (3), and (3) follows from (1). \(\square\)




The map \(T\) need not be optimal: the conclusion uses the genuine infimum over couplings, which can only be smaller than the cost of \(T\). With the normalized product distance of the paper, (3) is the strong geodesic convexity of the free energy of one datum.

\begin{leanbox}{Lean: \texttt{freeEnergy\_\allowbreak{}displacement\_\allowbreak{}strongConvex}}
\leansource{\texttt{freeEnergy\_\allowbreak{}displacement\_\allowbreak{}strongConvex} \quad LatentDiffusion/\allowbreak{}RegularTransport.lean, lines 333–349}
\begin{leancode}
\leanline{0}{theorem\ freeEnergy\_displacement\_strongConvex}
\leanline{4}{(D\ :\ RegularDisplacement\ ν\ μ\ d)\ (U\ :\ P\ →\ Z\ →\ ℝ)\ (p₀\ p₁\ :\ P)\ (rho\ :\ ℝ)}
\leanline{4}{(hrho\ :\ 0\ ≤\ rho)}
\leanline{4}{(hU\ :\ StrongConvexOn\ univ\ rho\ (fun\ w\ :\ PBGSO.ProductTransport.Joint\ P\ Z\ ↦\ U\ w.fst\ w.snd))}
\leanline{4}{(hmeas\ :\ ∀\ p,\ Measurable\ (U\ p))}
\leanline{4}{(hint\ :\ ∀\ s\ ∈\ Icc\ (0\ :\ ℝ)\ 1,}
\leanline{6}{Integrable\ (fun\ z\ ↦\ U\ (parameterInterpolate\ p₀\ p₁\ s)\ (interpolate\ D.transport\ s\ z))\ μ)}
\leanline{4}{(hcost\ :\ MemLp\ (fun\ z\ ↦\ z\ -\ D.transport\ z)\ 2\ μ)}
\leanline{4}{(s\ :\ ℝ)\ (hs\ :\ s\ ∈\ Icc\ (0\ :\ ℝ)\ 1)\ :}
\leanline{4}{freeEnergy\ D\ U\ p₀\ p₁\ s\ ≤}
\leanline{6}{(1\ -\ s)\ *\ freeEnergy\ D\ U\ p₀\ p₁\ 0\ +\ s\ *\ freeEnergy\ D\ U\ p₀\ p₁\ 1\ -}
\leanline{8}{rho\ /\ 2\ *\ s\ *\ (1\ -\ s)\ *}
\leanline{10}{(‖p₀\ -\ p₁‖\ \textasciicircum{}\ 2\ +\ ULA.wasserstein2Real\ μ\ (D.law\ 1)\ \textasciicircum{}\ 2)}
\end{leancode}
\leanlabel{How the Lean proof follows the argument}
\begin{enumerate}


\item \texttt{free\allowbreak{}Energy\_\allowbreak{}displacement\_\allowbreak{}strong\allowbreak{}Convex\_\allowbreak{}cost} adds \texttt{D.\allowbreak{}entropy\_\allowbreak{}displacement\_\allowbreak{}convex} and \texttt{potential\_\allowbreak{}displacement\_\allowbreak{}strong\allowbreak{}Convex}; \texttt{hmap} is the endpoint pushforward formula, and \texttt{linarith} combines the inequalities after \texttt{simp only} evaluates the interpolations at \(0\) and \(1\).

\item \texttt{wasserstein\_\allowbreak{}sq\_\allowbreak{}le\_\allowbreak{}displacement\_\allowbreak{}cost} applies \texttt{ULA.\allowbreak{}wasserstein2\_\allowbreak{}le\_\allowbreak{}random\_\allowbreak{}cost\_\allowbreak{}real} to the random variables \(\mathrm{id}\) and \(T\) on \((Z,\mu)\): the real Wasserstein distance between their laws is at most the \(L^2\) norm of their difference. Here \texttt{ULA.\allowbreak{}wasserstein2\allowbreak{}Real} is the real conversion of the coupling infimum \texttt{ULA.\allowbreak{}wasserstein2}. \texttt{sq\_\allowbreak{}le\_\allowbreak{}sq\ensuremath{{}_{0}}} squares the inequality.

\item \texttt{free\allowbreak{}Energy\_\allowbreak{}displacement\_\allowbreak{}strong\allowbreak{}Convex} takes \texttt{hrho : 0 \ensuremath{\le} rho}, chains the first result with the second, and uses \texttt{hc}, nonnegativity of \(\frac\lambda2s(1-s)\), in \texttt{mul\_\allowbreak{}le\_\allowbreak{}mul\_\allowbreak{}of\_\allowbreak{}nonneg\_\allowbreak{}left} and \texttt{sub\_\allowbreak{}le\_\allowbreak{}sub\_\allowbreak{}left}.

\end{enumerate}
\end{leanbox}
\section{\texorpdfstring{Extended log-Sobolev and Talagrand inequalities}{Extended log-Sobolev and Talagrand inequalities}}


This section proves the two functional inequalities that drive the convergence of the classical flow: the extended log-Sobolev inequality (the paper's Theorem stating that strong log-concavity implies extended LSI) and a Talagrand-type quadratic growth inequality, both for the normalized free energy of \(M\) data. The setting is that of the previous section, now for each datum \(m\): a probability law \(Q_m\) on \(Z\) with positive continuously differentiable density \(\rho_{m}\), regular displacement data \(T_m\) from \(Q_m\), and a jointly \(\lambda\)-strongly convex potential \(U_m\) on \(P\times Z\). With \(p_0\) the current parameter, \(p_1\) the target parameter and \(\mathcal F_m\) the free energy of datum \(m\) along the displacement, put \(\mathcal F=M^{-1}\sum_m\mathcal F_m\), and define


\[\begin{aligned}G&=\frac1M\sum_m\int\nabla_pU_m(p_0,z)\,Q_m(dz),\\ S_m&=\nabla_z\bigl(\log\rho_m+U_m(p_0,\cdot)\bigr),\\I&=\|G\|^2+\frac1M\sum_m\int\|S_m\|^2\,dQ_m.\end{aligned}\]


The main result is \(2\lambda\bigl(\mathcal F(0)-\mathcal F(1)\bigr)\le I\). In the convergence theorem \(\mathcal F(0)\) is the free energy of the current state and \(\mathcal F(1)\) that of the posterior optimizer, so this is \(2\lambda(F-F_\star)\le I\). The proof sums two tangent inequalities, one for the entropy (from the Jacobian formula, \(\log x\le x-1\) and integration by parts) and one for the potential (from strong convexity), and then completes a square in which the parameter forces are averaged over the data before the norm is squared. The last step proves \(\frac\lambda2d^2\le\mathcal F(0)-\mathcal F(1)\) when the endpoint minimizes \(\mathcal F\) along the displacement, where \(d\) is the normalized product distance. No first-order transport inequality or log-Sobolev inequality is assumed.



\textbf{Claim.} Let \(M\ge1\), \(\lambda>0\), and assume the transport hypotheses for the data set. Then


\[2\lambda\bigl(\mathcal F(0)-\mathcal F(1)\bigr)\le I.\]


\textbf{Proof.} For each datum \(m\) put


\[\begin{aligned}\gamma_m&=\mathcal F_m(0)-\mathcal F_m(1),&g_m&=\int G_{m,P}\,dQ_m,\\\kappa_m&=\int\langle S_m,T_m-\mathrm{id}\rangle\,dQ_m,\\b_m&=\|S_m\|_{L^2(Q_m)},&c_m&=\|\mathrm{id}-T_m\|_{L^2(Q_m)}.\end{aligned}\]


\emph{Tangent inequality.} The free-energy tangent inequality for datum \(m\), with its left side rewritten by the splitting step, reads


\[\langle g_m,p_1-p_0\rangle+\kappa_m+\frac\lambda2\bigl(\|p_0-p_1\|^2+c_m^2\bigr)\le\mathcal F_m(1)-\mathcal F_m(0),\]


that is, \(\gamma_m+\langle g_m,w\rangle+\kappa_m+\frac\lambda2(\|w\|^2+c_m^2)\le0\) with \(w=p_1-p_0\), since \(\|p_0-p_1\|=\|w\|\).




\emph{Pairing bound.} By the Cauchy–Schwarz bound of the previous section applied to \(-S_m\) and \(T_m-\mathrm{id}\), \(-\kappa_m=\int\langle-S_m,T_m-\mathrm{id}\rangle\,dQ_m\le b_mc_m\); here \(\|T_m-\mathrm{id}\|_{L^2}=\|\mathrm{id}-T_m\|_{L^2}\).




\emph{Completion of the square.} The previous step gives \(2\lambda M^{-1}\sum_m\gamma_m\le\|M^{-1}\sum_mg_m\|^2+M^{-1}\sum_mb_m^2\). The left side is \(2\lambda(\mathcal F(0)-\mathcal F(1))\). On the right, \(M^{-1}\sum_mg_m=G\) and \(b_m^2=\int\|S_m\|^2\,dQ_m\), so the right side is \(I\). \(\square\)

\begin{leanbox}{Lean: \texttt{extendedLSI}}
\leansource{\texttt{extendedLSI} \quad LatentDiffusion/\allowbreak{}RegularTransportFirstOrder.lean, lines 661–718}
\begin{leancode}
\leanline{0}{theorem\ extendedLSI\ (hM\ :\ 0\ <\ M)\ (hrho\ :\ 0\ <\ rho)\ :}
\leanline{4}{2\ *\ rho\ *\ (R.freeEnergy\ 0\ -\ R.freeEnergy\ 1)\ ≤\ R.fisherInformation}
\end{leancode}
\leanlabel{How the Lean proof follows the argument}
In \texttt{extended\allowbreak{}LSI} the local definitions \texttt{g}, \texttt{pair}, \texttt{b}, \texttt{c} and \texttt{gap} are \(g_m\), \(\kappa_m\), \(b_m\), \(c_m\) and \(\gamma_m\); \texttt{b} and \texttt{c} are real conversions of \texttt{e\allowbreak{}Lp\allowbreak{}Norm … 2}.



\begin{enumerate}


\item \texttt{htangent} applies \texttt{free\allowbreak{}Energy\_\allowbreak{}tangent} to the fields of \texttt{R} for datum \texttt{m}, rewrites its left side with \texttt{tangent\_\allowbreak{}pairing\_\allowbreak{}eq}, and uses \texttt{norm\_\allowbreak{}sub\_\allowbreak{}rev}.

\item \texttt{hpair} applies \texttt{integral\_\allowbreak{}inner\_\allowbreak{}le\_\allowbreak{}second\allowbreak{}Cost} to \texttt{-\allowbreak{}(R.\allowbreak{}latent\allowbreak{}Score m)} and the displacement, using \texttt{R.\allowbreak{}latent\allowbreak{}Score\_\allowbreak{}mem\allowbreak{}Lp}; \texttt{hvnorm} identifies the \(L^2\) norms of \(T-\mathrm{id}\) and \(\mathrm{id}-T\) with \texttt{e\allowbreak{}Lp\allowbreak{}Norm\_\allowbreak{}neg}.

\item \texttt{h} is \texttt{normalized\_\allowbreak{}extended\allowbreak{}LSI\_\allowbreak{}of\_\allowbreak{}tangents}, \texttt{hsq} replaces \(b_m^2\) by the integral through \texttt{second\allowbreak{}Cost\_\allowbreak{}sq\_\allowbreak{}eq\_\allowbreak{}integral}, and the final \texttt{simpa} unfolds \texttt{free\allowbreak{}Energy}, \texttt{fisher\allowbreak{}Information} and \texttt{parameter\allowbreak{}Force}.

\end{enumerate}
\end{leanbox}
\section{\texorpdfstring{Exponential convergence of the continuous flow}{Exponential convergence of the continuous flow}}


This section proves the paper's continuous-time convergence theorem (Theorem 1 and its appendix form) for classical solutions of the joint gradient flow. The flow consists of a parameter path \(a_t\) and positive densities \(q_t^1,\dots,q_t^M\) solving


\[\dot a_t=-\frac1M\sum_{m=1}^M\int\nabla_aU_m(a_t,z)\,q_t^m(z)\,dz,\qquad\partial_tq_t^m=\nabla\cdot\bigl(q_t^m\nabla(\log q_t^m+U_m(a_t,\cdot))\bigr),\]


The flow also satisfies extra assumptions R1–R4: local domination, integrable spatial tails, and, at each time, regular displacement data from the current laws to the laws of a fixed comparison state \(a_\star\).




Existence of such a classical solution is not constructed; the estimates hold for every solution in this class. Write \(E(t)=\mathcal E(a_t,q_t)-\mathcal E(a_\star,q_\star)\) for the energy gap to the comparison state and \(d_t\) for the normalized product distance


\[d_t^2=\|a_t-a_\star\|^2+\frac1M\sum_{m=1}^MW_2(q_t^m,\pi_{a_\star}^m)^2.\]


When \(a_\star\) maximizes the evidence and the comparison densities are the posteriors \(\pi_{a_\star}^m\), we prove


\[d_t\le\sqrt{\frac{2E(t)}\lambda}\le\sqrt{\frac{2E(0)}\lambda}\,e^{-\lambda t},\qquad\|a_t-a_\star\|\le\sqrt{\frac{2E(0)}\lambda}\,e^{-\lambda t}\qquad(t\ge0).\]


The ingredients are the de Bruijn identity \(E'=-I\), the extended log-Sobolev inequality \(2\lambda E\le I\), Grönwall's inequality, and the Talagrand inequality \(\frac\lambda2d_t^2\le E(t)\), whose hypothesis (minimality of the posterior state) follows from the Gibbs variational principle. The final theorem constructs \(a_\star\) as the unique maximizer of the evidence of the original potentials, so the only remaining hypotheses are the extra assumptions defining the flow.



\textbf{Claim (Theorem 1, corrected).} Let \(M\ge1\), \(\lambda>0\), and let \(U_1,\dots,U_M:\mathbb R^{d_\theta}\times\mathbb R^{d_z}\to\mathbb R\) be differentiable and \(\lambda\)-strongly convex (A3). Let \(\ell(a)=M^{-1}\sum_m\log\int e^{-U_m(a,z)}dz\). Then there is exactly one \(a_\star\in\mathbb R^{d_\theta}\) such that



\begin{enumerate}


\item \(\ell(a)\le\ell(a_\star)\) for every \(a\), and

\item for every classical flow with comparison parameter \(a_\star\), constant \(\lambda\) and potentials \(U_m\), whose comparison densities are the posteriors \(q_\star^m=e^{-U_m(a_\star,\cdot)}/\int e^{-U_m(a_\star,w)}dw\), and for every \(t\ge0\),
\[\|a_t-a_\star\|\le\sqrt{\frac{2(\mathcal E(0)-\mathcal E_\star)}\lambda}\,e^{-\lambda t}.\]

\end{enumerate}



\textbf{Proof.} By the density-model section, \(\ell=-V\), where \(V\) is the average marginal potential. By the marginal section, \(V\) has exactly one minimizer \(a_\star\); hence \(a_\star\) is the unique maximizer of \(\ell\), which proves (1). For (2), A3 gives integrability of \(e^{-U_m(a,\cdot)}\) and positivity of its integral for all \(a\) (marginal section), so the hypotheses of the previous steps hold, and the parameter bound of the previous step is (2). Uniqueness: if \(b\) satisfies (1) and (2), then in particular \(\ell(a)\le\ell(b)\) for all \(a\), so \(V(b)\le V(a)\) for all \(a\), and \(b=a_\star\) by uniqueness of the minimizer of \(V\). \(\square\)




By the second step, \(\mathcal E_\star=-\ell(a_\star)=F_\star\) is the minimal free energy, so the prefactor is \(\sqrt{2(F(a_0,q_0)-F_\star)/\lambda}\) as in the paper. The theorem contains no assumed dissipation identity, log-Sobolev inequality or transport inequality; the only hypotheses beyond A3 are the extra assumptions defining a classical flow.

\begin{leanbox}{Lean: \texttt{existsUnique\_\allowbreak{}optimizer\_\allowbreak{}and\_\allowbreak{}flow\_\allowbreak{}convergence}}
\leansource{\texttt{existsUnique\_\allowbreak{}optimizer\_\allowbreak{}and\_\allowbreak{}flow\_\allowbreak{}convergence} \quad LatentDiffusion/\allowbreak{}RegularConvergence.lean, lines 478–517}
\begin{leancode}
\leanline{0}{theorem\ existsUnique\_optimizer\_and\_flow\_convergence\ \{dθ\ dz\ :\ ℕ\}}
\leanline{4}{(U\ :\ Fin\ M\ →\ IPLA.Joint\ dθ\ dz\ →\ ℝ)\ (hM\ :\ 0\ <\ M)}
\leanline{4}{(rho\ :\ ℝ)\ (hrho\ :\ 0\ <\ rho)}
\leanline{4}{(hdiff\ :\ ∀\ i,\ Differentiable\ ℝ\ (U\ i))}
\leanline{4}{(hsc\ :\ ∀\ i,\ StrongConvexOn\ univ\ rho\ (U\ i))\ :}
\leanline{4}{let\ V\ :=\ fun\ i\ a\ z\ ↦\ U\ i\ (IPLA.pair\ a\ z)}
\leanline{4}{let\ ell\ :=\ fun\ a\ ↦\ (M\ :\ ℝ)⁻¹\ *\ ∑\ i,\ Real.log\ (∫\ z,\ Real.exp\ (-V\ i\ a\ z))}
\leanline{4}{∃!\ astar\ :\ EuclideanSpace\ ℝ\ (Fin\ dθ),}
\leanline{6}{(∀\ a,\ ell\ a\ ≤\ ell\ astar)\ ∧}
\leanline{6}{∀\ S\ :\ ClassicalFlow\ M\ dz\ rho\ astar\ V,}
\leanline{8}{(∀\ i\ z,\ S.targetDensity\ i\ z\ =\ Real.exp\ (-V\ i\ astar\ z)/}
\leanline{10}{(∫\ w,\ Real.exp\ (-V\ i\ astar\ w)))\ →}
\leanline{8}{∀\ t,\ 0\ ≤\ t\ →\ ‖S.parameter\ t-astar‖\ ≤}
\leanline{10}{Real.sqrt\ (2*(S.energy\ 0-S.targetEnergy)/rho)*Real.exp\ (-rho*t)}
\end{leancode}
\leanlabel{How the Lean proof follows the argument}
In \texttt{exists\allowbreak{}Unique\_\allowbreak{}optimizer\_\allowbreak{}and\_\allowbreak{}flow\_\allowbreak{}convergence}, the local definitions \texttt{V} and \texttt{ell} are the slice potentials \(U_m(a,z)\) (written through \texttt{IPLA.\allowbreak{}pair}) and the evidence. The conclusion \texttt{\ensuremath{\exists}! astar, (\ensuremath{\forall} a, ell a \ensuremath{\le} ell astar) ∧ \ensuremath{\forall} S : Classical\allowbreak{}Flow M dz rho astar V, …} is the claim.



\begin{enumerate}


\item \texttt{he} identifies \texttt{ell} with \texttt{-\allowbreak{}average\allowbreak{}Marginal\allowbreak{}Potential U} through \texttt{Regular\allowbreak{}Model.\allowbreak{}density\allowbreak{}Model\_\allowbreak{}log\allowbreak{}Evidence}.

\item \texttt{exists\allowbreak{}Unique\_\allowbreak{}average\allowbreak{}Marginal\allowbreak{}Minimizer} supplies \texttt{astar} and its uniqueness \texttt{hu}; \texttt{hmax} is part (1).

\item Part (2) is \texttt{S.\allowbreak{}parameter\_\allowbreak{}exponential\_\allowbreak{}convergence}, with the integrability and positivity inputs \texttt{IPLA.\allowbreak{}integrable\_\allowbreak{}marginal\_\allowbreak{}weight} and \texttt{IPLA.\allowbreak{}marginal\allowbreak{}Integral\_\allowbreak{}pos}.

\item Uniqueness applies \texttt{hu} to a competitor \texttt{b}, using only its first property, rewritten with \texttt{he} and closed by \texttt{linarith}.

\end{enumerate}
\end{leanbox}
\section{\texorpdfstring{The normalized particle-cloud force}{The normalized particle-cloud force}}


Let \(K>0\) and let \(U_i\) be differentiable potentials on a parameter-latent pair, each \(m\)-strongly convex with \(L\)-Lipschitz gradient. Use scaled cloud coordinates \(z=(a,x_1/\sqrt K,\ldots,x_K/\sqrt K)\). Write \(J_i z=(a,x_i)\), \(V(z)=\sum_{i=1}^KU_i(J_i z)\), and \(F(z)=K^{-1}\nabla V(z)\). This normalization prevents an artificial deterioration of the monotonicity and smoothness constants when the particle count grows.



\textbf{Claim.} The normalized force \(F\) is \(m\)-strongly monotone and \(L\)-Lipschitz.



\textbf{Proof.} The coordinate maps satisfy \(\sum_i\|J_i d\|^2=K\|d\|^2\): the common parameter square appears \(K\) times and each scaled latent square is multiplied by \(K\). Summing the strong-convexity inequalities for the potentials therefore makes \(V\) \(Km\)-strongly convex. Differentiability follows by the chain rule and finite summation. For the gradient, the chain rule gives

\[\langle\nabla V(x)-\nabla V(y),w\rangle=\sum_i\langle\nabla U_i(J_ix)-\nabla U_i(J_iy),J_iw\rangle.\]

Lipschitz continuity, Cauchy–Schwarz for the finite sum, and the coordinate identity bound the absolute value by \(KL\|x-y\|\|w\|\). Taking the supremum over unit vectors gives \(\|\nabla V(x)-\nabla V(y)\|\le KL\|x-y\|\). Strong convexity similarly gives \(\langle\nabla V(x)-\nabla V(y),x-y\rangle\ge Km\|x-y\|^2\). Multiplication by \(1/K>0\) proves both assertions for \(F\). The argument allows distinct potentials and arbitrary finite parameter and latent dimensions. \(\square\)

\begin{leanbox}{Lean: \texttt{meanForce\_\allowbreak{}regular}}
\leansource{\texttt{meanForce\_\allowbreak{}regular} \quad LatentDiffusion/\allowbreak{}Cloud.lean, lines 233–253}
\begin{leancode}
\leanline{0}{theorem\ meanForce\_regular\ \{N\ dθ\ dx\ :\ ℕ\}\ (hN\ :\ 0\ <\ N)}
\leanline{4}{(U\ :\ Fin\ N\ →\ Joint\ dθ\ dx\ →\ ℝ)\ (m\ :\ ℝ)\ (L\ :\ ℝ≥0)}
\leanline{4}{(hdiff\ :\ ∀\ i,\ Differentiable\ ℝ\ (U\ i))}
\leanline{4}{(hsc\ :\ ∀\ i,\ StrongConvexOn\ Set.univ\ m\ (U\ i))}
\leanline{4}{(hLip\ :\ ∀\ i,\ LipschitzWith\ L\ (gradient\ (U\ i)))\ :}
\leanline{4}{LatentDiffusion.ForceRegularity\ (meanForce\ U)\ m\ L}
\end{leancode}
\leanlabel{How the Lean proof follows the argument}
\texttt{joint\_\allowbreak{}norm\_\allowbreak{}sq\_\allowbreak{}sum} expands the Euclidean coordinate sums. \texttt{mean\allowbreak{}Force\_\allowbreak{}regular} constructs the two fields of \texttt{Force\allowbreak{}Regularity}: it obtains monotonicity from \texttt{potential\_\allowbreak{}strong\allowbreak{}Convex} and \texttt{ULA.\allowbreak{}strong\_\allowbreak{}gradient\_\allowbreak{}monotone}, and smoothness from \texttt{potential\_\allowbreak{}gradient\_\allowbreak{}lipschitz}. In each branch the proof multiplies by the nonnegative reciprocal of the particle count and cancels it using \texttt{inv\_\allowbreak{}mul\_\allowbreak{}cancel\ensuremath{{}_{0}}}. The hypothesis \texttt{h\allowbreak{}N} ensures this reciprocal exists.
\end{leanbox}
\section{\texorpdfstring{Euler contraction and local consistency}{Euler contraction and local consistency}}


Work in a complete real inner-product space. The positive force \(F\) is \(m\)-strongly monotone and \(L\)-Lipschitz: \(\langle F(x)-F(y),x-y\rangle\ge m\|x-y\|^2\) and \(\|F(x)-F(y)\|\le L\|x-y\|\). Take \(m,h>0\), \(mh\le1\), and \(hL^2\le m\). The implemented force \(A\) can differ from \(F\). The numerical path is \(x_{k+1}=x_k-hA(x_k)+\xi_k\). A reference segment solves its actual integral equation with the same noise increment; its drift and noise sizes are explicit hypotheses, rather than an assumed numerical-error recurrence.



\textbf{Claim.} Under the force and step assumptions above, for any common increment \(\xi\), \(\|(x-hF(x)+\xi)-(y-hF(y)+\xi)\|\le q\|x-y\|\), where \(q=1-mh/2\in[0,1)\).



\textbf{Proof.} Put \(d=x-y\) and \(v=F(x)-F(y)\). The common increment cancels. Expanding the square and applying monotonicity and Lipschitz continuity gives

\[\|d-hv\|^2=\|d\|^2-2h\langle v,d\rangle+h^2\|v\|^2\le(1-2mh+h^2L^2)\|d\|^2\le(1-mh)\|d\|^2.\]

The last inequality uses \(hL^2\le m\), multiplied by \(h\ge0\). Since \(1-mh\le(1-mh/2)^2\) and \(q\ge0\), taking nonnegative square roots proves the claim. Positivity of \(mh\) gives \(q<1\). \(\square\)



\textbf{Claim.} Suppose every reference segment satisfies \(Y_k(0)=y_k\), \(Y_k(h)=y_{k+1}\), \(W_k(h)=\xi_k\), and

\[Y_k(t)-y_k=-\int_0^tF(Y_k(s))\,ds+W_k(t).\]

Assume the drift is interval integrable, \(\|F(Y_k(t))\|\le G\), and \(\|W_k(t)\|\le N\sqrt t\) for \(0\le t\le h\), with \(G,N\ge0\). Then the defect of a reference Euler step is at most \(R=L(h^2G/2+(2/3)h^{3/2}N)\).



\textbf{Proof.} The integral equation and the triangle inequality give \(\|Y_k(t)-y_k\|\le tG+N\sqrt t\). Subtract \(y_k-hF(y_k)+\xi_k\) from the endpoint equation; the noise increments cancel. The remaining vector is \(-\int_0^h(F(Y_k(s))-F(y_k))\,ds\). Its norm is at most

\[L\int_0^h\|Y_k(s)-y_k\|\,ds\le L\int_0^h(sG+N\sqrt s)\,ds=L\left(\frac{h^2G}{2}+\frac23h^{3/2}N\right).\]

The elementary integrals are \(h^2/2\) and \(2h^{3/2}/3\). In a random-state application the ambient norm is an \(L^2\) norm, so these same estimates are root-mean-square moment bounds. \(\square\)



\textbf{Claim.} With the preceding hypotheses, suppose \(\|A(x_k)-F(x_k)\|\le B\) for every \(k\), where \(B\ge0\). Then

\[\|x_n-y_n\|\le q^n\|x_0-y_0\|+\frac{2B}{m}+\frac Lm\left(hG+\frac43N\sqrt h\right).\]

\textbf{Proof.} Insert two intermediate points, namely the Euler updates with force \(F\) at \(x_k\) and \(y_k\). The triangle inequality, the force error, the contraction and the local defect imply \(e_{k+1}\le qe_k+hB+R\), where \(e_k=\|x_k-y_k\|\). Induction yields \(e_n\le q^ne_0+(hB+R)\sum_{j=0}^{n-1}q^j\). Since \(0\le q<1\) and the forcing term is nonnegative, the finite sum is at most \(1/(1-q)=2/(mh)\). Hence

\[\frac{hB+R}{1-q}=\frac{2B}{m}+\frac Lm\left(hG+\frac43N\sqrt h\right).\]

The division is legitimate because \(m,h>0\). This proves the bound uniformly in the iteration count. It does not construct the reference path or establish its moment bounds for a new application. \(\square\)

\begin{leanbox}{Lean: \texttt{integral\_\allowbreak{}reference\_\allowbreak{}error\_\allowbreak{}bound}}
\leansource{\texttt{integral\_\allowbreak{}reference\_\allowbreak{}error\_\allowbreak{}bound} \quad LatentDiffusion/\allowbreak{}Numerics.lean, lines 310–337}
\begin{leancode}
\leanline{0}{theorem\ integral\_reference\_error\_bound\ (F\ A\ :\ H\ →\ H)\ (m\ h\ B\ G\ N\ :\ ℝ)\ (L\ :\ ℝ≥0)}
\leanline{4}{(hF\ :\ ForceRegularity\ F\ m\ L)\ (hs\ :\ StableStep\ m\ L\ h)}
\leanline{4}{(hB\ :\ 0\ ≤\ B)\ (hG\ :\ 0\ ≤\ G)\ (hN\ :\ 0\ ≤\ N)}
\leanline{4}{(y\ :\ ℕ\ →\ H)\ (ξ\ :\ ℕ\ →\ H)\ (Y\ W\ :\ ℕ\ →\ ℝ\ →\ H)\ (x₀\ :\ H)}
\leanline{4}{(hr\ :\ IntegralReference\ F\ h\ G\ N\ y\ ξ\ Y\ W)}
\leanline{4}{(hforce\ :\ ∀\ n,\ ‖A\ (eulerTrajectory\ A\ h\ ξ\ x₀\ n)\ -\ F\ (eulerTrajectory\ A\ h\ ξ\ x₀\ n)‖\ ≤\ B)}
\leanline{4}{(n\ :\ ℕ)\ :}
\leanline{4}{‖eulerTrajectory\ A\ h\ ξ\ x₀\ n\ -\ y\ n‖\ ≤}
\leanline{6}{(1\ -\ m\ *\ h\ /\ 2)\ \textasciicircum{}\ n\ *\ ‖x₀\ -\ y\ 0‖\ +}
\leanline{8}{2\ *\ B\ /\ m\ +\ (L\ :\ ℝ)\ /\ m\ *\ (h\ *\ G\ +\ (4\ /\ 3\ :\ ℝ)\ *\ N\ *\ Real.sqrt\ h)}
\end{leancode}
\leanlabel{How the Lean proof follows the argument}
In \texttt{integral\_\allowbreak{}reference\_\allowbreak{}error\_\allowbreak{}bound}, \texttt{R} is the displayed local defect and \texttt{hstep} is the affine scalar recurrence obtained from \texttt{euler\_\allowbreak{}local\_\allowbreak{}comparison}. The call to \texttt{ULA.\allowbreak{}affine\_\allowbreak{}recurrence\_\allowbreak{}bound} sums that recurrence. Finally \texttt{halgebra} uses \texttt{field\_\allowbreak{}simp} and \texttt{ring} to replace its denominator by the three explicit error contributions.
\end{leanbox}
\section{\texorpdfstring{First-order condition for the evidence}{First-order condition for the evidence}}


For a reference measure \(\nu\), define \(Z_i(a)=\int e^{-U_i(a,z)}\,d\nu(z)\), \(\pi_{i,a}(dz)=Z_i(a)^{-1}e^{-U_i(a,z)}\nu(dz)\), and the averaged log evidence \(\ell(a)=M^{-1}\sum_i\log Z_i(a)\), where \(M>0\). The normalizers must be positive. Assume the densities are measurable at the point under consideration and almost everywhere strongly measurable for every parameter in some neighbourhood of that point; at the point they are integrable and their weighted gradients are almost everywhere strongly measurable; and on one common parameter neighbourhood the potentials are differentiable almost everywhere and \(e^{-U_i(b,z)}\|\nabla_bU_i(b,z)\|\) is bounded by an integrable function of \(z\). These domination hypotheses justify differentiation under the integral sign.



\textbf{Claim.} If \(a\) is a local maximizer of \(\ell\), then \(\sum_i\int\nabla_aU_i(a,z)\,d\pi_{i,a}(z)=0\).



\textbf{Proof.} Positivity of \(Z_i(a)\) permits the logarithmic chain rule. Dividing its derivative by \(Z_i(a)\) converts the weighted score integral into a posterior expectation, so \(\nabla\ell(a)=-M^{-1}\sum_i\int\nabla_aU_i(a,z)\,d\pi_{i,a}(z)\). A differentiable real function on the whole parameter space has zero derivative at a local maximum: for each direction its restriction to a real line has a local maximum at zero, whose derivative is zero. Thus the averaged score vanishes. Since \(M>0\), multiplication by \(M\) yields the unaveraged identity. Posterior integrability is furnished by the normalizer domination assumptions, rather than inferred from formal cancellation. \(\square\)

\begin{leanbox}{Lean: \texttt{sum\_\allowbreak{}posterior\_\allowbreak{}parameter\_\allowbreak{}gradient\_\allowbreak{}eq\_\allowbreak{}zero}}
\leansource{\texttt{sum\_\allowbreak{}posterior\_\allowbreak{}parameter\_\allowbreak{}gradient\_\allowbreak{}eq\_\allowbreak{}zero} \quad LatentDiffusion/\allowbreak{}RegularEvidence.lean, lines 185–193}
\begin{leancode}
\leanline{0}{theorem\ sum\_posterior\_parameter\_gradient\_eq\_zero\ \{M\ :\ ℕ\}\ (hM\ :\ 0\ <\ M)}
\leanline{4}{(ν\ :\ Measure\ Z)\ (U\ :\ Fin\ M\ →\ A\ →\ Z\ →\ ℝ)\ (a\ :\ A)}
\leanline{4}{(data\ :\ ∀\ i,\ LocalNormalizerRegularity\ ν\ (U\ i)\ a)}
\leanline{4}{(hmeas\ :\ ∀\ i,\ Measurable\ (fun\ z\ ↦\ Real.exp\ (-U\ i\ a\ z)))}
\leanline{4}{(hZ\ :\ ∀\ i,\ 0\ <\ normalizer\ ν\ (U\ i)\ a)}
\leanline{4}{(hmax\ :\ IsLocalMax\ (logEvidence\ ν\ U)\ a)\ :}
\leanline{4}{∑\ i,\ (∫\ z,\ gradient\ (fun\ b\ ↦\ U\ i\ b\ z)\ a\ ∂posterior\ ν\ (U\ i)\ a)\ =\ 0}
\end{leancode}
\leanlabel{How the Lean proof follows the argument}
\texttt{sum\_\allowbreak{}posterior\_\allowbreak{}parameter\_\allowbreak{}gradient\_\allowbreak{}eq\_\allowbreak{}zero} applies \texttt{averaged\_\allowbreak{}posterior\_\allowbreak{}parameter\_\allowbreak{}gradient\_\allowbreak{}eq\_\allowbreak{}zero} to the differentiable evidence. Its \texttt{smul\_\allowbreak{}eq\_\allowbreak{}zero} step removes the scalar reciprocal; \texttt{inv\_\allowbreak{}ne\_\allowbreak{}zero} and \texttt{h\allowbreak{}M} rule out a zero scalar. The local-maximality input is \texttt{hmax}; a zero-score hypothesis is not assumed.
\end{leanbox}
\section{\texorpdfstring{Euler recursion against a stationary posterior reference}{Euler recursion against a stationary posterior reference}}


Use the scaled cloud with \(K>0\), differentiable jointly \(m\)-strongly convex potentials, and \(L\)-Lipschitz gradients. Take \(m,h>0\), \(mh\le1\), \(hL^2\le m\), and put \(q=1-mh/2\). All random states lie in \(L^2(P)\), whose norm is \((\mathbb E\|X\|^2)^{1/2}\), for a probability measure \(P\). Supply a continuous stationary frozen-posterior integral solution \(S_t\), driven by the stated coordinate Brownian processes. Its parameter block at time zero is the constant \(a_\star\). Assume the initial per-particle parameter scores are independent and their expectations sum to zero. These are explicit stochastic and reference assumptions.



\textbf{Claim.} Let \(Q\) be extraction of the parameter block, with operator norm at most one. Put \(V=\sum_i\mathbb E\|s_i-\mathbb Es_i\|^2\), \(G=\|F_{\rm frozen}(S_0)\|_{L^2}\), and let \(W\) denote the physical Brownian increment scale. The coupled full-batch trajectory satisfies

\[\|QX_n-a_\star\|_{L^2}\le q^n\|X_0-S_0\|_{L^2}+\frac2m\frac{\sqrt V}{K}+\frac Lm\left(hG+\frac43W\sqrt h\right).\]

\textbf{Proof.} Write \(s_i=(s_i-\mathbb Es_i)+\mathbb Es_i\). The sum of the means vanishes. Expanding the squared norm of their average leaves diagonal terms plus cross inner products. Independence and centring make every cross expectation zero. Consequently the mean-square residual is exactly \(K^{-2}V\), whose nonnegative square root is \(\sqrt V/K\). Stationarity preserves this residual norm at every reference time. The frozen integral equation supplies the reference segments used in the numerical comparison; its drift norm is \(G\), and Brownian increments supply the square-root-time scale \(W\). Apply the equilibrium-reference Euler comparison with this residual bound. The parameter block of \(S_t\) is constant. Finally \(\|Q(X_n-S_{nh})\|_{L^2}\le\|X_n-S_{nh}\|_{L^2}\) transfers the cloud estimate to the parameter estimate. \(\square\)

\begin{leanbox}{Lean: \texttt{full\_\allowbreak{}batch\_\allowbreak{}parameter\_\allowbreak{}bound}}
\leansource{\texttt{full\_\allowbreak{}batch\_\allowbreak{}parameter\_\allowbreak{}bound} \quad LatentDiffusion/\allowbreak{}RegularStochastic.lean, lines 482–528}
\begin{leancode}
\leanline{0}{theorem\ full\_batch\_parameter\_bound}
\leanline{4}{(hK\ :\ 0\ <\ K)\ (U\ :\ Fin\ K\ →\ Cloud.Joint\ dθ\ dz\ →\ ℝ)}
\leanline{4}{(m\ h\ :\ ℝ)\ (L\ :\ ℝ≥0)\ (hdiff\ :\ ∀\ i,\ Differentiable\ ℝ\ (U\ i))}
\leanline{4}{(hsc\ :\ ∀\ i,\ StrongConvexOn\ Set.univ\ m\ (U\ i))}
\leanline{4}{(hLip\ :\ ∀\ i,\ LipschitzWith\ L\ (gradient\ (U\ i)))\ (hs\ :\ StableStep\ m\ L\ h)}
\leanline{4}{(B\ :\ (Fin\ K\ ×\ Fin\ dz)\ →\ ℝ≥0\ →\ Ω\ →\ ℝ)}
\leanline{4}{(hB\ :\ ∀\ i,\ IsPreBrownianReal\ (B\ i)\ P)}
\leanline{4}{(S\ :\ FrozenPosteriorSolution\ P\ (Cloud.meanForce\_regular\ hK\ U\ m\ L\ hdiff\ hsc\ hLip)\ B\ hB)}
\leanline{4}{(hind\ :\ iIndepFun\ (fun\ i\ ω\ =>\ scoreLp\ P\ U\ L\ hLip\ (S.path\ 0)\ i\ ω)\ P)}
\leanline{4}{(hopt\ :\ ∑\ i,\ (∫\ ω,\ scoreLp\ P\ U\ L\ hLip\ (S.path\ 0)\ i\ ω\ ∂P)\ =\ 0)}
\leanline{4}{(x₀\ :\ Lp\ (Cloud.State\ K\ dθ\ dz)\ 2\ P)\ (astar\ :\ Parameter\ dθ)}
\leanline{4}{(hstar\ :\ parameterOutput\ P\ (S.path\ 0)\ =\ constantL2\ P\ astar)\ (n\ :\ ℕ)\ :}
\leanline{4}{‖parameterOutput\ P}
\leanline{6}{(eulerTrajectory\ (particleForce\ P\ hK\ U\ m\ L\ hdiff\ hsc\ hLip)\ h}
\leanline{8}{(fun\ k\ =>\ physicalNoise\ P\ B\ hB\ ((k\ :\ ℝ)*h+h)\ -\ physicalNoise\ P\ B\ hB\ ((k\ :\ ℝ)*h))\ x₀\ n)\ -}
\leanline{6}{constantL2\ P\ astar‖\ ≤}
\leanline{6}{(1\ -\ m*h/2)\textasciicircum{}n\ *\ ‖x₀\ -\ S.path\ 0‖\ +}
\leanline{6}{(2/m)\ *\ Real.sqrt\ ((1/(K\ :\ ℝ))\textasciicircum{}2\ *\ scoreVarianceBudget\ P\ U\ L\ hLip\ (S.path\ 0))\ +}
\leanline{6}{(L\ :\ ℝ)/m\ *\ (h\ *\ ‖frozenForce\ P\ (Cloud.meanForce\_regular\ hK\ U\ m\ L\ hdiff\ hsc\ hLip)}
\leanline{8}{(S.path\ 0)‖\ +\ (4/3\ :\ ℝ)*physicalNoiseScale\ K\ dz*Real.sqrt\ h)}
\end{leancode}
\leanlabel{How the Lean proof follows the argument}
\texttt{full\_\allowbreak{}batch\_\allowbreak{}parameter\_\allowbreak{}bound} defines the residual size \texttt{b}. The block \texttt{hres} uses \texttt{residual\_\allowbreak{}norm\_\allowbreak{}eq\_\allowbreak{}of\_\allowbreak{}law}, \texttt{lifted\_\allowbreak{}mean\allowbreak{}Force\_\allowbreak{}residual} and \texttt{independent\_\allowbreak{}average\_\allowbreak{}norm\_\allowbreak{}sq} to prove the variance identity. It then calls \texttt{equilibrium\_\allowbreak{}reference\_\allowbreak{}error}. The equality \texttt{ht} identifies the reference parameter with \texttt{constant\allowbreak{}L2} of the optimizer; \texttt{parameter\allowbreak{}Output\_\allowbreak{}norm\_\allowbreak{}le} proves that the projection cannot enlarge the error. The displayed norm is therefore an actual root-mean-square parameter error.
\end{leanbox}
\section{\texorpdfstring{Mean-field and particle comparisons}{Mean-field and particle comparisons}}


Let the heterogeneous particle potentials be differentiable and jointly \(m\)-strongly convex, with \(L\)-Lipschitz gradients, and let the particle count be positive. Consider a continuous \(L^2(P)\) mean-field integral solution \(Y_t\) and a continuous stationary frozen-posterior solution \(S_t\) driven by the same physical Brownian path. Assume the posterior scores have zero total expectation, and the initial parameter of \(Y\) and the reference parameter of \(S\) are deterministic. The statement holds for any real \(m\); it gives decay when \(m>0\). Existence of a new solution is not asserted by a comparison of supplied solutions.



\textbf{Claim.} If the reference parameter is \(a_\star\), then for every \(t\ge0\), \(\|QY_t-a_\star\|_{L^2}\le e^{-mt}\|Y_0-S_0\|_{L^2}\).



\textbf{Proof.} Apply the parameter projection to the integral equation of \(Y\). The physical noise has zero parameter block, so its parameter remains the deterministic initial value minus the integral of the expected parameter force. The zero-score identity makes the frozen reference an admissible mean-field solution with constant parameter \(a_\star\). Subtract their two equations. Their common additive noise cancels. Strong monotonicity, including the deterministic-parameter expectation identities, gives the energy inequality \(\frac d{dt}\|Y_t-S_t\|_{L^2}^2\le-2m\|Y_t-S_t\|_{L^2}^2\). Multiply by \(e^{2mt}\); the derivative of the product is nonpositive. Integrating from zero to \(t\) gives \(\|Y_t-S_t\|_{L^2}^2\le e^{-2mt}\|Y_0-S_0\|_{L^2}^2\). Nonnegative square roots and the norm-one parameter projection finish the proof. Continuity and the supplied integral equations provide the time-integrability needed for this energy argument. \(\square\)

\begin{leanbox}{Lean: \texttt{meanField\_\allowbreak{}optimizer\_\allowbreak{}convergence}}
\leansource{\texttt{meanField\_\allowbreak{}optimizer\_\allowbreak{}convergence} \quad LatentDiffusion/\allowbreak{}RegularMeanField.lean, lines 560–587}
\begin{leancode}
\leanline{0}{theorem\ meanField\_optimizer\_convergence}
\leanline{4}{(hK\ :\ 0\ <\ K)\ (U\ :\ Fin\ K\ →\ Cloud.Joint\ dθ\ dz\ →\ ℝ)}
\leanline{4}{(m\ :\ ℝ)\ (L\ :\ ℝ≥0)\ (hdiff\ :\ ∀\ i,\ Differentiable\ ℝ\ (U\ i))}
\leanline{4}{(hsc\ :\ ∀\ i,\ StrongConvexOn\ Set.univ\ m\ (U\ i))}
\leanline{4}{(hLip\ :\ ∀\ i,\ LipschitzWith\ L\ (gradient\ (U\ i)))}
\leanline{4}{(B\ :\ (Fin\ K\ ×\ Fin\ dz)\ →\ ℝ≥0\ →\ Ω\ →\ ℝ)}
\leanline{4}{(hB\ :\ ∀\ i,\ IsPreBrownianReal\ (B\ i)\ P)}
\leanline{4}{(S\ :\ FrozenPosteriorSolution\ P\ (Cloud.meanForce\_regular\ hK\ U\ m\ L\ hdiff\ hsc\ hLip)\ B\ hB)}
\leanline{4}{(hopt\ :\ ∑\ i,\ (∫\ ω,\ scoreLp\ P\ U\ L\ hLip\ (S.path\ 0)\ i\ ω\ ∂P)=0)}
\leanline{4}{(astar\ :\ Parameter\ dθ)\ (hstar\ :\ parameterOutput\ P\ (S.path\ 0)=constantL2\ P\ astar)}
\leanline{4}{(Y\ :\ ℝ\ →\ Lp\ (Cloud.State\ K\ dθ\ dz)\ 2\ P)\ (hY\ :\ Continuous\ Y)}
\leanline{4}{(hYe\ :\ ∀\ t,\ 0\ ≤\ t\ →\ Y\ t-Y\ 0\ =}
\leanline{6}{-(∫\ r\ in\ 0..t,\ force\ P\ (Cloud.meanForce\_regular\ hK\ U\ m\ L\ hdiff\ hsc\ hLip)\ (Y\ r))+}
\leanline{8}{physicalNoise\ P\ B\ hB\ t)}
\leanline{4}{(a₀\ :\ Parameter\ dθ)\ (ha₀\ :\ parameterOutput\ P\ (Y\ 0)=constantL2\ P\ a₀)}
\leanline{4}{(t\ :\ ℝ)\ (ht\ :\ 0\ ≤\ t)\ :}
\leanline{4}{‖parameterOutput\ P\ (Y\ t)-constantL2\ P\ astar‖\ ≤\ Real.exp\ (-m*t)*‖Y\ 0-S.path\ 0‖}
\end{leancode}
\leanlabel{How the Lean proof follows the argument}
\texttt{mean\allowbreak{}Field\_\allowbreak{}optimizer\_\allowbreak{}convergence} first invokes \texttt{mean\allowbreak{}Field\_\allowbreak{}parameter\_\allowbreak{}deterministic} to establish \texttt{hpa}. The block \texttt{hpb} identifies the reference parameter; \texttt{frozen\_\allowbreak{}is\_\allowbreak{}mean\allowbreak{}Field} verifies its equation using \texttt{hopt}. The actual synchronous energy estimate is \texttt{mean\allowbreak{}Field\_\allowbreak{}contraction}. The final two inequalities apply \texttt{parameter\allowbreak{}Output\_\allowbreak{}norm\_\allowbreak{}le} to its norm estimate. In particular the proof does not replace stochastic independence or the reference equation by an assumed error bound.
\end{leanbox}
\section{\texorpdfstring{Root-mean-square error of the full-batch scheme}{Root-mean-square error of the full-batch scheme}}


Let \(K=MN>0\). Assume \(m>0\), differentiable jointly \(m\)-strongly convex heterogeneous potentials with \(L\)-Lipschitz gradients, and a stable positive step. The optimizer \(a_\star\) is the actual marginal optimizer already constructed from these potentials.




Supply an \(L^2\) initial latent cloud with the stated frozen Gibbs law at \(a_\star\), the correct individual posterior marginals, the local domination assumptions for evidence differentiation, and independent initial posterior scores. The coordinate Brownian processes have independent scalar Euler increments, independent of the initial latent cloud. Suppose each centred score has second moment at most \(V_0\). These extra initialization, independence and differentiation assumptions are part of the theorem.



\textbf{Claim.} For the actual pointwise Euler recursion,

\[\left(\mathbb E\|a_n-a_\star\|^2\right)^{1/2}\le q^nR_0+\frac2m\sqrt{\frac{V_0}{MN}}+\frac Lm\left(h\sqrt{L d_z}+\frac43\sqrt{2d_z}\sqrt h\right),\qquad q=1-mh/2,\]

where \(R_0\) is the \(L^2\) distance from the initial cloud to the scaled posterior cloud at \(a_\star\), and \(d_z\) is the dimension of one latent vector.



\textbf{Proof.} The evidence derivative at the optimizer centres the sum of posterior scores. The stipulated Gibbs initialization and independent Brownian increments furnish the stationary frozen reference used by the preceding stochastic comparison. Its frozen drift moment is bounded by \(\sqrt{Ld_z}\), and the scaled physical noise has increment norm \(\sqrt{2d_z}\sqrt t\). These facts give the three-term comparison with sampling residual \(K^{-1}\sqrt{\sum_i\operatorname{Var}(s_i)}\). The per-score bound gives \(\sum_i\operatorname{Var}(s_i)\le KV_0\), hence

\[\frac1K\sqrt{\sum_i\operatorname{Var}(s_i)}\le\frac{\sqrt{KV_0}}K=\sqrt{\frac{V_0}K}.\]

All factors are nonnegative and \(K>0\), so taking square roots and dividing preserves the inequality. Substitute \(K=MN\) and the explicit moment constants. The pointwise and lifted recursions agree in \(L^2\), which converts the projected norm to the displayed integral. No assertion of uniform complexity follows unless the constants stay bounded independently of particle count and step size. \(\square\)

\begin{leanbox}{Lean: \texttt{full\_\allowbreak{}batch\_\allowbreak{}optimizer\_\allowbreak{}rms\_\allowbreak{}rate}}
\leansource{\texttt{full\_\allowbreak{}batch\_\allowbreak{}optimizer\_\allowbreak{}rms\_\allowbreak{}rate} \quad LatentDiffusion/\allowbreak{}RegularSourceTheorem.lean, lines 167–209}
\begin{leancode}
\leanline{0}{theorem\ full\_batch\_optimizer\_rms\_rate}
\leanline{4}{\{M\ N\ :\ ℕ\}\ (hMN\ :\ 0\ <\ M*N)}
\leanline{4}{(U\ :\ Fin\ (M*N)\ →\ Cloud.Joint\ dθ\ dz\ →\ ℝ)}
\leanline{4}{(m\ h\ :\ ℝ)\ (L\ :\ ℝ≥0)\ (hm\ :\ 0\ <\ m)}
\leanline{4}{(hdiff\ :\ ∀\ i,\ Differentiable\ ℝ\ (U\ i))}
\leanline{4}{(hsc\ :\ ∀\ i,\ StrongConvexOn\ Set.univ\ m\ (U\ i))}
\leanline{4}{(hLip\ :\ ∀\ i,\ LipschitzWith\ L\ (gradient\ (U\ i)))\ (hs\ :\ StableStep\ m\ L\ h)}
\leanline{4}{(initial\ :\ Lp\ (Latent\ (M*N)\ dz)\ 2\ P)}
\leanline{4}{(hinitial\ :\ P.map\ initial\ =\ ULA.gibbsMeasure}
\leanline{6}{(frozenPotential\ U\ (optimizerParameter\ hMN\ U\ m\ hm\ hdiff\ hsc)))}
\leanline{4}{(hposterior\ :\ ∀\ i,\ P.map\ (fun\ ω\ =>\ physicalParticleCLM\ i}
\leanline{8}{((packLp\ P\ (optimizerParameter\ hMN\ U\ m\ hm\ hdiff\ hsc)\ initial)\ ω))\ =}
\leanline{6}{RegularEvidence.posterior\ volume\ (latentPotential\ U\ i)\ (optimizerParameter\ hMN\ U\ m\ hm\ hdiff\ hsc))}
\leanline{4}{(data\ :\ ∀\ i,\ RegularEvidence.LocalNormalizerRegularity\ volume\ (latentPotential\ U\ i)}
\leanline{6}{(optimizerParameter\ hMN\ U\ m\ hm\ hdiff\ hsc))}
\leanline{4}{(hindScores\ :\ iIndepFun\ (fun\ i\ ω\ =>\ scoreLp\ P\ U\ L\ hLip}
\leanline{6}{(packLp\ P\ (optimizerParameter\ hMN\ U\ m\ hm\ hdiff\ hsc)\ initial)\ i\ ω)\ P)}
\leanline{4}{(B\ :\ (Fin\ (M*N)\ ×\ Fin\ dz)\ →\ ℝ≥0\ →\ Ω\ →\ ℝ)}
\leanline{4}{(hB\ :\ ∀\ i,\ IsPreBrownianReal\ (B\ i)\ P)}
\leanline{4}{(hscalar\ :\ ∀\ s\ :\ ℝ,\ 0\ <\ s\ →\ iIndepFun\ (fun\ p\ :\ ℕ\ ×\ (Fin\ (M*N)\ ×\ Fin\ dz)\ =>\ fun\ ω\ =>}
\leanline{6}{ULA.Hilbert.brownianEulerNoise\ P\ B\ hB\ s\ p.1\ ω\ p.2)\ P)}
\leanline{4}{(hind\ :\ ∀\ s\ :\ ℝ,\ 0\ <\ s\ →\ IndepFun\ initial}
\leanline{6}{(fun\ ω\ n\ =>\ ULA.Hilbert.brownianEulerNoise\ P\ B\ hB\ s\ n\ ω)\ P)}
\leanline{4}{(V₀\ :\ ℝ)}
\leanline{4}{(hvar\ :\ ∀\ i,\ ‖scoreLp\ P\ U\ L\ hLip}
\leanline{8}{(packLp\ P\ (optimizerParameter\ hMN\ U\ m\ hm\ hdiff\ hsc)\ initial)\ i-}
\leanline{6}{constantL2\ P\ (∫\ ω,\ scoreLp\ P\ U\ L\ hLip}
\leanline{8}{(packLp\ P\ (optimizerParameter\ hMN\ U\ m\ hm\ hdiff\ hsc)\ initial)\ i\ ω\ ∂P)‖\textasciicircum{}2\ ≤\ V₀)}
\leanline{4}{(x₀\ :\ Lp\ (Cloud.State\ (M*N)\ dθ\ dz)\ 2\ P)\ (n\ :\ ℕ)\ :}
\leanline{4}{Real.sqrt\ (∫\ ω,\ ‖IPLA.parameterCLM}
\leanline{6}{(eulerTrajectory\ (Cloud.meanForce\ U)\ h}
\leanline{8}{(fun\ k\ =>\ noiseIncrement\ P\ B\ hB\ h\ k\ ω)\ (x₀\ ω)\ n)-}
\leanline{6}{optimizerParameter\ hMN\ U\ m\ hm\ hdiff\ hsc‖\textasciicircum{}2\ ∂P)\ ≤}
\leanline{6}{(1-m*h/2)\textasciicircum{}n\ *\ ‖x₀-packLp\ P\ (optimizerParameter\ hMN\ U\ m\ hm\ hdiff\ hsc)\ initial‖\ +}
\leanline{6}{(2/m)\ *\ Real.sqrt\ (V₀/((M*N\ :\ ℕ)\ :\ ℝ))\ +}
\leanline{6}{(L\ :\ ℝ)/m\ *\ (h*Real.sqrt\ ((L\ :\ ℝ)*(dz\ :\ ℝ))+}
\leanline{8}{(4/3\ :\ ℝ)*Real.sqrt\ (2*(dz\ :\ ℝ))*Real.sqrt\ h)}
\end{leancode}
\leanlabel{How the Lean proof follows the argument}
\texttt{full\_\allowbreak{}batch\_\allowbreak{}optimizer\_\allowbreak{}rms\_\allowbreak{}rate} obtains the actual root-mean-square comparison from \texttt{full\_\allowbreak{}batch\_\allowbreak{}optimizer\_\allowbreak{}rms}. It obtains the finite-particle square-root rate from \texttt{score\_\allowbreak{}sampling\_\allowbreak{}rate} and multiplies by the nonnegative factor \texttt{2/\allowbreak{}m} in \texttt{hv'}. The final \texttt{linarith} composes these estimates. The long hypothesis list explicitly identifies the posterior law, its differentiability data, the independent scores, Brownian increments, and their independence from the initial cloud; none is omitted by the short final arithmetic proof.
\end{leanbox}
\section{\texorpdfstring{Explicit accuracy choices}{Explicit accuracy choices}}


Let \(\varepsilon,m,K>0\), \(h,R_0,V,L,G,W\ge0\), and \(\sigma\in\mathbb R\). Suppose \(q=1-mh/2\ge0\) and an error \(E\) has the five-term upper bound \(q^nR_0+(2/m)\sqrt{V/K}+(L/m)hG+(4L/(3m))W\sqrt h+\sqrt{h\sigma^2/m}\). The following arithmetic allocates \(\varepsilon/5\) to each term. It applies after the stochastic/\allowbreak{}reference theorem supplies this upper bound; the arithmetic does not establish those hypotheses.



\textbf{Claim.} The error satisfies \(E\le\varepsilon\) if

\[\begin{gathered}\log\max\{1,5R_0/\varepsilon\}\le mhn/2,\qquad 100V\le m^2\varepsilon^2K,\\5LhG\le m\varepsilon,\qquad400L^2W^2h\le9m^2\varepsilon^2,\qquad25h\sigma^2\le m\varepsilon^2.\end{gathered}\]

\textbf{Proof.} The inequality \(1-u\le e^{-u}\), with \(u=mh/2\), and \(q\ge0\) imply \(q^nR_0\le e^{-mhn/2}R_0\). The logarithmic restriction makes this at most \(\varepsilon/5\); taking a maximum with one also covers \(R_0=0\) and an initial error already below tolerance. For the particle term, square its desired nonnegative inequality:

\[\left(\frac2m\sqrt{\frac VK}\right)^2\le\frac{\varepsilon^2}{25}\quad\Longleftrightarrow\quad100V\le m^2\varepsilon^2K.\]

The drift inequality is equivalent to its desired bound after multiplication by \(5m>0\). Squaring the nonnegative Brownian term gives exactly \(400L^2W^2h\le9m^2\varepsilon^2\). Squaring the oracle term gives exactly \(25h\sigma^2\le m\varepsilon^2\). Nonnegativity permits the square-root comparisons in each case. Thus the sum of all five terms is at most \(5\varepsilon/5=\varepsilon\). The division-free restrictions also make sense when variance, drift or noise constants vanish. \(\square\)

\begin{leanbox}{Lean: \texttt{budget\_\allowbreak{}le\_\allowbreak{}epsilon}}
\leansource{\texttt{budget\_\allowbreak{}le\_\allowbreak{}epsilon} \quad LatentDiffusion/\allowbreak{}Accuracy.lean, lines 53–82}
\begin{leancode}
\leanline{0}{theorem\ budget\_le\_epsilon}
\leanline{4}{(m\ h\ R0\ V\ L\ G\ W\ sigma\ epsilon\ K\ E\ :\ ℝ)\ (n\ :\ ℕ)}
\leanline{4}{(hm\ :\ 0\ <\ m)\ (hh\ :\ 0\ ≤\ h)\ (hR\ :\ 0\ ≤\ R0)\ (hV\ :\ 0\ ≤\ V)}
\leanline{4}{(hL\ :\ 0\ ≤\ L)\ (\_hG\ :\ 0\ ≤\ G)\ (hW\ :\ 0\ ≤\ W)\ (hK\ :\ 0\ <\ K)\ (he\ :\ 0\ <\ epsilon)}
\leanline{4}{(hq\ :\ 0\ ≤\ 1\ -\ m\ *\ h\ /\ 2)}
\leanline{4}{(hn\ :\ Real.log\ (max\ 1\ (5\ *\ R0\ /\ epsilon))\ ≤\ m\ *\ h\ *\ n\ /\ 2)}
\leanline{4}{(hparticles\ :\ 100\ *\ V\ ≤\ m\ \textasciicircum{}\ 2\ *\ epsilon\ \textasciicircum{}\ 2\ *\ K)}
\leanline{4}{(hdrift\ :\ 5\ *\ L\ *\ h\ *\ G\ ≤\ m\ *\ epsilon)}
\leanline{4}{(hnoise\ :\ 400\ *\ L\ \textasciicircum{}\ 2\ *\ W\ \textasciicircum{}\ 2\ *\ h\ ≤\ 9\ *\ m\ \textasciicircum{}\ 2\ *\ epsilon\ \textasciicircum{}\ 2)}
\leanline{4}{(horacle\ :\ 25\ *\ h\ *\ sigma\ \textasciicircum{}\ 2\ ≤\ m\ *\ epsilon\ \textasciicircum{}\ 2)}
\leanline{4}{(hbudget\ :\ E\ ≤\ (1\ -\ m\ *\ h\ /\ 2)\ \textasciicircum{}\ n\ *\ R0\ +\ (2\ /\ m)\ *\ Real.sqrt\ (V\ /\ K)\ +}
\leanline{6}{L\ /\ m\ *\ (h\ *\ G\ +\ (4\ /\ 3\ :\ ℝ)\ *\ W\ *\ Real.sqrt\ h)\ +}
\leanline{6}{Real.sqrt\ (h\ *\ sigma\ \textasciicircum{}\ 2\ /\ m))\ :\ E\ ≤\ epsilon}
\end{leancode}
\leanlabel{How the Lean proof follows the argument}
\texttt{budget\_\allowbreak{}le\_\allowbreak{}epsilon} names the five bounds \texttt{hb}, \texttt{hp}, \texttt{hd}, \texttt{hw} and \texttt{ho}. \texttt{burnin} proves the logarithmic estimate. \texttt{scaled\_\allowbreak{}sqrt\_\allowbreak{}le} handles the particle and Brownian terms by comparing squares; \texttt{Real.\allowbreak{}sqrt\_\allowbreak{}le\_\allowbreak{}left} handles the oracle term. Positive denominators are discharged by \texttt{positivity} and \texttt{field\_\allowbreak{}simp}. The last \texttt{nlinarith} adds the five allocations to \texttt{hbudget}.
\end{leanbox}
\chapter{Supporting proofs}
The following sections prove the results behind the boxes of the proof roadmap, in roadmap order.

\section{\texorpdfstring{The normalized variational problem}{The normalized variational problem}}


This section proves the variational characterization of the evidence. Let \(M\ge1\) be the number of data, let \(\mathcal A\) be an arbitrary set of parameters and let \(Z\) be a measurable space of latent variables. A \emph{posterior family} consists of a real function \(\ell:\mathcal A\to\mathbb R\), the log evidence, and, for each parameter \(a\in\mathcal A\) and each datum \(m\in\{1,\dots,M\}\), a probability measure \(\pi_a^m\) on \(Z\), the posterior of datum \(m\). For a collection \(q=(q^1,\dots,q^M)\) of probability measures on \(Z\), the free energy is


\[F(a,q)=-\ell(a)+\frac1M\sum_{m=1}^M\mathrm{KL}(q^m\|\pi_a^m),\]


where \(\mathrm{KL}\) denotes relative entropy. We prove that \(F(a,q)\ge-\ell(a)\), with equality exactly when \(q^m=\pi_a^m\) for every \(m\), and deduce the corrected form of the paper's Proposition 1: a pair \((a_\star,q_\star)\) minimizes \(F\) among all pairs with finite relative entropies if and only if \(a_\star\) maximizes \(\ell\) and \(q_\star^m=\pi_{a_\star}^m\) for every \(m\). No existence of a maximizer is asserted; existence and uniqueness are proved later from strong convexity of the marginal potential.




In the application \(\mathcal A\) is the space of parameters \(a=(\theta,\phi)\), \(\pi_a^m\) is the normalized tilted posterior \(Z_m(a)^{-1}e^{-U_m(a,z)}\,dz\) of datum \(m\), and \(\ell(a)=M^{-1}\sum_m\log Z_m(a)\) is the average tilted log evidence. The results hold for every parameter set and every measurable latent space, in particular for tilted models. The last step shows that a zero tilting penalty at a maximizer of the tilted evidence does not identify a maximizer of the original evidence, and gives a sufficient condition that does.

\subsection{\texorpdfstring{Relative entropy, the free energy and its admissible domain}{Relative entropy, the free energy and its admissible domain}}


\textbf{Definitions.} Let \(\mu\) and \(\nu\) be probability measures on \(Z\). If \(\mu\) is absolutely continuous with respect to \(\nu\) (written \(\mu\ll\nu\)), let \(d\mu/d\nu\) denote its Radon–Nikodym derivative. The relative entropy of \(\mu\) with respect to \(\nu\) is


\[\mathrm{KL}(\mu\|\nu)=\begin{cases}\displaystyle\int_Z\log\frac{d\mu}{d\nu}\,d\mu, & \text{if }\mu\ll\nu\text{ and }\log\frac{d\mu}{d\nu}\in L^1(\mu),\\ +\infty, & \text{otherwise}.\end{cases}\]


It takes values in \([0,\infty]\). Gibbs' inequality states that the integral is nonnegative, and that \(\mathrm{KL}(\mu\|\nu)=0\) holds if and only if \(\mu=\nu\).




A posterior family is a pair \((\ell,\pi)\) as in the summary. A collection \(q=(q^m)_{m=1}^M\) of probability measures on \(Z\) is called \emph{admissible at} \(a\in\mathcal A\) if


\[\mathrm{KL}(q^m\|\pi_a^m)<\infty\qquad\text{for every }m\in\{1,\dots,M\}.\]


To obtain a real-valued function on all pairs, write \(k(\mu\|\nu)=\mathrm{KL}(\mu\|\nu)\) when this number is finite and \(k(\mu\|\nu)=0\) otherwise. The free energy is


\[F(a,q)=-\ell(a)+\frac1M\sum_{m=1}^Mk(q^m\|\pi_a^m),\qquad(1)\]


with \(1/M\) read as \(0\) when \(M=0\). On admissible pairs \(k\) coincides with \(\mathrm{KL}\), so (1) is the free energy of the summary. Every minimization statement below is restricted to admissible pairs; the convention for infinite values therefore never enters a conclusion about minimizers. Finally, \((a,q)\) is a \emph{global minimizer} if \(q\) is admissible at \(a\) and \(F(a,q)\le F(b,r)\) for every \(b\in\mathcal A\) and every collection \(r\) admissible at \(b\).

\begin{leanbox}{Lean: \texttt{PosteriorFamily}, \texttt{Admissible}, \texttt{freeEnergy} and 1 more}
\leansource{\texttt{PosteriorFamily} \quad LatentDiffusion/\allowbreak{}Variational.lean, lines 22–24}
\begin{leancode}
\leanline{0}{structure\ PosteriorFamily\ (A\ Z\ :\ Type*)\ [MeasurableSpace\ Z]\ (M\ :\ ℕ)\ where}
\leanline{2}{logEvidence\ :\ A\ →\ ℝ}
\leanline{2}{posterior\ :\ A\ →\ Fin\ M\ →\ ProbabilityMeasure\ Z}
\end{leancode}
\leansource{\texttt{Admissible} \quad LatentDiffusion/\allowbreak{}Variational.lean, lines 30–31}
\begin{leancode}
\leanline{0}{def\ Admissible\ (a\ :\ A)\ (q\ :\ Fin\ M\ →\ ProbabilityMeasure\ Z)\ :\ Prop\ :=}
\leanline{2}{∀\ m,\ InformationTheory.klDiv\ (q\ m\ :\ Measure\ Z)\ (P.posterior\ a\ m\ :\ Measure\ Z)\ ≠\ ⊤}
\end{leancode}
\leansource{\texttt{freeEnergy} \quad LatentDiffusion/\allowbreak{}Variational.lean, lines 35–37}
\begin{leancode}
\leanline{0}{def\ freeEnergy\ (a\ :\ A)\ (q\ :\ Fin\ M\ →\ ProbabilityMeasure\ Z)\ :\ ℝ\ :=}
\leanline{2}{-P.logEvidence\ a\ +\ (M\ :\ ℝ)⁻¹\ *}
\leanline{4}{∑\ m,\ (InformationTheory.klDiv\ (q\ m\ :\ Measure\ Z)\ (P.posterior\ a\ m\ :\ Measure\ Z)).toReal}
\end{leancode}
\leansource{\texttt{IsGlobalMin} \quad LatentDiffusion/\allowbreak{}Variational.lean, lines 39–40}
\begin{leancode}
\leanline{0}{def\ IsGlobalMin\ (a\ :\ A)\ (q\ :\ Fin\ M\ →\ ProbabilityMeasure\ Z)\ :\ Prop\ :=}
\leanline{2}{P.Admissible\ a\ q\ ∧\ ∀\ b\ r,\ P.Admissible\ b\ r\ →\ P.freeEnergy\ a\ q\ ≤\ P.freeEnergy\ b\ r}
\end{leancode}
\leanlabel{How the Lean proof follows the argument}
\texttt{Posterior\allowbreak{}Family A Z M} is a structure with two fields: \texttt{log\allowbreak{}Evidence : A \ensuremath{\to} \ensuremath{\mathbb R}} is \(\ell\), and \texttt{posterior : A \ensuremath{\to} Fin M \ensuremath{\to} Probability\allowbreak{}Measure Z} assigns to \(a\) and \(m\) the posterior \(\pi_a^m\). Here \texttt{Fin M} is the index set \(\{0,\dots,M-1\}\), which plays the role of \(\{1,\dots,M\}\), and \texttt{Probability\allowbreak{}Measure Z} is the type of probability measures on \(Z\); the coercion \texttt{(q m : Measure Z)} forgets the normalization property. The relative entropy is Mathlib's \texttt{Information\allowbreak{}Theory.\allowbreak{}kl\allowbreak{}Div}, valued in \texttt{\ensuremath{\mathbb R}\ensuremath{\ge}0\ensuremath{\infty}} \(=[0,\infty]\). It is defined as \(\infty\) unless \(\mu\ll\nu\) and the log-likelihood ratio is \(\mu\)-integrable, in which case it is the integral of the log-likelihood ratio plus \(\nu(Z)-\mu(Z)\); the last two terms cancel for probability measures, which gives the displayed formula.



\begin{enumerate}


\item \texttt{Admissible} is the condition \texttt{kl\allowbreak{}Div (q m) (P.\allowbreak{}posterior a m) ≠ \ensuremath{\top}} for every \texttt{m}, where \texttt{\ensuremath{\top}} is \(+\infty\).

\item \texttt{free\allowbreak{}Energy} is formula (1). The conversion \texttt{.\allowbreak{}to\allowbreak{}Real} sends \(+\infty\) to \(0\) and every finite value to itself, so it is exactly the function \(k\). In \texttt{(M : \ensuremath{\mathbb R})\ensuremath{{}^{-1}}} the natural number is cast to a real, and Lean's convention \(0^{-1}=0\) is the reading of \(1/M\) for \(M=0\).

\item \texttt{Is\allowbreak{}Global\allowbreak{}Min} is the conjunction of admissibility of \texttt{q} at \texttt{a} and the inequality against every admissible pair.

\end{enumerate}
\end{leanbox}
\subsection{\texorpdfstring{The posterior attains the variational lower bound}{The posterior attains the variational lower bound}}


\textbf{Claim.} Let \((\ell,\pi)\) be a posterior family with \(M\in\mathbb N\) data and let \(a\in\mathcal A\). Then



\begin{enumerate}


\item the posterior collection \(\pi_a=(\pi_a^m)_m\) is admissible at \(a\);

\item \(F(a,\pi_a)=-\ell(a)\);

\item \(-\ell(a)\le F(a,q)\) for every collection \(q\) of probability measures, admissible or not.

\end{enumerate}



\textbf{Proof.} (1) Let \(\nu\) be a probability measure. Then \(\nu\ll\nu\) and \(d\nu/d\nu=1\) holds \(\nu\)-almost everywhere, so \(\log(d\nu/d\nu)=0\) almost everywhere. This function is integrable and its integral is \(0\). Hence \(\mathrm{KL}(\nu\|\nu)=0<\infty\). Applying this with \(\nu=\pi_a^m\) for each \(m\) proves admissibility.




(2) By (1) every summand of (1) of the previous step vanishes, so \(F(a,\pi_a)=-\ell(a)+M^{-1}\cdot0=-\ell(a)\).




(3) Each number \(k(q^m\|\pi_a^m)\) is nonnegative: it is either a finite value of \(\mathrm{KL}\), which lies in \([0,\infty)\), or \(0\). The factor \(1/M\) is nonnegative, also when it is read as \(0\). Hence the second term of \(F(a,q)\) is nonnegative and \(F(a,q)\ge-\ell(a)\). \(\square\)




Rearranged, (3) is the evidence lower bound \(\ell(a)\ge-F(a,q)\), and (2) shows that it is attained at the posterior.

\begin{leanbox}{Lean: \texttt{posterior\_\allowbreak{}admissible}, \texttt{posterior\_\allowbreak{}freeEnergy}, \texttt{lower\_\allowbreak{}bound}}
\leansource{\texttt{posterior\_\allowbreak{}admissible} \quad LatentDiffusion/\allowbreak{}Variational.lean, lines 42–44}
\begin{leancode}
\leanline{0}{theorem\ posterior\_admissible\ (a\ :\ A)\ :\ P.Admissible\ a\ (P.posterior\ a)}
\end{leancode}
\leansource{\texttt{posterior\_\allowbreak{}freeEnergy} \quad LatentDiffusion/\allowbreak{}Variational.lean, lines 46–47}
\begin{leancode}
\leanline{0}{theorem\ posterior\_freeEnergy\ (a\ :\ A)\ :\ P.freeEnergy\ a\ (P.posterior\ a)\ =\ -P.logEvidence\ a}
\end{leancode}
\leansource{\texttt{lower\_\allowbreak{}bound} \quad LatentDiffusion/\allowbreak{}Variational.lean, lines 49–56}
\begin{leancode}
\leanline{0}{theorem\ lower\_bound\ (a\ :\ A)\ (q\ :\ Fin\ M\ →\ ProbabilityMeasure\ Z)\ :}
\leanline{4}{-P.logEvidence\ a\ ≤\ P.freeEnergy\ a\ q}
\end{leancode}
\leanlabel{How the Lean proof follows the argument}
\begin{enumerate}


\item \texttt{posterior\_\allowbreak{}admissible}: after \texttt{intro m}, \texttt{simp} rewrites \texttt{kl\allowbreak{}Div \ensuremath{\pi} \ensuremath{\pi}} to \texttt{0} using Mathlib's simplification lemma \texttt{Information\allowbreak{}Theory.\allowbreak{}kl\allowbreak{}Div\_\allowbreak{}self}, and closes the goal \(0\ne\infty\).

\item \texttt{posterior\_\allowbreak{}free\allowbreak{}Energy}: \texttt{simp [free\allowbreak{}Energy]} applies the same lemma to every summand, then removes the zero sum, the product with \(M^{-1}\) and the addition of zero.

\item \texttt{lower\_\allowbreak{}bound}: the term \texttt{h} proves that the averaged entropy term is nonnegative. \texttt{inv\_\allowbreak{}nonneg.\allowbreak{}mpr (Nat.\allowbreak{}cast\_\allowbreak{}nonneg \_\allowbreak{})} gives \(M^{-1}\ge0\), \texttt{Finset.\allowbreak{}sum\_\allowbreak{}nonneg} with \texttt{ENNReal.\allowbreak{}to\allowbreak{}Real\_\allowbreak{}nonneg} gives nonnegativity of the sum, and \texttt{mul\_\allowbreak{}nonneg} combines them. \texttt{linarith} then adds \(-\ell(a)\). No admissibility hypothesis appears, in agreement with (3).

\end{enumerate}
\end{leanbox}
\subsection{\texorpdfstring{Equality in the variational bound}{Equality in the variational bound}}


\textbf{Claim.} Let \(M\ge1\), let \(a\in\mathcal A\) and let \(q\) be admissible at \(a\). Then


\[F(a,q)=-\ell(a)\iff q^m=\pi_a^m\text{ for every }m\in\{1,\dots,M\}.\]


\textbf{Proof.} (\(\Leftarrow\)) If \(q=\pi_a\), the previous step gives \(F(a,q)=-\ell(a)\).




(\(\Rightarrow\)) Write \(k_m=\mathrm{KL}(q^m\|\pi_a^m)\). By admissibility each \(k_m\) is a finite number in \([0,\infty)\), and \(F(a,q)=-\ell(a)+M^{-1}\sum_mk_m\). If \(F(a,q)=-\ell(a)\), then \(M^{-1}\sum_mk_m=0\). Because \(M\ge1\), the factor \(M^{-1}\) is nonzero, so \(\sum_mk_m=0\). A finite sum of nonnegative real numbers vanishes only if every summand vanishes; hence \(k_m=0\) for every \(m\). By the equality case of Gibbs' inequality, \(\mathrm{KL}(q^m\|\pi_a^m)=0\) implies \(q^m=\pi_a^m\). \(\square\)




The hypothesis \(M\ge1\) is used only to cancel the factor \(1/M\). Admissibility is needed because an infinite relative entropy contributes \(0\) to the real-valued formula (1); without it, the vanishing of the real sum would not imply the vanishing of the relative entropies.

\begin{leanbox}{Lean: \texttt{equality\_\allowbreak{}iff}}
\leansource{\texttt{equality\_\allowbreak{}iff} \quad LatentDiffusion/\allowbreak{}Variational.lean, lines 60–86}
\begin{leancode}
\leanline{0}{theorem\ equality\_iff\ (hM\ :\ 0\ <\ M)\ (a\ :\ A)\ (q\ :\ Fin\ M\ →\ ProbabilityMeasure\ Z)}
\leanline{4}{(hq\ :\ P.Admissible\ a\ q)\ :}
\leanline{4}{P.freeEnergy\ a\ q\ =\ -P.logEvidence\ a\ ↔\ q\ =\ P.posterior\ a}
\end{leancode}
\leanlabel{How the Lean proof follows the argument}
In \texttt{equality\_\allowbreak{}iff} the hypothesis \texttt{h\allowbreak{}M : 0 < M} is \(M\ge1\) and \texttt{hq} is admissibility. The forward direction proceeds as in the proof.



\begin{enumerate}


\item \texttt{hsum} derives that the sum of the real relative entropies is zero: \texttt{hh} isolates \(M^{-1}\sum_mk_m=0\) from the hypothesis by \texttt{linarith}, and \texttt{mul\_\allowbreak{}eq\_\allowbreak{}zero} with \texttt{inv\_\allowbreak{}ne\_\allowbreak{}zero hm} removes the nonzero factor \(M^{-1}\).

\item \texttt{funext m} reduces equality of the collections to equality at each index. \texttt{Finset.\allowbreak{}sum\_\allowbreak{}eq\_\allowbreak{}zero\_\allowbreak{}iff\_\allowbreak{}of\_\allowbreak{}nonneg} turns the vanishing sum of nonnegative terms into the vanishing of the term at \texttt{m}.

\item \texttt{ENNReal.\allowbreak{}to\allowbreak{}Real\_\allowbreak{}eq\_\allowbreak{}zero\_\allowbreak{}iff} states that the real conversion of \(x\in[0,\infty]\) is \(0\) exactly when \(x=0\) or \(x=\infty\); \texttt{.\allowbreak{}resolve\_\allowbreak{}right (hq m)} excludes \(x=\infty\) by admissibility, so \texttt{hz} is \(\mathrm{KL}(q^m\|\pi_a^m)=0\).

\item The two \texttt{let\allowbreak{}I} lines record that both measures are probability measures, and Mathlib's \texttt{Information\allowbreak{}Theory.\allowbreak{}kl\allowbreak{}Div\_\allowbreak{}eq\_\allowbreak{}zero\_\allowbreak{}iff} (the equality case of Gibbs' inequality for finite measures) gives equality of the underlying measures. \texttt{Subtype.\allowbreak{}ext} lifts it to equality in \texttt{Probability\allowbreak{}Measure Z}.

\end{enumerate}



The reverse direction substitutes \(q=\pi_a\) (\texttt{subst q}) and applies \texttt{posterior\_\allowbreak{}free\allowbreak{}Energy}.
\end{leanbox}
\subsection{\texorpdfstring{Characterization of the global minimizers}{Characterization of the global minimizers}}


\textbf{Claim.} Let \(M\ge1\), \(a\in\mathcal A\), and let \(q\) be a collection of \(M\) probability measures on \(Z\). Then \((a,q)\) is a global minimizer of \(F\) if and only if


\[\ell(b)\le\ell(a)\ \text{ for every }b\in\mathcal A\qquad\text{and}\qquad q^m=\pi_a^m\ \text{ for every }m.\]


\textbf{Proof.} (\(\Rightarrow\)) Let \((a,q)\) be a global minimizer; in particular \(q\) is admissible at \(a\). The posterior collection \(\pi_a\) is admissible at \(a\), so minimality and the posterior identity of the second step give


\[F(a,q)\le F(a,\pi_a)=-\ell(a).\]


The lower bound of the second step gives \(F(a,q)\ge-\ell(a)\); hence \(F(a,q)=-\ell(a)\), and the previous step yields \(q=\pi_a\). Now let \(b\in\mathcal A\). Since \(\pi_b\) is admissible at \(b\), minimality gives


\[-\ell(a)=F(a,q)\le F(b,\pi_b)=-\ell(b),\]


that is, \(\ell(b)\le\ell(a)\).




(\(\Leftarrow\)) Suppose \(a\) maximizes \(\ell\) and \(q=\pi_a\). Then \(q\) is admissible at \(a\). For every \(b\) and every \(r\) admissible at \(b\), the posterior identity, the maximality of \(a\) and the lower bound at \(b\) give


\[F(a,\pi_a)=-\ell(a)\le-\ell(b)\le F(b,r).\]


Thus \((a,\pi_a)\) is a global minimizer. \(\square\)




The statement asserts neither that \(\ell\) has a maximizer nor that a maximizer is unique; both questions are settled later for strongly convex potentials. The normalized posterior collection is the unique optimal variational collection at each maximizer.

\begin{leanbox}{Lean: \texttt{global\_\allowbreak{}minimizers\_\allowbreak{}iff}}
\leansource{\texttt{global\_\allowbreak{}minimizers\_\allowbreak{}iff} \quad LatentDiffusion/\allowbreak{}Variational.lean, lines 90–106}
\begin{leancode}
\leanline{0}{theorem\ global\_minimizers\_iff\ (hM\ :\ 0\ <\ M)\ (a\ :\ A)}
\leanline{4}{(q\ :\ Fin\ M\ →\ ProbabilityMeasure\ Z)\ :}
\leanline{4}{P.IsGlobalMin\ a\ q\ ↔}
\leanline{6}{(∀\ b,\ P.logEvidence\ b\ ≤\ P.logEvidence\ a)\ ∧\ q\ =\ P.posterior\ a}
\end{leancode}
\leanlabel{How the Lean proof follows the argument}
\texttt{global\_\allowbreak{}minimizers\_\allowbreak{}iff} states the equivalence for \texttt{P.\allowbreak{}Is\allowbreak{}Global\allowbreak{}Min a q}; the right-hand side is the conjunction of \texttt{\ensuremath{\forall} b, P.\allowbreak{}log\allowbreak{}Evidence b \ensuremath{\le} P.\allowbreak{}log\allowbreak{}Evidence a} and \texttt{q = P.\allowbreak{}posterior a}.



\begin{enumerate}


\item Forward direction: \texttt{rintro \ensuremath{\langle}hq, hmin\ensuremath{\rangle}} splits the hypothesis into admissibility and minimality. \texttt{he} is the identity \(F(a,q)=-\ell(a)\), obtained by \texttt{le\_\allowbreak{}antisymm} from \texttt{hmin a (P.\allowbreak{}posterior a) (P.\allowbreak{}posterior\_\allowbreak{}admissible a)}, rewritten with \texttt{P.\allowbreak{}posterior\_\allowbreak{}free\allowbreak{}Energy}, and from \texttt{P.\allowbreak{}lower\_\allowbreak{}bound a q}. The second component is \texttt{(P.\allowbreak{}equality\_\allowbreak{}iff h\allowbreak{}M a q hq).\allowbreak{}mp he}. For the first component, \texttt{hb} applies minimality against \((b,\pi_b)\), and after rewriting with \texttt{he} and \texttt{P.\allowbreak{}posterior\_\allowbreak{}free\allowbreak{}Energy}, \texttt{linarith} turns \(-\ell(a)\le-\ell(b)\) into \(\ell(b)\le\ell(a)\).

\item Reverse direction: \texttt{rintro \ensuremath{\langle}hmax, rfl\ensuremath{\rangle}} substitutes \(q=\pi_a\). Admissibility is \texttt{P.\allowbreak{}posterior\_\allowbreak{}admissible a}. For an admissible pair \((b,r)\) the chain of the proof is \texttt{(neg\_\allowbreak{}le\_\allowbreak{}neg (hmax b)).\allowbreak{}trans (P.\allowbreak{}lower\_\allowbreak{}bound b r)}, after rewriting \(F(a,\pi_a)\) by \texttt{P.\allowbreak{}posterior\_\allowbreak{}free\allowbreak{}Energy}.

\end{enumerate}
\end{leanbox}
\subsection{\texorpdfstring{Evidence maximizers are the parameter components of minimizers}{Evidence maximizers are the parameter components of minimizers}}


\textbf{Claim.} Let \(M\ge1\) and \(a\in\mathcal A\). Then \(a\) maximizes \(\ell\) if and only if there is a collection \(q\) such that \((a,q)\) is a global minimizer of \(F\).




\textbf{Proof.} If \(a\) maximizes \(\ell\), the previous step applied with \(q=\pi_a\) shows that \((a,\pi_a)\) is a global minimizer. Conversely, if \((a,q)\) is a global minimizer for some \(q\), the forward direction of the previous step shows that \(a\) maximizes \(\ell\). \(\square\)

\begin{leanbox}{Lean: \texttt{maximizer\_\allowbreak{}iff\_\allowbreak{}exists\_\allowbreak{}minimizer}}
\leansource{\texttt{maximizer\_\allowbreak{}iff\_\allowbreak{}exists\_\allowbreak{}minimizer} \quad LatentDiffusion/\allowbreak{}Variational.lean, lines 108–114}
\begin{leancode}
\leanline{0}{theorem\ maximizer\_iff\_exists\_minimizer\ (hM\ :\ 0\ <\ M)\ (a\ :\ A)\ :}
\leanline{4}{(∀\ b,\ P.logEvidence\ b\ ≤\ P.logEvidence\ a)\ ↔\ ∃\ q,\ P.IsGlobalMin\ a\ q}
\end{leancode}
\leanlabel{How the Lean proof follows the argument}
In \texttt{maximizer\_\allowbreak{}iff\_\allowbreak{}exists\_\allowbreak{}minimizer} the forward direction supplies the witness \texttt{P.\allowbreak{}posterior a} and proves its minimality by \texttt{(P.\allowbreak{}global\_\allowbreak{}minimizers\_\allowbreak{}iff h\allowbreak{}M a \_\allowbreak{}).\allowbreak{}mpr \ensuremath{\langle}h, rfl\ensuremath{\rangle}}. The reverse direction destructures the witness with \texttt{rintro \ensuremath{\langle}q, hq\ensuremath{\rangle}} and takes the first component of \texttt{(P.\allowbreak{}global\_\allowbreak{}minimizers\_\allowbreak{}iff h\allowbreak{}M a q).\allowbreak{}mp hq}.
\end{leanbox}
\subsection{\texorpdfstring{A zero tilting gap at a tilted maximizer is not sufficient}{A zero tilting gap at a tilted maximizer is not sufficient}}


The paper optimizes a tilted evidence \(\tilde\ell\le\ell\). This step compares maximizers of \(\tilde\ell\) with those of the original evidence \(\ell\).




\textbf{Claim 1 (sufficient condition).} Let \(\ell,\tilde\ell:\mathcal A\to\mathbb R\) satisfy \(\tilde\ell(x)\le\ell(x)\) for all \(x\). Let \(b\) maximize \(\ell\) with \(\tilde\ell(b)=\ell(b)\), and let \(a\) maximize \(\tilde\ell\). Then \(\ell(a)=\ell(b)\); in particular \(a\) maximizes \(\ell\).




\textbf{Proof.} Maximality of \(a\) for \(\tilde\ell\), the zero gap at \(b\), the comparison \(\tilde\ell\le\ell\) at \(a\), and maximality of \(b\) for \(\ell\) give


\[\ell(b)=\tilde\ell(b)\le\tilde\ell(a)\le\ell(a)\le\ell(b).\]


All inequalities are therefore equalities, and \(\ell(a)=\ell(b)\). \(\square\)




\textbf{Claim 2 (counterexample).} On the two-point set \(\mathcal A=\{0,1\}\) let \(\ell(0)=0\), \(\ell(1)=1\) and let the gap \(g=\ell-\tilde\ell\) be \(g(0)=0\), \(g(1)=2\). Then \(g\ge0\), \(g(0)=0\), the point \(0\) maximizes \(\tilde\ell=\ell-g\), and \(\ell(0)<\ell(1)\).




\textbf{Proof.} The tilted values are \(\tilde\ell(0)=0\) and \(\tilde\ell(1)=1-2=-1\), so \(0\) maximizes \(\tilde\ell\). The remaining assertions are the numerical facts \(0\le0\), \(0\le2\) and \(0<1\). \(\square\)




Thus a zero penalty at the tilted maximizer itself does not imply that it maximizes the original evidence. Claim 1 shows what does: a maximizer of the original evidence at which the tilting gap vanishes.

\begin{leanbox}{Lean: \texttt{zero\_\allowbreak{}gap\_\allowbreak{}comparator}, \texttt{zero\_\allowbreak{}gap\_\allowbreak{}at\_\allowbreak{}tilted\_\allowbreak{}max\_\allowbreak{}not\_\allowbreak{}sufficient}}
\leansource{\texttt{zero\_\allowbreak{}gap\_\allowbreak{}comparator} \quad LatentDiffusion/\allowbreak{}Variational.lean, lines 119–126}
\begin{leancode}
\leanline{0}{theorem\ zero\_gap\_comparator\ \{ell\ tilted\ :\ A\ →\ ℝ\}\ (a\ b\ :\ A)}
\leanline{4}{(hupper\ :\ ∀\ x,\ tilted\ x\ ≤\ ell\ x)}
\leanline{4}{(htrue\ :\ ∀\ x,\ ell\ x\ ≤\ ell\ b)\ (hzero\ :\ tilted\ b\ =\ ell\ b)}
\leanline{4}{(hmax\ :\ ∀\ x,\ tilted\ x\ ≤\ tilted\ a)\ :}
\leanline{4}{ell\ a\ =\ ell\ b}
\end{leancode}
\leansource{\texttt{zero\_\allowbreak{}gap\_\allowbreak{}at\_\allowbreak{}tilted\_\allowbreak{}max\_\allowbreak{}not\_\allowbreak{}sufficient} \quad LatentDiffusion/\allowbreak{}Variational.lean, lines 130–136}
\begin{leancode}
\leanline{0}{theorem\ zero\_gap\_at\_tilted\_max\_not\_sufficient\ :}
\leanline{4}{let\ ell\ :\ Bool\ →\ ℝ\ :=\ fun\ b\ =>\ if\ b\ then\ 1\ else\ 0}
\leanline{4}{let\ gap\ :\ Bool\ →\ ℝ\ :=\ fun\ b\ =>\ if\ b\ then\ 2\ else\ 0}
\leanline{4}{(∀\ b,\ 0\ ≤\ gap\ b)\ ∧\ gap\ false\ =\ 0\ ∧}
\leanline{6}{(∀\ b,\ ell\ b\ -\ gap\ b\ ≤\ ell\ false\ -\ gap\ false)\ ∧\ ell\ false\ <\ ell\ true}
\end{leancode}
\leanlabel{How the Lean proof follows the argument}
\begin{enumerate}


\item \texttt{zero\_\allowbreak{}gap\_\allowbreak{}comparator}: \texttt{hupper} is \(\tilde\ell\le\ell\), \texttt{htrue} and \texttt{hzero} are maximality of \(b\) and the zero gap, and \texttt{hmax} is maximality of \(a\) for the tilted evidence. The proof rewrites \texttt{hmax b} with \texttt{hzero} to obtain \(\ell(b)\le\tilde\ell(a)\), then concludes by \texttt{le\_\allowbreak{}antisymm (htrue a) (h.\allowbreak{}trans (hupper a))}, which is the displayed chain.

\item \texttt{zero\_\allowbreak{}gap\_\allowbreak{}at\_\allowbreak{}tilted\_\allowbreak{}max\_\allowbreak{}not\_\allowbreak{}sufficient} uses the type \texttt{Bool} as the two-point set, with \texttt{false} in the role of \(0\) and \texttt{true} in the role of \(1\); the functions are defined by \texttt{if b then … else …}. After \texttt{dsimp}, the four conjuncts are proved by \texttt{norm\_\allowbreak{}num}, with a case split \texttt{cases b} for the two universally quantified ones.

\end{enumerate}
\end{leanbox}
\section{\texorpdfstring{Differentiation of the free energy along the classical flow}{Differentiation of the free energy along the classical flow}}


This section proves the de Bruijn identity for the joint parameter–density flow of the paper,.




The extra assumptions are R2 and R3.




The latent space \(E\) is a finite-dimensional real inner product space with Lebesgue measure, and the parameter space \(A\) is a complete real inner product space. For probability densities \(q^1,\dots,q^M\) on \(E\), potentials \(U_m:A\times E\to\mathbb R\) and a parameter \(a\), define the normalized density energy, the averaged parameter force and the Fisher information


\[\mathcal E(a,q)=\frac1M\sum_{m=1}^M\int_Eq^m(\log q^m+U_m(a,\cdot))\,dz,\qquad G(a,q)=\frac1M\sum_{m=1}^M\int_Eq^m\,\nabla_aU_m(a,\cdot)\,dz,\]
\[I(a,q)=\|G(a,q)\|^2+\frac1M\sum_{m=1}^M\int_Eq^m\,\bigl\|\nabla_z\bigl(\log q^m+U_m(a,\cdot)\bigr)\bigr\|^2\,dz.\]


When \(q^m\) is the density of a probability measure \(Q^m\) and the integrals are finite, \(\int q^m(\log q^m+U_m(a,\cdot))\,dz=\mathrm{KL}(Q^m\|\pi_a^m)-\log Z_m(a)\), so \(\mathcal E(a,q)\) is the free energy \(F(a,Q)\) of the variational section, written as the paper's literal integral. The flow is the parameter equation \(\dot a_t=-G(a_t,q_t)\) together with the Fokker–Planck equations \(\partial_tq_t^m=\nabla\cdot\bigl(q_t^m\nabla(\log q_t^m+U_m(a_t,\cdot))\bigr)\). The main result is


\[\frac{d}{dt}\,\mathcal E(a_t,q_t)=-I(a_t,q_t).\]


Its proof has three ingredients: differentiation under the integral sign, justified by local domination; Lebesgue integration by parts on the unbounded space \(E\), justified by integrability of three explicit products; and the chain rule in the parameter. Along the way we obtain the gradient of \(\mathcal E\) for the normalized product metric \(d^2=\|a-b\|^2+M^{-1}\sum_mW_2(q^m,r^m)^2\): its parameter component is \(G\) and its \(m\)-th density component is \(\nabla(\log q^m+U_m)\), with no factor \(1/M\). The identity is used in the convergence theorem for the classical flow; the integration-by-parts identities are also used for the entropy tangent inequality of the transport section.

\subsection{\texorpdfstring{Pointwise derivative of the energy integrand}{Pointwise derivative of the energy integrand}}


\textbf{Claim.} Let \(q,U:\mathbb R\to\mathbb R\) be differentiable at \(t\) with derivatives \(\dot q\) and \(\dot U\), and let \(q(t)>0\). Then \(s\mapsto q(s)\bigl(\log q(s)+U(s)\bigr)\) is differentiable at \(t\), with derivative


\[\dot q\,\bigl(\log q(t)+U(t)+1\bigr)+q(t)\,\dot U.\]


\textbf{Proof.} Since \(q(t)>0\), the chain rule for the logarithm gives that \(s\mapsto\log q(s)\) is differentiable at \(t\) with derivative \(\dot q/q(t)\). The sum rule gives the derivative \(\dot q/q(t)+\dot U\) of \(\log q+U\), and the product rule gives


\[\dot q\,\bigl(\log q(t)+U(t)\bigr)+q(t)\Bigl(\frac{\dot q}{q(t)}+\dot U\Bigr)=\dot q\,\bigl(\log q(t)+U(t)\bigr)+\dot q+q(t)\dot U,\]


which is the stated expression. \(\square\)




The term \(+1\) comes from differentiating \(q\log q\). It does not appear in the paper's first variation; it disappears against directions of zero total mass, and in the proofs below it is absorbed into the test function \(\log q+U+1\).

\begin{leanbox}{Lean: \texttt{energy\_\allowbreak{}integrand\_\allowbreak{}hasDerivAt}}
\leansource{\texttt{energy\_\allowbreak{}integrand\_\allowbreak{}hasDerivAt} \quad LatentDiffusion/\allowbreak{}RegularFlowCalculus.lean, lines 31–40}
\begin{leancode}
\leanline{0}{theorem\ energy\_integrand\_hasDerivAt}
\leanline{4}{\{q\ U\ :\ ℝ\ →\ ℝ\}\ \{qd\ Ud\ t\ :\ ℝ\}}
\leanline{4}{(hq\ :\ HasDerivAt\ q\ qd\ t)\ (hU\ :\ HasDerivAt\ U\ Ud\ t)\ (hpos\ :\ 0\ <\ q\ t)\ :}
\leanline{4}{HasDerivAt\ (fun\ s\ ↦\ q\ s\ *\ (Real.log\ (q\ s)\ +\ U\ s))}
\leanline{6}{(qd\ *\ (Real.log\ (q\ t)\ +\ U\ t\ +\ 1)\ +\ q\ t\ *\ Ud)\ t}
\end{leancode}
\leanlabel{How the Lean proof follows the argument}
\texttt{energy\_\allowbreak{}integrand\_\allowbreak{}has\allowbreak{}Deriv\allowbreak{}At} takes \texttt{hq : Has\allowbreak{}Deriv\allowbreak{}At q qd t} and \texttt{h\allowbreak{}U : Has\allowbreak{}Deriv\allowbreak{}At U Ud t}; \texttt{Has\allowbreak{}Deriv\allowbreak{}At f f' t} means that \(f\) is differentiable at \(t\) with derivative \(f'\). The proof applies the product rule \texttt{hq.\allowbreak{}mul} to the chain rule \texttt{hq.\allowbreak{}log hpos.\allowbreak{}ne'} added to \texttt{h\allowbreak{}U}. The identity \texttt{hc}, \(q(t)\cdot(\dot q/q(t))=\dot q\), uses \(q(t)\ne0\); after it, \texttt{ring} rearranges the derivative into the stated form.
\end{leanbox}
\subsection{\texorpdfstring{Differentiation of the density energy under the integral sign}{Differentiation of the density energy under the integral sign}}


\textbf{Setting.} Let \(\nu\) be a measure on a measurable space \(X\), and let \(q,\dot q,U,\dot U:\mathbb R\times X\to\mathbb R\), written \(q_u(x)\) and so on. The density energy at time \(u\) is


\[\mathcal E(u)=\int_Xq_u(x)\bigl(\log q_u(x)+U_u(x)\bigr)\,\nu(dx),\]


and the candidate derivative integrand is \(\varphi_u(x)=\dot q_u(x)\bigl(\log q_u(x)+U_u(x)+1\bigr)+q_u(x)\dot U_u(x)\).






Local domination at \(t\) consists of a neighbourhood \(N\) of \(t\) and a function \(g:X\to\mathbb R\) such that



\begin{enumerate}


\item for all \(u\) near \(t\), the integrand \(x\mapsto q_u(x)(\log q_u(x)+U_u(x))\) is measurable, and it is \(\nu\)-integrable for \(u=t\);

\item \(\varphi_t\) is measurable;

\item for \(\nu\)-almost every \(x\), \(|\varphi_u(x)|\le g(x)\) for all \(u\in N\), and \(g\) is \(\nu\)-integrable;

\item for \(\nu\)-almost every \(x\) and all \(u\in N\), \(q_u(x)>0\) and \(u\mapsto q_u(x)\) has derivative \(\dot q_u(x)\) at \(u\).

\end{enumerate}





The local differentiation assumptions additionally require: (5) for \(\nu\)-almost every \(x\) and all \(u\in N\), \(u\mapsto U_u(x)\) has derivative \(\dot U_u(x)\) at \(u\). Measurability is understood up to \(\nu\)-null sets.




\textbf{Claim.} Under the local differentiation assumptions at \(t\), \(\varphi_t\) is \(\nu\)-integrable and \(\mathcal E\) is differentiable at \(t\) with


\[\mathcal E'(t)=\int_X\varphi_t\,d\nu.\]


\textbf{Proof.} We use the theorem on differentiation under the integral sign: let \(F(u,x)\) be measurable in \(x\) for \(u\) near \(t\), with \(F(t,\cdot)\) integrable; suppose that for almost every \(x\) the function \(u\mapsto F(u,x)\) has derivative \(F'(u,x)\) for every \(u\) in a neighbourhood \(N\) of \(t\), with \(|F'(u,x)|\le g(x)\) for \(u\in N\) and \(g\) integrable, and that \(F'(t,\cdot)\) is measurable. Then \(F'(t,\cdot)\) is integrable and \(u\mapsto\int F(u,x)\,dx\) has derivative \(\int F'(t,x)\,dx\) at \(t\). (The proof: by the mean value inequality, \(|F(u,x)-F(t,x)|\le g(x)|u-t|\) for \(u\) in a ball inside \(N\), so the difference quotients are dominated by \(g\); dominated convergence passes to the limit.)




Apply it with \(F(u,x)=q_u(x)(\log q_u(x)+U_u(x))\) and \(F'=\varphi\). Hypotheses (1)–(3) are its measurability, integrability and domination conditions. For almost every \(x\) and every \(u\in N\), hypotheses (4) and (5) give \(q_u(x)>0\) and the derivatives of \(u\mapsto q_u(x)\) and \(u\mapsto U_u(x)\); the previous step then shows that \(u\mapsto F(u,x)\) has derivative \(\varphi_u(x)\). \(\square\)




None of the hypotheses is a statement about the derivative of the integral itself. Domination is needed because \(E\) is unbounded and the densities need not be compactly supported.

\begin{leanbox}{Lean: \texttt{LocalEnergyDomination}, \texttt{LocalEnergyRegularity}, \texttt{LocalEnergyRegularity.\allowbreak{}hasDerivAt\_\allowbreak{}energy}}
\leansource{\texttt{LocalEnergyDomination} \quad LatentDiffusion/\allowbreak{}RegularFlowCalculus.lean, lines 45–60}
\begin{leancode}
\leanline{0}{structure\ LocalEnergyDomination\ \{X\ :\ Type*\}\ [MeasurableSpace\ X]}
\leanline{4}{(ν\ :\ Measure\ X)\ (q\ qdot\ U\ Udot\ :\ ℝ\ →\ X\ →\ ℝ)\ (t\ :\ ℝ)\ where}
\leanline{2}{neighborhood\ :\ Set\ ℝ}
\leanline{2}{neighborhood\_mem\ :\ neighborhood\ ∈\ {\leanfallback 𝓝}\ t}
\leanline{2}{bound\ :\ X\ →\ ℝ}
\leanline{2}{measurable\ :\ ∀ᶠ\ u\ in\ {\leanfallback 𝓝}\ t,}
\leanline{4}{AEStronglyMeasurable\ (fun\ x\ ↦\ q\ u\ x\ *\ (Real.log\ (q\ u\ x)\ +\ U\ u\ x))\ ν}
\leanline{2}{integrable\ :\ Integrable\ (fun\ x\ ↦\ q\ t\ x\ *\ (Real.log\ (q\ t\ x)\ +\ U\ t\ x))\ ν}
\leanline{2}{derivative\_measurable\ :\ AEStronglyMeasurable}
\leanline{4}{(fun\ x\ ↦\ qdot\ t\ x\ *\ (Real.log\ (q\ t\ x)\ +\ U\ t\ x\ +\ 1)\ +\ q\ t\ x\ *\ Udot\ t\ x)\ ν}
\leanline{2}{derivative\_bound\ :\ ∀ᵐ\ x\ ∂ν,\ ∀\ u\ ∈\ neighborhood,}
\leanline{4}{‖qdot\ u\ x\ *\ (Real.log\ (q\ u\ x)\ +\ U\ u\ x\ +\ 1)\ +\ q\ u\ x\ *\ Udot\ u\ x‖\ ≤\ bound\ x}
\leanline{2}{bound\_integrable\ :\ Integrable\ bound\ ν}
\leanline{2}{positive\ :\ ∀ᵐ\ x\ ∂ν,\ ∀\ u\ ∈\ neighborhood,\ 0\ <\ q\ u\ x}
\leanline{2}{density\_derivative\ :\ ∀ᵐ\ x\ ∂ν,\ ∀\ u\ ∈\ neighborhood,}
\leanline{4}{HasDerivAt\ (fun\ v\ ↦\ q\ v\ x)\ (qdot\ u\ x)\ u}
\end{leancode}
\leansource{\texttt{LocalEnergyRegularity} \quad LatentDiffusion/\allowbreak{}RegularFlowCalculus.lean, lines 64–68}
\begin{leancode}
\leanline{0}{structure\ LocalEnergyRegularity\ \{X\ :\ Type*\}\ [MeasurableSpace\ X]}
\leanline{4}{(ν\ :\ Measure\ X)\ (q\ qdot\ U\ Udot\ :\ ℝ\ →\ X\ →\ ℝ)\ (t\ :\ ℝ)}
\leanline{4}{extends\ LocalEnergyDomination\ ν\ q\ qdot\ U\ Udot\ t\ where}
\leanline{2}{potential\_derivative\ :\ ∀ᵐ\ x\ ∂ν,\ ∀\ u\ ∈\ neighborhood,}
\leanline{4}{HasDerivAt\ (fun\ v\ ↦\ U\ v\ x)\ (Udot\ u\ x)\ u}
\end{leancode}
\leansource{\texttt{LocalEnergyRegularity.\allowbreak{}hasDerivAt\_\allowbreak{}energy} \quad LatentDiffusion/\allowbreak{}RegularFlowCalculus.lean, lines 71–84}
\begin{leancode}
\leanline{0}{theorem\ LocalEnergyRegularity.hasDerivAt\_energy}
\leanline{4}{\{X\ :\ Type*\}\ [MeasurableSpace\ X]\ \{ν\ :\ Measure\ X\}}
\leanline{4}{\{q\ qdot\ U\ Udot\ :\ ℝ\ →\ X\ →\ ℝ\}\ \{t\ :\ ℝ\}}
\leanline{4}{(data\ :\ LocalEnergyRegularity\ ν\ q\ qdot\ U\ Udot\ t)\ :}
\leanline{4}{Integrable}
\leanline{6}{(fun\ x\ ↦\ qdot\ t\ x\ *\ (Real.log\ (q\ t\ x)\ +\ U\ t\ x\ +\ 1)\ +\ q\ t\ x\ *\ Udot\ t\ x)\ ν\ ∧}
\leanline{4}{HasDerivAt\ (fun\ u\ ↦\ SourceCalculus.densityFreeEnergy\ ν\ (q\ u)\ (U\ u))}
\leanline{6}{(∫\ x,\ qdot\ t\ x\ *\ (Real.log\ (q\ t\ x)\ +\ U\ t\ x\ +\ 1)\ +\ q\ t\ x\ *\ Udot\ t\ x\ ∂ν)\ t}
\end{leancode}
\leanlabel{How the Lean proof follows the argument}
\texttt{Local\allowbreak{}Energy\allowbreak{}Domination \ensuremath{\nu} q qdot U Udot t} is a structure whose fields are the neighbourhood \texttt{neighborhood} with \texttt{neighborhood\_\allowbreak{}mem : neighborhood \ensuremath{\in} {\leanfallback 𝓝} t}, the bound \texttt{bound}, and conditions (1)–(4): \texttt{measurable} (stated with \texttt{\ensuremath{\forall}ᶠ u in {\leanfallback 𝓝} t}, “for all \(u\) near \(t\)”), \texttt{integrable}, \texttt{derivative\_\allowbreak{}measurable}, \texttt{derivative\_\allowbreak{}bound} and \texttt{bound\_\allowbreak{}integrable}, \texttt{positive}, and \texttt{density\_\allowbreak{}derivative}. The quantifier \texttt{\ensuremath{\forall}ᵐ x \ensuremath{\partial}\ensuremath{\nu}} means “for \(\nu\)-almost every \(x\)”, and \texttt{AEStrongly\allowbreak{}Measurable} is measurability up to a null set. \texttt{Local\allowbreak{}Energy\allowbreak{}Regularity} extends it by the field \texttt{potential\_\allowbreak{}derivative}, condition (5).




\texttt{Local\allowbreak{}Energy\allowbreak{}Regularity.\allowbreak{}has\allowbreak{}Deriv\allowbreak{}At\_\allowbreak{}energy} states integrability of \(\varphi_t\) and \texttt{Has\allowbreak{}Deriv\allowbreak{}At} of \texttt{fun u \ensuremath{\mapsto} Source\allowbreak{}Calculus.\allowbreak{}density\allowbreak{}Free\allowbreak{}Energy \ensuremath{\nu} (q u) (U u)}, where \texttt{Source\allowbreak{}Calculus.\allowbreak{}density\allowbreak{}Free\allowbreak{}Energy \ensuremath{\nu} q U} is the integral \(\int q(\log q+U)\,d\nu\). The proof applies Mathlib's \texttt{has\allowbreak{}Deriv\allowbreak{}At\_\allowbreak{}integral\_\allowbreak{}of\_\allowbreak{}dominated\_\allowbreak{}loc\_\allowbreak{}of\_\allowbreak{}deriv\_\allowbreak{}le}, the theorem quoted in the proof, whose conclusion is exactly this pair. Its hypotheses are supplied by the fields of \texttt{data}; the pointwise derivative is obtained by \texttt{filter\_\allowbreak{}upwards} on the three almost-everywhere fields and \texttt{energy\_\allowbreak{}integrand\_\allowbreak{}has\allowbreak{}Deriv\allowbreak{}At}.
\end{leanbox}
\subsection{\texorpdfstring{Divergence in coordinates and integration by parts}{Divergence in coordinates and integration by parts}}


\textbf{Definitions.} Let \(e_1,\dots,e_d\) be a fixed orthonormal basis of \(E\). For a vector field \(V:E\to E\) the components are \(V_i(x)=\langle V(x),e_i\rangle\), and the divergence is


\[\nabla\cdot V(x)=\sum_{i=1}^d\partial_iV_i(x),\qquad\partial_if(x)=Df(x)\,e_i,\]


where \(Df(x)\) is the Fréchet derivative. If \(V\) has Fréchet derivative \(A\) at \(x\), then \(V_i\) has derivative \(h\mapsto\langle Ah,e_i\rangle\), so \(\nabla\cdot V(x)=\sum_i\langle Ae_i,e_i\rangle=\operatorname{tr}A\). Components and divergence are linear in \(V\); in particular \((-V)_i=-V_i\) and \(\nabla\cdot(-V)=-\nabla\cdot V\).




For \(V:E\to E\) and \(\psi:E\to\mathbb R\), the \emph{no-boundary-flux} condition requires that each component \(V_i\) is differentiable and that, for each \(i\), the three functions


\[(\partial_iV_i)\,\psi,\qquad V_i\,\partial_i\psi,\qquad V_i\,\psi\]


are Lebesgue integrable on \(E\). These are the hypotheses of Lebesgue integration by parts on \(E\); they replace the boundary terms that vanish on a bounded domain. The condition is unchanged when \(V\) is replaced by \(-V\).




\textbf{Claim.} Let \(\psi\) be differentiable and let \((V,\psi)\) satisfy the no-boundary-flux condition. Then \(x\mapsto\langle V(x),\nabla\psi(x)\rangle\) and \(x\mapsto\nabla\cdot V(x)\,\psi(x)\) are integrable, and


\[\int_E\nabla\cdot V\,\psi\,dx=-\int_E\langle V,\nabla\psi\rangle\,dx.\qquad(1)\]


\textbf{Proof.} Expanding in the orthonormal basis, \(\langle u,w\rangle=\sum_i\langle u,e_i\rangle\langle w,e_i\rangle\); with \(\langle\nabla\psi(x),e_i\rangle=D\psi(x)e_i=\partial_i\psi(x)\) this gives


\[\langle V(x),\nabla\psi(x)\rangle=\sum_iV_i(x)\,\partial_i\psi(x),\qquad\nabla\cdot V(x)\,\psi(x)=\sum_i\partial_iV_i(x)\,\psi(x).\]


Both are finite sums of integrable functions, hence integrable, and the integrals split into sums over \(i\). For each \(i\) we use Lebesgue integration by parts in the direction \(e_i\): if \(f,g:E\to\mathbb R\) are differentiable and \(f\,\partial_ig\), \((\partial_if)\,g\) and \(fg\) are integrable, then \(\int f\,\partial_ig\,dx=-\int(\partial_if)\,g\,dx\). (On each line parallel to \(e_i\) this is the one-dimensional formula; integrability of \(fg\) forces the boundary values to vanish along a sequence of endpoints, and Fubini's theorem integrates over the lines.) With \(f=V_i\) and \(g=\psi\), the three integrability conditions are those of the no-boundary-flux condition, so \(\int V_i\,\partial_i\psi=-\int\partial_iV_i\,\psi\). Summing over \(i\) gives (1). \(\square\)

\begin{leanbox}{Lean: \texttt{fluxComponent}, \texttt{fluxDivergence}, \texttt{fluxComponent\_\allowbreak{}neg} and 7 more}
\leansource{\texttt{fluxComponent} \quad LatentDiffusion/\allowbreak{}RegularFlowCalculus.lean, lines 157–158}
\begin{leancode}
\leanline{0}{def\ fluxComponent\ (V\ :\ E\ →\ E)\ (i\ :\ Fin\ (Module.finrank\ ℝ\ E))\ :\ E\ →\ ℝ\ :=}
\leanline{2}{fun\ x\ ↦\ inner\ ℝ\ (V\ x)\ (stdOrthonormalBasis\ ℝ\ E\ i)}
\end{leancode}
\leansource{\texttt{fluxDivergence} \quad LatentDiffusion/\allowbreak{}RegularFlowCalculus.lean, lines 161–163}
\begin{leancode}
\leanline{0}{def\ fluxDivergence\ (V\ :\ E\ →\ E)\ (x\ :\ E)\ :\ ℝ\ :=}
\leanline{2}{∑\ i\ :\ Fin\ (Module.finrank\ ℝ\ E),}
\leanline{4}{fderiv\ ℝ\ (fluxComponent\ V\ i)\ x\ (stdOrthonormalBasis\ ℝ\ E\ i)}
\end{leancode}
\leansource{\texttt{fluxComponent\_\allowbreak{}neg} \quad LatentDiffusion/\allowbreak{}RegularFlowCalculus.lean, lines 166–169}
\begin{leancode}
\leanline{0}{theorem\ fluxComponent\_neg\ (V\ :\ E\ →\ E)\ (i\ :\ Fin\ (Module.finrank\ ℝ\ E))\ :}
\leanline{4}{fluxComponent\ (fun\ y\ ↦\ -V\ y)\ i\ =\ -fluxComponent\ V\ i}
\end{leancode}
\leansource{\texttt{fluxDivergence\_\allowbreak{}neg} \quad LatentDiffusion/\allowbreak{}RegularFlowCalculus.lean, lines 172–175}
\begin{leancode}
\leanline{0}{theorem\ fluxDivergence\_neg\ (V\ :\ E\ →\ E)\ (x\ :\ E)\ :}
\leanline{4}{fluxDivergence\ (fun\ y\ ↦\ -V\ y)\ x\ =\ -fluxDivergence\ V\ x}
\end{leancode}
\leansource{\texttt{fluxDivergence\_\allowbreak{}eq\_\allowbreak{}trace} \quad LatentDiffusion/\allowbreak{}RegularFlowCalculus.lean, lines 180–193}
\begin{leancode}
\leanline{0}{theorem\ fluxDivergence\_eq\_trace\ (V\ :\ E\ →\ E)\ (A\ :\ E\ →L[ℝ]\ E)\ (x\ :\ E)}
\leanline{4}{(hV\ :\ HasFDerivAt\ V\ A\ x)\ :}
\leanline{4}{fluxDivergence\ V\ x\ =\ LinearMap.trace\ ℝ\ E\ A.toLinearMap}
\end{leancode}
\leansource{\texttt{NoBoundaryFlux} \quad LatentDiffusion/\allowbreak{}RegularFlowCalculus.lean, lines 199–207}
\begin{leancode}
\leanline{0}{structure\ NoBoundaryFlux\ (V\ :\ E\ →\ E)\ (ψ\ :\ E\ →\ ℝ)\ :\ Prop\ where}
\leanline{2}{flux\_differentiable\ :\ ∀\ i,\ Differentiable\ ℝ\ (fluxComponent\ V\ i)}
\leanline{2}{derivative\_flux\_integrable\ :\ ∀\ i,\ Integrable}
\leanline{4}{(fun\ x\ ↦\ fderiv\ ℝ\ (fluxComponent\ V\ i)\ x\ (stdOrthonormalBasis\ ℝ\ E\ i)\ *\ ψ\ x)}
\leanline{4}{volume}
\leanline{2}{derivative\_test\_integrable\ :\ ∀\ i,\ Integrable}
\leanline{4}{(fun\ x\ ↦\ fluxComponent\ V\ i\ x\ *\ fderiv\ ℝ\ ψ\ x\ (stdOrthonormalBasis\ ℝ\ E\ i))}
\leanline{4}{volume}
\leanline{2}{product\_integrable\ :\ ∀\ i,\ Integrable\ (fun\ x\ ↦\ fluxComponent\ V\ i\ x\ *\ ψ\ x)\ volume}
\end{leancode}
\leansource{\texttt{NoBoundaryFlux.\allowbreak{}neg} \quad LatentDiffusion/\allowbreak{}RegularFlowCalculus.lean, lines 210–225}
\begin{leancode}
\leanline{0}{theorem\ NoBoundaryFlux.neg\ \{V\ :\ E\ →\ E\}\ \{ψ\ :\ E\ →\ ℝ\}}
\leanline{4}{(htail\ :\ NoBoundaryFlux\ V\ ψ)\ :\ NoBoundaryFlux\ (fun\ y\ ↦\ -V\ y)\ ψ}
\end{leancode}
\leansource{\texttt{NoBoundaryFlux.\allowbreak{}integrable\_\allowbreak{}flux\_\allowbreak{}gradient} \quad LatentDiffusion/\allowbreak{}RegularFlowCalculus.lean, lines 229–240}
\begin{leancode}
\leanline{0}{theorem\ NoBoundaryFlux.integrable\_flux\_gradient}
\leanline{4}{\{V\ :\ E\ →\ E\}\ \{ψ\ :\ E\ →\ ℝ\}\ (htail\ :\ NoBoundaryFlux\ V\ ψ)\ :}
\leanline{4}{Integrable\ (fun\ x\ ↦\ inner\ ℝ\ (V\ x)\ (gradient\ ψ\ x))\ volume}
\end{leancode}
\leansource{\texttt{NoBoundaryFlux.\allowbreak{}integrable\_\allowbreak{}divergence\_\allowbreak{}mul} \quad LatentDiffusion/\allowbreak{}RegularFlowCalculus.lean, lines 243–247}
\begin{leancode}
\leanline{0}{theorem\ NoBoundaryFlux.integrable\_divergence\_mul}
\leanline{4}{\{V\ :\ E\ →\ E\}\ \{ψ\ :\ E\ →\ ℝ\}\ (htail\ :\ NoBoundaryFlux\ V\ ψ)\ :}
\leanline{4}{Integrable\ (fun\ x\ ↦\ fluxDivergence\ V\ x\ *\ ψ\ x)\ volume}
\end{leancode}
\leansource{\texttt{integral\_\allowbreak{}divergence\_\allowbreak{}mul} \quad LatentDiffusion/\allowbreak{}RegularFlowCalculus.lean, lines 251–273}
\begin{leancode}
\leanline{0}{theorem\ integral\_divergence\_mul}
\leanline{4}{(V\ :\ E\ →\ E)\ (ψ\ :\ E\ →\ ℝ)\ (hψ\ :\ Differentiable\ ℝ\ ψ)}
\leanline{4}{(htail\ :\ NoBoundaryFlux\ V\ ψ)\ :}
\leanline{4}{(∫\ x,\ fluxDivergence\ V\ x\ *\ ψ\ x\ ∂volume)\ =}
\leanline{6}{-(∫\ x,\ inner\ ℝ\ (V\ x)\ (gradient\ ψ\ x)\ ∂volume)}
\end{leancode}
\leanlabel{How the Lean proof follows the argument}
The basis is Mathlib's \texttt{std\allowbreak{}Orthonormal\allowbreak{}Basis \ensuremath{\mathbb R} E}, indexed by \texttt{Fin (Module.\allowbreak{}finrank \ensuremath{\mathbb R} E)}. \texttt{flux\allowbreak{}Component V i} is \(V_i\) and \texttt{flux\allowbreak{}Divergence V x} is \(\sum_i\partial_iV_i(x)\), where the directional derivative is \texttt{fderiv \ensuremath{\mathbb R} (flux\allowbreak{}Component V i) x (std\allowbreak{}Orthonormal\allowbreak{}Basis \ensuremath{\mathbb R} E i)}. \texttt{flux\allowbreak{}Component\_\allowbreak{}neg} and \texttt{flux\allowbreak{}Divergence\_\allowbreak{}neg} are the linearity identities, proved by \texttt{inner\_\allowbreak{}neg\_\allowbreak{}left} and \texttt{fderiv\_\allowbreak{}neg}. \texttt{flux\allowbreak{}Divergence\_\allowbreak{}eq\_\allowbreak{}trace} identifies the divergence with \texttt{Linear\allowbreak{}Map.\allowbreak{}trace} of the derivative: it expands the trace in the orthonormal basis (\texttt{Linear\allowbreak{}Map.\allowbreak{}trace\_\allowbreak{}eq\_\allowbreak{}sum\_\allowbreak{}inner}) and computes the derivative of each component by the chain rule for \texttt{inner\allowbreak{}SL}.



\begin{enumerate}


\item \texttt{No\allowbreak{}Boundary\allowbreak{}Flux V \ensuremath{\psi}} is a structure with the four fields \texttt{flux\_\allowbreak{}differentiable}, \texttt{derivative\_\allowbreak{}flux\_\allowbreak{}integrable}, \texttt{derivative\_\allowbreak{}test\_\allowbreak{}integrable} and \texttt{product\_\allowbreak{}integrable}, which are the differentiability of the components and the three integrability conditions. \texttt{No\allowbreak{}Boundary\allowbreak{}Flux.\allowbreak{}neg} transfers each field to \(-V\).

\item \texttt{No\allowbreak{}Boundary\allowbreak{}Flux.\allowbreak{}integrable\_\allowbreak{}flux\_\allowbreak{}gradient} and \texttt{No\allowbreak{}Boundary\allowbreak{}Flux.\allowbreak{}integrable\_\allowbreak{}divergence\_\allowbreak{}mul} are the two integrability statements. The pointwise expansion \texttt{hpoint} uses \texttt{sum\_\allowbreak{}inner\_\allowbreak{}mul\_\allowbreak{}inner} of the orthonormal basis together with \texttt{inner\_\allowbreak{}gradient\_\allowbreak{}left}, which identifies \(\langle\nabla\psi(x),w\rangle\) with \(D\psi(x)w\); \texttt{integrable\_\allowbreak{}finset\allowbreak{}Sum} handles the finite sums.

\item \texttt{integral\_\allowbreak{}divergence\_\allowbreak{}mul} is (1). After the expansion, \texttt{integral\_\allowbreak{}finset\allowbreak{}Sum} splits both integrals into sums over \(i\), and for each \(i\) Mathlib's \texttt{integral\_\allowbreak{}mul\_\allowbreak{}fderiv\_\allowbreak{}eq\_\allowbreak{}neg\_\allowbreak{}fderiv\_\allowbreak{}mul\_\allowbreak{}of\_\allowbreak{}integrable} is the integration-by-parts formula quoted in the proof, applied with the three integrability fields and the differentiability of \(V_i\) and \(\psi\).

\end{enumerate}
\end{leanbox}
\subsection{\texorpdfstring{The score pairing under a density law}{The score pairing under a density law}}


\textbf{Claim.} Let \(q:E\to\mathbb R\) be continuously differentiable with \(q(x)>0\) for every \(x\), let \(V:E\to E\), and let \((V,q)\) satisfy the no-boundary-flux condition. Let \(Q\) be the measure \(Q(dx)=q(x)\,dx\). Then \(\nabla\cdot V\) and \(\langle V,\nabla\log q\rangle\) are \(Q\)-integrable and


\[\int_E\nabla\cdot V\,dQ=-\int_E\langle V,\nabla\log q\rangle\,dQ.\]


\textbf{Proof.} For a measurable function \(f\), integrability of \(f\) with respect to \(Q\) is equivalent to Lebesgue integrability of \(qf\), and then \(\int f\,dQ=\int qf\,dx\); this is the definition of a measure with a density, valid because \(q\) is measurable (continuous) and nonnegative. Since \(q>0\) is differentiable, the chain rule gives \(\nabla\log q=q^{-1}\nabla q\), hence


\[q\,\langle V,\nabla\log q\rangle=\langle V,\nabla q\rangle.\]


By the previous step with \(\psi=q\), the functions \(\nabla\cdot V\,q\) and \(\langle V,\nabla q\rangle\) are Lebesgue integrable, and \(\int\nabla\cdot V\,q\,dx=-\int\langle V,\nabla q\rangle\,dx\). Translating back to \(Q\) with the identity above proves both integrability statements and the formula. \(\square\)




When \(V(x)=T(x)-x\) is a transport displacement and \(Q\) is a probability law, this is the identity that relates the divergence of the displacement to the score \(\nabla\log q\).

\begin{leanbox}{Lean: \texttt{transport\_\allowbreak{}score\_\allowbreak{}pairing}, \texttt{transport\_\allowbreak{}score\_\allowbreak{}integrable}}
\leansource{\texttt{transport\_\allowbreak{}score\_\allowbreak{}pairing} \quad LatentDiffusion/\allowbreak{}RegularFlowCalculus.lean, lines 278–298}
\begin{leancode}
\leanline{0}{theorem\ transport\_score\_pairing}
\leanline{4}{(q\ :\ E\ →\ ℝ)\ (V\ :\ E\ →\ E)\ (hq\ :\ ContDiff\ ℝ\ 1\ q)\ (hpos\ :\ ∀\ x,\ 0\ <\ q\ x)}
\leanline{4}{(htail\ :\ NoBoundaryFlux\ V\ q)\ :}
\leanline{4}{(∫\ x,\ fluxDivergence\ V\ x\ ∂(volume.withDensity\ (fun\ x\ ↦\ ENNReal.ofReal\ (q\ x))))\ =}
\leanline{6}{-(∫\ x,\ inner\ ℝ\ (V\ x)\ (gradient\ (fun\ y\ ↦\ Real.log\ (q\ y))\ x)}
\leanline{8}{∂(volume.withDensity\ (fun\ x\ ↦\ ENNReal.ofReal\ (q\ x))))}
\end{leancode}
\leansource{\texttt{transport\_\allowbreak{}score\_\allowbreak{}integrable} \quad LatentDiffusion/\allowbreak{}RegularFlowCalculus.lean, lines 302–324}
\begin{leancode}
\leanline{0}{theorem\ transport\_score\_integrable}
\leanline{4}{(q\ :\ E\ →\ ℝ)\ (V\ :\ E\ →\ E)\ (hq\ :\ ContDiff\ ℝ\ 1\ q)\ (hpos\ :\ ∀\ x,\ 0\ <\ q\ x)}
\leanline{4}{(htail\ :\ NoBoundaryFlux\ V\ q)\ :}
\leanline{4}{Integrable\ (fluxDivergence\ V)\ (volume.withDensity\ (fun\ x\ ↦\ ENNReal.ofReal\ (q\ x)))\ ∧}
\leanline{4}{Integrable\ (fun\ x\ ↦\ inner\ ℝ\ (V\ x)\ (gradient\ (fun\ y\ ↦\ Real.log\ (q\ y))\ x))}
\leanline{6}{(volume.withDensity\ (fun\ x\ ↦\ ENNReal.ofReal\ (q\ x)))}
\end{leancode}
\leanlabel{How the Lean proof follows the argument}
The measure \(Q\) is \texttt{volume.\allowbreak{}with\allowbreak{}Density (fun x \ensuremath{\mapsto} ENNReal.\allowbreak{}of\allowbreak{}Real (q x))}: Lebesgue measure with density \(q\), converted to an extended nonnegative number by \texttt{ENNReal.\allowbreak{}of\allowbreak{}Real}.



\begin{enumerate}


\item \texttt{transport\_\allowbreak{}score\_\allowbreak{}pairing}: \texttt{hconvert} is \(\int f\,dQ=\int qf\,dx\), from \texttt{integral\_\allowbreak{}with\allowbreak{}Density\_\allowbreak{}eq\_\allowbreak{}integral\_\allowbreak{}to\allowbreak{}Real\_\allowbreak{}smul} and \texttt{ENNReal.\allowbreak{}to\allowbreak{}Real\_\allowbreak{}of\allowbreak{}Real} (valid since \(q\ge0\)). \texttt{hpoint} is the identity \(q\langle V,\nabla\log q\rangle=\langle V,\nabla q\rangle\); inside it, \texttt{hscore} computes \(\nabla\log q=q^{-1}\nabla q\) with \texttt{fderiv.\allowbreak{}log}. The formula is then \texttt{integral\_\allowbreak{}divergence\_\allowbreak{}mul V q}.

\item \texttt{transport\_\allowbreak{}score\_\allowbreak{}integrable}: the integrability analogue \texttt{integrable\_\allowbreak{}with\allowbreak{}Density\_\allowbreak{}iff\_\allowbreak{}integrable\_\allowbreak{}smul'} reduces both claims to Lebesgue integrability of \(q\,\nabla\cdot V\) and \(\langle V,\nabla q\rangle\), which are \texttt{htail.\allowbreak{}integrable\_\allowbreak{}divergence\_\allowbreak{}mul} and \texttt{htail.\allowbreak{}integrable\_\allowbreak{}flux\_\allowbreak{}gradient}.

\end{enumerate}
\end{leanbox}
\subsection{\texorpdfstring{First variation along a continuity equation}{First variation along a continuity equation}}


The \emph{chemical potential} of a density \(q>0\) and a potential \(U\) is \(\mu_{q,U}=\log q+U\). Adding a constant does not change a gradient: \(\nabla(\mu+1)=\nabla\mu\).




\textbf{Claim 1.} Let \(q,\dot q,U:E\to\mathbb R\) and \(v:E\to E\). Assume \(\mu=\log q+U\) is differentiable, the continuity equation \(\dot q=-\nabla\cdot(qv)\) holds almost everywhere, and \((qv,\mu+1)\) satisfies the no-boundary-flux condition. Then


\[\int_E\dot q\,(\mu+1)\,dx=\int_Eq\,\langle v,\nabla\mu\rangle\,dx.\]


\textbf{Proof.} Integration by parts (1) with \(V=qv\) and \(\psi=\mu+1\), together with \(\nabla(\mu+1)=\nabla\mu\), gives \(\int\nabla\cdot(qv)(\mu+1)=-\int q\langle v,\nabla\mu\rangle\). Since \(\dot q=-\nabla\cdot(qv)\) almost everywhere, the integrals of \(\dot q(\mu+1)\) and \(-\nabla\cdot(qv)(\mu+1)\) agree. \(\square\)




\textbf{Claim 2.} Let \(q,\dot q,U,\dot U:\mathbb R\times E\to\mathbb R\) satisfy the local differentiation assumptions at \(t\) for Lebesgue measure (second step). Assume moreover that \(q_t(x)>0\) for every \(x\), that \(q_t\) and \(U_t\) are continuously differentiable, that \(q_t\dot U_t\) is integrable, that \(\dot q_t=-\nabla\cdot(q_tv)\) almost everywhere for a vector field \(v\), and that \((q_tv,\mu_t+1)\) satisfies the no-boundary-flux condition, where \(\mu_t=\log q_t+U_t\). Then


\[\mathcal E'(t)=\int_Eq_t\,\langle v,\nabla\mu_t\rangle\,dx+\int_Eq_t\,\dot U_t\,dx.\]


\textbf{Proof.} Since \(q_t>0\) and \(q_t,U_t\) are continuously differentiable, \(\mu_t\) is continuously differentiable. By the second step, \(\varphi_t=\dot q_t(\mu_t+1)+q_t\dot U_t\) is integrable and \(\mathcal E'(t)=\int\varphi_t\). As \(q_t\dot U_t\) is integrable, so is \(\dot q_t(\mu_t+1)=\varphi_t-q_t\dot U_t\), and the integral of \(\varphi_t\) splits as \(\int\dot q_t(\mu_t+1)+\int q_t\dot U_t\). Claim 1 evaluates the first integral. \(\square\)

\begin{leanbox}{Lean: \texttt{chemicalPotential}, \texttt{gradient\_\allowbreak{}add\_\allowbreak{}one}, \texttt{continuity\_\allowbreak{}density\_\allowbreak{}variation} and 1 more}
\leansource{\texttt{chemicalPotential} \quad LatentDiffusion/\allowbreak{}RegularFlowCalculus.lean, lines 91–91}
\begin{leancode}
\leanline{0}{def\ chemicalPotential\ (q\ U\ :\ E\ →\ ℝ)\ :\ E\ →\ ℝ\ :=\ fun\ x\ ↦\ Real.log\ (q\ x)\ +\ U\ x}
\end{leancode}
\leansource{\texttt{gradient\_\allowbreak{}add\_\allowbreak{}one} \quad LatentDiffusion/\allowbreak{}RegularFlowCalculus.lean, lines 106–108}
\begin{leancode}
\leanline{0}{theorem\ gradient\_add\_one\ (f\ :\ E\ →\ ℝ)\ (x\ :\ E)\ :}
\leanline{4}{gradient\ (fun\ y\ ↦\ f\ y\ +\ 1)\ x\ =\ gradient\ f\ x}
\end{leancode}
\leansource{\texttt{continuity\_\allowbreak{}density\_\allowbreak{}variation} \quad LatentDiffusion/\allowbreak{}RegularFlowCalculus.lean, lines 352–373}
\begin{leancode}
\leanline{0}{theorem\ continuity\_density\_variation}
\leanline{4}{(q\ qdot\ U\ :\ E\ →\ ℝ)\ (v\ :\ E\ →\ E)}
\leanline{4}{(hW\ :\ Differentiable\ ℝ\ (chemicalPotential\ q\ U))}
\leanline{4}{(hcontinuity\ :\ ∀ᵐ\ x\ ∂volume,}
\leanline{6}{qdot\ x\ =\ -fluxDivergence\ (fun\ y\ ↦\ q\ y\ •\ v\ y)\ x)}
\leanline{4}{(htail\ :\ NoBoundaryFlux\ (fun\ y\ ↦\ q\ y\ •\ v\ y)}
\leanline{6}{(fun\ x\ ↦\ chemicalPotential\ q\ U\ x\ +\ 1))\ :}
\leanline{4}{(∫\ x,\ qdot\ x\ *\ (chemicalPotential\ q\ U\ x\ +\ 1)\ ∂volume)\ =}
\leanline{6}{∫\ x,\ q\ x\ *\ inner\ ℝ\ (v\ x)\ (gradient\ (chemicalPotential\ q\ U)\ x)\ ∂volume}
\end{leancode}
\leansource{\texttt{energy\_\allowbreak{}derivative\_\allowbreak{}of\_\allowbreak{}regular\_\allowbreak{}continuity} \quad LatentDiffusion/\allowbreak{}RegularFlowCalculus.lean, lines 378–410}
\begin{leancode}
\leanline{0}{theorem\ energy\_derivative\_of\_regular\_continuity}
\leanline{4}{(q\ qdot\ U\ Udot\ :\ ℝ\ →\ E\ →\ ℝ)\ (t\ :\ ℝ)\ (v\ :\ E\ →\ E)}
\leanline{4}{(data\ :\ LocalEnergyRegularity\ volume\ q\ qdot\ U\ Udot\ t)}
\leanline{4}{(hpos\ :\ ∀\ x,\ 0\ <\ q\ t\ x)\ (hq\ :\ ContDiff\ ℝ\ 1\ (q\ t))}
\leanline{4}{(hU\ :\ ContDiff\ ℝ\ 1\ (U\ t))}
\leanline{4}{(hpotential\ :\ Integrable\ (fun\ x\ ↦\ q\ t\ x\ *\ Udot\ t\ x)\ volume)}
\leanline{4}{(hcontinuity\ :\ ∀ᵐ\ x\ ∂volume,}
\leanline{6}{qdot\ t\ x\ =\ -fluxDivergence\ (fun\ y\ ↦\ q\ t\ y\ •\ v\ y)\ x)}
\leanline{4}{(htail\ :\ NoBoundaryFlux\ (fun\ y\ ↦\ q\ t\ y\ •\ v\ y)}
\leanline{6}{(fun\ x\ ↦\ chemicalPotential\ (q\ t)\ (U\ t)\ x\ +\ 1))\ :}
\leanline{4}{HasDerivAt\ (fun\ u\ ↦\ SourceCalculus.densityFreeEnergy\ volume\ (q\ u)\ (U\ u))}
\leanline{6}{((∫\ x,\ q\ t\ x\ *\ inner\ ℝ\ (v\ x)}
\leanline{8}{(gradient\ (chemicalPotential\ (q\ t)\ (U\ t))\ x)\ ∂volume)\ +}
\leanline{8}{∫\ x,\ q\ t\ x\ *\ Udot\ t\ x\ ∂volume)\ t}
\end{leancode}
\leanlabel{How the Lean proof follows the argument}
\begin{enumerate}


\item \texttt{chemical\allowbreak{}Potential q U} is the function \texttt{fun x \ensuremath{\mapsto} Real.\allowbreak{}log (q x) + U x}, and \texttt{gradient\_\allowbreak{}add\_\allowbreak{}one} is \(\nabla(f+1)=\nabla f\), from \texttt{fderiv\_\allowbreak{}add\_\allowbreak{}const}.

\item \texttt{continuity\_\allowbreak{}density\_\allowbreak{}variation}: \texttt{hcontinuity} is the almost-everywhere continuity equation and \texttt{htail} the no-boundary-flux condition. \texttt{hip} is \texttt{integral\_\allowbreak{}divergence\_\allowbreak{}mul} for \texttt{fun y \ensuremath{\mapsto} q y \ensuremath{\bullet} v y}, simplified by \texttt{gradient\_\allowbreak{}add\_\allowbreak{}one} and \texttt{inner\_\allowbreak{}smul\_\allowbreak{}left}; the \texttt{calc} block replaces \(\dot q\) almost everywhere (\texttt{integral\_\allowbreak{}congr\_\allowbreak{}ae}) and applies \texttt{hip}.

\item \texttt{energy\_\allowbreak{}derivative\_\allowbreak{}of\_\allowbreak{}regular\_\allowbreak{}continuity}: \texttt{h\allowbreak{}W} is continuous differentiability of \(\mu_t\) (\texttt{hq.\allowbreak{}log} plus \texttt{h\allowbreak{}U}). \texttt{hd} is the second step, \texttt{hint} derives integrability of \(\dot q_t(\mu_t+1)\) as \texttt{hd.\allowbreak{}1.\allowbreak{}sub hpotential}, and \texttt{heq} splits the integral by \texttt{integral\_\allowbreak{}add} and inserts \texttt{continuity\_\allowbreak{}density\_\allowbreak{}variation}.

\end{enumerate}
\end{leanbox}
\subsection{\texorpdfstring{The parameter contribution}{The parameter contribution}}


For a density \(q\) on a measure space \((X,\nu)\), a potential \(U:A\times X\to\mathbb R\) and \(a\in A\), the \emph{parameter force} is the Bochner integral


\[P(a)=\int_Xq(x)\,\nabla_aU(a,x)\,\nu(dx)\in A,\]


where \(\nabla_aU(a,x)\in A\) is the gradient of \(b\mapsto U(b,x)\) at \(a\), defined by \(\langle\nabla_aU(a,x),w\rangle=D_aU(a,x)[w]\).




\textbf{Claim 1.} If \(x\mapsto q(x)\nabla_aU(a,x)\) is \(\nu\)-integrable, then for every \(w\in A\) the function \(x\mapsto q(x)\,D_aU(a,x)[w]\) is integrable and


\[\int_Xq(x)\,D_aU(a,x)[w]\,\nu(dx)=\langle w,P(a)\rangle.\]


\textbf{Proof.} By definition of the gradient, \(q(x)D_aU(a,x)[w]=\langle w,q(x)\nabla_aU(a,x)\rangle\). The map \(y\mapsto\langle w,y\rangle\) is a continuous linear functional on \(A\). Composing an integrable function with a continuous linear functional gives an integrable function, and the functional commutes with the Bochner integral. \(\square\)




\textbf{Claim 2 (chain rule along a parameter path).} Let \(p:\mathbb R\to A\) and put \(U_u(x)=U(p(u),x)\) and \(\dot U_u(x)=D_aU(p(u),x)[\dot p(u)]\). Suppose local domination at \(t\) holds for \((q,\dot q,U_\cdot,\dot U_\cdot)\) with neighbourhood \(N\), that \(p\) has derivative \(\dot p(u)\) at every \(u\in N\), and that \(b\mapsto U(b,x)\) is differentiable at \(p(u)\) for every \(x\) and every \(u\in N\). Then all local differentiation assumptions hold at \(t\).




\textbf{Proof.} Only condition (5) is new. For every \(x\) and \(u\in N\), the chain rule for the composition of \(u\mapsto p(u)\) with \(b\mapsto U(b,x)\) gives the derivative \(D_aU(p(u),x)[\dot p(u)]=\dot U_u(x)\). \(\square\)

\begin{leanbox}{Lean: \texttt{parameterForce}, \texttt{parameter\_\allowbreak{}rate\_\allowbreak{}integral}, \texttt{LocalEnergyDomination.\allowbreak{}along\_\allowbreak{}parameter\_\allowbreak{}path}}
\leansource{\texttt{parameterForce} \quad LatentDiffusion/\allowbreak{}RegularFlowCalculus.lean, lines 460–461}
\begin{leancode}
\leanline{0}{def\ parameterForce\ (ν\ :\ Measure\ X)\ (q\ :\ X\ →\ ℝ)\ (U\ :\ A\ →\ X\ →\ ℝ)\ (a\ :\ A)\ :\ A\ :=}
\leanline{2}{∫\ x,\ q\ x\ •\ gradient\ (fun\ b\ ↦\ U\ b\ x)\ a\ ∂ν}
\end{leancode}
\leansource{\texttt{parameter\_\allowbreak{}rate\_\allowbreak{}integral} \quad LatentDiffusion/\allowbreak{}RegularFlowCalculus.lean, lines 465–475}
\begin{leancode}
\leanline{0}{theorem\ parameter\_rate\_integral}
\leanline{4}{(ν\ :\ Measure\ X)\ (q\ :\ X\ →\ ℝ)\ (U\ :\ A\ →\ X\ →\ ℝ)\ (a\ v\ :\ A)}
\leanline{4}{(hgrad\ :\ Integrable\ (fun\ x\ ↦\ q\ x\ •\ gradient\ (fun\ b\ ↦\ U\ b\ x)\ a)\ ν)\ :}
\leanline{4}{Integrable\ (fun\ x\ ↦\ q\ x\ *\ fderiv\ ℝ\ (fun\ b\ ↦\ U\ b\ x)\ a\ v)\ ν\ ∧}
\leanline{4}{(∫\ x,\ q\ x\ *\ fderiv\ ℝ\ (fun\ b\ ↦\ U\ b\ x)\ a\ v\ ∂ν)\ =}
\leanline{6}{inner\ ℝ\ v\ (parameterForce\ ν\ q\ U\ a)}
\end{leancode}
\leansource{\texttt{LocalEnergyDomination.\allowbreak{}along\_\allowbreak{}parameter\_\allowbreak{}path} \quad LatentDiffusion/\allowbreak{}RegularFlowCalculus.lean, lines 480–491}
\begin{leancode}
\leanline{0}{def\ LocalEnergyDomination.along\_parameter\_path}
\leanline{4}{(ν\ :\ Measure\ X)\ (q\ qdot\ :\ ℝ\ →\ X\ →\ ℝ)\ (U\ :\ A\ →\ X\ →\ ℝ)}
\leanline{4}{(p\ pdot\ :\ ℝ\ →\ A)\ (t\ :\ ℝ)}
\leanline{4}{(data\ :\ LocalEnergyDomination\ ν\ q\ qdot\ (fun\ u\ ↦\ U\ (p\ u))}
\leanline{6}{(fun\ u\ x\ ↦\ fderiv\ ℝ\ (fun\ a\ ↦\ U\ a\ x)\ (p\ u)\ (pdot\ u))\ t)}
\leanline{4}{(hp\ :\ ∀\ u\ ∈\ data.neighborhood,\ HasDerivAt\ p\ (pdot\ u)\ u)}
\leanline{4}{(hU\ :\ ∀\ x\ u,\ u\ ∈\ data.neighborhood\ →\ DifferentiableAt\ ℝ\ (fun\ a\ ↦\ U\ a\ x)\ (p\ u))\ :}
\leanline{4}{LocalEnergyRegularity\ ν\ q\ qdot\ (fun\ u\ ↦\ U\ (p\ u))}
\leanline{6}{(fun\ u\ x\ ↦\ fderiv\ ℝ\ (fun\ a\ ↦\ U\ a\ x)\ (p\ u)\ (pdot\ u))\ t\ where}
\leanline{2}{toLocalEnergyDomination\ :=\ data}
\leanline{2}{potential\_derivative\ :=\ ae\_of\_all\ \_\ fun\ x\ u\ hu\ ↦}
\leanline{4}{(hU\ x\ u\ hu).hasFDerivAt.comp\_hasDerivAt\ u\ (hp\ u\ hu)}
\end{leancode}
\leanlabel{How the Lean proof follows the argument}
\begin{enumerate}


\item \texttt{parameter\allowbreak{}Force \ensuremath{\nu} q U a} is the Bochner integral \texttt{\ensuremath{\int} x, q x \ensuremath{\bullet} gradient (fun b \ensuremath{\mapsto} U b x) a \ensuremath{\partial}\ensuremath{\nu}}; \texttt{gradient} is Mathlib's Riesz representative of the Fréchet derivative.

\item \texttt{parameter\_\allowbreak{}rate\_\allowbreak{}integral}: \texttt{hpoint} rewrites the integrand with \texttt{inner\_\allowbreak{}smul\_\allowbreak{}right} and \texttt{inner\_\allowbreak{}gradient\_\allowbreak{}right}; integrability is \texttt{hgrad.\allowbreak{}const\_\allowbreak{}inner v} and the integral identity is Mathlib's \texttt{integral\_\allowbreak{}inner}, which moves the integral through an inner product with a fixed vector.

\item \texttt{Local\allowbreak{}Energy\allowbreak{}Domination.\allowbreak{}along\_\allowbreak{}parameter\_\allowbreak{}path} is a definition producing a \texttt{Local\allowbreak{}Energy\allowbreak{}Regularity} from \texttt{data}: the field \texttt{to\allowbreak{}Local\allowbreak{}Energy\allowbreak{}Domination} is \texttt{data}, and \texttt{potential\_\allowbreak{}derivative} is the chain rule \texttt{(h\allowbreak{}U x u hu).\allowbreak{}has\allowbreak{}FDeriv\allowbreak{}At.\allowbreak{}comp\_\allowbreak{}has\allowbreak{}Deriv\allowbreak{}At u (hp u hu)}, valid for every \(x\) (so in particular almost everywhere).

\end{enumerate}
\end{leanbox}
\subsection{\texorpdfstring{Representation of the normalized product gradient}{Representation of the normalized product gradient}}


For densities \(q=(q^1,\dots,q^M)\) on \(E\), potentials \(U_m:A\times E\to\mathbb R\) and \(a\in A\), recall \(\mathcal E(a,q)\) and \(G(a,q)\) from the summary; \(G\) is the average of the parameter forces of the data. For a parameter velocity \(\dot a\in A\) and latent velocity fields \(v=(v_1,\dots,v_M)\), define the pairing


\[\Pi(a,q;\dot a,v)=\langle\dot a,G(a,q)\rangle+\frac1M\sum_{m=1}^M\int_Eq^m\,\langle v_m,\nabla\mu_m\rangle\,dz,\qquad\mu_m=\log q^m+U_m(a,\cdot).\]


By the first claim of the previous step, averaged over \(m\), \(M^{-1}\sum_m\int q^mD_aU_m(a,\cdot)[\dot a]\,dz=\langle\dot a,G(a,q)\rangle\) whenever every \(q^m\nabla_aU_m(a,\cdot)\) is integrable.




\textbf{Claim.} Let \(q^m_u\) and \(\dot q^m_u\) (\(m=1,\dots,M\)) be density paths and rates, \(p:\mathbb R\to A\) a path with derivative \(\dot p(u)\) at every \(u\), and \(t\in\mathbb R\). Assume:



\begin{enumerate}


\item for each \(m\), local domination at \(t\) for \(q^m\), \(\dot q^m\), \(U_m(p(u),\cdot)\) and \(D_aU_m(p(u),\cdot)[\dot p(u)]\);

\item \(b\mapsto U_m(b,z)\) is differentiable at every \(b\), for all \(m\) and \(z\);

\item \(q^m_t>0\) everywhere, and \(q^m_t\) and \(U_m(p(t),\cdot)\) are continuously differentiable;

\item \(q^m_t\,\nabla_aU_m(p(t),\cdot)\) is integrable;

\item the continuity equations \(\dot q^m_t=-\nabla\cdot(q^m_tv_m)\) hold almost everywhere, and \((q^m_tv_m,\mu_m+1)\) satisfies the no-boundary-flux condition.

\end{enumerate}



Then \(u\mapsto\mathcal E(p(u),q_u)\) is differentiable at \(t\) with derivative \(\Pi(p(t),q_t;\dot p(t),v)\).




\textbf{Proof.} Fix \(m\). Claim 2 of the previous step proves all local differentiation assumptions from hypothesis (1), using (2) and the differentiability of \(p\). Claim 1 of the previous step and (4) give integrability of \(q^m_t\,D_aU_m(p(t),\cdot)[\dot p(t)]\), which is the term \(q_t\dot U_t\). The second claim of the first-variation step therefore applies and gives the derivative of the \(m\)-th energy:


\[\int_Eq^m_t\langle v_m,\nabla\mu_m\rangle\,dz+\int_Eq^m_t\,D_aU_m(p(t),\cdot)[\dot p(t)]\,dz.\]


A finite sum of differentiable functions is differentiable, with the sum of the derivatives, and so is a constant multiple. Multiplying the sum by \(1/M\) and using the averaged identity for the parameter terms gives \(\Pi\). \(\square\)




With the normalized product metric, whose tangent inner product is \[\langle\dot a,\dot b\rangle+M^{-1}\sum_m\int q^m\langle v_m,w_m\rangle\], the claim says that the gradient of \(\mathcal E\) has parameter component \(G\) and latent components \(\nabla\mu_m\). The factor \(1/M\) of the energy is cancelled by the factor \(M\) of the inverse metric; there is no \(1/M\) in the individual density drift.

\begin{leanbox}{Lean: \texttt{averagedParameterForce}, \texttt{normalizedDensityEnergy}, \texttt{normalizedProductGradientPairing} and 2 more}
\leansource{\texttt{averagedParameterForce} \quad LatentDiffusion/\allowbreak{}RegularFlowCalculus.lean, lines 512–514}
\begin{leancode}
\leanline{0}{def\ averagedParameterForce\ (ν\ :\ Measure\ E)\ (q\ :\ Fin\ M\ →\ E\ →\ ℝ)}
\leanline{4}{(U\ :\ Fin\ M\ →\ A\ →\ E\ →\ ℝ)\ (a\ :\ A)\ :\ A\ :=}
\leanline{2}{(M\ :\ ℝ)⁻¹\ •\ ∑\ i,\ parameterForce\ ν\ (q\ i)\ (U\ i)\ a}
\end{leancode}
\leansource{\texttt{normalizedDensityEnergy} \quad LatentDiffusion/\allowbreak{}RegularFlowCalculus.lean, lines 517–519}
\begin{leancode}
\leanline{0}{def\ normalizedDensityEnergy\ (ν\ :\ Measure\ E)\ (q\ :\ Fin\ M\ →\ E\ →\ ℝ)}
\leanline{4}{(U\ :\ Fin\ M\ →\ A\ →\ E\ →\ ℝ)\ (a\ :\ A)\ :\ ℝ\ :=}
\leanline{2}{(M\ :\ ℝ)⁻¹\ *\ ∑\ i,\ SourceCalculus.densityFreeEnergy\ ν\ (q\ i)\ (U\ i\ a)}
\end{leancode}
\leansource{\texttt{normalizedProductGradientPairing} \quad LatentDiffusion/\allowbreak{}RegularFlowCalculus.lean, lines 524–528}
\begin{leancode}
\leanline{0}{def\ normalizedProductGradientPairing\ (ν\ :\ Measure\ E)\ (q\ :\ Fin\ M\ →\ E\ →\ ℝ)}
\leanline{4}{(U\ :\ Fin\ M\ →\ A\ →\ E\ →\ ℝ)\ (a\ da\ :\ A)\ (v\ :\ Fin\ M\ →\ E\ →\ E)\ :\ ℝ\ :=}
\leanline{2}{inner\ ℝ\ da\ (averagedParameterForce\ ν\ q\ U\ a)\ +}
\leanline{4}{(M\ :\ ℝ)⁻¹\ *\ ∑\ i,\ ∫\ x,\ q\ i\ x\ *\ inner\ ℝ\ (v\ i\ x)}
\leanline{6}{(gradient\ (chemicalPotential\ (q\ i)\ (U\ i\ a))\ x)\ ∂ν}
\end{leancode}
\leansource{\texttt{averaged\_\allowbreak{}parameter\_\allowbreak{}rate} \quad LatentDiffusion/\allowbreak{}RegularFlowCalculus.lean, lines 532–539}
\begin{leancode}
\leanline{0}{theorem\ averaged\_parameter\_rate}
\leanline{4}{(ν\ :\ Measure\ E)\ (q\ :\ Fin\ M\ →\ E\ →\ ℝ)\ (U\ :\ Fin\ M\ →\ A\ →\ E\ →\ ℝ)}
\leanline{4}{(a\ da\ :\ A)}
\leanline{4}{(hgrad\ :\ ∀\ i,\ Integrable\ (fun\ x\ ↦\ q\ i\ x\ •\ gradient\ (fun\ b\ ↦\ U\ i\ b\ x)\ a)\ ν)\ :}
\leanline{4}{(M\ :\ ℝ)⁻¹\ *\ ∑\ i,\ (∫\ x,\ q\ i\ x\ *\ fderiv\ ℝ\ (fun\ b\ ↦\ U\ i\ b\ x)\ a\ da\ ∂ν)\ =}
\leanline{6}{inner\ ℝ\ da\ (averagedParameterForce\ ν\ q\ U\ a)}
\end{leancode}
\leansource{\texttt{normalized\_\allowbreak{}product\_\allowbreak{}gradient\_\allowbreak{}representation} \quad LatentDiffusion/\allowbreak{}RegularFlowCalculus.lean, lines 548–590}
\begin{leancode}
\leanline{0}{theorem\ normalized\_product\_gradient\_representation}
\leanline{4}{(q\ qdot\ :\ Fin\ M\ →\ ℝ\ →\ E\ →\ ℝ)\ (U\ :\ Fin\ M\ →\ A\ →\ E\ →\ ℝ)}
\leanline{4}{(p\ pdot\ :\ ℝ\ →\ A)\ (t\ :\ ℝ)\ (v\ :\ Fin\ M\ →\ E\ →\ E)}
\leanline{4}{(data\ :\ ∀\ i,\ LocalEnergyDomination\ volume\ (q\ i)\ (qdot\ i)\ (fun\ u\ ↦\ U\ i\ (p\ u))}
\leanline{6}{(fun\ u\ x\ ↦\ fderiv\ ℝ\ (fun\ a\ ↦\ U\ i\ a\ x)\ (p\ u)\ (pdot\ u))\ t)}
\leanline{4}{(hp\ :\ ∀\ u,\ HasDerivAt\ p\ (pdot\ u)\ u)}
\leanline{4}{(hparam\ :\ ∀\ i\ x\ a,\ DifferentiableAt\ ℝ\ (fun\ b\ ↦\ U\ i\ b\ x)\ a)}
\leanline{4}{(hpos\ :\ ∀\ i\ x,\ 0\ <\ q\ i\ t\ x)\ (hq\ :\ ∀\ i,\ ContDiff\ ℝ\ 1\ (q\ i\ t))}
\leanline{4}{(hU\ :\ ∀\ i,\ ContDiff\ ℝ\ 1\ (U\ i\ (p\ t)))}
\leanline{4}{(hgrad\ :\ ∀\ i,\ Integrable}
\leanline{6}{(fun\ x\ ↦\ q\ i\ t\ x\ •\ gradient\ (fun\ b\ ↦\ U\ i\ b\ x)\ (p\ t))\ volume)}
\leanline{4}{(hcontinuity\ :\ ∀\ i,\ ∀ᵐ\ x\ ∂volume,}
\leanline{6}{qdot\ i\ t\ x\ =\ -fluxDivergence\ (fun\ y\ ↦\ q\ i\ t\ y\ •\ v\ i\ y)\ x)}
\leanline{4}{(htail\ :\ ∀\ i,\ NoBoundaryFlux\ (fun\ y\ ↦\ q\ i\ t\ y\ •\ v\ i\ y)}
\leanline{6}{(fun\ x\ ↦\ chemicalPotential\ (q\ i\ t)\ (U\ i\ (p\ t))\ x\ +\ 1))\ :}
\leanline{4}{HasDerivAt\ (fun\ u\ ↦\ normalizedDensityEnergy\ volume\ (fun\ i\ ↦\ q\ i\ u)\ U\ (p\ u))}
\leanline{6}{(normalizedProductGradientPairing\ volume\ (fun\ i\ ↦\ q\ i\ t)\ U\ (p\ t)\ (pdot\ t)\ v)\ t}
\end{leancode}
\leanlabel{How the Lean proof follows the argument}
\begin{enumerate}


\item \texttt{averaged\allowbreak{}Parameter\allowbreak{}Force \ensuremath{\nu} q U a} is \texttt{(M : \ensuremath{\mathbb R})\ensuremath{{}^{-1}} \ensuremath{\bullet} ∑ i, parameter\allowbreak{}Force \ensuremath{\nu} (q i) (U i) a}, \texttt{normalized\allowbreak{}Density\allowbreak{}Energy} is \(M^{-1}\sum_i\) of \texttt{Source\allowbreak{}Calculus.\allowbreak{}density\allowbreak{}Free\allowbreak{}Energy}, and \texttt{normalized\allowbreak{}Product\allowbreak{}Gradient\allowbreak{}Pairing} is \(\Pi\).

\item \texttt{averaged\_\allowbreak{}parameter\_\allowbreak{}rate} is the averaged identity; it rewrites each summand with \texttt{parameter\_\allowbreak{}rate\_\allowbreak{}integral} and collects the average with \texttt{inner\_\allowbreak{}smul\_\allowbreak{}right} and \texttt{inner\_\allowbreak{}sum}.

\item \texttt{normalized\_\allowbreak{}product\_\allowbreak{}gradient\_\allowbreak{}representation}: \texttt{heach} is the per-datum derivative; \texttt{d} is \texttt{Local\allowbreak{}Energy\allowbreak{}Domination.\allowbreak{}along\_\allowbreak{}parameter\_\allowbreak{}path} applied to \texttt{data i}, and the derivative comes from \texttt{energy\_\allowbreak{}derivative\_\allowbreak{}of\_\allowbreak{}regular\_\allowbreak{}continuity}, whose integrability input is \texttt{(parameter\_\allowbreak{}rate\_\allowbreak{}integral …).\allowbreak{}1}. \texttt{Has\allowbreak{}Deriv\allowbreak{}At.\allowbreak{}fun\_\allowbreak{}sum} differentiates the finite sum and \texttt{.\allowbreak{}const\_\allowbreak{}mul} the factor \(M^{-1}\). \texttt{hrate} regroups the result into \(\Pi\) with \texttt{Finset.\allowbreak{}sum\_\allowbreak{}add\_\allowbreak{}distrib}, \texttt{mul\_\allowbreak{}add} and \texttt{averaged\_\allowbreak{}parameter\_\allowbreak{}rate}.

\end{enumerate}
\end{leanbox}
\subsection{\texorpdfstring{The de Bruijn identity}{The de Bruijn identity}}


The Fokker–Planck flux of a density \(q\) and a potential \(U\) is \(J_{q,U}=q\,\nabla(\log q+U)\). The Fisher information \(I(a,q)\) is defined in the summary; it is nonnegative whenever every \(q^m\ge0\), since it is a squared norm plus an average of integrals of nonnegative functions.




\textbf{Claim.} Assume hypotheses (1)–(4) of the previous step. Assume further the parameter equation and the Fokker–Planck equations at time \(t\),


\[\dot p(t)=-G(p(t),q_t),\qquad\dot q^m_t=\nabla\cdot J_{q^m_t,U_m(p(t),\cdot)}\ \text{ almost everywhere }(m=1,\dots,M),\]


and the no-boundary-flux condition for \(\bigl(J_{q^m_t,U_m(p(t),\cdot)},\mu_m+1\bigr)\) for every \(m\). Then



\begin{enumerate}


\item every function \(q^m_t\,\|\nabla\mu_m\|^2\) is integrable, and

\item \(u\mapsto\mathcal E(p(u),q_u)\) is differentiable at \(t\) with derivative \(-I(p(t),q_t)\).

\end{enumerate}



\textbf{Proof.} Put \(v_m=-\nabla\mu_m\). Then \(q^m_tv_m=-J^m\), where \(J^m\) is the flux of datum \(m\). By linearity of the divergence, \(-\nabla\cdot(q^m_tv_m)=\nabla\cdot J^m=\dot q^m_t\) almost everywhere, so the continuity equation of the previous step holds. The no-boundary-flux condition is invariant under \(V\mapsto-V\), so it holds for \((q^m_tv_m,\mu_m+1)\). The previous step therefore gives the derivative


\[\Pi=\langle\dot p(t),G\rangle+\frac1M\sum_m\int q^m_t\langle-\nabla\mu_m,\nabla\mu_m\rangle\,dz=-\|G\|^2-\frac1M\sum_m\int q^m_t\|\nabla\mu_m\|^2\,dz=-I,\]


using the parameter equation in the first term. For (1), the integration-by-parts step shows that \(\langle J^m,\nabla(\mu_m+1)\rangle\) is integrable; since \(\nabla(\mu_m+1)=\nabla\mu_m\), this function is \(q^m_t\langle\nabla\mu_m,\nabla\mu_m\rangle=q^m_t\|\nabla\mu_m\|^2\). \(\square\)




The identity is proved from the equations; no entropy-dissipation inequality is assumed. Statement (1) shows that the Fisher information is a finite integral at every time where the hypotheses hold.

\begin{leanbox}{Lean: \texttt{fokkerPlanckFlux}, \texttt{normalizedSlopeSquared}, \texttt{normalizedSlopeSquared\_\allowbreak{}nonneg} and 1 more}
\leansource{\texttt{fokkerPlanckFlux} \quad LatentDiffusion/\allowbreak{}RegularFlowCalculus.lean, lines 328–329}
\begin{leancode}
\leanline{0}{def\ fokkerPlanckFlux\ (q\ U\ :\ E\ →\ ℝ)\ :\ E\ →\ E\ :=}
\leanline{2}{fun\ x\ ↦\ q\ x\ •\ gradient\ (chemicalPotential\ q\ U)\ x}
\end{leancode}
\leansource{\texttt{normalizedSlopeSquared} \quad LatentDiffusion/\allowbreak{}RegularFlowCalculus.lean, lines 594–597}
\begin{leancode}
\leanline{0}{def\ normalizedSlopeSquared\ (ν\ :\ Measure\ E)\ (q\ :\ Fin\ M\ →\ E\ →\ ℝ)}
\leanline{4}{(U\ :\ Fin\ M\ →\ A\ →\ E\ →\ ℝ)\ (a\ :\ A)\ :\ ℝ\ :=}
\leanline{2}{‖averagedParameterForce\ ν\ q\ U\ a‖\ \textasciicircum{}\ 2\ +\ (M\ :\ ℝ)⁻¹\ *\ ∑\ i,}
\leanline{4}{∫\ x,\ q\ i\ x\ *\ ‖gradient\ (chemicalPotential\ (q\ i)\ (U\ i\ a))\ x‖\ \textasciicircum{}\ 2\ ∂ν}
\end{leancode}
\leansource{\texttt{normalizedSlopeSquared\_\allowbreak{}nonneg} \quad LatentDiffusion/\allowbreak{}RegularFlowCalculus.lean, lines 601–605}
\begin{leancode}
\leanline{0}{theorem\ normalizedSlopeSquared\_nonneg\ (ν\ :\ Measure\ E)\ (q\ :\ Fin\ M\ →\ E\ →\ ℝ)}
\leanline{4}{(U\ :\ Fin\ M\ →\ A\ →\ E\ →\ ℝ)\ (a\ :\ A)\ (hq\ :\ ∀\ i\ x,\ 0\ ≤\ q\ i\ x)\ :}
\leanline{4}{0\ ≤\ normalizedSlopeSquared\ ν\ q\ U\ a}
\end{leancode}
\leansource{\texttt{normalized\_\allowbreak{}product\_\allowbreak{}deBruijn} \quad LatentDiffusion/\allowbreak{}RegularFlowCalculus.lean, lines 613–663}
\begin{leancode}
\leanline{0}{theorem\ normalized\_product\_deBruijn}
\leanline{4}{(q\ qdot\ :\ Fin\ M\ →\ ℝ\ →\ E\ →\ ℝ)\ (U\ :\ Fin\ M\ →\ A\ →\ E\ →\ ℝ)}
\leanline{4}{(p\ pdot\ :\ ℝ\ →\ A)\ (t\ :\ ℝ)}
\leanline{4}{(data\ :\ ∀\ i,\ LocalEnergyDomination\ volume\ (q\ i)\ (qdot\ i)\ (fun\ u\ ↦\ U\ i\ (p\ u))}
\leanline{6}{(fun\ u\ x\ ↦\ fderiv\ ℝ\ (fun\ a\ ↦\ U\ i\ a\ x)\ (p\ u)\ (pdot\ u))\ t)}
\leanline{4}{(hp\ :\ ∀\ u,\ HasDerivAt\ p\ (pdot\ u)\ u)}
\leanline{4}{(hparam\ :\ ∀\ i\ x\ a,\ DifferentiableAt\ ℝ\ (fun\ b\ ↦\ U\ i\ b\ x)\ a)}
\leanline{4}{(hpos\ :\ ∀\ i\ x,\ 0\ <\ q\ i\ t\ x)\ (hq\ :\ ∀\ i,\ ContDiff\ ℝ\ 1\ (q\ i\ t))}
\leanline{4}{(hU\ :\ ∀\ i,\ ContDiff\ ℝ\ 1\ (U\ i\ (p\ t)))}
\leanline{4}{(hgrad\ :\ ∀\ i,\ Integrable}
\leanline{6}{(fun\ x\ ↦\ q\ i\ t\ x\ •\ gradient\ (fun\ b\ ↦\ U\ i\ b\ x)\ (p\ t))\ volume)}
\leanline{4}{(hode\ :\ pdot\ t\ =\ -averagedParameterForce\ volume\ (fun\ i\ ↦\ q\ i\ t)\ U\ (p\ t))}
\leanline{4}{(hpde\ :\ ∀\ i,\ ∀ᵐ\ x\ ∂volume,\ qdot\ i\ t\ x\ =}
\leanline{6}{fluxDivergence\ (fokkerPlanckFlux\ (q\ i\ t)\ (U\ i\ (p\ t)))\ x)}
\leanline{4}{(htail\ :\ ∀\ i,\ NoBoundaryFlux\ (fokkerPlanckFlux\ (q\ i\ t)\ (U\ i\ (p\ t)))}
\leanline{6}{(fun\ x\ ↦\ chemicalPotential\ (q\ i\ t)\ (U\ i\ (p\ t))\ x\ +\ 1))\ :}
\leanline{4}{(∀\ i,\ Integrable}
\leanline{6}{(fun\ x\ ↦\ q\ i\ t\ x\ *\ ‖gradient\ (chemicalPotential\ (q\ i\ t)\ (U\ i\ (p\ t)))\ x‖\ \textasciicircum{}\ 2)}
\leanline{6}{volume)\ ∧}
\leanline{4}{HasDerivAt\ (fun\ u\ ↦\ normalizedDensityEnergy\ volume\ (fun\ i\ ↦\ q\ i\ u)\ U\ (p\ u))}
\leanline{6}{(-normalizedSlopeSquared\ volume\ (fun\ i\ ↦\ q\ i\ t)\ U\ (p\ t))\ t}
\end{leancode}
\leanlabel{How the Lean proof follows the argument}
\begin{enumerate}


\item \texttt{fokker\allowbreak{}Planck\allowbreak{}Flux q U} is \texttt{fun x \ensuremath{\mapsto} q x \ensuremath{\bullet} gradient (chemical\allowbreak{}Potential q U) x}. \texttt{normalized\allowbreak{}Slope\allowbreak{}Squared} is \(I\), and \texttt{normalized\allowbreak{}Slope\allowbreak{}Squared\_\allowbreak{}nonneg} proves \(I\ge0\) from \texttt{sq\_\allowbreak{}nonneg} and \texttt{integral\_\allowbreak{}nonneg}.

\item \texttt{normalized\_\allowbreak{}product\_\allowbreak{}de\allowbreak{}Bruijn}: \texttt{hode} is the parameter equation, \texttt{hpde} the Fokker–Planck equations and \texttt{htail} the no-boundary-flux conditions. The velocity \texttt{v} is \(-\nabla\mu_m\); \texttt{hvel} is \(q v_m=-J^m\), \texttt{hcont} derives the continuity equation with \texttt{flux\allowbreak{}Divergence\_\allowbreak{}neg}, and \texttt{htailv} is \texttt{(htail i).\allowbreak{}neg}.

\item \texttt{hd} applies \texttt{normalized\_\allowbreak{}product\_\allowbreak{}gradient\_\allowbreak{}representation}, and \texttt{hrate} evaluates the pairing as \(-I\) by substituting \texttt{hode} and simplifying \(\langle-x,x\rangle=-\|x\|^2\) (\texttt{inner\_\allowbreak{}neg\_\allowbreak{}left}, \texttt{real\_\allowbreak{}inner\_\allowbreak{}self\_\allowbreak{}eq\_\allowbreak{}norm\_\allowbreak{}sq}).

\item Integrability (1) is \texttt{(htail i).\allowbreak{}integrable\_\allowbreak{}flux\_\allowbreak{}gradient}, rewritten with \texttt{gradient\_\allowbreak{}add\_\allowbreak{}one}, \texttt{inner\_\allowbreak{}smul\_\allowbreak{}left} and \texttt{real\_\allowbreak{}inner\_\allowbreak{}self\_\allowbreak{}eq\_\allowbreak{}norm\_\allowbreak{}sq}.

\end{enumerate}
\end{leanbox}
\section{\texorpdfstring{Contractive Euler approximation and accumulated error}{Contractive Euler approximation and accumulated error}}


Work in a complete real inner-product space. The positive force \(F\) is \(m\)-strongly monotone and \(L\)-Lipschitz: \(\langle F(x)-F(y),x-y\rangle\ge m\|x-y\|^2\) and \(\|F(x)-F(y)\|\le L\|x-y\|\). Take \(m,h>0\), \(mh\le1\), and \(hL^2\le m\). The implemented force \(A\) can differ from \(F\). The numerical path is \(x_{k+1}=x_k-hA(x_k)+\xi_k\). A reference segment solves its actual integral equation with the same noise increment; its drift and noise sizes are explicit hypotheses, rather than an assumed numerical-error recurrence.

\subsection{\texorpdfstring{Shared noise cancels in the contractive update}{Shared noise cancels in the contractive update}}


\textbf{Claim.} Under the force and step assumptions above, for any common increment \(\xi\), \(\|(x-hF(x)+\xi)-(y-hF(y)+\xi)\|\le q\|x-y\|\), where \(q=1-mh/2\in[0,1)\).



\textbf{Proof.} Put \(d=x-y\) and \(v=F(x)-F(y)\). The common increment cancels. Expanding the square and applying monotonicity and Lipschitz continuity gives

\[\|d-hv\|^2=\|d\|^2-2h\langle v,d\rangle+h^2\|v\|^2\le(1-2mh+h^2L^2)\|d\|^2\le(1-mh)\|d\|^2.\]

The last inequality uses \(hL^2\le m\), multiplied by \(h\ge0\). Since \(1-mh\le(1-mh/2)^2\) and \(q\ge0\), taking nonnegative square roots proves the claim. Positivity of \(mh\) gives \(q<1\). \(\square\)

\begin{leanbox}{Lean: \texttt{euler\_\allowbreak{}contraction}}
\leansource{\texttt{euler\_\allowbreak{}contraction} \quad LatentDiffusion/\allowbreak{}Numerics.lean, lines 69–79}
\begin{leancode}
\leanline{0}{theorem\ euler\_contraction\ (F\ :\ H\ →\ H)\ (m\ h\ :\ ℝ)\ (L\ :\ ℝ≥0)}
\leanline{4}{(hF\ :\ ForceRegularity\ F\ m\ L)\ (hs\ :\ StableStep\ m\ L\ h)\ (x\ y\ ξ\ :\ H)\ :}
\leanline{4}{‖eulerStep\ F\ h\ x\ ξ\ -\ eulerStep\ F\ h\ y\ ξ‖\ ≤}
\leanline{6}{(1\ -\ m\ *\ h\ /\ 2)\ *\ ‖x\ -\ y‖}
\end{leancode}
\leanlabel{How the Lean proof follows the argument}
\texttt{euler\_\allowbreak{}contraction} first applies \texttt{euler\_\allowbreak{}sq\_\allowbreak{}contraction\_\allowbreak{}stable}. Its \texttt{hsq} is the squared estimate above; \texttt{hq} and \texttt{hr} establish nonnegativity before \texttt{nlinarith} compares the squares. The four fields of \texttt{Stable\allowbreak{}Step} encode precisely the two positivity assumptions and the two step restrictions.
\end{leanbox}
\subsection{\texorpdfstring{The local defect follows from the integral equation}{The local defect follows from the integral equation}}


\textbf{Claim.} Suppose every reference segment satisfies \(Y_k(0)=y_k\), \(Y_k(h)=y_{k+1}\), \(W_k(h)=\xi_k\), and

\[Y_k(t)-y_k=-\int_0^tF(Y_k(s))\,ds+W_k(t).\]

Assume the drift is interval integrable, \(\|F(Y_k(t))\|\le G\), and \(\|W_k(t)\|\le N\sqrt t\) for \(0\le t\le h\), with \(G,N\ge0\). Then the defect of a reference Euler step is at most \(R=L(h^2G/2+(2/3)h^{3/2}N)\).



\textbf{Proof.} The integral equation and the triangle inequality give \(\|Y_k(t)-y_k\|\le tG+N\sqrt t\). Subtract \(y_k-hF(y_k)+\xi_k\) from the endpoint equation; the noise increments cancel. The remaining vector is \(-\int_0^h(F(Y_k(s))-F(y_k))\,ds\). Its norm is at most

\[L\int_0^h\|Y_k(s)-y_k\|\,ds\le L\int_0^h(sG+N\sqrt s)\,ds=L\left(\frac{h^2G}{2}+\frac23h^{3/2}N\right).\]

The elementary integrals are \(h^2/2\) and \(2h^{3/2}/3\). In a random-state application the ambient norm is an \(L^2\) norm, so these same estimates are root-mean-square moment bounds. \(\square\)

\begin{leanbox}{Lean: \texttt{integral\_\allowbreak{}reference\_\allowbreak{}local\_\allowbreak{}defect}}
\leansource{\texttt{integral\_\allowbreak{}reference\_\allowbreak{}local\_\allowbreak{}defect} \quad LatentDiffusion/\allowbreak{}Numerics.lean, lines 286–303}
\begin{leancode}
\leanline{0}{theorem\ integral\_reference\_local\_defect\ (F\ :\ H\ →\ H)\ (m\ h\ G\ N\ :\ ℝ)\ (L\ :\ ℝ≥0)}
\leanline{4}{(hF\ :\ ForceRegularity\ F\ m\ L)\ (hh\ :\ 0\ ≤\ h)}
\leanline{4}{(y\ :\ ℕ\ →\ H)\ (ξ\ :\ ℕ\ →\ H)\ (Y\ W\ :\ ℕ\ →\ ℝ\ →\ H)}
\leanline{4}{(hr\ :\ IntegralReference\ F\ h\ G\ N\ y\ ξ\ Y\ W)\ (n\ :\ ℕ)\ :}
\leanline{4}{‖y\ (n\ +\ 1)\ -\ eulerStep\ F\ h\ (y\ n)\ (ξ\ n)‖\ ≤}
\leanline{6}{(L\ :\ ℝ)\ *\ (h\ \textasciicircum{}\ 2\ *\ G\ /\ 2\ +\ (2\ /\ 3\ :\ ℝ)\ *\ h\ *\ Real.sqrt\ h\ *\ N)}
\end{leancode}
\leanlabel{How the Lean proof follows the argument}
\texttt{Integral\allowbreak{}Reference} records \texttt{starts}, \texttt{ends}, \texttt{noise\_\allowbreak{}end}, \texttt{integrable}, \texttt{drift\_\allowbreak{}radius}, \texttt{noise\_\allowbreak{}growth} and \texttt{equation}. The proof of \texttt{integral\_\allowbreak{}reference\_\allowbreak{}local\_\allowbreak{}defect} uses these primitive data with \texttt{ULA.\allowbreak{}local\allowbreak{}Defect\_\allowbreak{}euler\_\allowbreak{}identity} and the Lipschitz force estimate. Thus the defect bound is derived from a supplied integral solution, rather than inserted as a conclusion-shaped hypothesis.
\end{leanbox}
\subsection{\texorpdfstring{Iteration separates initial error, force bias and time discretisation}{Iteration separates initial error, force bias and time discretisation}}


\textbf{Claim.} With the preceding hypotheses, suppose \(\|A(x_k)-F(x_k)\|\le B\) for every \(k\), where \(B\ge0\). Then

\[\|x_n-y_n\|\le q^n\|x_0-y_0\|+\frac{2B}{m}+\frac Lm\left(hG+\frac43N\sqrt h\right).\]

\textbf{Proof.} Insert two intermediate points, namely the Euler updates with force \(F\) at \(x_k\) and \(y_k\). The triangle inequality, the force error, the contraction and the local defect imply \(e_{k+1}\le qe_k+hB+R\), where \(e_k=\|x_k-y_k\|\). Induction yields \(e_n\le q^ne_0+(hB+R)\sum_{j=0}^{n-1}q^j\). Since \(0\le q<1\) and the forcing term is nonnegative, the finite sum is at most \(1/(1-q)=2/(mh)\). Hence

\[\frac{hB+R}{1-q}=\frac{2B}{m}+\frac Lm\left(hG+\frac43N\sqrt h\right).\]

The division is legitimate because \(m,h>0\). This proves the bound uniformly in the iteration count. It does not construct the reference path or establish its moment bounds for a new application. \(\square\)

\begin{leanbox}{Lean: \texttt{integral\_\allowbreak{}reference\_\allowbreak{}error\_\allowbreak{}bound}}
\leansource{\texttt{integral\_\allowbreak{}reference\_\allowbreak{}error\_\allowbreak{}bound} \quad LatentDiffusion/\allowbreak{}Numerics.lean, lines 310–337}
\begin{leancode}
\leanline{0}{theorem\ integral\_reference\_error\_bound\ (F\ A\ :\ H\ →\ H)\ (m\ h\ B\ G\ N\ :\ ℝ)\ (L\ :\ ℝ≥0)}
\leanline{4}{(hF\ :\ ForceRegularity\ F\ m\ L)\ (hs\ :\ StableStep\ m\ L\ h)}
\leanline{4}{(hB\ :\ 0\ ≤\ B)\ (hG\ :\ 0\ ≤\ G)\ (hN\ :\ 0\ ≤\ N)}
\leanline{4}{(y\ :\ ℕ\ →\ H)\ (ξ\ :\ ℕ\ →\ H)\ (Y\ W\ :\ ℕ\ →\ ℝ\ →\ H)\ (x₀\ :\ H)}
\leanline{4}{(hr\ :\ IntegralReference\ F\ h\ G\ N\ y\ ξ\ Y\ W)}
\leanline{4}{(hforce\ :\ ∀\ n,\ ‖A\ (eulerTrajectory\ A\ h\ ξ\ x₀\ n)\ -\ F\ (eulerTrajectory\ A\ h\ ξ\ x₀\ n)‖\ ≤\ B)}
\leanline{4}{(n\ :\ ℕ)\ :}
\leanline{4}{‖eulerTrajectory\ A\ h\ ξ\ x₀\ n\ -\ y\ n‖\ ≤}
\leanline{6}{(1\ -\ m\ *\ h\ /\ 2)\ \textasciicircum{}\ n\ *\ ‖x₀\ -\ y\ 0‖\ +}
\leanline{8}{2\ *\ B\ /\ m\ +\ (L\ :\ ℝ)\ /\ m\ *\ (h\ *\ G\ +\ (4\ /\ 3\ :\ ℝ)\ *\ N\ *\ Real.sqrt\ h)}
\end{leancode}
\leanlabel{How the Lean proof follows the argument}
In \texttt{integral\_\allowbreak{}reference\_\allowbreak{}error\_\allowbreak{}bound}, \texttt{R} is the displayed local defect and \texttt{hstep} is the affine scalar recurrence obtained from \texttt{euler\_\allowbreak{}local\_\allowbreak{}comparison}. The call to \texttt{ULA.\allowbreak{}affine\_\allowbreak{}recurrence\_\allowbreak{}bound} sums that recurrence. Finally \texttt{halgebra} uses \texttt{field\_\allowbreak{}simp} and \texttt{ring} to replace its denominator by the three explicit error contributions.
\end{leanbox}
\section{\texorpdfstring{The marginal evidence and its unique maximizer}{The marginal evidence and its unique maximizer}}


For each datum \(m\in\{1,\dots,M\}\) let \(U_m:\mathbb R^{d_\theta}\times\mathbb R^{d_z}\to\mathbb R\) be the joint potential, the negative logarithm of the tilted joint density of the parameter \(a\) and the latent variable \(z\). The pair space carries the Euclidean norm \(\|(a,z)\|^2=\|a\|^2+\|z\|^2\). Assumption A3 states that each \(U_m\) is differentiable and \(\lambda\)-strongly convex with one constant \(\lambda>0\):


\[U_m(sx+(1-s)y)\le sU_m(x)+(1-s)U_m(y)-\frac\lambda2s(1-s)\|x-y\|^2\qquad(x,y\in\mathbb R^{d_\theta+d_z},\ s\in[0,1]).\]


The normalizer and the marginal potential of datum \(m\) are


\[Z_m(a)=\int_{\mathbb R^{d_z}}e^{-U_m(a,z)}\,dz,\qquad V_m(a)=-\log Z_m(a),\qquad V(a)=\frac1M\sum_{m=1}^MV_m(a)=-\ell(a),\]


so \(V\) is minus the average log evidence \(\ell\). This section proves that \(V\) is \(\lambda\)-strongly convex, tends to \(+\infty\) at infinity and has exactly one minimizer \(a_\star\); that for every posterior family with log evidence \(-V\) the free energy has exactly one global minimizer \((a_\star,\pi_{a_\star})\); and that \(V\) grows quadratically away from \(a_\star\). The last step quantifies the effect of the tilting on the optimizer. The unique minimizer is the target of all convergence results of this report.

\subsection{\texorpdfstring{Strong convexity of the average marginal potential}{Strong convexity of the average marginal potential}}


Two facts about a single joint potential are used; they are proved in the companion report on interacting particle Langevin algorithms. Let \(U\) be differentiable and \(\lambda\)-strongly convex on \(\mathbb R^{d_\theta+d_z}\) with \(\lambda>0\), and put \(Z_U(a)=\int e^{-U(a,z)}dz\).



\begin{itemize}


\item \emph{(Integrability.)} For each \(a\), the function \(z\mapsto e^{-U(a,z)}\) is Lebesgue integrable and \(Z_U(a)>0\). Indeed, strong convexity gives \(U(x)\ge U(0)-\|\nabla U(0)\|^2/\lambda+\frac\lambda4\|x\|^2\), so \(e^{-U(a,z)}\le Ce^{-\lambda\|z\|^2/4}\), and an integral of a positive continuous function is positive.

\item \emph{(Prékopa–Leindler.)} The marginal potential \(a\mapsto-\log Z_U(a)\) is \(\lambda\)-strongly convex, with the same constant \(\lambda\).

\end{itemize}



\textbf{Claim.} Let \(M\ge1\), \(\lambda>0\), and let every \(U_m\) be differentiable and \(\lambda\)-strongly convex. Then \(V=M^{-1}\sum_mV_m\) is \(\lambda\)-strongly convex: for all \(x,y\in\mathbb R^{d_\theta}\) and all \(\alpha,\beta\ge0\) with \(\alpha+\beta=1\),


\[V(\alpha x+\beta y)\le\alpha V(x)+\beta V(y)-\alpha\beta\,\frac\lambda2\|x-y\|^2.\]


\textbf{Proof.} By the Prékopa–Leindler fact applied to \(U=U_m\), for each \(m\),


\[V_m(\alpha x+\beta y)\le\alpha V_m(x)+\beta V_m(y)-\alpha\beta\,\frac\lambda2\|x-y\|^2.\]


Summing over the \(M\) data, the last term contributes \(M\alpha\beta\frac\lambda2\|x-y\|^2\). Multiplying by \(1/M\ge0\) preserves the inequality and gives


\[V(\alpha x+\beta y)\le\alpha V(x)+\beta V(y)-\frac MM\,\alpha\beta\,\frac\lambda2\|x-y\|^2.\]


Since \(M\ge1\), \(M/M=1\). \(\square\)




The constant is not divided by \(M\) because the potentials are averaged. Averaging is the normalization used throughout this report.

\begin{leanbox}{Lean: \texttt{averageMarginalPotential}, \texttt{averageMarginalPotential\_\allowbreak{}strongConvex}}
\leansource{\texttt{averageMarginalPotential} \quad LatentDiffusion/\allowbreak{}Marginal.lean, lines 13–15}
\begin{leancode}
\leanline{0}{def\ averageMarginalPotential\ (U\ :\ Fin\ M\ →\ IPLA.Joint\ dθ\ dz\ →\ ℝ)}
\leanline{4}{(θ\ :\ EuclideanSpace\ ℝ\ (Fin\ dθ))\ :\ ℝ\ :=}
\leanline{2}{(M\ :\ ℝ)⁻¹\ *\ ∑\ i,\ IPLA.marginalPotential\ (U\ i)\ θ}
\end{leancode}
\leansource{\texttt{averageMarginalPotential\_\allowbreak{}strongConvex} \quad LatentDiffusion/\allowbreak{}Marginal.lean, lines 17–32}
\begin{leancode}
\leanline{0}{theorem\ averageMarginalPotential\_strongConvex}
\leanline{4}{(U\ :\ Fin\ M\ →\ IPLA.Joint\ dθ\ dz\ →\ ℝ)\ (hM\ :\ 0\ <\ M)}
\leanline{4}{(rho\ :\ ℝ)\ (hrho\ :\ 0\ <\ rho)\ (hdiff\ :\ ∀\ i,\ Differentiable\ ℝ\ (U\ i))}
\leanline{4}{(hsc\ :\ ∀\ i,\ StrongConvexOn\ univ\ rho\ (U\ i))\ :}
\leanline{4}{StrongConvexOn\ univ\ rho\ (averageMarginalPotential\ U)}
\end{leancode}
\leanlabel{How the Lean proof follows the argument}
\texttt{average\allowbreak{}Marginal\allowbreak{}Potential U \ensuremath{\theta}} is \(M^{-1}\sum_iV_i(\theta)\), with \(V_i\) given by \texttt{IPLA.\allowbreak{}marginal\allowbreak{}Potential (U i)}, the negative logarithm of \texttt{IPLA.\allowbreak{}marginal\allowbreak{}Integral}, the integral of \(e^{-U_i(\theta,x)}\) over the latent variable. Here \texttt{IPLA.\allowbreak{}Joint d\ensuremath{\theta} dz} is the Euclidean space of pairs, so its norm is the Euclidean norm of the summary, and the parameter is called \texttt{\ensuremath{\theta}} in the Lean statements. Mathlib's \texttt{Strong\allowbreak{}Convex\allowbreak{}On univ rho f} is exactly the displayed inequality with weights \(\alpha,\beta\ge0\), \(\alpha+\beta=1\), and modulus \texttt{rho} \(=\lambda\).



\begin{enumerate}


\item \texttt{refine \ensuremath{\langle}convex\_\allowbreak{}univ, ?\_\allowbreak{}\ensuremath{\rangle}} reduces strong convexity on the whole space to the inequality.

\item \texttt{hs} sums the per-datum inequalities supplied by \texttt{IPLA.\allowbreak{}marginal\allowbreak{}Potential\_\allowbreak{}strong\allowbreak{}Convex} (the Prékopa–Leindler fact) with \texttt{Finset.\allowbreak{}sum\_\allowbreak{}le\_\allowbreak{}sum}, and \texttt{hh} multiplies by \(M^{-1}\ge0\).

\item The \texttt{simp only} call distributes the sum over the three terms and evaluates the sum of the constant term as \(M\) times that term (\texttt{Finset.\allowbreak{}sum\_\allowbreak{}const}, \texttt{Fintype.\allowbreak{}card\_\allowbreak{}fin}, \texttt{nsmul\_\allowbreak{}eq\_\allowbreak{}mul}). After \texttt{convert hh using 1}, \texttt{field\_\allowbreak{}simp} uses \(M\ne0\) to cancel \(M^{-1}M\).

\end{enumerate}
\end{leanbox}
\subsection{\texorpdfstring{Coercivity of the marginal potentials}{Coercivity of the marginal potentials}}


A third fact from the companion report is used: under the hypotheses of the previous step, \(Z_U(a)\to0\) as \(\|a\|\to\infty\).




\textbf{Claim.} (1) If \(U\) is differentiable and \(\lambda\)-strongly convex with \(\lambda>0\), then \(-\log Z_U(a)\to+\infty\) as \(\|a\|\to\infty\). (2) If in addition \(M\ge1\) and every \(U_m\) satisfies these hypotheses, then \(V(a)\to+\infty\) as \(\|a\|\to\infty\).




\textbf{Proof.} (1) By the cited fact and positivity of \(Z_U\), \(Z_U(a)\) tends to \(0\) through positive values. Since \(\log t\to-\infty\) as \(t\downarrow0\), the composition \(\log Z_U(a)\to-\infty\), and its negative tends to \(+\infty\).




(2) Let \(B\in\mathbb R\). By (1), for each \(m\) there is \(R_m\) with \(V_m(a)\ge B\) whenever \(\|a\|\ge R_m\). Put \(R=\max_mR_m\), a maximum over finitely many numbers. For \(\|a\|\ge R\), all \(M\) inequalities hold, and averaging them gives


\[V(a)=\frac1M\sum_{m=1}^MV_m(a)\ge\frac1M\cdot MB=B.\]


Since \(B\) was arbitrary, \(V(a)\to+\infty\). \(\square\)

\begin{leanbox}{Lean: \texttt{marginalPotential\_\allowbreak{}coercive}, \texttt{averageMarginalPotential\_\allowbreak{}coercive}}
\leansource{\texttt{marginalPotential\_\allowbreak{}coercive} \quad LatentDiffusion/\allowbreak{}Marginal.lean, lines 34–41}
\begin{leancode}
\leanline{0}{theorem\ marginalPotential\_coercive\ (U\ :\ IPLA.Joint\ dθ\ dz\ →\ ℝ)}
\leanline{4}{(rho\ :\ ℝ)\ (hrho\ :\ 0\ <\ rho)\ (hdiff\ :\ Differentiable\ ℝ\ U)}
\leanline{4}{(hsc\ :\ StrongConvexOn\ univ\ rho\ U)\ :}
\leanline{4}{Tendsto\ (IPLA.marginalPotential\ U)\ (cocompact\ \_)\ atTop}
\end{leancode}
\leansource{\texttt{averageMarginalPotential\_\allowbreak{}coercive} \quad LatentDiffusion/\allowbreak{}Marginal.lean, lines 43–60}
\begin{leancode}
\leanline{0}{theorem\ averageMarginalPotential\_coercive}
\leanline{4}{(U\ :\ Fin\ M\ →\ IPLA.Joint\ dθ\ dz\ →\ ℝ)\ (hM\ :\ 0\ <\ M)}
\leanline{4}{(rho\ :\ ℝ)\ (hrho\ :\ 0\ <\ rho)\ (hdiff\ :\ ∀\ i,\ Differentiable\ ℝ\ (U\ i))}
\leanline{4}{(hsc\ :\ ∀\ i,\ StrongConvexOn\ univ\ rho\ (U\ i))\ :}
\leanline{4}{Tendsto\ (averageMarginalPotential\ U)\ (cocompact\ \_)\ atTop}
\end{leancode}
\leanlabel{How the Lean proof follows the argument}
Mathlib expresses “as \(\|a\|\to\infty\)” with the filter \texttt{cocompact}, whose sets are the complements of compact sets; in finite dimensions these contain the complements of large balls. \texttt{Tendsto f (cocompact \_\allowbreak{}) at\allowbreak{}Top} therefore means \(f(a)\to+\infty\) as \(\|a\|\to\infty\).



\begin{enumerate}


\item \texttt{marginal\allowbreak{}Potential\_\allowbreak{}coercive}: \texttt{ht} states that the normalizer tends to \(0\) within \((0,\infty)\), written \texttt{{\leanfallback 𝓝}[>] 0}; it combines \texttt{IPLA.\allowbreak{}tendsto\_\allowbreak{}marginal\_\allowbreak{}integral\_\allowbreak{}cocompact} with the positivity \texttt{IPLA.\allowbreak{}marginal\allowbreak{}Integral\_\allowbreak{}pos} through \texttt{tendsto\_\allowbreak{}nhds\allowbreak{}Within\_\allowbreak{}iff}. Composing with \texttt{Real.\allowbreak{}tendsto\_\allowbreak{}log\_\allowbreak{}nhds\allowbreak{}GT\_\allowbreak{}zero} (\(\log t\to-\infty\) as \(t\downarrow0\)) and \texttt{tendsto\_\allowbreak{}neg\_\allowbreak{}at\allowbreak{}Bot\_\allowbreak{}at\allowbreak{}Top} (negation) gives the claim.

\item \texttt{average\allowbreak{}Marginal\allowbreak{}Potential\_\allowbreak{}coercive}: \texttt{tendsto\_\allowbreak{}at\allowbreak{}Top.\allowbreak{}2} reduces the claim to “eventually \(V\ge b\)” for each \texttt{b}. \texttt{hh} uses \texttt{eventually\_\allowbreak{}all}, valid for the finite index set, to obtain all \(M\) inequalities simultaneously; this is the maximum \(R=\max_mR_m\) of the proof. \texttt{filter\_\allowbreak{}upwards} fixes such a point, the inequalities are summed, multiplied by \(M^{-1}\), and \texttt{simpa} cancels \(M^{-1}M\) using \texttt{hp.\allowbreak{}ne'}.

\end{enumerate}
\end{leanbox}
\subsection{\texorpdfstring{Existence and uniqueness of the minimizer}{Existence and uniqueness of the minimizer}}


\textbf{Claim.} Let \(M\ge1\), \(\lambda>0\), and let every \(U_m\) be differentiable and \(\lambda\)-strongly convex. Then there is exactly one \(a_\star\in\mathbb R^{d_\theta}\) with \(V(a_\star)\le V(a)\) for all \(a\). Equivalently, the average log evidence \(\ell=-V\) has exactly one maximizer.




\textbf{Proof.} \emph{Continuity.} Each \(V_m\) is continuous (companion report: the normalizer is continuous and positive, and \(\log\) is continuous on \((0,\infty)\)); hence \(V\), a constant multiple of a finite sum, is continuous.




\emph{Existence.} Fix \(a_0\). By the previous step there is \(R\) with \(V(a)>V(a_0)\) for \(\|a\|\ge R\). The closed ball \(\bar B_R=\{\|a\|\le R\}\) is compact (Heine–Borel) and nonempty, since \(a_0\in\bar B_R\). A continuous function on a nonempty compact set attains its minimum, say at \(a_\star\in\bar B_R\). Then \(V(a_\star)\le V(a)\) for \(a\in\bar B_R\), and for \(\|a\|\ge R\) we have \(V(a)>V(a_0)\ge V(a_\star)\). Thus \(a_\star\) is a global minimizer.




\emph{Uniqueness.} Let \(a\) and \(b\) both be global minimizers, so \(V(a)=V(b)=\min V\). Strong convexity with \(\alpha=\beta=\frac12\) gives


\[V\Bigl(\frac{a+b}2\Bigr)\le\frac{V(a)+V(b)}2-\frac\lambda8\|a-b\|^2=\min V-\frac\lambda8\|a-b\|^2.\]


If \(a\ne b\), the right side is strictly smaller than \(\min V\) because \(\lambda>0\), which is impossible. Hence \(a=b\). \(\square\)

\begin{leanbox}{Lean: \texttt{existsUnique\_\allowbreak{}averageMarginalMinimizer}}
\leansource{\texttt{existsUnique\_\allowbreak{}averageMarginalMinimizer} \quad LatentDiffusion/\allowbreak{}Marginal.lean, lines 62–75}
\begin{leancode}
\leanline{0}{theorem\ existsUnique\_averageMarginalMinimizer}
\leanline{4}{(U\ :\ Fin\ M\ →\ IPLA.Joint\ dθ\ dz\ →\ ℝ)\ (hM\ :\ 0\ <\ M)}
\leanline{4}{(rho\ :\ ℝ)\ (hrho\ :\ 0\ <\ rho)\ (hdiff\ :\ ∀\ i,\ Differentiable\ ℝ\ (U\ i))}
\leanline{4}{(hsc\ :\ ∀\ i,\ StrongConvexOn\ univ\ rho\ (U\ i))\ :}
\leanline{4}{∃!\ θ,\ ∀\ θ',\ averageMarginalPotential\ U\ θ\ ≤\ averageMarginalPotential\ U\ θ'}
\end{leancode}
\leanlabel{How the Lean proof follows the argument}
The conclusion \texttt{\ensuremath{\exists}! \ensuremath{\theta}, \ensuremath{\forall} \ensuremath{\theta}', average\allowbreak{}Marginal\allowbreak{}Potential U \ensuremath{\theta} \ensuremath{\le} average\allowbreak{}Marginal\allowbreak{}Potential U \ensuremath{\theta}'} is “there exists a unique global minimizer”.



\begin{enumerate}


\item \texttt{hc} proves continuity from \texttt{IPLA.\allowbreak{}marginal\allowbreak{}Potential\_\allowbreak{}continuous}, \texttt{continuous\_\allowbreak{}finset\allowbreak{}Sum} and multiplication by a constant.

\item \texttt{hc.\allowbreak{}exists\_\allowbreak{}forall\_\allowbreak{}le} is Mathlib's \texttt{Continuous.\allowbreak{}exists\_\allowbreak{}forall\_\allowbreak{}le}: a continuous function that tends to \(+\infty\) along \texttt{cocompact} has a global minimum. Its coercivity input is the previous step. This packages the compactness argument of the proof.

\item Uniqueness: \texttt{Strong\allowbreak{}Convex\allowbreak{}On.\allowbreak{}strict\allowbreak{}Convex\allowbreak{}On} turns strong convexity with \texttt{hrho : 0 < rho} into strict convexity, and \texttt{Strict\allowbreak{}Convex\allowbreak{}On.\allowbreak{}eq\_\allowbreak{}of\_\allowbreak{}is\allowbreak{}Min\allowbreak{}On} states that a strictly convex function has at most one minimizer on a convex set; its proof is the midpoint argument above.

\end{enumerate}
\end{leanbox}
\subsection{\texorpdfstring{The unique variational optimizer}{The unique variational optimizer}}


Recall from the variational section: for a posterior family \((\ell,\pi)\) with \(M\ge1\) data, \((a,q)\) is a global minimizer of \(F(a,q)=-\ell(a)+M^{-1}\sum_m\mathrm{KL}(q^m\|\pi_a^m)\) over admissible pairs if and only if \(a\) maximizes \(\ell\) and \(q^m=\pi_a^m\) for all \(m\).




\textbf{Claim.} Let \(M\ge1\), \(\lambda>0\), and let every \(U_m\) be differentiable and \(\lambda\)-strongly convex. Let \(Z\) be any measurable space and \((\ell,\pi)\) any posterior family on \(\mathbb R^{d_\theta}\) and \(Z\) whose log evidence is \(\ell(a)=-V(a)\) for every \(a\). Then there is exactly one pair \((a,q)\) that is a global minimizer of \(F\), namely \((a_\star,\pi_{a_\star})\) with \(a_\star\) the minimizer of \(V\).




\textbf{Proof.} \emph{Existence.} Let \(a_\star\) be the minimizer of the previous step. For every \(b\), \(\ell(b)=-V(b)\le-V(a_\star)=\ell(a_\star)\), so \(a_\star\) maximizes \(\ell\). By the characterization, \((a_\star,\pi_{a_\star})\) is a global minimizer.




\emph{Uniqueness.} Let \((b,q)\) be a global minimizer. By the characterization, \(b\) maximizes \(\ell\) and \(q=\pi_b\). For every \(c\), \(-V(c)=\ell(c)\le\ell(b)=-V(b)\), so \(b\) minimizes \(V\). By uniqueness in the previous step, \(b=a_\star\), and then \(q=\pi_b=\pi_{a_\star}\). \(\square\)




The posterior family is arbitrary here; normalization of its posteriors is part of the family. The next section constructs the family from the densities \(e^{-U_m}\) and verifies the hypothesis \(\ell=-V\).

\begin{leanbox}{Lean: \texttt{existsUnique\_\allowbreak{}variationalOptimizer}}
\leansource{\texttt{existsUnique\_\allowbreak{}variationalOptimizer} \quad LatentDiffusion/\allowbreak{}Marginal.lean, lines 79–102}
\begin{leancode}
\leanline{0}{theorem\ existsUnique\_variationalOptimizer}
\leanline{4}{\{Z\ :\ Type*\}\ [MeasurableSpace\ Z]}
\leanline{4}{(U\ :\ Fin\ M\ →\ IPLA.Joint\ dθ\ dz\ →\ ℝ)\ (hM\ :\ 0\ <\ M)}
\leanline{4}{(rho\ :\ ℝ)\ (hrho\ :\ 0\ <\ rho)\ (hdiff\ :\ ∀\ i,\ Differentiable\ ℝ\ (U\ i))}
\leanline{4}{(hsc\ :\ ∀\ i,\ StrongConvexOn\ univ\ rho\ (U\ i))}
\leanline{4}{(P\ :\ PosteriorFamily\ (EuclideanSpace\ ℝ\ (Fin\ dθ))\ Z\ M)}
\leanline{4}{(hlog\ :\ ∀\ θ,\ P.logEvidence\ θ\ =\ -averageMarginalPotential\ U\ θ)\ :}
\leanline{4}{∃!\ p\ :\ EuclideanSpace\ ℝ\ (Fin\ dθ)\ ×\ (Fin\ M\ →\ ProbabilityMeasure\ Z),}
\leanline{6}{P.IsGlobalMin\ p.1\ p.2}
\end{leancode}
\leanlabel{How the Lean proof follows the argument}
In \texttt{exists\allowbreak{}Unique\_\allowbreak{}variational\allowbreak{}Optimizer} the family is \texttt{P : Posterior\allowbreak{}Family (Euclidean\allowbreak{}Space \ensuremath{\mathbb R} (Fin d\ensuremath{\theta})) Z M} and \texttt{hlog} is \(\ell=-V\). The conclusion is unique existence of a pair \texttt{p} with \texttt{P.\allowbreak{}Is\allowbreak{}Global\allowbreak{}Min p.\allowbreak{}1 p.\allowbreak{}2}.



\begin{enumerate}


\item \texttt{obtain \ensuremath{\langle}\ensuremath{\theta}, h\ensuremath{\theta}, hu\ensuremath{\rangle}} extracts the minimizer and its uniqueness from \texttt{exists\allowbreak{}Unique\_\allowbreak{}average\allowbreak{}Marginal\allowbreak{}Minimizer}.

\item Existence: the witness is \texttt{(\ensuremath{\theta}, P.\allowbreak{}posterior \ensuremath{\theta})}; \texttt{(P.\allowbreak{}global\_\allowbreak{}minimizers\_\allowbreak{}iff h\allowbreak{}M \ensuremath{\theta} \_\allowbreak{}).\allowbreak{}mpr} reduces minimality to maximality of \(\ell\), which follows from \texttt{h\ensuremath{\theta} b} after rewriting both values with \texttt{hlog} and applying \texttt{neg\_\allowbreak{}le\_\allowbreak{}neg}.

\item Uniqueness: for a minimizer \texttt{\ensuremath{\langle}b, q\ensuremath{\rangle}}, \texttt{(P.\allowbreak{}global\_\allowbreak{}minimizers\_\allowbreak{}iff h\allowbreak{}M b q).\allowbreak{}mp hq} gives maximality of \texttt{b} and, through the pattern \texttt{rfl}, the identity \texttt{q = P.\allowbreak{}posterior b}. \texttt{he : b = \ensuremath{\theta}} follows from \texttt{hu} once \(V(b)\le V(c)\) is derived by \texttt{linarith} from the rewritten maximality. \texttt{subst b} and \texttt{rfl} finish.

\end{enumerate}
\end{leanbox}
\subsection{\texorpdfstring{Quadratic growth at a minimizer}{Quadratic growth at a minimizer}}


\textbf{Claim.} Let \(H\) be a real normed vector space, \(\lambda\in\mathbb R\), and \(F:H\to\mathbb R\) \(\lambda\)-strongly convex on \(H\). If \(z\) minimizes \(F\), then for every \(x\in H\),


\[\frac\lambda2\|x-z\|^2\le F(x)-F(z).\]


No differentiability of \(F\) is assumed.




\textbf{Proof.} Fix \(x\) and let \(0<s<1\). Strong convexity with the weights \(1-s\) and \(s\) at the points \(z\) and \(x\) gives


\[F((1-s)z+sx)\le(1-s)F(z)+sF(x)-(1-s)s\,\frac\lambda2\|z-x\|^2.\]


Minimality of \(z\) gives \(F(z)\le F((1-s)z+sx)\). Combining the two inequalities and subtracting \((1-s)F(z)\),


\[s\cdot\frac\lambda2(1-s)\|x-z\|^2\le s\,\bigl(F(x)-F(z)\bigr).\]


Dividing by \(s>0\) gives \(\frac\lambda2(1-s)\|x-z\|^2\le F(x)-F(z)\) for every \(s\in(0,1)\). The left side is a continuous function of \(s\) with limit \(\frac\lambda2\|x-z\|^2\) as \(s\downarrow0\), and a limit of numbers bounded above by \(F(x)-F(z)\) is bounded by the same number. \(\square\)

\begin{leanbox}{Lean: \texttt{strongConvex\_\allowbreak{}quadratic\_\allowbreak{}growth}}
\leansource{\texttt{strongConvex\_\allowbreak{}quadratic\_\allowbreak{}growth} \quad LatentDiffusion/\allowbreak{}Marginal.lean, lines 105–127}
\begin{leancode}
\leanline{0}{theorem\ strongConvex\_quadratic\_growth}
\leanline{4}{\{H\ :\ Type*\}\ [NormedAddCommGroup\ H]\ [NormedSpace\ ℝ\ H]}
\leanline{4}{(F\ :\ H\ →\ ℝ)\ (rho\ :\ ℝ)\ (hsc\ :\ StrongConvexOn\ univ\ rho\ F)}
\leanline{4}{(z\ :\ H)\ (hz\ :\ ∀\ x,\ F\ z\ ≤\ F\ x)\ (x\ :\ H)\ :}
\leanline{4}{rho\ /\ 2\ *\ ‖x\ -\ z‖\ \textasciicircum{}\ 2\ ≤\ F\ x\ -\ F\ z}
\end{leancode}
\leanlabel{How the Lean proof follows the argument}
In \texttt{strong\allowbreak{}Convex\_\allowbreak{}quadratic\_\allowbreak{}growth} the hypothesis \texttt{hz : \ensuremath{\forall} x, F z \ensuremath{\le} F x} is minimality; \texttt{rho} is \(\lambda\) and has no sign condition.



\begin{enumerate}


\item \texttt{hb} is the inequality \(\frac\lambda2(1-s)\|x-z\|^2\le F(x)-F(z)\) for all \(s\) in a right neighbourhood of \(0\), written \texttt{\ensuremath{\forall}ᶠ s in {\leanfallback 𝓝}[>] 0}. \texttt{filter\_\allowbreak{}upwards} restricts to \(0<s<1\), using \texttt{self\_\allowbreak{}mem\_\allowbreak{}nhds\allowbreak{}Within} and \texttt{eventually\_\allowbreak{}lt\_\allowbreak{}nhds}.

\item For such \(s\), \texttt{hh} is strong convexity at \(z\) and \(x\) with weights \(1-s\) and \(s\), \texttt{hm} is minimality at the interpolated point, and \texttt{norm\_\allowbreak{}sub\_\allowbreak{}rev} rewrites \(\|z-x\|=\|x-z\|\). \texttt{(mul\_\allowbreak{}le\_\allowbreak{}mul\_\allowbreak{}iff\_\allowbreak{}right\ensuremath{{}_{0}} hsp).\allowbreak{}mp} cancels the positive factor \(s\), reducing the goal to the multiplied inequality, which \texttt{nlinarith} proves from \texttt{hh} and \texttt{hm}.

\item \texttt{ht} is the limit of the continuous left side as \(s\downarrow0\), and \texttt{le\_\allowbreak{}of\_\allowbreak{}tendsto ht hb} passes the bound to the limit.

\end{enumerate}
\end{leanbox}
\subsection{\texorpdfstring{Simultaneous control of parameter and posterior errors}{Simultaneous control of parameter and posterior errors}}


\textbf{Claim.} Assume the hypotheses of the unique-optimizer step: \(M\ge1\), \(\lambda>0\), every \(U_m\) differentiable and \(\lambda\)-strongly convex, and \((\ell,\pi)\) a posterior family with \(\ell=-V\). Let \(a_\star\) minimize \(V\). Then for every parameter \(a\) and every collection \(q\) of probability measures,


\[\frac\lambda2\|a-a_\star\|^2+\frac1M\sum_{m=1}^Mk(q^m\|\pi_a^m)\le F(a,q)-F(a_\star,\pi_{a_\star}),\]


where \(k\) is the relative entropy with infinite values replaced by \(0\), as in the definition of \(F\). For admissible \(q\) the second term is the average relative entropy to the posteriors.




\textbf{Proof.} By the posterior identity of the variational section, \(F(a_\star,\pi_{a_\star})=-\ell(a_\star)=V(a_\star)\). By definition \(F(a,q)=-\ell(a)+M^{-1}\sum_mk(q^m\|\pi_a^m)=V(a)+M^{-1}\sum_mk(q^m\|\pi_a^m)\). Hence


\[F(a,q)-F(a_\star,\pi_{a_\star})=V(a)-V(a_\star)+\frac1M\sum_mk(q^m\|\pi_a^m).\]


By strong convexity of \(V\) and the quadratic growth step, \(V(a)-V(a_\star)\ge\frac\lambda2\|a-a_\star\|^2\). Substituting this bound proves the claim. \(\square\)




Thus a small free-energy excess forces both a small parameter error and small posterior relative entropies.

\begin{leanbox}{Lean: \texttt{parameter\_\allowbreak{}posterior\_\allowbreak{}control}}
\leansource{\texttt{parameter\_\allowbreak{}posterior\_\allowbreak{}control} \quad LatentDiffusion/\allowbreak{}Marginal.lean, lines 131–150}
\begin{leancode}
\leanline{0}{theorem\ parameter\_posterior\_control}
\leanline{4}{\{Z\ :\ Type*\}\ [MeasurableSpace\ Z]}
\leanline{4}{(U\ :\ Fin\ M\ →\ IPLA.Joint\ dθ\ dz\ →\ ℝ)\ (hM\ :\ 0\ <\ M)}
\leanline{4}{(rho\ :\ ℝ)\ (hrho\ :\ 0\ <\ rho)\ (hdiff\ :\ ∀\ i,\ Differentiable\ ℝ\ (U\ i))}
\leanline{4}{(hsc\ :\ ∀\ i,\ StrongConvexOn\ univ\ rho\ (U\ i))}
\leanline{4}{(P\ :\ PosteriorFamily\ (EuclideanSpace\ ℝ\ (Fin\ dθ))\ Z\ M)}
\leanline{4}{(hlog\ :\ ∀\ θ,\ P.logEvidence\ θ\ =\ -averageMarginalPotential\ U\ θ)}
\leanline{4}{(θstar\ θ\ :\ EuclideanSpace\ ℝ\ (Fin\ dθ))}
\leanline{4}{(hmin\ :\ ∀\ b,\ averageMarginalPotential\ U\ θstar\ ≤\ averageMarginalPotential\ U\ b)}
\leanline{4}{(q\ :\ Fin\ M\ →\ ProbabilityMeasure\ Z)\ :}
\leanline{4}{rho\ /\ 2\ *\ ‖θ\ -\ θstar‖\ \textasciicircum{}\ 2\ +}
\leanline{6}{(M\ :\ ℝ)⁻¹\ *\ ∑\ i,\ (InformationTheory.klDiv\ (q\ i\ :\ Measure\ Z)}
\leanline{8}{(P.posterior\ θ\ i\ :\ Measure\ Z)).toReal\ ≤}
\leanline{4}{P.freeEnergy\ θ\ q\ -\ P.freeEnergy\ θstar\ (P.posterior\ θstar)}
\end{leancode}
\leanlabel{How the Lean proof follows the argument}
In \texttt{parameter\_\allowbreak{}posterior\_\allowbreak{}control}, \texttt{hmin} states that \texttt{\ensuremath{\theta}star} minimizes \(V\). \texttt{hh} is the quadratic growth of \texttt{average\allowbreak{}Marginal\allowbreak{}Potential U} at \texttt{\ensuremath{\theta}star}, obtained from \texttt{strong\allowbreak{}Convex\_\allowbreak{}quadratic\_\allowbreak{}growth} and \texttt{average\allowbreak{}Marginal\allowbreak{}Potential\_\allowbreak{}strong\allowbreak{}Convex}. The proof rewrites \(F(a_\star,\pi_{a_\star})\) with \texttt{P.\allowbreak{}posterior\_\allowbreak{}free\allowbreak{}Energy}, unfolds \texttt{Posterior\allowbreak{}Family.\allowbreak{}free\allowbreak{}Energy}, replaces both log evidences by \texttt{hlog}, and closes the resulting linear inequality with \texttt{linarith}.
\end{leanbox}
\subsection{\texorpdfstring{Robustness of the optimizer to the tilting bias}{Robustness of the optimizer to the tilting bias}}


Let \(\ell\) be the original evidence and \(\tilde\ell\) the tilted evidence, both real functions on a set \(\mathcal A\).




\textbf{Claim 1.} Let \(\delta,\varepsilon\in\mathbb R\). Assume a one-sided uniform gap \(0\le\ell(x)-\tilde\ell(x)\le\delta\) for all \(x\), and that \(\hat a\) is an \(\varepsilon\)-maximizer of \(\tilde\ell\): \(\tilde\ell(x)\le\tilde\ell(\hat a)+\varepsilon\) for all \(x\). Then \(\ell(a)-\ell(\hat a)\le\delta+\varepsilon\) for every \(a\in\mathcal A\).




\textbf{Proof.} The gap bound at \(a\), the \(\varepsilon\)-maximality, and the lower gap bound at \(\hat a\) give


\[\ell(a)\le\tilde\ell(a)+\delta\le\tilde\ell(\hat a)+\varepsilon+\delta\le\ell(\hat a)+\varepsilon+\delta.\qquad\square\]


\textbf{Claim 2.} Let \(H\) be a real normed vector space, \(\lambda>0\), \(\delta,\varepsilon\ge0\), and suppose \(-\ell\) is \(\lambda\)-strongly convex on \(H\). Assume the gap and \(\varepsilon\)-maximality conditions of Claim 1 and let \(a_\star\) maximize \(\ell\). Then


\[\|\hat a-a_\star\|\le\sqrt{\frac{2(\delta+\varepsilon)}\lambda}.\]


\textbf{Proof.} Quadratic growth of \(-\ell\) at its minimizer \(a_\star\) gives \(\frac\lambda2\|\hat a-a_\star\|^2\le\ell(a_\star)-\ell(\hat a)\), and Claim 1 bounds the right side by \(\delta+\varepsilon\). Dividing by \(\lambda/2>0\) gives \(\|\hat a-a_\star\|^2\le2(\delta+\varepsilon)/\lambda\). Both sides are nonnegative, so the inequality between their square roots follows from monotonicity of the square root. \(\square\)




The target here is a maximizer of the original evidence, not of the tilted one; the gap \(\delta\) is a model hypothesis and is not inferred from a pointwise training penalty.

\begin{leanbox}{Lean: \texttt{tilted\_\allowbreak{}objective\_\allowbreak{}bias}, \texttt{tilted\_\allowbreak{}parameter\_\allowbreak{}bias}}
\leansource{\texttt{tilted\_\allowbreak{}objective\_\allowbreak{}bias} \quad LatentDiffusion/\allowbreak{}Marginal.lean, lines 213–221}
\begin{leancode}
\leanline{0}{theorem\ tilted\_objective\_bias}
\leanline{4}{\{A\ :\ Type*\}\ (ell\ tilted\ :\ A\ →\ ℝ)\ (delta\ epsilon\ :\ ℝ)}
\leanline{4}{(hgap\ :\ ∀\ x,\ 0\ ≤\ ell\ x\ -\ tilted\ x\ ∧\ ell\ x\ -\ tilted\ x\ ≤\ delta)}
\leanline{4}{(star\ approx\ :\ A)\ (happrox\ :\ ∀\ x,\ tilted\ x\ ≤\ tilted\ approx\ +\ epsilon)\ :}
\leanline{4}{ell\ star\ -\ ell\ approx\ ≤\ delta\ +\ epsilon}
\end{leancode}
\leansource{\texttt{tilted\_\allowbreak{}parameter\_\allowbreak{}bias} \quad LatentDiffusion/\allowbreak{}Marginal.lean, lines 226–244}
\begin{leancode}
\leanline{0}{theorem\ tilted\_parameter\_bias}
\leanline{4}{\{H\ :\ Type*\}\ [NormedAddCommGroup\ H]\ [NormedSpace\ ℝ\ H]}
\leanline{4}{(ell\ tilted\ :\ H\ →\ ℝ)\ (rho\ delta\ epsilon\ :\ ℝ)\ (hrho\ :\ 0\ <\ rho)}
\leanline{4}{(hdelta\ :\ 0\ ≤\ delta)\ (hepsilon\ :\ 0\ ≤\ epsilon)}
\leanline{4}{(hsc\ :\ StrongConvexOn\ univ\ rho\ (fun\ x\ =>\ -ell\ x))}
\leanline{4}{(hgap\ :\ ∀\ x,\ 0\ ≤\ ell\ x\ -\ tilted\ x\ ∧\ ell\ x\ -\ tilted\ x\ ≤\ delta)}
\leanline{4}{(star\ approx\ :\ H)\ (hstar\ :\ ∀\ x,\ ell\ x\ ≤\ ell\ star)}
\leanline{4}{(happrox\ :\ ∀\ x,\ tilted\ x\ ≤\ tilted\ approx\ +\ epsilon)\ :}
\leanline{4}{‖approx\ -\ star‖\ ≤\ Real.sqrt\ (2\ *\ (delta\ +\ epsilon)\ /\ rho)}
\end{leancode}
\leanlabel{How the Lean proof follows the argument}
\begin{enumerate}


\item \texttt{tilted\_\allowbreak{}objective\_\allowbreak{}bias}: \texttt{hgap} is the two-sided gap condition and \texttt{happrox} the \(\varepsilon\)-maximality. The point \texttt{star} is arbitrary. The proof takes the upper gap bound at \texttt{star}, the lower bound at \texttt{approx} and \texttt{happrox star}, and \texttt{linarith} chains them.

\item \texttt{tilted\_\allowbreak{}parameter\_\allowbreak{}bias}: \texttt{hsc} is strong convexity of \texttt{fun x => -\allowbreak{}ell x} and \texttt{hstar} maximality of \texttt{star}. \texttt{hg} is quadratic growth from \texttt{strong\allowbreak{}Convex\_\allowbreak{}quadratic\_\allowbreak{}growth}, \texttt{hb} is Claim 1, and \texttt{hs} divides by \(\lambda\) through \texttt{le\_\allowbreak{}div\_\allowbreak{}iff\ensuremath{{}_{0}}}. The final \texttt{nlinarith} passes to square roots using \texttt{Real.\allowbreak{}sq\_\allowbreak{}sqrt} and nonnegativity of both sides.

\end{enumerate}
\end{leanbox}
\section{\texorpdfstring{Displacement convexity of the free energy}{Displacement convexity of the free energy}}


This section proves the strong geodesic convexity of the tilted free energy along displacement interpolations, the paper's Lemma on geodesic convexity,.




The extra transport assumptions are R4.




Let \(Z\) be a finite-dimensional real inner product space of dimension \(d\), let \(\mu\) be a probability measure on \(Z\) (the initial law) and let \(T:Z\to Z\) be a transport map. The displacement interpolation is \(T_s(z)=(1-s)z+sT(z)\), \(s\in[0,1]\), and its law is the pushforward \(\mu_s=\mu\circ T_s^{-1}\). For one datum, a parameter interpolation \(p_s=(1-s)p_0+sp_1\) and a joint potential \(U(p,z)\), the free energy along the interpolation is


\[\mathcal F(s)=H(s)+\int_ZU(p_s,z)\,\mu_s(dz),\qquad H(s)=\int_Z\log\rho_s\,d\mu_s,\]


where \(\rho_s\) is the Lebesgue density of \(\mu_s\) and \(H\) is the negative Boltzmann entropy.




The extra transport assumptions are positive densities, positive eigenvalues of the Jacobian \(DT\) in an orthonormal eigenbasis, the change-of-variables formula \(\rho_s(T_s(z))=\rho_0(z)/\det\bigl((1-s)I+sDT(z)\bigr)\), and integrable logarithms; no convexity or transport inequality is among them.




If \(U\) is jointly \(\lambda\)-strongly convex, we prove


\[\mathcal F(s)\le(1-s)\mathcal F(0)+s\mathcal F(1)-\frac\lambda2s(1-s)\bigl(\|p_0-p_1\|^2+W_2(\mu_0,\mu_1)^2\bigr),\]


where \(W_2\) is the Wasserstein distance. The entropy part follows from concavity of \(\log\det\); the potential part from joint strong convexity. Consequences are a Talagrand-type quadratic growth at a minimizer and the tangent inequality used for the extended log-Sobolev inequality in the next section.

\subsection{\texorpdfstring{The interpolated Jacobian and concavity of its logarithm}{The interpolated Jacobian and concavity of its logarithm}}


For real numbers \(\lambda_1,\dots,\lambda_d\) and \(s\in\mathbb R\) put


\[J(s)=\prod_{i=1}^d\bigl((1-s)+s\lambda_i\bigr).\]


If a linear map \(A\) has an orthonormal basis of eigenvectors with eigenvalues \(\lambda_i\), then \((1-s)I+sA\) has the same eigenvectors with eigenvalues \((1-s)+s\lambda_i\), so \(J(s)=\det((1-s)I+sA)\).




\textbf{Claim.} Let \(\lambda_i>0\) for all \(i\) and \(s\in[0,1]\). Then



\begin{enumerate}


\item \(J(s)>0\), \(J(0)=1\) and \(J(1)=\prod_i\lambda_i\);

\item \(s\log J(1)\le\log J(s)\).

\end{enumerate}



\textbf{Proof.} (1) If \(s=0\) every factor equals \(1\). If \(s>0\), every factor is the sum of \(1-s\ge0\) and \(s\lambda_i>0\), hence positive. A product of positive numbers is positive. The values at \(s=0\) and \(s=1\) are immediate.




(2) The logarithm is concave on \((0,\infty)\). For the points \(1\) and \(\lambda_i\) and the weights \(1-s\) and \(s\),


\[\log\bigl((1-s)\cdot1+s\lambda_i\bigr)\ge(1-s)\log1+s\log\lambda_i=s\log\lambda_i.\]


Since all factors are positive, the logarithm of each product is the sum of the logarithms, and summing over \(i\) gives \(\log J(s)\ge s\sum_i\log\lambda_i=s\log J(1)\). \(\square\)

\begin{leanbox}{Lean: \texttt{spectralJacobian}, \texttt{spectralJacobian\_\allowbreak{}pos}, \texttt{spectralJacobian\_\allowbreak{}zero} and 2 more}
\leansource{\texttt{spectralJacobian} \quad LatentDiffusion/\allowbreak{}RegularTransport.lean, lines 37–38}
\begin{leancode}
\leanline{0}{def\ spectralJacobian\ \{d\ :\ ℕ\}\ (eigenvalue\ :\ Fin\ d\ →\ ℝ)\ (s\ :\ ℝ)\ :\ ℝ\ :=}
\leanline{2}{∏\ i,\ ((1\ -\ s)\ +\ s\ *\ eigenvalue\ i)}
\end{leancode}
\leansource{\texttt{spectralJacobian\_\allowbreak{}pos} \quad LatentDiffusion/\allowbreak{}RegularTransport.lean, lines 40–48}
\begin{leancode}
\leanline{0}{theorem\ spectralJacobian\_pos\ \{d\ :\ ℕ\}\ (eigenvalue\ :\ Fin\ d\ →\ ℝ)}
\leanline{4}{(hpos\ :\ ∀\ i,\ 0\ <\ eigenvalue\ i)\ (s\ :\ ℝ)\ (hs\ :\ s\ ∈\ Icc\ (0\ :\ ℝ)\ 1)\ :}
\leanline{4}{0\ <\ spectralJacobian\ eigenvalue\ s}
\end{leancode}
\leansource{\texttt{spectralJacobian\_\allowbreak{}zero} \quad LatentDiffusion/\allowbreak{}RegularTransport.lean, lines 50–51}
\begin{leancode}
\leanline{0}{@[simp]\ theorem\ spectralJacobian\_zero\ \{d\ :\ ℕ\}\ (eigenvalue\ :\ Fin\ d\ →\ ℝ)\ :}
\leanline{4}{spectralJacobian\ eigenvalue\ 0\ =\ 1}
\end{leancode}
\leansource{\texttt{spectralJacobian\_\allowbreak{}one} \quad LatentDiffusion/\allowbreak{}RegularTransport.lean, lines 53–54}
\begin{leancode}
\leanline{0}{@[simp]\ theorem\ spectralJacobian\_one\ \{d\ :\ ℕ\}\ (eigenvalue\ :\ Fin\ d\ →\ ℝ)\ :}
\leanline{4}{spectralJacobian\ eigenvalue\ 1\ =\ ∏\ i,\ eigenvalue\ i}
\end{leancode}
\leansource{\texttt{log\_\allowbreak{}spectralJacobian\_\allowbreak{}ge} \quad LatentDiffusion/\allowbreak{}RegularTransport.lean, lines 58–74}
\begin{leancode}
\leanline{0}{theorem\ log\_spectralJacobian\_ge\ \{d\ :\ ℕ\}\ (eigenvalue\ :\ Fin\ d\ →\ ℝ)}
\leanline{4}{(hpos\ :\ ∀\ i,\ 0\ <\ eigenvalue\ i)\ (s\ :\ ℝ)\ (hs\ :\ s\ ∈\ Icc\ (0\ :\ ℝ)\ 1)\ :}
\leanline{4}{s\ *\ Real.log\ (spectralJacobian\ eigenvalue\ 1)\ ≤}
\leanline{6}{Real.log\ (spectralJacobian\ eigenvalue\ s)}
\end{leancode}
\leanlabel{How the Lean proof follows the argument}
\begin{enumerate}


\item \texttt{spectral\allowbreak{}Jacobian eigenvalue s} is the product \texttt{∏ i, ((1 -\allowbreak{} s) + s * eigenvalue i)} over \texttt{Fin d}. \texttt{spectral\allowbreak{}Jacobian\_\allowbreak{}pos} uses \texttt{Finset.\allowbreak{}prod\_\allowbreak{}pos} and the case split \texttt{by\_\allowbreak{}cases h : s = 0} of the proof; \texttt{spectral\allowbreak{}Jacobian\_\allowbreak{}zero} and \texttt{spectral\allowbreak{}Jacobian\_\allowbreak{}one} are \texttt{simp} evaluations.

\item \texttt{log\_\allowbreak{}spectral\allowbreak{}Jacobian\_\allowbreak{}ge}: \texttt{he} is the concavity inequality for each eigenvalue, from \texttt{strict\allowbreak{}Concave\allowbreak{}On\_\allowbreak{}log\_\allowbreak{}Ioi} (strict concavity of \(\log\) on \((0,\infty)\)) weakened to concavity. \texttt{Real.\allowbreak{}log\_\allowbreak{}prod} turns logarithms of products into sums of logarithms, using positivity of the factors from \texttt{spectral\allowbreak{}Jacobian\_\allowbreak{}pos} with \(d=1\); \texttt{Finset.\allowbreak{}mul\_\allowbreak{}sum} and \texttt{Finset.\allowbreak{}sum\_\allowbreak{}le\_\allowbreak{}sum} sum the inequalities.

\end{enumerate}
\end{leanbox}
\subsection{\texorpdfstring{Regular displacement data}{Regular displacement data}}


\textbf{Definition.} Let \(\nu\) and \(\mu\) be measures on \(Z\) (in applications \(\nu\) is Lebesgue measure and \(\mu\) a probability measure). For a map \(T:Z\to Z\) the displacement interpolation is \(T_s(z)=(1-s)z+sT(z)\); thus \(T_0=\mathrm{id}\), \(T_1=T\), and \(T_s\) is measurable when \(T\) is. Regular displacement data consist of



\begin{enumerate}


\item a measurable map \(T\);

\item functions \(\rho_s:Z\to\mathbb R\), measurable and strictly positive for \(s\in[0,1]\), such that the pushforward of \(\mu\) under \(T_s\) has density \(\rho_s\) with respect to \(\nu\): \(\mu\circ T_s^{-1}=\rho_s\,\nu\);

\item for each \(z\), an orthonormal basis \((b_i(z))_{i=1}^d\) of \(Z\) and positive numbers \(\lambda_i(z)\) such that \(T\) is Fréchet differentiable at \(z\) with derivative \(DT(z)\) satisfying \(DT(z)b_i(z)=\lambda_i(z)b_i(z)\);

\item the change-of-variables formula: for \(s\in[0,1]\) and every \(z\),
\[\rho_s(T_s(z))=\frac{\rho_0(z)}{J_z(s)},\qquad J_z(s)=\prod_i\bigl((1-s)+s\lambda_i(z)\bigr);\]

\item for \(s\in[0,1]\), \(\mu\)-integrability of \(z\mapsto\log\rho_s(T_s(z))\).

\end{enumerate}



Condition (4) is the local conservation of mass for the smooth positive-Jacobian map \(T_s\); it is an assumption on the data, not an entropy inequality. Arbitrary pairs of probability laws need not admit such data, which is why they form the explicit hypothesis R4.




\textbf{Claim.} (1) \(\det DT(z)=J_z(1)=\prod_i\lambda_i(z)\) and \(\operatorname{tr}DT(z)=\sum_i\lambda_i(z)\). (2) The law \(\mu_s=\mu\circ T_s^{-1}\) satisfies \(\mu_0=\mu\) and \(\mu_1=\mu\circ T^{-1}\), and it is a probability measure for every \(s\) if \(\mu\) is.




\textbf{Proof.} (1) In the orthonormal basis \((b_i(z))\) the matrix of \(DT(z)\) is diagonal with entries \(\lambda_i(z)\), because \(DT(z)b_j=\lambda_jb_j\). The determinant and the trace are independent of the basis; for a diagonal matrix they are the product and the sum of the diagonal entries. (2) \(T_0\) is the identity and \(T_1=T\). The pushforward of a probability measure by a measurable map has total mass \(\mu(T_s^{-1}(Z))=\mu(Z)=1\). \(\square\)




Part (1) checks that the product in (4) is the Jacobian determinant of the actual map \(T\), and the trace formula is used for the divergence of the displacement \(T-\mathrm{id}\).

\begin{leanbox}{Lean: \texttt{interpolate}, \texttt{interpolate\_\allowbreak{}zero}, \texttt{interpolate\_\allowbreak{}one} and 8 more}
\leansource{\texttt{interpolate} \quad LatentDiffusion/\allowbreak{}RegularTransport.lean, lines 82–82}
\begin{leancode}
\leanline{0}{def\ interpolate\ (T\ :\ Z\ →\ Z)\ (s\ :\ ℝ)\ (z\ :\ Z)\ :\ Z\ :=\ (1\ -\ s)\ •\ z\ +\ s\ •\ T\ z}
\end{leancode}
\leansource{\texttt{interpolate\_\allowbreak{}zero} \quad LatentDiffusion/\allowbreak{}RegularTransport.lean, lines 84–86}
\begin{leancode}
\leanline{0}{@[simp]\ theorem\ interpolate\_zero\ (T\ :\ Z\ →\ Z)\ :\ interpolate\ T\ 0\ =\ id}
\end{leancode}
\leansource{\texttt{interpolate\_\allowbreak{}one} \quad LatentDiffusion/\allowbreak{}RegularTransport.lean, lines 88–90}
\begin{leancode}
\leanline{0}{@[simp]\ theorem\ interpolate\_one\ (T\ :\ Z\ →\ Z)\ :\ interpolate\ T\ 1\ =\ T}
\end{leancode}
\leansource{\texttt{interpolate\_\allowbreak{}measurable} \quad LatentDiffusion/\allowbreak{}RegularTransport.lean, lines 92–93}
\begin{leancode}
\leanline{0}{@[fun\_prop]\ theorem\ interpolate\_measurable\ (T\ :\ Z\ →\ Z)\ (hT\ :\ Measurable\ T)\ (s\ :\ ℝ)\ :}
\leanline{4}{Measurable\ (interpolate\ T\ s)}
\end{leancode}
\leansource{\texttt{RegularDisplacement} \quad LatentDiffusion/\allowbreak{}RegularTransport.lean, lines 100–119}
\begin{leancode}
\leanline{0}{structure\ RegularDisplacement\ (ν\ μ\ :\ Measure\ Z)\ (d\ :\ ℕ)\ where}
\leanline{2}{transport\ :\ Z\ →\ Z}
\leanline{2}{measurable\_transport\ :\ Measurable\ transport}
\leanline{2}{density\ :\ ℝ\ →\ Z\ →\ ℝ}
\leanline{2}{measurable\_density\ :\ ∀\ s\ ∈\ Icc\ (0\ :\ ℝ)\ 1,\ Measurable\ (density\ s)}
\leanline{2}{positive\_density\ :\ ∀\ s\ ∈\ Icc\ (0\ :\ ℝ)\ 1,\ ∀\ z,\ 0\ <\ density\ s\ z}
\leanline{2}{density\_law\ :\ ∀\ s\ ∈\ Icc\ (0\ :\ ℝ)\ 1,}
\leanline{4}{μ.map\ (interpolate\ transport\ s)\ =\ ν.withDensity\ (fun\ z\ ↦\ ENNReal.ofReal\ (density\ s\ z))}
\leanline{2}{eigenvalue\ :\ Z\ →\ Fin\ d\ →\ ℝ}
\leanline{2}{positive\_eigenvalue\ :\ ∀\ z\ i,\ 0\ <\ eigenvalue\ z\ i}
\leanline{2}{jacobian\ :\ Z\ →\ Z\ →L[ℝ]\ Z}
\leanline{2}{derivative\_transport\ :\ ∀\ z,\ HasFDerivAt\ transport\ (jacobian\ z)\ z}
\leanline{2}{spectral\_basis\ :\ Z\ →\ OrthonormalBasis\ (Fin\ d)\ ℝ\ Z}
\leanline{2}{jacobian\_eigenvector\ :\ ∀\ z\ i,}
\leanline{4}{jacobian\ z\ (spectral\_basis\ z\ i)\ =\ eigenvalue\ z\ i\ •\ spectral\_basis\ z\ i}
\leanline{2}{density\_jacobian\ :\ ∀\ s\ ∈\ Icc\ (0\ :\ ℝ)\ 1,\ ∀\ z,}
\leanline{4}{density\ s\ (interpolate\ transport\ s\ z)\ =}
\leanline{6}{density\ 0\ z\ /\ spectralJacobian\ (eigenvalue\ z)\ s}
\leanline{2}{integrable\_log\ :\ ∀\ s\ ∈\ Icc\ (0\ :\ ℝ)\ 1,}
\leanline{4}{Integrable\ (fun\ z\ ↦\ Real.log\ (density\ s\ (interpolate\ transport\ s\ z)))\ μ}
\end{leancode}
\leansource{\texttt{det\_\allowbreak{}jacobian} \quad LatentDiffusion/\allowbreak{}RegularTransport.lean, lines 127–141}
\begin{leancode}
\leanline{0}{theorem\ det\_jacobian\ (z\ :\ Z)\ :}
\leanline{4}{(D.jacobian\ z).toLinearMap.det\ =\ spectralJacobian\ (D.eigenvalue\ z)\ 1}
\end{leancode}
\leansource{\texttt{trace\_\allowbreak{}jacobian} \quad LatentDiffusion/\allowbreak{}RegularTransport.lean, lines 145–152}
\begin{leancode}
\leanline{0}{theorem\ trace\_jacobian\ (z\ :\ Z)\ :}
\leanline{4}{LinearMap.trace\ ℝ\ Z\ (D.jacobian\ z).toLinearMap\ =\ ∑\ i,\ D.eigenvalue\ z\ i}
\end{leancode}
\leansource{\texttt{law} \quad LatentDiffusion/\allowbreak{}RegularTransport.lean, lines 155–155}
\begin{leancode}
\leanline{0}{def\ law\ (s\ :\ ℝ)\ :\ Measure\ Z\ :=\ μ.map\ (interpolate\ D.transport\ s)}
\end{leancode}
\leansource{\texttt{law\_\allowbreak{}zero} \quad LatentDiffusion/\allowbreak{}RegularTransport.lean, lines 157–157}
\begin{leancode}
\leanline{0}{@[simp]\ theorem\ law\_zero\ :\ D.law\ 0\ =\ μ}
\end{leancode}
\leansource{\texttt{law\_\allowbreak{}one} \quad LatentDiffusion/\allowbreak{}RegularTransport.lean, lines 159–159}
\begin{leancode}
\leanline{0}{@[simp]\ theorem\ law\_one\ :\ D.law\ 1\ =\ μ.map\ D.transport}
\end{leancode}
\leansource{\texttt{law\_\allowbreak{}probability} \quad LatentDiffusion/\allowbreak{}RegularTransport.lean, lines 161–164}
\begin{leancode}
\leanline{0}{instance\ law\_probability\ [IsProbabilityMeasure\ μ]\ (s\ :\ ℝ)\ :}
\leanline{4}{IsProbabilityMeasure\ (D.law\ s)\ :=}
\leanline{2}{(Measure.isProbabilityMeasure\_map\_iff}
\leanline{4}{(interpolate\_measurable\ \_\ D.measurable\_transport\ s).aemeasurable).mpr\ inferInstance}
\end{leancode}
\leanlabel{How the Lean proof follows the argument}
\begin{enumerate}


\item \texttt{interpolate T s z} is \texttt{(1 -\allowbreak{} s) \ensuremath{\bullet} z + s \ensuremath{\bullet} T z}; \texttt{interpolate\_\allowbreak{}zero}, \texttt{interpolate\_\allowbreak{}one} and \texttt{interpolate\_\allowbreak{}measurable} are the three facts about \(T_s\).

\item \texttt{Regular\allowbreak{}Displacement \ensuremath{\nu} \ensuremath{\mu} d} is a structure with fields \texttt{transport}, \texttt{measurable\_\allowbreak{}transport}, \texttt{density}, \texttt{measurable\_\allowbreak{}density}, \texttt{positive\_\allowbreak{}density}, \texttt{density\_\allowbreak{}law} (condition 2, with \texttt{\ensuremath{\mu}.\allowbreak{}map} the pushforward and \texttt{\ensuremath{\nu}.\allowbreak{}with\allowbreak{}Density} the measure with density), \texttt{eigenvalue}, \texttt{positive\_\allowbreak{}eigenvalue}, \texttt{jacobian}, \texttt{derivative\_\allowbreak{}transport} (Fréchet differentiability, \texttt{Has\allowbreak{}FDeriv\allowbreak{}At}), \texttt{spectral\_\allowbreak{}basis} (an \texttt{Orthonormal\allowbreak{}Basis (Fin d) \ensuremath{\mathbb R} Z}), \texttt{jacobian\_\allowbreak{}eigenvector}, \texttt{density\_\allowbreak{}jacobian} (condition 4) and \texttt{integrable\_\allowbreak{}log} (condition 5). The quantifiers \texttt{\ensuremath{\forall} s \ensuremath{\in} Icc (0 : \ensuremath{\mathbb R}) 1} restrict to \(s\in[0,1]\).

\item \texttt{det\_\allowbreak{}jacobian} computes \texttt{Linear\allowbreak{}Map.\allowbreak{}det} through the matrix in the spectral basis (\texttt{Linear\allowbreak{}Map.\allowbreak{}det\_\allowbreak{}to\allowbreak{}Matrix}), shows the matrix is \texttt{Matrix.\allowbreak{}diagonal} of the eigenvalues, and applies \texttt{Matrix.\allowbreak{}det\_\allowbreak{}diagonal}. \texttt{trace\_\allowbreak{}jacobian} uses \texttt{Linear\allowbreak{}Map.\allowbreak{}trace\_\allowbreak{}eq\_\allowbreak{}matrix\_\allowbreak{}trace} in the same basis.

\item \texttt{law s} is \texttt{\ensuremath{\mu}.\allowbreak{}map (interpolate D.\allowbreak{}transport s)}; \texttt{law\_\allowbreak{}zero} and \texttt{law\_\allowbreak{}one} follow by \texttt{simp}, and the instance \texttt{law\_\allowbreak{}probability} is \texttt{Measure.\allowbreak{}is\allowbreak{}Probability\allowbreak{}Measure\_\allowbreak{}map\_\allowbreak{}iff}.

\end{enumerate}
\end{leanbox}
\subsection{\texorpdfstring{Displacement convexity of the entropy}{Displacement convexity of the entropy}}


The entropy at time \(s\) is \(H(s)=\int\log\rho_s\,d\mu_s\). Since \(\mu_s\) has density \(\rho_s\), this is \(\int\rho_s\log\rho_s\,d\nu\), the negative Boltzmann entropy of \(\mu_s\).




\textbf{Claim.} For regular displacement data and \(s\in[0,1]\),



\begin{enumerate}


\item for every \(z\), \(\log\rho_s(T_s(z))\le(1-s)\log\rho_0(z)+s\log\rho_1(T(z))\);

\item \(H(s)=\int\log\rho_s(T_s(z))\,\mu(dz)\);

\item \(H(s)\le(1-s)H(0)+sH(1)\).

\end{enumerate}



\textbf{Proof.} (1) By the change-of-variables formula at times \(s\) and \(1\) (with \(T_1=T\)),


\[\log\rho_s(T_s(z))=\log\rho_0(z)-\log J_z(s),\qquad\log\rho_1(T(z))=\log\rho_0(z)-\log J_z(1),\]


where the logarithms of quotients split because \(\rho_0(z)>0\) and \(J_z(s),J_z(1)>0\) (first step). The first step gives \(\log J_z(s)\ge s\log J_z(1)\), hence


\[\log\rho_s(T_s(z))\le\log\rho_0(z)-s\log J_z(1)=(1-s)\log\rho_0(z)+s\bigl(\log\rho_0(z)-\log J_z(1)\bigr),\]


which is (1).




(2) For a measurable map \(\Phi\) and a measurable function \(f\), \(\int f\,d(\mu\circ\Phi^{-1})=\int f\circ\Phi\,d\mu\) (change of variables for pushforward measures). Apply it with \(\Phi=T_s\) and \(f=\log\rho_s\), which is measurable because \(\rho_s\) is.




(3) By (2), \(H(s)\), \(H(0)\) and \(H(1)\) are the \(\mu\)-integrals of \(\log\rho_s\circ T_s\), \(\log\rho_0\) and \(\log\rho_1\circ T\); these three functions are \(\mu\)-integrable by condition (5) at \(s\), \(0\) and \(1\). Integrating the pointwise inequality (1) against \(\mu\) and using linearity and monotonicity of the integral gives (3). \(\square\)




No convexity hypothesis on the laws is used: displacement convexity of the entropy is a consequence of the change-of-variables formula and the concavity of \(\log\det\) on positive definite spectra.

\begin{leanbox}{Lean: \texttt{log\_\allowbreak{}density\_\allowbreak{}interpolate\_\allowbreak{}le}, \texttt{entropy}, \texttt{entropy\_\allowbreak{}eq\_\allowbreak{}pullback} and 1 more}
\leansource{\texttt{log\_\allowbreak{}density\_\allowbreak{}interpolate\_\allowbreak{}le} \quad LatentDiffusion/\allowbreak{}RegularTransport.lean, lines 168–181}
\begin{leancode}
\leanline{0}{theorem\ log\_density\_interpolate\_le\ (s\ :\ ℝ)\ (hs\ :\ s\ ∈\ Icc\ (0\ :\ ℝ)\ 1)\ (z\ :\ Z)\ :}
\leanline{4}{Real.log\ (D.density\ s\ (interpolate\ D.transport\ s\ z))\ ≤}
\leanline{6}{(1\ -\ s)\ *\ Real.log\ (D.density\ 0\ z)\ +}
\leanline{8}{s\ *\ Real.log\ (D.density\ 1\ (D.transport\ z))}
\end{leancode}
\leansource{\texttt{entropy} \quad LatentDiffusion/\allowbreak{}RegularTransport.lean, lines 184–184}
\begin{leancode}
\leanline{0}{def\ entropy\ (s\ :\ ℝ)\ :\ ℝ\ :=\ ∫\ z,\ Real.log\ (D.density\ s\ z)\ ∂D.law\ s}
\end{leancode}
\leansource{\texttt{entropy\_\allowbreak{}eq\_\allowbreak{}pullback} \quad LatentDiffusion/\allowbreak{}RegularTransport.lean, lines 186–189}
\begin{leancode}
\leanline{0}{theorem\ entropy\_eq\_pullback\ (s\ :\ ℝ)\ (hs\ :\ s\ ∈\ Icc\ (0\ :\ ℝ)\ 1)\ :}
\leanline{4}{D.entropy\ s\ =\ ∫\ z,\ Real.log\ (D.density\ s\ (interpolate\ D.transport\ s\ z))\ ∂μ}
\end{leancode}
\leansource{\texttt{entropy\_\allowbreak{}displacement\_\allowbreak{}convex} \quad LatentDiffusion/\allowbreak{}RegularTransport.lean, lines 193–204}
\begin{leancode}
\leanline{0}{theorem\ entropy\_displacement\_convex\ (s\ :\ ℝ)\ (hs\ :\ s\ ∈\ Icc\ (0\ :\ ℝ)\ 1)\ :}
\leanline{4}{D.entropy\ s\ ≤\ (1\ -\ s)\ *\ D.entropy\ 0\ +\ s\ *\ D.entropy\ 1}
\end{leancode}
\leanlabel{How the Lean proof follows the argument}
\begin{enumerate}


\item \texttt{log\_\allowbreak{}density\_\allowbreak{}interpolate\_\allowbreak{}le}: \texttt{hmass} and \texttt{hmass1} are \texttt{D.\allowbreak{}density\_\allowbreak{}jacobian} at \(s\) and \(1\); \texttt{Real.\allowbreak{}log\_\allowbreak{}div} splits the logarithms, and \texttt{nlinarith} combines the result with \texttt{log\_\allowbreak{}spectral\allowbreak{}Jacobian\_\allowbreak{}ge}.

\item \texttt{entropy s} is \texttt{\ensuremath{\int} z, Real.\allowbreak{}log (D.\allowbreak{}density s z) \ensuremath{\partial}D.\allowbreak{}law s}, and \texttt{entropy\_\allowbreak{}eq\_\allowbreak{}pullback} is Mathlib's \texttt{integral\_\allowbreak{}map} for the pushforward.

\item \texttt{entropy\_\allowbreak{}displacement\_\allowbreak{}convex} rewrites the three entropies by \texttt{entropy\_\allowbreak{}eq\_\allowbreak{}pullback}, combines the right side into one integral with \texttt{integral\_\allowbreak{}const\_\allowbreak{}mul} and \texttt{integral\_\allowbreak{}add} (the integrability inputs are \texttt{D.\allowbreak{}integrable\_\allowbreak{}log} at \(0\) and \(1\)), and concludes with \texttt{integral\_\allowbreak{}mono} applied to the pointwise bound.

\end{enumerate}
\end{leanbox}
\subsection{\texorpdfstring{Strong convexity of the potential energy along the displacement}{Strong convexity of the potential energy along the displacement}}


Let \(P\) be a complete real inner product space of parameters, and give \(P\times Z\) the Euclidean product norm \(\|(p,z)\|^2=\|p\|^2+\|z\|^2\). For \(p_0,p_1\in P\) put \(p_s=(1-s)p_0+sp_1\), and for \(U:P\times Z\to\mathbb R\) define the free energy of one datum along the interpolation by \(\mathcal F(s)=H(s)+\int U(p_s,\cdot)\,d\mu_s\). For a square-integrable vector function \(X\) write \(\|X\|_{L^2(\mu)}=(\int\|X\|^2d\mu)^{1/2}\); the identity \(\|X\|_{L^2(\mu)}^2=\int\|X\|^2\,d\mu\) holds by definition of this norm.




\textbf{Claim.} Let \(\mu\) be a probability measure, \(\lambda\in\mathbb R\), and let \(U\) be \(\lambda\)-strongly convex on \(P\times Z\). Assume that \(U(p,\cdot)\) is measurable for every \(p\), that \(z\mapsto U(p_s,T_s(z))\) is \(\mu\)-integrable for every \(s\in[0,1]\), and that \(z\mapsto z-T(z)\) is square integrable under \(\mu\). Then for \(s\in[0,1]\),


\[\begin{aligned}\int U(p_s,\cdot)\,d\mu_s\le{}&(1-s)\int U(p_0,z)\,\mu(dz)+s\int U(p_1,T(z))\,\mu(dz)\\&-\frac\lambda2s(1-s)\bigl(\|p_0-p_1\|^2+\|\mathrm{id}-T\|_{L^2(\mu)}^2\bigr).\end{aligned}\]


\textbf{Proof.} Fix \(z\). The points \((p_0,z)\) and \((p_1,T(z))\) have the convex combination \((1-s)(p_0,z)+s(p_1,T(z))=(p_s,T_s(z))\) and squared distance \(\|p_0-p_1\|^2+\|z-T(z)\|^2\). Strong convexity with the weights \(1-s\) and \(s\) gives


\[U(p_s,T_s(z))\le(1-s)U(p_0,z)+sU(p_1,T(z))-\frac\lambda2s(1-s)\bigl(\|p_0-p_1\|^2+\|z-T(z)\|^2\bigr).\]


All functions of \(z\) on both sides are \(\mu\)-integrable: the potential terms by hypothesis at \(s\), \(0\) and \(1\), the constant because \(\mu\) is finite, and \(\|z-T(z)\|^2\) by square integrability. Integrating against \(\mu\) and using \(\mu(Z)=1\) and \(\int\|z-T(z)\|^2d\mu=\|\mathrm{id}-T\|_{L^2(\mu)}^2\) gives the claim, after writing \(\int U(p_s,\cdot)\,d\mu_s=\int U(p_s,T_s(z))\,\mu(dz)\) by the pushforward formula. \(\square\)

\begin{leanbox}{Lean: \texttt{parameterInterpolate}, \texttt{freeEnergy}, \texttt{secondCost\_\allowbreak{}sq\_\allowbreak{}eq\_\allowbreak{}integral} and 1 more}
\leansource{\texttt{parameterInterpolate} \quad LatentDiffusion/\allowbreak{}RegularTransport.lean, lines 220–220}
\begin{leancode}
\leanline{0}{def\ parameterInterpolate\ (p₀\ p₁\ :\ P)\ (s\ :\ ℝ)\ :\ P\ :=\ (1\ -\ s)\ •\ p₀\ +\ s\ •\ p₁}
\end{leancode}
\leansource{\texttt{freeEnergy} \quad LatentDiffusion/\allowbreak{}RegularTransport.lean, lines 223–225}
\begin{leancode}
\leanline{0}{def\ freeEnergy\ (D\ :\ RegularDisplacement\ ν\ μ\ d)\ (U\ :\ P\ →\ Z\ →\ ℝ)}
\leanline{4}{(p₀\ p₁\ :\ P)\ (s\ :\ ℝ)\ :\ ℝ\ :=}
\leanline{2}{D.entropy\ s\ +\ ∫\ z,\ U\ (parameterInterpolate\ p₀\ p₁\ s)\ z\ ∂D.law\ s}
\end{leancode}
\leansource{\texttt{secondCost\_\allowbreak{}sq\_\allowbreak{}eq\_\allowbreak{}integral} \quad LatentDiffusion/\allowbreak{}RegularTransport.lean, lines 228–234}
\begin{leancode}
\leanline{0}{theorem\ secondCost\_sq\_eq\_integral\ \{Ω\ H\ :\ Type*\}\ [MeasurableSpace\ Ω]}
\leanline{4}{[NormedAddCommGroup\ H]\ [InnerProductSpace\ ℝ\ H]}
\leanline{4}{(Q\ :\ Measure\ Ω)\ (X\ :\ Ω\ →\ H)\ (hX\ :\ MemLp\ X\ 2\ Q)\ :}
\leanline{4}{(eLpNorm\ X\ 2\ Q).toReal\ \textasciicircum{}\ 2\ =\ ∫\ ω,\ ‖X\ ω‖\ \textasciicircum{}\ 2\ ∂Q}
\end{leancode}
\leansource{\texttt{potential\_\allowbreak{}displacement\_\allowbreak{}strongConvex} \quad LatentDiffusion/\allowbreak{}RegularTransport.lean, lines 238–294}
\begin{leancode}
\leanline{0}{theorem\ potential\_displacement\_strongConvex}
\leanline{4}{(D\ :\ RegularDisplacement\ ν\ μ\ d)\ (U\ :\ P\ →\ Z\ →\ ℝ)\ (p₀\ p₁\ :\ P)\ (rho\ :\ ℝ)}
\leanline{4}{(hU\ :\ StrongConvexOn\ univ\ rho\ (fun\ w\ :\ PBGSO.ProductTransport.Joint\ P\ Z\ ↦\ U\ w.fst\ w.snd))}
\leanline{4}{(hmeas\ :\ ∀\ p,\ Measurable\ (U\ p))}
\leanline{4}{(hint\ :\ ∀\ s\ ∈\ Icc\ (0\ :\ ℝ)\ 1,}
\leanline{6}{Integrable\ (fun\ z\ ↦\ U\ (parameterInterpolate\ p₀\ p₁\ s)\ (interpolate\ D.transport\ s\ z))\ μ)}
\leanline{4}{(hcost\ :\ MemLp\ (fun\ z\ ↦\ z\ -\ D.transport\ z)\ 2\ μ)}
\leanline{4}{(s\ :\ ℝ)\ (hs\ :\ s\ ∈\ Icc\ (0\ :\ ℝ)\ 1)\ :}
\leanline{4}{(∫\ z,\ U\ (parameterInterpolate\ p₀\ p₁\ s)\ z\ ∂D.law\ s)\ ≤}
\leanline{6}{(1\ -\ s)\ *\ (∫\ z,\ U\ p₀\ z\ ∂μ)\ +\ s\ *\ (∫\ z,\ U\ p₁\ (D.transport\ z)\ ∂μ)\ -}
\leanline{8}{rho\ /\ 2\ *\ s\ *\ (1\ -\ s)\ *}
\leanline{10}{(‖p₀\ -\ p₁‖\ \textasciicircum{}\ 2\ +\ (eLpNorm\ (fun\ z\ ↦\ z\ -\ D.transport\ z)\ 2\ μ).toReal\ \textasciicircum{}\ 2)}
\end{leancode}
\leanlabel{How the Lean proof follows the argument}
\begin{enumerate}


\item \texttt{parameter\allowbreak{}Interpolate p\ensuremath{{}_{0}} p\ensuremath{{}_{1}} s} is \(p_s\), and \texttt{free\allowbreak{}Energy D U p\ensuremath{{}_{0}} p\ensuremath{{}_{1}} s} is \texttt{D.\allowbreak{}entropy s + \ensuremath{\int} z, U (parameter\allowbreak{}Interpolate p\ensuremath{{}_{0}} p\ensuremath{{}_{1}} s) z \ensuremath{\partial}D.\allowbreak{}law s}.

\item \texttt{second\allowbreak{}Cost\_\allowbreak{}sq\_\allowbreak{}eq\_\allowbreak{}integral} states \texttt{(e\allowbreak{}Lp\allowbreak{}Norm X 2 Q).\allowbreak{}to\allowbreak{}Real \textasciicircum{} 2 = \ensuremath{\int} \ensuremath{\omega}, \ensuremath{\Vert}X \ensuremath{\omega}\ensuremath{\Vert} \textasciicircum{} 2 \ensuremath{\partial}Q} for \texttt{Mem\allowbreak{}Lp X 2 Q}; \texttt{e\allowbreak{}Lp\allowbreak{}Norm X 2 Q} is the \(L^2\) norm with values in \([0,\infty]\), and square integrability makes its real conversion meaningful. It is derived from the corresponding identity for elements of the space \texttt{Lp}.

\item \texttt{potential\_\allowbreak{}displacement\_\allowbreak{}strong\allowbreak{}Convex}: the joint space is \texttt{PBGSO.\allowbreak{}Product\allowbreak{}Transport.\allowbreak{}Joint P Z}, the product with the Euclidean norm, and \texttt{PBGSO.\allowbreak{}Product\allowbreak{}Transport.\allowbreak{}embed p z} is the pair. \texttt{hpoint} is the pointwise inequality; it unfolds the strong convexity of \texttt{h\allowbreak{}U} at the two pairs and rewrites the norm with \texttt{With\allowbreak{}Lp.\allowbreak{}prod\_\allowbreak{}norm\_\allowbreak{}sq\_\allowbreak{}eq\_\allowbreak{}of\_\allowbreak{}L2}. \texttt{hbound := integral\_\allowbreak{}mono (hint s hs) hrhs hpoint} integrates it, \texttt{hmap} is the pushforward formula \texttt{integral\_\allowbreak{}map}, and the final rewriting evaluates the integral of the right side by linearity and \texttt{second\allowbreak{}Cost\_\allowbreak{}sq\_\allowbreak{}eq\_\allowbreak{}integral}.

\end{enumerate}
\end{leanbox}
\subsection{\texorpdfstring{Strong displacement convexity of the free energy}{Strong displacement convexity of the free energy}}


\textbf{Wasserstein distance.} For probability measures \(\mu,\eta\) on \(Z\), a coupling is a probability measure \(\gamma\) on \(Z\times Z\) with marginals \(\mu\) and \(\eta\), and


\[W_2(\mu,\eta)=\inf_\gamma\Bigl(\int\|x-y\|^2\,\gamma(dx,dy)\Bigr)^{1/2}\in[0,\infty].\]


\textbf{Claim.} Under the hypotheses of the previous step, for \(s\in[0,1]\):



\begin{enumerate}


\item \(\mathcal F(s)\le(1-s)\mathcal F(0)+s\mathcal F(1)-\frac\lambda2s(1-s)\bigl(\|p_0-p_1\|^2+\|\mathrm{id}-T\|_{L^2(\mu)}^2\bigr)\);

\item \(W_2(\mu,\mu_1)^2\le\|\mathrm{id}-T\|_{L^2(\mu)}^2\), and in particular \(W_2(\mu,\mu_1)<\infty\);

\item if moreover \(\lambda\ge0\), then \(\mathcal F(s)\le(1-s)\mathcal F(0)+s\mathcal F(1)-\frac\lambda2s(1-s)\bigl(\|p_0-p_1\|^2+W_2(\mu,\mu_1)^2\bigr)\).

\end{enumerate}



\textbf{Proof.} (1) Add the entropy inequality of the third step to the potential inequality of the previous step. At the endpoints, \(\mathcal F(0)=H(0)+\int U(p_0,\cdot)\,d\mu\) because \(\mu_0=\mu\), and \(\mathcal F(1)=H(1)+\int U(p_1,T(z))\,\mu(dz)\) by the pushforward formula for \(\mu_1=\mu\circ T^{-1}\).




(2) The law of the pair \((z,T(z))\) under \(\mu\) is a coupling of \(\mu\) and \(\mu_1\), and its cost is \(\int\|z-T(z)\|^2\,\mu(dz)\). The infimum over couplings is at most this value, which is finite by square integrability.




(3) When \(\lambda\ge0\) and \(s\in[0,1]\), the coefficient \(c=\frac\lambda2s(1-s)\) is nonnegative. By (2), \(c\bigl(\|p_0-p_1\|^2+W_2(\mu,\mu_1)^2\bigr)\le c\bigl(\|p_0-p_1\|^2+\|\mathrm{id}-T\|_{L^2(\mu)}^2\bigr)\), so the right side of (1) is at most the right side of (3), and (3) follows from (1). \(\square\)




The map \(T\) need not be optimal: the conclusion uses the genuine infimum over couplings, which can only be smaller than the cost of \(T\). With the normalized product distance of the paper, (3) is the strong geodesic convexity of the free energy of one datum.

\begin{leanbox}{Lean: \texttt{freeEnergy\_\allowbreak{}displacement\_\allowbreak{}strongConvex\_\allowbreak{}cost}, \texttt{wasserstein\_\allowbreak{}sq\_\allowbreak{}le\_\allowbreak{}displacement\_\allowbreak{}cost}, \texttt{freeEnergy\_\allowbreak{}displacement\_\allowbreak{}strongConvex}}
\leansource{\texttt{freeEnergy\_\allowbreak{}displacement\_\allowbreak{}strongConvex\_\allowbreak{}cost} \quad LatentDiffusion/\allowbreak{}RegularTransport.lean, lines 298–318}
\begin{leancode}
\leanline{0}{theorem\ freeEnergy\_displacement\_strongConvex\_cost}
\leanline{4}{(D\ :\ RegularDisplacement\ ν\ μ\ d)\ (U\ :\ P\ →\ Z\ →\ ℝ)\ (p₀\ p₁\ :\ P)\ (rho\ :\ ℝ)}
\leanline{4}{(hU\ :\ StrongConvexOn\ univ\ rho\ (fun\ w\ :\ PBGSO.ProductTransport.Joint\ P\ Z\ ↦\ U\ w.fst\ w.snd))}
\leanline{4}{(hmeas\ :\ ∀\ p,\ Measurable\ (U\ p))}
\leanline{4}{(hint\ :\ ∀\ s\ ∈\ Icc\ (0\ :\ ℝ)\ 1,}
\leanline{6}{Integrable\ (fun\ z\ ↦\ U\ (parameterInterpolate\ p₀\ p₁\ s)\ (interpolate\ D.transport\ s\ z))\ μ)}
\leanline{4}{(hcost\ :\ MemLp\ (fun\ z\ ↦\ z\ -\ D.transport\ z)\ 2\ μ)}
\leanline{4}{(s\ :\ ℝ)\ (hs\ :\ s\ ∈\ Icc\ (0\ :\ ℝ)\ 1)\ :}
\leanline{4}{freeEnergy\ D\ U\ p₀\ p₁\ s\ ≤}
\leanline{6}{(1\ -\ s)\ *\ freeEnergy\ D\ U\ p₀\ p₁\ 0\ +\ s\ *\ freeEnergy\ D\ U\ p₀\ p₁\ 1\ -}
\leanline{8}{rho\ /\ 2\ *\ s\ *\ (1\ -\ s)\ *}
\leanline{10}{(‖p₀\ -\ p₁‖\ \textasciicircum{}\ 2\ +\ (eLpNorm\ (fun\ z\ ↦\ z\ -\ D.transport\ z)\ 2\ μ).toReal\ \textasciicircum{}\ 2)}
\end{leancode}
\leansource{\texttt{wasserstein\_\allowbreak{}sq\_\allowbreak{}le\_\allowbreak{}displacement\_\allowbreak{}cost} \quad LatentDiffusion/\allowbreak{}RegularTransport.lean, lines 321–328}
\begin{leancode}
\leanline{0}{theorem\ wasserstein\_sq\_le\_displacement\_cost}
\leanline{4}{(D\ :\ RegularDisplacement\ ν\ μ\ d)\ (hcost\ :\ MemLp\ (fun\ z\ ↦\ z\ -\ D.transport\ z)\ 2\ μ)\ :}
\leanline{4}{ULA.wasserstein2Real\ μ\ (D.law\ 1)\ \textasciicircum{}\ 2\ ≤}
\leanline{6}{(eLpNorm\ (fun\ z\ ↦\ z\ -\ D.transport\ z)\ 2\ μ).toReal\ \textasciicircum{}\ 2}
\end{leancode}
\leansource{\texttt{freeEnergy\_\allowbreak{}displacement\_\allowbreak{}strongConvex} \quad LatentDiffusion/\allowbreak{}RegularTransport.lean, lines 333–349}
\begin{leancode}
\leanline{0}{theorem\ freeEnergy\_displacement\_strongConvex}
\leanline{4}{(D\ :\ RegularDisplacement\ ν\ μ\ d)\ (U\ :\ P\ →\ Z\ →\ ℝ)\ (p₀\ p₁\ :\ P)\ (rho\ :\ ℝ)}
\leanline{4}{(hrho\ :\ 0\ ≤\ rho)}
\leanline{4}{(hU\ :\ StrongConvexOn\ univ\ rho\ (fun\ w\ :\ PBGSO.ProductTransport.Joint\ P\ Z\ ↦\ U\ w.fst\ w.snd))}
\leanline{4}{(hmeas\ :\ ∀\ p,\ Measurable\ (U\ p))}
\leanline{4}{(hint\ :\ ∀\ s\ ∈\ Icc\ (0\ :\ ℝ)\ 1,}
\leanline{6}{Integrable\ (fun\ z\ ↦\ U\ (parameterInterpolate\ p₀\ p₁\ s)\ (interpolate\ D.transport\ s\ z))\ μ)}
\leanline{4}{(hcost\ :\ MemLp\ (fun\ z\ ↦\ z\ -\ D.transport\ z)\ 2\ μ)}
\leanline{4}{(s\ :\ ℝ)\ (hs\ :\ s\ ∈\ Icc\ (0\ :\ ℝ)\ 1)\ :}
\leanline{4}{freeEnergy\ D\ U\ p₀\ p₁\ s\ ≤}
\leanline{6}{(1\ -\ s)\ *\ freeEnergy\ D\ U\ p₀\ p₁\ 0\ +\ s\ *\ freeEnergy\ D\ U\ p₀\ p₁\ 1\ -}
\leanline{8}{rho\ /\ 2\ *\ s\ *\ (1\ -\ s)\ *}
\leanline{10}{(‖p₀\ -\ p₁‖\ \textasciicircum{}\ 2\ +\ ULA.wasserstein2Real\ μ\ (D.law\ 1)\ \textasciicircum{}\ 2)}
\end{leancode}
\leanlabel{How the Lean proof follows the argument}
\begin{enumerate}


\item \texttt{free\allowbreak{}Energy\_\allowbreak{}displacement\_\allowbreak{}strong\allowbreak{}Convex\_\allowbreak{}cost} adds \texttt{D.\allowbreak{}entropy\_\allowbreak{}displacement\_\allowbreak{}convex} and \texttt{potential\_\allowbreak{}displacement\_\allowbreak{}strong\allowbreak{}Convex}; \texttt{hmap} is the endpoint pushforward formula, and \texttt{linarith} combines the inequalities after \texttt{simp only} evaluates the interpolations at \(0\) and \(1\).

\item \texttt{wasserstein\_\allowbreak{}sq\_\allowbreak{}le\_\allowbreak{}displacement\_\allowbreak{}cost} applies \texttt{ULA.\allowbreak{}wasserstein2\_\allowbreak{}le\_\allowbreak{}random\_\allowbreak{}cost\_\allowbreak{}real} to the random variables \(\mathrm{id}\) and \(T\) on \((Z,\mu)\): the real Wasserstein distance between their laws is at most the \(L^2\) norm of their difference. Here \texttt{ULA.\allowbreak{}wasserstein2\allowbreak{}Real} is the real conversion of the coupling infimum \texttt{ULA.\allowbreak{}wasserstein2}. \texttt{sq\_\allowbreak{}le\_\allowbreak{}sq\ensuremath{{}_{0}}} squares the inequality.

\item \texttt{free\allowbreak{}Energy\_\allowbreak{}displacement\_\allowbreak{}strong\allowbreak{}Convex} takes \texttt{hrho : 0 \ensuremath{\le} rho}, chains the first result with the second, and uses \texttt{hc}, nonnegativity of \(\frac\lambda2s(1-s)\), in \texttt{mul\_\allowbreak{}le\_\allowbreak{}mul\_\allowbreak{}of\_\allowbreak{}nonneg\_\allowbreak{}left} and \texttt{sub\_\allowbreak{}le\_\allowbreak{}sub\_\allowbreak{}left}.

\end{enumerate}
\end{leanbox}
\subsection{\texorpdfstring{Quadratic growth along the displacement from a minimizer}{Quadratic growth along the displacement from a minimizer}}


\textbf{Claim.} Assume the hypotheses of the previous step with \(\lambda\ge0\), and assume that the starting point minimizes the free energy along the interpolation: \(\mathcal F(0)\le\mathcal F(s)\) for every \(s\in[0,1]\). Then


\[\frac\lambda2\bigl(\|p_0-p_1\|^2+W_2(\mu,\mu_1)^2\bigr)\le\mathcal F(1)-\mathcal F(0).\]


\textbf{Proof.} We use the following elementary fact. Let \(f:[0,1]\to\mathbb R\), \(f_1,\lambda,D\in\mathbb R\), with \(f(0)\le f(s)\) and \(f(s)\le(1-s)f(0)+sf_1-\frac\lambda2s(1-s)D^2\) for \(s\in[0,1]\). Then \(\frac\lambda2D^2\le f_1-f(0)\). Indeed, for \(0<s<1\) the two inequalities give \(s\cdot\frac\lambda2(1-s)D^2\le s(f_1-f(0))\); dividing by \(s\) and letting \(s\downarrow0\) gives the bound.




Apply the fact to \(f=\mathcal F\), \(f_1=\mathcal F(1)\) and \(D=\bigl(\|p_0-p_1\|^2+W_2(\mu,\mu_1)^2\bigr)^{1/2}\); the convexity hypothesis is part (3) of the previous step, since \(D^2\) is the bracket there. \(\square\)

\begin{leanbox}{Lean: \texttt{talagrand\_\allowbreak{}product}}
\leansource{\texttt{talagrand\_\allowbreak{}product} \quad LatentDiffusion/\allowbreak{}RegularTransport.lean, lines 353–370}
\begin{leancode}
\leanline{0}{theorem\ talagrand\_product}
\leanline{4}{(D\ :\ RegularDisplacement\ ν\ μ\ d)\ (U\ :\ P\ →\ Z\ →\ ℝ)\ (p₀\ p₁\ :\ P)\ (rho\ :\ ℝ)}
\leanline{4}{(hrho\ :\ 0\ ≤\ rho)}
\leanline{4}{(hU\ :\ StrongConvexOn\ univ\ rho\ (fun\ w\ :\ PBGSO.ProductTransport.Joint\ P\ Z\ ↦\ U\ w.fst\ w.snd))}
\leanline{4}{(hmeas\ :\ ∀\ p,\ Measurable\ (U\ p))}
\leanline{4}{(hint\ :\ ∀\ s\ ∈\ Icc\ (0\ :\ ℝ)\ 1,}
\leanline{6}{Integrable\ (fun\ z\ ↦\ U\ (parameterInterpolate\ p₀\ p₁\ s)\ (interpolate\ D.transport\ s\ z))\ μ)}
\leanline{4}{(hcost\ :\ MemLp\ (fun\ z\ ↦\ z\ -\ D.transport\ z)\ 2\ μ)}
\leanline{4}{(hmin\ :\ ∀\ s\ ∈\ Icc\ (0\ :\ ℝ)\ 1,}
\leanline{6}{freeEnergy\ D\ U\ p₀\ p₁\ 0\ ≤\ freeEnergy\ D\ U\ p₀\ p₁\ s)\ :}
\leanline{4}{rho\ /\ 2\ *\ (‖p₀\ -\ p₁‖\ \textasciicircum{}\ 2\ +\ ULA.wasserstein2Real\ μ\ (D.law\ 1)\ \textasciicircum{}\ 2)\ ≤}
\leanline{6}{freeEnergy\ D\ U\ p₀\ p₁\ 1\ -\ freeEnergy\ D\ U\ p₀\ p₁\ 0}
\end{leancode}
\leanlabel{How the Lean proof follows the argument}
\texttt{talagrand\_\allowbreak{}product} applies \texttt{Source\allowbreak{}Results.\allowbreak{}geodesic\_\allowbreak{}quadratic\_\allowbreak{}growth}, the elementary fact, to the function \texttt{free\allowbreak{}Energy D U p\ensuremath{{}_{0}} p\ensuremath{{}_{1}}} with distance the square root of the bracket; \texttt{Real.\allowbreak{}sq\_\allowbreak{}sqrt} converts the square of that root back to the bracket, both in the convexity hypothesis and in the conclusion.
\end{leanbox}
\subsection{\texorpdfstring{The tangent inequality and the Cauchy–Schwarz bound}{The tangent inequality and the Cauchy–Schwarz bound}}


\textbf{Claim 1.} Let \(H\) be a real normed space, \(\lambda\in\mathbb R\), and \(U:H\to\mathbb R\) \(\lambda\)-strongly convex. If \(U\) has Fréchet derivative \(A\) at \(x\), then for every \(y\in H\),


\[U(x)+A(y-x)+\frac\lambda2\|x-y\|^2\le U(y).\]


\textbf{Proof.} Put \(f(s)=U(x+s(y-x))\). By the chain rule \(f\) is differentiable at \(0\) with \(f'(0)=A(y-x)\). Since \(x+s(y-x)=(1-s)x+sy\), strong convexity gives \(f(s)\le(1-s)f(0)+sU(y)-\frac\lambda2s(1-s)\|x-y\|^2\) for \(s\in[0,1]\). For \(0<s<1\) this is equivalent to


\[\frac{f(s)-f(0)}s\le U(y)-f(0)-\frac\lambda2(1-s)\|x-y\|^2.\]


As \(s\downarrow0\) the left side tends to \(f'(0)\) and the right side to \(U(y)-U(x)-\frac\lambda2\|x-y\|^2\); limits preserve the inequality. \(\square\)




\textbf{Claim 2.} Let \(Q\) be a measure and let \(A,B\) be square-integrable functions with values in a real inner product space. Then \(\int\langle A,B\rangle\,dQ\le\|A\|_{L^2(Q)}\|B\|_{L^2(Q)}\).




\textbf{Proof.} The integral is the inner product of \(A\) and \(B\) in the Hilbert space \(L^2(Q)\), and the Cauchy-Schwarz inequality bounds it by the product of the norms. \(\square\)

\begin{leanbox}{Lean: \texttt{strongConvex\_\allowbreak{}tangent}, \texttt{integral\_\allowbreak{}inner\_\allowbreak{}le\_\allowbreak{}secondCost}}
\leansource{\texttt{strongConvex\_\allowbreak{}tangent} \quad LatentDiffusion/\allowbreak{}RegularTransport.lean, lines 380–400}
\begin{leancode}
\leanline{0}{theorem\ strongConvex\_tangent\ (U\ :\ H\ →\ ℝ)\ (rho\ :\ ℝ)}
\leanline{4}{(hU\ :\ StrongConvexOn\ univ\ rho\ U)\ (x\ y\ :\ H)\ (A\ :\ H\ →L[ℝ]\ ℝ)}
\leanline{4}{(hA\ :\ HasFDerivAt\ U\ A\ x)\ :}
\leanline{4}{U\ x\ +\ A\ (y\ -\ x)\ +\ rho\ /\ 2\ *\ ‖x\ -\ y‖\ \textasciicircum{}\ 2\ ≤\ U\ y}
\end{leancode}
\leansource{\texttt{integral\_\allowbreak{}inner\_\allowbreak{}le\_\allowbreak{}secondCost} \quad LatentDiffusion/\allowbreak{}RegularTransport.lean, lines 404–414}
\begin{leancode}
\leanline{0}{theorem\ integral\_inner\_le\_secondCost\ \{Ω\ :\ Type*\}\ [MeasurableSpace\ Ω]}
\leanline{4}{[InnerProductSpace\ ℝ\ H]\ (Q\ :\ Measure\ Ω)\ (A\ B\ :\ Ω\ →\ H)}
\leanline{4}{(hA\ :\ MemLp\ A\ 2\ Q)\ (hB\ :\ MemLp\ B\ 2\ Q)\ :}
\leanline{4}{(∫\ ω,\ inner\ ℝ\ (A\ ω)\ (B\ ω)\ ∂Q)\ ≤}
\leanline{6}{(eLpNorm\ A\ 2\ Q).toReal\ *\ (eLpNorm\ B\ 2\ Q).toReal}
\end{leancode}
\leanlabel{How the Lean proof follows the argument}
\begin{enumerate}


\item \texttt{strong\allowbreak{}Convex\_\allowbreak{}tangent}: \texttt{hl} and \texttt{hd} compute \(f'(0)\) by the chain rule \texttt{comp\_\allowbreak{}has\allowbreak{}Deriv\allowbreak{}At}; \texttt{hconv} is the convexity inequality along the segment. \texttt{Source\allowbreak{}Results.\allowbreak{}geodesic\_\allowbreak{}convexity\_\allowbreak{}first\_\allowbreak{}order} is the scalar statement “right derivative at \(0\) is at most \(f_1-f(0)-\frac\lambda2D^2\)”, proved there by the limit of difference quotients described above.

\item \texttt{integral\_\allowbreak{}inner\_\allowbreak{}le\_\allowbreak{}second\allowbreak{}Cost} converts \(A\) and \(B\) into elements of \texttt{Lp} with \texttt{to\allowbreak{}Lp}, applies \texttt{real\_\allowbreak{}inner\_\allowbreak{}le\_\allowbreak{}norm}, and identifies the inner product with the integral by \texttt{L2.\allowbreak{}inner\_\allowbreak{}def} and the norms with \texttt{Lp.\allowbreak{}norm\_\allowbreak{}to\allowbreak{}Lp}.

\end{enumerate}
\end{leanbox}
\section{\texorpdfstring{Particle scaling preserves the original force constants}{Particle scaling preserves the original force constants}}


Let \(K>0\) and let \(U_i\) be differentiable potentials on a parameter-latent pair, each \(m\)-strongly convex with \(L\)-Lipschitz gradient. Use scaled cloud coordinates \(z=(a,x_1/\sqrt K,\ldots,x_K/\sqrt K)\). Write \(J_i z=(a,x_i)\), \(V(z)=\sum_{i=1}^KU_i(J_i z)\), and \(F(z)=K^{-1}\nabla V(z)\). This normalization prevents an artificial deterioration of the monotonicity and smoothness constants when the particle count grows.

\subsection{\texorpdfstring{A replicated-coordinate identity gives the scaled force estimate}{A replicated-coordinate identity gives the scaled force estimate}}


\textbf{Claim.} The normalized force \(F\) is \(m\)-strongly monotone and \(L\)-Lipschitz.



\textbf{Proof.} The coordinate maps satisfy \(\sum_i\|J_i d\|^2=K\|d\|^2\): the common parameter square appears \(K\) times and each scaled latent square is multiplied by \(K\). Summing the strong-convexity inequalities for the potentials therefore makes \(V\) \(Km\)-strongly convex. Differentiability follows by the chain rule and finite summation. For the gradient, the chain rule gives

\[\langle\nabla V(x)-\nabla V(y),w\rangle=\sum_i\langle\nabla U_i(J_ix)-\nabla U_i(J_iy),J_iw\rangle.\]

Lipschitz continuity, Cauchy-Schwarz for the finite sum, and the coordinate identity bound the absolute value by \(KL\|x-y\|\|w\|\). Taking the supremum over unit vectors gives \(\|\nabla V(x)-\nabla V(y)\|\le KL\|x-y\|\). Strong convexity similarly gives \(\langle\nabla V(x)-\nabla V(y),x-y\rangle\ge Km\|x-y\|^2\). Multiplication by \(1/K>0\) proves both assertions for \(F\). The argument allows distinct potentials and arbitrary finite parameter and latent dimensions. \(\square\)

\begin{leanbox}{Lean: \texttt{joint\_\allowbreak{}norm\_\allowbreak{}sq\_\allowbreak{}sum}, \texttt{meanForce\_\allowbreak{}regular}}
\leansource{\texttt{joint\_\allowbreak{}norm\_\allowbreak{}sq\_\allowbreak{}sum} \quad LatentDiffusion/\allowbreak{}Cloud.lean, lines 79–86}
\begin{leancode}
\leanline{0}{theorem\ joint\_norm\_sq\_sum\ \{N\ dθ\ dx\ :\ ℕ\}\ (z\ :\ State\ N\ dθ\ dx)\ :}
\leanline{4}{∑\ i\ :\ Fin\ N,\ ‖joint\ i\ z‖\ \textasciicircum{}\ 2\ =\ (N\ :\ ℝ)\ *\ ‖z‖\ \textasciicircum{}\ 2}
\end{leancode}
\leansource{\texttt{meanForce\_\allowbreak{}regular} \quad LatentDiffusion/\allowbreak{}Cloud.lean, lines 233–253}
\begin{leancode}
\leanline{0}{theorem\ meanForce\_regular\ \{N\ dθ\ dx\ :\ ℕ\}\ (hN\ :\ 0\ <\ N)}
\leanline{4}{(U\ :\ Fin\ N\ →\ Joint\ dθ\ dx\ →\ ℝ)\ (m\ :\ ℝ)\ (L\ :\ ℝ≥0)}
\leanline{4}{(hdiff\ :\ ∀\ i,\ Differentiable\ ℝ\ (U\ i))}
\leanline{4}{(hsc\ :\ ∀\ i,\ StrongConvexOn\ Set.univ\ m\ (U\ i))}
\leanline{4}{(hLip\ :\ ∀\ i,\ LipschitzWith\ L\ (gradient\ (U\ i)))\ :}
\leanline{4}{LatentDiffusion.ForceRegularity\ (meanForce\ U)\ m\ L}
\end{leancode}
\leanlabel{How the Lean proof follows the argument}
\texttt{joint\_\allowbreak{}norm\_\allowbreak{}sq\_\allowbreak{}sum} expands the Euclidean coordinate sums. \texttt{mean\allowbreak{}Force\_\allowbreak{}regular} constructs the two fields of \texttt{Force\allowbreak{}Regularity}: it obtains monotonicity from \texttt{potential\_\allowbreak{}strong\allowbreak{}Convex} and \texttt{ULA.\allowbreak{}strong\_\allowbreak{}gradient\_\allowbreak{}monotone}, and smoothness from \texttt{potential\_\allowbreak{}gradient\_\allowbreak{}lipschitz}. In each branch the proof multiplies by the nonnegative reciprocal of the particle count and cancels it using \texttt{inv\_\allowbreak{}mul\_\allowbreak{}cancel\ensuremath{{}_{0}}}. The hypothesis \texttt{h\allowbreak{}N} ensures this reciprocal exists.
\end{leanbox}
\section{\texorpdfstring{The density model built from the joint potentials}{The density model built from the joint potentials}}


This section constructs the tilted latent model from the joint potentials \(U_1,\dots,U_M\) of assumption A3 and identifies its optimizer. For a parameter \(a\in\mathbb R^{d_\theta}\) and a datum \(m\), the unnormalized posterior density on \(\mathbb R^{d_z}\) is \(\rho_a^m(z)=e^{-U_m(a,z)}\), its normalizer is \(Z_m(a)=\int\rho_a^m\,dz\), and the posterior is the probability measure


\[\pi_a^m(dz)=\frac{\rho_a^m(z)}{Z_m(a)}\,dz.\]


The log evidence is \(\ell(a)=M^{-1}\sum_m\log Z_m(a)=-V(a)\), with \(V\) the average marginal potential of the previous section. We show that measurability, positivity, integrability and positivity of the normalizers follow from A3, so the construction requires no further hypothesis; that the free energy \(F(a,q)=-\ell(a)+M^{-1}\sum_m\mathrm{KL}(q^m\|\pi_a^m)\) of this model has exactly one global minimizer \((a_\star,\pi_{a_\star})\); and that every posterior has a finite second moment. The last step transfers uniqueness to the paper's literal density integral


\[\widetilde F(a,q)=\frac1M\sum_{m=1}^M\int_{\mathbb R^{d_z}}\tilde q^m(z)\log\frac{\tilde q^m(z)}{\rho_a^m(z)}\,dz\]


on the domain of collections \(q\) with finite relative entropies and finite second moments; here \(\tilde q^m\) is the Lebesgue density of \(q^m\).




This domain is the corrected form of the paper's \(\mathcal P_2\) domain (extra assumption R1).




The optimizer \(a_\star\) is the target of both main convergence theorems.

\subsection{\texorpdfstring{Latent slices of the joint potentials}{Latent slices of the joint potentials}}


\textbf{Claim.} Let \(U:\mathbb R^{d_\theta}\times\mathbb R^{d_z}\to\mathbb R\) and fix \(a\in\mathbb R^{d_\theta}\). Write \(U_a(z)=U(a,z)\).



\begin{enumerate}


\item If \(U\) is differentiable, then \(U_a\) is differentiable.

\item If \(U\) is \(\lambda\)-strongly convex for the Euclidean norm of pairs, \(\lambda\in\mathbb R\), then \(U_a\) is \(\lambda\)-strongly convex on \(\mathbb R^{d_z}\), with the same constant.

\end{enumerate}



\textbf{Proof.} Let \(\iota:\mathbb R^{d_z}\to\mathbb R^{d_\theta}\times\mathbb R^{d_z}\), \(\iota(z)=(0,z)\); it is linear and continuous. Then \((a,z)=(a,0)+\iota(z)\), so \(U_a\) is the composition of \(U\) with an affine map, which is differentiable; the chain rule proves (1).




For (2), let \(x,y\in\mathbb R^{d_z}\) and \(\alpha,\beta\ge0\) with \(\alpha+\beta=1\). Since \(\alpha a+\beta a=a\),


\[(a,\alpha x+\beta y)=\alpha(a,x)+\beta(a,y),\qquad\|(a,x)-(a,y)\|^2=\|0\|^2+\|x-y\|^2=\|x-y\|^2.\]


Strong convexity of \(U\) at the points \((a,x)\) and \((a,y)\) therefore reads


\[U_a(\alpha x+\beta y)\le\alpha U_a(x)+\beta U_a(y)-\alpha\beta\,\frac\lambda2\|x-y\|^2,\]


which is \(\lambda\)-strong convexity of \(U_a\). \(\square\)

\begin{leanbox}{Lean: \texttt{latentInsert}, \texttt{slice\_\allowbreak{}differentiable}, \texttt{slice\_\allowbreak{}strongConvex}}
\leansource{\texttt{latentInsert} \quad LatentDiffusion/\allowbreak{}RegularModel.lean, lines 13–17}
\begin{leancode}
\leanline{0}{def\ latentInsert\ :\ EuclideanSpace\ ℝ\ (Fin\ dz)\ →L[ℝ]\ IPLA.Joint\ dθ\ dz\ :=}
\leanline{2}{LinearMap.toContinuousLinearMap}
\leanline{4}{\{\ toFun\ :=\ fun\ z\ ↦\ IPLA.pair\ 0\ z}
\leanline{6}{map\_add'\ :=\ by\ intro\ x\ y;\ ext\ j;\ cases\ j\ <;>\ simp\ [IPLA.pair]}
\leanline{6}{map\_smul'\ :=\ by\ intro\ c\ x;\ ext\ j;\ cases\ j\ <;>\ simp\ [IPLA.pair]\ \}}
\end{leancode}
\leansource{\texttt{slice\_\allowbreak{}differentiable} \quad LatentDiffusion/\allowbreak{}RegularModel.lean, lines 19–28}
\begin{leancode}
\leanline{0}{theorem\ slice\_differentiable\ (U\ :\ IPLA.Joint\ dθ\ dz\ →\ ℝ)}
\leanline{4}{(hU\ :\ Differentiable\ ℝ\ U)\ (a\ :\ EuclideanSpace\ ℝ\ (Fin\ dθ))\ :}
\leanline{4}{Differentiable\ ℝ\ (fun\ z\ ↦\ U\ (IPLA.pair\ a\ z))}
\end{leancode}
\leansource{\texttt{slice\_\allowbreak{}strongConvex} \quad LatentDiffusion/\allowbreak{}RegularModel.lean, lines 30–40}
\begin{leancode}
\leanline{0}{theorem\ slice\_strongConvex\ (U\ :\ IPLA.Joint\ dθ\ dz\ →\ ℝ)\ (m\ :\ ℝ)}
\leanline{4}{(hU\ :\ StrongConvexOn\ univ\ m\ U)\ (a\ :\ EuclideanSpace\ ℝ\ (Fin\ dθ))\ :}
\leanline{4}{StrongConvexOn\ univ\ m\ (fun\ z\ ↦\ U\ (IPLA.pair\ a\ z))}
\end{leancode}
\leanlabel{How the Lean proof follows the argument}
\texttt{latent\allowbreak{}Insert} is the continuous linear map \(\iota\), built from the linear map \texttt{z \ensuremath{\mapsto} IPLA.\allowbreak{}pair 0 z} by \texttt{Linear\allowbreak{}Map.\allowbreak{}to\allowbreak{}Continuous\allowbreak{}Linear\allowbreak{}Map}; \texttt{IPLA.\allowbreak{}pair a z} is the pair \((a,z)\) in the Euclidean space \texttt{IPLA.\allowbreak{}Joint d\ensuremath{\theta} dz}.



\begin{enumerate}


\item \texttt{slice\_\allowbreak{}differentiable}: \texttt{hp} is the identity \((a,z)=(a,0)+\iota(z)\), proved coordinatewise by \texttt{ext j; cases j}. The composition \texttt{h\allowbreak{}U.\allowbreak{}comp (latent\allowbreak{}Insert.\allowbreak{}differentiable.\allowbreak{}const\_\allowbreak{}add \_\allowbreak{})} is the chain rule.

\item \texttt{slice\_\allowbreak{}strong\allowbreak{}Convex}: \texttt{h} is strong convexity of \texttt{U} at the two pairs. \texttt{ha} is \(r a+s a=a\) for weights summing to one. The rewrites \texttt{IPLA.\allowbreak{}pair\_\allowbreak{}affine} (affine combinations of pairs) and \texttt{IPLA.\allowbreak{}pair\_\allowbreak{}sub} (differences of pairs) put the interpolated point and the difference in pair form, and \texttt{IPLA.\allowbreak{}norm\_\allowbreak{}pair\_\allowbreak{}sq} is \(\|(u,v)\|^2=\|u\|^2+\|v\|^2\). The final \texttt{simpa} removes the zero parameter difference.

\end{enumerate}
\end{leanbox}
\subsection{\texorpdfstring{The density model and its evidence}{The density model and its evidence}}


\textbf{Density models.} Let \(\nu\) be a measure on a measurable space \(Z\). A density model with \(M\) data consists of functions \(\rho_a^m:Z\to\mathbb R\), indexed by parameters \(a\) and data \(m\), such that each \(\rho_a^m\) is measurable, strictly positive and \(\nu\)-integrable with \(Z_m(a)=\int\rho_a^m\,d\nu>0\). Its posteriors are \(\pi_a^m(dz)=Z_m(a)^{-1}\rho_a^m(z)\,\nu(dz)\); these are probability measures, because the density \(\rho_a^m/Z_m(a)\) is nonnegative and has integral \(Z_m(a)/Z_m(a)=1\). Its posterior family has log evidence \(\ell(a)=M^{-1}\sum_m\log Z_m(a)\) and posteriors \(\pi_a^m\).




\textbf{Claim.} Let \(\lambda>0\) and let every \(U_m\) be differentiable and \(\lambda\)-strongly convex. Then \(\rho_a^m(z)=e^{-U_m(a,z)}\), with \(\nu\) Lebesgue measure on \(\mathbb R^{d_z}\), is a density model, and its log evidence satisfies


\[\ell(a)=\frac1M\sum_{m=1}^M\log Z_m(a)=-V(a)\qquad\text{for every }a.\]


\textbf{Proof.} \emph{Measurability.} By the previous step the slice \(z\mapsto U_m(a,z)\) is differentiable, hence continuous; composing with the continuous functions \(t\mapsto-t\) and \(\exp\) gives a continuous, hence Borel measurable, function. \emph{Positivity.} The exponential is positive. \emph{Integrability and positive normalizer.} These are the facts recalled in the marginal section: strong convexity gives the Gaussian bound \(e^{-U_m(a,z)}\le Ce^{-\lambda\|z\|^2/4}\), and the integral of a positive continuous integrable function is positive.




\emph{Evidence.} By definition \(V(a)=M^{-1}\sum_m(-\log Z_m(a))=-M^{-1}\sum_m\log Z_m(a)\), so \(\ell(a)=-V(a)\). \(\square\)




Thus normalizability of the posteriors is a consequence of A3, not an additional assumption.

\begin{leanbox}{Lean: \texttt{densityModel}, \texttt{densityModel\_\allowbreak{}logEvidence}}
\leansource{\texttt{densityModel} \quad LatentDiffusion/\allowbreak{}RegularModel.lean, lines 44–54}
\begin{leancode}
\leanline{0}{def\ densityModel\ (U\ :\ Fin\ M\ →\ IPLA.Joint\ dθ\ dz\ →\ ℝ)\ (m\ :\ ℝ)\ (hm\ :\ 0\ <\ m)}
\leanline{4}{(hdiff\ :\ ∀\ i,\ Differentiable\ ℝ\ (U\ i))}
\leanline{4}{(hsc\ :\ ∀\ i,\ StrongConvexOn\ univ\ m\ (U\ i))\ :}
\leanline{4}{RegularVariational.DensityModel\ (EuclideanSpace\ ℝ\ (Fin\ dθ))}
\leanline{6}{(EuclideanSpace\ ℝ\ (Fin\ dz))\ M\ volume\ where}
\leanline{2}{density\ a\ i\ z\ :=\ Real.exp\ (-U\ i\ (IPLA.pair\ a\ z))}
\leanline{2}{measurable\ a\ i\ :=\ (Real.continuous\_exp.comp}
\leanline{4}{((slice\_differentiable\ (U\ i)\ (hdiff\ i)\ a).continuous.neg)).measurable}
\leanline{2}{positive\ \_\ \_\ \_\ :=\ Real.exp\_pos\ \_}
\leanline{2}{integrable\ a\ i\ :=\ IPLA.integrable\_marginal\_weight\ (U\ i)\ m\ hm\ (hdiff\ i)\ (hsc\ i)\ a}
\leanline{2}{mass\_positive\ a\ i\ :=\ IPLA.marginalIntegral\_pos\ (U\ i)\ m\ hm\ (hdiff\ i)\ (hsc\ i)\ a}
\end{leancode}
\leansource{\texttt{densityModel\_\allowbreak{}logEvidence} \quad LatentDiffusion/\allowbreak{}RegularModel.lean, lines 56–62}
\begin{leancode}
\leanline{0}{theorem\ densityModel\_logEvidence\ (U\ :\ Fin\ M\ →\ IPLA.Joint\ dθ\ dz\ →\ ℝ)\ (m\ :\ ℝ)\ (hm\ :\ 0\ <\ m)}
\leanline{4}{(hdiff\ :\ ∀\ i,\ Differentiable\ ℝ\ (U\ i))}
\leanline{4}{(hsc\ :\ ∀\ i,\ StrongConvexOn\ univ\ m\ (U\ i))\ (a\ :\ EuclideanSpace\ ℝ\ (Fin\ dθ))\ :}
\leanline{4}{(densityModel\ U\ m\ hm\ hdiff\ hsc).family.logEvidence\ a\ =\ -averageMarginalPotential\ U\ a}
\end{leancode}
\leanlabel{How the Lean proof follows the argument}
The structure \texttt{Regular\allowbreak{}Variational.\allowbreak{}Density\allowbreak{}Model A Z M \ensuremath{\nu}} has the fields \texttt{density}, \texttt{measurable}, \texttt{positive}, \texttt{integrable} and \texttt{mass\_\allowbreak{}positive}, matching the definition. Its \texttt{mass} is \(Z_m(a)\), its \texttt{posterior} is the probability measure with density \(\rho/Z\), and its \texttt{family} is the posterior family with \texttt{log\allowbreak{}Evidence a = (M : \ensuremath{\mathbb R})\ensuremath{{}^{-1}} * ∑ i, Real.\allowbreak{}log (D.\allowbreak{}mass a i)}.



\begin{enumerate}


\item \texttt{density\allowbreak{}Model} fills the fields: \texttt{density a i z := Real.\allowbreak{}exp (-\allowbreak{}U i (IPLA.\allowbreak{}pair a z))}; measurability from continuity of \texttt{Real.\allowbreak{}exp} composed with the slice of \texttt{slice\_\allowbreak{}differentiable}; positivity from \texttt{Real.\allowbreak{}exp\_\allowbreak{}pos}; integrability and positivity of the mass from \texttt{IPLA.\allowbreak{}integrable\_\allowbreak{}marginal\_\allowbreak{}weight} and \texttt{IPLA.\allowbreak{}marginal\allowbreak{}Integral\_\allowbreak{}pos}. The reference measure is Lebesgue measure \texttt{volume}.

\item \texttt{density\allowbreak{}Model\_\allowbreak{}log\allowbreak{}Evidence} unfolds the family, the mass, \texttt{average\allowbreak{}Marginal\allowbreak{}Potential}, \texttt{IPLA.\allowbreak{}marginal\allowbreak{}Potential} and \texttt{IPLA.\allowbreak{}marginal\allowbreak{}Integral}; the two sides then agree after \texttt{Finset.\allowbreak{}sum\_\allowbreak{}neg\_\allowbreak{}distrib}, \texttt{mul\_\allowbreak{}neg} and \texttt{neg\_\allowbreak{}neg}.

\end{enumerate}
\end{leanbox}
\subsection{\texorpdfstring{The unique variational optimizer of the model}{The unique variational optimizer of the model}}


\textbf{Claim.} Let \(M\ge1\), \(\lambda>0\), and let every \(U_m\) be differentiable and \(\lambda\)-strongly convex. For the density model of the previous step there is exactly one pair \((a,q)\), with \(a\in\mathbb R^{d_\theta}\) and \(q\) a collection of \(M\) probability measures on \(\mathbb R^{d_z}\), that is a global minimizer of


\[F(a,q)=-\ell(a)+\frac1M\sum_{m=1}^M\mathrm{KL}(q^m\|\pi_a^m)\]


over admissible pairs (pairs with finite relative entropies). This pair is \((a_\star,\pi_{a_\star})\), where \(a_\star\) is the unique minimizer of \(V\).




\textbf{Proof.} The marginal section proves: if \(M\ge1\), \(\lambda>0\), every \(U_m\) is differentiable and \(\lambda\)-strongly convex, and \((\ell,\pi)\) is a posterior family with \(\ell=-V\), then \(F\) has exactly one global minimizer, \((a_\star,\pi_{a_\star})\). The posterior family of the density model satisfies \(\ell=-V\) by the previous step. \(\square\)

\begin{leanbox}{Lean: \texttt{existsUnique\_\allowbreak{}variationalOptimizer}}
\leansource{\texttt{existsUnique\_\allowbreak{}variationalOptimizer} \quad LatentDiffusion/\allowbreak{}RegularModel.lean, lines 66–74}
\begin{leancode}
\leanline{0}{theorem\ existsUnique\_variationalOptimizer\ (U\ :\ Fin\ M\ →\ IPLA.Joint\ dθ\ dz\ →\ ℝ)}
\leanline{4}{(hM\ :\ 0\ <\ M)\ (m\ :\ ℝ)\ (hm\ :\ 0\ <\ m)}
\leanline{4}{(hdiff\ :\ ∀\ i,\ Differentiable\ ℝ\ (U\ i))}
\leanline{4}{(hsc\ :\ ∀\ i,\ StrongConvexOn\ univ\ m\ (U\ i))\ :}
\leanline{4}{∃!\ p\ :\ EuclideanSpace\ ℝ\ (Fin\ dθ)\ ×}
\leanline{8}{(Fin\ M\ →\ ProbabilityMeasure\ (EuclideanSpace\ ℝ\ (Fin\ dz))),}
\leanline{6}{(densityModel\ U\ m\ hm\ hdiff\ hsc).family.IsGlobalMin\ p.1\ p.2}
\end{leancode}
\leanlabel{How the Lean proof follows the argument}
\texttt{exists\allowbreak{}Unique\_\allowbreak{}variational\allowbreak{}Optimizer} is a one-line application of \texttt{Latent\allowbreak{}Diffusion.\allowbreak{}exists\allowbreak{}Unique\_\allowbreak{}variational\allowbreak{}Optimizer} from the marginal section to the family \texttt{(density\allowbreak{}Model U m hm hdiff hsc).\allowbreak{}family}, with the identification \texttt{density\allowbreak{}Model\_\allowbreak{}log\allowbreak{}Evidence U m hm hdiff hsc} as the hypothesis \(\ell=-V\). The strong-convexity constant is named \texttt{m} in this file.
\end{leanbox}
\subsection{\texorpdfstring{The posteriors have finite second moments}{The posteriors have finite second moments}}


\textbf{Claim.} Under the hypotheses of the previous step, for every \(a\) and every \(m\),


\[\int_{\mathbb R^{d_z}}\|z\|^2\,\pi_a^m(dz)<\infty.\]


\textbf{Proof.} Put \(W(z)=U_m(a,z)\). For a function \(W\) with \(0<\int e^{-W}dz<\infty\), the Gibbs measure of \(W\) is \(\gamma_W(dz)=e^{-W(z)}dz/\int e^{-W}\). We first show \(\pi_a^m=\gamma_W\). Both measures have a density with respect to Lebesgue measure: \(\pi_a^m\) has density \(e^{-W(z)}/Z_m(a)\), and \(\gamma_W\) has density \(e^{-W(z)}\) divided by \(\int e^{-W}dz\), which equals \(Z_m(a)\). The two densities coincide at every point, so the measures coincide.




By the first step, \(W\) is differentiable and \(\lambda\)-strongly convex with \(\lambda>0\). The report on the unadjusted Langevin algorithm proves that the Gibbs measure of such a potential has a finite second moment; the reason is the bound \(W(z)\ge W(0)-\|\nabla W(0)\|^2/\lambda+\frac\lambda4\|z\|^2\), which dominates \(\|z\|^2e^{-W(z)}\) by a multiple of \(\|z\|^2e^{-\lambda\|z\|^2/4}\), an integrable function. Hence \(\int\|z\|^2\,d\pi_a^m=\int\|z\|^2\,d\gamma_W<\infty\). \(\square\)




The posteriors therefore belong to the paper's finite-second-moment domain \(\mathcal P_2\), as a consequence of A3.

\begin{leanbox}{Lean: \texttt{posterior\_\allowbreak{}memLp}}
\leansource{\texttt{posterior\_\allowbreak{}memLp} \quad LatentDiffusion/\allowbreak{}RegularModel.lean, lines 78–104}
\begin{leancode}
\leanline{0}{theorem\ posterior\_memLp\ (U\ :\ Fin\ M\ →\ IPLA.Joint\ dθ\ dz\ →\ ℝ)\ (m\ :\ ℝ)\ (hm\ :\ 0\ <\ m)}
\leanline{4}{(hdiff\ :\ ∀\ i,\ Differentiable\ ℝ\ (U\ i))}
\leanline{4}{(hsc\ :\ ∀\ i,\ StrongConvexOn\ univ\ m\ (U\ i))\ (a\ :\ EuclideanSpace\ ℝ\ (Fin\ dθ))\ (i\ :\ Fin\ M)\ :}
\leanline{4}{MemLp\ (fun\ z\ :\ EuclideanSpace\ ℝ\ (Fin\ dz)\ ↦\ z)\ 2}
\leanline{6}{((densityModel\ U\ m\ hm\ hdiff\ hsc).posterior\ a\ i\ :\ Measure\ \_)}
\end{leancode}
\leanlabel{How the Lean proof follows the argument}
The conclusion \texttt{Mem\allowbreak{}Lp (fun z \ensuremath{\mapsto} z) 2 (posterior a i)} states that the identity map is square integrable for the posterior, which is the claim.



\begin{enumerate}


\item \texttt{hid} identifies the posterior with \texttt{ULA.\allowbreak{}gibbs\allowbreak{}Measure V}. Mathlib's Gibbs measure is written with an extended-valued normalizer: \texttt{ULA.\allowbreak{}gibbs\allowbreak{}Measure V} is \texttt{(gibbs\allowbreak{}Partition V)\ensuremath{{}^{-1}} \ensuremath{\bullet} volume.\allowbreak{}with\allowbreak{}Density (gibbs\allowbreak{}Weight V)}, with \texttt{gibbs\allowbreak{}Weight V z = ENNReal.\allowbreak{}of\allowbreak{}Real (exp (-\allowbreak{}V z))} and \texttt{gibbs\allowbreak{}Partition V} its lower integral \texttt{\ensuremath{\int}\ensuremath{{}^{-}}}. After unfolding, \texttt{with\allowbreak{}Density\_\allowbreak{}smul} moves the scalar into the density and \texttt{with\allowbreak{}Density\_\allowbreak{}congr\_\allowbreak{}ae} reduces the identity to a pointwise identity of densities.

\item Pointwise, \texttt{ENNReal.\allowbreak{}of\allowbreak{}Real\_\allowbreak{}div\_\allowbreak{}of\_\allowbreak{}pos h\allowbreak{}Z} splits \(\mathrm{ofReal}(e^{-V}/Z)\) into a quotient, and \texttt{of\allowbreak{}Real\_\allowbreak{}integral\_\allowbreak{}eq\_\allowbreak{}lintegral\_\allowbreak{}of\allowbreak{}Real} converts the real normalizer into the lower integral of the extended weight; this uses integrability and nonnegativity of \(e^{-V}\). The remaining identity is rearranged by \texttt{simp only}.

\item The moment bound is \texttt{ULA.\allowbreak{}mem\allowbreak{}Lp\_\allowbreak{}id\_\allowbreak{}gibbs\allowbreak{}Measure V m hm hdiff\allowbreak{}V hsc\allowbreak{}V}, with \texttt{hdiff\allowbreak{}V} and \texttt{hsc\allowbreak{}V} from the first step.

\end{enumerate}
\end{leanbox}
\subsection{\texorpdfstring{The literal density integral on the finite-second-moment domain}{The literal density integral on the finite-second-moment domain}}


The paper writes the free energy as the literal density integral \(\widetilde F\) of the summary. The variational section on regular density models proves the following facts for every density model with \(M\ge1\) data.



\begin{itemize}


\item If \(q\) is admissible at \(a\), then each \(q^m\) has a Lebesgue density \(\tilde q^m\), constructed from the Radon–Nikodym derivative \(r=dq^m/d\pi_a^m\) as \(\tilde q^m=r\,\rho_a^m/Z_m(a)\); the function \(\tilde q^m\log(\tilde q^m/\rho_a^m)\) is integrable; and \(\int\tilde q^m\log(\tilde q^m/\rho_a^m)\,dz=\mathrm{KL}(q^m\|\pi_a^m)-\log Z_m(a)\). Consequently \(\widetilde F(a,q)=F(a,q)\) on admissible pairs.

\item Let \(\mathcal D\) be any set of collections that contains every posterior collection \(\pi_b\). Then \(q\) is admissible at \(a\), \(q\in\mathcal D\), and \(\widetilde F(a,q)\le\widetilde F(b,r)\) for every admissible \((b,r)\) with \(r\in\mathcal D\), if and only if \(a\) maximizes \(\ell\) and \(q=\pi_a\).

\end{itemize}



\textbf{Claim.} Let \(M\ge1\), \(\lambda>0\), and let every \(U_m\) be differentiable and \(\lambda\)-strongly convex. Then there is exactly one pair \((a,q)\) such that



\begin{enumerate}


\item \(q\) is admissible at \(a\) and every \(q^m\) has a finite second moment, and

\item \(\widetilde F(a,q)\le\widetilde F(b,r)\) for every \(b\) and every \(r\) admissible at \(b\) whose measures have finite second moments.

\end{enumerate}



\textbf{Proof.} Let \(\mathcal D\) be the set of collections whose measures all have finite second moments. By the previous step every posterior collection lies in \(\mathcal D\). By the second cited fact, a pair satisfies (1) and (2) if and only if \(a\) maximizes \(\ell\) and \(q=\pi_a\). By the variational characterization this holds if and only if \((a,q)\) is a global minimizer of \(F\). By the third step there is exactly one such pair. \(\square\)




In particular the paper's Proposition 1 holds on its finite-second-moment domain with the literal integral: the unique optimizer is \((a_\star,\pi_{a_\star})\).

\begin{leanbox}{Lean: \texttt{existsUnique\_\allowbreak{}p2\_\allowbreak{}literalOptimizer}}
\leansource{\texttt{existsUnique\_\allowbreak{}p2\_\allowbreak{}literalOptimizer} \quad LatentDiffusion/\allowbreak{}RegularModel.lean, lines 108–133}
\begin{leancode}
\leanline{0}{theorem\ existsUnique\_p2\_literalOptimizer\ (U\ :\ Fin\ M\ →\ IPLA.Joint\ dθ\ dz\ →\ ℝ)}
\leanline{4}{(hM\ :\ 0\ <\ M)\ (m\ :\ ℝ)\ (hm\ :\ 0\ <\ m)}
\leanline{4}{(hdiff\ :\ ∀\ i,\ Differentiable\ ℝ\ (U\ i))}
\leanline{4}{(hsc\ :\ ∀\ i,\ StrongConvexOn\ univ\ m\ (U\ i))\ :}
\leanline{4}{let\ D\ :=\ densityModel\ U\ m\ hm\ hdiff\ hsc}
\leanline{4}{∃!\ p\ :\ EuclideanSpace\ ℝ\ (Fin\ dθ)\ ×}
\leanline{8}{(Fin\ M\ →\ ProbabilityMeasure\ (EuclideanSpace\ ℝ\ (Fin\ dz))),}
\leanline{6}{D.family.Admissible\ p.1\ p.2\ ∧}
\leanline{6}{(∀\ i,\ MemLp\ (fun\ z\ :\ EuclideanSpace\ ℝ\ (Fin\ dz)\ ↦\ z)\ 2\ (p.2\ i\ :\ Measure\ \_))\ ∧}
\leanline{6}{∀\ b\ q,\ D.family.Admissible\ b\ q\ →}
\leanline{8}{(∀\ i,\ MemLp\ (fun\ z\ :\ EuclideanSpace\ ℝ\ (Fin\ dz)\ ↦\ z)\ 2\ (q\ i\ :\ Measure\ \_))\ →}
\leanline{8}{D.literalEnergy\ p.1\ p.2\ ≤\ D.literalEnergy\ b\ q}
\end{leancode}
\leanlabel{How the Lean proof follows the argument}
In \texttt{exists\allowbreak{}Unique\_\allowbreak{}p2\_\allowbreak{}literal\allowbreak{}Optimizer}, \texttt{D.\allowbreak{}literal\allowbreak{}Energy} is \(\widetilde F\); the trial density used inside it is the Radon–Nikodym construction of the first cited fact (the definition trialDensity of \texttt{Regular\allowbreak{}Variational.\allowbreak{}Density\allowbreak{}Model}), so no density of the trial law is assumed.



\begin{enumerate}


\item \texttt{domain} is the predicate “every measure has finite second moment”, written with \texttt{Mem\allowbreak{}Lp (fun z \ensuremath{\mapsto} z) 2}, and \texttt{hpost} proves that posteriors satisfy it by \texttt{posterior\_\allowbreak{}mem\allowbreak{}Lp}.

\item \texttt{heq p} chains two equivalences: \texttt{D.\allowbreak{}literal\_\allowbreak{}minimizers\_\allowbreak{}on\_\allowbreak{}iff h\allowbreak{}M domain hpost} (the second cited fact) and the reversed \texttt{D.\allowbreak{}family.\allowbreak{}global\_\allowbreak{}minimizers\_\allowbreak{}iff h\allowbreak{}M} (the variational characterization).

\item \texttt{obtain \ensuremath{\langle}p, hp, hu\ensuremath{\rangle}} takes the unique global minimizer from \texttt{exists\allowbreak{}Unique\_\allowbreak{}variational\allowbreak{}Optimizer}, and the final term transports existence and uniqueness through \texttt{heq}.

\end{enumerate}
\end{leanbox}
\section{\texorpdfstring{The extended log-Sobolev and Talagrand inequalities}{The extended log-Sobolev and Talagrand inequalities}}


This section proves the two functional inequalities that drive the convergence of the classical flow: the extended log-Sobolev inequality (the paper's Theorem stating that strong log-concavity implies extended LSI) and a Talagrand-type quadratic growth inequality, both for the normalized free energy of \(M\) data. The setting is that of the previous section, now for each datum \(m\): a probability law \(Q_m\) on \(Z\) with positive continuously differentiable density \(\rho_{m}\), regular displacement data \(T_m\) from \(Q_m\), and a jointly \(\lambda\)-strongly convex potential \(U_m\) on \(P\times Z\). With \(p_0\) the current parameter, \(p_1\) the target parameter and \(\mathcal F_m\) the free energy of datum \(m\) along the displacement, put \(\mathcal F=M^{-1}\sum_m\mathcal F_m\), and define


\[\begin{aligned}G&=\frac1M\sum_m\int\nabla_pU_m(p_0,z)\,Q_m(dz),\\ S_m&=\nabla_z\bigl(\log\rho_m+U_m(p_0,\cdot)\bigr),\\I&=\|G\|^2+\frac1M\sum_m\int\|S_m\|^2\,dQ_m.\end{aligned}\]


The main result is \(2\lambda\bigl(\mathcal F(0)-\mathcal F(1)\bigr)\le I\). In the convergence theorem \(\mathcal F(0)\) is the free energy of the current state and \(\mathcal F(1)\) that of the posterior optimizer, so this is \(2\lambda(F-F_\star)\le I\). The proof sums two tangent inequalities, one for the entropy (from the Jacobian formula, \(\log x\le x-1\) and integration by parts) and one for the potential (from strong convexity), and then completes a square in which the parameter forces are averaged over the data before the norm is squared. The last step proves \(\frac\lambda2d^2\le\mathcal F(0)-\mathcal F(1)\) when the endpoint minimizes \(\mathcal F\) along the displacement, where \(d\) is the normalized product distance. No first-order transport inequality or log-Sobolev inequality is assumed.

\subsection{\texorpdfstring{The entropy tangent inequality}{The entropy tangent inequality}}


Let \(D\) be regular displacement data from a probability law \(\mu\) on \(Z\), with Lebesgue measure as reference: transport \(T\), interpolation \(T_s\), densities \(\rho_s\), Jacobian eigenvalues \(\lambda_i(z)>0\), and entropies \(H(s)=\int\log\rho_s\,d\mu_s=\int\log\rho_s(T_s(z))\,\mu(dz)\).




\textbf{Claim.} (1) \(\mu(dz)=\rho_0(z)\,dz\). (2) The displacement \(V(z)=T(z)-z\) has divergence \(\nabla\cdot V(z)=\sum_i(\lambda_i(z)-1)\). (3) For every \(z\), \(\log\rho_0(z)-\sum_i(\lambda_i(z)-1)\le\log\rho_1(T(z))\). (4) If moreover \(\rho_0\) is continuously differentiable and \((V,\rho_0)\) satisfies the no-boundary-flux condition, then


\[\int_Z\langle T(z)-z,\nabla\log\rho_0(z)\rangle\,\mu(dz)\le H(1)-H(0).\]


\textbf{Proof.} (1) is the law condition at \(s=0\), where \(T_0\) is the identity. (2) \(V\) has derivative \(DT(z)-I\) at \(z\), so its divergence is \(\operatorname{tr}(DT(z)-I)\). In the orthonormal eigenbasis \((b_i(z))\), \(\langle(DT(z)-I)b_i,b_i\rangle=\lambda_i(z)-1\), and the trace is the sum of these numbers.




(3) By the change-of-variables formula at \(s=1\), \(\rho_1(T(z))=\rho_0(z)/\prod_i\lambda_i(z)\), so \(\log\rho_1(T(z))=\log\rho_0(z)-\sum_i\log\lambda_i(z)\). The inequality \(\log x\le x-1\) for \(x>0\) gives \(\sum_i\log\lambda_i(z)\le\sum_i(\lambda_i(z)-1)\), which proves (3).




(4) By (1) and the score pairing of the flow-calculus section, applied with \(q=\rho_0\) and \(V=T-\mathrm{id}\), both \(\nabla\cdot V\) and \(\langle V,\nabla\log\rho_0\rangle\) are \(\mu\)-integrable and


\[\int\langle T(z)-z,\nabla\log\rho_0(z)\rangle\,\mu(dz)=-\int\nabla\cdot V\,d\mu=-\int\sum_i(\lambda_i-1)\,d\mu,\]


using (2). The functions \(\log\rho_0\) and \(\log\rho_1\circ T\) are \(\mu\)-integrable by the extra transport assumptions. Integrating (3) against \(\mu\) gives \(H(0)-\int\sum_i(\lambda_i-1)\,d\mu\le H(1)\), which is (4) after substituting the displayed identity. \(\square\)

\begin{leanbox}{Lean: \texttt{initial\_\allowbreak{}density\_\allowbreak{}law}, \texttt{displacement\_\allowbreak{}divergence}, \texttt{log\_\allowbreak{}density\_\allowbreak{}tangent} and 1 more}
\leansource{\texttt{initial\_\allowbreak{}density\_\allowbreak{}law} \quad LatentDiffusion/\allowbreak{}RegularTransportFirstOrder.lean, lines 34–35}
\begin{leancode}
\leanline{0}{theorem\ initial\_density\_law\ :\ μ\ =\ volume.withDensity\ (fun\ z\ ↦\ ENNReal.ofReal\ (D.density\ 0\ z))}
\end{leancode}
\leansource{\texttt{displacement\_\allowbreak{}divergence} \quad LatentDiffusion/\allowbreak{}RegularTransportFirstOrder.lean, lines 38–50}
\begin{leancode}
\leanline{0}{theorem\ displacement\_divergence\ (z\ :\ Z)\ :}
\leanline{4}{RegularFlowCalculus.fluxDivergence\ (fun\ z\ ↦\ D.transport\ z\ -\ z)\ z\ =}
\leanline{6}{∑\ i,\ (D.eigenvalue\ z\ i\ -\ 1)}
\end{leancode}
\leansource{\texttt{log\_\allowbreak{}density\_\allowbreak{}tangent} \quad LatentDiffusion/\allowbreak{}RegularTransportFirstOrder.lean, lines 54–66}
\begin{leancode}
\leanline{0}{theorem\ log\_density\_tangent\ (z\ :\ Z)\ :}
\leanline{4}{Real.log\ (D.density\ 0\ z)\ -\ ∑\ i,\ (D.eigenvalue\ z\ i\ -\ 1)\ ≤}
\leanline{6}{Real.log\ (D.density\ 1\ (D.transport\ z))}
\end{leancode}
\leansource{\texttt{entropy\_\allowbreak{}tangent} \quad LatentDiffusion/\allowbreak{}RegularTransportFirstOrder.lean, lines 70–92}
\begin{leancode}
\leanline{0}{theorem\ entropy\_tangent}
\leanline{4}{(hq\ :\ ContDiff\ ℝ\ 1\ (D.density\ 0))}
\leanline{4}{(htail\ :\ RegularFlowCalculus.NoBoundaryFlux\ (fun\ z\ ↦\ D.transport\ z\ -\ z)\ (D.density\ 0))\ :}
\leanline{4}{(∫\ z,\ inner\ ℝ\ (D.transport\ z\ -\ z)\ (gradient\ (fun\ z\ ↦\ Real.log\ (D.density\ 0\ z))\ z)\ ∂μ)\ ≤}
\leanline{6}{D.entropy\ 1\ -\ D.entropy\ 0}
\end{leancode}
\leanlabel{How the Lean proof follows the argument}
\begin{enumerate}


\item \texttt{initial\_\allowbreak{}density\_\allowbreak{}law} is \texttt{D.\allowbreak{}density\_\allowbreak{}law 0} simplified by \(T_0=\mathrm{id}\).

\item \texttt{displacement\_\allowbreak{}divergence} applies \texttt{Regular\allowbreak{}Flow\allowbreak{}Calculus.\allowbreak{}flux\allowbreak{}Divergence\_\allowbreak{}eq\_\allowbreak{}trace} to the derivative \texttt{(D.\allowbreak{}derivative\_\allowbreak{}transport z).\allowbreak{}sub (has\allowbreak{}FDeriv\allowbreak{}At\_\allowbreak{}id z)} and evaluates the trace in the spectral basis by \texttt{Linear\allowbreak{}Map.\allowbreak{}trace\_\allowbreak{}eq\_\allowbreak{}matrix\_\allowbreak{}trace} and \texttt{D.\allowbreak{}jacobian\_\allowbreak{}eigenvector}.

\item \texttt{log\_\allowbreak{}density\_\allowbreak{}tangent} rewrites with \texttt{D.\allowbreak{}density\_\allowbreak{}jacobian 1}, \texttt{Real.\allowbreak{}log\_\allowbreak{}div}, \texttt{spectral\allowbreak{}Jacobian\_\allowbreak{}one} and \texttt{Real.\allowbreak{}log\_\allowbreak{}prod}, and sums \texttt{Real.\allowbreak{}log\_\allowbreak{}le\_\allowbreak{}sub\_\allowbreak{}one\_\allowbreak{}of\_\allowbreak{}pos} over the eigenvalues.

\item \texttt{entropy\_\allowbreak{}tangent}: \texttt{h\allowbreak{}IBP} and \texttt{h\allowbreak{}Int} are \texttt{Regular\allowbreak{}Flow\allowbreak{}Calculus.\allowbreak{}transport\_\allowbreak{}score\_\allowbreak{}pairing} and \texttt{Regular\allowbreak{}Flow\allowbreak{}Calculus.\allowbreak{}transport\_\allowbreak{}score\_\allowbreak{}integrable}, transported to \(\mu\) by rewriting with \texttt{D.\allowbreak{}initial\_\allowbreak{}density\_\allowbreak{}law}; \texttt{simp\_\allowbreak{}rw [D.\allowbreak{}displacement\_\allowbreak{}divergence]} inserts (2). \texttt{hbound} integrates (3) with \texttt{integral\_\allowbreak{}mono}, and the entropies are rewritten by \texttt{D.\allowbreak{}entropy\_\allowbreak{}eq\_\allowbreak{}pullback} before \texttt{linarith}.

\end{enumerate}
\end{leanbox}
\subsection{\texorpdfstring{The potential tangent inequality}{The potential tangent inequality}}


The joint space \(P\times Z\) carries the Euclidean product norm. The displacement field of the transport and of the parameter move \(p_0\to p_1\) is \(\Delta(z)=(p_1-p_0,\,T(z)-z)\in P\times Z\).




\textbf{Claim.} (1) If \(z\mapsto z-T(z)\) is square integrable under \(\mu\), so is \(\Delta\). (2) If \(A\) and \(B\) are square integrable under \(\mu\), then \(\langle A,B\rangle\) is \(\mu\)-integrable. (3) Let \(\lambda\in\mathbb R\) and let \(U\) be \(\lambda\)-strongly convex on \(P\times Z\). Assume that for every \(z\), \(U\) has gradient \(G(z)\) at the point \((p_0,z)\); that \(G\) and \(z\mapsto z-T(z)\) are square integrable under \(\mu\); and that \(U(p_0,\cdot)\) and \(z\mapsto U(p_1,T(z))\) are \(\mu\)-integrable. Then


\[\int\langle G(z),\Delta(z)\rangle\,\mu(dz)+\frac\lambda2\bigl(\|p_0-p_1\|^2+\|\mathrm{id}-T\|_{L^2(\mu)}^2\bigr)\le\int U(p_1,T(z))\,\mu(dz)-\int U(p_0,z)\,\mu(dz).\]


\textbf{Proof.} (1) \(\Delta\) is the image of \(z-T(z)\) under the map \(w\mapsto(p_1-p_0,-w)\), which is Lipschitz with constant \(1\); a Lipschitz function of a square-integrable function is square integrable for a finite measure, since \(\|f(w)\|\le\|f(0)\|+\|w\|\). (2) By Cauchy-Schwarz, \(|\langle A,B\rangle|\le\|A\|\|B\|\le\frac12(\|A\|^2+\|B\|^2)\), which is integrable.




(3) Fix \(z\) and apply the tangent inequality of the previous section at \(x=(p_0,z)\) and \(y=(p_1,T(z))\), with \(y-x=\Delta(z)\) and \(\|x-y\|^2=\|p_0-p_1\|^2+\|z-T(z)\|^2\):


\[\langle G(z),\Delta(z)\rangle+\frac\lambda2\bigl(\|p_0-p_1\|^2+\|z-T(z)\|^2\bigr)\le U(p_1,T(z))-U(p_0,z).\]


Both sides are \(\mu\)-integrable by (1), (2) and the hypotheses. Integrating against the probability measure \(\mu\) and using \(\int\|z-T(z)\|^2\,d\mu=\|\mathrm{id}-T\|^2_{L^2(\mu)}\) proves (3). \(\square\)

\begin{leanbox}{Lean: \texttt{Joint}, \texttt{displacementField}, \texttt{displacementField\_\allowbreak{}memLp} and 2 more}
\leansource{\texttt{Joint} \quad LatentDiffusion/\allowbreak{}RegularTransportFirstOrder.lean, lines 147–147}
\begin{leancode}
\leanline{0}{abbrev\ Joint\ :=\ PBGSO.ProductTransport.Joint\ P\ Z}
\end{leancode}
\leansource{\texttt{displacementField} \quad LatentDiffusion/\allowbreak{}RegularTransportFirstOrder.lean, lines 150–152}
\begin{leancode}
\leanline{0}{def\ displacementField\ (D\ :\ RegularDisplacement\ (volume\ :\ Measure\ Z)\ μ\ d)}
\leanline{4}{(p₀\ p₁\ :\ P)\ (z\ :\ Z)\ :\ Joint\ (P\ :=\ P)\ (Z\ :=\ Z)\ :=}
\leanline{2}{PBGSO.ProductTransport.embed\ (p₁\ -\ p₀)\ (D.transport\ z\ -\ z)}
\end{leancode}
\leansource{\texttt{displacementField\_\allowbreak{}memLp} \quad LatentDiffusion/\allowbreak{}RegularTransportFirstOrder.lean, lines 154–162}
\begin{leancode}
\leanline{0}{theorem\ displacementField\_memLp}
\leanline{4}{(D\ :\ RegularDisplacement\ (volume\ :\ Measure\ Z)\ μ\ d)\ (p₀\ p₁\ :\ P)}
\leanline{4}{(hc\ :\ MemLp\ (fun\ z\ ↦\ z\ -\ D.transport\ z)\ 2\ μ)\ :}
\leanline{4}{MemLp\ (displacementField\ D\ p₀\ p₁)\ 2\ μ}
\end{leancode}
\leansource{\texttt{integrable\_\allowbreak{}inner\_\allowbreak{}memLp} \quad LatentDiffusion/\allowbreak{}RegularTransportFirstOrder.lean, lines 164–169}
\begin{leancode}
\leanline{0}{theorem\ integrable\_inner\_memLp\ \{H\ :\ Type*\}\ [NormedAddCommGroup\ H]\ [InnerProductSpace\ ℝ\ H]}
\leanline{4}{(A\ B\ :\ Z\ →\ H)\ (hA\ :\ MemLp\ A\ 2\ μ)\ (hB\ :\ MemLp\ B\ 2\ μ)\ :}
\leanline{4}{Integrable\ (fun\ z\ ↦\ inner\ ℝ\ (A\ z)\ (B\ z))\ μ}
\end{leancode}
\leansource{\texttt{potential\_\allowbreak{}tangent} \quad LatentDiffusion/\allowbreak{}RegularTransportFirstOrder.lean, lines 173–206}
\begin{leancode}
\leanline{0}{theorem\ potential\_tangent}
\leanline{4}{(D\ :\ RegularDisplacement\ (volume\ :\ Measure\ Z)\ μ\ d)}
\leanline{4}{(U\ :\ P\ →\ Z\ →\ ℝ)\ (p₀\ p₁\ :\ P)\ (rho\ :\ ℝ)\ (G\ :\ Z\ →\ Joint\ (P\ :=\ P)\ (Z\ :=\ Z))}
\leanline{4}{(hU\ :\ StrongConvexOn\ univ\ rho\ (fun\ w\ :\ Joint\ (P\ :=\ P)\ (Z\ :=\ Z)\ ↦\ U\ w.fst\ w.snd))}
\leanline{4}{(hG\ :\ ∀\ z,\ HasGradientAt\ (fun\ w\ :\ Joint\ (P\ :=\ P)\ (Z\ :=\ Z)\ ↦\ U\ w.fst\ w.snd)}
\leanline{6}{(G\ z)\ (PBGSO.ProductTransport.embed\ p₀\ z))}
\leanline{4}{(hG2\ :\ MemLp\ G\ 2\ μ)\ (hc\ :\ MemLp\ (fun\ z\ ↦\ z\ -\ D.transport\ z)\ 2\ μ)}
\leanline{4}{(hU0\ :\ Integrable\ (U\ p₀)\ μ)}
\leanline{4}{(hU1\ :\ Integrable\ (fun\ z\ ↦\ U\ p₁\ (D.transport\ z))\ μ)\ :}
\leanline{4}{(∫\ z,\ inner\ ℝ\ (G\ z)\ (displacementField\ D\ p₀\ p₁\ z)\ ∂μ)\ +}
\leanline{6}{rho\ /\ 2\ *\ (‖p₀\ -\ p₁‖\ \textasciicircum{}\ 2\ +\ (eLpNorm\ (fun\ z\ ↦\ z\ -\ D.transport\ z)\ 2\ μ).toReal\ \textasciicircum{}\ 2)\ ≤}
\leanline{6}{(∫\ z,\ U\ p₁\ (D.transport\ z)\ ∂μ)\ -\ ∫\ z,\ U\ p₀\ z\ ∂μ}
\end{leancode}
\leanlabel{How the Lean proof follows the argument}
\begin{enumerate}


\item \texttt{Joint} abbreviates \texttt{PBGSO.\allowbreak{}Product\allowbreak{}Transport.\allowbreak{}Joint P Z}, the product with Euclidean norm, and \texttt{displacement\allowbreak{}Field D p\ensuremath{{}_{0}} p\ensuremath{{}_{1}} z} is \texttt{PBGSO.\allowbreak{}Product\allowbreak{}Transport.\allowbreak{}embed (p\ensuremath{{}_{1}} -\allowbreak{} p\ensuremath{{}_{0}}) (D.\allowbreak{}transport z -\allowbreak{} z)}.

\item \texttt{displacement\allowbreak{}Field\_\allowbreak{}mem\allowbreak{}Lp} uses \texttt{Proximal\allowbreak{}Langevin.\allowbreak{}mem\allowbreak{}Lp\_\allowbreak{}lipschitz\_\allowbreak{}observable} with the Lipschitz map \texttt{PBGSO.\allowbreak{}Product\allowbreak{}Transport.\allowbreak{}embed\_\allowbreak{}lipschitz}. \texttt{integrable\_\allowbreak{}inner\_\allowbreak{}mem\allowbreak{}Lp} converts both functions into elements of \texttt{Lp} and uses \texttt{L2.\allowbreak{}integrable\_\allowbreak{}inner}.

\item \texttt{potential\_\allowbreak{}tangent}: \texttt{h\allowbreak{}G} states \texttt{Has\allowbreak{}Gradient\allowbreak{}At} of the joint potential at \texttt{embed p\ensuremath{{}_{0}} z}. \texttt{hp} is the pointwise inequality, from \texttt{strong\allowbreak{}Convex\_\allowbreak{}tangent} and \texttt{With\allowbreak{}Lp.\allowbreak{}prod\_\allowbreak{}norm\_\allowbreak{}sq\_\allowbreak{}eq\_\allowbreak{}of\_\allowbreak{}L2}. \texttt{hb} integrates it by \texttt{integral\_\allowbreak{}mono}, and the integrals are split by linearity and \texttt{second\allowbreak{}Cost\_\allowbreak{}sq\_\allowbreak{}eq\_\allowbreak{}integral}.

\end{enumerate}
\end{leanbox}
\subsection{\texorpdfstring{The free-energy tangent inequality}{The free-energy tangent inequality}}


\textbf{Claim.} Let \(\mathcal F(s)=H(s)+\int U(p_s,\cdot)\,d\mu_s\) be the free energy of one datum along the displacement. Assume the hypotheses of the two previous steps, together with measurability of \(U(p,\cdot)\) for every \(p\). Then


\[\int\langle G,\Delta\rangle\,d\mu+\int\langle T(z)-z,\nabla\log\rho_0(z)\rangle\,\mu(dz)+\frac\lambda2\bigl(\|p_0-p_1\|^2+\|\mathrm{id}-T\|_{L^2(\mu)}^2\bigr)\le\mathcal F(1)-\mathcal F(0).\]


\textbf{Proof.} Since \(\mu_0=\mu\), \(p_0\) is the parameter at \(s=0\), and \(\int U(p_1,\cdot)\,d\mu_1=\int U(p_1,T(z))\,\mu(dz)\) by the pushforward formula,


\[\mathcal F(1)-\mathcal F(0)=H(1)-H(0)+\int U(p_1,T(z))\,\mu(dz)-\int U(p_0,z)\,\mu(dz).\]


The entropy tangent inequality bounds \(H(1)-H(0)\) from below and the potential tangent inequality bounds the potential difference from below; adding the two lower bounds gives the claim. \(\square\)

\begin{leanbox}{Lean: \texttt{freeEnergy\_\allowbreak{}tangent}}
\leansource{\texttt{freeEnergy\_\allowbreak{}tangent} \quad LatentDiffusion/\allowbreak{}RegularTransportFirstOrder.lean, lines 261–283}
\begin{leancode}
\leanline{0}{theorem\ freeEnergy\_tangent}
\leanline{4}{(D\ :\ RegularDisplacement\ (volume\ :\ Measure\ Z)\ μ\ d)}
\leanline{4}{(U\ :\ P\ →\ Z\ →\ ℝ)\ (p₀\ p₁\ :\ P)\ (rho\ :\ ℝ)\ (G\ :\ Z\ →\ Joint\ (P\ :=\ P)\ (Z\ :=\ Z))}
\leanline{4}{(hU\ :\ StrongConvexOn\ univ\ rho\ (fun\ w\ :\ Joint\ (P\ :=\ P)\ (Z\ :=\ Z)\ ↦\ U\ w.fst\ w.snd))}
\leanline{4}{(hG\ :\ ∀\ z,\ HasGradientAt\ (fun\ w\ :\ Joint\ (P\ :=\ P)\ (Z\ :=\ Z)\ ↦\ U\ w.fst\ w.snd)}
\leanline{6}{(G\ z)\ (PBGSO.ProductTransport.embed\ p₀\ z))}
\leanline{4}{(hG2\ :\ MemLp\ G\ 2\ μ)\ (hc\ :\ MemLp\ (fun\ z\ ↦\ z\ -\ D.transport\ z)\ 2\ μ)}
\leanline{4}{(hUmeas\ :\ ∀\ p,\ Measurable\ (U\ p))}
\leanline{4}{(hU0\ :\ Integrable\ (U\ p₀)\ μ)\ (hU1\ :\ Integrable\ (fun\ z\ ↦\ U\ p₁\ (D.transport\ z))\ μ)}
\leanline{4}{(hq\ :\ ContDiff\ ℝ\ 1\ (D.density\ 0))}
\leanline{4}{(htail\ :\ RegularFlowCalculus.NoBoundaryFlux\ (fun\ z\ ↦\ D.transport\ z\ -\ z)\ (D.density\ 0))\ :}
\leanline{4}{(∫\ z,\ inner\ ℝ\ (G\ z)\ (displacementField\ D\ p₀\ p₁\ z)\ ∂μ)\ +}
\leanline{6}{(∫\ z,\ inner\ ℝ\ (D.transport\ z\ -\ z)\ (gradient\ (fun\ z\ ↦\ Real.log\ (D.density\ 0\ z))\ z)\ ∂μ)\ +}
\leanline{6}{rho\ /\ 2\ *\ (‖p₀\ -\ p₁‖\ \textasciicircum{}\ 2\ +\ (eLpNorm\ (fun\ z\ ↦\ z\ -\ D.transport\ z)\ 2\ μ).toReal\ \textasciicircum{}\ 2)\ ≤}
\leanline{6}{freeEnergy\ D\ U\ p₀\ p₁\ 1\ -\ freeEnergy\ D\ U\ p₀\ p₁\ 0}
\end{leancode}
\leanlabel{How the Lean proof follows the argument}
\texttt{free\allowbreak{}Energy\_\allowbreak{}tangent} combines \texttt{hp} (\texttt{potential\_\allowbreak{}tangent}) and \texttt{he} (\texttt{D.\allowbreak{}entropy\_\allowbreak{}tangent}). \texttt{hmap} is the pushforward formula at \(s=1\) (\texttt{D.\allowbreak{}law\_\allowbreak{}one} and \texttt{integral\_\allowbreak{}map}), the \texttt{simp only} call evaluates \texttt{free\allowbreak{}Energy} at \(0\) and \(1\), and \texttt{linarith} adds the inequalities.
\end{leanbox}
\subsection{\texorpdfstring{Splitting the tangent pairing}{Splitting the tangent pairing}}


Write \(G(z)=(G_P(z),G_Z(z))\in P\times Z\) for the components of the gradient.




\textbf{Claim.} If \(G\) is square integrable under \(\mu\), so are \(G_P\) and \(G_Z\). If moreover \(z\mapsto z-T(z)\) and \(\nabla\log\rho_0\) are square integrable, then


\[\int\langle G,\Delta\rangle\,d\mu+\int\langle T-\mathrm{id},\nabla\log\rho_0\rangle\,d\mu=\Bigl\langle\int G_P\,d\mu,\;p_1-p_0\Bigr\rangle+\int\bigl\langle G_Z+\nabla\log\rho_0,\;T-\mathrm{id}\bigr\rangle\,d\mu.\]


\textbf{Proof.} The components satisfy \(\|G_P(z)\|\le\|G(z)\|\) and \(\|G_Z(z)\|\le\|G(z)\|\) and are measurable, so they are square integrable. By definition of the product inner product, \(\langle G(z),\Delta(z)\rangle=\langle G_P(z),p_1-p_0\rangle+\langle G_Z(z),T(z)-z\rangle\). Each term is integrable: the first because \(G_P\) is integrable (a square-integrable function is integrable for a probability measure) and \(y\mapsto\langle y,p_1-p_0\rangle\) is a continuous linear functional, the others by the integrability of inner products of square-integrable functions. A continuous linear functional commutes with the Bochner integral, so \(\int\langle G_P,p_1-p_0\rangle\,d\mu=\langle\int G_P\,d\mu,p_1-p_0\rangle\). Adding \(\int\langle T-\mathrm{id},\nabla\log\rho_0\rangle\,d\mu\) and combining the two latent integrals into one by linearity gives the claim. \(\square\)

\begin{leanbox}{Lean: \texttt{joint\_\allowbreak{}fst\_\allowbreak{}memLp}, \texttt{joint\_\allowbreak{}snd\_\allowbreak{}memLp}, \texttt{tangent\_\allowbreak{}pairing\_\allowbreak{}eq}}
\leansource{\texttt{joint\_\allowbreak{}fst\_\allowbreak{}memLp} \quad LatentDiffusion/\allowbreak{}RegularTransportFirstOrder.lean, lines 320–324}
\begin{leancode}
\leanline{0}{theorem\ joint\_fst\_memLp\ (G\ :\ Z\ →\ PBGSO.ProductTransport.Joint\ P\ Z)\ (hG\ :\ MemLp\ G\ 2\ μ)\ :}
\leanline{4}{MemLp\ (fun\ z\ ↦\ (G\ z).fst)\ 2\ μ}
\end{leancode}
\leansource{\texttt{joint\_\allowbreak{}snd\_\allowbreak{}memLp} \quad LatentDiffusion/\allowbreak{}RegularTransportFirstOrder.lean, lines 326–330}
\begin{leancode}
\leanline{0}{theorem\ joint\_snd\_memLp\ (G\ :\ Z\ →\ PBGSO.ProductTransport.Joint\ P\ Z)\ (hG\ :\ MemLp\ G\ 2\ μ)\ :}
\leanline{4}{MemLp\ (fun\ z\ ↦\ (G\ z).snd)\ 2\ μ}
\end{leancode}
\leansource{\texttt{tangent\_\allowbreak{}pairing\_\allowbreak{}eq} \quad LatentDiffusion/\allowbreak{}RegularTransportFirstOrder.lean, lines 349–383}
\begin{leancode}
\leanline{0}{theorem\ tangent\_pairing\_eq}
\leanline{4}{(D\ :\ RegularDisplacement\ (volume\ :\ Measure\ Z)\ μ\ d)\ (p₀\ p₁\ :\ P)}
\leanline{4}{(G\ :\ Z\ →\ PBGSO.ProductTransport.Joint\ P\ Z)\ (hG\ :\ MemLp\ G\ 2\ μ)}
\leanline{4}{(hc\ :\ MemLp\ (fun\ z\ ↦\ z\ -\ D.transport\ z)\ 2\ μ)}
\leanline{4}{(hs\ :\ MemLp\ (fun\ z\ ↦\ gradient\ (fun\ z\ ↦\ Real.log\ (D.density\ 0\ z))\ z)\ 2\ μ)\ :}
\leanline{4}{(∫\ z,\ inner\ ℝ\ (G\ z)\ (displacementField\ D\ p₀\ p₁\ z)\ ∂μ)\ +}
\leanline{6}{(∫\ z,\ inner\ ℝ\ (D.transport\ z\ -\ z)\ (gradient\ (fun\ z\ ↦\ Real.log\ (D.density\ 0\ z))\ z)\ ∂μ)\ =}
\leanline{4}{inner\ ℝ\ (∫\ z,\ (G\ z).fst\ ∂μ)\ (p₁\ -\ p₀)\ +}
\leanline{6}{∫\ z,\ inner\ ℝ\ ((G\ z).snd\ +\ gradient\ (fun\ z\ ↦\ Real.log\ (D.density\ 0\ z))\ z)}
\leanline{8}{(D.transport\ z\ -\ z)\ ∂μ}
\end{leancode}
\leanlabel{How the Lean proof follows the argument}
\begin{enumerate}


\item \texttt{joint\_\allowbreak{}fst\_\allowbreak{}mem\allowbreak{}Lp} and \texttt{joint\_\allowbreak{}snd\_\allowbreak{}mem\allowbreak{}Lp} use \texttt{h\allowbreak{}G.\allowbreak{}mono}, monotonicity of square integrability under a pointwise norm bound, with the pointwise bounds \texttt{With\allowbreak{}Lp.\allowbreak{}norm\_\allowbreak{}fst\_\allowbreak{}le} and \texttt{With\allowbreak{}Lp.\allowbreak{}norm\_\allowbreak{}snd\_\allowbreak{}le}.

\item \texttt{tangent\_\allowbreak{}pairing\_\allowbreak{}eq}: \texttt{hpar\allowbreak{}Eq} is the commutation of the integral with the inner product, from \texttt{integral\_\allowbreak{}inner}; the integrability facts \texttt{hpar}, \texttt{hip} and \texttt{his} come from \texttt{integrable\_\allowbreak{}inner\_\allowbreak{}mem\allowbreak{}Lp}. After expanding the product inner product (\texttt{With\allowbreak{}Lp.\allowbreak{}prod\_\allowbreak{}inner\_\allowbreak{}apply}), \texttt{integral\_\allowbreak{}add} splits and recombines the integrals, and \texttt{ring} regroups the terms.

\end{enumerate}
\end{leanbox}
\subsection{\texorpdfstring{Completing the square over the data}{Completing the square over the data}}


\textbf{Claim.} Let \(M\ge1\), \(\lambda>0\), \(w\in P\), and for \(m=1,\dots,M\) let \(\gamma_m,\kappa_m,b_m,c_m\in\mathbb R\) and \(g_m\in P\) satisfy


\[\gamma_m+\langle g_m,w\rangle+\kappa_m+\frac\lambda2\bigl(\|w\|^2+c_m^2\bigr)\le0\qquad\text{and}\qquad-\kappa_m\le b_mc_m.\]


Then


\[2\lambda\cdot\frac1M\sum_m\gamma_m\le\Bigl\|\frac1M\sum_mg_m\Bigr\|^2+\frac1M\sum_mb_m^2.\]


\textbf{Proof.} Multiply the first hypothesis by \(2\lambda>0\) and use the second:


\[2\lambda\gamma_m\le-2\lambda\langle g_m,w\rangle-\lambda^2\|w\|^2+\bigl(2\lambda b_mc_m-\lambda^2c_m^2\bigr)\le-2\lambda\langle g_m,w\rangle-\lambda^2\|w\|^2+b_m^2,\]


where the last inequality is \((b_m-\lambda c_m)^2\ge0\). Average over \(m\) and write \(\bar g=M^{-1}\sum_mg_m\); since \(M^{-1}\cdot M=1\),


\[2\lambda\cdot\frac1M\sum_m\gamma_m\le-2\lambda\langle\bar g,w\rangle-\lambda^2\|w\|^2+\frac1M\sum_mb_m^2.\]


Finally \(0\le\|\bar g+\lambda w\|^2=\|\bar g\|^2+2\lambda\langle\bar g,w\rangle+\lambda^2\|w\|^2\), so \(-2\lambda\langle\bar g,w\rangle-\lambda^2\|w\|^2\le\|\bar g\|^2\). \(\square\)




The parameter vectors are averaged before the square is completed. Averaging \(M\) separate one-datum inequalities would produce \(M^{-1}\sum_m\|g_m\|^2\), which is larger than \(\|\bar g\|^2\) and is not the Fisher information of the paper.

\begin{leanbox}{Lean: \texttt{normalized\_\allowbreak{}extendedLSI\_\allowbreak{}of\_\allowbreak{}tangents}}
\leansource{\texttt{normalized\_\allowbreak{}extendedLSI\_\allowbreak{}of\_\allowbreak{}tangents} \quad LatentDiffusion/\allowbreak{}RegularTransportFirstOrder.lean, lines 388–422}
\begin{leancode}
\leanline{0}{theorem\ normalized\_extendedLSI\_of\_tangents}
\leanline{4}{\{M\ :\ ℕ\}\ (hM\ :\ 0\ <\ M)\ (rho\ :\ ℝ)\ (hrho\ :\ 0\ <\ rho)\ (dp\ :\ P)}
\leanline{4}{(gap\ pair\ b\ c\ :\ Fin\ M\ →\ ℝ)\ (g\ :\ Fin\ M\ →\ P)}
\leanline{4}{(htangent\ :\ ∀\ m,\ gap\ m\ +\ inner\ ℝ\ (g\ m)\ dp\ +\ pair\ m\ +}
\leanline{6}{rho\ /\ 2\ *\ (‖dp‖\ \textasciicircum{}\ 2\ +\ c\ m\ \textasciicircum{}\ 2)\ ≤\ 0)}
\leanline{4}{(hpair\ :\ ∀\ m,\ -pair\ m\ ≤\ b\ m\ *\ c\ m)\ :}
\leanline{4}{2\ *\ rho\ *\ ((M\ :\ ℝ)⁻¹\ *\ ∑\ m,\ gap\ m)\ ≤}
\leanline{6}{‖(M\ :\ ℝ)⁻¹\ •\ ∑\ m,\ g\ m‖\ \textasciicircum{}\ 2\ +\ (M\ :\ ℝ)⁻¹\ *\ ∑\ m,\ b\ m\ \textasciicircum{}\ 2}
\end{leancode}
\leanlabel{How the Lean proof follows the argument}
In \texttt{normalized\_\allowbreak{}extended\allowbreak{}LSI\_\allowbreak{}of\_\allowbreak{}tangents} the variables \texttt{gap}, \texttt{pair}, \texttt{b}, \texttt{c}, \texttt{g} and \texttt{dp} are \(\gamma\), \(\kappa\), \(b\), \(c\), \(g\) and \(w\). \texttt{hpoint} is the per-datum inequality, proved by \texttt{nlinarith} from the two hypotheses and \texttt{sq\_\allowbreak{}nonneg (b m -\allowbreak{} rho * c m)}. \texttt{hsum} sums and multiplies by \(w=M^{-1}\), \texttt{hmean} moves the average inside the inner product, \texttt{hquad} is \(\|\bar g+\lambda w\|^2\ge0\) expanded with \texttt{norm\_\allowbreak{}add\_\allowbreak{}sq\_\allowbreak{}real}, and \texttt{hconst} is \(M^{-1}\cdot M\lambda^2\|w\|^2=\lambda^2\|w\|^2\). The final \texttt{nlinarith} combines them.
\end{leanbox}
\subsection{\texorpdfstring{Transport assumptions for the data set}{Transport assumptions for the data set}}


\textbf{Definition.} Fix parameters \(p_0,p_1\in P\), \(\lambda\in\mathbb R\) and \(M\) data. The transport hypotheses combine the paper's strong convexity assumption with the extra density, integrability and map assumptions R4.




For each datum \(m\), the extra conditions are the density/\allowbreak{}map data and integrability requirements below; strong convexity is assumption A3. Supply:



\begin{enumerate}


\item a probability law \(Q_m\) on \(Z\) and regular displacement data \(D_m\) from \(Q_m\) with Lebesgue reference measure (transport \(T_m\), densities \(\rho_{m,s}\));

\item a potential \(U_m:P\times Z\to\mathbb R\) that is \(\lambda\)-strongly convex on \(P\times Z\), with \(U_m(p,\cdot)\) measurable for every \(p\);

\item integrability under \(Q_m\) of \(U_m(p_0,\cdot)\), of \(z\mapsto U_m(p_1,T_m(z))\) and of \(z\mapsto U_m(p_s,T_{m,s}(z))\) for \(s\in[0,1]\);

\item gradients \(G_m(z)\) of \(U_m\) at \((p_0,z)\), with \(G_m\) square integrable under \(Q_m\);

\item square integrability of \(z\mapsto z-T_m(z)\) and of \(\nabla\log\rho_{m,0}\) under \(Q_m\);

\item continuous differentiability of \(\rho_{m,0}\), and the no-boundary-flux condition for \((T_m-\mathrm{id},\rho_{m,0})\).

\end{enumerate}



From these data define the averaged free energy \(\mathcal F(s)=M^{-1}\sum_m\mathcal F_m(s)\), the parameter force \(G=M^{-1}\sum_m\int G_{m,P}\,dQ_m\), the latent scores \(S_m=G_{m,Z}+\nabla\log\rho_{m,0}\), and the Fisher information


\[I=\|G\|^2+\frac1M\sum_m\int\|S_m\|^2\,dQ_m.\]


Since \(G_{m,Z}(z)\) is the latent component of the gradient of \(U_m\) at \((p_0,z)\), \(S_m=\nabla_z(\log\rho_{m,0}+U_m(p_0,\cdot))\) is the latent score of the summary.




\textbf{Claim.} Each \(S_m\) is square integrable under \(Q_m\).




\textbf{Proof.} \(G_{m,Z}\) is square integrable by the previous step, \(\nabla\log\rho_{m,0}\) by (5), and a sum of square-integrable functions is square integrable. \(\square\)

\begin{leanbox}{Lean: \texttt{DatasetRegularity}, \texttt{freeEnergy}, \texttt{parameterForce} and 3 more}
\leansource{\texttt{DatasetRegularity} \quad LatentDiffusion/\allowbreak{}RegularTransportFirstOrder.lean, lines 528–553}
\begin{leancode}
\leanline{0}{structure\ DatasetRegularity\ (M\ d\ :\ ℕ)\ (p₀\ p₁\ :\ P)\ (rho\ :\ ℝ)\ where}
\leanline{2}{law\ :\ Fin\ M\ →\ ProbabilityMeasure\ Z}
\leanline{2}{displacement\ :\ (m\ :\ Fin\ M)\ →\ RegularDisplacement\ volume\ (law\ m\ :\ Measure\ Z)\ d}
\leanline{2}{potential\ :\ Fin\ M\ →\ P\ →\ Z\ →\ ℝ}
\leanline{2}{potential\_strong\ :\ ∀\ m,\ StrongConvexOn\ univ\ rho}
\leanline{4}{(fun\ w\ :\ PBGSO.ProductTransport.Joint\ P\ Z\ ↦\ potential\ m\ w.fst\ w.snd)}
\leanline{2}{potential\_measurable\ :\ ∀\ m\ p,\ Measurable\ (potential\ m\ p)}
\leanline{2}{potential\_integrable\_initial\ :\ ∀\ m,\ Integrable\ (potential\ m\ p₀)\ (law\ m\ :\ Measure\ Z)}
\leanline{2}{potential\_integrable\_final\ :\ ∀\ m,}
\leanline{4}{Integrable\ (fun\ z\ ↦\ potential\ m\ p₁\ ((displacement\ m).transport\ z))\ (law\ m\ :\ Measure\ Z)}
\leanline{2}{potential\_integrable\_path\ :\ ∀\ m\ s,\ s\ ∈\ Icc\ (0\ :\ ℝ)\ 1\ →}
\leanline{4}{Integrable\ (fun\ z\ ↦\ potential\ m\ (parameterInterpolate\ p₀\ p₁\ s)}
\leanline{6}{(interpolate\ (displacement\ m).transport\ s\ z))\ (law\ m\ :\ Measure\ Z)}
\leanline{2}{potential\_gradient\ :\ Fin\ M\ →\ Z\ →\ PBGSO.ProductTransport.Joint\ P\ Z}
\leanline{2}{potential\_hasGradientAt\ :\ ∀\ m\ z,\ HasGradientAt}
\leanline{4}{(fun\ w\ :\ PBGSO.ProductTransport.Joint\ P\ Z\ ↦\ potential\ m\ w.fst\ w.snd)}
\leanline{4}{(potential\_gradient\ m\ z)\ (PBGSO.ProductTransport.embed\ p₀\ z)}
\leanline{2}{potential\_gradient\_memLp\ :\ ∀\ m,\ MemLp\ (potential\_gradient\ m)\ 2\ (law\ m\ :\ Measure\ Z)}
\leanline{2}{displacement\_memLp\ :\ ∀\ m,}
\leanline{4}{MemLp\ (fun\ z\ ↦\ z\ -\ (displacement\ m).transport\ z)\ 2\ (law\ m\ :\ Measure\ Z)}
\leanline{2}{score\_memLp\ :\ ∀\ m,}
\leanline{4}{MemLp\ (fun\ z\ ↦\ gradient\ (fun\ z\ ↦\ Real.log\ ((displacement\ m).density\ 0\ z))\ z)}
\leanline{6}{2\ (law\ m\ :\ Measure\ Z)}
\leanline{2}{density\_C1\ :\ ∀\ m,\ ContDiff\ ℝ\ 1\ ((displacement\ m).density\ 0)}
\leanline{2}{integrable\_tails\ :\ ∀\ m,\ RegularFlowCalculus.NoBoundaryFlux}
\leanline{4}{(fun\ z\ ↦\ (displacement\ m).transport\ z\ -\ z)\ ((displacement\ m).density\ 0)}
\end{leancode}
\leansource{\texttt{freeEnergy} \quad LatentDiffusion/\allowbreak{}RegularTransportFirstOrder.lean, lines 560–561}
\begin{leancode}
\leanline{0}{def\ freeEnergy\ (s\ :\ ℝ)\ :\ ℝ\ :=\ (M\ :\ ℝ)⁻¹\ *\ ∑\ m,}
\leanline{2}{RegularTransport.freeEnergy\ (R.displacement\ m)\ (R.potential\ m)\ p₀\ p₁\ s}
\end{leancode}
\leansource{\texttt{parameterForce} \quad LatentDiffusion/\allowbreak{}RegularTransportFirstOrder.lean, lines 563–564}
\begin{leancode}
\leanline{0}{def\ parameterForce\ :\ P\ :=\ (M\ :\ ℝ)⁻¹\ •\ ∑\ m,}
\leanline{2}{∫\ z,\ (R.potential\_gradient\ m\ z).fst\ ∂(R.law\ m\ :\ Measure\ Z)}
\end{leancode}
\leansource{\texttt{latentScore} \quad LatentDiffusion/\allowbreak{}RegularTransportFirstOrder.lean, lines 566–568}
\begin{leancode}
\leanline{0}{def\ latentScore\ (m\ :\ Fin\ M)\ (z\ :\ Z)\ :\ Z\ :=}
\leanline{2}{(R.potential\_gradient\ m\ z).snd\ +}
\leanline{4}{gradient\ (fun\ z\ ↦\ Real.log\ ((R.displacement\ m).density\ 0\ z))\ z}
\end{leancode}
\leansource{\texttt{fisherInformation} \quad LatentDiffusion/\allowbreak{}RegularTransportFirstOrder.lean, lines 570–571}
\begin{leancode}
\leanline{0}{def\ fisherInformation\ :\ ℝ\ :=\ ‖R.parameterForce‖\ \textasciicircum{}\ 2\ +\ (M\ :\ ℝ)⁻¹\ *\ ∑\ m,}
\leanline{2}{∫\ z,\ ‖R.latentScore\ m\ z‖\ \textasciicircum{}\ 2\ ∂(R.law\ m\ :\ Measure\ Z)}
\end{leancode}
\leansource{\texttt{latentScore\_\allowbreak{}memLp} \quad LatentDiffusion/\allowbreak{}RegularTransportFirstOrder.lean, lines 582–585}
\begin{leancode}
\leanline{0}{theorem\ latentScore\_memLp\ (m\ :\ Fin\ M)\ :\ MemLp\ (R.latentScore\ m)\ 2\ (R.law\ m\ :\ Measure\ Z)}
\end{leancode}
\leanlabel{How the Lean proof follows the argument}
\texttt{Dataset\allowbreak{}Regularity M d p\ensuremath{{}_{0}} p\ensuremath{{}_{1}} rho} is a structure with the fields \texttt{law}, \texttt{displacement}, \texttt{potential}, \texttt{potential\_\allowbreak{}strong}, \texttt{potential\_\allowbreak{}measurable}, \texttt{potential\_\allowbreak{}integrable\_\allowbreak{}initial}, \texttt{potential\_\allowbreak{}integrable\_\allowbreak{}final}, \texttt{potential\_\allowbreak{}integrable\_\allowbreak{}path}, \texttt{potential\_\allowbreak{}gradient}, \texttt{potential\_\allowbreak{}has\allowbreak{}Gradient\allowbreak{}At}, \texttt{potential\_\allowbreak{}gradient\_\allowbreak{}mem\allowbreak{}Lp}, \texttt{displacement\_\allowbreak{}mem\allowbreak{}Lp}, \texttt{score\_\allowbreak{}mem\allowbreak{}Lp}, \texttt{density\_\allowbreak{}C1} and \texttt{integrable\_\allowbreak{}tails}, in the order of conditions (1)–(6). \texttt{free\allowbreak{}Energy}, \texttt{parameter\allowbreak{}Force}, \texttt{latent\allowbreak{}Score} and \texttt{fisher\allowbreak{}Information} are the four derived quantities; in \texttt{fisher\allowbreak{}Information} the norm of the averaged force is squared, as in the definition of \(I\). \texttt{latent\allowbreak{}Score\_\allowbreak{}mem\allowbreak{}Lp} adds \texttt{joint\_\allowbreak{}snd\_\allowbreak{}mem\allowbreak{}Lp} of the gradient to \texttt{score\_\allowbreak{}mem\allowbreak{}Lp}.
\end{leanbox}
\subsection{\texorpdfstring{The extended log-Sobolev inequality}{The extended log-Sobolev inequality}}


\textbf{Claim.} Let \(M\ge1\), \(\lambda>0\), and assume the transport hypotheses for the data set. Then


\[2\lambda\bigl(\mathcal F(0)-\mathcal F(1)\bigr)\le I.\]


\textbf{Proof.} For each datum \(m\) put


\[\begin{aligned}\gamma_m&=\mathcal F_m(0)-\mathcal F_m(1),&g_m&=\int G_{m,P}\,dQ_m,\\\kappa_m&=\int\langle S_m,T_m-\mathrm{id}\rangle\,dQ_m,\\b_m&=\|S_m\|_{L^2(Q_m)},&c_m&=\|\mathrm{id}-T_m\|_{L^2(Q_m)}.\end{aligned}\]


\emph{Tangent inequality.} The free-energy tangent inequality for datum \(m\), with its left side rewritten by the splitting step, reads


\[\langle g_m,p_1-p_0\rangle+\kappa_m+\frac\lambda2\bigl(\|p_0-p_1\|^2+c_m^2\bigr)\le\mathcal F_m(1)-\mathcal F_m(0),\]


that is, \(\gamma_m+\langle g_m,w\rangle+\kappa_m+\frac\lambda2(\|w\|^2+c_m^2)\le0\) with \(w=p_1-p_0\), since \(\|p_0-p_1\|=\|w\|\).




\emph{Pairing bound.} By the Cauchy-Schwarz bound of the previous section applied to \(-S_m\) and \(T_m-\mathrm{id}\), \(-\kappa_m=\int\langle-S_m,T_m-\mathrm{id}\rangle\,dQ_m\le b_mc_m\); here \(\|T_m-\mathrm{id}\|_{L^2}=\|\mathrm{id}-T_m\|_{L^2}\).




\emph{Completion of the square.} The previous step gives \(2\lambda M^{-1}\sum_m\gamma_m\le\|M^{-1}\sum_mg_m\|^2+M^{-1}\sum_mb_m^2\). The left side is \(2\lambda(\mathcal F(0)-\mathcal F(1))\). On the right, \(M^{-1}\sum_mg_m=G\) and \(b_m^2=\int\|S_m\|^2\,dQ_m\), so the right side is \(I\). \(\square\)

\begin{leanbox}{Lean: \texttt{extendedLSI}}
\leansource{\texttt{extendedLSI} \quad LatentDiffusion/\allowbreak{}RegularTransportFirstOrder.lean, lines 661–718}
\begin{leancode}
\leanline{0}{theorem\ extendedLSI\ (hM\ :\ 0\ <\ M)\ (hrho\ :\ 0\ <\ rho)\ :}
\leanline{4}{2\ *\ rho\ *\ (R.freeEnergy\ 0\ -\ R.freeEnergy\ 1)\ ≤\ R.fisherInformation}
\end{leancode}
\leanlabel{How the Lean proof follows the argument}
In \texttt{extended\allowbreak{}LSI} the local definitions \texttt{g}, \texttt{pair}, \texttt{b}, \texttt{c} and \texttt{gap} are \(g_m\), \(\kappa_m\), \(b_m\), \(c_m\) and \(\gamma_m\); \texttt{b} and \texttt{c} are real conversions of \texttt{e\allowbreak{}Lp\allowbreak{}Norm … 2}.



\begin{enumerate}


\item \texttt{htangent} applies \texttt{free\allowbreak{}Energy\_\allowbreak{}tangent} to the fields of \texttt{R} for datum \texttt{m}, rewrites its left side with \texttt{tangent\_\allowbreak{}pairing\_\allowbreak{}eq}, and uses \texttt{norm\_\allowbreak{}sub\_\allowbreak{}rev}.

\item \texttt{hpair} applies \texttt{integral\_\allowbreak{}inner\_\allowbreak{}le\_\allowbreak{}second\allowbreak{}Cost} to \texttt{-\allowbreak{}(R.\allowbreak{}latent\allowbreak{}Score m)} and the displacement, using \texttt{R.\allowbreak{}latent\allowbreak{}Score\_\allowbreak{}mem\allowbreak{}Lp}; \texttt{hvnorm} identifies the \(L^2\) norms of \(T-\mathrm{id}\) and \(\mathrm{id}-T\) with \texttt{e\allowbreak{}Lp\allowbreak{}Norm\_\allowbreak{}neg}.

\item \texttt{h} is \texttt{normalized\_\allowbreak{}extended\allowbreak{}LSI\_\allowbreak{}of\_\allowbreak{}tangents}, \texttt{hsq} replaces \(b_m^2\) by the integral through \texttt{second\allowbreak{}Cost\_\allowbreak{}sq\_\allowbreak{}eq\_\allowbreak{}integral}, and the final \texttt{simpa} unfolds \texttt{free\allowbreak{}Energy}, \texttt{fisher\allowbreak{}Information} and \texttt{parameter\allowbreak{}Force}.

\end{enumerate}
\end{leanbox}
\subsection{\texorpdfstring{Displacement convexity over the data set and the Talagrand inequality}{Displacement convexity over the data set and the Talagrand inequality}}


The normalized product distance between \((p_0,(Q_m)_m)\) and \((p_1,(Q_m\circ T_m^{-1})_m)\) is given by


\[d^2=\|p_0-p_1\|^2+\frac1M\sum_mW_2\bigl(Q_m,Q_m\circ T_m^{-1}\bigr)^2\ge0.\]


\textbf{Claim.} Let \(M\ge1\) and \(\lambda\ge0\). (1) For every \(s\in[0,1]\),


\[\mathcal F(s)\le(1-s)\mathcal F(0)+s\mathcal F(1)-\frac\lambda2s(1-s)\,d^2.\]


(2) If the endpoint minimizes along the displacement, \(\mathcal F(1)\le\mathcal F(s)\) for all \(s\in[0,1]\), then


\[\frac\lambda2\,d^2\le\mathcal F(0)-\mathcal F(1).\]


\textbf{Proof.} (1) For each \(m\), the strong displacement convexity of the previous section gives the inequality for \(\mathcal F_m\) with the bracket \(\|p_0-p_1\|^2+W_2(Q_m,Q_m\circ T_m^{-1})^2\). Averaging over \(m\) and using \(M^{-1}\sum_m\|p_0-p_1\|^2=\|p_0-p_1\|^2\) (here \(M\ge1\)) gives (1).




(2) Put \(f(s)=\mathcal F(1-s)\) for \(s\in[0,1]\). Then \(f(0)=\mathcal F(1)\le f(s)\) by hypothesis, and (1) at the time \(1-s\in[0,1]\) gives


\[f(s)\le s\,\mathcal F(0)+(1-s)\mathcal F(1)-\frac\lambda2(1-s)s\,d^2=(1-s)f(0)+s\,\mathcal F(0)-\frac\lambda2s(1-s)\,d^2.\]


The elementary quadratic-growth fact of the previous section (if \(f(0)\le f(s)\le(1-s)f(0)+sf_1-\frac\lambda2s(1-s)D^2\) on \([0,1]\), then \(\frac\lambda2D^2\le f_1-f(0)\)) applies with \(f_1=\mathcal F(0)\) and \(D=d\). \(\square\)




In the convergence theorem, the endpoint is the posterior optimizer, whose minimality along the displacement follows from the Gibbs variational principle; (2) then bounds the distance to the optimizer by the free-energy gap.

\begin{leanbox}{Lean: \texttt{distanceSq}, \texttt{distanceSq\_\allowbreak{}nonneg}, \texttt{displacement\_\allowbreak{}strongConvex} and 1 more}
\leansource{\texttt{distanceSq} \quad LatentDiffusion/\allowbreak{}RegularTransportFirstOrder.lean, lines 722–723}
\begin{leancode}
\leanline{0}{def\ distanceSq\ :\ ℝ\ :=\ ‖p₀\ -\ p₁‖\ \textasciicircum{}\ 2\ +\ (M\ :\ ℝ)⁻¹\ *\ ∑\ m,}
\leanline{2}{ULA.wasserstein2Real\ (R.law\ m\ :\ Measure\ Z)\ ((R.displacement\ m).law\ 1)\ \textasciicircum{}\ 2}
\end{leancode}
\leansource{\texttt{distanceSq\_\allowbreak{}nonneg} \quad LatentDiffusion/\allowbreak{}RegularTransportFirstOrder.lean, lines 725–727}
\begin{leancode}
\leanline{0}{theorem\ distanceSq\_nonneg\ :\ 0\ ≤\ R.distanceSq}
\end{leancode}
\leansource{\texttt{displacement\_\allowbreak{}strongConvex} \quad LatentDiffusion/\allowbreak{}RegularTransportFirstOrder.lean, lines 731–747}
\begin{leancode}
\leanline{0}{theorem\ displacement\_strongConvex\ (hM\ :\ 0\ <\ M)\ (hrho\ :\ 0\ ≤\ rho)}
\leanline{4}{(s\ :\ ℝ)\ (hs\ :\ s\ ∈\ Icc\ (0\ :\ ℝ)\ 1)\ :}
\leanline{4}{R.freeEnergy\ s\ ≤\ (1\ -\ s)\ *\ R.freeEnergy\ 0\ +\ s\ *\ R.freeEnergy\ 1-}
\leanline{6}{rho\ /\ 2\ *\ s\ *\ (1\ -\ s)\ *\ R.distanceSq}
\end{leancode}
\leansource{\texttt{talagrand\_\allowbreak{}final} \quad LatentDiffusion/\allowbreak{}RegularTransportFirstOrder.lean, lines 762–782}
\begin{leancode}
\leanline{0}{theorem\ talagrand\_final\ (hM\ :\ 0\ <\ M)\ (hrho\ :\ 0\ ≤\ rho)}
\leanline{4}{(hmin\ :\ ∀\ s\ ∈\ Icc\ (0\ :\ ℝ)\ 1,\ R.freeEnergy\ 1\ ≤\ R.freeEnergy\ s)\ :}
\leanline{4}{rho\ /\ 2\ *\ R.distanceSq\ ≤\ R.freeEnergy\ 0\ -\ R.freeEnergy\ 1}
\end{leancode}
\leanlabel{How the Lean proof follows the argument}
\begin{enumerate}


\item \texttt{distance\allowbreak{}Sq} uses \texttt{ULA.\allowbreak{}wasserstein2\allowbreak{}Real}, the real value of the infimum over couplings, with the target law \texttt{(R.\allowbreak{}displacement m).\allowbreak{}law 1}; \texttt{distance\allowbreak{}Sq\_\allowbreak{}nonneg} is \texttt{add\_\allowbreak{}nonneg} of squares.

\item \texttt{displacement\_\allowbreak{}strong\allowbreak{}Convex} sums \texttt{free\allowbreak{}Energy\_\allowbreak{}displacement\_\allowbreak{}strong\allowbreak{}Convex} over the data, multiplies by \(M^{-1}\), and rearranges with \texttt{field\_\allowbreak{}simp}, which uses \(M\ne0\).

\item \texttt{talagrand\_\allowbreak{}final}: \texttt{hmin'} and \texttt{hconv} are the two hypotheses for the reversed function \texttt{fun s \ensuremath{\mapsto} R.\allowbreak{}free\allowbreak{}Energy (1 -\allowbreak{} s)}, obtained from \texttt{hmin} and \texttt{R.\allowbreak{}displacement\_\allowbreak{}strong\allowbreak{}Convex} at \(1-s\) (the identity of the two right sides is checked by \texttt{ring}). \texttt{Source\allowbreak{}Results.\allowbreak{}geodesic\_\allowbreak{}quadratic\_\allowbreak{}growth} is the elementary fact, and \texttt{Real.\allowbreak{}sq\_\allowbreak{}sqrt} converts \((\sqrt{d^2})^2\) back to \(d^2\).

\end{enumerate}
\end{leanbox}
\section{\texorpdfstring{Full-batch Euler error against a stationary posterior reference}{Full-batch Euler error against a stationary posterior reference}}


Use the scaled cloud with \(K>0\), differentiable jointly \(m\)-strongly convex potentials, and \(L\)-Lipschitz gradients. Take \(m,h>0\), \(mh\le1\), \(hL^2\le m\), and put \(q=1-mh/2\). All random states lie in \(L^2(P)\), whose norm is \((\mathbb E\|X\|^2)^{1/2}\), for a probability measure \(P\). Supply a continuous stationary frozen-posterior integral solution \(S_t\), driven by the stated coordinate Brownian processes. Its parameter block at time zero is the constant \(a_\star\). Assume the initial per-particle parameter scores are independent and their expectations sum to zero. These are explicit stochastic and reference assumptions.

\subsection{\texorpdfstring{Independent score fluctuations yield the sampling term}{Independent score fluctuations yield the sampling term}}


\textbf{Claim.} Let \(Q\) be extraction of the parameter block, with operator norm at most one. Put \(V=\sum_i\mathbb E\|s_i-\mathbb Es_i\|^2\), \(G=\|F_{\rm frozen}(S_0)\|_{L^2}\), and let \(W\) denote the physical Brownian increment scale. The coupled full-batch trajectory satisfies

\[\|QX_n-a_\star\|_{L^2}\le q^n\|X_0-S_0\|_{L^2}+\frac2m\frac{\sqrt V}{K}+\frac Lm\left(hG+\frac43W\sqrt h\right).\]

\textbf{Proof.} Write \(s_i=(s_i-\mathbb Es_i)+\mathbb Es_i\). The sum of the means vanishes. Expanding the squared norm of their average leaves diagonal terms plus cross inner products. Independence and centring make every cross expectation zero. Consequently the mean-square residual is exactly \(K^{-2}V\), whose nonnegative square root is \(\sqrt V/K\). Stationarity preserves this residual norm at every reference time. The frozen integral equation supplies the reference segments used in the numerical comparison; its drift norm is \(G\), and Brownian increments supply the square-root-time scale \(W\). Apply the equilibrium-reference Euler comparison with this residual bound. The parameter block of \(S_t\) is constant. Finally \(\|Q(X_n-S_{nh})\|_{L^2}\le\|X_n-S_{nh}\|_{L^2}\) transfers the cloud estimate to the parameter estimate. \(\square\)

\begin{leanbox}{Lean: \texttt{full\_\allowbreak{}batch\_\allowbreak{}parameter\_\allowbreak{}bound}}
\leansource{\texttt{full\_\allowbreak{}batch\_\allowbreak{}parameter\_\allowbreak{}bound} \quad LatentDiffusion/\allowbreak{}RegularStochastic.lean, lines 482–528}
\begin{leancode}
\leanline{0}{theorem\ full\_batch\_parameter\_bound}
\leanline{4}{(hK\ :\ 0\ <\ K)\ (U\ :\ Fin\ K\ →\ Cloud.Joint\ dθ\ dz\ →\ ℝ)}
\leanline{4}{(m\ h\ :\ ℝ)\ (L\ :\ ℝ≥0)\ (hdiff\ :\ ∀\ i,\ Differentiable\ ℝ\ (U\ i))}
\leanline{4}{(hsc\ :\ ∀\ i,\ StrongConvexOn\ Set.univ\ m\ (U\ i))}
\leanline{4}{(hLip\ :\ ∀\ i,\ LipschitzWith\ L\ (gradient\ (U\ i)))\ (hs\ :\ StableStep\ m\ L\ h)}
\leanline{4}{(B\ :\ (Fin\ K\ ×\ Fin\ dz)\ →\ ℝ≥0\ →\ Ω\ →\ ℝ)}
\leanline{4}{(hB\ :\ ∀\ i,\ IsPreBrownianReal\ (B\ i)\ P)}
\leanline{4}{(S\ :\ FrozenPosteriorSolution\ P\ (Cloud.meanForce\_regular\ hK\ U\ m\ L\ hdiff\ hsc\ hLip)\ B\ hB)}
\leanline{4}{(hind\ :\ iIndepFun\ (fun\ i\ ω\ =>\ scoreLp\ P\ U\ L\ hLip\ (S.path\ 0)\ i\ ω)\ P)}
\leanline{4}{(hopt\ :\ ∑\ i,\ (∫\ ω,\ scoreLp\ P\ U\ L\ hLip\ (S.path\ 0)\ i\ ω\ ∂P)\ =\ 0)}
\leanline{4}{(x₀\ :\ Lp\ (Cloud.State\ K\ dθ\ dz)\ 2\ P)\ (astar\ :\ Parameter\ dθ)}
\leanline{4}{(hstar\ :\ parameterOutput\ P\ (S.path\ 0)\ =\ constantL2\ P\ astar)\ (n\ :\ ℕ)\ :}
\leanline{4}{‖parameterOutput\ P}
\leanline{6}{(eulerTrajectory\ (particleForce\ P\ hK\ U\ m\ L\ hdiff\ hsc\ hLip)\ h}
\leanline{8}{(fun\ k\ =>\ physicalNoise\ P\ B\ hB\ ((k\ :\ ℝ)*h+h)\ -\ physicalNoise\ P\ B\ hB\ ((k\ :\ ℝ)*h))\ x₀\ n)\ -}
\leanline{6}{constantL2\ P\ astar‖\ ≤}
\leanline{6}{(1\ -\ m*h/2)\textasciicircum{}n\ *\ ‖x₀\ -\ S.path\ 0‖\ +}
\leanline{6}{(2/m)\ *\ Real.sqrt\ ((1/(K\ :\ ℝ))\textasciicircum{}2\ *\ scoreVarianceBudget\ P\ U\ L\ hLip\ (S.path\ 0))\ +}
\leanline{6}{(L\ :\ ℝ)/m\ *\ (h\ *\ ‖frozenForce\ P\ (Cloud.meanForce\_regular\ hK\ U\ m\ L\ hdiff\ hsc\ hLip)}
\leanline{8}{(S.path\ 0)‖\ +\ (4/3\ :\ ℝ)*physicalNoiseScale\ K\ dz*Real.sqrt\ h)}
\end{leancode}
\leanlabel{How the Lean proof follows the argument}
\texttt{full\_\allowbreak{}batch\_\allowbreak{}parameter\_\allowbreak{}bound} defines the residual size \texttt{b}. The block \texttt{hres} uses \texttt{residual\_\allowbreak{}norm\_\allowbreak{}eq\_\allowbreak{}of\_\allowbreak{}law}, \texttt{lifted\_\allowbreak{}mean\allowbreak{}Force\_\allowbreak{}residual} and \texttt{independent\_\allowbreak{}average\_\allowbreak{}norm\_\allowbreak{}sq} to prove the variance identity. It then calls \texttt{equilibrium\_\allowbreak{}reference\_\allowbreak{}error}. The equality \texttt{ht} identifies the reference parameter with \texttt{constant\allowbreak{}L2} of the optimizer; \texttt{parameter\allowbreak{}Output\_\allowbreak{}norm\_\allowbreak{}le} proves that the projection cannot enlarge the error. The displayed norm is therefore an actual root-mean-square parameter error.
\end{leanbox}
\section{\texorpdfstring{Evidence differentiation and centred posterior scores}{Evidence differentiation and centred posterior scores}}


For a reference measure \(\nu\), define \(Z_i(a)=\int e^{-U_i(a,z)}\,d\nu(z)\), \(\pi_{i,a}(dz)=Z_i(a)^{-1}e^{-U_i(a,z)}\nu(dz)\), and the averaged log evidence \(\ell(a)=M^{-1}\sum_i\log Z_i(a)\), where \(M>0\). The normalizers must be positive. Assume the densities are measurable at the point under consideration and almost everywhere strongly measurable for every parameter in some neighbourhood of that point; at the point they are integrable and their weighted gradients are almost everywhere strongly measurable; and on one common parameter neighbourhood the potentials are differentiable almost everywhere and \(e^{-U_i(b,z)}\|\nabla_bU_i(b,z)\|\) is bounded by an integrable function of \(z\). These domination hypotheses justify differentiation under the integral sign.

\subsection{\texorpdfstring{The derivative of a normalizer is its negative weighted score}{The derivative of a normalizer is its negative weighted score}}


\textbf{Claim.} Under the stated local domination assumptions, \[DZ_i(a)[v]=-\int e^{-U_i(a,z)}\langle\nabla_aU_i(a,z),v\rangle\,d\nu(z).\].



\textbf{Proof.} For almost every \(z\), the chain rule differentiates \(a\mapsto e^{-U_i(a,z)}\) to the displayed integrand. The operator norm of that derivative is bounded by the common integrable dominating function throughout the neighbourhood. The dominated differentiation theorem therefore exchanges differentiation and integration. The weighted gradient is Bochner integrable by the same bound. Continuity of the inner-product functional permits it to be moved outside the Bochner integral, giving the claimed derivative for every direction \(v\). \(\square\)

\begin{leanbox}{Lean: \texttt{normalizer\_\allowbreak{}hasFDerivAt}}
\leansource{\texttt{normalizer\_\allowbreak{}hasFDerivAt} \quad LatentDiffusion/\allowbreak{}RegularEvidence.lean, lines 68–100}
\begin{leancode}
\leanline{0}{theorem\ normalizer\_hasFDerivAt}
\leanline{4}{(ν\ :\ Measure\ Z)\ (U\ :\ A\ →\ Z\ →\ ℝ)\ (a\ :\ A)}
\leanline{4}{(data\ :\ LocalNormalizerRegularity\ ν\ U\ a)\ :}
\leanline{4}{HasFDerivAt\ (normalizer\ ν\ U)}
\leanline{6}{(-(InnerProductSpace.toDual\ ℝ\ A)\ (unnormalizedScore\ ν\ U\ a))\ a}
\end{leancode}
\leanlabel{How the Lean proof follows the argument}
\texttt{normalizer\_\allowbreak{}has\allowbreak{}FDeriv\allowbreak{}At} defines the derivative-valued integrand \texttt{D}, checks its measurable and dominated hypotheses, and applies \texttt{has\allowbreak{}FDeriv\allowbreak{}At\_\allowbreak{}integral\_\allowbreak{}of\_\allowbreak{}dominated\_\allowbreak{}of\_\allowbreak{}fderiv\_\allowbreak{}le}. The final \texttt{heq} moves the continuous linear inner-product map through the integral. \texttt{Has\allowbreak{}FDeriv\allowbreak{}At} expresses the full Fréchet derivative, not merely directional differentiability.
\end{leanbox}
\subsection{\texorpdfstring{A local evidence maximizer has zero total posterior score}{A local evidence maximizer has zero total posterior score}}


\textbf{Claim.} If \(a\) is a local maximizer of \(\ell\), then \(\sum_i\int\nabla_aU_i(a,z)\,d\pi_{i,a}(z)=0\).



\textbf{Proof.} Positivity of \(Z_i(a)\) permits the logarithmic chain rule. Dividing its derivative by \(Z_i(a)\) converts the weighted score integral into a posterior expectation, so \(\nabla\ell(a)=-M^{-1}\sum_i\int\nabla_aU_i(a,z)\,d\pi_{i,a}(z)\). A differentiable real function on the whole parameter space has zero derivative at a local maximum: for each direction its restriction to a real line has a local maximum at zero, whose derivative is zero. Thus the averaged score vanishes. Since \(M>0\), multiplication by \(M\) yields the unaveraged identity. Posterior integrability is furnished by the normalizer domination assumptions, rather than inferred from formal cancellation. \(\square\)

\begin{leanbox}{Lean: \texttt{sum\_\allowbreak{}posterior\_\allowbreak{}parameter\_\allowbreak{}gradient\_\allowbreak{}eq\_\allowbreak{}zero}}
\leansource{\texttt{sum\_\allowbreak{}posterior\_\allowbreak{}parameter\_\allowbreak{}gradient\_\allowbreak{}eq\_\allowbreak{}zero} \quad LatentDiffusion/\allowbreak{}RegularEvidence.lean, lines 185–193}
\begin{leancode}
\leanline{0}{theorem\ sum\_posterior\_parameter\_gradient\_eq\_zero\ \{M\ :\ ℕ\}\ (hM\ :\ 0\ <\ M)}
\leanline{4}{(ν\ :\ Measure\ Z)\ (U\ :\ Fin\ M\ →\ A\ →\ Z\ →\ ℝ)\ (a\ :\ A)}
\leanline{4}{(data\ :\ ∀\ i,\ LocalNormalizerRegularity\ ν\ (U\ i)\ a)}
\leanline{4}{(hmeas\ :\ ∀\ i,\ Measurable\ (fun\ z\ ↦\ Real.exp\ (-U\ i\ a\ z)))}
\leanline{4}{(hZ\ :\ ∀\ i,\ 0\ <\ normalizer\ ν\ (U\ i)\ a)}
\leanline{4}{(hmax\ :\ IsLocalMax\ (logEvidence\ ν\ U)\ a)\ :}
\leanline{4}{∑\ i,\ (∫\ z,\ gradient\ (fun\ b\ ↦\ U\ i\ b\ z)\ a\ ∂posterior\ ν\ (U\ i)\ a)\ =\ 0}
\end{leancode}
\leanlabel{How the Lean proof follows the argument}
\texttt{sum\_\allowbreak{}posterior\_\allowbreak{}parameter\_\allowbreak{}gradient\_\allowbreak{}eq\_\allowbreak{}zero} applies \texttt{averaged\_\allowbreak{}posterior\_\allowbreak{}parameter\_\allowbreak{}gradient\_\allowbreak{}eq\_\allowbreak{}zero} to the differentiable evidence. Its \texttt{smul\_\allowbreak{}eq\_\allowbreak{}zero} step removes the scalar reciprocal; \texttt{inv\_\allowbreak{}ne\_\allowbreak{}zero} and \texttt{h\allowbreak{}M} rule out a zero scalar. The local-maximality input is \texttt{hmax}; a zero-score hypothesis is not assumed.
\end{leanbox}
\section{\texorpdfstring{Synchronous mean-field contraction and optimizer comparison}{Synchronous mean-field contraction and optimizer comparison}}


Let the heterogeneous particle potentials be differentiable and jointly \(m\)-strongly convex, with \(L\)-Lipschitz gradients, and let the particle count be positive. Consider a continuous \(L^2(P)\) mean-field integral solution \(Y_t\) and a continuous stationary frozen-posterior solution \(S_t\) driven by the same physical Brownian path. Assume the posterior scores have zero total expectation, and the initial parameter of \(Y\) and the reference parameter of \(S\) are deterministic. The statement holds for any real \(m\); it gives decay when \(m>0\). Existence of a new solution is not asserted by a comparison of supplied solutions.

\subsection{\texorpdfstring{Cancellation of common noise gives exponential parameter control}{Cancellation of common noise gives exponential parameter control}}


\textbf{Claim.} If the reference parameter is \(a_\star\), then for every \(t\ge0\), \(\|QY_t-a_\star\|_{L^2}\le e^{-mt}\|Y_0-S_0\|_{L^2}\).



\textbf{Proof.} Apply the parameter projection to the integral equation of \(Y\). The physical noise has zero parameter block, so its parameter remains the deterministic initial value minus the integral of the expected parameter force. The zero-score identity makes the frozen reference an admissible mean-field solution with constant parameter \(a_\star\). Subtract their two equations. Their common additive noise cancels. Strong monotonicity, including the deterministic-parameter expectation identities, gives the energy inequality \(\frac d{dt}\|Y_t-S_t\|_{L^2}^2\le-2m\|Y_t-S_t\|_{L^2}^2\). Multiply by \(e^{2mt}\); the derivative of the product is nonpositive. Integrating from zero to \(t\) gives \(\|Y_t-S_t\|_{L^2}^2\le e^{-2mt}\|Y_0-S_0\|_{L^2}^2\). Nonnegative square roots and the norm-one parameter projection finish the proof. Continuity and the supplied integral equations provide the time-integrability needed for this energy argument. \(\square\)

\begin{leanbox}{Lean: \texttt{meanField\_\allowbreak{}optimizer\_\allowbreak{}convergence}}
\leansource{\texttt{meanField\_\allowbreak{}optimizer\_\allowbreak{}convergence} \quad LatentDiffusion/\allowbreak{}RegularMeanField.lean, lines 560–587}
\begin{leancode}
\leanline{0}{theorem\ meanField\_optimizer\_convergence}
\leanline{4}{(hK\ :\ 0\ <\ K)\ (U\ :\ Fin\ K\ →\ Cloud.Joint\ dθ\ dz\ →\ ℝ)}
\leanline{4}{(m\ :\ ℝ)\ (L\ :\ ℝ≥0)\ (hdiff\ :\ ∀\ i,\ Differentiable\ ℝ\ (U\ i))}
\leanline{4}{(hsc\ :\ ∀\ i,\ StrongConvexOn\ Set.univ\ m\ (U\ i))}
\leanline{4}{(hLip\ :\ ∀\ i,\ LipschitzWith\ L\ (gradient\ (U\ i)))}
\leanline{4}{(B\ :\ (Fin\ K\ ×\ Fin\ dz)\ →\ ℝ≥0\ →\ Ω\ →\ ℝ)}
\leanline{4}{(hB\ :\ ∀\ i,\ IsPreBrownianReal\ (B\ i)\ P)}
\leanline{4}{(S\ :\ FrozenPosteriorSolution\ P\ (Cloud.meanForce\_regular\ hK\ U\ m\ L\ hdiff\ hsc\ hLip)\ B\ hB)}
\leanline{4}{(hopt\ :\ ∑\ i,\ (∫\ ω,\ scoreLp\ P\ U\ L\ hLip\ (S.path\ 0)\ i\ ω\ ∂P)=0)}
\leanline{4}{(astar\ :\ Parameter\ dθ)\ (hstar\ :\ parameterOutput\ P\ (S.path\ 0)=constantL2\ P\ astar)}
\leanline{4}{(Y\ :\ ℝ\ →\ Lp\ (Cloud.State\ K\ dθ\ dz)\ 2\ P)\ (hY\ :\ Continuous\ Y)}
\leanline{4}{(hYe\ :\ ∀\ t,\ 0\ ≤\ t\ →\ Y\ t-Y\ 0\ =}
\leanline{6}{-(∫\ r\ in\ 0..t,\ force\ P\ (Cloud.meanForce\_regular\ hK\ U\ m\ L\ hdiff\ hsc\ hLip)\ (Y\ r))+}
\leanline{8}{physicalNoise\ P\ B\ hB\ t)}
\leanline{4}{(a₀\ :\ Parameter\ dθ)\ (ha₀\ :\ parameterOutput\ P\ (Y\ 0)=constantL2\ P\ a₀)}
\leanline{4}{(t\ :\ ℝ)\ (ht\ :\ 0\ ≤\ t)\ :}
\leanline{4}{‖parameterOutput\ P\ (Y\ t)-constantL2\ P\ astar‖\ ≤\ Real.exp\ (-m*t)*‖Y\ 0-S.path\ 0‖}
\end{leancode}
\leanlabel{How the Lean proof follows the argument}
\texttt{mean\allowbreak{}Field\_\allowbreak{}optimizer\_\allowbreak{}convergence} first invokes \texttt{mean\allowbreak{}Field\_\allowbreak{}parameter\_\allowbreak{}deterministic} to establish \texttt{hpa}. The block \texttt{hpb} identifies the reference parameter; \texttt{frozen\_\allowbreak{}is\_\allowbreak{}mean\allowbreak{}Field} verifies its equation using \texttt{hopt}. The actual synchronous energy estimate is \texttt{mean\allowbreak{}Field\_\allowbreak{}contraction}. The final two inequalities apply \texttt{parameter\allowbreak{}Output\_\allowbreak{}norm\_\allowbreak{}le} to its norm estimate. In particular the proof does not replace stochastic independence or the reference equation by an assumed error bound.
\end{leanbox}
\section{\texorpdfstring{Exponential convergence of the classical flow}{Exponential convergence of the classical flow}}


This section proves the paper's continuous-time convergence theorem (Theorem 1 and its appendix form) for classical solutions of the joint gradient flow. The flow consists of a parameter path \(a_t\) and positive densities \(q_t^1,\dots,q_t^M\) solving


\[\dot a_t=-\frac1M\sum_{m=1}^M\int\nabla_aU_m(a_t,z)\,q_t^m(z)\,dz,\qquad\partial_tq_t^m=\nabla\cdot\bigl(q_t^m\nabla(\log q_t^m+U_m(a_t,\cdot))\bigr),\]


The flow also satisfies extra assumptions R1–R4: local domination, integrable spatial tails, and, at each time, regular displacement data from the current laws to the laws of a fixed comparison state \(a_\star\).




Existence of such a classical solution is not constructed; the estimates hold for every solution in this class. Write \(E(t)=\mathcal E(a_t,q_t)-\mathcal E(a_\star,q_\star)\) for the energy gap to the comparison state and \(d_t\) for the normalized product distance


\[d_t^2=\|a_t-a_\star\|^2+\frac1M\sum_{m=1}^MW_2(q_t^m,\pi_{a_\star}^m)^2.\]


When \(a_\star\) maximizes the evidence and the comparison densities are the posteriors \(\pi_{a_\star}^m\), we prove


\[d_t\le\sqrt{\frac{2E(t)}\lambda}\le\sqrt{\frac{2E(0)}\lambda}\,e^{-\lambda t},\qquad\|a_t-a_\star\|\le\sqrt{\frac{2E(0)}\lambda}\,e^{-\lambda t}\qquad(t\ge0).\]


The ingredients are the de Bruijn identity \(E'=-I\), the extended log-Sobolev inequality \(2\lambda E\le I\), Grönwall's inequality, and the Talagrand inequality \(\frac\lambda2d_t^2\le E(t)\), whose hypothesis (minimality of the posterior state) follows from the Gibbs variational principle. The final theorem constructs \(a_\star\) as the unique maximizer of the evidence of the original potentials, so the only remaining hypotheses are the extra assumptions defining the flow.

\subsection{\texorpdfstring{The Gibbs variational lower bound}{The Gibbs variational lower bound}}


\textbf{Claim 1.} If \(q:Z\to\mathbb R\) is measurable and nonnegative and \(Q(dz)=q(z)\,dz\), then \(\int f\,dQ=\int q\,f\,dz\) for every function \(f\) with values in a normed space, with both sides read as Bochner integrals (equal to \(0\) for nonintegrable functions).




\textbf{Claim 2.} Let \(Q\) be a probability law with a measurable, strictly positive Lebesgue density \(q\). Let \(V:Z\to\mathbb R\) be such that \(e^{-V}\) is integrable with \(Z_V=\int e^{-V}dz>0\), and such that \(\log q\) and \(V\) are \(Q\)-integrable. Then


\[-\log Z_V\le\int\log q\,dQ+\int V\,dQ.\]


\textbf{Proof.} Claim 1 is the definition of the integral against a measure with density, together with the fact that \(q\) is real and nonnegative. For Claim 2, define \(r(z)=\dfrac{e^{-V(z)}/Z_V}{q(z)}>0\), so that \(q\,r=e^{-V}/Z_V\). By Claim 1, \(r\) is \(Q\)-integrable if and only if \(qr\) is Lebesgue integrable, which holds, and \(\int r\,dQ=\int e^{-V}dz/Z_V=1\).




For \(x>0\), \(\log x\le x-1\). With \(x=r(z)\) and \(\log r=-V-\log Z_V-\log q\) this gives, for every \(z\),


\[1-r(z)\le\log q(z)+V(z)+\log Z_V.\]


Both sides are \(Q\)-integrable. Integrating against the probability measure \(Q\) gives \(1-1\le\int\log q\,dQ+\int V\,dQ+\log Z_V\). \(\square\)




Claim 2 is the nonnegativity of the relative entropy of \(Q\) with respect to the Gibbs law \(e^{-V}dz/Z_V\), proved directly from the pointwise inequality.

\begin{leanbox}{Lean: \texttt{integral\_\allowbreak{}density}, \texttt{gibbs\_\allowbreak{}variational\_\allowbreak{}lower\_\allowbreak{}bound}}
\leansource{\texttt{integral\_\allowbreak{}density} \quad LatentDiffusion/\allowbreak{}RegularConvergence.lean, lines 30–36}
\begin{leancode}
\leanline{0}{theorem\ integral\_density\ \{H\ :\ Type*\}\ [NormedAddCommGroup\ H]\ [NormedSpace\ ℝ\ H]}
\leanline{4}{(q\ :\ Z\ →\ ℝ)\ (hq\ :\ Measurable\ q)\ (hpos\ :\ ∀\ z,\ 0\ ≤\ q\ z)\ (f\ :\ Z\ →\ H)\ :}
\leanline{4}{(∫\ z,\ f\ z\ ∂(volume.withDensity\ (fun\ z\ ↦\ ENNReal.ofReal\ (q\ z))))\ =}
\leanline{6}{∫\ z,\ q\ z\ •\ f\ z}
\end{leancode}
\leansource{\texttt{gibbs\_\allowbreak{}variational\_\allowbreak{}lower\_\allowbreak{}bound} \quad LatentDiffusion/\allowbreak{}RegularConvergence.lean, lines 40–82}
\begin{leancode}
\leanline{0}{theorem\ gibbs\_variational\_lower\_bound\ (Q\ :\ ProbabilityMeasure\ Z)\ (q\ V\ :\ Z\ →\ ℝ)}
\leanline{4}{(hq\ :\ Measurable\ q)\ (hpos\ :\ ∀\ z,\ 0\ <\ q\ z)}
\leanline{4}{(hlaw\ :\ (Q\ :\ Measure\ Z)\ =\ volume.withDensity\ (fun\ z\ ↦\ ENNReal.ofReal\ (q\ z)))}
\leanline{4}{(hweight\ :\ Integrable\ (fun\ z\ ↦\ Real.exp\ (-V\ z))\ volume)}
\leanline{4}{(hZ\ :\ 0\ <\ ∫\ z,\ Real.exp\ (-V\ z))}
\leanline{4}{(hlog\ :\ Integrable\ (fun\ z\ ↦\ Real.log\ (q\ z))\ (Q\ :\ Measure\ Z))}
\leanline{4}{(hV\ :\ Integrable\ V\ (Q\ :\ Measure\ Z))\ :}
\leanline{4}{-(Real.log\ (∫\ z,\ Real.exp\ (-V\ z)))\ ≤}
\leanline{6}{(∫\ z,\ Real.log\ (q\ z)\ ∂(Q\ :\ Measure\ Z))\ +\ ∫\ z,\ V\ z\ ∂(Q\ :\ Measure\ Z)}
\end{leancode}
\leanlabel{How the Lean proof follows the argument}
\begin{enumerate}


\item \texttt{integral\_\allowbreak{}density} is \texttt{integral\_\allowbreak{}with\allowbreak{}Density\_\allowbreak{}eq\_\allowbreak{}integral\_\allowbreak{}to\allowbreak{}Real\_\allowbreak{}smul} with \texttt{ENNReal.\allowbreak{}to\allowbreak{}Real\_\allowbreak{}of\allowbreak{}Real}, which removes the conversion of \(q\) to an extended number.

\item \texttt{gibbs\_\allowbreak{}variational\_\allowbreak{}lower\_\allowbreak{}bound}: \texttt{hlaw} states \(Q=q\,dz\) as \texttt{volume.\allowbreak{}with\allowbreak{}Density (fun z \ensuremath{\mapsto} ENNReal.\allowbreak{}of\allowbreak{}Real (q z))}. \texttt{hprod} is \(qr=e^{-V}/Z_V\), \texttt{hr} and \texttt{hrint} are integrability and \(\int r\,dQ=1\), and \texttt{hpoint} is the pointwise inequality, from \texttt{Real.\allowbreak{}log\_\allowbreak{}le\_\allowbreak{}sub\_\allowbreak{}one\_\allowbreak{}of\_\allowbreak{}pos} and \texttt{Real.\allowbreak{}log\_\allowbreak{}div}. \texttt{hbound} integrates it by \texttt{integral\_\allowbreak{}mono}; the remaining rewrites evaluate the integrals of sums and of constants (\texttt{measure\_\allowbreak{}univ}), and \texttt{linarith} concludes.

\end{enumerate}
\end{leanbox}
\subsection{\texorpdfstring{Minimality of the posterior state along a displacement}{Minimality of the posterior state along a displacement}}


Let \(R\) collect the transport assumptions for the data set (previous section) from \((p_0,(Q_m)_m)=(a,(Q_m)_m)\) to the parameter \(p_1=b\), with potentials \(U_m(p,z)\), displacement laws \(\mu_{m,s}\) with densities \(\rho_{m,s}\), and averaged free energy \(\mathcal F(s)\). Define the evidence \(\ell(p)=M^{-1}\sum_m\log\int e^{-U_m(p,z)}dz\).




\textbf{Claim.} (1) For \(s\in[0,1]\), \(\log\rho_{m,s}\) and \(U_m(p_s,\cdot)\) are \(\mu_{m,s}\)-integrable. (2) If \(e^{-U_m(p,\cdot)}\) is integrable with positive integral for all \(m\) and \(p\), then \(-\ell(p_s)\le\mathcal F(s)\) for \(s\in[0,1]\). (3) If the endpoint densities are the posteriors at \(b\), \(\rho_{m,1}=e^{-U_m(b,\cdot)}/\int e^{-U_m(b,w)}dw\), and the normalizers at \(b\) are positive, then \(\mathcal F(1)=-\ell(b)\). (4) Under the hypotheses of (2) and (3), if \(b\) maximizes \(\ell\), then \(\mathcal F(1)\le\mathcal F(s)\) for every \(s\in[0,1]\).




\textbf{Proof.} (1) By the pushforward formula, integrability of \(f\) under \(\mu_{m,s}=Q_m\circ T_{m,s}^{-1}\) is integrability of \(f\circ T_{m,s}\) under \(Q_m\); for \(f=\log\rho_{m,s}\) and \(f=U_m(p_s,\cdot)\) this is one of the extra transport assumptions.




(2) Fix \(m\). The law \(\mu_{m,s}\) is a probability measure with positive measurable density \(\rho_{m,s}\). By the previous step with \(Q=\mu_{m,s}\), \(q=\rho_{m,s}\) and \(V=U_m(p_s,\cdot)\), and by (1),


\[-\log\int e^{-U_m(p_s,z)}dz\le\int\log\rho_{m,s}\,d\mu_{m,s}+\int U_m(p_s,\cdot)\,d\mu_{m,s}=\mathcal F_m(s).\]


Averaging over \(m\) gives \(-\ell(p_s)\le\mathcal F(s)\).




(3) At \(s=1\), \(p_1=b\) and \(\log\rho_{m,1}(z)+U_m(b,z)=-\log\int e^{-U_m(b,w)}dw\) is constant in \(z\). Its integral against the probability measure \(\mu_{m,1}\) is the same constant, so \(\mathcal F_m(1)=-\log\int e^{-U_m(b,w)}dw\), and averaging gives \(\mathcal F(1)=-\ell(b)\).




(4) By (3), maximality of \(b\) and (2): \(\mathcal F(1)=-\ell(b)\le-\ell(p_s)\le\mathcal F(s)\). \(\square\)

\begin{leanbox}{Lean: \texttt{log\_\allowbreak{}integrable}, \texttt{potential\_\allowbreak{}integrable}, \texttt{LatentDiffusion.\allowbreak{}RegularConvergence.\allowbreak{}RegularDataset.\allowbreak{}evidence} and 3 more}
\leansource{\texttt{log\_\allowbreak{}integrable} \quad LatentDiffusion/\allowbreak{}RegularConvergence.lean, lines 88–95}
\begin{leancode}
\leanline{0}{theorem\ log\_integrable\ (i\ :\ Fin\ M)\ (s\ :\ ℝ)\ (hs\ :\ s\ ∈\ Icc\ (0\ :\ ℝ)\ 1)\ :}
\leanline{4}{Integrable\ (fun\ z\ ↦\ Real.log\ ((R.displacement\ i).density\ s\ z))}
\leanline{6}{((R.displacement\ i).law\ s)}
\end{leancode}
\leansource{\texttt{potential\_\allowbreak{}integrable} \quad LatentDiffusion/\allowbreak{}RegularConvergence.lean, lines 97–102}
\begin{leancode}
\leanline{0}{theorem\ potential\_integrable\ (i\ :\ Fin\ M)\ (s\ :\ ℝ)\ (hs\ :\ s\ ∈\ Icc\ (0\ :\ ℝ)\ 1)\ :}
\leanline{4}{Integrable\ (R.potential\ i\ (parameterInterpolate\ a\ b\ s))\ ((R.displacement\ i).law\ s)}
\end{leancode}
\leansource{\texttt{LatentDiffusion.\allowbreak{}RegularConvergence.\allowbreak{}RegularDataset.\allowbreak{}evidence} \quad LatentDiffusion/\allowbreak{}RegularConvergence.lean, lines 104–105}
\begin{leancode}
\leanline{0}{def\ evidence\ (p\ :\ A)\ :\ ℝ\ :=\ (M\ :\ ℝ)⁻¹\ *\ ∑\ i,}
\leanline{2}{Real.log\ (∫\ z,\ Real.exp\ (-R.potential\ i\ p\ z))}
\end{leancode}
\leansource{\texttt{energy\_\allowbreak{}lower\_\allowbreak{}bound} \quad LatentDiffusion/\allowbreak{}RegularConvergence.lean, lines 107–126}
\begin{leancode}
\leanline{0}{theorem\ energy\_lower\_bound}
\leanline{4}{(hweight\ :\ ∀\ i\ p,\ Integrable\ (fun\ z\ ↦\ Real.exp\ (-R.potential\ i\ p\ z))\ volume)}
\leanline{4}{(hZ\ :\ ∀\ i\ p,\ 0\ <\ ∫\ z,\ Real.exp\ (-R.potential\ i\ p\ z))}
\leanline{4}{(s\ :\ ℝ)\ (hs\ :\ s\ ∈\ Icc\ (0\ :\ ℝ)\ 1)\ :}
\leanline{4}{-evidence\ R\ (parameterInterpolate\ a\ b\ s)\ ≤\ R.freeEnergy\ s}
\end{leancode}
\leansource{\texttt{posterior\_\allowbreak{}energy} \quad LatentDiffusion/\allowbreak{}RegularConvergence.lean, lines 128–151}
\begin{leancode}
\leanline{0}{theorem\ posterior\_energy}
\leanline{4}{(hZ\ :\ ∀\ i,\ 0\ <\ ∫\ z,\ Real.exp\ (-R.potential\ i\ b\ z))}
\leanline{4}{(hpost\ :\ ∀\ i\ z,\ (R.displacement\ i).density\ 1\ z\ =}
\leanline{6}{Real.exp\ (-R.potential\ i\ b\ z)/(∫\ w,\ Real.exp\ (-R.potential\ i\ b\ w)))\ :}
\leanline{4}{R.freeEnergy\ 1\ =\ -evidence\ R\ b}
\end{leancode}
\leansource{\texttt{posterior\_\allowbreak{}minimal} \quad LatentDiffusion/\allowbreak{}RegularConvergence.lean, lines 153–162}
\begin{leancode}
\leanline{0}{theorem\ posterior\_minimal}
\leanline{4}{(hweight\ :\ ∀\ i\ p,\ Integrable\ (fun\ z\ ↦\ Real.exp\ (-R.potential\ i\ p\ z))\ volume)}
\leanline{4}{(hZ\ :\ ∀\ i\ p,\ 0\ <\ ∫\ z,\ Real.exp\ (-R.potential\ i\ p\ z))}
\leanline{4}{(hmax\ :\ ∀\ p,\ evidence\ R\ p\ ≤\ evidence\ R\ b)}
\leanline{4}{(hpost\ :\ ∀\ i\ z,\ (R.displacement\ i).density\ 1\ z\ =}
\leanline{6}{Real.exp\ (-R.potential\ i\ b\ z)/(∫\ w,\ Real.exp\ (-R.potential\ i\ b\ w)))\ :}
\leanline{4}{∀\ s\ ∈\ Icc\ (0\ :\ ℝ)\ 1,\ R.freeEnergy\ 1\ ≤\ R.freeEnergy\ s}
\end{leancode}
\leanlabel{How the Lean proof follows the argument}
\begin{enumerate}


\item \texttt{log\_\allowbreak{}integrable} and \texttt{potential\_\allowbreak{}integrable} use \texttt{integrable\_\allowbreak{}map\_\allowbreak{}measure}, which transfers integrability through the pushforward, with the fields \texttt{integrable\_\allowbreak{}log} and \texttt{potential\_\allowbreak{}integrable\_\allowbreak{}path}.

\item \texttt{Regular\allowbreak{}Dataset.\allowbreak{}evidence} is \(\ell\). In \texttt{energy\_\allowbreak{}lower\_\allowbreak{}bound}, \texttt{hi} applies \texttt{gibbs\_\allowbreak{}variational\_\allowbreak{}lower\_\allowbreak{}bound} to the probability measure \texttt{Q} built from the law \(\mu_{m,s}\); the averaged inequality \texttt{hh} is \texttt{Finset.\allowbreak{}sum\_\allowbreak{}le\_\allowbreak{}sum} multiplied by \(M^{-1}\).

\item \texttt{posterior\_\allowbreak{}energy}: \texttt{hp} is the constancy of \(\log\rho_{m,1}+U_m(b,\cdot)\), from \texttt{hpost}, \texttt{Real.\allowbreak{}log\_\allowbreak{}div} and \texttt{Real.\allowbreak{}log\_\allowbreak{}exp}; the integral of the constant is evaluated by \texttt{simp}.

\item \texttt{posterior\_\allowbreak{}minimal} chains \texttt{posterior\_\allowbreak{}energy}, \texttt{neg\_\allowbreak{}le\_\allowbreak{}neg (hmax \_\allowbreak{})} and \texttt{energy\_\allowbreak{}lower\_\allowbreak{}bound}.

\end{enumerate}
\end{leanbox}
\subsection{\texorpdfstring{Classical flows}{Classical flows}}




\textbf{Definition.} Let \(A\) be a complete real inner product space of parameters, \(Z\) a finite-dimensional latent space, \(U_m:A\times Z\to\mathbb R\) potentials, \(\lambda\in\mathbb R\) and \(a_\star\in A\) a comparison parameter. A classical flow consists of a parameter path \(a_t\) with derivatives \(\dot a_t\), densities \(q_t^m\) with rates \(\dot q_t^m\), comparison densities \(q_\star^m\), and, for every time \(t\), the transport assumptions for the data set from \((a_t,(q_t^m\,dz)_m)\) to \(a_\star\), such that:



\begin{enumerate}


\item the data at time \(t\) use the potentials \(U_m\), start from the densities \(q_t^m\) and end at the densities \(q_\star^m\);

\item \(a\) has derivative \(\dot a_t\) at every \(t\), and \(b\mapsto U_m(b,z)\) is differentiable everywhere;

\item \(q_t^m>0\), and \(q_t^m\) and \(U_m(a_t,\cdot)\) are continuously differentiable;

\item the gradient fields supplied by the extra transport assumptions at time \(t\) are the actual gradients: parameter component \(\nabla_aU_m(a_t,z)\) and latent component \(\nabla_zU_m(a_t,z)\);

\item \(q_t^m\nabla_aU_m(a_t,\cdot)\) is integrable, and local domination holds at every \(t\) for the energy integrand;

\item for \(t\ge0\), the parameter equation \(\dot a_t=-G(a_t,q_t)\) and the Fokker–Planck equations \(\dot q_t^m=\nabla\cdot(q_t^m\nabla(\log q_t^m+U_m(a_t,\cdot)))\) almost everywhere;

\item the no-boundary-flux condition for the Fokker–Planck flux and the test \(\log q_t^m+U_m(a_t,\cdot)+1\), at every \(t\).

\end{enumerate}





With \(\mathcal E\) and \(I\) of the flow-calculus section, define the energy \(\mathcal E(t)=\mathcal E(a_t,q_t)\), the Fisher information \(I(t)=I(a_t,q_t)\), the target energy \(\mathcal E_\star=\mathcal E(a_\star,q_\star)\), and the evidence \(\ell(a)=M^{-1}\sum_m\log\int e^{-U_m(a,z)}dz\).




All conditions concern densities, derivatives, integrability and the equations themselves. Neither the energy identity, nor a log-Sobolev or transport inequality, nor a convergence estimate is a condition.

\begin{leanbox}{Lean: \texttt{ClassicalFlow}, \texttt{energy}, \texttt{fisher} and 2 more}
\leansource{\texttt{ClassicalFlow} \quad LatentDiffusion/\allowbreak{}RegularConvergence.lean, lines 170–201}
\begin{leancode}
\leanline{0}{structure\ ClassicalFlow\ (M\ d\ :\ ℕ)\ (rho\ :\ ℝ)\ (astar\ :\ A)}
\leanline{4}{(U\ :\ Fin\ M\ →\ A\ →\ Z\ →\ ℝ)\ where}
\leanline{2}{parameter\ :\ ℝ\ →\ A}
\leanline{2}{parameterDerivative\ :\ ℝ\ →\ A}
\leanline{2}{density\ :\ Fin\ M\ →\ ℝ\ →\ Z\ →\ ℝ}
\leanline{2}{densityDerivative\ :\ Fin\ M\ →\ ℝ\ →\ Z\ →\ ℝ}
\leanline{2}{targetDensity\ :\ Fin\ M\ →\ Z\ →\ ℝ}
\leanline{2}{transport\ :\ (t\ :\ ℝ)\ →\ DatasetRegularity\ (Z\ :=\ Z)\ M\ d\ (parameter\ t)\ astar\ rho}
\leanline{2}{potential\_eq\ :\ ∀\ t,\ (transport\ t).potential\ =\ U}
\leanline{2}{density\_eq\ :\ ∀\ t\ i,\ ((transport\ t).displacement\ i).density\ 0\ =\ density\ i\ t}
\leanline{2}{target\_density\_eq\ :\ ∀\ t\ i,\ ((transport\ t).displacement\ i).density\ 1\ =\ targetDensity\ i}
\leanline{2}{parameter\_differentiable\ :\ ∀\ t,\ HasDerivAt\ parameter\ (parameterDerivative\ t)\ t}
\leanline{2}{potential\_parameter\_differentiable\ :\ ∀\ i\ z\ a,\ DifferentiableAt\ ℝ\ (fun\ b\ ↦\ U\ i\ b\ z)\ a}
\leanline{2}{positive\ :\ ∀\ i\ t\ z,\ 0\ <\ density\ i\ t\ z}
\leanline{2}{density\_C1\ :\ ∀\ i\ t,\ ContDiff\ ℝ\ 1\ (density\ i\ t)}
\leanline{2}{potential\_C1\ :\ ∀\ i\ t,\ ContDiff\ ℝ\ 1\ (U\ i\ (parameter\ t))}
\leanline{2}{parameter\_gradient\ :\ ∀\ t\ i\ z,\ HasGradientAt\ (fun\ a\ ↦\ U\ i\ a\ z)}
\leanline{4}{(((transport\ t).potential\_gradient\ i\ z).fst)\ (parameter\ t)}
\leanline{2}{latent\_gradient\ :\ ∀\ t\ i\ z,\ HasGradientAt\ (U\ i\ (parameter\ t))}
\leanline{4}{(((transport\ t).potential\_gradient\ i\ z).snd)\ z}
\leanline{2}{parameter\_gradient\_integrable\ :\ ∀\ t\ i,}
\leanline{4}{Integrable\ (fun\ z\ ↦\ density\ i\ t\ z\ •\ gradient\ (fun\ a\ ↦\ U\ i\ a\ z)\ (parameter\ t))\ volume}
\leanline{2}{domination\ :\ ∀\ t\ i,\ LocalEnergyDomination\ volume\ (density\ i)\ (densityDerivative\ i)}
\leanline{4}{(fun\ u\ ↦\ U\ i\ (parameter\ u))}
\leanline{4}{(fun\ u\ z\ ↦\ fderiv\ ℝ\ (fun\ a\ ↦\ U\ i\ a\ z)\ (parameter\ u)\ (parameterDerivative\ u))\ t}
\leanline{2}{parameter\_equation\ :\ ∀\ t,\ 0\ ≤\ t\ →\ parameterDerivative\ t\ =}
\leanline{4}{-averagedParameterForce\ volume\ (fun\ i\ ↦\ density\ i\ t)\ U\ (parameter\ t)}
\leanline{2}{density\_equation\ :\ ∀\ t,\ 0\ ≤\ t\ →\ ∀\ i,\ ∀ᵐ\ z\ ∂volume,\ densityDerivative\ i\ t\ z\ =}
\leanline{4}{fluxDivergence\ (fokkerPlanckFlux\ (density\ i\ t)\ (U\ i\ (parameter\ t)))\ z}
\leanline{2}{integrable\_tails\ :\ ∀\ t\ i,\ NoBoundaryFlux}
\leanline{4}{(fokkerPlanckFlux\ (density\ i\ t)\ (U\ i\ (parameter\ t)))}
\leanline{4}{(fun\ z\ ↦\ chemicalPotential\ (density\ i\ t)\ (U\ i\ (parameter\ t))\ z\ +\ 1)}
\end{leancode}
\leansource{\texttt{energy} \quad LatentDiffusion/\allowbreak{}RegularConvergence.lean, lines 208–209}
\begin{leancode}
\leanline{0}{def\ energy\ (t\ :\ ℝ)\ :\ ℝ\ :=}
\leanline{2}{normalizedDensityEnergy\ volume\ (fun\ i\ ↦\ S.density\ i\ t)\ U\ (S.parameter\ t)}
\end{leancode}
\leansource{\texttt{fisher} \quad LatentDiffusion/\allowbreak{}RegularConvergence.lean, lines 211–212}
\begin{leancode}
\leanline{0}{def\ fisher\ (t\ :\ ℝ)\ :\ ℝ\ :=}
\leanline{2}{normalizedSlopeSquared\ volume\ (fun\ i\ ↦\ S.density\ i\ t)\ U\ (S.parameter\ t)}
\end{leancode}
\leansource{\texttt{targetEnergy} \quad LatentDiffusion/\allowbreak{}RegularConvergence.lean, lines 214–214}
\begin{leancode}
\leanline{0}{def\ targetEnergy\ :\ ℝ\ :=\ normalizedDensityEnergy\ volume\ S.targetDensity\ U\ astar}
\end{leancode}
\leansource{\texttt{evidence} \quad LatentDiffusion/\allowbreak{}RegularConvergence.lean, lines 216–217}
\begin{leancode}
\leanline{0}{def\ evidence\ (\_S\ :\ ClassicalFlow\ M\ d\ rho\ astar\ U)\ (a\ :\ A)\ :\ ℝ\ :=\ (M\ :\ ℝ)⁻¹\ *\ ∑\ i,}
\leanline{2}{Real.log\ (∫\ z,\ Real.exp\ (-U\ i\ a\ z))}
\end{leancode}
\leanlabel{How the Lean proof follows the argument}
\texttt{Classical\allowbreak{}Flow M d rho astar U} is a structure. Its fields are \texttt{parameter}, \texttt{parameter\allowbreak{}Derivative}, \texttt{density}, \texttt{density\allowbreak{}Derivative}, \texttt{target\allowbreak{}Density}, \texttt{transport} (a \texttt{Dataset\allowbreak{}Regularity} for every time) and the conditions (1)–(7): \texttt{potential\_\allowbreak{}eq}, \texttt{density\_\allowbreak{}eq}, \texttt{target\_\allowbreak{}density\_\allowbreak{}eq}; \texttt{parameter\_\allowbreak{}differentiable}, \texttt{potential\_\allowbreak{}parameter\_\allowbreak{}differentiable}; \texttt{positive}, \texttt{density\_\allowbreak{}C1}, \texttt{potential\_\allowbreak{}C1} (with \texttt{Cont\allowbreak{}Diff \ensuremath{\mathbb R} 1} for continuous differentiability); \texttt{parameter\_\allowbreak{}gradient} and \texttt{latent\_\allowbreak{}gradient} (\texttt{Has\allowbreak{}Gradient\allowbreak{}At} statements); \texttt{parameter\_\allowbreak{}gradient\_\allowbreak{}integrable}, \texttt{domination}; \texttt{parameter\_\allowbreak{}equation}, \texttt{density\_\allowbreak{}equation}; \texttt{integrable\_\allowbreak{}tails}. The definitions \texttt{energy}, \texttt{fisher}, \texttt{target\allowbreak{}Energy} and \texttt{evidence} use \texttt{normalized\allowbreak{}Density\allowbreak{}Energy} and \texttt{normalized\allowbreak{}Slope\allowbreak{}Squared} from the flow-calculus section.
\end{leanbox}
\subsection{\texorpdfstring{Identification of energies and Fisher information}{Identification of energies and Fisher information}}


For a classical flow and a time \(t\), let \(\mathcal F^{(t)}\), \(G^{(t)}\), \(S^{(t)}_m\), \(I^{(t)}\) and \(d_t^2\) denote the averaged free energy, parameter force, latent scores, Fisher information and squared distance under the extra transport assumptions at time \(t\) (previous section).




\textbf{Claim.} For every \(t\):



\begin{enumerate}


\item \(\mathcal F^{(t)}(0)=\mathcal E(t)\) and \(\mathcal F^{(t)}(1)=\mathcal E_\star\);

\item the endpoint law of datum \(m\) is \(q_\star^m(z)\,dz\); if \(q_\star^m=e^{-U_m(a_\star,\cdot)}/\int e^{-U_m(a_\star,w)}dw\), it is the posterior \(\pi_{a_\star}^m\), and then \(d_t^2=\|a_t-a_\star\|^2+M^{-1}\sum_mW_2(q_t^m\,dz,\pi^m_{a_\star})^2\);

\item \(G^{(t)}=G(a_t,q_t)\), \(S^{(t)}_m=\nabla(\log q_t^m+U_m(a_t,\cdot))\), and \(I^{(t)}=I(t)\).

\end{enumerate}



\textbf{Proof.} (1) At \(s=0\) the law of datum \(m\) is \(q_t^m\,dz\) and the parameter is \(a_t\). Hence \(\mathcal F^{(t)}_m(0)=\int\log q_t^m\,dQ_m+\int U_m(a_t,\cdot)\,dQ_m\); both functions are integrable by the extra transport assumptions, so the sum of the integrals is the integral of the sum, and by Claim 1 of the first step it equals \(\int q_t^m(\log q_t^m+U_m(a_t,\cdot))\,dz\). Averaging gives \(\mathcal E(t)\). The same computation at \(s=1\), with the law \(q_\star^m\,dz\) and the parameter \(a_\star\), gives \(\mathcal E_\star\).




(2) The first statement is the law condition at \(s=1\); the second substitutes the formula for \(q_\star^m\); the third inserts it into the definition of the distance.




(3) A gradient is unique: the gradient field of the data has, by condition (4) of the definition, the actual gradients as components. Hence \(\int G_{m,P}\,dQ_m=\int q_t^m\nabla_aU_m(a_t,\cdot)\,dz\) by Claim 1, and averaging gives \(G(a_t,q_t)\). Since \(q_t^m>0\) is continuously differentiable, \(\log q_t^m\) is differentiable and \(\nabla(\log q_t^m+U_m(a_t,\cdot))=\nabla\log q_t^m+\nabla_zU_m(a_t,\cdot)=S^{(t)}_m\). Inserting both identities into the definition of \(I^{(t)}\) and converting the \(Q_m\)-integrals by Claim 1 gives \(I(t)\). \(\square\)

\begin{leanbox}{Lean: \texttt{initial\_\allowbreak{}energy}, \texttt{terminal\_\allowbreak{}energy}, \texttt{terminal\_\allowbreak{}law} and 5 more}
\leansource{\texttt{initial\_\allowbreak{}energy} \quad LatentDiffusion/\allowbreak{}RegularConvergence.lean, lines 219–240}
\begin{leancode}
\leanline{0}{theorem\ initial\_energy\ (t\ :\ ℝ)\ :\ (S.transport\ t).freeEnergy\ 0\ =\ S.energy\ t}
\end{leancode}
\leansource{\texttt{terminal\_\allowbreak{}energy} \quad LatentDiffusion/\allowbreak{}RegularConvergence.lean, lines 242–273}
\begin{leancode}
\leanline{0}{theorem\ terminal\_energy\ (t\ :\ ℝ)\ :\ (S.transport\ t).freeEnergy\ 1\ =\ S.targetEnergy}
\end{leancode}
\leansource{\texttt{terminal\_\allowbreak{}law} \quad LatentDiffusion/\allowbreak{}RegularConvergence.lean, lines 275–279}
\begin{leancode}
\leanline{0}{theorem\ terminal\_law\ (t\ :\ ℝ)\ (i\ :\ Fin\ M)\ :}
\leanline{4}{((S.transport\ t).displacement\ i).law\ 1\ =}
\leanline{6}{volume.withDensity\ (fun\ z\ ↦\ ENNReal.ofReal\ (S.targetDensity\ i\ z))}
\end{leancode}
\leansource{\texttt{terminal\_\allowbreak{}posterior} \quad LatentDiffusion/\allowbreak{}RegularConvergence.lean, lines 283–290}
\begin{leancode}
\leanline{0}{theorem\ terminal\_posterior}
\leanline{4}{(hpost\ :\ ∀\ i\ z,\ S.targetDensity\ i\ z\ =}
\leanline{6}{Real.exp\ (-U\ i\ astar\ z)\ /\ (∫\ w,\ Real.exp\ (-U\ i\ astar\ w)))\ (t\ :\ ℝ)\ (i\ :\ Fin\ M)\ :}
\leanline{4}{((S.transport\ t).displacement\ i).law\ 1\ =}
\leanline{6}{Density.posteriorMeasure\ volume\ (fun\ z\ ↦\ Real.exp\ (-U\ i\ astar\ z))}
\leanline{8}{(∫\ w,\ Real.exp\ (-U\ i\ astar\ w))}
\end{leancode}
\leansource{\texttt{distanceSq\_\allowbreak{}eq\_\allowbreak{}posterior} \quad LatentDiffusion/\allowbreak{}RegularConvergence.lean, lines 292–299}
\begin{leancode}
\leanline{0}{theorem\ distanceSq\_eq\_posterior}
\leanline{4}{(hpost\ :\ ∀\ i\ z,\ S.targetDensity\ i\ z\ =}
\leanline{6}{Real.exp\ (-U\ i\ astar\ z)\ /\ (∫\ w,\ Real.exp\ (-U\ i\ astar\ w)))\ (t\ :\ ℝ)\ :}
\leanline{4}{(S.transport\ t).distanceSq\ =\ ‖S.parameter\ t-astar‖\ \textasciicircum{}\ 2\ +\ (M\ :\ ℝ)⁻¹\ *\ ∑\ i,}
\leanline{6}{ULA.wasserstein2Real\ ((S.transport\ t).law\ i\ :\ Measure\ Z)}
\leanline{8}{(Density.posteriorMeasure\ volume\ (fun\ z\ ↦\ Real.exp\ (-U\ i\ astar\ z))}
\leanline{10}{(∫\ w,\ Real.exp\ (-U\ i\ astar\ w)))\ \textasciicircum{}\ 2}
\end{leancode}
\leansource{\texttt{parameterForce\_\allowbreak{}eq} \quad LatentDiffusion/\allowbreak{}RegularConvergence.lean, lines 301–315}
\begin{leancode}
\leanline{0}{theorem\ parameterForce\_eq\ (t\ :\ ℝ)\ :\ (S.transport\ t).parameterForce\ =}
\leanline{4}{averagedParameterForce\ volume\ (fun\ i\ ↦\ S.density\ i\ t)\ U\ (S.parameter\ t)}
\end{leancode}
\leansource{\texttt{latentScore\_\allowbreak{}eq} \quad LatentDiffusion/\allowbreak{}RegularConvergence.lean, lines 317–332}
\begin{leancode}
\leanline{0}{theorem\ latentScore\_eq\ (t\ :\ ℝ)\ (i\ :\ Fin\ M)\ (z\ :\ Z)\ :}
\leanline{4}{(S.transport\ t).latentScore\ i\ z\ =}
\leanline{6}{gradient\ (chemicalPotential\ (S.density\ i\ t)\ (U\ i\ (S.parameter\ t)))\ z}
\end{leancode}
\leansource{\texttt{fisher\_\allowbreak{}eq} \quad LatentDiffusion/\allowbreak{}RegularConvergence.lean, lines 334–345}
\begin{leancode}
\leanline{0}{theorem\ fisher\_eq\ (t\ :\ ℝ)\ :\ (S.transport\ t).fisherInformation\ =\ S.fisher\ t}
\end{leancode}
\leanlabel{How the Lean proof follows the argument}
\begin{enumerate}


\item \texttt{initial\_\allowbreak{}energy} and \texttt{terminal\_\allowbreak{}energy} unfold \texttt{Dataset\allowbreak{}Regularity.\allowbreak{}free\allowbreak{}Energy} and \texttt{normalized\allowbreak{}Density\allowbreak{}Energy}, compare the summands, combine the two integrals by \texttt{integral\_\allowbreak{}add} and convert them with \texttt{integral\_\allowbreak{}density}; at the endpoint the integrability of \(\log q_\star^m\) and \(U_m(a_\star,\cdot)\) is transferred through the pushforward with \texttt{integrable\_\allowbreak{}map\_\allowbreak{}measure}.

\item \texttt{terminal\_\allowbreak{}law} is the field \texttt{density\_\allowbreak{}law} at \(s=1\); \texttt{terminal\_\allowbreak{}posterior} substitutes \texttt{hpost} into \texttt{Density.\allowbreak{}posterior\allowbreak{}Measure}; \texttt{distance\allowbreak{}Sq\_\allowbreak{}eq\_\allowbreak{}posterior} rewrites \texttt{Dataset\allowbreak{}Regularity.\allowbreak{}distance\allowbreak{}Sq} with it.

\item \texttt{parameter\allowbreak{}Force\_\allowbreak{}eq} uses \texttt{Has\allowbreak{}Gradient\allowbreak{}At.\allowbreak{}unique} to replace the parameter component, then \texttt{integral\_\allowbreak{}density}. \texttt{latent\allowbreak{}Score\_\allowbreak{}eq} computes the gradient of the sum by \texttt{hder}, the derivative of \(\log q+U\), and \texttt{map\_\allowbreak{}add} for the Riesz isomorphism \texttt{Inner\allowbreak{}Product\allowbreak{}Space.\allowbreak{}to\allowbreak{}Dual}, which turns derivatives into gradients. \texttt{fisher\_\allowbreak{}eq} combines the two.

\end{enumerate}
\end{leanbox}
\subsection{\texorpdfstring{Dissipation and extended log-Sobolev inequality along the flow}{Dissipation and extended log-Sobolev inequality along the flow}}


\textbf{Claim.} (1) For \(t\ge0\), the energy is differentiable at \(t\) and \(\mathcal E'(t)=-I(t)\). (2) If \(M\ge1\) and \(\lambda>0\), then for every \(t\), \(2\lambda\bigl(\mathcal E(t)-\mathcal F^{(t)}(1)\bigr)\le I(t)\).




\textbf{Proof.} (1) The de Bruijn identity of the flow-calculus section requires local domination, differentiability of the parameter path and of the potentials in the parameter, positivity and continuous differentiability of \(q_t^m\) and \(U_m(a_t,\cdot)\), integrability of \(q_t^m\nabla_aU_m(a_t,\cdot)\), the parameter equation, the Fokker–Planck equations and the no-boundary-flux conditions. These are conditions (2), (3), (5), (6) and (7) of a classical flow; the equations hold because \(t\ge0\). Its conclusion is \(\mathcal E'(t)=-I(t)\).




(2) The extended log-Sobolev inequality of the previous section, under the extra transport assumptions at time \(t\), gives \(2\lambda(\mathcal F^{(t)}(0)-\mathcal F^{(t)}(1))\le I^{(t)}\). By the previous step, \(\mathcal F^{(t)}(0)=\mathcal E(t)\) and \(I^{(t)}=I(t)\). \(\square\)

\begin{leanbox}{Lean: \texttt{deBruijn}, \texttt{extendedLSI}}
\leansource{\texttt{deBruijn} \quad LatentDiffusion/\allowbreak{}RegularConvergence.lean, lines 349–355}
\begin{leancode}
\leanline{0}{theorem\ deBruijn\ (t\ :\ ℝ)\ (ht\ :\ 0\ ≤\ t)\ :\ HasDerivAt\ S.energy\ (-S.fisher\ t)\ t}
\end{leancode}
\leansource{\texttt{extendedLSI} \quad LatentDiffusion/\allowbreak{}RegularConvergence.lean, lines 359–362}
\begin{leancode}
\leanline{0}{theorem\ extendedLSI\ (hM\ :\ 0\ <\ M)\ (hrho\ :\ 0\ <\ rho)\ (t\ :\ ℝ)\ :}
\leanline{4}{2*rho*(S.energy\ t-(S.transport\ t).freeEnergy\ 1)\ ≤\ S.fisher\ t}
\end{leancode}
\leanlabel{How the Lean proof follows the argument}
\begin{enumerate}


\item \texttt{de\allowbreak{}Bruijn} is the second component of \texttt{normalized\_\allowbreak{}product\_\allowbreak{}de\allowbreak{}Bruijn}, applied to the fields of the flow at time \texttt{t}; the hypothesis \texttt{ht : 0 \ensuremath{\le} t} is passed to \texttt{parameter\_\allowbreak{}equation} and \texttt{density\_\allowbreak{}equation}.

\item \texttt{extended\allowbreak{}LSI} rewrites \texttt{(S.\allowbreak{}transport t).\allowbreak{}extended\allowbreak{}LSI h\allowbreak{}M hrho} with \texttt{S.\allowbreak{}initial\_\allowbreak{}energy} and \texttt{S.\allowbreak{}fisher\_\allowbreak{}eq}.

\end{enumerate}
\end{leanbox}
\subsection{\texorpdfstring{Exponential decay of the energy gap}{Exponential decay of the energy gap}}


\textbf{Claim.} Let \(M\ge1\), \(\lambda>0\) and \(T\ge0\). For every \(t\in[0,T]\),


\[\mathcal E(t)-\mathcal E_\star\le\bigl(\mathcal E(0)-\mathcal E_\star\bigr)e^{-2\lambda t}.\]


\textbf{Proof.} Put \(E(u)=\mathcal E(u)-\mathcal E_\star\). By the previous step, \(E\) is differentiable at every \(u\ge0\) with \(E'(u)=-I(u)\); in particular \(E\) is continuous on \([0,T]\). By the extended log-Sobolev inequality and \(\mathcal F^{(u)}(1)=\mathcal E_\star\),


\[E'(u)=-I(u)\le-2\lambda E(u)\qquad(u\ge0).\]


We use Grönwall's inequality in the following form: if \(f\) is continuous on \([0,T]\), has at every \(u\in[0,T)\) a right derivative \(f'(u)\) with \(f'(u)\le Kf(u)\), and \(f(0)\le\delta\), then \(f(t)\le\delta e^{Kt}\) for \(t\in[0,T]\). (One proof compares \(f\) with \(\delta e^{K t}+\varepsilon'(e^{Kt}-1)/K\) for \(\varepsilon'>0\): at a first time where \(f\) would exceed it, the right derivatives would contradict the strict inequality; then \(\varepsilon'\downarrow0\).) With \(f=E\), \(K=-2\lambda\) and \(\delta=E(0)\) we obtain the claim. \(\square\)

\begin{leanbox}{Lean: \texttt{energy\_\allowbreak{}decay}}
\leansource{\texttt{energy\_\allowbreak{}decay} \quad LatentDiffusion/\allowbreak{}RegularConvergence.lean, lines 368–381}
\begin{leancode}
\leanline{0}{theorem\ energy\_decay\ (hM\ :\ 0\ <\ M)\ (hrho\ :\ 0\ <\ rho)\ (T\ :\ ℝ)}
\leanline{4}{(t\ :\ ℝ)\ (ht\ :\ t\ ∈\ Icc\ 0\ T)\ :}
\leanline{4}{S.energy\ t-S.targetEnergy\ ≤\ (S.energy\ 0-S.targetEnergy)*Real.exp\ (-2*rho*t)}
\end{leancode}
\leanlabel{How the Lean proof follows the argument}
\texttt{energy\_\allowbreak{}decay}: \texttt{hd} is the derivative of the gap, \texttt{(S.\allowbreak{}de\allowbreak{}Bruijn u hu).\allowbreak{}sub\_\allowbreak{}const S.\allowbreak{}target\allowbreak{}Energy}, and \texttt{hc} its continuity on \texttt{Icc 0 T}. Mathlib's \texttt{le\_\allowbreak{}gronwall\allowbreak{}Bound\_\allowbreak{}of\_\allowbreak{}liminf\_\allowbreak{}deriv\_\allowbreak{}right\_\allowbreak{}le} is the Grönwall inequality with constants \texttt{K := -\allowbreak{}2*rho} and \texttt{\ensuremath{\varepsilon} := 0}; its derivative hypothesis is supplied in the weaker form of a bound on right difference quotients by \texttt{liminf\_\allowbreak{}right\_\allowbreak{}slope\_\allowbreak{}le}, and its growth hypothesis \(f'\le Kf\) by \texttt{S.\allowbreak{}extended\allowbreak{}LSI} rewritten with \texttt{S.\allowbreak{}terminal\_\allowbreak{}energy}. \texttt{gronwall\allowbreak{}Bound\_\allowbreak{}\ensuremath{\varepsilon}0} evaluates the Grönwall bound as \(\delta e^{Kt}\).
\end{leanbox}
\subsection{\texorpdfstring{The Talagrand bound for the posterior optimizer}{The Talagrand bound for the posterior optimizer}}


\textbf{Claim.} Assume that \(e^{-U_m(a,\cdot)}\) is integrable with positive integral for all \(m\) and \(a\), that \(a_\star\) maximizes the evidence \(\ell\), and that the comparison densities are the posteriors, \(q_\star^m=e^{-U_m(a_\star,\cdot)}/\int e^{-U_m(a_\star,w)}dw\). Then (1) for every \(t\), \(\mathcal F^{(t)}(1)\le\mathcal F^{(t)}(s)\) for all \(s\in[0,1]\); and (2) if moreover \(M\ge1\) and \(\lambda>0\), then for every \(t\),


\[\frac\lambda2\,d_t^2\le\mathcal E(t)-\mathcal E_\star.\]


\textbf{Proof.} (1) is part (4) of the second step, applied under the extra transport assumptions at time \(t\), whose endpoint parameter is \(a_\star\) and whose endpoint densities are the posteriors. (2) By the Talagrand inequality of the previous section, (1) implies \(\frac\lambda2d_t^2\le\mathcal F^{(t)}(0)-\mathcal F^{(t)}(1)\), and the identification step gives \(\mathcal F^{(t)}(0)=\mathcal E(t)\) and \(\mathcal F^{(t)}(1)=\mathcal E_\star\). \(\square\)

\begin{leanbox}{Lean: \texttt{target\_\allowbreak{}minimal}, \texttt{distance\_\allowbreak{}control}}
\leansource{\texttt{target\_\allowbreak{}minimal} \quad LatentDiffusion/\allowbreak{}RegularConvergence.lean, lines 386–397}
\begin{leancode}
\leanline{0}{theorem\ target\_minimal}
\leanline{4}{(hweight\ :\ ∀\ i\ a,\ Integrable\ (fun\ z\ ↦\ Real.exp\ (-U\ i\ a\ z))\ volume)}
\leanline{4}{(hZ\ :\ ∀\ i\ a,\ 0\ <\ ∫\ z,\ Real.exp\ (-U\ i\ a\ z))}
\leanline{4}{(hmax\ :\ ∀\ a,\ S.evidence\ a\ ≤\ S.evidence\ astar)}
\leanline{4}{(hpost\ :\ ∀\ i\ z,\ S.targetDensity\ i\ z\ =}
\leanline{6}{Real.exp\ (-U\ i\ astar\ z)/(∫\ w,\ Real.exp\ (-U\ i\ astar\ w)))\ (t\ :\ ℝ)\ :}
\leanline{4}{∀\ s\ ∈\ Icc\ (0\ :\ ℝ)\ 1,\ (S.transport\ t).freeEnergy\ 1\ ≤\ (S.transport\ t).freeEnergy\ s}
\end{leancode}
\leansource{\texttt{distance\_\allowbreak{}control} \quad LatentDiffusion/\allowbreak{}RegularConvergence.lean, lines 401–410}
\begin{leancode}
\leanline{0}{theorem\ distance\_control\ (hM\ :\ 0\ <\ M)\ (hrho\ :\ 0\ <\ rho)}
\leanline{4}{(hweight\ :\ ∀\ i\ a,\ Integrable\ (fun\ z\ ↦\ Real.exp\ (-U\ i\ a\ z))\ volume)}
\leanline{4}{(hZ\ :\ ∀\ i\ a,\ 0\ <\ ∫\ z,\ Real.exp\ (-U\ i\ a\ z))}
\leanline{4}{(hmax\ :\ ∀\ a,\ S.evidence\ a\ ≤\ S.evidence\ astar)}
\leanline{4}{(hpost\ :\ ∀\ i\ z,\ S.targetDensity\ i\ z\ =}
\leanline{6}{Real.exp\ (-U\ i\ astar\ z)/(∫\ w,\ Real.exp\ (-U\ i\ astar\ w)))\ (t\ :\ ℝ)\ :}
\leanline{4}{rho/2*(S.transport\ t).distanceSq\ ≤\ S.energy\ t-S.targetEnergy}
\end{leancode}
\leanlabel{How the Lean proof follows the argument}
\begin{enumerate}


\item \texttt{target\_\allowbreak{}minimal} applies \texttt{Regular\allowbreak{}Dataset.\allowbreak{}posterior\_\allowbreak{}minimal} to \texttt{S.\allowbreak{}transport t}, rewriting each hypothesis with \texttt{S.\allowbreak{}potential\_\allowbreak{}eq t} and \texttt{S.\allowbreak{}target\_\allowbreak{}density\_\allowbreak{}eq t}.

\item \texttt{distance\_\allowbreak{}control} applies \texttt{(S.\allowbreak{}transport t).\allowbreak{}talagrand\_\allowbreak{}final h\allowbreak{}M hrho.\allowbreak{}le} with the minimality of the first part and rewrites with \texttt{S.\allowbreak{}initial\_\allowbreak{}energy} and \texttt{S.\allowbreak{}terminal\_\allowbreak{}energy}.

\end{enumerate}
\end{leanbox}
\subsection{\texorpdfstring{Product-state and parameter convergence}{Product-state and parameter convergence}}


\textbf{Claim.} Under the hypotheses of the previous step, for every \(t\ge0\),


\[d_t\le\sqrt{\frac{2E(t)}\lambda},\qquad\sqrt{\frac{2E(t)}\lambda}\le\sqrt{\frac{2E(0)}\lambda}\,e^{-\lambda t},\qquad\|a_t-a_\star\|\le\sqrt{\frac{2E(0)}\lambda}\,e^{-\lambda t},\]


where \(E(t)=\mathcal E(t)-\mathcal E_\star\) and \(d_t=\sqrt{d_t^2}\).




\textbf{Proof.} By the previous step, \(E(t)\ge\frac\lambda2d_t^2\ge0\), and likewise \(E(0)\ge0\). Dividing the Talagrand bound by \(\lambda/2>0\) gives \(d_t^2\le2E(t)/\lambda\), and the square root is increasing on \([0,\infty)\). For the second inequality, the decay step with \(T=t\) gives \(E(t)\le E(0)e^{-2\lambda t}\); multiplying by \(2/\lambda>0\),


\[\frac{2E(t)}\lambda\le\frac{2E(0)}\lambda e^{-2\lambda t}=\Bigl(\sqrt{\frac{2E(0)}\lambda}\,e^{-\lambda t}\Bigr)^2,\]


and taking square roots of these nonnegative numbers gives the bound. Finally \(\|a_t-a_\star\|^2\le d_t^2\), because \(d_t^2\) is \(\|a_t-a_\star\|^2\) plus a nonnegative average of squared distances; combining with the first two inequalities gives the parameter bound. \(\square\)

\begin{leanbox}{Lean: \texttt{full\_\allowbreak{}product\_\allowbreak{}exponential\_\allowbreak{}convergence}, \texttt{parameter\_\allowbreak{}exponential\_\allowbreak{}convergence}}
\leansource{\texttt{full\_\allowbreak{}product\_\allowbreak{}exponential\_\allowbreak{}convergence} \quad LatentDiffusion/\allowbreak{}RegularConvergence.lean, lines 416–450}
\begin{leancode}
\leanline{0}{theorem\ full\_product\_exponential\_convergence\ (hM\ :\ 0\ <\ M)\ (hrho\ :\ 0\ <\ rho)}
\leanline{4}{(hweight\ :\ ∀\ i\ a,\ Integrable\ (fun\ z\ ↦\ Real.exp\ (-U\ i\ a\ z))\ volume)}
\leanline{4}{(hZ\ :\ ∀\ i\ a,\ 0\ <\ ∫\ z,\ Real.exp\ (-U\ i\ a\ z))}
\leanline{4}{(hmax\ :\ ∀\ a,\ S.evidence\ a\ ≤\ S.evidence\ astar)}
\leanline{4}{(hpost\ :\ ∀\ i\ z,\ S.targetDensity\ i\ z\ =}
\leanline{6}{Real.exp\ (-U\ i\ astar\ z)/(∫\ w,\ Real.exp\ (-U\ i\ astar\ w)))}
\leanline{4}{(t\ :\ ℝ)\ (ht\ :\ 0\ ≤\ t)\ :}
\leanline{4}{Real.sqrt\ (S.transport\ t).distanceSq\ ≤\ Real.sqrt\ (2*(S.energy\ t-S.targetEnergy)/rho)\ ∧}
\leanline{4}{Real.sqrt\ (2*(S.energy\ t-S.targetEnergy)/rho)\ ≤}
\leanline{6}{Real.sqrt\ (2*(S.energy\ 0-S.targetEnergy)/rho)*Real.exp\ (-rho*t)}
\end{leancode}
\leansource{\texttt{parameter\_\allowbreak{}exponential\_\allowbreak{}convergence} \quad LatentDiffusion/\allowbreak{}RegularConvergence.lean, lines 454–469}
\begin{leancode}
\leanline{0}{theorem\ parameter\_exponential\_convergence\ (hM\ :\ 0\ <\ M)\ (hrho\ :\ 0\ <\ rho)}
\leanline{4}{(hweight\ :\ ∀\ i\ a,\ Integrable\ (fun\ z\ ↦\ Real.exp\ (-U\ i\ a\ z))\ volume)}
\leanline{4}{(hZ\ :\ ∀\ i\ a,\ 0\ <\ ∫\ z,\ Real.exp\ (-U\ i\ a\ z))}
\leanline{4}{(hmax\ :\ ∀\ a,\ S.evidence\ a\ ≤\ S.evidence\ astar)}
\leanline{4}{(hpost\ :\ ∀\ i\ z,\ S.targetDensity\ i\ z\ =}
\leanline{6}{Real.exp\ (-U\ i\ astar\ z)/(∫\ w,\ Real.exp\ (-U\ i\ astar\ w)))}
\leanline{4}{(t\ :\ ℝ)\ (ht\ :\ 0\ ≤\ t)\ :}
\leanline{4}{‖S.parameter\ t-astar‖\ ≤}
\leanline{6}{Real.sqrt\ (2*(S.energy\ 0-S.targetEnergy)/rho)*Real.exp\ (-rho*t)}
\end{leancode}
\leanlabel{How the Lean proof follows the argument}
\begin{enumerate}


\item \texttt{full\_\allowbreak{}product\_\allowbreak{}exponential\_\allowbreak{}convergence}: \texttt{hctrl} is the Talagrand bound, \texttt{hgap} and \texttt{hgap0} the nonnegativity of the gaps (from \texttt{distance\allowbreak{}Sq\_\allowbreak{}nonneg}). The first inequality uses \texttt{Real.\allowbreak{}sqrt\_\allowbreak{}le\_\allowbreak{}sqrt} and \texttt{le\_\allowbreak{}div\_\allowbreak{}iff\ensuremath{{}_{0}}}; the second uses \texttt{S.\allowbreak{}energy\_\allowbreak{}decay}, the identity \texttt{hsquare} for the square of the right side, and \texttt{nlinarith} with \texttt{Real.\allowbreak{}sq\_\allowbreak{}sqrt}.

\item \texttt{parameter\_\allowbreak{}exponential\_\allowbreak{}convergence}: \texttt{hp} is \(\|a_t-a_\star\|\le d_t\), from \texttt{Real.\allowbreak{}le\_\allowbreak{}sqrt} and \texttt{le\_\allowbreak{}add\_\allowbreak{}of\_\allowbreak{}nonneg\_\allowbreak{}right}; it is chained with both parts of the previous theorem.

\end{enumerate}
\end{leanbox}
\subsection{\texorpdfstring{Unique optimizer and exponential convergence of the flow}{Unique optimizer and exponential convergence of the flow}}


\textbf{Claim (Theorem 1, corrected).} Let \(M\ge1\), \(\lambda>0\), and let \(U_1,\dots,U_M:\mathbb R^{d_\theta}\times\mathbb R^{d_z}\to\mathbb R\) be differentiable and \(\lambda\)-strongly convex (A3). Let \(\ell(a)=M^{-1}\sum_m\log\int e^{-U_m(a,z)}dz\). Then there is exactly one \(a_\star\in\mathbb R^{d_\theta}\) such that



\begin{enumerate}


\item \(\ell(a)\le\ell(a_\star)\) for every \(a\), and

\item for every classical flow with comparison parameter \(a_\star\), constant \(\lambda\) and potentials \(U_m\), whose comparison densities are the posteriors \(q_\star^m=e^{-U_m(a_\star,\cdot)}/\int e^{-U_m(a_\star,w)}dw\), and for every \(t\ge0\),
\[\|a_t-a_\star\|\le\sqrt{\frac{2(\mathcal E(0)-\mathcal E_\star)}\lambda}\,e^{-\lambda t}.\]

\end{enumerate}



\textbf{Proof.} By the density-model section, \(\ell=-V\), where \(V\) is the average marginal potential. By the marginal section, \(V\) has exactly one minimizer \(a_\star\); hence \(a_\star\) is the unique maximizer of \(\ell\), which proves (1). For (2), A3 gives integrability of \(e^{-U_m(a,\cdot)}\) and positivity of its integral for all \(a\) (marginal section), so the hypotheses of the previous steps hold, and the parameter bound of the previous step is (2). Uniqueness: if \(b\) satisfies (1) and (2), then in particular \(\ell(a)\le\ell(b)\) for all \(a\), so \(V(b)\le V(a)\) for all \(a\), and \(b=a_\star\) by uniqueness of the minimizer of \(V\). \(\square\)




By the second step, \(\mathcal E_\star=-\ell(a_\star)=F_\star\) is the minimal free energy, so the prefactor is \(\sqrt{2(F(a_0,q_0)-F_\star)/\lambda}\) as in the paper. The theorem contains no assumed dissipation identity, log-Sobolev inequality or transport inequality; the only hypotheses beyond A3 are the extra assumptions defining a classical flow.

\begin{leanbox}{Lean: \texttt{existsUnique\_\allowbreak{}optimizer\_\allowbreak{}and\_\allowbreak{}flow\_\allowbreak{}convergence}}
\leansource{\texttt{existsUnique\_\allowbreak{}optimizer\_\allowbreak{}and\_\allowbreak{}flow\_\allowbreak{}convergence} \quad LatentDiffusion/\allowbreak{}RegularConvergence.lean, lines 478–517}
\begin{leancode}
\leanline{0}{theorem\ existsUnique\_optimizer\_and\_flow\_convergence\ \{dθ\ dz\ :\ ℕ\}}
\leanline{4}{(U\ :\ Fin\ M\ →\ IPLA.Joint\ dθ\ dz\ →\ ℝ)\ (hM\ :\ 0\ <\ M)}
\leanline{4}{(rho\ :\ ℝ)\ (hrho\ :\ 0\ <\ rho)}
\leanline{4}{(hdiff\ :\ ∀\ i,\ Differentiable\ ℝ\ (U\ i))}
\leanline{4}{(hsc\ :\ ∀\ i,\ StrongConvexOn\ univ\ rho\ (U\ i))\ :}
\leanline{4}{let\ V\ :=\ fun\ i\ a\ z\ ↦\ U\ i\ (IPLA.pair\ a\ z)}
\leanline{4}{let\ ell\ :=\ fun\ a\ ↦\ (M\ :\ ℝ)⁻¹\ *\ ∑\ i,\ Real.log\ (∫\ z,\ Real.exp\ (-V\ i\ a\ z))}
\leanline{4}{∃!\ astar\ :\ EuclideanSpace\ ℝ\ (Fin\ dθ),}
\leanline{6}{(∀\ a,\ ell\ a\ ≤\ ell\ astar)\ ∧}
\leanline{6}{∀\ S\ :\ ClassicalFlow\ M\ dz\ rho\ astar\ V,}
\leanline{8}{(∀\ i\ z,\ S.targetDensity\ i\ z\ =\ Real.exp\ (-V\ i\ astar\ z)/}
\leanline{10}{(∫\ w,\ Real.exp\ (-V\ i\ astar\ w)))\ →}
\leanline{8}{∀\ t,\ 0\ ≤\ t\ →\ ‖S.parameter\ t-astar‖\ ≤}
\leanline{10}{Real.sqrt\ (2*(S.energy\ 0-S.targetEnergy)/rho)*Real.exp\ (-rho*t)}
\end{leancode}
\leanlabel{How the Lean proof follows the argument}
In \texttt{exists\allowbreak{}Unique\_\allowbreak{}optimizer\_\allowbreak{}and\_\allowbreak{}flow\_\allowbreak{}convergence}, the local definitions \texttt{V} and \texttt{ell} are the slice potentials \(U_m(a,z)\) (written through \texttt{IPLA.\allowbreak{}pair}) and the evidence. The conclusion \texttt{\ensuremath{\exists}! astar, (\ensuremath{\forall} a, ell a \ensuremath{\le} ell astar) ∧ \ensuremath{\forall} S : Classical\allowbreak{}Flow M dz rho astar V, …} is the claim.



\begin{enumerate}


\item \texttt{he} identifies \texttt{ell} with \texttt{-\allowbreak{}average\allowbreak{}Marginal\allowbreak{}Potential U} through \texttt{Regular\allowbreak{}Model.\allowbreak{}density\allowbreak{}Model\_\allowbreak{}log\allowbreak{}Evidence}.

\item \texttt{exists\allowbreak{}Unique\_\allowbreak{}average\allowbreak{}Marginal\allowbreak{}Minimizer} supplies \texttt{astar} and its uniqueness \texttt{hu}; \texttt{hmax} is part (1).

\item Part (2) is \texttt{S.\allowbreak{}parameter\_\allowbreak{}exponential\_\allowbreak{}convergence}, with the integrability and positivity inputs \texttt{IPLA.\allowbreak{}integrable\_\allowbreak{}marginal\_\allowbreak{}weight} and \texttt{IPLA.\allowbreak{}marginal\allowbreak{}Integral\_\allowbreak{}pos}.

\item Uniqueness applies \texttt{hu} to a competitor \texttt{b}, using only its first property, rewritten with \texttt{he} and closed by \texttt{linarith}.

\end{enumerate}
\end{leanbox}
\section{\texorpdfstring{A full-batch root-mean-square rate to the actual optimizer}{A full-batch root-mean-square rate to the actual optimizer}}


Let \(K=MN>0\). Assume \(m>0\), differentiable jointly \(m\)-strongly convex heterogeneous potentials with \(L\)-Lipschitz gradients, and a stable positive step. The optimizer \(a_\star\) is the actual marginal optimizer already constructed from these potentials.




Supply an \(L^2\) initial latent cloud with the stated frozen Gibbs law at \(a_\star\), the correct individual posterior marginals, the local domination assumptions for evidence differentiation, and independent initial posterior scores. The coordinate Brownian processes have independent scalar Euler increments, independent of the initial latent cloud. Suppose each centred score has second moment at most \(V_0\). These extra initialization, independence and differentiation assumptions are part of the theorem.

\subsection{\texorpdfstring{Identify the optimizer, then insert explicit variance and Brownian constants}{Identify the optimizer, then insert explicit variance and Brownian constants}}


\textbf{Claim.} For the actual pointwise Euler recursion,

\[\left(\mathbb E\|a_n-a_\star\|^2\right)^{1/2}\le q^nR_0+\frac2m\sqrt{\frac{V_0}{MN}}+\frac Lm\left(h\sqrt{L d_z}+\frac43\sqrt{2d_z}\sqrt h\right),\qquad q=1-mh/2,\]

where \(R_0\) is the \(L^2\) distance from the initial cloud to the scaled posterior cloud at \(a_\star\), and \(d_z\) is the dimension of one latent vector.



\textbf{Proof.} The evidence derivative at the optimizer centres the sum of posterior scores. The stipulated Gibbs initialization and independent Brownian increments furnish the stationary frozen reference used by the preceding stochastic comparison. Its frozen drift moment is bounded by \(\sqrt{Ld_z}\), and the scaled physical noise has increment norm \(\sqrt{2d_z}\sqrt t\). These facts give the three-term comparison with sampling residual \(K^{-1}\sqrt{\sum_i\operatorname{Var}(s_i)}\). The per-score bound gives \(\sum_i\operatorname{Var}(s_i)\le KV_0\), hence

\[\frac1K\sqrt{\sum_i\operatorname{Var}(s_i)}\le\frac{\sqrt{KV_0}}K=\sqrt{\frac{V_0}K}.\]

All factors are nonnegative and \(K>0\), so taking square roots and dividing preserves the inequality. Substitute \(K=MN\) and the explicit moment constants. The pointwise and lifted recursions agree in \(L^2\), which converts the projected norm to the displayed integral. No assertion of uniform complexity follows unless the constants stay bounded independently of particle count and step size. \(\square\)

\begin{leanbox}{Lean: \texttt{full\_\allowbreak{}batch\_\allowbreak{}optimizer\_\allowbreak{}rms\_\allowbreak{}rate}}
\leansource{\texttt{full\_\allowbreak{}batch\_\allowbreak{}optimizer\_\allowbreak{}rms\_\allowbreak{}rate} \quad LatentDiffusion/\allowbreak{}RegularSourceTheorem.lean, lines 167–209}
\begin{leancode}
\leanline{0}{theorem\ full\_batch\_optimizer\_rms\_rate}
\leanline{4}{\{M\ N\ :\ ℕ\}\ (hMN\ :\ 0\ <\ M*N)}
\leanline{4}{(U\ :\ Fin\ (M*N)\ →\ Cloud.Joint\ dθ\ dz\ →\ ℝ)}
\leanline{4}{(m\ h\ :\ ℝ)\ (L\ :\ ℝ≥0)\ (hm\ :\ 0\ <\ m)}
\leanline{4}{(hdiff\ :\ ∀\ i,\ Differentiable\ ℝ\ (U\ i))}
\leanline{4}{(hsc\ :\ ∀\ i,\ StrongConvexOn\ Set.univ\ m\ (U\ i))}
\leanline{4}{(hLip\ :\ ∀\ i,\ LipschitzWith\ L\ (gradient\ (U\ i)))\ (hs\ :\ StableStep\ m\ L\ h)}
\leanline{4}{(initial\ :\ Lp\ (Latent\ (M*N)\ dz)\ 2\ P)}
\leanline{4}{(hinitial\ :\ P.map\ initial\ =\ ULA.gibbsMeasure}
\leanline{6}{(frozenPotential\ U\ (optimizerParameter\ hMN\ U\ m\ hm\ hdiff\ hsc)))}
\leanline{4}{(hposterior\ :\ ∀\ i,\ P.map\ (fun\ ω\ =>\ physicalParticleCLM\ i}
\leanline{8}{((packLp\ P\ (optimizerParameter\ hMN\ U\ m\ hm\ hdiff\ hsc)\ initial)\ ω))\ =}
\leanline{6}{RegularEvidence.posterior\ volume\ (latentPotential\ U\ i)\ (optimizerParameter\ hMN\ U\ m\ hm\ hdiff\ hsc))}
\leanline{4}{(data\ :\ ∀\ i,\ RegularEvidence.LocalNormalizerRegularity\ volume\ (latentPotential\ U\ i)}
\leanline{6}{(optimizerParameter\ hMN\ U\ m\ hm\ hdiff\ hsc))}
\leanline{4}{(hindScores\ :\ iIndepFun\ (fun\ i\ ω\ =>\ scoreLp\ P\ U\ L\ hLip}
\leanline{6}{(packLp\ P\ (optimizerParameter\ hMN\ U\ m\ hm\ hdiff\ hsc)\ initial)\ i\ ω)\ P)}
\leanline{4}{(B\ :\ (Fin\ (M*N)\ ×\ Fin\ dz)\ →\ ℝ≥0\ →\ Ω\ →\ ℝ)}
\leanline{4}{(hB\ :\ ∀\ i,\ IsPreBrownianReal\ (B\ i)\ P)}
\leanline{4}{(hscalar\ :\ ∀\ s\ :\ ℝ,\ 0\ <\ s\ →\ iIndepFun\ (fun\ p\ :\ ℕ\ ×\ (Fin\ (M*N)\ ×\ Fin\ dz)\ =>\ fun\ ω\ =>}
\leanline{6}{ULA.Hilbert.brownianEulerNoise\ P\ B\ hB\ s\ p.1\ ω\ p.2)\ P)}
\leanline{4}{(hind\ :\ ∀\ s\ :\ ℝ,\ 0\ <\ s\ →\ IndepFun\ initial}
\leanline{6}{(fun\ ω\ n\ =>\ ULA.Hilbert.brownianEulerNoise\ P\ B\ hB\ s\ n\ ω)\ P)}
\leanline{4}{(V₀\ :\ ℝ)}
\leanline{4}{(hvar\ :\ ∀\ i,\ ‖scoreLp\ P\ U\ L\ hLip}
\leanline{8}{(packLp\ P\ (optimizerParameter\ hMN\ U\ m\ hm\ hdiff\ hsc)\ initial)\ i-}
\leanline{6}{constantL2\ P\ (∫\ ω,\ scoreLp\ P\ U\ L\ hLip}
\leanline{8}{(packLp\ P\ (optimizerParameter\ hMN\ U\ m\ hm\ hdiff\ hsc)\ initial)\ i\ ω\ ∂P)‖\textasciicircum{}2\ ≤\ V₀)}
\leanline{4}{(x₀\ :\ Lp\ (Cloud.State\ (M*N)\ dθ\ dz)\ 2\ P)\ (n\ :\ ℕ)\ :}
\leanline{4}{Real.sqrt\ (∫\ ω,\ ‖IPLA.parameterCLM}
\leanline{6}{(eulerTrajectory\ (Cloud.meanForce\ U)\ h}
\leanline{8}{(fun\ k\ =>\ noiseIncrement\ P\ B\ hB\ h\ k\ ω)\ (x₀\ ω)\ n)-}
\leanline{6}{optimizerParameter\ hMN\ U\ m\ hm\ hdiff\ hsc‖\textasciicircum{}2\ ∂P)\ ≤}
\leanline{6}{(1-m*h/2)\textasciicircum{}n\ *\ ‖x₀-packLp\ P\ (optimizerParameter\ hMN\ U\ m\ hm\ hdiff\ hsc)\ initial‖\ +}
\leanline{6}{(2/m)\ *\ Real.sqrt\ (V₀/((M*N\ :\ ℕ)\ :\ ℝ))\ +}
\leanline{6}{(L\ :\ ℝ)/m\ *\ (h*Real.sqrt\ ((L\ :\ ℝ)*(dz\ :\ ℝ))+}
\leanline{8}{(4/3\ :\ ℝ)*Real.sqrt\ (2*(dz\ :\ ℝ))*Real.sqrt\ h)}
\end{leancode}
\leanlabel{How the Lean proof follows the argument}
\texttt{full\_\allowbreak{}batch\_\allowbreak{}optimizer\_\allowbreak{}rms\_\allowbreak{}rate} obtains the actual root-mean-square comparison from \texttt{full\_\allowbreak{}batch\_\allowbreak{}optimizer\_\allowbreak{}rms}. It obtains the finite-particle square-root rate from \texttt{score\_\allowbreak{}sampling\_\allowbreak{}rate} and multiplies by the nonnegative factor \texttt{2/\allowbreak{}m} in \texttt{hv'}. The final \texttt{linarith} composes these estimates. The long hypothesis list explicitly identifies the posterior law, its differentiability data, the independent scores, Brownian increments, and their independence from the initial cloud; none is omitted by the short final arithmetic proof.
\end{leanbox}
\section{\texorpdfstring{Five explicit allocations of the requested error tolerance}{Five explicit allocations of the requested error tolerance}}


Let \(\varepsilon,m,K>0\), \(h,R_0,V,L,G,W\ge0\), and \(\sigma\in\mathbb R\). Suppose \(q=1-mh/2\ge0\) and an error \(E\) has the five-term upper bound \(q^nR_0+(2/m)\sqrt{V/K}+(L/m)hG+(4L/(3m))W\sqrt h+\sqrt{h\sigma^2/m}\). The following arithmetic allocates \(\varepsilon/5\) to each term. It applies after the stochastic/\allowbreak{}reference theorem supplies this upper bound; the arithmetic does not establish those hypotheses.

\subsection{\texorpdfstring{Logarithmic burn-in and four algebraic restrictions}{Logarithmic burn-in and four algebraic restrictions}}


\textbf{Claim.} The error satisfies \(E\le\varepsilon\) if

\[\begin{gathered}\log\max\{1,5R_0/\varepsilon\}\le mhn/2,\qquad 100V\le m^2\varepsilon^2K,\\5LhG\le m\varepsilon,\qquad400L^2W^2h\le9m^2\varepsilon^2,\qquad25h\sigma^2\le m\varepsilon^2.\end{gathered}\]

\textbf{Proof.} The inequality \(1-u\le e^{-u}\), with \(u=mh/2\), and \(q\ge0\) imply \(q^nR_0\le e^{-mhn/2}R_0\). The logarithmic restriction makes this at most \(\varepsilon/5\); taking a maximum with one also covers \(R_0=0\) and an initial error already below tolerance. For the particle term, square its desired nonnegative inequality:

\[\left(\frac2m\sqrt{\frac VK}\right)^2\le\frac{\varepsilon^2}{25}\quad\Longleftrightarrow\quad100V\le m^2\varepsilon^2K.\]

The drift inequality is equivalent to its desired bound after multiplication by \(5m>0\). Squaring the nonnegative Brownian term gives exactly \(400L^2W^2h\le9m^2\varepsilon^2\). Squaring the oracle term gives exactly \(25h\sigma^2\le m\varepsilon^2\). Nonnegativity permits the square-root comparisons in each case. Thus the sum of all five terms is at most \(5\varepsilon/5=\varepsilon\). The division-free restrictions also make sense when variance, drift or noise constants vanish. \(\square\)

\begin{leanbox}{Lean: \texttt{budget\_\allowbreak{}le\_\allowbreak{}epsilon}}
\leansource{\texttt{budget\_\allowbreak{}le\_\allowbreak{}epsilon} \quad LatentDiffusion/\allowbreak{}Accuracy.lean, lines 53–82}
\begin{leancode}
\leanline{0}{theorem\ budget\_le\_epsilon}
\leanline{4}{(m\ h\ R0\ V\ L\ G\ W\ sigma\ epsilon\ K\ E\ :\ ℝ)\ (n\ :\ ℕ)}
\leanline{4}{(hm\ :\ 0\ <\ m)\ (hh\ :\ 0\ ≤\ h)\ (hR\ :\ 0\ ≤\ R0)\ (hV\ :\ 0\ ≤\ V)}
\leanline{4}{(hL\ :\ 0\ ≤\ L)\ (\_hG\ :\ 0\ ≤\ G)\ (hW\ :\ 0\ ≤\ W)\ (hK\ :\ 0\ <\ K)\ (he\ :\ 0\ <\ epsilon)}
\leanline{4}{(hq\ :\ 0\ ≤\ 1\ -\ m\ *\ h\ /\ 2)}
\leanline{4}{(hn\ :\ Real.log\ (max\ 1\ (5\ *\ R0\ /\ epsilon))\ ≤\ m\ *\ h\ *\ n\ /\ 2)}
\leanline{4}{(hparticles\ :\ 100\ *\ V\ ≤\ m\ \textasciicircum{}\ 2\ *\ epsilon\ \textasciicircum{}\ 2\ *\ K)}
\leanline{4}{(hdrift\ :\ 5\ *\ L\ *\ h\ *\ G\ ≤\ m\ *\ epsilon)}
\leanline{4}{(hnoise\ :\ 400\ *\ L\ \textasciicircum{}\ 2\ *\ W\ \textasciicircum{}\ 2\ *\ h\ ≤\ 9\ *\ m\ \textasciicircum{}\ 2\ *\ epsilon\ \textasciicircum{}\ 2)}
\leanline{4}{(horacle\ :\ 25\ *\ h\ *\ sigma\ \textasciicircum{}\ 2\ ≤\ m\ *\ epsilon\ \textasciicircum{}\ 2)}
\leanline{4}{(hbudget\ :\ E\ ≤\ (1\ -\ m\ *\ h\ /\ 2)\ \textasciicircum{}\ n\ *\ R0\ +\ (2\ /\ m)\ *\ Real.sqrt\ (V\ /\ K)\ +}
\leanline{6}{L\ /\ m\ *\ (h\ *\ G\ +\ (4\ /\ 3\ :\ ℝ)\ *\ W\ *\ Real.sqrt\ h)\ +}
\leanline{6}{Real.sqrt\ (h\ *\ sigma\ \textasciicircum{}\ 2\ /\ m))\ :\ E\ ≤\ epsilon}
\end{leancode}
\leanlabel{How the Lean proof follows the argument}
\texttt{budget\_\allowbreak{}le\_\allowbreak{}epsilon} names the five bounds \texttt{hb}, \texttt{hp}, \texttt{hd}, \texttt{hw} and \texttt{ho}. \texttt{burnin} proves the logarithmic estimate. \texttt{scaled\_\allowbreak{}sqrt\_\allowbreak{}le} handles the particle and Brownian terms by comparing squares; \texttt{Real.\allowbreak{}sqrt\_\allowbreak{}le\_\allowbreak{}left} handles the oracle term. Positive denominators are discharged by \texttt{positivity} and \texttt{field\_\allowbreak{}simp}. The last \texttt{nlinarith} adds the five allocations to \texttt{hbudget}.
\end{leanbox}

\end{document}
